\documentclass[11pt,letterpaper]{article}
\usepackage[margin=1.1in,headheight=14pt]{geometry}
\usepackage[T1]{fontenc}
\usepackage{amsmath,amssymb,amsthm,mathtools}
\usepackage{mathpazo}
\usepackage{microtype}
\usepackage{enumitem}
\usepackage{booktabs,longtable,array}
\usepackage[numbers,sort&compress]{natbib}
\usepackage[hidelinks,pdfencoding=auto]{hyperref}
\usepackage{fancyhdr}
\pagestyle{fancy}
\fancyhf{}
\fancyhead[L]{\small J. Adriazola}
\fancyhead[R]{\small The Alp\"oge proof of the Riemann hypothesis}
\fancyfoot[C]{\thepage}
\renewcommand{\headrulewidth}{0.3pt}

\theoremstyle{plain}
\newtheorem{theorem}{Theorem}[section]
\newtheorem{proposition}[theorem]{Proposition}
\newtheorem{lemma}[theorem]{Lemma}
\newtheorem{corollary}[theorem]{Corollary}
\newtheorem*{theoremA}{Theorem A}
\theoremstyle{definition}
\newtheorem{definition}[theorem]{Definition}
\newtheorem{hypothesis}[theorem]{Hypothesis}
\theoremstyle{remark}
\newtheorem{remark}[theorem]{Remark}
\numberwithin{equation}{section}

\DeclareMathOperator{\Res}{Res}
\DeclareMathOperator{\rk}{rk}
\DeclareMathOperator{\covol}{covol}
\DeclareMathOperator{\ph}{ph}
\newcommand{\degh}{\widehat{\deg}}
\DeclareMathOperator{\ord}{ord}
\DeclareMathOperator{\len}{len}
\DeclareMathOperator{\Hom}{Hom}
\DeclareMathOperator{\GL}{GL}
\DeclareMathOperator{\tr}{tr}
\DeclareMathOperator{\Tr}{Tr}
\DeclareMathOperator{\vol}{vol}
\DeclareMathOperator{\spn}{span}
\DeclareMathOperator{\hgt}{ht}
\DeclareMathOperator{\lcv}{lcv}
\DeclareMathOperator{\Syl}{Syl}
\DeclareMathOperator{\Gal}{Gal}
\DeclareMathOperator{\Cl}{Cl}
\DeclareMathOperator{\Art}{Art}
\DeclareMathOperator{\cond}{cond}
\newcommand{\Int}{\mathrm{Int}}
\newcommand{\HS}{\mathrm{HS}}
\newcommand{\OK}{\mathcal O_K}
\newcommand{\NP}{N\!P}
\newcommand{\cO}{\mathcal O}
\newcommand{\cK}{\mathcal K}
\newcommand{\cL}{\mathcal L}
\newcommand{\cX}{\mathcal X}
\newcommand{\cF}{\mathcal F}
\newcommand{\cV}{\mathcal V}
\newcommand{\cT}{\mathcal T}
\newcommand{\Wl}{\mathsf W}
\newcommand{\bK}{\mathbf K}
\newcommand{\bC}{\mathbb C}
\newcommand{\bF}{\mathbb F}
\newcommand{\bQ}{\mathbb Q}
\newcommand{\bZ}{\mathbb Z}
\newcommand{\A}{\mathbb A}
\newcommand{\Af}{\mathbb A_f}
\newcommand{\F}{\mathbb F}
\newcommand{\C}{\mathbb C}
\newcommand{\R}{\mathbb R}
\newcommand{\Q}{\mathbb Q}
\newcommand{\Z}{\mathbb Z}
\newcommand{\Ql}{\mathbb Q_l}
\newcommand{\Zl}{\mathbb Z_l}
\newcommand{\sq}{\mathrm{sq}}
\newcommand{\per}{\mathrm{per}}
\font\Shafont=wncyr10
\newcommand{\Sha}{\text{\Shafont SH}}


\setlength{\emergencystretch}{2em}
\setcounter{tocdepth}{2}
\makeatletter
\renewcommand*\l@subsection{\@dottedtocline{2}{1.5em}{2.9em}}
\makeatother
\allowdisplaybreaks[2]

\hypersetup{pdftitle={The Alpöge Proof of the Riemann Hypothesis},pdfauthor={Jimmie Adriazola}}
\title{\bfseries The Claude-Alp\"oge Proof of the Riemann Hypothesis}
\author{Jimmie Adriazola$^{1,2}$\thanks{\texttt{jimmie.adriazola@asu.edu}}\\[6pt]
\small $^1$School of Mathematical and Statistical Sciences, Arizona State University,\\[-2pt]
\small Tempe, AZ 85287, USA\\[3pt]
\small $^2$Simon A. Levin Mathematical, Computational and Modeling Sciences Center,\\[-2pt]
\small Arizona State University, Tempe, AZ 85287, USA}
\date{}

% Every step whose proof is not given is wrapped in \notgiven{...}.
\newcommand{\notgiven}[1]{#1}

\begin{document}
\maketitle

\begin{abstract}
On November 20, 2026, Levent Alp\"oge announced a proof of the Riemann hypothesis that he and Claude had found together. He released the proof in Lean. Shortly thereafter, the Lean and Mathlib communities certified the formal statement and proof.

Alp\"oge's argument is as follows. First, a quantitative form of Landau's theorem on Mellin transforms reduces the zero-free half-plane $\operatorname{Re}s>1/2$ for Dirichlet $L$-functions, including the Riemann zeta function, to a one-sided square-root bound for smoothed prime sums on suitable intervals. Alp\"oge begins with this reduction and then adapts Stepanov's point-counting method over finite fields. In this arithmetic version of Stepanov's method, the increasing divisibility of successive theta lifts at inert primes replaces high-order vanishing at rational points. Each unit of $p$-adic divisibility contributes $\log p$ to the gain term in the slope inequality. For a sum of length $X$, a chain of about $\sqrt X$ theta lifts balances the errors in that inequality and gives the one-sided square-root bound that Landau's criterion requires.

This manuscript presents a human-readable reconstruction of Alp\"oge's argument. It gives the argument in conventional mathematical language and identifies the mathematical inputs and manuscript-to-Lean correspondences that still require verification.
\end{abstract}


\noindent\textbf{2020 Mathematics Subject Classification.} 11M26, 11M06, 11G05, 11F27, 11F70, 11F80, 14G40.

\tableofcontents

\section{Introduction}\label{sec:intro}

This paper reconstructs the argument for the following theorem.

\begin{theorem}\label{thm:intro}
For every primitive Dirichlet character $\chi$, the meromorphic function $L(s,\chi)$ has no zeros in the half-plane $\operatorname{Re}s>\frac12$. In particular, the Riemann zeta function has no zeros in $\operatorname{Re}s>\frac12$.
\end{theorem}

The proof is due to Levent Alpöge and Claude. On November 20, 2026, Alpöge released a four-page note~\cite{note} together with a proof of Theorem~\ref{thm:intro} formalized in Lean~\cite{lean4} on top of Mathlib~\cite{mathlib}, and on November 21, 2026, the Lean and Mathlib communities, with the endorsement of the Lean FRO and the Mathlib Initiative, certified that the formal statement and the formal proof are correct~\cite{leanstatement}. The note describes the argument as a descendant of Stepanov's method~\cite{stepanov} for curves over finite fields. It credits Alpöge with the plan and with the recognition that an auxiliary-function argument of Siegel, which appeared in mathematical work by researchers at OpenAI, is Stepanov's method with multiplicity one. Alpöge describes the turning point as the appearance of multiplicity two in an automorphic determinant on $\mathrm O(2,n)$, which suggested that further theta lifts could increase the multiplicity. The note credits Claude with the proposal that a chain of theta lifts of length about $\sqrt x$ should play the role of multiplicity, and with the arguments on automorphic forms on groups of higher rank. It thanks David Dohan, Ben Michel, Ralph Furman, Eric Price and Cora Vinther. This paper writes the proof as conventional mathematics that a reader can check without the formal development. The reconstruction is not complete. The paper proves Theorem~\ref{thm:intro} from the published theorems it cites and from a finite list of internal statements, each stated precisely where it is used and summarized in Subsection~\ref{sec:frontier}, whose conventional proofs it does not yet give. The formal development proves each of these statements. A few of the cited published theorems have not yet been compared line by line with their sources, and Subsection~\ref{sec:frontier} lists them as well.

\subsection{The architecture of the proof}

For $x\ge1$ let
\[
G(x)=\sum_{n\le x}\Lambda(n)\Big(\frac xn-1\Big)-x\log x+(1+\gamma)x ,
\]
and let $G_\chi$ be the analogous function for a Dirichlet character $\chi$, defined in \eqref{Gchi}. The Mellin transform of $G$ has a simple pole at each zero of $\zeta$ in the critical strip, with residue $-1/(\rho(\rho-1))$ times the multiplicity. The proof has two parts.

\emph{Part I} reduces Theorem~\ref{thm:intro} to one arithmetic statement. Section~\ref{sec:criterion} proves a quantitative form of Landau's theorem on Dirichlet integrals of one sign. If a positive combination $G+\sum_iw_iG_{\chi_i}$ is at most $Cx^\theta$ for every $x$ in windows $Y\le x\le Y^{3/2}$, for arbitrarily large $Y$, then $\zeta$ and the $L(s,\chi_i)$ do not vanish for $\operatorname{Re}s>\theta$. So it suffices to prove such a bound with $\theta=\frac12+O(\varepsilon)$, on windows, for the powers of a fixed primitive character $\psi_0$. The bound comes from an abelian CM field $E$ of degree $4n$ with $n$ about $x^{2+8\varepsilon}$, cut out by $\psi_0$ and two auxiliary primes $q$ and $d$. Section~\ref{sec:accounting} proves two identities. The first, Proposition~\ref{acct}, shows that seven inequalities between slopes of hermitian lattices, heights of period maps and local terms bound a sum $S$ over the primes inert in $\mathbb Q(\sqrt{-d})$ by an explicit main term. The second, Proposition~\ref{excess}, identifies $S$ minus that main term with a combination of the functions $G_\psi$ over the characters $\psi$ of $\mathrm{Gal}(E/\mathbb Q)$. Section~\ref{sec:point} turns a bound for this combination, summed over many weights and many nodes, into a bound for a positive combination of the $G_\psi$ at one point $x$. An average over translations removes the characters outside a small family, and an equidistributed choice of weights gives every remaining character a positive weight. The step leaves the values $\lvert G_\psi(x)\rvert$ for $3o$ auxiliary characters, where $o$ is the order of $\psi_0$. Section~\ref{sec:auxiliary} chooses $q$ and $d$ by a mean square over families of prime moduli so that these values are at most $x^{1/2+\varepsilon}$ on a whole window, without assuming any zero-free region beyond the classical one, and constructs the field $E$ and the remaining data. Section~\ref{sec:reduction} states the arithmetic theorem, Theorem~A, and deduces Theorem~\ref{thm:intro} from it.

\emph{Part II} concerns Theorem~A. Bombieri's form of Stepanov's argument~\cite{bombieri} bounds the number of rational points of a curve by comparing an auxiliary function of controlled pole order with its order of vanishing $M$ at the rational points, and the two errors balance at $M=\sqrt q$. Over a number field the count of zeros against poles becomes a slope inequality of Bost's type~\cite{bost} for a family of maps between hermitian lattices, and the order of vanishing becomes the number $m$ of steps in a chain of theta lifts between unitary groups of growing rank. To each step the argument attaches a hermitian lattice over the ring of integers of $E$ and a period map into the dual of a lattice of automorphic forms. At a prime $p$ inert in $\mathbb Q(\sqrt{-d})$, an elliptic curve with complex multiplication by $\mathbb Q(\sqrt{-d})$ has a formal group of height two, and the theta kernel of one lift is divisible by a power of $p$ whose exponent at weight $\alpha$ has the main term $\sum_{i\ \mathrm{odd}}\alpha/p^i=\alpha p/(p^2-1)$. This main term matches the contribution of $p$ to the sum $S$. Hecke operators at the inert primes cost a loss that a nondegeneracy theorem for integral representations of $\mathrm{GL}_k$ controls. The archimedean size of the period map is $m$ times a quantity close to $\sum_t\log\Gamma(a_t)$, measured against the period of the CM curve, and arithmetic intersection theory on unitary Shimura varieties bounds the slope of the lattice of automorphic forms. Bessel Poincar\'e series on a unitary group realize the period maps at one complex place, and their near orthogonality gives the slope inequality for the family. At the places where the local type of the forms is twisted, lengths of Tate--Shafarevich groups of the CM field bound the congruence defects that enter the heights, through Galois representations attached to automorphic representations of unitary groups and $p$-adic Hodge theory. The two error terms balance at $m=\sqrt x$, as in Stepanov's argument. Section~\ref{sec:organization} describes the sections of Part~II.

\subsection{Contemporaneous accounts}

Several groups began translating, mapping, explaining or distilling the formal development within days of its release. Zhang~\cite{zhang} wrote a preliminary mathematical account of the proof and of its principal ingredients, with an emphasis on their sources and history. Rubinstein~\cite{rubinsteinpage,rubinsteinzeta} maintains a guide to the proof together with a detailed exposition of the argument specialized to the Riemann zeta function, and his verification run compiled the formal development independently. Naskręcki~\cite{naskrecki} built a map of the formal proof, its structure and dependencies, and a record of its verification status, together with an atlas of the route. Epstein~\cite{epstein} wrote a human-readable blueprint that identifies the top-level statements, separates the principal mathematical crux from library material, and draws the dependency structure. Brox~\cite{brox} leads a program to distill and simplify the formal proof into shorter human-readable proofs of the Riemann hypothesis and of the generalized Riemann hypothesis for Dirichlet characters, which is in progress. Slijepčević~\cite{slijepcevic} wrote an outline of the proof in standard language, including the positivity argument, the theta-lift and lattice estimates, the arithmetic losses and the Galois-representation input. The present paper writes the argument as one continuous conventional proof with its open steps stated, so that a reader can judge its proofs and its cited external inputs without appeal to the formal development as mathematical authority.

\subsection{How this account was obtained}

The formal development has about $1.07$ million declarations, of which about $767{,}000$ lie on the dependency route of the final theorem, and a readable proof cannot arise by translating them one at a time. We began at the formal theorem \texttt{riemannHypothesis}, traced its dependencies backward, and identified the mathematical statements on which it depends. Whenever a region of the development implemented a recognizable argument, we stated and proved that argument at its natural mathematical level, and whenever the route reached an established theorem of the literature, we cited that theorem and did not use the formal development as its justification. This process does not repeat the formal verification of~\cite{leanstatement}, and it did not involve reading every declaration. The author used Claude (Fable 5.2 and Opus 5.5), GPT (including GPT 5.6 Pro) and Kimi K3.1 Pro in the reconstruction. The formal certification shows that the Lean term proves the formal statement. The reconstruction writes the content of the internal dependencies as mathematics, and a separate comparison of each cited theorem with its published statement checks hypotheses, normalizations and quantifiers.

Two findings of the reconstruction concern the strength of the internal statements. The bound for Tate--Shafarevich groups on windows of scales, Hypothesis~\ref{w4:thm:vwin}, implies a bound in the direction of Vandiver's conjecture that is stronger than the published bounds of Soul\'e~\cite{soule99} and of de~la~Torre~Dur\'an~\cite{delatorre} (Proposition~\ref{w4:prop:need}). One branch of the formal proof uses an automorphy statement for residually reducible Galois representations of unitary type in three variables, (G3) of Subsection~\ref{tw:sec:vandiver}, which lies outside the hypotheses of the published automorphy lifting theorems for residually reducible representations. Second, the bound for the covolumes of the lattices of automorphic forms needs a constant independent of the level of a deep cover of a unitary Shimura variety, and the argument obtains it from the arithmetic Hilbert--Samuel theorem on each model of a tower of models and an exact identity at the good primes, without a uniform version of the Hilbert--Samuel theorem (Subsection~\ref{w3:sec:unif}).

\subsection{Organization}\label{sec:organization}

Sections~\ref{sec:criterion} to~\ref{sec:reduction} form Part~I. Part~II begins in Section~\ref{sec:chain} with the theta kernels between hermitian lattices over the CM field and their covariance. Section~\ref{sec:inert} studies the theta function of a CM elliptic curve at an inert prime and proves the form of the kernels in the parameters of a formal group of height two. Section~\ref{sec:gain} classifies the semilinear pairings on that formal group and proves the valuation estimate for the kernels. Section~\ref{sec:chaindiv} turns this estimate into divisibility of the composite of a chain of lifts. Section~\ref{sec:lifts} constructs the lifts and shows that they form a chain datum. Section~\ref{r2:sec:loss} proves the nondegeneracy of Hecke operators that controls the loss, and hypotheses \eqref{A1} and \eqref{A2}. Section~\ref{r3:sec} proves the slope inequality \eqref{A3} for the family of period maps, and Section~\ref{sec:bessel} proves its automorphic inputs for Bessel Poincar\'e series. Section~\ref{r4:sec:R4} computes the period of the CM curve and proves \eqref{A4}, \eqref{A6} and \eqref{A7}. Section~\ref{r5:sec:localterms} proves \eqref{A5}. Section~\ref{sec:twisted} treats the twisted places, the Tate--Shafarevich bound and the congruence bound. Section~\ref{sec:proofA} constructs the data, the weights and the nodes, and proves Theorem~A from the statements of the preceding sections. Subsection~\ref{sec:frontier} collects the hypotheses used in the final deduction.


\part{Reduction to an arithmetic theorem}
\section{One-sided bounds on windows}\label{sec:criterion}

This section proves the analytic criterion of the paper. An upper bound $G_+(x)\le Cx^\theta$ for a positive combination of smoothed prime sums, known for $x$ in a sparse sequence of windows $Y\le x\le Y^\lambda$, excludes zeros of $\zeta$ and of the corresponding Dirichlet $L$-functions from the half-plane $\operatorname{Re}s>\theta$. The proof is a quantitative form of Landau's theorem on Dirichlet integrals of functions of one sign~\cite{landau}.

\subsection{The smoothed prime sums}\label{sec:functions}

Let $\Lambda$ denote the von Mangoldt function and $\gamma$ Euler's constant. For $x\ge1$ put
\begin{equation}\label{G}
G(x)=\sum_{n\le x}\Lambda(n)\Big(\frac xn-1\Big)-x\log x+(1+\gamma)x ,
\end{equation}
and for a nontrivial Dirichlet character $\chi$ put
\begin{equation}\label{Gchi}
G_\chi(x)=\sum_{n\le x}\Lambda(n)\operatorname{Re}\chi(n)\Big(\frac xn-1\Big)+\operatorname{Re}\frac{L'(1,\chi)}{L(1,\chi)}\,x .
\end{equation}
For $\operatorname{Re}s>1$ and every integer $n\ge1$,
\begin{gather*}
\int_n^\infty\Big(\frac xn-1\Big)x^{-s-1}\,dx=\frac{n^{-s}}{s(s-1)},\\
\int_1^\infty x\log x\;x^{-s-1}\,dx=\frac1{(s-1)^2},\qquad
\int_1^\infty x\;x^{-s-1}\,dx=\frac1{s-1}.
\end{gather*}
Summation over $n$ gives, for $\operatorname{Re}s>1$,
\begin{equation}\label{F}
\int_1^\infty G(x)x^{-s-1}\,dx=F(s),\qquad \int_1^\infty G_\chi(x)x^{-s-1}\,dx=\tfrac12\big(F_\chi(s)+F_{\bar\chi}(s)\big),
\end{equation}
where
\begin{equation}\label{Fchi}
F(s)=-\frac{\zeta'(s)}{\zeta(s)}\,\frac1{s(s-1)}-\frac1{(s-1)^2}+\frac{1+\gamma}{s-1},\qquad
F_\chi(s)=-\frac{L'(s,\chi)}{L(s,\chi)}\,\frac1{s(s-1)}+\frac{L'(1,\chi)}{L(1,\chi)}\,\frac1{s-1}.
\end{equation}
The expansions $-\zeta'(s)/\zeta(s)=(s-1)^{-1}-\gamma+O(s-1)$ and $1/(s(s-1))=(s-1)^{-1}-1+O(s-1)$ show that $F$ is analytic at $s=1$. The same computation, together with $L(1,\chi)\ne0$, shows that $F_\chi$ is analytic at $s=1$. Both functions are meromorphic on $\mathbb C$. In the half-plane $\operatorname{Re}s>0$ their poles are simple and lie at the zeros $\rho$ of $\zeta$ and of $L(\cdot,\chi)$, with residues
\begin{equation}\label{res}
\Res_{s=\rho}F=-\frac{m_\zeta(\rho)}{\rho(\rho-1)},\qquad \Res_{s=\rho}F_\chi=-\frac{m_\chi(\rho)}{\rho(\rho-1)},
\end{equation}
where $m_\zeta(\rho)$ and $m_\chi(\rho)$ denote the orders of vanishing of $\zeta$ and of $L(\cdot,\chi)$ at $\rho$.

We use four classical facts~\cite{davenport}. The functions $\zeta$ and $L(\cdot,\chi)$ do not vanish on the half-plane $\operatorname{Re}s\ge1$, and $\zeta$ does not vanish on the interval $0<s<1$. Mertens' estimate $\sum_{n\le x}\Lambda(n)/n=\log x+O(1)$ and Chebyshev's estimate $\sum_{n\le x}\Lambda(n)=O(x)$ hold, and for nontrivial $\chi$ the partial sums $\sum_{n\le x}\Lambda(n)\chi(n)/n$ are bounded. These estimates give a constant $B$ with
\begin{equation}\label{triv}
\lvert G(x)\rvert\le Bx,\qquad \lvert G_\chi(x)\rvert\le Bx\qquad(x\ge1).
\end{equation}
The function $G$ is continuous, because the term of index $n$ in \eqref{G} vanishes at $x=n$, and away from the integers $G'(x)=\sum_{n\le x}\Lambda(n)/n-\log x+\gamma$ is bounded. Hence there is a constant $K$ with
\begin{equation}\label{lip}
\lvert G(x)-G(y)\rvert\le K\lvert x-y\rvert\qquad(x,y\ge1),
\end{equation}
and the same argument gives \eqref{lip} for $G_\chi$.

\subsection{Concentration of the kernel and the nearest poles}

\begin{lemma}\label{tails}
Let $I(t)=t-1-\log t$. For $n\ge1$ an integer, the density
$e^{-v}v^{n-1}/(n-1)!$ on $v>0$ has mass at most $e^{-nI(t)}$
on $v\le nt$ if $0<t<1$, and on $v\ge nt$ if $t>1$.
\end{lemma}

\begin{proof}
Write $V$ for a variable with this density. Substitution in the gamma
integral gives $\mathbb E e^{sV}=(1-s)^{-n}$ for $s<1$.
Take $s=1-1/t$. On either indicated tail $sV\ge snt$, so its
probability is at most $e^{-snt}(1-s)^{-n}=e^{-nI(t)}$.
\end{proof}

\begin{lemma}\label{rot}
Let $0<\varphi<\pi/2$ and $\alpha\in\mathbb R$. Every block of
$\lceil2\pi/\varphi\rceil$ consecutive integers contains a $j$
with $\cos(\alpha+j\varphi)\ge\tfrac12$.
\end{lemma}

\begin{proof}
Around the circle, the angles in the block leave gaps of length at most
$\varphi$, including the gap back to the first angle. The arc on which
cosine is at least $1/2$ has length $2\pi/3>\varphi$.
\end{proof}

\begin{lemma}\label{near}
Suppose a closed discrete set $P\subset\mathbb C$, stable under conjugation,
contains $\rho_0$ with $\operatorname{Re}\rho_0>\theta$.
For arbitrarily large real $\sigma$ the nearest points of $P$ to $\sigma$
form a single real point or a conjugate pair, at distance $r<\sigma-\theta$.
\end{lemma}

\begin{proof}
The difference $(\sigma-\theta)^2-|\sigma-\rho_0|^2$ grows linearly
with positive leading coefficient $2(\operatorname{Re}\rho_0-\theta)$.
For large $\sigma$ the nearest points therefore lie within $\sigma-\theta$.
Two distinct points of $P$ can be equidistant from a real centre for
at most one value of $\sigma$, unless they are conjugates. Discard those
exceptional centres for the countably many pairs.
\end{proof}

\subsection{The window theorem}\label{sec:win}

\begin{theorem}\label{engine}
Let $H\colon[1,\infty)\to\mathbb R$ be measurable with $\lvert H(x)\rvert\le Bx$ for a constant $B$, and suppose that the function $A(s)=\int_1^\infty H(x)x^{-s-1}\,dx$, defined for $\operatorname{Re}s>1$, extends to a meromorphic function on $\mathbb C$. Let $0\le\theta<1$. Suppose that every pole $p$ of $A$ with $\operatorname{Re}p>\theta$ is simple and satisfies $\operatorname{Re}p<1$, and that the residue of $A$ at every real pole $p>\theta$ is positive. Suppose that there are constants $\lambda>1$ and $C$ such that, for arbitrarily large $Y$, the inequality $H(x)\le Cx^\theta$ holds for every $x$ with $Y\le x\le Y^\lambda$. Then $A$ does not have a pole in the half-plane $\operatorname{Re}s>\theta$.
\end{theorem}

\begin{proof}
Suppose the set $P$ of poles contains a point to the right of $\theta$.
Since $H$ is real, the poles and residues occur in conjugate pairs
away from the real axis. Fix $1<\alpha<\beta<\lambda$ and put
$\delta=\min\{I(1/\alpha),I(\lambda/\beta)\}>0$, with $I$ from
Lemma~\ref{tails}. Choose a centre $\sigma>1$ as in Lemma~\ref{near}
and write $a=\sigma-1$. All poles have real part less than $1$, so
\[
  1<r/a<(\sigma-\theta)/(\sigma-1).
\]
We can thus fix $\sigma$ large enough that $\log(r/a)<\delta$.

\emph{A lower bound from the poles.}
For $n\ge1$, differentiating the Mellin integral at $\sigma$ gives
\[
 T_n=\frac{(-1)^{n-1}A^{(n-1)}(\sigma)}{(n-1)!}
     =\int_0^\infty H(e^u)e^{-\sigma u}\frac{u^{n-1}}{(n-1)!}\,du .
\]
The growth bound justifies differentiation under the integral.
Subtract the principal parts of the nearest poles. The remainder is
analytic on a closed disc of radius $R>r$, so Cauchy's estimate gives
\[
 r^n T_n=\sum_{|\sigma-p|=r}
   (\Res_{s=p} A)\left(\frac r{\sigma-p}\right)^n+O((r/R)^n).
\]
For a real nearest pole the sum is its positive residue. For a pair, write
$\sigma-p=re^{-i\varphi}$ for the pole in the upper half-plane and
$\Res_{s=p}A=|c|e^{i\omega}$. The sum is
$2|c|\cos(\omega+n\varphi)$, with $0<\varphi<\pi/2$.
By Lemma~\ref{rot}, in either case there is a fixed $c_0>0$ such that
$r^n T_n\ge c_0$ for a sequence of $n$ with bounded gaps.

\emph{An upper bound from a window.}
Put $L=\log Y$ and take $\alpha aL\le n\le\beta aL$. On
$L\le u\le\lambda L$ the window bound gives a contribution to $T_n$
of at most $\max(C,0)(\sigma-\theta)^{-n}$.
On the two tails use $|H(e^u)|\le Be^u$ and change variables to $v=au$.
The resulting integrals are bounded by $Ba^{-n}$ times the mass of the
density in Lemma~\ref{tails} on $v\le aL$ or $v\ge a\lambda L$.
Since $aL\le n/\alpha$ and $a\lambda L\ge n\lambda/\beta$, we obtain
\[
 r^n T_n\le \max(C,0)\left(\frac r{\sigma-\theta}\right)^n
             +2B\exp\bigl(-n(\delta-\log(r/a))\bigr).
\]
Both terms tend to zero. As $Y$ increases along the given windows,
the interval $[\alpha aL,\beta aL]$ grows in length and eventually
contains an $n$ from the lower-bound sequence, a contradiction.
\end{proof}

\begin{corollary}\label{corA}
Let $\chi_1,\dots,\chi_N$ be nontrivial Dirichlet characters, let $w_1,\dots,w_N$ be positive real numbers, and put $G_+=G+\sum_{i=1}^Nw_iG_{\chi_i}$ with $G$ and $G_\chi$ as in \eqref{G} and \eqref{Gchi}. Let $0\le\theta<1$, and suppose that there are constants $\lambda>1$ and $C$ such that, for arbitrarily large $Y$, the inequality $G_+(x)\le Cx^\theta$ holds for every $x$ with $Y\le x\le Y^\lambda$. Then $\zeta(s)\ne0$ and $L(s,\chi_i)\ne0$ for every $i$ and every $s$ with $\operatorname{Re}s>\theta$.
\end{corollary}

\begin{proof}
By \eqref{F}, \eqref{Fchi} and \eqref{triv}, the function $H=G_+$ satisfies $\lvert H(x)\rvert\le(1+\sum_iw_i)Bx$, and its transform $A=F+\frac12\sum_iw_i(F_{\chi_i}+F_{\bar\chi_i})$ is meromorphic on $\mathbb C$. By \eqref{res}, the poles of $A$ in $\operatorname{Re}s>\theta$ are simple, they lie at the zeros $\rho$ of $\zeta$, of the $L(\cdot,\chi_i)$ and of the $L(\cdot,\bar\chi_i)$, and
\[
\Res_{s=\rho}A=-\frac{c_\rho}{\rho(\rho-1)},\qquad c_\rho=m_\zeta(\rho)+\tfrac12\sum_{i=1}^Nw_i\big(m_{\chi_i}(\rho)+m_{\bar\chi_i}(\rho)\big)\ge0 .
\]
These zeros satisfy $\operatorname{Re}\rho<1$. At a real zero $\rho\in(\theta,1)$ the residue is positive because $\rho(\rho-1)<0$. Theorem~\ref{engine} shows that $A$ does not have a pole with $\operatorname{Re}s>\theta$, so $c_\rho=0$ for every $\rho$ in that half-plane. Each term of $c_\rho$ is nonnegative and each $w_i$ is positive, so each multiplicity vanishes.
\end{proof}

If $E$ is an abelian number field with group $X$ of primitive Dirichlet characters, then $\zeta_E(s)=\prod_{\chi\in X}L(s,\chi)$, and the weights $w_\chi=1$ for the nontrivial $\chi\in X$ make $G_+$ the smoothed sum over the prime ideals of $E$. Part~II proves a bound of this kind for a positive combination of characters attached to one CM field, and Corollary~\ref{corA} then treats all of these characters at once.

\input{sec3_accounting}
\input{sec4_point}
\section{The auxiliary characters and the host field}\label{sec:auxiliary}

Corollary~\ref{asm} leaves two kinds of terms. The first is the bound \eqref{L3} at very large scales, and the second is the value $\lvert G_\psi(x)\rvert$ for the characters $\psi\in\Psi_0\setminus K$. In the application the characters of $\Psi_0\setminus K$ are products of characters of a fixed modulus with the quadratic characters of two auxiliary primes $q$ and $d$. The conductors of these characters exceed every power of $x$, so the required bound at the single point $x$ does not follow from any zero-free region for a fixed $L$-function. This section chooses $q$ and $d$ so that the bound holds on a whole window of values of $x$, by a mean square over a family of prime moduli. It then constructs the CM field $E$ cut out by $q$, $d$ and the target character, together with the remaining data that the arithmetic part of the argument uses.

\subsection{Notation and classical statements}\label{sec:classical}

For a Dirichlet character $\chi$ put $\Theta(u,\chi)=\sum_{n\le u}\Lambda(n)\chi(n)$ and $\Theta(u)=\sum_{n\le u}\Lambda(n)$. For a primitive character $\chi$ modulo $N\ge3$ and $y\ge1$ put
\begin{equation}\label{Gstar}
G^\ast_\chi(y)=\sum_{n\le y}\Lambda(n)\chi(n)\Big(\frac yn-1\Big)+\frac{L'(1,\chi)}{L(1,\chi)}\,y ,
\end{equation}
so that the function \eqref{Gchi} is $G_\chi=\operatorname{Re}G^\ast_\chi$ and $\lvert G_\chi(y)\rvert\le\lvert G^\ast_\chi(y)\rvert$. A zero of $L(s,\chi)$ means a zero in the strip $0<\operatorname{Re}s<1$. Write $N(\sigma,T,\chi)$ for the number of zeros $\rho$ with $\operatorname{Re}\rho\ge\sigma$ and $\lvert\operatorname{Im}\rho\rvert\le T$, and $\operatorname{li}(X)=\int_2^Xdt/\log t$. We use the following eight statements, in which every constant is absolute.
\begin{enumerate}
\item[(C1)] Chebyshev~\cite{chebyshev}. There is $c_0\ge1$ with $\Theta(u)\le c_0u$ and $\pi(u)\log u\le c_0u$ for every $u\ge1$.
\item[(C2)] The prime number theorem for characters~\cite[Theorem~5.27]{ik}, \cite[Ch.~20]{davenport}. There are $0<c_1\le\frac12$ and $C_1\ge1$ such that for every primitive character $\chi$ modulo $N\ge3$ and every $u\ge2$
\[
\Big\lvert\Theta(u,\chi)+\frac{u^{\beta}}{\beta}\Big\rvert\le C_1\,u\,(\log N)^4\exp\Big(-\frac{c_1\log u}{\sqrt{\log u}+\log N}\Big),
\]
where the term with $\beta$ is present when $\chi$ is real and $L(s,\chi)$ has a real zero $\beta>1-c_1/\log N$, and is absent otherwise. Moreover $\lvert\Theta(u)-u\rvert\le C_1u\exp(-c_1\sqrt{\log u})$ for $u\ge2$, and the series $\sum_n\Lambda(n)\chi(n)/n$ converges to $-L'(1,\chi)/L(1,\chi)$.
\item[(C3)] Page~\cite{page}, \cite[Ch.~14]{davenport}. There is $c_2>0$ such that for every $R\ge3$ at most one primitive real character of conductor at most $R$ has a real zero $\beta>1-c_2/\log R$.
\item[(C4)] Zeros and the explicit formula~\cite[Ch.~16 and~19]{davenport}. Let $\chi$ be a primitive character modulo $N\ge3$. The map $\rho\mapsto1-\bar\rho$ permutes the zeros of $L(s,\chi)$, and for every real $t$ the number of zeros with $\lvert\operatorname{Im}\rho-t\rvert\le1$ is at most $C_2\log(N(\lvert t\rvert+2))$. For $2\le T\le u$,
\[
\Big\lvert\Theta(u,\chi)+\frac{u^{\beta_1}}{\beta_1}+\sum_{\lvert\operatorname{Im}\rho\rvert<T}\frac{u^\rho}{\rho}\Big\rvert\le C_2\Big(\frac{u\log^2(Nu)}T+u^{1/4}\log u\Big),
\]
where the term with $\beta_1$ is present when the zero $\beta$ of (C2) exists, in which case $\beta_1=\beta$ and the sum omits $\beta_1$ and $1-\beta_1$.
\item[(C5)] Montgomery's density theorem~\cite[Theorem~12.2]{montgomerytopics}. There is $C_3$ such that for $Q,T\ge2$ and $\frac45\le\sigma\le1$
\[
\sum_{N\le Q}\ \sum_{\chi}N(\sigma,T,\chi)\le C_3\,(Q^2T)^{5(1-\sigma)/2}\log^{14}(QT),
\]
where the inner sum runs over the primitive characters modulo $N$.
\item[(C6)] Primes in the class $3$ modulo $8$~\cite[Ch.~20]{davenport}. There is $U_0$ such that for $U\ge U_0$ the number of primes $q\equiv3\pmod8$ with $\sqrt U<q\le U$ is at least $U/(5\log U)$.
\item[(C7)] Primes in progressions of large modulus. There are $c_6>0$, $A_0'$ and $C_0$ with the following property. For every integer $M$ divisible by $8$, consider the pairs $(\chi_1,\beta_1)$ in which $\chi_1$ is a real nonprincipal character modulo $M$ and $\beta_1$ is a real zero of $L(s,\chi_1)$ with $1-1/(8\log M)<\beta_1<1$. If such pairs exist, one of them can be chosen so that the following estimate holds with the chosen pair, and if no such pair exists the estimate holds with the second form of the main term. For every $X\ge2$ with $C_0\log M\cdot\log(\log M+2)\le\log X$ and every $b$ coprime to $M$, the number $\pi(X;M,b)$ of primes $p\le X$ with $p\equiv b\pmod M$ satisfies
\[
\big\lvert\pi(X;M,b)-\mathfrak m(X;M,b)\big\rvert\le A_0'\Big(e^{-c_6\log X/\log M}+e^{-\sqrt{c_6\log X}}\Big)\,\mathfrak m(X;M,b),
\]
where $\mathfrak m(X;M,b)=\big(\operatorname{li}(X)-\chi_1(b)\operatorname{li}(X^{\beta_1})\big)/\varphi(M)$ for the chosen pair, and $\mathfrak m(X;M,b)=\operatorname{li}(X)/\varphi(M)$ when no such pair exists.
\item[(C8)] The distance of a real zero from $1$~\cite[Ch.~14]{davenport}. There is $0<c_7\le1$ such that every real zero $\beta$ of $L(s,\chi)$, for a real nonprincipal character $\chi$ modulo $M\ge3$, satisfies $\beta\le1-c_7/(M^{1/2}\log^2M)$.
\end{enumerate}

\begin{proof}[Proof of (C7)]
Apply the Chebotarev theorem of Thorner and Zaman~\cite[Theorem~1.4]{thornerzaman19}
to $L=\mathbb Q(\zeta_M)$ over $\mathbb Q$. Take the abelian subgroup to be the
whole Galois group $G=(\mathbb Z/M\mathbb Z)^\times$ and the conjugacy class
to be $\{b\}$. The fixed field is $\mathbb Q$. At $p\nmid M$ the arithmetic
Frobenius sends $\zeta_M$ to $\zeta_M^p$, so the theorem counts precisely
$\pi(X;M,b)$, with density $1/\varphi(M)$.

The conductor parameter in their formula~(1-11) is the largest conductor
of a primitive character inducing a character modulo $M$. This maximum
equals $M$. Indeed, for each $p^a\parallel M$ the reduction from units
modulo $p^a$ to units modulo $p^{a-1}$ has nontrivial kernel
(for $p=2$ we have $a\ge3$). Characters of a finite abelian group separate
points, so we can choose a character nontrivial on this kernel. Its conductor
is $p^a$, and the product of the local characters has conductor $M$.

Thus $D_K\mathcal Q n_K^{n_K}=M$ in the theorem. Its exceptional-zero
interval is exactly $(1-1/(8\log M),1)$. Passing from a primitive character
to the character it induces modulo $M$ removes Euler factors which are
nonzero on this interval. The principal character has no zero there, since
$\zeta$ has none. Hence the exceptional pair, if present, gives the pair
in (C7) for all $X$ and $b$, with coefficient $\chi_1(b)$. In its absence
the correction term vanishes.

Theorem~1.4 now gives the stated relative error with $c_6=c_3$ of that
theorem and an absolute implied constant, for $X\ge M^{c_2}$.
The normalization of $\operatorname{li}$ at $2$ gives the same difference
when $\chi_1(b)=1$. In the other cases an additive constant in the
logarithmic integral costs $O(\log X/X)$ relative to the main term,
which the displayed error absorbs. Taking $C_0\ge\max(1,c_2)$ suffices
since $M\ge8$ and $\log(\log M+2)>1$.
\end{proof}

\begin{lemma}\label{boxpsi}
Let $\chi$ be a primitive character modulo $N\ge3$ and let $Q\ge2$. Suppose that $L(s,\chi)\ne0$ for every $s$ with $\operatorname{Re}s>1-\frac1{400}$ and $\lvert\operatorname{Im}s\rvert\le Q$. Then for every $u\ge2$ and every $T$ with $2\le T\le\min(u,Q)$
\[
\lvert\Theta(u,\chi)\rvert\le C_4\Big(u^{1-1/400}\log^2(NT)+\frac{u\log^2(Nu)}T\Big),
\]
with an absolute constant $C_4$.
\end{lemma}

\begin{proof}
Let $\rho=\beta+i\gamma$ be a zero with $\lvert\gamma\rvert<T$. Then $\beta\le1-\frac1{400}$ by hypothesis, and the same bound for the zero $1-\bar\rho$ gives $\beta\ge\frac1{400}$, so $\lvert u^\rho/\rho\rvert\le u^{1-1/400}/\max(\frac1{400},\lvert\gamma\rvert)$. By the count of zeros in (C4),
\[
\sum_{\lvert\gamma\rvert<T}\frac1{\max(\frac1{400},\lvert\gamma\rvert)}\le400\,C_2\log(3N)+\sum_{1\le k\le T}\frac{2C_2\log(N(k+3))}{k}\le C\log^2(NT).
\]
If the term with $\beta_1$ is present, then $\frac12\le\beta_1\le1-\frac1{400}$ and $u^{\beta_1}/\beta_1\le2u^{1-1/400}$. Finally $u^{1/4}\log u\le Cu^{1-1/400}$, and the explicit formula of (C4) gives the claim.
\end{proof}

\subsection{Identities and the thin scales}

\begin{lemma}\label{gchisplit}
Let $\chi$ be a primitive character modulo $N\ge3$, and let $1\le y\le y'\le Z$. Then
\begin{align}
G^\ast_\chi(y)&=-\Theta(y,\chi)-y\sum_{n>y}\frac{\Lambda(n)\chi(n)}{n}=-y\int_y^\infty\frac{\Theta(u,\chi)}{u^2}\,du,\label{id1}\\
G^\ast_\chi(y)&=-S_Z(y)-y\,T_Z,\label{id2}\\
\frac{S_Z(y')}{y'}-\frac{S_Z(y)}{y}&=-\int_y^{y'}\frac{\Theta(u,\chi)}{u^2}\,du,\label{id3}
\end{align}
where
\[
S_Z(y)=\sum_{n\le Z}\Lambda(n)\chi(n)\min\Big(1,\frac yn\Big),\qquad T_Z=\sum_{n>Z}\frac{\Lambda(n)\chi(n)}{n}=-\frac{\Theta(Z,\chi)}{Z}+\int_Z^\infty\frac{\Theta(u,\chi)}{u^2}\,du .
\]
The identity \eqref{id3} holds with $\Lambda(n)\chi(n)$ replaced by arbitrary complex coefficients in $S_Z$ and in $\Theta$. Moreover $G(y)=-y\int_y^\infty(\Theta(u)-u)\,u^{-2}\,du$ for the function $G$ of \eqref{G}.
\end{lemma}

\begin{proof}
By (C2) the series $\sum_n\Lambda(n)\chi(n)/n$ converges to $-L'(1,\chi)/L(1,\chi)$, so
\[
G^\ast_\chi(y)=y\sum_{n\le y}\frac{\Lambda(n)\chi(n)}n-\Theta(y,\chi)-y\sum_{n\ge1}\frac{\Lambda(n)\chi(n)}n ,
\]
which is the first equality of \eqref{id1}. For $V>y$, summation by parts gives
\begin{equation}\label{parts}
\sum_{y<n\le V}\frac{\Lambda(n)\chi(n)}n=\frac{\Theta(V,\chi)}V-\frac{\Theta(y,\chi)}y+\int_y^V\frac{\Theta(u,\chi)}{u^2}\,du .
\end{equation}
By (C2), $\Theta(V,\chi)/V$ tends to zero as $V\to\infty$ and the integral converges, which gives the second equality of \eqref{id1} and, with $y=Z$, the formula for $T_Z$. Splitting the sum over $n>y$ at $Z$ gives \eqref{id2}, because $\min(1,y/n)$ equals $1$ for $n\le y$ and $y/n$ for $y<n\le Z$. The same splitting and \eqref{parts} with $V=Z$ give
\[
\frac{S_Z(y)}y=\frac{\Theta(y,\chi)}y+\sum_{y<n\le Z}\frac{\Lambda(n)\chi(n)}n=\frac{\Theta(Z,\chi)}Z+\int_y^Z\frac{\Theta(u,\chi)}{u^2}\,du ,
\]
and the difference of this identity at $y'$ and at $y$ is \eqref{id3}. This step is a summation by parts, valid for arbitrary coefficients. For $G$, the same summation by parts gives $\sum_{n\le y}\Lambda(n)/n=\Theta(y)/y+\int_1^y\Theta(u)u^{-2}\,du$, and (C2) together with the expansion of $\zeta'/\zeta$ at $s=1$ gives $\int_1^\infty(\Theta(u)-u)u^{-2}\,du=-1-\gamma$. Hence
\[
y\sum_{n\le y}\frac{\Lambda(n)}n-\Theta(y)=y\log y-(1+\gamma)y-y\int_y^\infty\frac{\Theta(u)-u}{u^2}\,du ,
\]
and \eqref{G} gives the claim.
\end{proof}

\begin{proposition}\label{thinprop}
Let $\chi$ be a primitive character modulo $N\ge3$, let $\lambda\ge\max(1,\log N)$ and $0<c_\ast\le\frac12$, and suppose that either $\chi$ is not real, or $L(\beta,\chi)\ne0$ for every real $\beta$ with $1-c_\ast<\beta<1$. Then for every $y\ge2$, with $t_0=\log y$,
\[
\frac{\lvert G^\ast_\chi(y)\rvert}{y}\le\frac{2e^{-c_\ast t_0}}{c_\ast}+C_1\lambda^4\Big(\frac{2\lambda}{c_1}\,e^{-c_1t_0/(2\lambda)}+\frac8{c_1^2}\Big(1+\frac{c_1\sqrt{t_0}}2\Big)e^{-c_1\sqrt{t_0}/2}\Big).
\]
\end{proposition}

\begin{proof}
If the term with $\beta$ in (C2) is present, then $\frac12\le\beta\le1-c_\ast$, so $u^\beta/\beta\le2u^{1-c_\ast}$. Since $\log N\le\lambda$, the bound of (C2) gives, for every $u\ge2$,
\[
\lvert\Theta(u,\chi)\rvert\le2u^{1-c_\ast}+C_1\,u\,\lambda^4\exp\Big(-\frac{c_1\log u}{\sqrt{\log u}+\lambda}\Big).
\]
For $t>0$ the quotient $t/(\sqrt t+\lambda)$ is at least $\sqrt t/2$ when $\sqrt t\ge\lambda$ and at least $t/(2\lambda)$ otherwise, so with $t=\log u$ the exponential is at most $e^{-c_1t/(2\lambda)}+e^{-c_1\sqrt t/2}$. By \eqref{id1} and the substitution $u=e^t$,
\[
\frac{\lvert G^\ast_\chi(y)\rvert}y\le\int_y^\infty\frac{2\,du}{u^{1+c_\ast}}+C_1\lambda^4\int_{t_0}^\infty\Big(e^{-c_1t/(2\lambda)}+e^{-c_1\sqrt t/2}\Big)dt ,
\]
and the three integrals equal $2y^{-c_\ast}/c_\ast$, $(2\lambda/c_1)e^{-c_1t_0/(2\lambda)}$ and $(8/c_1^2)(1+c_1\sqrt{t_0}/2)e^{-c_1\sqrt{t_0}/2}$.
\end{proof}

For the trivial character, the last identity of Lemma~\ref{gchisplit} and the second bound of (C2) give the same decay for $G$.

\subsection{A mean square over prime moduli}

\begin{proposition}\label{meansq}
Let $Z\ge1$ and $\eta\ge0$, let $P$ be a finite set of primes $q>Z$, let $E_0$ be a finite set of squarefree integers $s>1$, and let $c_n$ for $1\le n\le Z$ be complex numbers. Suppose that $\lvert\sum_{q\in P}(\frac sq)\rvert\le\eta\,\lvert P\rvert$ for every squarefree integer $s$ with $1<s\le Z^2$ and $s\notin E_0$. Then
\[
\sum_{q\in P}\Big\lvert\sum_{n\le Z}c_n\Big(\frac nq\Big)\Big\rvert^2\le\lvert P\rvert\Big(\sum_{\substack{n,n'\le Z\\ nn'\ \mathrm{a\ square}}}\lvert c_nc_{n'}\rvert+\sum_{s\in E_0}\ \sum_{\substack{n,n'\le Z\\ nn'/s\ \mathrm{a\ square}}}\lvert c_nc_{n'}\rvert+\eta\Big(\sum_{n\le Z}\lvert c_n\rvert\Big)^2\Big).
\]
\end{proposition}

\begin{proof}
Expanding the square gives $\sum_{n,n'}c_n\overline{c_{n'}}\sum_{q\in P}(\frac{nn'}q)$. Each $q\in P$ exceeds $Z$, so $q$ does not divide $nn'$. Write $nn'=sr^2$ with $s$ squarefree, so that $s\le Z^2$ and the symbol equals $(\frac sq)$. If $s=1$ the inner sum is $\lvert P\rvert$, if $s\in E_0$ it has absolute value at most $\lvert P\rvert$, and otherwise it has absolute value at most $\eta\lvert P\rvert$.
\end{proof}

In both choices of auxiliary primes the sum of the symbols $(\frac sq)$ over the candidates is small for most squarefree $s$ and large for the $s$ in a set of the following kind. For $z\ge2$ and an integer $s^\ast\ge1$ without prime factors $p\le z$, let $\mathcal S(z,s^\ast)$ be the set of squarefree integers $s>1$ whose largest divisor without prime factors $p\le z$ equals $1$ or $s^\ast$.

\begin{lemma}\label{coef}
Let $3\le y\le Z$, let $\kappa(n)$ be complex numbers with $\lvert\kappa(n)\rvert\le1$, and put $c_n=\Lambda(n)\kappa(n)\min(1,y/n)$ for $n\le Z$. Let $z\ge2$ with $z\log y\le\sqrt y$, and let $s^\ast\ge1$ be an integer without prime factors $p\le z$. With the constant $c_0$ of \textup{(C1)},
\begin{align*}
\sum_{\substack{n,n'\le Z\\ nn'\ \mathrm{a\ square}}}\lvert c_nc_{n'}\rvert&\le60\,c_0^2\,y\log y,\\
\sum_{s\in\mathcal S(z,s^\ast)}\ \sum_{\substack{n,n'\le Z\\ nn'/s\ \mathrm{a\ square}}}\lvert c_nc_{n'}\rvert&\le81\,c_0^2\,y,\\
\sum_{n\le Z}\lvert c_n\rvert&\le3\,c_0\,y\log Z .
\end{align*}
\end{lemma}

\begin{proof}
For a prime $p$ put $W_p=\sum_{a\ge1}\log p\,\min(1,y/p^a)$, and put $W_\square=\sum_m\Lambda(m)\min(1,y/m^2)$, which bounds the sum of $\lvert c_n\rvert$ over the squares $n$.

If $p\le y$, the powers $p^a\le y$ contribute at most $\log y$ to $W_p$, and the powers $p^a>y$ contribute at most $2\log p$, because the numbers $y/p^a$ are less than one and decrease by the factor $p\ge2$. If $p>y$, then $W_p\le2y\log p/p\le2\log y$. Hence $W_p\le3\log y$ for every $p$, and by (C1)
\[
\sum_{p\le y}W_p^2\le9\,\pi(y)\log^2y\le9c_0\,y\log y .
\]
For $p>y$ we group the primes in the intervals $2^ky<p\le2^{k+1}y$ and use (C1) again,
\[
\sum_{p>y}W_p^2\le4y^2\sum_{k\ge0}\pi(2^{k+1}y)\,\frac{\log^2(2^{k+1}y)}{4^ky^2}\le8c_0\,y\sum_{k\ge0}\frac{\log y+(k+1)\log2}{2^k}\le40c_0\,y\log y ,
\]
since $\sum_k2^{-k}=2$, $\sum_k(k+1)2^{-k}=4$ and $4\log2\le3\log y$. By (C1) and a summation by parts,
\[
W_\square\le\Theta(\sqrt y)+y\sum_{m>\sqrt y}\frac{\Lambda(m)}{m^2}\le c_0\sqrt y+2y\int_{\sqrt y}^\infty\frac{\Theta(u)}{u^3}\,du\le3c_0\sqrt y .
\]
Let $n=p^a$ and $n'=p'^{\,b}$ be prime powers. If $nn'$ is a square, then either $p=p'$, or $n$ and $n'$ are both squares, so the first sum is at most $\sum_pW_p^2+W_\square^2\le49c_0y\log y+9c_0^2y\le60c_0^2y\log y$.

For the second sum, write $nn'=sr^2$ with $s>1$ squarefree, so that the prime factors of $s$ are the primes that occur in $nn'$ to an odd power. Let $\mathcal A$ be the set of the $n\le Z$ that are squares or powers of a prime $p\le z$. The sum of $\lvert c_n\rvert$ over $\mathcal A$ is at most
\[
W_\square+\sum_{p\le z}W_p\le3c_0\sqrt y+3z\log y\le6c_0\sqrt y .
\]
If every prime factor of $s$ is at most $z$, then $n$ and $n'$ both lie in $\mathcal A$, because a prime $p>z$ that occurs in $n$ to an odd power must occur in $n'$ to an odd power, and then $nn'$ is a square. These pairs contribute at most $36c_0^2y$. Otherwise $s\in\mathcal S(z,s^\ast)$ forces the primes $p>z$ dividing $s$ to be the prime factors of $s^\ast$, and $nn'$ has at most two prime factors. If $s^\ast$ is a prime $p^\ast$, then one of $n$, $n'$ is an odd power of $p^\ast$ and the other lies in $\mathcal A$, which contributes at most $2W_{p^\ast}\cdot6c_0\sqrt y\le36c_0\sqrt y\log y$. If $s^\ast$ is a product of two primes, then $n$ and $n'$ are odd powers of these primes, which contributes at most $18\log^2y$. The integer $s^\ast$ is fixed, so at most one of the last two cases occurs, and $18\log^2y\le36c_0\sqrt y\log y\le36c_0^2y$ because $\log y\le\sqrt y$ and $c_0\ge1$. The second sum is therefore at most $72c_0^2y$. Finally, by (C1) and a summation by parts,
\[
\sum_{n\le Z}\lvert c_n\rvert\le\Theta(y)+y\sum_{y<n\le Z}\frac{\Lambda(n)}n\le c_0y+y\Big(c_0+c_0\log\frac Zy\Big)\le3c_0\,y\log Z .\qedhere
\]
\end{proof}

\subsection{Uniformity on a window}

A mean square bounds a sum at one point $y$ for most $q$, and Corollary~\ref{asm} needs the bound at every point of a window. The identity \eqref{id3} controls a dyadic block of values of $y$ by its left endpoint and one integral.

\begin{proposition}\label{window}
Let $0<\tau\le\frac12$, $8\le Y\le X$ and $Z\ge2X$, and let $\eta\ge0$ with $\eta\,X\log^2Z\le1$. Let $z\ge2$ with $z\log Y\le\sqrt Y$, and let $s^\ast\ge1$ be an integer without prime factors $p\le z$. Let $P$ be a finite set of primes $q>Z$, and suppose that $\lvert\sum_{q\in P}(\frac sq)\rvert\le\eta\lvert P\rvert$ for every squarefree $s$ with $1<s\le Z^2$ and $s\notin\mathcal S(z,s^\ast)$. Let $\lvert\kappa(n)\rvert\le1$, and for $q\in P$ put
\[
S_q(y)=\sum_{n\le Z}\Lambda(n)\kappa(n)\Big(\frac nq\Big)\min\Big(1,\frac yn\Big).
\]
Then the number of $q\in P$ for which $\lvert S_q(y)\rvert>\frac12y^{1/2+\tau}$ for some $y$ with $Y\le y\le X$ is at most
\[
\frac{16000\,c_0^2}{(1-4^{-\tau})^2}\;\lvert P\rvert\;Y^{-2\tau}\log(2Y).
\]
\end{proposition}

\begin{proof}
Put $\Theta_q(u)=\sum_{n\le u}\Lambda(n)\kappa(n)(\frac nq)$. The function $\sqrt y/\log y$ increases for $y\ge8$, so $z\log y\le\sqrt y$ for every $y\ge Y$. For $Y\le y\le2X$, Proposition~\ref{meansq} with $E_0$ the set of the $s\le Z^2$ in $\mathcal S(z,s^\ast)$, together with Lemma~\ref{coef}, gives
\[
\sum_{q\in P}\lvert S_q(y)\rvert^2\le\lvert P\rvert\,c_0^2\Big(60\,y\log y+81\,y+9\,\eta\,y^2\log^2Z\Big)\le159\,c_0^2\,\lvert P\rvert\,y\log y ,
\]
because $\eta y^2\log^2Z\le2y$. The same two results with $Z$ replaced by $u$ and $y=u$ give $\sum_{q\in P}\lvert\Theta_q(u)\rvert^2\le159c_0^2\lvert P\rvert u\log u$ for $Y\le u\le2X$. Put $J_q(y)=\int_y^{2y}\lvert\Theta_q(u)\rvert^2u^{-2}\,du$, so that for $Y\le y\le X$
\[
\sum_{q\in P}J_q(y)\le159\,c_0^2\,\lvert P\rvert\int_y^{2y}\frac{\log u}u\,du\le159\,c_0^2\,\lvert P\rvert\log(2y).
\]
Fix $y$ with $Y\le y\le X$ and let $y\le y'\le2y$. By \eqref{id3} and the Cauchy--Schwarz inequality,
\begin{equation}\label{block}
\begin{aligned}
\lvert S_q(y')\rvert&\le\frac{y'}y\,\lvert S_q(y)\rvert+y'\int_y^{2y}\frac{\lvert\Theta_q(u)\rvert}{u^2}\,du\\
&\le2\,\lvert S_q(y)\rvert+2y\,J_q(y)^{1/2}\Big(\int_y^{2y}\frac{du}{u^2}\Big)^{1/2}=2\,\lvert S_q(y)\rvert+\sqrt{2y}\,J_q(y)^{1/2}.
\end{aligned}
\end{equation}
By Chebyshev's inequality, the number of $q$ with $\lvert S_q(y)\rvert>\frac18y^{1/2+\tau}$ is at most $64\cdot159\,c_0^2\lvert P\rvert y^{-2\tau}\log y$, and the number of $q$ with $J_q(y)>\frac1{32}y^{2\tau}$ is at most $32\cdot159\,c_0^2\lvert P\rvert y^{-2\tau}\log(2y)$. For every other $q$, \eqref{block} gives
\[
\lvert S_q(y')\rvert\le\tfrac14y^{1/2+\tau}+\sqrt{2y}\,\frac{y^{\tau}}{\sqrt{32}}=\tfrac12y^{1/2+\tau}\le\tfrac12y'^{\,1/2+\tau}\qquad(y\le y'\le2y).
\]
The blocks with left endpoints $y_i=2^iY$, $0\le i\le\log_2(X/Y)$, cover the interval from $Y$ to $X$. The number of exceptional $q$ is therefore at most
\[
96\cdot159\,c_0^2\,\lvert P\rvert\sum_{i\ge0}(2^iY)^{-2\tau}\big(\log(2Y)+i\log2\big)\le\frac{15264\,c_0^2}{(1-4^{-\tau})^2}\,\lvert P\rvert\,Y^{-2\tau}\log(2Y),
\]
since $\sum_ir^i+\sum_iir^i=(1-r)^{-2}$ for $r=4^{-\tau}$.
\end{proof}

\subsection{The first auxiliary prime}

\begin{proposition}\label{legendre}
There are absolute constants $c_3>0$ and $z_0$ with the following property. Let $z\ge z_0$, let $P$ be the set of primes $q\equiv3\pmod8$ with $e^{z/2}<q\le e^z$, and let $R\ge3$ with $\log(8R)\le\sqrt{z/2}$. Then there is an odd integer $s^\ast\ge1$ such that
\[
\Big\lvert\sum_{q\in P}\Big(\frac sq\Big)\Big\rvert\le e^{-c_3\sqrt z}\,\lvert P\rvert\qquad\text{for every squarefree }s\text{ with }1<s\le R\text{ and }s\notin\mathcal S(2,s^\ast).
\]
\end{proposition}

\begin{proof}
Put $U=e^z$. Let $r^\ast$ be the conductor of the character excluded in (C3) at the level $8R$, and let $s^\ast$ be the odd part of $r^\ast$. If that character does not exist, put $s^\ast=1$. Let $s\notin\mathcal S(2,s^\ast)$ be squarefree with $1<s\le R$, and write $s=2^es_1$ with $s_1$ odd, so that $s_1\ne1$ and $s_1\ne s^\ast$. Every $q\in P$ exceeds $s$ and satisfies $q\equiv3\pmod8$, so $(\frac2q)=-1$, and the reciprocity law for the Jacobi symbol gives $(\frac{s_1}q)=(-1)^{(s_1-1)/2}(\frac q{s_1})$. We detect the class of $q$ modulo $8$ with the four characters $\varpi$ modulo $8$ and obtain
\[
\sum_{q\in P}\Big(\frac sq\Big)=\pm\frac14\sum_{\varpi}\varpi(3)\sum_{\substack{\sqrt U<q\le U\\ q\ \mathrm{prime}}}\chi'_\varpi(q),
\]
where $\chi'_\varpi$ is the primitive real character that induces $\varpi\cdot(\frac{\cdot}{s_1})$. Its conductor is $s_1$ times the conductor of $\varpi$, so it is at least $3$ and at most $8R$, its odd part is $s_1\ne s^\ast$, and $\chi'_\varpi$ is not the character excluded in (C3).

Let $\chi'$ be one of these four characters and let $\sqrt U\le V\le U$, so that $\log V\ge z/2$ and $\log(8R)\le\sqrt{\log V}$. If the term with $\beta$ in (C2) is present, then $\frac12\le\beta\le1-c_2/\log(8R)$ by (C3), and $V^\beta/\beta\le2V\exp(-c_2\sqrt{z/2})$. The exponent in (C2) is at least $c_1\sqrt{\log V}/2\ge c_1\sqrt z/3$, and $(\log8R)^4\le z^2/4$. Hence
\[
\lvert\Theta(V,\chi')\rvert\le V\epsilon_1,\qquad\epsilon_1=2e^{-c_2\sqrt{z/2}}+\tfrac14C_1z^2e^{-c_1\sqrt z/3}.
\]
The prime powers $p^a\le V$ with $a\ge2$ contribute at most $2c_0\sqrt V\log V$ to $\Theta(V,\chi')$ by (C1), so $\vartheta(V)=\sum_{p\le V}\chi'(p)\log p$ satisfies $\lvert\vartheta(V)\rvert\le V(\epsilon_1+2c_0ze^{-z/4})$ for $\sqrt U\le V\le U$. A summation by parts gives
\[
\Big\lvert\sum_{\sqrt U<q\le U}\chi'(q)\Big\rvert=\Big\lvert\frac{\vartheta(U)}{\log U}-\frac{\vartheta(\sqrt U)}{\log\sqrt U}+\int_{\sqrt U}^U\frac{\vartheta(V)}{V\log^2V}\,dV\Big\rvert\le\frac{4U}z\big(\epsilon_1+2c_0ze^{-z/4}\big)
\]
for $z\ge4$. By (C6), $\lvert P\rvert\ge U/(5z)$ for large $z$, and therefore
\[
\Big\lvert\sum_{q\in P}\Big(\frac sq\Big)\Big\rvert\le20\,\lvert P\rvert\,\big(\epsilon_1+2c_0ze^{-z/4}\big)\le e^{-c_3\sqrt z}\,\lvert P\rvert
\]
for $z\ge z_0$, with $c_3=\frac14\min(c_1,c_2)$.
\end{proof}

\begin{proposition}\label{tailprop}
There is an absolute constant $C_5$ with the following property. Let $\chi$ be a primitive character modulo $N\ge3$, put $\lambda=\log N$, and let $Q\ge N$. Suppose that $L(s,\chi)\ne0$ for every $s$ with $\operatorname{Re}s>1-\frac1{400}$ and $\lvert\operatorname{Im}s\rvert\le Q$. Then for every $Z$ with $e^{800}\le Z\le\exp(\lambda^4)$ the number $T_Z=\sum_{n>Z}\Lambda(n)\chi(n)/n$ satisfies
\[
\lvert T_Z\rvert\le C_5\,\lambda^{12}\Big(Z^{-1/400}+\frac1Q+e^{-c_1\lambda^2/2}\Big),\qquad\Big\lvert\frac{L'(1,\chi)}{L(1,\chi)}\Big\rvert\le c_0(1+\log Z)+\lvert T_Z\rvert .
\]
\end{proposition}

\begin{proof}
By Lemma~\ref{gchisplit}, $T_Z=-\Theta(Z,\chi)/Z+\int_Z^\infty\Theta(u,\chi)u^{-2}\,du$. Put $V_1=\exp(\lambda^4)$. For $Z\le u\le V_1$ we apply Lemma~\ref{boxpsi} with $T=\min(u^{1/400},Q)$, which satisfies $2\le T\le\min(u,Q)$ and $1/T\le u^{-1/400}+1/Q$. Since $\log(Nu)\le2\lambda^4$,
\[
\lvert\Theta(u,\chi)\rvert\le4C_4\lambda^8\Big(2u^{1-1/400}+\frac uQ\Big)\qquad(Z\le u\le V_1).
\]
This bounds $\lvert\Theta(Z,\chi)\rvert/Z$ by $4C_4\lambda^8(2Z^{-1/400}+1/Q)$, and
\[
\int_Z^{V_1}\frac{\lvert\Theta(u,\chi)\rvert}{u^2}\,du\le4C_4\lambda^8\Big(800\,Z^{-1/400}+\frac{\lambda^4}Q\Big).
\]
For $u\ge V_1$ we use (C2). If the term with $\beta$ is present, then $\frac12\le\beta\le1-\frac1{400}$ by hypothesis, and $\log u\ge\lambda^2$ gives the exponent $c_1\sqrt{\log u}/2$, so
\[
\lvert\Theta(u,\chi)\rvert\le2u^{1-1/400}+C_1\,u\,\lambda^4\,e^{-c_1\sqrt{\log u}/2}\qquad(u\ge V_1).
\]
The integral of the first term against $u^{-2}du$ over $u\ge V_1$ is $800V_1^{-1/400}\le800Z^{-1/400}$. With $u=e^t$ the integral of the second term is
\[
C_1\lambda^4\int_{\lambda^4}^\infty e^{-c_1\sqrt t/2}\,dt=\frac{8C_1\lambda^4}{c_1^2}\Big(1+\frac{c_1\lambda^2}2\Big)e^{-c_1\lambda^2/2}.
\]
The sum of these bounds gives the first claim. For the second, $-L'(1,\chi)/L(1,\chi)=\sum_{n\le Z}\Lambda(n)\chi(n)/n+T_Z$, and (C1) with a summation by parts gives $\sum_{n\le Z}\Lambda(n)/n\le c_0(1+\log Z)$.
\end{proof}

\begin{proposition}\label{boxes}
Let $N_0\ge1$ and $Q\ge2$, and let $B$ be a set of odd primes $q>N_0$ with $N_0q\le Q$. Suppose that for every $q\in B$ there is a primitive character $\xi$ of conductor $r$ dividing $N_0$ such that $L(s,\xi\cdot(\frac{\cdot}{q}))$ has a zero with $\operatorname{Re}s>1-\frac1{400}$ and $\lvert\operatorname{Im}s\rvert\le Q$. Then
\[
\lvert B\rvert\le C_3\,Q^{3/160}\,(2\log Q)^{14}.
\]
\end{proposition}

\begin{proof}
The character $\xi\cdot(\frac{\cdot}{q})$ is primitive of conductor $rq\le Q$, because $q$ does not divide $N_0$. Since $q$ is the largest prime factor of $rq$, different $q\in B$ give different primitive characters, each with $N(1-\frac1{400},Q,\chi)\ge1$. Statement (C5) with $T=Q$ and $\sigma=1-\frac1{400}$ bounds their number by $C_3(Q^3)^{1/160}\log^{14}(Q^2)$.
\end{proof}

\begin{theorem}\label{stage1}
Let $N_0\ge1$ be an integer, and let $K\ge1$, $0<\varepsilon\le\frac1{52}$ and $0<\tau\le\frac12$. There is $Y_0$ such that the following holds for every $Y\ge Y_0$. Put $z=Y^{\varepsilon/5}$, and let $P$ be the set of primes $q\equiv3\pmod8$ with $e^{z/2}<q\le e^z$. Then $P$ is nonempty, and for all $q\in P$ outside a set of at most $Y^{-\tau}\lvert P\rvert$ elements, every primitive character $\xi$ of conductor dividing $N_0$ and the character $\chi=\xi\cdot(\frac{\cdot}{q})$ satisfy
\begin{enumerate}
\item[(i)] $L(s,\chi)\ne0$ for every $s$ with $\operatorname{Re}s>1-\frac1{400}$ and $\lvert\operatorname{Im}s\rvert\le N_0e^z$,
\item[(ii)] $\lvert L'(1,\chi)/L(1,\chi)\rvert\le401\,c_0\,K\log Y$,
\item[(iii)] $\lvert G^\ast_\chi(y)\rvert\le y^{1/2+\tau}$ for every $y$ with $Y\le y\le Y^K$.
\end{enumerate}
\end{theorem}

\begin{proof}
Put $X=Y^K$, $Z=X^{400}$, $Q=N_0e^z$ and $U=e^z$. Every inequality that we use between powers of $Y$, powers of $\log Y$ and exponentials of powers of $Y$ holds for $Y\ge Y_0$, with $Y_0$ depending on $N_0$, $K$, $\varepsilon$ and $\tau$. In particular $e^{z/2}>Z^2\ge N_0$, $2X\le Z$, and $\log(8Z^2)\le\sqrt{z/2}$.

\emph{The Legendre symbols.} By (C6) the set $P$ is nonempty and $\lvert P\rvert\ge e^z/(5z)$. Proposition~\ref{legendre} with $R=Z^2$ gives an odd integer $s^\ast$ for which the hypothesis of Proposition~\ref{window} holds with the parameter $2$ in place of $z$ and with $\eta=e^{-c_3\sqrt z}$, and this $\eta$ satisfies $\eta X\log^2Z\le1$.

\emph{The truncated sums.} Let $\xi$ be a primitive character of conductor $r$ dividing $N_0$. For $q\in P$ the character $\chi=\xi\cdot(\frac{\cdot}{q})$ is primitive modulo $rq$, and the sum $S_Z(y)$ of Lemma~\ref{gchisplit} for $\chi$ is the sum $S_q(y)$ of Proposition~\ref{window} with $\kappa=\xi$. There are at most $N_0$ characters $\xi$, so Proposition~\ref{window} gives $\lvert S_Z(y)\rvert\le\frac12y^{1/2+\tau}$ for all $\xi$ and all $Y\le y\le X$, for all $q\in P$ outside a set with at most
\[
\frac{16000\,c_0^2\,N_0}{(1-4^{-\tau})^2}\;\lvert P\rvert\;Y^{-2\tau}\log(2Y)\le\tfrac12Y^{-\tau}\lvert P\rvert
\]
elements.

\emph{The boxes.} By Proposition~\ref{boxes}, the number of $q\in P$ for which (i) fails for some $\xi$ is at most $C_3Q^{3/160}(2\log Q)^{14}\le Q^{1/4}\le\frac12Y^{-\tau}\lvert P\rvert$, because $\lvert P\rvert\ge e^z/(5z)$ and $Q^{1/4}=N_0^{1/4}e^{z/4}$.

\emph{The tail.} Let $q$ lie outside both exceptional sets, and let $\chi$ be as above, with conductor $N=rq$ and $\lambda=\log N$, so that $z/2<\lambda\le2z$ and $N\le Q$. Then $e^{800}\le Z\le\exp(\lambda^4)$, and Proposition~\ref{tailprop} gives
\[
\lvert T_Z\rvert\le C_5(2z)^{12}\Big(\frac1X+\frac1{N_0e^z}+e^{-c_1z^2/8}\Big)\le\frac{2^{13}C_5\,Y^{12\varepsilon/5}}X .
\]
Since $12\varepsilon/5<\frac12$, this gives $y\lvert T_Z\rvert\le2^{13}C_5Y^{12\varepsilon/5}\le\frac12y^{1/2+\tau}$ for $Y\le y\le X$, and \eqref{id2} gives (iii). It also gives $\lvert T_Z\rvert\le1$, and the second bound of Proposition~\ref{tailprop} gives
\[
\Big\lvert\frac{L'(1,\chi)}{L(1,\chi)}\Big\rvert\le c_0\big(1+400K\log Y\big)+1\le401\,c_0\,K\log Y ,
\]
which is (ii).
\end{proof}

The proof of Theorem~\ref{stage1} uses the classical zero-free region through (C2), the possible exceptional character through (C3), and the density of zeros through (C5). It does not use a zero-free half-plane for any character. The square root in (iii) comes from the mean square over the family of moduli.

\subsection{The second auxiliary prime}\label{sec:stage2}

The second auxiliary prime $d$ is chosen in one residue class of a large modulus. Let $z\ge2$, let $q>z$ be a prime, and put
\[
m_0=16\,q\prod_{2<p\le z}p .
\]
Let $a_0$ be the residue class modulo $m_0$ with $a_0\equiv7\pmod{16}$ and $a_0\equiv-1$ modulo $q$ and modulo every odd prime $p\le z$. For a prime $d\equiv a_0\pmod{m_0}$ these congruences give $(\frac{-d}p)=1$ for the odd primes $p\le z$ and for $p=q$. They also fix the symbols $(\frac sd)$ for every $s$ whose prime factors are at most $z$, and this is the reason for the sets $\mathcal S(z,s^\ast)$.

\begin{proposition}\label{legendre2}
There is an absolute constant $U_1$ with the following property. Let $z\ge2$, let $R\ge3$, let $q$ be a prime with $q>\max(z,R)$, and let $m_0$ and $a_0$ be as above. Put $\ell=\log(m_0R)$, let $U\ge U_1$ with
\[
\log U\ge16\,\ell\qquad\text{and}\qquad\log U\ge C_0\,\ell\,\log(\ell+2),
\]
and let $P$ be the set of primes $d\equiv a_0\pmod{m_0}$ with $\sqrt U<d\le U$. Put
\[
\eta=4\log U\,e^{-c_2\log U/\ell}+2A_0'\Big(e^{-c_6\log U/\ell}+e^{-\sqrt{c_6\log U}}\Big)+2U^{-3/8},
\]
and suppose that $A_0'\big(e^{-c_6\log U/\ell}+e^{-\sqrt{c_6\log U}}\big)\le\frac14$. Then $\lvert P\rvert\ge U^{3/4}$, and there is an integer $s^\ast\ge1$ without prime factors $p\le z$ such that
\[
\Big\lvert\sum_{d\in P}\Big(\frac sd\Big)\Big\rvert\le\eta\,\lvert P\rvert\qquad\text{for every squarefree }s\text{ with }1<s\le R\text{ and }s\notin\mathcal S(z,s^\ast).
\]
\end{proposition}

\begin{proof}
\emph{The symbols.} Let $d\in P$. Then $d\equiv7\pmod8$, so $(\frac2d)=1$ and $d\equiv3\pmod4$. For an odd prime $p\le z$ the reciprocity law and $d\equiv-1\pmod p$ give $(\frac pd)=(-1)^{(p-1)/2}(\frac{-1}p)=1$. Let $s$ be squarefree with $1<s\le R$, and let $s_b$ be its largest divisor without prime factors $p\le z$. Since $s\le R<q$, the integer $s_b$ is coprime to $m_0$, and
\[
\Big(\frac sd\Big)=\Big(\frac{s_b}d\Big)=(-1)^{(s_b-1)/2}\,\chi(d),\qquad\chi=\Big(\frac{\cdot}{s_b}\Big).
\]
Let $\chi^\ast$ be the character excluded in (C3) at the level $m_0R$, let $s^\ast$ be the largest divisor of its conductor that is coprime to $m_0$, and put $s^\ast=1$ if $\chi^\ast$ does not exist. From now on let $s\notin\mathcal S(z,s^\ast)$, so that $s_b\ne1$ and $s_b\ne s^\ast$, and $\chi$ is a primitive real character modulo $s_b$.

\emph{The progression.} Put $M=m_0s_b$, so that $8$ divides $M$ and $\log M\le\ell$. For $b$ coprime to $s_b$ let $c_b$ be the class modulo $M$ with $c_b\equiv a_0\pmod{m_0}$ and $c_b\equiv b\pmod{s_b}$. Write $\pi_b=\pi(U;M,c_b)$ and $\mathfrak m_b=\mathfrak m(U;M,c_b)$ in the notation of (C7), and put $\mathfrak M=\sum_b\mathfrak m_b$ and $\delta=A_0'(e^{-c_6\log U/\ell}+e^{-\sqrt{c_6\log U}})$. Each $\mathfrak m_b$ is nonnegative, and (C7) gives $\lvert\pi_b-\mathfrak m_b\rvert\le\delta\,\mathfrak m_b$. The primes $d\le U$ in the class $a_0$ modulo $m_0$ that do not divide $s_b$ are counted by $\sum_b\pi_b$, so
\begin{equation}\label{st2a}
\Big\lvert\sum_{d\in P}\chi(d)-\sum_b\chi(b)\,\mathfrak m_b\Big\rvert\le\delta\,\mathfrak M+\sqrt U,\qquad\lvert P\rvert\ge(1-\delta)\,\mathfrak M-\sqrt U .
\end{equation}

\emph{The main terms.} If the pair $(\chi_1,\beta_1)$ of (C7) for the modulus $M$ does not exist, then $\sum_b\chi(b)\mathfrak m_b=0$ and $\mathfrak M=\operatorname{li}(U)/\varphi(m_0)$. If it exists, write $\chi_1=\chi_1'\chi_1''$ with characters $\chi_1'$ modulo $m_0$ and $\chi_1''$ modulo $s_b$, so that $\chi_1(c_b)=\chi_1'(a_0)\chi_1''(b)$. The sum $\sum_b\chi(b)\chi_1''(b)$ equals $\varphi(s_b)$ when $\chi_1''=\chi$ and vanishes otherwise, and $\sum_b\chi_1''(b)$ equals $\varphi(s_b)$ when $\chi_1''$ is principal and vanishes otherwise. Hence
\begin{equation}\label{st2b}
\Big\lvert\sum_b\chi(b)\,\mathfrak m_b\Big\rvert\le\frac{\operatorname{li}(U^{\beta_1})}{\varphi(m_0)}\ \text{ if }\chi_1''=\chi,\qquad\sum_b\chi(b)\,\mathfrak m_b=0\ \text{ otherwise},
\end{equation}
and $\mathfrak M\ge(\operatorname{li}(U)-\operatorname{li}(U^{\beta_1}))/\varphi(m_0)$ in every case.

\emph{A lower bound for $\mathfrak M$.} By (C8), $1-\beta_1\ge c_7/(M^{1/2}\log^2M)$. Since $1-e^{-t}\ge\frac12\min(1,t)$ for $t\ge0$,
\[
\operatorname{li}(U)-\operatorname{li}(U^{\beta_1})\ge\frac{U-U^{\beta_1}}{\log U}\ge\frac U{2\log U}\min\big(1,(1-\beta_1)\log U\big)\ge\frac{c_7\,U}{2M^{1/2}\log^2M\,\log U}
\]
for large $U$. With $\varphi(m_0)\le M\le U^{1/16}$ this gives $\mathfrak M\ge U^{7/8}$ for $U\ge U_1$, and the same bound holds when the pair does not exist. Since $\delta\le\frac14$ and $\sqrt U\le U^{-3/8}\mathfrak M$, the second inequality of \eqref{st2a} gives $\lvert P\rvert\ge\frac12\mathfrak M\ge U^{3/4}$.

\emph{The exceptional term.} Suppose that the pair exists and $\chi_1''=\chi$. The primitive character that induces $\chi_1$ is real, its conductor is at most $M\le m_0R$, and the largest divisor of its conductor that is coprime to $m_0$ is $s_b$. Since $s_b\ne s^\ast$, this character is not $\chi^\ast$, and (C3) gives $\beta_1\le1-c_2/\ell$. Then $\operatorname{li}(U^{\beta_1})\le U^{\beta_1}\le Ue^{-c_2\log U/\ell}$ and $\operatorname{li}(U)\ge U/(2\log U)$, and in this case $\mathfrak M=\operatorname{li}(U)/\varphi(m_0)$, so \eqref{st2b} gives
\[
\Big\lvert\sum_b\chi(b)\,\mathfrak m_b\Big\rvert\le2\log U\,e^{-c_2\log U/\ell}\;\mathfrak M .
\]
This bound also holds in the other cases, where the left side vanishes.

\emph{Conclusion.} The first inequality of \eqref{st2a} and $\mathfrak M\le2\lvert P\rvert$ give
\[
\Big\lvert\sum_{d\in P}\chi(d)\Big\rvert\le\Big(2\log U\,e^{-c_2\log U/\ell}+\delta+U^{-3/8}\Big)\mathfrak M\le\eta\,\lvert P\rvert .\qedhere
\]
\end{proof}

\begin{theorem}\label{stage2}
Let $N_0\ge1$ be an integer, and let $K\ge1$, $0<\varepsilon\le\frac1{52}$ and $0<\tau\le\frac12$. There is $Y_0$ such that the following holds for every $Y\ge Y_0$. Put $z=Y^{\varepsilon/5}$, let $q\equiv3\pmod8$ be a prime with $e^{z/2}<q\le e^z$, let $z'$ satisfy $z\log^2Y\le z'\le K^2z\log^2Y$, and let $P$ be the set of primes $d$ with $e^{z'/2}<d\le e^{z'}$, $d\equiv7\pmod{16}$ and $d\equiv-1$ modulo $q$ and modulo every odd prime $p\le z$. Then $P$ is nonempty, and for all $d\in P$ outside a set of at most $Y^{-\tau}\lvert P\rvert$ elements, every primitive character $\xi$ of conductor dividing $N_0$, every $a\in\{0,1\}$ and the character $\chi=\xi\cdot(\frac{\cdot}{q})^a(\frac{\cdot}{d})$ satisfy
\begin{enumerate}
\item[(i)] $L(s,\chi)\ne0$ for every $s$ with $\operatorname{Re}s>1-\frac1{400}$ and $\lvert\operatorname{Im}s\rvert\le N_0e^{z+z'}$,
\item[(ii)] $\lvert L'(1,\chi)/L(1,\chi)\rvert\le401\,c_0\,K\log Y$,
\item[(iii)] $\lvert G^\ast_\chi(y)\rvert\le y^{1/2+\tau}$ for every $y$ with $Y\le y\le Y^K$.
\end{enumerate}
\end{theorem}

\begin{proof}
Put $X=Y^K$, $Z=X^{400}$, $R=Z^2$, $U=e^{z'}$ and $Q=N_0e^{z+z'}$, and let $m_0$ and $a_0$ be as above, so that $P$ is the set of Proposition~\ref{legendre2}. As in the proof of Theorem~\ref{stage1}, the inequalities between powers of $Y$, powers of $\log Y$ and exponentials that we use hold for $Y\ge Y_0$.

\emph{The Legendre symbols.} By (C1), $\log m_0\le\log16+z+c_0z$, so $\ell=\log(m_0R)\le(2+c_0)z$ and
\[
\frac{\log U}{\ell}\ge\frac{\log^2Y}{2+c_0}.
\]
The hypotheses of Proposition~\ref{legendre2} therefore hold, with $q>e^{z/2}>\max(z,R)$, and its number $\eta$ is at most $\exp(-c\log^2Y)$ for a constant $c>0$, so that $\eta X\log^2Z\le1$. The proposition gives $\lvert P\rvert\ge e^{3z'/4}$ and an integer $s^\ast$ for which the hypothesis of Proposition~\ref{window} holds with this $\eta$ and this $z$. The condition $z\log Y\le\sqrt Y$ of Proposition~\ref{window} holds because $\varepsilon/5<\frac12$, and every $d\in P$ exceeds $Z$.

\emph{The truncated sums.} Let $\xi$ be a primitive character of conductor $r$ dividing $N_0$ and let $a\in\{0,1\}$. For $d\in P$ the character $\chi=\xi\cdot(\frac{\cdot}{q})^a(\frac{\cdot}{d})$ is primitive modulo $rq^ad$, and its sum $S_Z(y)$ is the sum $S_d(y)$ of Proposition~\ref{window} with $\kappa=\xi\cdot(\frac{\cdot}{q})^a$. There are at most $2N_0$ pairs $(\xi,a)$, so $\lvert S_Z(y)\rvert\le\frac12y^{1/2+\tau}$ for all pairs and all $Y\le y\le X$, for all $d\in P$ outside a set with at most
\[
\frac{32000\,c_0^2\,N_0}{(1-4^{-\tau})^2}\;\lvert P\rvert\;Y^{-2\tau}\log(2Y)\le\tfrac12Y^{-\tau}\lvert P\rvert
\]
elements.

\emph{The boxes.} The characters $\chi$ are primitive of conductor at most $Q$, and different $d$ give different characters because $d$ is the largest prime factor of the conductor. As in Proposition~\ref{boxes}, (C5) bounds the number of $d\in P$ for which (i) fails for some pair $(\xi,a)$ by $C_3Q^{3/160}(2\log Q)^{14}\le Q^{1/4}$, and $Q^{1/4}=N_0^{1/4}e^{(z+z')/4}\le\frac12Y^{-\tau}e^{3z'/4}$.

\emph{The tail.} Let $d$ lie outside both exceptional sets and let $\chi$ be as above, with conductor $N=rq^ad$ and $\lambda=\log N$, so that $z'/2<\lambda\le2z'$ and $N\le Q$. Then $e^{800}\le Z\le\exp(\lambda^4)$, and Proposition~\ref{tailprop} gives
\[
\lvert T_Z\rvert\le C_5(2z')^{12}\Big(\frac1X+\frac1Q+e^{-c_1z'^2/8}\Big)\le\frac{2^{13}C_5\,K^{24}\,Y^{12\varepsilon/5}\log^{24}Y}X .
\]
Since $12\varepsilon/5<\frac12$, this gives $y\lvert T_Z\rvert\le\frac12y^{1/2+\tau}$ for $Y\le y\le X$, and \eqref{id2} gives (iii). It also gives $\lvert T_Z\rvert\le1$, and the second bound of Proposition~\ref{tailprop} gives $\lvert L'(1,\chi)/L(1,\chi)\rvert\le c_0(1+400K\log Y)+1\le401\,c_0\,K\log Y$, which is (ii).
\end{proof}

The modulus $m_0$ has $\log m_0$ of order $z$, while $\log d$ has order $z\log^2Y$. This ratio produces the factor $\exp(-c\log^2Y)$ in the proof, and it takes the argument outside the range of (C2) and (C3) into the range of (C7).

\begin{corollary}\label{realbox}
Let $M_1\ge1$ be an odd integer and $N_0=8M_1$, let $z,z'>0$, and let $q\le e^z$ and $d\le e^{z'}$ be distinct primes larger than $N_0$. Suppose that $L(s,\xi\cdot(\frac{\cdot}{q}))\ne0$ for $\operatorname{Re}s>1-\frac1{400}$ and $\lvert\operatorname{Im}s\rvert\le N_0e^z$, and that $L(s,\xi\cdot(\frac{\cdot}{q})^a(\frac{\cdot}{d}))\ne0$ for $\operatorname{Re}s>1-\frac1{400}$ and $\lvert\operatorname{Im}s\rvert\le N_0e^{z+z'}$, for every primitive character $\xi$ of conductor dividing $N_0$ and every $a\in\{0,1\}$. Let $\chi$ be a real primitive Dirichlet character whose conductor $N$ divides $2^jM_1qd$ for some $j\ge0$. If $d$ divides $N$, then $N\le N_0e^{z+z'}$ and $L(s,\chi)\ne0$ for $\operatorname{Re}s>1-\frac1{400}$ and $\lvert\operatorname{Im}s\rvert\le N_0e^{z+z'}$. If $q$ divides $N$ and $d$ does not, the same holds with $N_0e^z$ in place of $N_0e^{z+z'}$.
\end{corollary}

\begin{proof}
A real primitive character is the product of real primitive characters whose conductors are the prime powers that divide its conductor exactly. The conductor at $2$ is $1$, $4$ or $8$, and the conductor at an odd prime $p$ is $p$, with the character $(\frac{\cdot}{p})$~\cite[Ch.~5]{davenport}. Hence $\chi=\xi\cdot(\frac{\cdot}{q})^a(\frac{\cdot}{d})^b$, where $a=1$ when $q$ divides $N$, $b=1$ when $d$ divides $N$, and $\xi$ is a real primitive character whose conductor divides $8M_1=N_0$. The conductor is $N\le N_0q^ad^b$, and the two hypotheses give the two claims.
\end{proof}

\subsection{The host field}\label{sec:field}

The CM field of the argument is cut out by the target character $\psi_0$ and by the primes $q$ and $d$.

\begin{proposition}\label{field}
Let $\psi_0$ be a Dirichlet character of conductor $M_\psi$ and order $o$, let $e\ge0$ be an integer, and let $q\ne d$ be primes with $q\equiv d\equiv3\pmod4$ that do not divide $M_\psi$. Put $N_1=\operatorname{lcm}(2^{e+2},M_\psi)$ and $N=N_1qd$, let $U_1$ be the group of the classes $a\in(\mathbb Z/N_1)^\times$ with $a\equiv\pm1\pmod{2^{e+2}}$ and $\psi_0(a)=1$, and let $U$ be the group of the classes $a\in(\mathbb Z/N)^\times$ whose reduction modulo $N_1$ lies in $U_1$ and that are squares modulo $q$ and modulo $d$. Let $E$ be the subfield of $\mathbb Q(\zeta_N)$ fixed by the automorphisms $\zeta_N\mapsto\zeta_N^a$ with $a\in U$, and put $H=(\mathbb Z/N_1)^\times/U_1$.
\begin{enumerate}
\item[(i)] The field $E$ is Galois over $\mathbb Q$ with group isomorphic to $(\mathbb Z/2)^2\times H$.
\item[(ii)] The field $E$ is a CM field, and it contains $\sqrt{-q}$ and $\sqrt{-d}$.
\item[(iii)] Reduction modulo $2^{e+2}$ induces a surjection of $H$ onto the cyclic group $(\mathbb Z/2^{e+2})^\times/\{\pm1\}$ of order $2^e$. Its kernel has order $c$ for a divisor $c$ of $o$, so that $[E:\mathbb Q]=4n$ with $n=c\,2^e$.
\item[(iv)] The image $A_0$ in $H$ of the kernel of $\psi_0$ is a cyclic group of index $o$ and order $n/o$, the surjection of (iii) is injective on $A_0$, and the characters of $H$ that are trivial on $A_0$ are the powers of $\psi_0$.
\end{enumerate}
\end{proposition}

\begin{proof}
(i) The integers $N_1$, $q$ and $d$ are pairwise coprime, so $(\mathbb Z/N)^\times$ is the product of $(\mathbb Z/N_1)^\times$, $(\mathbb Z/q)^\times$ and $(\mathbb Z/d)^\times$, and $U$ is the product of $U_1$ and the two groups of squares. The Galois group of $\mathbb Q(\zeta_N)$ over $\mathbb Q$ is $(\mathbb Z/N)^\times$, which is abelian, so $E$ is Galois over $\mathbb Q$ with group $(\mathbb Z/N)^\times/U$. This quotient is the product of $H$ with the quotients of $(\mathbb Z/q)^\times$ and $(\mathbb Z/d)^\times$ by their squares, each of order $2$.

(ii) The subfield of $\mathbb Q(\zeta_q)$ fixed by the squares modulo $q$ is the quadratic field of discriminant $(-1)^{(q-1)/2}q$, which is $\mathbb Q(\sqrt{-q})$ because $q\equiv3\pmod4$. Every element of $U$ is a square modulo $q$, so $\sqrt{-q}\in E$, and $\sqrt{-d}\in E$ for the same reason. Complex conjugation is the automorphism with $a=-1$, and it maps $\sqrt{-q}$ to $-\sqrt{-q}$. Since $E$ is abelian over $\mathbb Q$, complex conjugation commutes with every embedding of $E$ into $\mathbb C$, so its fixed field is totally real and $E$ is a totally imaginary quadratic extension of it.

(iii) The group $(\mathbb Z/2^{e+2})^\times$ is the product of $\{\pm1\}$ and the cyclic group of order $2^e$ generated by $5$, so its quotient by $\{\pm1\}$ is cyclic of order $2^e$. Reduction from $(\mathbb Z/N_1)^\times$ to $(\mathbb Z/2^{e+2})^\times$ is surjective and maps $U_1$ into $\{\pm1\}$, which gives the surjection. Let $V$ be the group of the classes $a\in(\mathbb Z/N_1)^\times$ with $a\equiv\pm1\pmod{2^{e+2}}$. The kernel of the surjection is $V/U_1$, and $\psi_0$ maps $V$ to the group of the $o$-th roots of unity with kernel $U_1$. Hence $V/U_1$ is isomorphic to a subgroup of a cyclic group of order $o$, and its order $c$ divides $o$. Then $\lvert H\rvert=c\,2^e$, and (i) gives the degree.

(iv) The kernel of $\psi_0$ in $(\mathbb Z/N_1)^\times$ contains $U_1$ and has index $o$, because $\psi_0$ maps onto the $o$-th roots of unity. So $A_0$ has index $o$ in $H$ and order $n/o$. If a class $a$ in the kernel of $\psi_0$ satisfies $a\equiv\pm1\pmod{2^{e+2}}$, then $a\in U_1$, so the surjection of (iii) is injective on $A_0$, and $A_0$ is cyclic as a subgroup of a cyclic group. A character of $H$ that is trivial on $A_0$ is a character of the quotient of $(\mathbb Z/N_1)^\times$ by the kernel of $\psi_0$, which is cyclic of order $o$, and the $o$ characters of this quotient are the powers of $\psi_0$.
\end{proof}

We identify the Galois group of $E$ with $\mathcal G=(\mathbb Z/2)^2\times H$ as follows. An element $(s_1,s_2,h)\in\{\pm1\}^2\times H$, where $s_1$ is the sign of the action on $\sqrt{-q}$ and $s_2$ the sign of the action on $\sqrt{-d}$, corresponds to $(e_1,e_2,h)\in\mathcal G$ with $(-1)^{e_1}=s_2$ and $(-1)^{e_2}=s_1$. The set $\Phi$ of Section~\ref{sec:excess} is then the CM type of the elements with $s_2=1$. Each character $\psi$ of $\mathcal G$ gives a primitive Dirichlet character $\chi_\psi$ through the Artin map of $\mathbb Q(\zeta_N)$, and we put $\ell_\psi=L'(1,\chi_\psi)/L(1,\chi_\psi)$ in Proposition~\ref{excess}. The character $\chi_1$ of Proposition~\ref{excess} then corresponds to the quadratic character $(\frac{-d}{\cdot})$ of $\mathbb Q(\sqrt{-d})$, and $\lambda_1=\ell_{\chi_1}=L'(1,\chi_{-d})/L(1,\chi_{-d})$ is the number that enters the Chowla--Selberg period $\log\Omega$ of the CM elliptic curve in Part~II.

\subsection{The remaining data}\label{sec:data}

This subsection constructs the sizes, the nodes and the sets $T$ and $T_2$ that enter Theorem~\ref{step5}, and two families of auxiliary primes that Part~II uses.

\paragraph{The scale and the ranks.} A scale consists of a real number $x>1$ and a number $0<\varepsilon<\frac1{52}$. To it we attach
\begin{equation}\label{scale}
L_\varepsilon=x^{\varepsilon},\qquad b=e^{L_\varepsilon},\qquad a=x\,b,\qquad X'=\exp\big(x^{\varepsilon/3}\big),\qquad b'=\frac a{X'},\qquad\delta_1=x^{-\varepsilon/8}.
\end{equation}
The ranks of the chain are integers $n$, $k$, $m$ and $t_r$, where $n=\lvert H\rvert$, $m$ is the number of lifts, the hermitian space of the $j$-th step of the chain has signature $(k+j-1,m-j)$, and the lattices have rank $r+jt_r$ with $r=k+m-1$. Put $f=2n$, $T_O=\lceil x^{377}\rceil$ and
\[
n_W=8\Big\lceil\frac{10\,r}8\Big(f+\frac{f^2(k-1)^3L_\varepsilon^2}{m}\Big)\Big\rceil .
\]

\begin{proposition}\label{ranks}
Let $0<\varepsilon<\frac1{52}$. There is $x_0$ such that for every $x\ge x_0$ and every integer $n$ with $x^{2+8\varepsilon}\le n\le2x^{2+8\varepsilon}$ the following holds.
\begin{enumerate}
\item[(i)] There are even integers $m$ and $k$ and an odd integer $t_r$ with
\begin{gather*}
\sqrt{x\delta_1}\le m\le2\sqrt{x\delta_1},\qquad x^{9/2+10\varepsilon}\le k\le2x^{9/2+10\varepsilon},\qquad r+t_r=8fr+4+T_O\,n_W,\\
2\le m,\qquad m+3\le k,\qquad k+m+4\le t_r,\\
3m<t_r,\qquad m\le2n\le k,\qquad m^2\le k+m(t_r+1),
\end{gather*}
where $r=k+m-1$.
\item[(ii)] For every such $m$ and $k$ and all integers $N_T$ and $N_2$ with $N_T\ge\frac12\delta_1n$ and $x^{-\varepsilon}n\le N_2\le n$, the hypotheses \eqref{sz1} and \eqref{sz2} of Proposition~\ref{sizes} hold.
\end{enumerate}
\end{proposition}

\begin{proof}
Let $x_0$ be so large that $\sqrt{x\delta_1}\ge2$, $x^{9\varepsilon/16}\ge2$, $x^{5\varepsilon}\ge2$, $x^{3\varepsilon/8}\ge26$ and $x^{5/2}\ge4$ for $x\ge x_0$.

(i) Each of the intervals from $\sqrt{x\delta_1}$ to $2\sqrt{x\delta_1}$ and from $x^{9/2+10\varepsilon}$ to $2x^{9/2+10\varepsilon}$ has length at least $2$ and contains an even integer, which gives $m$ and $k$. Define $t_r$ by the equation $r+t_r=8fr+4+T_On_W$. The integer $r=k+m-1$ is odd, and $8fr+4$ and $n_W$ are even, so $t_r$ is odd. Since $f\ge1$ and $T_On_W\ge0$, we get $t_r\ge7r+4\ge k+m+4$. Since $r\ge m$, this gives $t_r\ge7m+4$, so that $3m<t_r$ and $m^2\le m(t_r+1)$. The inequalities $2\le m$ and $m+3\le k$ are clear, $m\le2\sqrt x\le2n$, and $2n\le4x^{2+8\varepsilon}\le x^{9/2+10\varepsilon}\le k$.

(ii) The inequalities $a\le2xb$ and $N_T\ge\frac12x^{-\varepsilon/8}n$ hold because $a=xb$ and $\delta_1=x^{-\varepsilon/8}$. Since $\sqrt{x\delta_1}=x^{1/2-\varepsilon/16}$, we have $m\ge x^{1/2-\varepsilon}$ and $m\le2x^{1/2-\varepsilon/16}\le x^{1/2+\varepsilon/2}$. Then $xm\le2x^{3/2}\le k$, and
\[
k\le2x^{9/2+10\varepsilon}\le x^{9/2+15\varepsilon}=\big(x^{2+8\varepsilon}\big)^2x^{1/2-\varepsilon}\le n^2x^{1/2-\varepsilon}.\qedhere
\]
\end{proof}

\paragraph{The nodes.} Each node has the form $b_1+(j+\frac12)h_0$ with $0\le j<m$, where $h_0=2\lceil x^{c_h}\rceil$ for a constant $c_h\ge0$ and the integer $b_1$ depends on the place. Let $W$ be the number of roots of unity in $E$.

\begin{proposition}\label{nodes}
Let $0<\varepsilon<\frac1{52}$ and $c_h\ge0$. There is $x_0$ such that the following holds for every $x\ge x_0$. Let $b$, $a$, $X'$ and $b'$ be as in \eqref{scale}. Let $n$, $m$, $k$ and $W$ be positive integers with $n\le2x^{9/4}$, $m\le2\sqrt x$ and $W\le32n^2$, put $f=2n$ and $h_0=2\lceil x^{c_h}\rceil$, let $b_1$ be an integer, and put $w_j=b_1+(j+\frac12)h_0$ for $0\le j<m$.
\begin{enumerate}
\item[(i)] If $\lceil b\rceil\le b_1<\lceil b\rceil+W$, then $b\le w_j\le b+\eta'$ with $\eta'=W+mh_0$, and $\eta'\le b$ and $6f(x+1)\eta'\le b/x$.
\item[(ii)] If $b_1=\lceil b'/(xn)\rceil$, or if $b'\le b_1-\lfloor b'/x\rfloor\le\frac54b'$, then $0<w_j\le2b'$.
\item[(iii)] If $a\le a_t\le a+a/(kx^{13})$, then $\lvert a_t/b-x\rvert\le1$.
\item[(iv)] $b'X'=xb$, and $\omega=\exp(-x^{\varepsilon/8})$ satisfies $2f(x+1)\omega\le1/x$.
\end{enumerate}
\end{proposition}

\begin{proof}
The numbers $n$, $m$, $h_0$ and $W$ are bounded by powers of $x$, while $b$, $b'$ and $1/\omega$ grow faster than every power of $x$. Let $x_0$ be so large that the inequalities between these quantities used below hold for $x\ge x_0$.

(i) We have $w_j\ge\lceil b\rceil\ge b$ and $w_j\le\lceil b\rceil+W-1+(m-\frac12)h_0\le b+W+mh_0$. The number $12n(x+1)x\,\eta'$ is bounded by a power of $x$, so it is at most $b$.

(ii) In the first case $w_j\le b'/(xn)+1+mh_0\le2b'$. In the second case $b_1\le\frac54b'+b'/x$, so $w_j\le b'(\frac54+\frac1x)+mh_0\le2b'$. In both cases $w_j>0$.

(iii) The hypothesis gives $0\le a_t/b-x\le x/(kx^{13})\le1$.

(iv) The first claim is the definition of $b'$, and $4n(x+1)x\le\exp(x^{\varepsilon/8})$ for large $x$.
\end{proof}

The bound $W\le32n^2$ holds because the field of $W$-th roots of unity lies in $E$, so that $\varphi(W)\le4n$, while $\varphi(W)\ge\sqrt{W/2}$. The three cases of Proposition~\ref{nodes} are the thick places, the thin places and the places attached to the elements of $T$. Under the identification of Section~\ref{sec:field}, the places attached to $T$ are the elements $(1,1,t)$ with $t\in T$, whose images in $\mathcal G$ are $(0,0,t)$. The thick places are the elements $(-1,1,s\,a_0+u_0)$, where $a_0$ is a generator of $A_0$ and $s$ runs over the set of integers $T_2$ of Proposition~\ref{behrend} below. The image of such an element is $(0,1,s\,a_0+u_0)$, so the thick places are the classes that are thick with respect to $u_0$ for the subset $\{s\,a_0:s\in T_2\}$ of $A_0$. The map $s\mapsto s\,a_0$ is injective on $T_2$ because every element of $T_2$ is smaller than the order of $a_0$, so this subset has $\lvert T_2\rvert$ elements.

\paragraph{The set $T_2$.}

\begin{proposition}\label{behrend}
For all integers $n_0\ge1$ and $L\ge2$ there is a set $T_2$ of nonnegative integers with the following properties. Every $s\in T_2$ satisfies $4Ls<n_0$. There are no distinct $s_1,s_2,s_3\in T_2$ and integers $i\ne j$ with $0<\lvert i\rvert,\lvert j\rvert\le L$ and $j(s_2-s_1)=i(s_3-s_1)$. Finally
\[
\lvert T_2\rvert\ge\frac{n_0}L\exp\Big(-36\sqrt{\log n_0\,\log L}\Big).
\]
\end{proposition}

\begin{proof}
Put $N=n_0/(4L)$ and $\lambda=\log(8L)$, so that $2<\lambda\le4\log L$. The set $T_2=\{0\}$ has the first two properties, and it has the third when the right side of the claim is at most $1$. This is the case when $n_0<L$. It is also the case when $n_0\ge L$ and $\log N\le8\lambda$, because then $\log(n_0/L)\le9\lambda\le36\log L\le36\sqrt{\log n_0\log L}$.

Let $\log N>8\lambda$, so that $n_0>L$, and put $D=\lfloor\sqrt{2\log N/\lambda}\rfloor$, so that $D\ge4$ and $\frac34\sqrt{2\log N/\lambda}\le D\le\sqrt{2\log N/\lambda}$. Then $\log N/D\ge\sqrt{\lambda\log N/2}\ge2\lambda$, so $N^{1/D}\ge(8L)^2$. Put $B=\lfloor N^{1/D}\rfloor$ and $g=\lfloor B/(2L)\rfloor$, so that $B^D\le N$, $B\ge\frac12N^{1/D}$ and $g\ge B/(4L)\ge N^{1/D}/(8L)$. For a vector $v=(v_0,\dots,v_{D-1})$ of integers with $0\le v_\nu<g$ put $s(v)=\sum_\nu v_\nu B^\nu$. The map $s$ is injective, and $4L\,s(v)<4LB^D\le n_0$. The number $\sum_\nu v_\nu^2$ takes at most $Dg^2$ values, so there is an integer $\varrho$ such that the set $V$ of the vectors $v$ with $\sum_\nu v_\nu^2=\varrho$ has at least $g^{D-2}/D$ elements. Let $T_2=s(V)$.

Suppose that $s_k=s(v^{(k)})$ for $k=1,2,3$ are distinct elements of $T_2$ and that $j(s_2-s_1)=i(s_3-s_1)$ with $0<\lvert i\rvert,\lvert j\rvert\le L$. The integers $c_\nu=j(v^{(2)}_\nu-v^{(1)}_\nu)-i(v^{(3)}_\nu-v^{(1)}_\nu)$ satisfy $\lvert c_\nu\rvert\le2L(g-1)<B$ and $\sum_\nu c_\nu B^\nu=0$, and an induction on $\nu$ gives $c_\nu=0$ for every $\nu$. Hence $v^{(3)}-v^{(1)}=\frac ji(v^{(2)}-v^{(1)})$, so the three distinct points $v^{(1)},v^{(2)},v^{(3)}$ lie on one line and on the sphere $\sum_\nu v_\nu^2=\varrho$. A line meets a sphere in at most two points, which is a contradiction.

For the size, $g^{D-2}/D\ge N^{1-2/D}(8L)^{-(D-2)}/D\ge N\exp\big(-\frac2D\log N-D\lambda-\log D\big)$. The bounds for $D$ give $\frac2D\log N\le\frac8{3\sqrt2}\sqrt{\lambda\log N}$, $D\lambda\le\sqrt{2\lambda\log N}$ and $\log D\le\sqrt{\lambda\log N}$, so the exponent is at least $-5\sqrt{\lambda\log N}\ge-10\sqrt{\log n_0\log L}$. Since $\sqrt{\log n_0\log L}\ge\log L\ge\log2$, we get $\lvert T_2\rvert\ge\frac{n_0}{4L}e^{-10\sqrt{\log n_0\log L}}\ge\frac{n_0}Le^{-36\sqrt{\log n_0\log L}}$.
\end{proof}

Proposition~\ref{behrend} is the construction of Behrend~\cite{behrend} with digits below $B/(2L)$ in place of digits below $B/2$. We apply it with $n_0=n/o$, where the order $o$ of $\psi_0$ is fixed, and with $L=L_w$ of order $\log x$. Then $\sqrt{\log n_0\log L_w}$ has smaller order than $\log x$, and the set has at least $x^{-\varepsilon}n$ elements for large $x$.

\paragraph{The set $T$.} A set $T\subset H$ indexes the weights $a_t$, and it must have two properties. The first is a condition on its shares on the cosets of $A_0$, which produces the positive constant $c$ in \eqref{equi}.

\begin{lemma}\label{shares}
Let $H$ be a finite abelian group, $A_0$ a subgroup of index $o$, $T$ a nonempty subset of $H$ with $N_T$ elements, and $\delta\ge0$. Put $\lambda_\gamma=4/(3o+1)$ for $\gamma\in A_0$ and $\lambda_\gamma=3/(3o+1)$ for $\gamma\in H\setminus A_0$, and suppose that for every $\gamma\in H$
\[
\Big\lvert\#\{t\in T:t-\gamma\in A_0\}-\lambda_\gamma N_T\Big\rvert\le\delta\,N_T .
\]
Then every nontrivial character $\chi$ of $H$ that is trivial on $A_0$ satisfies
\[
\Big\lvert\sum_{t\in T}\chi(t)-\frac{N_T}{3o+1}\Big\rvert\le o\,\delta\,N_T .
\]
\end{lemma}

\begin{proof}
Let $\gamma_1=0,\gamma_2,\dots,\gamma_o$ represent the cosets of $A_0$, and let $n_i$ be the number of $t\in T$ with $t-\gamma_i\in A_0$. The character $\chi$ is constant on each coset, so $\sum_{t\in T}\chi(t)=\sum_i\chi(\gamma_i)\,n_i$. It induces a nontrivial character of $H/A_0$, so $\sum_i\chi(\gamma_i)=0$ and
\[
\sum_{i=1}^o\chi(\gamma_i)\,\lambda_{\gamma_i}=\frac{4+3\sum_{i\ge2}\chi(\gamma_i)}{3o+1}=\frac{4-3}{3o+1}=\frac1{3o+1}.
\]
Therefore $\sum_{t\in T}\chi(t)-N_T/(3o+1)=\sum_i\chi(\gamma_i)(n_i-\lambda_{\gamma_i}N_T)$, and each of the $o$ terms of the last sum has absolute value at most $\delta N_T$.
\end{proof}

With $\delta=x^{-1+\varepsilon}$, Lemma~\ref{shares} gives \eqref{equi} with $c=1/(3o+1)$ and $\nu=o\,x^{-1+\varepsilon}$. The second property of $T$ is a regularity condition along the orbits of Frobenius elements, which the gain at inert primes in Section~\ref{sec:gain} uses.

\begin{definition}\label{def:goodorbit}
Let $h\in H$ have order $n_h$, let $t\in H$, and let $P_t\subset\mathbb Z/n_h\mathbb Z$ be the set of the residues $i$ with $t_0+ih\in T$, where $t_0$ is a fixed base point of the orbit $t+\mathbb Zh$. Let $K_\varepsilon=\lceil8(4+8\varepsilon)/\varepsilon\rceil$, $r_0=K_\varepsilon+1$, $u=\lceil1/\delta_1\rceil$, $M_0=\lceil(5+8\varepsilon)\,10^{6r_0+1}\log x\rceil$ and $G'=(M_0+5)(6r_0+1)u$, and let $L_2\ge1$ be an integer. The orbit of $t$ under $h$ is \emph{good for $T$} when the following hold.
\begin{enumerate}
\item[(i)] Every window of $L_2$ consecutive residues contains at most $K_\varepsilon$ elements of $P_t$.
\item[(ii)] Every window of $G'$ consecutive residues contains at most $\lceil\max(6e\,G'\delta_1,(5+8\varepsilon)\log_2x)\rceil+1$ elements of $P_t$.
\item[(iii)] Starting at $0$, take the $\lfloor n_h/((6r_0+1)u)\rfloor$ full tiles of length $(6r_0+1)u$ in $[0,n_h)$. Divide each into $6r_0+1$ cells of length $u$. Among every $M_0+1$ cyclically consecutive tiles there is one in which the cells in odd position are disjoint from $P_t$ and each cell in even position contains exactly one element of $P_t$, and that element is an even integer.
\end{enumerate}
\end{definition}

\begin{proposition}\label{tops}
Let $0<\varepsilon<\frac1{52}$. There are an absolute constant $C_b$ and, for every integer $o\ge1$, a number $x_0$ with the following property. Let $x\ge x_0$, let $H$ and $A_0$ be as in Proposition~\ref{field}, with $A_0$ of index $o$, with $x^{2+8\varepsilon}\le n\le2x^{2+8\varepsilon}$, let $\mathcal P$ be a finite set of primes, and for $p\in\mathcal P$ let $h_p\in H$ have order at least $x^{\varepsilon/4}$ and let $L_{2,p}\le400\log x$ be a positive integer. Then there is a set $T\subset H$ with $\frac12\delta_1n\le\lvert T\rvert\le2\delta_1n$ that satisfies the hypothesis of Lemma~\ref{shares} with $\delta=x^{-1+\varepsilon}$ and such that the number $\beta_p(T)$ of the $t\in T$ whose orbit under $h_p$ is not good for $T$, with $L_2=L_{2,p}$, satisfies
\[
\sum_{p\in\mathcal P}\frac{\log p}p\,\beta_p(T)\le4C_b\,x^{-2}\,\delta_1n\sum_{p\in\mathcal P}\frac{\log p}p .
\]
\end{proposition}

\begin{proof}
Choose each element of $H$ independently, with probability $\delta_1$
on $A_0$ and $3\delta_1/4$ on its other cosets. For the count $X_i$ in
coset $i$, write $\mu_i=\mathbb E X_i$ and $\mu=\mathbb E|T|$.
The cosets all have size $n/o$, so
\[
 \mu=\frac{3o+1}{4o}\,\delta_1 n,\qquad
 \mu_i=\lambda_i\mu,\qquad \operatorname{Var}(X_i)\le\mu_i,
\]
where $\lambda_i$ is the weight of that coset in Lemma~\ref{shares}.

\emph{Size and shares.} Put $\delta=x^{-1+\varepsilon}$ and
$\eta=\delta\delta_1 n/(2(o+1))$. By Chebyshev's inequality the probability
that any $|X_i-\mu_i|>\eta$ is at most $\mu/\eta^2$, which tends to zero
because $\delta^2\delta_1 n\ge x^{(10-1/8)\varepsilon}$ and $o$ is fixed.
Otherwise $|N_T-\mu|\le o\eta$. Since
$3\delta_1 n/4\le\mu\le\delta_1 n$ and $o\eta/(\delta_1 n)\to0$ as
$x\to\infty$, we have $\delta_1 n/2\le N_T\le2\delta_1 n$. Also
$|X_i-\lambda_iN_T|\le(1+o)\eta\le\delta N_T$ in every coset.

\emph{One orbit.} Fix $p,t$ and condition on $t\in T$. On its orbit the
residue $i_t$ of $t$ is now selected. All other residues are still chosen
independently with probabilities between $3\delta_1/4$ and $\delta_1$.
Write $K=K_\varepsilon$ and $L_2=L_{2,p}$. For large $x$ both $L_2$
and $G'$ are smaller than $n_h$, since $G'=O_\varepsilon(x^{\varepsilon/8}\log x)$
and $n_h\ge x^{\varepsilon/4}$. Thus no window repeats a residue.

For (i), a failing window needs $K+1$ successes, or $K$ further successes
if it contains $i_t$. There are at most $n_h$ windows of the first kind
and $L_2$ of the second. The union bound over choices of residues gives
\[
 \Pr(\text{(i) fails}\mid t\in T)
 \le n_h(L_2\delta_1)^{K+1}+L_2(L_2\delta_1)^K\le x^{-2}
\]
for large $x$, using $\delta_1^K\le x^{-4-8\varepsilon}$,
$n_h\le2x^{2+8\varepsilon}$ and $L_2\le400\log x$.

For (ii), put $N=1+\lceil\max(6eG'\delta_1,(5+8\varepsilon)\log_2x)\rceil$.
A failing window contains at least $N$ successes other than $i_t$.
Their probability in a given window is at most
$\binom{G'}{N}\delta_1^N\le(eG'\delta_1/N)^N\le2^{-N}$.
Summing over the $n_h$ windows gives at most $x^{-3}$.

For (iii), consider a tile avoiding $i_t$. A cell of length
$u=\lceil1/\delta_1\rceil$ is empty with probability at least
$(1-\delta_1)^u$. The probability that its only selected residue is
even is at least
\[
 \lfloor u/2\rfloor(3\delta_1/4)(1-\delta_1)^{u-1}.
\]
These lower bounds tend to $1/e$ and $3/(8e)$, respectively, both
greater than $1/10$. The cells are disjoint, so the prescribed tile
pattern has probability at least $\rho=10^{-(6r_0+1)}$. For large $x$ the
orbit contains at least $M_0+1$ full tiles, since $n_h\ge x^{\varepsilon/4}$
and $(M_0+1)(6r_0+1)u=O_\varepsilon(x^{\varepsilon/8}\log x)$.
Any run in (iii) therefore contains at least $M_0$ distinct tiles
avoiding $i_t$. By independence the run fails with probability at most
$e^{-\rho M_0}\le x^{-5-8\varepsilon}$. There are at most $n_h$ runs.

The three failure probabilities sum to at most $C_b x^{-2}$ with
$C_b=2$ for large $x$, uniformly in $p,t$. Multiplying by $\Pr(t\in T)$
and summing over $t$ gives $\mathbb E\beta_p(T)\le C_b x^{-2}\mu$.
Since $\mu\le\delta_1 n$, Markov's inequality gives the claimed weighted
bound with probability at least $3/4$ (also when all weights vanish).
The size and share conditions hold with probability tending to one, so
some $T$ satisfies both requirements.
\end{proof}

\paragraph{Two families of auxiliary primes.} Part~II uses three primes congruent to $1$ modulo $Q_Sqd$ and two primes $g_0\ne g_1$ that split in $\mathbb Q(\sqrt{-d})$ into principal primes with a generator congruent to $1$ modulo $Q_Sqd$, where $Q_S$ is the part of the conductor of $E$ prime to $qd$. Linnik's theorem~\cite{linnik} gives a constant $A_L\ge1$ such that for every integer $M\ge2$ there is a prime $p\equiv1\pmod M$ with $p\le M^{A_L}$.

\begin{lemma}\label{three}
Let $A\ge1$ be a number such that for every integer $M\ge2$ there is a prime $p\equiv1\pmod M$ with $p\le M^A$. Then for every integer $M\ge2$ there are three distinct primes $p_1,p_2,p_3\equiv1\pmod M$ with $p_i\le M^{A(1+A)^2}$.
\end{lemma}

\begin{proof}
Choose a prime $p_1\equiv1\pmod M$ with $p_1\le M^A$, then a prime $p_2\equiv1\pmod{Mp_1}$ with $p_2\le(Mp_1)^A\le M^{A(1+A)}$, and then a prime $p_3\equiv1\pmod{Mp_1p_2}$ with $p_3\le(Mp_1p_2)^A\le M^{A(1+A)^2}$. Then $p_2>Mp_1$ and $p_3>Mp_1p_2$, so the three primes are distinct, and each is congruent to $1$ modulo $M$.
\end{proof}

For the primes $g_0$ and $g_1$ we use a prime number theorem for a ray class of an imaginary quadratic field, uniform in the conductor.
\begin{enumerate}
\item[(C9)] There are constants $c,c'>0$ and $A_1\ge1$ such that for every prime $d\equiv3\pmod4$, every integer $M\ge1$ and every $X\ge2$ with $\log(4dM^2)\le c\sqrt{\log X}$, and for every $e^{-c\sqrt{\log X}}\le h\le1$, the number of primes $g$ with $X<g\le X(1+h)$ of the form $g=\pi\bar\pi$ with $\pi\in\mathbb Z[\frac{1+\sqrt{-d}}2]$ and $\pi\equiv1\pmod M$ is at least $c'hX/((4dM^2)^{A_1}\log X)$.
\end{enumerate}
Statement (C9) follows from Theorem~1.4 of Thorner and Zaman~\cite{thornerzaman19}, applied to the ray class field $L$ of $K=\mathbb Q(\sqrt{-d})$ modulo $M\mathcal O_K$ with $F=\mathbb Q$, $H=\mathrm{Gal}(L/K)$ and $C=\{1\}$, for which $D_K\mathcal Q n_K^{n_K}\le4dM^2$. A prime that splits completely in $L$ has the form $\pi\bar\pi$ with $\pi\equiv1\pmod M$. The possible exceptional zero is a zero of $\zeta_{K'}$ for a quadratic extension $K'/K$ with discriminant at most $d^2M^2$, so Stark's bound~\cite[Theorem~1$'$]{stark74}, in the form~(1-4) of~\cite{thornerzaman19}, gives $1-\beta_1\gg(4dM^2)^{-1}$. With the ray class number at most a constant times $dM^2$ this gives (C9) with $A_1=2$, for every prime $d\equiv3\pmod4$ and every $M\ge1$. With $X=(kb)^8$ the estimate gives at least $X^{9/10}$ such primes between $X$ and $2X$.

\begin{lemma}\label{disjoint}
For integers $g\ge2$ and $N\ge1$ let $B_N(g)$ be the set that consists of $g$ and of the primes $l$ for which there is an integer $1\le c\le N$ with $l^{4v}>g^{c}$, where $l^{v}$ is the exact power of $l$ that divides $g^{c}-1$. Let $Y\ge2$ be a real number, let $N\ge1$ be an integer, let $\mathcal P$ be a finite set of primes $g$ with $Y\le g\le2Y$, and suppose that
\[
\lvert\mathcal P\rvert>(3N+1)\big(1+2N^2(Y^{3/4}+1)\big).
\]
Then for every $g_1\in\mathcal P$ there is $g_0\in\mathcal P$ with $B_N(g_0)\cap B_N(g_1)=\emptyset$, and such a $g_0$ differs from $g_1$.
\end{lemma}

\begin{proof}
Fix $g\ge2$ and $1\le c\le N$. If four distinct primes $l_1,\dots,l_4$ satisfied $l_i^{4v_i}>g^{c}$ with $l_i^{v_i}$ the exact power dividing $g^{c}-1$, then the product of the four numbers $l_i^{v_i}$ would exceed $g^{c}$ and divide the positive integer $g^{c}-1$. Hence at most three primes belong to each $c$, and $\lvert B_N(g)\rvert\le3N+1$.

Fix a prime $l$ and $1\le c\le N$, and let $w\ge1$ be the least integer with $l^{w}>Y^{c/4}$. Let $g$ be an integer with $Y\le g\le2Y$ such that $l^{4v}>g^{c}$ for the exact power $l^{v}$ dividing $g^{c}-1$. Then $l^{v}>g^{c/4}\ge Y^{c/4}$, so $v\ge w$ and $g^{c}\equiv1\pmod{l^{w}}$. The group of units modulo $l^{w}$ is cyclic when $l$ is odd, and it is the product of a group of order at most $2$ and a cyclic group when $l=2$, so the congruence $y^{c}\equiv1\pmod{l^{w}}$ has at most $2c$ solutions modulo $l^{w}$. Each residue class modulo $l^{w}$ contains at most $Y/l^{w}+1$ integers between $Y$ and $2Y$, and $Y/l^{w}<Y^{1-c/4}\le Y^{3/4}$. So at most $2c\,(Y^{3/4}+1)$ integers $g$ between $Y$ and $2Y$ satisfy the condition for this $c$, and the sum over $c\le N$ is at most $N(N+1)(Y^{3/4}+1)\le2N^2(Y^{3/4}+1)$. Together with the possibility $g=l$, the prime $l$ lies in $B_N(g)$ for at most $1+2N^2(Y^{3/4}+1)$ integers $g$ between $Y$ and $2Y$.

Let $g_1\in\mathcal P$. Every element of $B_N(g_1)$ is a prime, so the two counts bound the number of $g\in\mathcal P$ for which $B_N(g)$ meets $B_N(g_1)$ by $(3N+1)\big(1+2N^2(Y^{3/4}+1)\big)$, which is less than $\lvert\mathcal P\rvert$. This gives $g_0$, and $g_0\ne g_1$ because $g_1\in B_N(g_1)$.
\end{proof}

We apply Lemma~\ref{disjoint} with $N=r$, $Y=(kb)^8$ and $\mathcal P$ the set of primes given by (C9). The rank $r$ is at most a power of $x$, while $Y^{9/10-3/4}\ge\exp(x^{\varepsilon})$, so $\lvert\mathcal P\rvert\ge Y^{9/10}$ exceeds the bound of the lemma for large $x$.

\section{The arithmetic theorem and the proof of the Riemann hypothesis}\label{sec:reduction}

This section collects the data of Section~\ref{sec:auxiliary} into one definition, states the arithmetic theorem that Part~II proves, and deduces the Riemann hypothesis and the generalized Riemann hypothesis for Dirichlet $L$-functions from it.

\subsection{Admissible data}

Fix a primitive Dirichlet character $\psi_0$ of conductor $M_\psi$ and order $o$, and fix $0<\varepsilon<\frac1{52}$. Put $N_0=8M_\psi$.

\begin{definition}\label{def:admissible}
Let $x\ge3$, and let $Y$ be a real number with $Y\le x\le Y^{3/2}$. \emph{Admissible data at the scale $x$ for the window $[Y,Y^{3/2}]$} consist of the following objects.
\begin{enumerate}
\item[(D1)] Primes $q$ and $d$ that satisfy the conclusions of Theorems~\ref{stage1} and~\ref{stage2} with $N_0$, $K=\frac32$, $\varepsilon$, $\tau=\varepsilon$ and $z'=z\log^2Y$. An integer $e\ge0$ with $x^{2+8\varepsilon}\le c\,2^e\le2x^{2+8\varepsilon}$, where $c$ is the integer of Proposition~\ref{field}(iii). The field $E$, the groups $H$ and $A_0$ and the integer $n=\lvert H\rvert$ of Proposition~\ref{field} for $\psi_0$, $e$, $q$ and $d$, with the identification of $\mathrm{Gal}(E/\mathbb Q)$ with $\mathcal G=(\mathbb Z/2)^2\times H$ of Section~\ref{sec:field}.
\item[(D2)] Integers $m$, $k$ and $t_r$ as in Proposition~\ref{ranks}(i).
\item[(D3)] A set $T\subset H$ with the properties of Proposition~\ref{tops}, for the set $\mathcal P$ and the elements $h_p$ specified in Section~\ref{sec:gain}.
\item[(D4)] A set $T_2\subset A_0$ of the form $\{s\,a_0:s\in T_2'\}$, where $a_0$ generates $A_0$ and $T_2'$ is a set given by Proposition~\ref{behrend} with $n_0=n/o$ and $L=L_w$, with $\lvert T_2\rvert\ge x^{-\varepsilon}n$.
\item[(D5)] Two primes $g_0\ne g_1$ with $(kb)^8\le g_i\le2(kb)^8$, each of the form $g_i=\pi_i\bar\pi_i$ with $\pi_i\in\mathbb Z[\frac{1+\sqrt{-d}}2]$ and $\pi_i\equiv1$ modulo $Q_Sqd$, with $B_r(g_0)\cap B_r(g_1)=\emptyset$ in the notation of Lemma~\ref{disjoint}. Three distinct primes congruent to $1$ modulo $15\,Q_Sqd$ and at most $(15\,Q_Sqd)^{A_L(1+A_L)^2}$.
\end{enumerate}
\end{definition}


\begin{proposition}\label{dataexist}
There is $Y_0$ such that for every $Y\ge Y_0$ and every $x$ with $Y\le x\le Y^{3/2}$, admissible data at the scale $x$ for the window $[Y,Y^{3/2}]$ exist.
\end{proposition}

\begin{proof}
Theorems~\ref{stage1} and~\ref{stage2} with $K=\frac32$ and $\tau=\varepsilon$ give the primes $q$ and $d$ for every large $Y$. For large $Y$ the primes $q$ and $d$ exceed $M_\psi$, and $q\equiv d\equiv3\pmod4$. The integer $c$ of Proposition~\ref{field}(iii) divides $o$ and does not depend on $e$ once $2^{e+2}$ is divisible by the power of $2$ in $M_\psi$, and the interval $[x^{2+8\varepsilon},2x^{2+8\varepsilon}]$ has ratio $2$, so it contains a number $c\,2^e$. Proposition~\ref{field} gives (D1). Propositions~\ref{ranks}, \ref{tops} and~\ref{behrend} give (D2), (D3) and (D4); for (D4) we use $n_0=n/o$, $L_w$ of order $\log x$, and the remark after Proposition~\ref{behrend}. Statement (C9) and Lemma~\ref{disjoint} give the primes $g_0$ and $g_1$, and Linnik's theorem with Lemma~\ref{three} for $M=15\,Q_Sqd$ gives the three primes of (D5).
\end{proof}

\subsection{The arithmetic theorem}

For admissible data let $\Phi$, $f=2n$ and $p_t$ be as in Section~\ref{sec:excess}, let $G_\psi$ be the functions of Proposition~\ref{excess} with $\ell_\psi=L'(1,\chi_\psi)/L(1,\chi_\psi)$, and let $\mathrm{De}$ be the excess \eqref{De}. Let $b$, $a$, $X'$ and $b'$ be as in \eqref{scale}, and let $\mathcal B$ be the number \eqref{calB} with $N_T=\lvert T\rvert$ and $N_2=\lvert T_2\rvert$.

Section~\ref{sec:proofA} refines Definition~\ref{def:admissible} by a base scale $Y_0\le Y$, at which the primes $q$ and $d$ are chosen, and by a condition (D6) on Tate--Shafarevich groups of $E$ (Definition~\ref{r6:R6:def:adm}). Theorem~A asserts the existence of such data at every scale of a window, for windows $[Y,Y^{3/2}]$ with $Y$ arbitrarily large.

\begin{theoremA}
There is an absolute constant $C\ge0$ with the following property. For every primitive Dirichlet character $\psi_0$ there are $0<\varepsilon_1\le\frac1{52}$, and for every $0<\varepsilon<\varepsilon_1$ a number $J\ge2$, such that for every $y$ there are real numbers $Y_0\le Y$ with $Y\ge y$ and the following property. For every $x$ with $Y\le x\le Y^{3/2}$ there are admissible data at the scale $x$ for the window $[Y,Y^{3/2}]$ with base $Y_0$, weights $a_t$ for $t\in T$ with $a\le a_t\le a+a/(kx^{13})$, and for every $u_0\in A_0$ and $1\le j\le m$ functions $w_{u_0,j}\colon\Phi\to(0,\infty)$, such that the following hold.
\begin{enumerate}
\item[(i)] For every $u_0\in A_0$, every $j$ and every $\rho\in\Phi$, the node $w_{u_0,j}(\rho)$ satisfies the bounds of Proposition~\ref{nodes}(i) when $\rho$ is thick with respect to $u_0$ for the set $T_2$, and the bounds of Proposition~\ref{nodes}(ii) otherwise.
\item[(ii)] For every $u_0\in A_0$,
\[
\sum_{t\in T}\sum_{j=1}^m\operatorname{Re}\mathrm{De}(a_t,t,w_{u_0,j})\le x^{C\varepsilon}\,\mathcal B .
\]
\end{enumerate}
One can take $C=5$.
\end{theoremA}

Part~II proves Theorem~A from the hypotheses that Subsection~\ref{sec:frontier} lists. For each $u_0\in A_0$ it constructs hermitian lattices and period maps whose slopes, heights and local terms satisfy the inequalities \eqref{A1} to \eqref{A7} of Proposition~\ref{acct}, with $\lambda_1=\ell_{\chi_1}$, with the sum $S$ of Corollary~\ref{acctexcess}, and with
\[
e_1+e_2+e_3+e_4+e_5+Ae_6+e_7\le x^{C\varepsilon}\,\mathcal B .
\]
Corollary~\ref{acctexcess} then gives (ii). Section~\ref{sec:proofA} gives the assembly.

\subsection{The Riemann hypothesis}

\begin{theorem}\label{mainthm}
Assume Theorem~A. Then for every primitive Dirichlet character $\psi_0$ the function $L(s,\psi_0)$ does not vanish in the half-plane $\operatorname{Re}s>\frac12$. In particular $\zeta(s)\ne0$ for $\operatorname{Re}s>\frac12$.
\end{theorem}

\begin{proof}
Fix $\psi_0$ of conductor $M_\psi$ and order $o$, let $\varepsilon_1$ and $C$ be as in Theorem~A, fix $0<\varepsilon<\varepsilon_1$, and let $J$ be the number of Theorem~A. Let $Y_0\le Y$ be numbers given by Theorem~A with $Y$ large, let $Y\le x\le Y^{3/2}$, and choose admissible data at the scale $x$ for the window $[Y,Y^{3/2}]$ with base $Y_0$, weights and nodes by Theorem~A.

\emph{The characters.} The characters of $\mathcal G$ that are trivial on $\{0\}\times\{0\}\times A_0$ form the set $\Psi_0$, of order $4o$. By Proposition~\ref{field}(iv) and the identification of Section~\ref{sec:field}, the Dirichlet characters $\chi_\psi$ for $\psi\in\Psi_0$ are the primitive characters that induce $\psi_0^i(\frac{\cdot}{q})^a(\frac{\cdot}{d})^{a'}$ with $0\le i<o$ and $a,a'\in\{0,1\}$. The class $(1,0,0)$ acts trivially on $\sqrt{-q}$ and on $H$, and by $-1$ on $\sqrt{-d}$, so $\psi(1,0,0)=1$ exactly when $a'=0$, and $\psi$ is trivial on $\{0\}\times\{0\}\times H$ exactly when $i=0$. Let $K$ be the set of the $\psi\in\Psi_0$ with $a=a'=0$, so that $\chi_\psi$ runs over the primitive characters $\chi^{(i)}$ inducing $\psi_0^i$. Then $K$ contains the trivial character, $\psi(1,0,0)=1$ for $\psi\in K$, and the characters of $\Psi_0\setminus K$ are the characters with $(a,a')\ne(0,0)$.

\emph{The functions $G_\psi$.} Let $S_0$ be the set of primes dividing $2qM_\psi$. For $p\in S_0$ the inertia group $I_p$ acts trivially on $\sqrt{-d}$, so $I_p\subset\Phi$, and $p$ splits in $\mathbb Q(\sqrt{-d})$, because $-d\equiv1\pmod8$, $-d\equiv1\pmod q$ and $-d\equiv1\pmod p$ for the odd primes $p\le Y^{\varepsilon/5}$, which include the prime factors of $M_\psi$ for large $Y$. So some Frobenius element $\sigma_p$ lies in $\Phi$. The inertia group at $d$ is $\{0,(1,0,0)\}$, and $E$ is unramified outside $2qdM_\psi$. Proposition~\ref{excess} therefore applies. For a prime $p$ and a character $\psi$, the average $\psi(p^i)$ of Proposition~\ref{excess} equals $\chi_\psi(p)^i$ when $\psi$ is trivial on $I_p$ and vanishes otherwise, and this is $\chi_\psi(p^i)$ because $\chi_\psi$ is primitive. Hence $G_\psi=G^\ast_{\chi_\psi}$ for nontrivial $\psi$ and $G_\psi=G$ for trivial $\psi$, and $\operatorname{Re}G_\psi=G_{\chi_\psi}$.

\emph{The hypotheses of Theorem~\ref{step5}.} The set $T$ satisfies \eqref{equi} with $c=1/(3o+1)$ and $\nu=o\,x^{-1+\varepsilon}$ by Lemma~\ref{shares}. Theorem~A(i) and Proposition~\ref{nodes} give the bounds for the nodes, with $\eta'=W+mh_0$ and $b'X'=xb$, and the weights satisfy $\lvert a_t/b-x\rvert\le1$ by Proposition~\ref{nodes}(iii). Lemma~\ref{reg}(i) gives \eqref{L1}. Lemma~\ref{reg}(ii) gives \eqref{L2} with $M=\log(x+1)+C_d+\max_{\psi\in\Psi_0}\lvert\ell_\psi\rvert$. For $\psi\in K$ the number $\ell_\psi$ depends on $\psi_0$ alone, and for $\psi\in\Psi_0\setminus K$ Theorems~\ref{stage1}(ii) and~\ref{stage2}(ii) at the scale $Y_0$ with $K=J$ give $\lvert\ell_\psi\rvert\le401\,c_0J\log Y_0\le401\,c_0J\log x$, so $M\le C(\psi_0)\log x$ with a constant that depends on $\psi_0$ and $J$.

\emph{The bound \eqref{L3}.} Let $\psi$ be a character of $\mathcal G$ and $N_\psi$ the conductor of $\chi_\psi$. It divides $N_1qd$ with $N_1\le2^{e+2}M_\psi\le8M_\psi n$, and $\log q\le Y_0^{\varepsilon/5}$, $\log d\le Y_0^{\varepsilon/5}\log^2Y_0$ with $Y_0\le x$, so $\log N_\psi\le\lambda:=A\,x^{\varepsilon/5}\log^2x$ for a constant $A$ and large $x$. If $\chi_\psi$ is real and $q$ or $d$ divides $N_\psi$, Corollary~\ref{realbox} with $M_1$ the odd part of $M_\psi$ shows that $L(\beta,\chi_\psi)\ne0$ for $1-\frac1{400}<\beta<1$. If $\chi_\psi$ is real and $N_\psi$ divides $N_1$, then $\chi_\psi$ is one of finitely many characters that depend on $M_\psi$ alone, because the real characters of $2$-power conductor have conductor at most $8$, and there is $c_\ast(M_\psi)>0$ such that none of these finitely many $L$-functions vanishes on $1-c_\ast<\beta<1$. Proposition~\ref{thinprop} with $c_\ast=\min(\frac1{400},c_\ast(M_\psi))$, and the remark after it for the trivial character, gives for $y\ge X'/3$, where $t_0=\log y\ge\frac12x^{\varepsilon/3}$,
\[
\frac{\lvert G_\psi(y)\rvert}{y}\le\frac2{c_\ast}e^{-c_\ast t_0}+C_1\lambda^4\Big(\frac{2\lambda}{c_1}e^{-c_1t_0/(2\lambda)}+\frac8{c_1^2}\Big(1+\frac{c_1\sqrt{t_0}}2\Big)e^{-c_1\sqrt{t_0}/2}\Big).
\]
Since $\frac13-\frac15>\frac18$ and $\frac16>\frac18$, the exponents $c_1t_0/(2\lambda)$ and $c_1\sqrt{t_0}/2$ exceed $2x^{\varepsilon/8}$ for large $x$, and the factors in front are powers of $x$. Hence \eqref{L3} holds with $\omega=\exp(-x^{\varepsilon/8})$, and Proposition~\ref{nodes}(iv) gives $2f(x+1)\omega\le1/x$.

\emph{The conclusion at $x$.} By Theorem~A(ii) the hypothesis of Theorem~\ref{step5} on the excess holds with $\mathrm{Err}=x^{C\varepsilon}\mathcal B$. Proposition~\ref{ranks}(ii) gives the hypotheses of Proposition~\ref{sizes} with $N_T=\lvert T\rvert\ge\frac12\delta_1n$ and $N_2=\lvert T_2\rvert$. In the last hypothesis of Corollary~\ref{asm}, the term $\nu\sum_{\psi\in K}\lvert G_\psi(x)\rvert$ is at most $o\,x^{-1+\varepsilon}\cdot o\,B(\psi_0)x$, where $B(\psi_0)$ bounds $\lvert G^\ast_{\chi^{(i)}}(y)\rvert/y$ for the $o$ characters $\chi^{(i)}$, and $\lvert\Psi_0\rvert M\le4oC(\psi_0)\log x$. Both are smaller than $\frac3{26}x^{1/2}$ for large $x$. Corollary~\ref{asm} gives
\[
G(x)+\frac1{3o+1}\sum_{i=1}^{o-1}G_{\chi^{(i)}}(x)\le x^{1/2+(C+3)\varepsilon}+\sum_{\psi\in\Psi_0\setminus K}\lvert G_\psi(x)\rvert .
\]
Each $\psi\in\Psi_0\setminus K$ has $\chi_\psi=\xi(\frac{\cdot}{q})^a(\frac{\cdot}{d})^{a'}$ with $\xi$ primitive of conductor dividing $N_0$ and $(a,a')\ne(0,0)$, and $(\frac{\cdot}{q})$ and $(\frac{-q}{\cdot})$ induce the same primitive character because $q\equiv3\pmod4$. Theorems~\ref{stage1}(iii) and~\ref{stage2}(iii) at the scale $Y_0$ with $K=J$ give $\lvert G_\psi(x)\rvert\le\lvert G^\ast_{\chi_\psi}(x)\rvert\le x^{1/2+\varepsilon}$ for $Y_0\le x\le Y_0^J$, an interval that contains $[Y,Y^{3/2}]$. There are $3o$ such characters, so for large $x$
\begin{equation}\label{finalbound}
G(x)+\frac1{3o+1}\sum_{i=1}^{o-1}G_{\chi^{(i)}}(x)\le x^{1/2+(C+4)\varepsilon}\qquad(Y\le x\le Y^{3/2}).
\end{equation}

\emph{The zeros.} The bound \eqref{finalbound} holds on the windows $[Y,Y^{3/2}]$ for arbitrarily large $Y$, uniformly throughout each selected window. The characters $\chi^{(i)}$ with $1\le i<o$ are nontrivial, and the weights $1/(3o+1)$ are positive. Corollary~\ref{corA} with $\lambda=\frac32$ and $\theta=\frac12+(C+4)\varepsilon$ gives $\zeta(s)\ne0$ and $L(s,\chi^{(i)})\ne0$ for $\operatorname{Re}s>\frac12+(C+4)\varepsilon$. The constant $C$ of Theorem~A does not depend on $\varepsilon$, and $\varepsilon$ is an arbitrary number in $(0,\varepsilon_1)$, so these functions do not vanish for $\operatorname{Re}s>\frac12$.
When $o\ge2$ the character $\chi^{(1)}$ is $\psi_0$, and every choice of $\psi_0$ includes $\zeta$, so the theorem follows.
\end{proof}


\part{Theorem A}
\input{sec7_chain}
\input{sec8_kernel}
\input{sec8_gain}
\input{sec10_chaindiv}
\input{sec11_lifts}

\section{The loss and the hypotheses \texorpdfstring{\eqref{A1} and \eqref{A2}}{(3.1) and (3.2)}}\label{r2:sec:loss}

This section proves the hypotheses \eqref{A1} and \eqref{A2} of Proposition~\ref{acct} for every $u_0\in A_0$. Hypothesis \eqref{A1} compares the sum $S$ over primes with the total gain $\Gamma$, a number defined below from the valuation estimates of Section~\ref{sec:chaindiv}, the nondegeneracy exponents of the Hecke operators, and a choice of weights. Hypothesis \eqref{A2} compares the slope of a lattice $N_u$ with the slope of the base lattice $M_u$. The lattice $N_u$ is the sublattice of $M_u$ spanned locally by images of powers of Hecke operators, rescaled at each place by the gain. Its slope exceeds that of $M_u$ by the gain minus the colength of the sublattice, and the colength is controlled by a nondegeneracy theorem for Hecke operators on integral forms of representations of $\mathrm{GL}_k$ along the edges of the building.

Subsection~\ref{r2:sec:loss-setting} fixes the classes of primes, the orbits of Frobenius, and the local quantities $g^0$, $g^1$, $\mathrm{Pr}$, $\tau$, $\mathrm{io}$ and $\mathrm{gain}$. Subsection~\ref{r2:sec:lp} states the nondegeneracy theorem and reduces it to one edge. Subsection~\ref{r2:sec:colength} turns nondegeneracy into a bound for the colength of the sublattice $\NP_w$. Subsection~\ref{r2:sec:Nu} defines $N_u$ and proves \eqref{A2}. Subsection~\ref{r2:sec:gainprime} compares the gain with the sum over primes and proves \eqref{A1}. Subsection~\ref{r2:sec:lossinputs} lists what this section takes from the other steps of the proof of Theorem~A.

\subsection{Primes, orbits and the local quantities}\label{r2:sec:loss-setting}

Fix admissible data (Definition~\ref{def:admissible}) at the scale $x$, and let $b$, $a$, $X'$, $b'$ and $\delta_1$ be as in \eqref{scale}. Put
\begin{equation}\label{r2:scales}
w''=\frac{b'}{xn},\qquad X''=\frac a{w''}=xnX',\qquad p_*=k^2x^{26},\qquad \theta=x^{-\varepsilon},\qquad \gamma_p=\frac p{p^2-1}.
\end{equation}
Let $T\subset H$ be the set of tops of the data and $N_T=\lvert T\rvert$. Fix $u_0\in A_0$, weights $a_t$ for $t\in T$ with $a\le a_t\le a+a/(kx^{13})$, and nodes $w_j=w_{u_0,j}\colon\Phi\to(0,\infty)$ for $1\le j\le m$. Let $S$ be the sum of Corollary~\ref{acctexcess}, that is, the sum over $t\in T$ and $1\le j\le m$ of the left side of \eqref{excessid} with $a=a_t$ and $w=w_j$.

\paragraph{Classes of primes.} Let $\mathcal P$ be the set of primes $p\notin S_0\cup\{d\}$ with $\sigma_p\notin\Phi$, the primes that are inert in $\mathbb Q(\sqrt{-d})$ and prime to $2qM_\psi$. For $p\in\mathcal P$ write $\sigma_p=(1,\epsilon_p,h_p)$ with $\epsilon_p\in\mathbb Z/2$ and $h_p\in H$. For odd $i$,
\[
p_t(i\sigma_p)=(0,\,1+\epsilon_p,\,t+ih_p).
\]
Let $\Phi_0$ and $\Phi_1$ be the sets of the elements $(0,0,h)$ and $(0,1,h)$ of $\Phi$. Under the identification of Section~\ref{sec:field} the places attached to the elements of $T$ are the elements $(0,0,t)$ of $\Phi_0$, which we call the \emph{top places}. We say that $p\in\mathcal P$ is of \emph{class I} when $\epsilon_p=1$, so that $p_t(i\sigma_p)\in\Phi_0$ for odd $i$, and of \emph{class II} when $\epsilon_p=0$, so that $p_t(i\sigma_p)\in\Phi_1$ for odd $i$. For $s\in H$ and odd $e$ put $\pi_p(s,e)=p_s(e\sigma_p)$. For $t\in T$, $1\le j\le m$ and $p\in\mathcal P$ put
\begin{equation}\label{r2:ask}
\mathrm{ask}_{j,t}(p)=\sum_{i\ \mathrm{odd}}\max\Big(\frac{a_t}{p^i}-w_j\big(\pi_p(t,i)\big),0\Big),
\end{equation}
and let $\mathrm{dt}_{j,t}=\log d\sum_{i\ge1}\max\big(a_t/d^i-w_j(p_tg_i),0\big)$ be the term of $d$ in \eqref{excessid}. Then
\begin{equation}\label{r2:Ssplit}
S=\sum_{t\in T}\sum_{j=1}^m\Big(\sum_{p\in\mathcal P}2\log p\;\mathrm{ask}_{j,t}(p)+\mathrm{dt}_{j,t}\Big).
\end{equation}
The terms with $p>2a$ vanish when the nodes are at least $1$, because then $a_t/p^i<1$.

\paragraph{Nodes and bottom weights.} The construction of the nodes (Section~\ref{sec:proofA}) attaches to each $t\in T$ an integer $b_t\ge b'$, the bottom weight of the base weight at $t$, and gives the following bounds, with a spread $\eta\ge1$ that satisfies the size conditions of Subsection~\ref{r2:sec:lossinputs}.
\begin{enumerate}
\item[(W1)] At a top place, $w_j(0,0,t)=b_t+o_j(t)$ with $b'/x\le o_j(t)\le b'/x+\eta$.
\item[(W2)] At a thick place, $b\le w_j(\rho)\le b+\eta+1$.
\item[(W3)] At every other place, $w''\le w_j(\rho)\le w''+\eta$.
\end{enumerate}
In particular $w_j(\rho)\ge w''\ge1$ for every $\rho$ when $x$ is large. In the quantities below the bottom weights $b_s$ are variables, and the value of $w_j$ at a top place is $b_s+o_j(s)$ for the variable $b_s$. For $1\le j\le m$ put $\kappa_j=\frac12\big(k+(2j-1)(t_r+1)\big)$, let $N_m=k+m(t_r+1)$, and let $c_G\ge0$ be a number with $\kappa_j\le c_G$ for every $j$. Put
\[
\hat w_j(\rho)=w_j(\rho)-\kappa_j,\qquad \mathrm{al}(s)=a_s-2c_G,\qquad \mathrm{cap}_j(s,e)=\hat w_j\big(\pi_p(s,e)\big)+c_G .
\]
The number $\kappa_j$ is the shift between a node and the parameter of the polynomial model of level $j$ in which the valuation estimates are written. The weights $a_s$, $b_s$ enter the quantities below as variables, and we write $(a,b)=(a_s,b_s)_{s\in T}$ for a vector of weights.

\paragraph{Orbits.} Let $p\in\mathcal P$. Put $s_p=h_p$ if $p$ is of class I and $s_p=2h_p$ if $p$ is of class II, let $n_p$ be the order of $h_p$ in the first case and twice the order of $2h_p$ in the second, and fix a set $R_p\subset H$ of representatives of the orbits of the translation by $s_p$. For $t_0\in R_p$ let $T(t_0)$ be the set of the elements of $T$ on the orbit of $t_0$ and $N(t_0)=\lvert T(t_0)\rvert$. If $p$ is of class I, the elements of the orbit of $t_0$ are $t_0+ih_p$ with $i\in\mathbb Z/n_p$, and we call a top $s=t_0+ih_p$ \emph{even} when $i$ is even. The orbit of $t_0$ is \emph{good} when it is good for $T$ in the sense of Definition~\ref{def:goodorbit}, with $h=h_p$ and $L_2=2\max_{p'}L_1(p')$, where $L_1(p')=\lceil\log_{p'}(kx^{13})^4\rceil$ and the maximum runs over the primes of class I. For a top $s$ and odd $e$ we write $\mathrm{Excl}(s,e)$ when $p$ is of class I, $\pi_p(s,e)$ is the top place of a top $s'\ne s$, and $e$ is the least positive integer with $s+eh_p=s'$.

\begin{lemma}\label{r2:evenorder}
Let $p\in\mathcal P$ be of class I with $n_p>o$. Then $n_p$ is even.
\end{lemma}

\begin{proof}
By Proposition~\ref{field}(iii) the group $H$ maps onto a cyclic group of order $2^e$ with a kernel of order $c$ dividing $o$. If $n_p$ were odd, the image of $h_p$ would be an element of odd order in a group of order $2^e$, hence trivial, so $h_p$ would lie in the kernel and $n_p\le c\le o$.
\end{proof}

When $x^{\varepsilon/4}>o$, every prime of class I whose $h_p$ has order at least $x^{\varepsilon/4}$ therefore has an even orbit length, and on each of its orbits the even and odd tops are well defined.

\begin{lemma}\label{r2:orbitsize}
Let $p$ be of class I with $n_p\ge x^{\varepsilon/4}$, and let $t_0\in R_p$ have a good orbit. Let $u$, $r_0$, $M_0$ be as in Definition~\ref{def:goodorbit}. Then
\[
1\le N(t_0),\qquad n_p\le(6r_0+1)\,u\,(M_0+1)\,\big(N(t_0)+1\big).
\]
In particular $n_p\,\delta_1\le C_{31}\,N(t_0)\log x$ for a constant $C_{31}$ depending on $\varepsilon$ and $x$ large.
\end{lemma}

\begin{proof}
Cut the residues modulo $n_p$ into $\lfloor n_p/((6r_0+1)u)\rfloor$ tiles as in Definition~\ref{def:goodorbit}(iii). Condition (iii) gives, in each run of $M_0+1$ cyclically consecutive tiles, a tile whose $3r_0+1$ even cells each contain an element of $P_t$, so $N(t_0)\ge1$, and the disjoint runs of $M_0+1$ consecutive tiles give at least $\lfloor n_p/((6r_0+1)u(M_0+1))\rfloor$ distinct elements of $P_t$. Since $u\le x^{\varepsilon/8}+1$ and $M_0\le C\log x$, the second claim gives the third.
\end{proof}

\paragraph{The valuation bounds at an orbit.} Fix $p\in\mathcal P$ with $p\le X''$, an orbit $t_0\in R_p$, a lift $1\le j\le m$ and weights $(a,b)$. The \emph{digit loss} at the orbit is
\[
\mathrm{dl}(t_0)=K_{d1}\,n_p\,N_m^2\log_p(2a)+K_{d2}\,N(t_0)\log_p(2a)^2,
\]
where $K_{d1},K_{d2}$ are the constants of the term $\Delta$ of Theorem~\ref{thm:valuation} in global parameters. The \emph{base bound} and the \emph{improved bound} are
\begin{align}
g^0_j(t_0)&=\sum_{s\in T(t_0)}\ \sum_{\substack{e\ \mathrm{odd}\\ \mathrm{not}\ \mathrm{Excl}(s,e)}}\max\Big(0,\frac{\mathrm{al}(s)}{p^e}-\mathrm{cap}_j(s,e)\Big)-\mathrm{dl}(t_0),\label{r2:g0}\\
g^1_j(t_0)&=\gamma_p\sum_{s\in T(t_0)}\mathrm{al}(s)-\sum_{\substack{s\in T(t_0)\\ s\ \mathrm{odd}}}\ \sum_{e\ \mathrm{odd}}\min\Big(\frac{\mathrm{al}(s)}{p^e},\mathrm{cap}_j(s,e)\Big)-\mathrm{dl}(t_0),\label{r2:g1}
\end{align}
where for $p$ of class II no top is called odd. The \emph{price} is
\begin{equation}\label{r2:price}
\mathrm{Pr}_j(t_0)=\begin{cases}\displaystyle\sum_{\substack{s\in T(t_0)\\ s\ \mathrm{odd}}}\big(\hat w_j(0,0,s)-b_s\big)+K_{\mathrm{Pr}}\,n_p\,(w''+\eta+c_G+N_m),&p\ge p_*,\\[2ex]
r_0\big(N(t_0)(2b'+c_G)+n_p(w''+\eta+c_G+N_m)\big),&p<p_*.\end{cases}
\end{equation}
These are the numbers $g^0$, $g^1$ and $\mathrm{PR}$ of Theorem~\ref{lowterm} at the orbit of $t_0$, written in the parameters of the chain (Section~\ref{sec:lifts}).

\paragraph{Nondegeneracy exponents and shares.} Subsection~\ref{r2:sec:lp} attaches to $p$ and $t_0$ a type $r_p$ of Hecke operator, $r_p=r_0$ for $p<p_*$ and $r_p=1$ for $p\ge p_*$, a depth $D_p$, a denominator $t_{t_0}\ge0$ and a share $\kappa_{t_0}(a,b)\ge0$ (Theorem~\ref{r2:LPthm}). Put
\[
\mathrm{Sig}(p)=m\sum_{t_0\ \mathrm{good}}\kappa_{t_0}(a,b),
\]
the sum over the representatives $t_0\in R_p$ with a good orbit. We call $p$ a \emph{candidate} when $p<X'$, $p$ is of class I, $n_p\ge x^{\varepsilon/4}$, and, if $p\ge p_*$, no power $p^i$ with $i\ge1$ satisfies $\lvert p^i-X'\rvert\le\rho_BX'$, where $\rho_B\le x^{-2}$ is the band width of Subsection~\ref{r2:sec:lossinputs}. The set of \emph{served orbits} at $p$ is
\[
\cV_p=\begin{cases}\{t_0\in R_p:\ \text{the orbit of }t_0\text{ is good}\},&p\text{ a candidate and }\mathrm{Sig}(p)\le\theta,\\ \emptyset&\text{otherwise.}\end{cases}
\]

\begin{definition}\label{r2:localq}
For $p\in\mathcal P$ with $p\le X''$ and weights $(a,b)$ put
\begin{align*}
S_{t_0}&=\Big\lceil\sum_{j=1}^m\max\big(0,\ g^1_j(t_0)-\mathrm{Pr}_j(t_0)-\max(g^0_j(t_0),0)\big)\Big\rceil\qquad(t_0\in\cV_p),\\
\tau(p)&=e_p+\sum_{t_0\in\cV_p}m\big(\lceil t_{t_0}\rceil+D_p+1\big)\quad\text{if }\cV_p\ne\emptyset,\qquad\tau(p)=0\quad\text{if }\cV_p=\emptyset,\\
\mathrm{io}(p)&=2\,\mathrm{Sig}(p)\max_{t_0\in\cV_p}S_{t_0}\quad\text{if }\cV_p\ne\emptyset,\qquad\mathrm{io}(p)=0\quad\text{if }\cV_p=\emptyset,\\
\mathrm{gain}_j(p)&=\sum_{t_0\in\cV_p}\max\big(g^1_j(t_0)-\mathrm{Pr}_j(t_0),\ \max(g^0_j(t_0),0)\big)+\sum_{t_0\in R_p\setminus\cV_p}\max\big(g^0_j(t_0),0\big),
\end{align*}
where $e_p$ is the exponent of the comparison of integral structures in (L2) of Subsection~\ref{r2:sec:lossinputs}. Put $\mathrm{gain}_j(p)=0$ when $p>X''$, or when $p$ is of class I and $h_p=0$.
\end{definition}

The number $S_{t_0}$ is the exponent $S(v)$ of Theorem~\ref{lowterm} at the served orbit $v=t_0$, and $\mathrm{gain}_j(p)$ is the contribution of the orbits at $p$ to the exponent $T$ of that theorem at the level $j$, with $B=\cV_p$ and $\mathrm{PR}=\mathrm{Pr}$.

\paragraph{Dropped primes and the total gain.} A prime $p\in\mathcal P$ \emph{has no gain} when $p>X''$, or when $p$ is of class I with $h_p$ of odd order and $h_p\ne0$, or when $p$ is of class I with $p<X'(1+C_r\rho_B)$ and one of the following holds: $n_p<x^{\varepsilon/4}$; $p\ge p_*$ and $\lvert p^i-X'\rvert\le\rho_BX'$ for some $i\ge1$; $p$ is a candidate with $\mathrm{Sig}(p)>\theta$. Here $C_r\ge1$ is a constant. For the auxiliary primes $T_u$ of Section~\ref{r2:sec:Nu} let $\mathcal D$ be the set of the primes that have no gain and differ from every $T_u$. Put $\mathrm{gain}^{\mathcal D}_j(p)=0$ for $p\in\mathcal D$ and $\mathrm{gain}^{\mathcal D}_j(p)=\mathrm{gain}_j(p)$ otherwise. If $p$ has no gain then $\cV_p=\emptyset$, because a prime with $\cV_p\ne\emptyset$ is a candidate with $\mathrm{Sig}(p)\le\theta$ and with $n_p$ even by Lemma~\ref{r2:evenorder}. So $\tau(p)=\mathrm{io}(p)=0$ for $p\in\mathcal D$. The \emph{total gain} is
\begin{equation}\label{r2:Gamma}
\Gamma=\sum_{p\in\mathcal P}2\log p\Big(\sum_{j=1}^m\mathrm{gain}^{\mathcal D}_j(p)-\tau(p)\Big),
\end{equation}
and the \emph{total correction} is $\mathrm{IO}=\sum_{p\in\mathcal P}\log p\;\mathrm{io}(p)$. Both depend on $u_0$ through the nodes, and on the weights $(a,b)$.

\begin{remark}\label{r2:placeform}
Let $\bK$ be the coefficient field of the hermitian lattices of Proposition~\ref{acct}, of degree $[\bK:\mathbb Q]$. For a finite place $w$ of $\bK$ with prime ideal $\mathfrak p_w$, ramification index $e_w$ and residue characteristic $p_w$ put $\ell_w=\log N\mathfrak p_w/[\bK:\mathbb Q]$ and $\sigma(w,y)=2\ell_we_wy$. Then $\sum_{w\mid p}\sigma(w,y)=2y\log p$, because $\sum_{w\mid p}e_wf_w=[\bK:\mathbb Q]$. So $\Gamma=\sum_w\sigma\big(w,\sum_j\mathrm{gain}^{\mathcal D}_j(p_w)-\tau(p_w)\big)$, which is the form in which $\Gamma$ enters the slope of $N_u$ in Subsection~\ref{r2:sec:Nu}.
\end{remark}
\subsection{Nondegeneracy of Hecke operators}\label{r2:sec:lp}

This subsection states the nondegeneracy theorem for Hecke operators at one Frobenius orbit, and reduces it to a statement on one edge of the building.

\paragraph{The local setting.} Let $p$ be a prime, $n'\ge2$ an even integer, $q=p^{n'}$, $W=W(\mathbb F_q)$ the ring of Witt vectors with fraction field $\mathbb Q_q$ and Frobenius $\sigma$. Let $\cO$ be a complete discrete valuation ring finite over $W$, with a uniformizer $\varpi$ such that $\varpi^2=p$, fraction field $\mathcal K$ and an embedding $\iota\colon\mathbb Q_q\to\mathcal K$ compatible with $W\to\cO$. Let $k\ge2$. A \emph{weight} is a family $\lambda=(\lambda_s)_{s\in\mathbb Z/n'}$ of dominant weights of $\mathrm{GL}_k$. Put $\lambda^*_s=\lambda_s$ for even $s$ and $\lambda^*_s(i)=-\lambda_s(k-1-i)$ for odd $s$, and let
\[
V=\bigotimes_{s\in\mathbb Z/n'}\nabla_{\mathcal K}(\lambda^*_s),
\]
on which $\mathrm{GL}_k(\mathbb Q_q)$ acts through $\iota\circ\sigma^s$ on the factor $s$. For a $W$-lattice $L\subset\mathbb Q_q^k$ let $L^{(s)}$ be the $\cO$-span of $\iota\sigma^s(L)$, and let $H(L)=V\cap\bigotimes_s\nabla_\cO(\lambda^*_s;L^{(s)})$ be the tensor product of the integral dual Weyl modules attached to the lattices $L^{(s)}$. Then $H(gL)=gH(L)$ for $g\in\mathrm{GL}_k(\mathbb Q_q)$. We normalize the action so that $\mathrm{diag}(c_1,\dots,c_k)$ multiplies a vector of weight $\mu$ by $\prod_ic_i^{-\mu_i}$, and we use that $H(\cO^k)$ has an $\cO$-basis of weight vectors. These conventions are those of the integral structures of Section~\ref{sec:lifts}.

Let $1\le r<k$. An \emph{$r$-edge} is a pair $e=(L''\subset L)$ of lattices with $pL\subset L''$ and $L/L''\cong\mathbb F_q^r$. We write $\mathrm{src}(e)=L''$ and $\mathrm{tgt}(e)=L$. The \emph{midpoint} of $e$ is the family of $\cO$-lattices $M_e^{(s)}=L''^{(s)}+\varpi L^{(s)}$, and $H(M_e)$ is defined like $H(L)$. Put $A_r=\sum_s\sum_{i<r}\lambda^*_s(i)$,
\[
R(e)=\varpi^{-A_r}H(M_e),\qquad R'(e)=\varpi^{-A_r}H(M_e)^\vee,\qquad \mathrm{be}(F,G)=G(p^{A_r}F)\quad(F\in V,\ G\in V^*),
\]
where $^\vee$ is the $\cO$-dual lattice in $V^*$. For $t\in\mathbb N$ the \emph{enlarged lattices} are
\[
N^{\mathrm{src}}_t(L'')=\varpi^{-t}H(L'')\cap\bigcap_{\mathrm{src}(e)=L''}R(e),\qquad N^{\mathrm{tgt}}_t(L)=\varpi^{-t}H(L)^\vee\cap\bigcap_{\mathrm{tgt}(e)=L}R'(e),
\]
the intersections over the $r$-edges with the given source or target.

\begin{lemma}\label{r2:perfect}
For every $r$-edge $e$ the lattices $R(e)$ and $R'(e)$ are dual under $\mathrm{be}$, and $H(\mathrm{src}\,e)\subset R(e)$ and $H(\mathrm{tgt}\,e)^\vee\subset R'(e)$.
\end{lemma}

\begin{proof}
All three statements are invariant under $\mathrm{GL}_k(\mathbb Q_q)$, which acts transitively on $r$-edges, so we may take $L''=\cO^k$ and $L=d^2\cO^k$ with $d=\mathrm{diag}(\varpi^{-1}I_r,I_{k-r})$, so that $M_e=d\cO^k$ in every factor. Let $(v_\mu)$ be a basis of $H(\cO^k)$ of weight vectors and $(v^*_\mu)$ the dual basis. Writing $|\mu|_r=\sum_{i<r}\mu_i$, the lattice $H(M_e)$ has the basis $\varpi^{|\mu|_r}v_\mu$ and $H(L)$ has the basis $p^{|\mu|_r}v_\mu$. Every weight $\mu$ of $V$ is a sum over $s$ of weights of $\nabla(\lambda^*_s)$, each of which is below $\lambda^*_s$ in the dominance order, so $|\mu|_r\le A_r$. Hence $R(e)$ has the basis $\varpi^{|\mu|_r-A_r}v_\mu$, which contains $H(\cO^k)$, and $R'(e)$ has the basis $\varpi^{-|\mu|_r-A_r}v^*_\mu$, which contains the basis $p^{-|\mu|_r}v^*_\mu$ of $H(L)^\vee$. Finally $\mathrm{be}(\varpi^{|\mu|_r-A_r}v_\mu,\varpi^{-|\nu|_r-A_r}v^*_\nu)=\delta_{\mu\nu}$.
\end{proof}

\begin{definition}\label{r2:LP}
Let $t,D,c^*\in\mathbb N$ and $\kappa\in\mathbb R$. The statement $\mathrm{LP}(r,\lambda;t,D,\kappa,c^*)$ holds when $4(D+1)\le c^*$ and for every lattice $L\subset\mathbb Q_q^k$ there are $r$-edges $e_1,\dots,e_M$ with target $L$ and vectors $F_1,\dots,F_M\in V$ such that
\begin{enumerate}
\item[(i)] $F_i\in N^{\mathrm{src}}_t(\mathrm{src}\,e_i)$;
\item[(ii)] $F_i\in\varpi^{c^*}R(e')$ for every $r$-edge $e'\ne e_i$ with $\mathrm{src}(e')=\mathrm{src}(e_i)$;
\item[(iii)] there are an integer $n\ge(1-\kappa)\dim V$, an $\cO$-linear map $u\colon\cO^n\to N^{\mathrm{tgt}}_t(L)$ and a matrix $A\in M_{n\times M}(\cO)$ with $\sum_iA_{ji}\,\mathrm{be}(F_i,u(z))=p^Dz_j$ for all $z\in\cO^n$ and all $j$.
\end{enumerate}
\end{definition}

Condition (ii) says that $F_i$ is concentrated on the edge $e_i$ among the edges out of its source. For $\kappa\ge1$ the statement holds with $M=n=0$, so its content is the bound $\kappa<1$. The \emph{least share} $\kappa^{\min}(r,\lambda;t,D,c^*)$ is the least $\kappa\in[0,1]$ for which $\mathrm{LP}(r,\lambda;t,D,\kappa,c^*)$ holds. It exists because the condition is monotone in $\kappa$ and is determined by the integer $\lceil(1-\kappa)\dim V\rceil$, which takes finitely many values for $\kappa\in[0,1]$.

\paragraph{The Hecke reading.} For a function $F$ on lattices with values in $V$, the Hecke operator of type $r$ is $(T_rF)(L)=\sum_{\mathrm{tgt}(e)=L}F(\mathrm{src}\,e)$. The vectors $F_i$ are local pieces of the operator $p^{A_r}T_r$, and condition (iii) says that $p^{A_r}T_r$, applied to forms that are integral in the enlarged sense, reaches $p^D$ times a share $1-\kappa$ of the dual lattice at $L$. Subsection~\ref{r2:sec:colength} uses the statement in the following form.

\begin{lemma}\label{r2:smithlocal}
Suppose that $\mathrm{LP}(r,\lambda;t,D,\kappa,c^*)$ holds. Let $L$ be a lattice and let $H^{\mathrm{du}}(L)=\{F\in V:G(F)\in\cO\text{ for all }G\in N^{\mathrm{tgt}}_t(L)\}$. Then there are $r$-edges $e_1,\dots,e_M$ with target $L$, vectors $F_k\in N^{\mathrm{src}}_t(\mathrm{src}\,e_k)$, an $\cO$-basis $(b_\alpha)$ of $H^{\mathrm{du}}(L)$ and a set $\mathrm{bad}$ of indices with $\lvert\mathrm{bad}\rvert\le\kappa\dim V$ such that
\begin{enumerate}
\item[(a)] $p^{A_r}F_k\in\varpi^{c^*}H^{\mathrm{du}}(L')$ for every $r$-edge $(\mathrm{src}\,e_k\subset L')$ with $L'\ne L$;
\item[(b)] $p^Db_\alpha\in\operatorname{span}_\cO\{p^{A_r}F_k\}$ for every $\alpha\notin\mathrm{bad}$.
\end{enumerate}
\end{lemma}

\begin{proof}
Take the edges and vectors of Definition~\ref{r2:LP}. For an $r$-edge $e'$ with target $L'$ we have $N^{\mathrm{tgt}}_t(L')\subset R'(e')$, and $R'(e')$ is the $\mathrm{be}$-dual of $R(e')$ by Lemma~\ref{r2:perfect}. Hence $F\in R(e')$ implies $G(p^{A_r}F)=\mathrm{be}(F,G)\in\cO$ for $G\in N^{\mathrm{tgt}}_t(L')$, that is, $p^{A_r}F\in H^{\mathrm{du}}(L')$. Applied to $e'=e_k$ and to the edges $e'\ne e_k$ of (ii), this gives $p^{A_r}F_k\in H^{\mathrm{du}}(L)$ and (a).

Put $\psi_j=\sum_kA_{jk}p^{A_r}F_k\in H^{\mathrm{du}}(L)$ and $u_l=u(\epsilon_l)$ for the unit vectors $\epsilon_l$ of $\cO^n$. Then $u_l(\psi_j)=p^D\delta_{jl}$. Let $\Psi\colon H^{\mathrm{du}}(L)\to\cO^n$ be $F\mapsto(u_l(F))_l$, so $\Psi(\psi_j)=p^D\epsilon_j$. The $\psi_j$ are linearly independent, and the saturation $\Sigma$ of their span in $H^{\mathrm{du}}(L)$ has rank $n$ and is a direct summand, $H^{\mathrm{du}}(L)=\Sigma\oplus\Sigma'$. The map $\Psi$ is injective on $\Sigma$, because $\Sigma\otimes\mathcal K$ is spanned by the $\psi_j$. For $F\in\Sigma$ the vector $\Psi(p^DF)$ lies in $p^D\cO^n=\Psi(\operatorname{span}\psi)$, so $p^DF=\psi$ for some $\psi\in\operatorname{span}\psi_j$ by injectivity. Take for $(b_\alpha)$ a basis of $\Sigma$ followed by a basis of $\Sigma'$, and let $\mathrm{bad}$ index the latter. Then $\lvert\mathrm{bad}\rvert=\dim V-n\le\kappa\dim V$.
\end{proof}

\paragraph{Reduction to one edge.} The next four statements reduce $\mathrm{LP}$ to families of sources and tests on single edges into one lattice.

\begin{lemma}\label{r2:onelattice}
If conditions (i) to (iii) of Definition~\ref{r2:LP} hold for one lattice $L_0$, they hold for every lattice.
\end{lemma}

\begin{proof}
$\mathrm{GL}_k(\mathbb Q_q)$ acts transitively on lattices. An element $g$ with $gL_0=L$ maps $r$-edges to $r$-edges and maps $H$, $R$, $R'$, $N^{\mathrm{src}}_t$ and $N^{\mathrm{tgt}}_t$ at $L_0$ to the same objects at $L$, and it preserves $\mathrm{be}$. Transport the data along $g$.
\end{proof}

\begin{lemma}\label{r2:smith}
Let $\Lambda$ be an $\cO$-module and $\varphi_1,\dots,\varphi_M\colon\Lambda\to\cO$ linear maps. Suppose that for every $c\in\cO^M$ with some $c_i$ a unit there is $G\in\Lambda$ with $\ord_\varpi\big(\sum_ic_i\varphi_i(G)\big)\le 2D$. Then there is a linear $u\colon\cO^M\to\Lambda$ with $\varphi_i(u(z))=p^Dz_i$ for all $i$.
\end{lemma}

\begin{proof}
Let $\Sigma\subset\cO^M$ be the image of $G\mapsto(\varphi_i(G))_i$. By the elementary divisor theorem there are a basis $f_1,\dots,f_M$ of $\cO^M$ and exponents $a_j\in\mathbb N\cup\{\infty\}$ with $\Sigma=\bigoplus_j\varpi^{a_j}\cO f_j$. The coordinate functional $f_j^*$ is given by a vector $c$ with a unit entry, so the hypothesis gives an element of $f_j^*(\Sigma)=\varpi^{a_j}\cO$ of valuation at most $2D$, and $a_j\le2D$. Hence $p^D\cO^M\subset\Sigma$. Choose $G_l\in\Lambda$ with $(\varphi_i(G_l))_i=p^D\epsilon_l$ and put $u(z)=\sum_lz_lG_l$.
\end{proof}

A \emph{source-test system} at $L$ is a set $E$ of $r$-edges into $L$ and, for each $e\in E$, vectors $F_{e,1},\dots,F_{e,n_e}$ that satisfy (i) and (ii) of Definition~\ref{r2:LP}, such that for every $c\in\cO^{n_e}$ with a unit entry there is $G\in N^{\mathrm{tgt}}_t(L)$ with $\ord_\varpi\mathrm{be}(\sum_ic_iF_{e,i},G)\le2D$ and $G\in\varpi^{c^*}R'(e')$ for every $e'\in E\setminus\{e\}$.

\begin{proposition}\label{r2:assembly}
Suppose $2(D+1)\le c^*$. If a source-test system exists at $L$ with $\sum_en_e\ge(1-\kappa)\dim V$, then (i) to (iii) of Definition~\ref{r2:LP} hold at $L$. If for every $r$-edge $e$ into a fixed lattice $L_0$ there is a source-test system on $\{e\}$ alone with $n_e\ge\rho\dim V/q^{r(k-r)}$, and the tests satisfy $G\in\varpi^{c^*}R'(e')$ for the other edges $e'$ into $L_0$, then $\mathrm{LP}(r,\lambda;t,D,\kappa,c^*)$ holds for every $\kappa$ with $1-\kappa\le\rho$ and $4(D+1)\le c^*$.
\end{proposition}

\begin{proof}
Take all the $F_{e,i}$ as the family of Definition~\ref{r2:LP}, put $\Lambda=N^{\mathrm{tgt}}_t(L)$ and $\varphi_{e,i}(G)=\mathrm{be}(F_{e,i},G)$. These are integral on $\Lambda$ because $\Lambda\subset R'(e)$ and $F_{e,i}\in R(e)$, which are dual by Lemma~\ref{r2:perfect}. Let $c$ have a unit entry in the block of $e_0$, and take the test $G$ of $e_0$ for the restriction of $c$ to that block. For $e\ne e_0$ we have $G\in\varpi^{c^*}R'(e)$, so $\varphi_{e,i}(G)\in\varpi^{c^*}\cO\subset\varpi^{2D+2}\cO$. So $\sum c_{e,i}\varphi_{e,i}(G)$ is an element of valuation at most $2D$ plus an element of valuation at least $2D+2$, and it has valuation at most $2D$. Lemma~\ref{r2:smith} gives (iii) with $n=\sum_en_e$ and $A$ the identity.

For the second statement, the $r$-edges into $L_0$ correspond to the subspaces of codimension $r$ of $L_0/pL_0\cong\mathbb F_q^k$, and their number is the Gaussian binomial coefficient $\binom kr_q\ge q^{r(k-r)}$. The union of the single-edge systems is a source-test system at $L_0$ with $\sum_en_e\ge\rho\dim V\ge(1-\kappa)\dim V$, and Lemma~\ref{r2:onelattice} passes from $L_0$ to every lattice.
\end{proof}

\paragraph{The theorem.} Fix the data of Subsection~\ref{r2:sec:loss-setting}, a prime $p\in\mathcal P$ of class I with $n_p\ge x^{\varepsilon/4}$, and a good orbit $t_0\in R_p$, and put $n'=n_p$. Let $\nu_1=\lceil\max(6eG'\delta_1,(5+8\varepsilon)\log_2x)\rceil+1$ be the bound of Definition~\ref{def:goodorbit}(ii), $Q_p=p^{G'}$, $N^+=40\,n'\delta_1$, $c_0=2b'+2c_G$, $\mathrm{La}=K_{\mathrm{La}}\nu_1^2\log a$, and let $K_\bullet$ denote constants. The base weight at a position $s\in\mathbb Z/n'$ of the orbit is a dominant weight $\lambda_s$ of $\mathrm{GL}_k$ whose middle entries have absolute value at most $c_G$, and whose extreme entries are $a_s$ and $-b_s$ when $t_0+sh_p$ is a top. Put
\begin{align*}
D^{\mathrm{II}}&=\Big\lceil K_{DII}N^+(\nu_1^3+r_0^5+k^2)\log_pa+2r_0n'c_G+r_0N^+c_0\Big\rceil,\\
t^{\mathrm{II}}&=K_{tII}\,r_0\,N(t_0)\Big(\frac{\mathrm{La}\,a}{pk}+a\,p^{-L_1(p)}+c_0+kQ_p\Big)+2r_0n'c_G,\\
\kappa^{\mathrm{II}}&=\min\big(1,\ K_{kII}\,N(t_0)\,r_0^2\,k\,x^{-26}\big),\\
D^{\mathrm I}&=\Big\lceil K_{DI}N^+x^{26}(\nu_1^5+k^3)\log_pa\Big\rceil,\\
t^{\mathrm I}(a,b)&=\sum_{y\ \mathrm{odd}}\min\Big(b_y,\frac{a_{y-d(y)}}{p^{d(y)}}\Big)+K_{\kappa1}N(t_0)\nu_1\Big(\frac a{p^2}+\frac{\mathrm{La}\,a}{pk}\Big)\\ &\quad {}+K_{\kappa2}N(t_0)\big(Q_pc_6+kQ_p+k\log_pa+c_G\big)+2n'c_G,
\end{align*}
where the sum runs over the odd tops $y$ of the orbit, $d(y)$ is the distance from $y$ back to the nearest even top (an odd number by Lemma~\ref{r2:evenorder}, and $0$ if there is none, in which case the term is $0$), and $c_6=8(D^{\mathrm I}+1)+K_{\kappa2}N^+k\log_pa+4N^+c_G+2n'c_G$.

\begin{theorem}\label{r2:LPthm}
Assume Hypothesis~\ref{r2:engines}. There are constants $K_\bullet$, depending on $\varepsilon$, with the following properties for $x$ large.
\begin{enumerate}
\item[(i)] \emph{Small primes.} If $13\le p<p_*$, the weights lie in the windows $a\le a_s\le a+a/(kx^{13})$ and $b'\le b_s\le b'+2n\delta_1x^{\varepsilon}\mathrm{La}\,b'/k$, and the size conditions (Z) of Subsection~\ref{r2:sec:lossinputs} hold, then $\mathrm{LP}(r_0,\lambda;\lceil t^{\mathrm{II}}\rceil,D^{\mathrm{II}},\kappa^{\mathrm{II}},8(D^{\mathrm{II}}+1))$ holds.
\item[(ii)] \emph{Medium primes.} Let the size conditions (Z) hold, and let $p_*\le p<X'$ with no power $p^i$ in the band $\lvert p^i-X'\rvert\le\rho_BX'$, and let $\kappa^{\mathrm I}(a,b)=\kappa^{\min}(1,\lambda;\lceil t^{\mathrm I}(a,b)\rceil,D^{\mathrm I},8(D^{\mathrm I}+1))$. Then the mean of $\kappa^{\mathrm I}(a,b)$ over the weights in the windows above is at most $K_m\,N(t_0)\,\nu_1\,x^{-13}$.
\end{enumerate}
\end{theorem}

\begin{proof}[Outline of the proof]
\notgiven{This proof is an outline. It uses Hypothesis~\ref{r2:engines}, and the Pieri splitting at the tops and the passage from segment shares to the share of $\mathrm{LP}$, statement (C11) of Subsection~\ref{w1:sec:Crest}, are not written out here.}
In both cases Proposition~\ref{r2:assembly} reduces the claim to the existence, for every $r$-edge $e$ into the lattice $W^k$, of a source-test system on $e$ of size at least $\rho\dim V/q^{r(k-r)}$, with $1-\rho$ equal to the share. For (i) the share $\kappa^{\mathrm{II}}$ comes from two counts. The Pieri rule splits the representation at the positions of the orbit into a flat part $\bigotimes_{s}\mathrm{Sym}^{d_s}(\mathcal K^k)$ at the tops and a remaining part, with a loss $(1-kc_0/a)^{N(t_0)}$ in dimension. The sources on $e$ are products of symbol polynomials in the flat part with an edge-localizing polynomial and a basis of the remaining part, and the rank of the span of the symbol polynomials is at least $q^{-r(k-r)}(1-N(t_0)r_0^2k/Y)\prod_s\binom{d_s+k-1}{k-1}$ with $Y=x^{13}$ (Hypothesis~\ref{r2:engines}(A)). The tests are pairings with the twisted Fock form after the Frobenius substitution adapted to $e$, and each source has a test at depth at most $D^{\mathrm{II}}$ (Hypothesis~\ref{r2:engines}(B)). For (ii) the edges are $1$-edges, the orbit is cut into segments between consecutive even tops, and the share is the sum of a model share $\kappa_P\le C_sN(t_0)\nu_1/Y$ and of the bad shares of the segment forms, whose mean over the windows is bounded by Hypothesis~\ref{r2:engines}(C). \notgiven{The three engines of Hypothesis~\ref{r2:engines} are not proved here; Subsections~\ref{w1:sec:C} to~\ref{w1:sec:A} prove parts of them.}
\end{proof}

\begin{hypothesis}[engines of Hecke nondegeneracy]\label{r2:engines}
In the setting of Theorem~\ref{r2:LPthm}, the following three statements hold.
\begin{enumerate}
\item[(A)] \emph{Rank of symbol spaces.} Let $\mathbb F$ be the residue field of $\cO$. For a frame system of the orbit with $r=r_0$ frame tops per segment, built from the marked tiles of Definition~\ref{def:goodorbit}(iii), and for degrees $a\le d_s\le2a$, the span of the reductions of the segment generators has dimension at least $q^{-rK}(1-N(t_0)r^2k/Y)\prod_s\binom{d_s+k-1}{k-1}$ over $\mathbb F$, where $K=k-r$, provided $C_kN_s^2\log a\le k$, $128k^2N_sY^2\le p^{L_1(p)-1}$ and $32pYk^2N_sQ_p^2\le a$, with $N_s$ the maximal number of tops of a segment.
\item[(B)] \emph{Detection for the twisted Fock form.} For every $F$ of multidegree $(d_s)$ satisfying the chain and frame conditions with balanced norm $0$ there is an integral $G$ satisfying the chain and frame conditions, with the Frobenius substitution $\mathrm{ph}$ adapted to the edge, such that $\ord_p\mathrm B(F,\mathrm{ph}(G))\le n_0+D_m+D_l+3\kappa_p$, where $n_0=\frac12\ord_p\prod_sd_s!$ and $D_m$, $D_l$, $\kappa_p$ are the logarithmic losses of the chain conditions, of the frame conditions and of the Pieri bases.
\item[(C)] \emph{Mean of the segment shares.} For a segment with $s$ tops and split index $\sigma$, the mean of its bad share over the top weights in the window and over the next top weight in the window is at most $C(s-\sigma+1)\nu_1/Y$.
\end{enumerate}
\end{hypothesis}

Subsections~\ref{w1:sec:C}, \ref{w1:sec:B} and~\ref{w1:sec:A} prove the parts of Hypothesis~\ref{r2:engines} listed there and state the remaining steps. \notgiven{The remaining steps, and hence Hypothesis~\ref{r2:engines}, are not proved here.} Statement (A) rests on Nakayama's lemma over $\cO$, the identity $f(X^{p^R})=f(X)^{p^R}$ over $\mathbb F_p$ with a discriminant that is a nonzero divisor modulo $\varpi$, the uniqueness of base $p^e$ expansions, an integral basis of monomials indexed by Pieri paths, which the literature calls standard monomials, and a concentration estimate for the weights of Pieri paths. Statement (B) rests on Fock-orthogonal projectors defined by the extremal projector of Asherova, Smirnov and Tolstoy, with denominators controlled by Kummer's theorem, and on the expansion of the unit parts of factorial digit splits in binomial polynomials. Statement (C) writes the Gram determinant of a segment form as a polynomial in the next top weight whose leading coefficient is the untwisted Fock Gram matrix, and bounds the mean valuation of such a polynomial over an arithmetic progression.

\begin{corollary}\label{r2:meankappa}
Put $\kappa_{t_0}(a,b)=\kappa^{\mathrm{II}}$ if $p<p_*$ and $\kappa_{t_0}(a,b)=\kappa^{\mathrm I}(a,b)$ if $p\ge p_*$, with the exponents $t_{t_0}$ and $D_p$ of the same case. For $x$ large, for every candidate $p$ and every good orbit $t_0\in R_p$, the mean of $\kappa_{t_0}$ over the weights in the windows is at most $N(t_0)x^{-12}$, and if $p<p_*$ then $\kappa_{t_0}\le N(t_0)x^{-12}$ for all weights.
\end{corollary}

\begin{proof}
For $p<p_*$, $\kappa^{\mathrm{II}}\le K_{kII}N(t_0)r_0^2kx^{-26}$ and $k\le2x^{9/2+10\varepsilon}$, so $K_{kII}r_0^2k\le x^{14}$ for large $x$. For $p\ge p_*$, Theorem~\ref{r2:LPthm}(ii) gives the mean bound $K_mN(t_0)\nu_1x^{-13}$, and $K_m\nu_1\le x$ for large $x$, since $G'\delta_1$ is at most a constant times $\log x$ and so $\nu_1\le C\log x$.
\end{proof}
\subsection{From nondegeneracy to the colength}\label{r2:sec:colength}

Let $w$ be a finite place of $\bK$ above a prime $p\in\mathcal P$ of class I with $p\le X''$ and $\cV_p\ne\emptyset$, and let $u$ be an index with $p\ne T_u$. This subsection bounds the colength of a sublattice $\NP_w$ of the base lattice at $w$, spanned by images of powers of the Hecke operators at the served orbits.

\paragraph{The forms at $w$.} Let $\cO_w$ be the completion of the ring of integers of $\bK$ at $w$, with fraction field $\bK_w$, and let $v_w(p)$ be the normalized valuation of $p$ in $\cO_w$. The chain datum of Section~\ref{sec:chaindiv} at level $0$ (Section~\ref{sec:lifts}) gives the following structure at $w$, in which the served orbits $\cV=\cV_p$ index the coordinates.
\begin{enumerate}
\item[(F1)] A set $\mathcal L$ of lattice tuples $L=\big((L_v)_{v\in\cV},L'\big)$, where $L_v$ is a lattice in $\mathbb Q_{q_v}^k$ for the unramified extension of degree $n_p$ of $\mathbb Q_p$, and a group $\Gamma$ that acts on $\mathcal L$ through elements $\gamma_v\in\mathrm{GL}_k(\mathbb Q_{q_v})$ on the components $L_v$, with finitely many orbits and trivial stabilizers.
\item[(F2)] A $\bK_w$-vector space $X=\bigotimes_{v\in\cV}V_v\otimes X'$ on which $\Gamma$ acts, where $V_v$ is the representation of Subsection~\ref{r2:sec:lp} attached to the base weights at the orbit $v$, with $\gamma$ acting on $V_v$ through $\gamma_v$, and for each $L$ an $\cO_w$-lattice $H'(L)\subset X'$ with $H'(\gamma L)=\gamma H'(L)$. The completion $\cO_w$ contains a uniformizer $\varpi$ with $\varpi^2=p$.
\item[(F3)] The space $\mathcal F$ of forms, the functions $F\colon\mathcal L\to X$ with $F(\gamma L)=\gamma F(L)$, and two lattices of forms $M\subset M^{\mathrm{na}}$, where $M^{\mathrm{na}}$ is the set of forms with $F(L)\in\bigotimes_vH(L_v)\otimes H'(L)$ for all $L$, and $p^{e_p}M^{\mathrm{na}}\subset M$.
\item[(F4)] For $v\in\cV$ the Hecke operator $y_v=p^{A_{r}}T_{r,v}$, where $(T_{r,v}F)(L)=\sum F(L[v\leftarrow L''])$ is the sum over the $r$-edges $(L''\subset L_v)$ and $L[v\leftarrow L'']$ is $L$ with $L_v$ replaced by $L''$, with $r=r_p$. The operators $y_v$ map $M$ and $M^{\mathrm{na}}$ into themselves.
\item[(F5)] A subspace $V^{\mathrm{ram}}\subset\mathcal F$, the forms of ramified type at $T_u$, such that $M^{\mathrm{ram}}=M\cap V^{\mathrm{ram}}$ is saturated in $M$.
\end{enumerate}
The operators $y_v$ and $y_{v'}$ commute, because they act on different components of the tuples.

For a family $\sigma=(\sigma_v)_{v\in\cV}$ of \emph{slots} $\sigma_v\in\{\mathrm{du},\mathrm{ba},\mathrm{ba}^+\}$ (dual, base, ball) let $M(\sigma)$ be the set of forms $F$ with
\[
F(L)\in\bigotimes_{v\in\cV}H^{\sigma_v}(L_v)\otimes H'(L)\qquad(L\in\mathcal L),
\]
where $H^{\mathrm{ba}}(L_v)=H(L_v)$, $H^{\mathrm{ba}^+}(L_v)=N^{\mathrm{src}}_{t_v}(L_v)$ and $H^{\mathrm{du}}(L_v)$ is the lattice of Lemma~\ref{r2:smithlocal}, all with $t_v=\lceil t_{v}\rceil$ the denominator of Theorem~\ref{r2:LPthm} at the orbit $v$. So $M(\mathrm{ba},\dots,\mathrm{ba})=M^{\mathrm{na}}$. For $B\subset\cV$ let $\sigma_B$ be the family with $\mathrm{du}$ on $B$ and $\mathrm{ba}$ off $B$, and write $\sigma[v\leftarrow s]$ for $\sigma$ with $\sigma_v$ replaced by $s$.

\begin{lemma}\label{r2:slots}
For every $\sigma$ and $v\in\cV$:
\begin{enumerate}
\item[(a)] $M(\sigma[v\leftarrow\mathrm{du}])\subset M(\sigma[v\leftarrow\mathrm{ba}])\subset M(\sigma[v\leftarrow\mathrm{ba}^+])$;
\item[(b)] $\varpi^{t_v}M(\sigma[v\leftarrow\mathrm{ba}^+])\subset M(\sigma[v\leftarrow\mathrm{ba}])$ and $\varpi^{t_v}M(\sigma[v\leftarrow\mathrm{ba}])\subset M(\sigma[v\leftarrow\mathrm{du}])$;
\item[(c)] $y_vM(\sigma[v\leftarrow\mathrm{ba}^+])\subset M(\sigma[v\leftarrow\mathrm{du}])$.
\end{enumerate}
\end{lemma}

\begin{proof}
It suffices to compare the lattices at one coordinate. (a) $H(L)\subset\varpi^{-t}H(L)$, and $H(L)\subset R(e)$ for every edge $e$ out of $L$ by Lemma~\ref{r2:perfect}, so $H\subset N^{\mathrm{src}}_t$. Also $H(L)^\vee\subset\varpi^{-t}H(L)^\vee$ and $H(L)^\vee\subset R'(e)$ for every edge into $L$, so $H(L)^\vee\subset N^{\mathrm{tgt}}_t(L)$, and $H^{\mathrm{du}}(L)\subset H(L)^{\vee\vee}=H(L)$. (b) $N^{\mathrm{src}}_t(L)\subset\varpi^{-t}H(L)$, and $N^{\mathrm{tgt}}_t(L)\subset\varpi^{-t}H(L)^\vee$ gives $\varpi^tH(L)\subset H^{\mathrm{du}}(L)$. (c) If $F(L[v\leftarrow L''])\in N^{\mathrm{src}}_t(L'')\otimes(\cdots)$ for an edge $e=(L''\subset L_v)$, then the factor at $v$ lies in $R(e)$, and $p^{A_r}R(e)\subset H^{\mathrm{du}}(L_v)$ by the first paragraph of the proof of Lemma~\ref{r2:smithlocal}. Summing over the edges gives the claim, and the other coordinates are unchanged.
\end{proof}

\paragraph{Delta forms.} For $L\in\mathcal L$ and $\xi\in X$ let $\delta_{L,\xi}$ be the form with $\delta_{L,\xi}(\gamma L)=\gamma\xi$ for $\gamma\in\Gamma$ and $\delta_{L,\xi}=0$ off $\Gamma L$. It is well defined because the stabilizer of $L$ is trivial. Let $\mathcal R\subset\mathcal L$ be a set of representatives of $\Gamma\backslash\mathcal L$. Every form satisfies $F=\sum_{L\in\mathcal R}\delta_{L,F(L)}$, so $F\mapsto(F(L))_{L\in\mathcal R}$ identifies $M(\sigma)$ with $\bigoplus_{L\in\mathcal R}\big(\bigotimes_vH^{\sigma_v}(L_v)\otimes H'(L)\big)$. In particular $M(\sigma)$ is free of finite rank, and
\[
d^{\mathrm{na}}:=\rk M^{\mathrm{na}}=\sum_{L\in\mathcal R}\rk H'(L)\prod_{v\in\cV}\dim V_v .
\]

\begin{proposition}\label{r2:lemmaJ}
Suppose that at every $v\in\cV$ the statement $\mathrm{LP}(r,\lambda_v;t_v,D,\kappa_v,c^*)$ holds with $c^*\ge4(D+1)$. Then for every nonempty $B\subset\cV$ there is a finite set $C_B\subset M(\sigma_B)$ with $\lvert C_B\rvert\le d^{\mathrm{na}}\prod_{v\in B}\kappa_v$ such that
\[
p^D\,M(\sigma_B)\subset\sum_{v\in B}y_v\big(M(\sigma_B[v\leftarrow\mathrm{ba}^+])\big)+\operatorname{span}_{\cO_w}C_B .
\]
\end{proposition}

\begin{proof}
Let $Q$ be the right side. It is a submodule of $M(\sigma_B)$ by Lemma~\ref{r2:slots}(c). Fix $L\in\mathcal R$. For each $v\in B$ apply Lemma~\ref{r2:smithlocal} to the lattice $L_v$. This gives edges $e^v_k=(L^v_k\subset L_v)$, vectors $F^v_k\in N^{\mathrm{src}}_{t_v}(L^v_k)$, a basis $(b^v_\alpha)$ of $H^{\mathrm{du}}(L_v)$ and a set $\mathrm{bad}_v$ with $\lvert\mathrm{bad}_v\rvert\le\kappa_v\dim V_v$. Choose bases of $H(L_v)$ for $v\notin B$ and of $H'(L)$. The pure tensors $\xi=\bigotimes_{v\in B}b^v_{\alpha_v}\otimes\eta$, with $\eta$ a product of the chosen basis vectors, form a basis of $\bigotimes_vH^{(\sigma_B)_v}(L_v)\otimes H'(L)$.

Let $\xi$ be such a tensor with $\alpha_v\notin\mathrm{bad}_v$ for some $v\in B$, and write $\xi=b^v_{\alpha_v}\otimes\xi'$. By Lemma~\ref{r2:smithlocal}(b) there are $c_k\in\cO_w$ with $p^Db^v_{\alpha_v}=\sum_kc_kp^{A_r}F^v_k$. Put $L^{(k)}=L[v\leftarrow L^v_k]$ and $\xi_k=F^v_k\otimes\xi'$. Then $\delta_{L^{(k)},\xi_k}\in M(\sigma_B[v\leftarrow\mathrm{ba}^+])$, because the other components of $L^{(k)}$ are those of $L$. We claim that
\[
y_v\,\delta_{L^{(k)},\xi_k}-\delta_{L,p^{A_r}\xi_k}\in\varpi^{c^*}M(\sigma_B).
\]
The value of $y_v\delta_{L^{(k)},\xi_k}$ at $L^*$ is $p^{A_r}$ times the sum of the values $\delta_{L^{(k)},\xi_k}(L^*[v\leftarrow L''])$ over the edges $(L''\subset L^*_v)$. A term is nonzero when $L^*[v\leftarrow L'']=\gamma L^{(k)}$ for some $\gamma\in\Gamma$, and then $\gamma$ is unique, $L^*=\gamma\big(L[v\leftarrow L']\big)$ with $L'=\gamma_v^{-1}L^*_v$, the pair $(L^v_k\subset L')$ is an $r$-edge, and the term is $p^{A_r}\gamma\xi_k$. The terms with $L'=L_v$ occur exactly at the tuples $\gamma L$, once each, and together they form $\delta_{L,p^{A_r}\xi_k}$. For the terms with $L'\ne L_v$, Lemma~\ref{r2:smithlocal}(a) gives $p^{A_r}F^v_k\in\varpi^{c^*}H^{\mathrm{du}}(L')$, so $p^{A_r}\gamma\xi_k$ lies in $\varpi^{c^*}$ times the lattice of $\sigma_B$ at $\gamma(L[v\leftarrow L'])=L^*$, by the equivariance $H^{\mathrm{du}}(\gamma_vL')=\gamma_vH^{\mathrm{du}}(L')$. The difference is a form, and its values lie in these lattices, which proves the claim. Hence
\[
p^D\delta_{L,\xi}=\sum_kc_k\delta_{L,p^{A_r}\xi_k}\in Q+\varpi^{c^*}M(\sigma_B).
\]
Let $C_B$ be the set of the forms $\delta_{L,\xi}$ for $L\in\mathcal R$ and the pure tensors $\xi$ with $\alpha_v\in\mathrm{bad}_v$ for all $v\in B$. Its size is at most $\sum_{L\in\mathcal R}\rk H'(L)\prod_{v\in B}\lvert\mathrm{bad}_v\rvert\prod_{v\notin B}\dim V_v\le d^{\mathrm{na}}\prod_{v\in B}\kappa_v$. Since the $\delta_{L,\xi}$ span $M(\sigma_B)$, we obtain $p^DM(\sigma_B)\subset Q+\varpi^{c^*}M(\sigma_B)\subset Q+p^{2D+2}M(\sigma_B)$.

Iterating, $p^DM(\sigma_B)\subset Q+p^{D+2}\cdot p^DM(\sigma_B)\subset Q+p^{j(D+2)}M(\sigma_B)$ for every $j\ge1$. Since $\cO_w$ is a complete Noetherian local ring and $M(\sigma_B)$ is finitely generated, the intersection of the submodules $Q+p^jM(\sigma_B)$ is $Q$ by Krull's intersection theorem applied to $M(\sigma_B)/Q$. So $p^DM(\sigma_B)\subset Q$.
\end{proof}

\paragraph{The walk over the served orbits.} For $B\subset\cV$ write $y_B^m=\prod_{v\in B}y_v^m$. For numbers $S_v\in\mathbb N$ and a lattice of forms $N$ that is stable under the $y_v$ put
\[
P(N)=\operatorname{span}_{\cO_w}\Big\{p^{\sum_{v\in\cV\setminus B}S_v}\,y_B^mF:\ B\subset\cV,\ F\in N\Big\}.
\]

\begin{proposition}\label{r2:walk}
Let $\pi$ be an element of a commutative ring $R$, let $\Lambda(\sigma)$ be $R$-submodules of an $R$-module $X$ indexed by the families of slots $\sigma\in\{\mathrm{du},\mathrm{ba},\mathrm{ba}^+\}^{\cV}$, let $y_v$ for $v\in\cV$ be commuting endomorphisms of $X$, and let $a_v,b_v,c_v,c_0,S_v\in\mathbb N$, $m\ge1$ and $\beta\ge0$. Suppose that for all $\sigma$ and $v$
\begin{gather*}
\Lambda(\sigma[v\leftarrow\mathrm{du}])\subset\Lambda(\sigma[v\leftarrow\mathrm{ba}])\subset\Lambda(\sigma[v\leftarrow\mathrm{ba}^+]),\\
\pi^{a_v}\Lambda(\sigma[v\leftarrow\mathrm{ba}^+])\subset\Lambda(\sigma[v\leftarrow\mathrm{ba}]),\\
\pi^{b_v}\Lambda(\sigma[v\leftarrow\mathrm{ba}])\subset\Lambda(\sigma[v\leftarrow\mathrm{du}]),\\
\pi^{c_v}\Lambda(\sigma[v\leftarrow\mathrm{ba}^+])\subset\Lambda(\sigma[v\leftarrow\mathrm{du}]),\\
y_v\Lambda(\sigma[v\leftarrow\mathrm{ba}^+])\subset\Lambda(\sigma[v\leftarrow\mathrm{du}]),
\end{gather*}
and that for every $B\subset\cV$ with $\lvert B\rvert=\beta+1$ there is a finite set $G_B\subset\Lambda(\sigma_B)$ with
\[
\pi^{c_0}\Lambda(\sigma_B)\subset\sum_{v\in B}y_v\Lambda(\sigma_B[v\leftarrow\mathrm{ba}^+])+\operatorname{span}_RG_B .
\]
Let $\Lambda_0=\Lambda(\mathrm{ba},\dots,\mathrm{ba})$, let $P_\pi(\Lambda_0)$ be the span of the elements $\pi^{\sum_{v\notin B}S_v}y_B^mF$ with $B\subset\cV$ and $F\in\Lambda_0$, and put $E=\sum_v\big(a_v+b_v+(m-1)c_v+mc_0\big)+\beta\max_vS_v$. Then there is a set $C\subset\Lambda_0$ with $\lvert C\rvert\le m^{\beta+1}\sum_{\lvert B\rvert=\beta+1}\lvert G_B\rvert$ and
\[
\pi^E\Lambda_0\subset P_\pi(\Lambda_0)+\operatorname{span}_RC .
\]
\end{proposition}

\begin{proof}
Order $\cV=\{v_1,\dots,v_K\}$. If $K\le\beta$, then $\pi^{\sum_vS_v}\Lambda_0\subset P_\pi(\Lambda_0)$ by the term $B=\emptyset$, and $\sum_vS_v\le E$. Let $K\ge\beta+1$. We run a branching process whose states are triples $(i,\sigma,z)$, where $i=(i_v)\in\{0,\dots,m\}^{\cV}$ counts the applications of each $y_v$, $\sigma$ is a family of slots, and $z\in\Lambda(\sigma)$. The element represented by a state is $\pi^{\epsilon}y^iz$, where $y^i=\prod_vy_v^{i_v}$ and $\epsilon$ is an exponent attached to the state. An orbit is \emph{finished} when $i_v=m$. Let $D(i)$ be the set of unfinished orbits and, when $\lvert D(i)\rvert\ge\beta+1$, let $B(i)$ be the set of its first $\beta+1$ elements. The process keeps the following invariant: $\sigma_v=\mathrm{ba}$ for finished $v$ and for $v\notin B(i)$, and $\sigma_v\in\{\mathrm{du},\mathrm{ba}^+\}$ for $v\in B(i)$ with $i_v\ge1$.

Start from $z\in\Lambda_0$ with $i=0$, $\sigma=(\mathrm{ba},\dots)$ and $\epsilon=0$. At a state with $\lvert D(i)\rvert\ge\beta+1$ and $B=B(i)$, multiply $z$ by $\pi^{b_v}$ for each $v\in B$ with $\sigma_v=\mathrm{ba}$ and by $\pi^{c_v}$ for each $v\in B$ with $\sigma_v=\mathrm{ba}^+$. The result lies in $\Lambda(\sigma_B)$. Multiply by $\pi^{c_0}$ and apply the hypothesis for $B$:
\[
\pi^{c_0}\pi^{(\cdots)}z=\sum_{v\in B}y_vz_v+g,\qquad z_v\in\Lambda(\sigma_B[v\leftarrow\mathrm{ba}^+]),\ g\in\operatorname{span}G_B .
\]
The state is replaced by the states $(i+\epsilon_v,\sigma_B[v\leftarrow\mathrm{ba}^+],z_v)$ for $v\in B$, where $\epsilon_v$ is the unit vector at $v$, and the element $\pi^{\epsilon'}y^ig$ is set aside as an exceptional term. If some $v$ becomes finished, multiply $z_v$ by $\pi^{a_v}$ to bring its slot at $v$ back to $\mathrm{ba}$. Since $B(i)$ consists of the first $\beta+1$ unfinished orbits, an orbit enters $B(i)$ once, when it has $i_v=0$ and slot $\mathrm{ba}$, and leaves it when it is finished. So along each branch an orbit $v$ that is finished has cost $b_v+(m-1)c_v+a_v$ in powers of $\pi$, and every application of the hypothesis costs $c_0$ and increases $\sum_vi_v$ by one.

Every branch reaches a state with $\lvert D(i)\rvert\le\beta$ after at most $mK$ steps. At such a state the element is $\pi^\epsilon y^iz$ with $\epsilon\le\sum_v(a_v+b_v+(m-1)c_v+mc_0)$. For $v\in D(i)$ the slot of $z$ at $v$ is $\mathrm{ba}$, or $\mathrm{du}\subset\mathrm{ba}$, or $\mathrm{ba}^+$ with $i_v\ge1$, in which case $y_v^{i_v}$ moves it into $\mathrm{du}$, or $\mathrm{ba}^+$ with $i_v=0$, which does not occur by the invariant. So $F=\prod_{v\in D(i)}y_v^{i_v}z\in\Lambda_0$, and
\[
\pi^{E}y^iz=\pi^{E-\epsilon-\sum_{v\in D(i)}S_v}\cdot\pi^{\epsilon}\Big(\pi^{\sum_{v\in D(i)}S_v}\,y_{\cV\setminus D(i)}^m\,F\Big)\in P_\pi(\Lambda_0),
\]
since $\lvert D(i)\rvert\le\beta$ gives $\epsilon+\sum_{v\in D(i)}S_v\le E$. The exceptional terms set aside at states with counter $i$ and set $B$ lie in $\pi^{(\cdots)}y^i\operatorname{span}G_B\subset\operatorname{span}(y^iG_B)$. At such a state $i_v=m$ for the orbits before $B$ in the order, $i_v=0$ for the orbits after $B$, and $0\le i_v<m$ for $v\in B$, so the counter is determined by $B$ and $(i_v)_{v\in B}\in\{0,\dots,m-1\}^B$. Let $C$ be the union of the sets $y^iG_B$ over these pairs $(B,i)$. Then $\lvert C\rvert\le m^{\beta+1}\sum_B\lvert G_B\rvert$, and $\pi^Ez$ is the sum of the elements of $P_\pi(\Lambda_0)$ obtained at the ends of the branches and of $\pi$-multiples of elements of $\operatorname{span}C$.
\end{proof}

\begin{lemma}\label{r2:layers}
Let $\cO$ be a discrete valuation ring with uniformizer $\varpi$, $\Lambda$ a free $\cO$-module of rank $d$, $P\subset\Lambda$ a submodule, $\pi=p$ an element with $v(p)\ge1$, and $\tau,S,S_{\mathrm{top}},\beta_0\in\mathbb N$. Suppose that for $0\le\beta\le\beta_0$ there is a set $C_\beta\subset\Lambda$ of $g_\beta$ elements with $p^{\tau+\beta S}\Lambda\subset P+\operatorname{span}C_\beta$, and that $p^{\tau+S_{\mathrm{top}}}\Lambda\subset P$. Then
\[
\operatorname{length}_\cO\big(\Lambda/(P\cap p^\tau\Lambda)\big)\le v(p)\Big(d\tau+S\sum_{\beta<\beta_0}g_\beta+S_{\mathrm{top}}\,g_{\beta_0}\Big).
\]
\end{lemma}

\begin{proof}
The module $p^\tau\Lambda/(P\cap p^\tau\Lambda)$ is isomorphic to $(P+p^\tau\Lambda)/P$, so the length on the left is $v(p)d\tau+\operatorname{length}\big((P+p^\tau\Lambda)/P\big)$. Put $\Phi_\beta=(P+p^{\tau+\beta S}\Lambda)/P$ for $0\le\beta\le\beta_0$. The module $\Phi_\beta$ is contained in $(P+\operatorname{span}C_\beta)/P$, which is generated by $g_\beta$ elements, and over a discrete valuation ring a submodule of a module generated by $g$ elements is generated by $g$ elements. The quotient $\Phi_\beta/\Phi_{\beta+1}$ is generated by $g_\beta$ elements and killed by $p^S$, so its length is at most $v(p)Sg_\beta$. The module $\Phi_{\beta_0}$ is generated by $g_{\beta_0}$ elements and killed by $p^{S_{\mathrm{top}}}$, because $p^{\tau+S_{\mathrm{top}}}\Lambda\subset P$, so its length is at most $v(p)S_{\mathrm{top}}g_{\beta_0}$. Summing over the filtration $\Phi_0\supset\Phi_1\supset\cdots\supset\Phi_{\beta_0}\supset0$ gives the claim.
\end{proof}

\paragraph{The sublattice and its colength.} Put $S_v=S_{t_0}$ for $v=t_0\in\cV$ (Definition~\ref{r2:localq}), $S=\max_vS_v$, and
\[
\NP_w=P(M)\cap p^{\lceil\tau(p)\rceil}M\cap V^{\mathrm{ram}} .
\]
When $\cV_p=\emptyset$ we put $\NP_w=M^{\mathrm{ram}}$.

\begin{theorem}\label{r2:colength}
In the setting (F1) to (F5), suppose that $\mathrm{LP}$ holds at every $v\in\cV_p$ with the exponents $\lceil t_v\rceil$, $D_p$, the share $\kappa_v$ and $c^*=8(D_p+1)$ of Corollary~\ref{r2:meankappa}, and that $\mathrm{Sig}(p)=m\sum_v\kappa_v\le\frac12$. Then
\begin{enumerate}
\item[(a)] $\NP_w\subset p^{\lceil\tau(p)\rceil}M^{\mathrm{ram}}$;
\item[(b)] $\operatorname{length}_{\cO_w}\big(M^{\mathrm{ram}}/\NP_w\big)\le v_w(p)\,\rk(M)\,\big(\tau(p)+\mathrm{io}(p)\big)$.
\end{enumerate}
\end{theorem}

\begin{proof}
(a) $\NP_w\subset p^{\lceil\tau\rceil}M\cap V^{\mathrm{ram}}$, and an element $p^{\lceil\tau\rceil}F$ with $F\in M$ that lies in $V^{\mathrm{ram}}$ has $F\in M^{\mathrm{ram}}$ because $M^{\mathrm{ram}}$ is saturated in $M$.

(b) Put $\tau_0=\sum_{v\in\cV}m(t_v+D_p+1)$, so that $\lceil\tau(p)\rceil=e_p+\tau_0$. Apply Proposition~\ref{r2:walk} with $R=\cO_w$, $\pi=\varpi$, $\Lambda(\sigma)=M(\sigma)$, $a_v=b_v=t_v$, $c_v=2t_v$, $c_0=2(D_p+1)$, and with $2S_v$ in place of $S_v$. The hypotheses on the slots are Lemma~\ref{r2:slots}, where $\varpi^{2t_v}\Lambda(\sigma[v\leftarrow\mathrm{ba}^+])\subset\Lambda(\sigma[v\leftarrow\mathrm{du}])$ is the composite of the two inclusions of (b) of that lemma. The hypothesis on the sets $B$ is Proposition~\ref{r2:lemmaJ}, since $p^D=\varpi^{2D}$ and $\varpi^{2D+2}\Lambda\subset\varpi^{2D}\Lambda$. The exponent is $E=2(\tau_0+\beta S)$, and $P_\varpi(M^{\mathrm{na}})=P(M^{\mathrm{na}})$. So for every $\beta\ge0$
\[
p^{\tau_0+\beta S}M^{\mathrm{na}}\subset P(M^{\mathrm{na}})+\operatorname{span}C_\beta,\qquad\lvert C_\beta\rvert\le d^{\mathrm{na}}m^{\beta+1}e_{\beta+1}(\kappa),
\]
where $e_{\beta+1}(\kappa)$ is the elementary symmetric function of the $\kappa_v$. Since $m^{\beta+1}e_{\beta+1}(\kappa)\le\mathrm{Sig}(p)^{\beta+1}/(\beta+1)!$, the set $C_\beta$ is empty for $\beta\ge\lvert\cV\rvert$. Now let $F\in M$. Then $F\in M^{\mathrm{na}}$, so $p^{\tau_0+\beta S}F=P_0+c$ with $P_0\in P(M^{\mathrm{na}})$ and $c\in\operatorname{span}C_\beta$. Multiplying by $p^{e_p}$ and using $p^{e_p}M^{\mathrm{na}}\subset M$ and the commutation of the $y_v$ with scalars gives $p^{e_p}P(M^{\mathrm{na}})\subset P(p^{e_p}M^{\mathrm{na}})\subset P(M)$, so
\[
p^{\lceil\tau\rceil+\beta S}M\subset P(M)+\operatorname{span}(p^{e_p}C_\beta),\qquad p^{e_p}C_\beta\subset M .
\]
The term $B=\emptyset$ of $P(M)$ gives $p^{\sum_vS_v}M\subset P(M)$. Lemma~\ref{r2:layers} with $\Lambda=M$, $\tau=\lceil\tau(p)\rceil$, $\beta_0=\lvert\cV\rvert$ and $S_{\mathrm{top}}=\sum_vS_v$ gives
\[
\begin{aligned}
\operatorname{length}\big(M/(P(M)\cap p^{\lceil\tau\rceil}M)\big)&\le v_w(p)\Big(\rk(M)\lceil\tau\rceil+S\,d^{\mathrm{na}}\sum_{\beta\ge0}\mathrm{Sig}(p)^{\beta+1}\Big)\\ &\le v_w(p)\rk(M)\big(\tau(p)+2S\,\mathrm{Sig}(p)\big),
\end{aligned}
\]
using $\mathrm{Sig}(p)\le\frac12$, $d^{\mathrm{na}}=\rk M$ (by (F3) the lattices $M\subset M^{\mathrm{na}}$ have the same rank), and $\lceil\tau(p)\rceil=\tau(p)$, an integer. The right side is $v_w(p)\rk(M)(\tau(p)+\mathrm{io}(p))$ by Definition~\ref{r2:localq}. Finally, the inclusion $M^{\mathrm{ram}}\to M$ induces an injection of $M^{\mathrm{ram}}/\NP_w$ into $M/(P(M)\cap p^{\lceil\tau\rceil}M)$, because $\NP_w=M^{\mathrm{ram}}\cap\big(P(M)\cap p^{\lceil\tau\rceil}M\big)$.
\end{proof}

The first claim of Theorem~\ref{r2:colength} gives that $\NP_w$ is spanned by forms of the shape $y_{B'}^mF_0$ with $F_0\in p^{\sum_{v\in\cV\setminus B'}S_v}M$, intersected with $p^{\tau}M^{\mathrm{ram}}$, which is the hypothesis of Theorem~\ref{lowterm} with $B=\cV_p$. This is the link between the loss and the divisibility of the lifts, which Section~\ref{r3:sec} uses.
\subsection{The lattice \texorpdfstring{$N_u$}{N\_u} and hypothesis \texorpdfstring{\eqref{A2}}{(A2)}}\label{r2:sec:Nu}

\paragraph{Hermitian lattices and slopes.} A hermitian lattice over $\cO_\bK$ is a finitely generated projective $\cO_\bK$-module $\Lambda$ with a hermitian norm on $\Lambda\otimes_{\sigma}\mathbb C$ for every embedding $\sigma\colon\bK\to\mathbb C$. Its normalized Arakelov degree is
\[
\degh\Lambda=\frac1{[\bK:\mathbb Q]}\Big(\log\#\big(\Lambda/(\cO_\bK s_1\oplus\dots\oplus\cO_\bK s_d)\big)-\sum_\sigma\log\lVert s_1\wedge\dots\wedge s_d\rVert_\sigma\Big)
\]
for any $\bK$-basis $s_1,\dots,s_d$ of $\Lambda\otimes\bK$ contained in $\Lambda$, and the value does not depend on the basis by the product formula. Following the normalization of Sections~\ref{r3:sec} and~\ref{r4:sec:R4}, the slope of a nonzero lattice is $\mu(\Lambda)=2\degh\Lambda/\rk\Lambda+c_\bK$ with a constant $c_\bK$ that depends on $\bK$, and $\mu_{\max}(\Lambda)$ is the maximum of $\mu(\Lambda')$ over the nonzero sublattices $\Lambda'\subset\Lambda$ with the induced metrics. The proofs below use differences of slopes of lattices over the same field, so the value of $c_\bK$ plays no role.

\begin{lemma}\label{r2:degree}
Let $\Lambda'\subset\Lambda$ be hermitian lattices of the same rank $d$ in the same $\bK$-space with the same metrics, and let $\mathfrak a=\prod_w\mathfrak p_w^{n_w}$ be a fractional ideal. For a finite place $w$ let $\ell^{\mathrm{len}}_w$ be the length of the $\cO_w$-module $\Lambda_w/\Lambda'_w$. Then
\[
\degh(\mathfrak a\Lambda')=\degh\Lambda-\sum_w\ell_w\,\ell^{\mathrm{len}}_w-d\sum_w\ell_w\,n_w,\qquad \ell_w=\frac{\log N\mathfrak p_w}{[\bK:\mathbb Q]} .
\]
\end{lemma}

\begin{proof}
Take the basis $s_i$ inside $\mathfrak a\Lambda'$. Then $\log\#(\Lambda/\mathfrak a\Lambda')$ is defined when $\mathfrak a\Lambda'\subset\Lambda$, and in general the index $[\Lambda:\mathfrak a\Lambda']$ is the generalized index $\prod_wN\mathfrak p_w^{\,\mathrm{len}_w(\Lambda_w/\Lambda'_w)+dn_w}$, which is multiplicative in towers. The index of $\cO s_1\oplus\dots\oplus\cO s_d$ in $\mathfrak a\Lambda'$ equals its index in $\Lambda$ divided by $[\Lambda:\mathfrak a\Lambda']$, and the archimedean terms do not change. The local index at $w$ of $\mathfrak p_w^{n_w}\Lambda'_w$ in $\Lambda_w$ has logarithm $\big(\mathrm{len}_w(\Lambda_w/\Lambda'_w)+dn_w\big)\log N\mathfrak p_w$, because $\Lambda'_w/\mathfrak p_w^{n_w}\Lambda'_w$ has length $dn_w$ for $n_w\ge0$, and for $n_w<0$ the module $\mathfrak p_w^{n_w}\Lambda'_w$ contains $\Lambda'_w$ with a quotient of length $-dn_w$, and lengths add along the chain $\mathfrak p_w^{n_w}\Lambda'_w\supset\Lambda'_w\subset\Lambda_w$ in the generalized index.
\end{proof}

\paragraph{The data at $u$.} Fix $u_0\in A_0$ and an index $u$ of the family $W$ of Proposition~\ref{acct}. Sections~\ref{sec:lifts}, \ref{r5:sec:localterms} and~\ref{sec:proofA} give the following objects, listed in Subsection~\ref{r2:sec:lossinputs} as (L4) to (L8): a hermitian lattice $M_u$ over $\cO_\bK$ of rank $d_M$ with slope $\mu^M_u=\mu(M_u)$, a saturated sublattice $M^{\mathrm{ram}}_u$ of rank $d_u$, an auxiliary prime $T_u$, the structure (F1) to (F5) at every place $w$ above a prime $p\in\mathcal P$ with $p\le X''$ and $p\ne T_u$, a sublattice $N^T_w\subset M^{\mathrm{ram}}_{u,w}$ at every place $w\mid T_u$, and real functions $F_j(w)$ and $E_{u,j}(w)$ of the finite places, for $1\le j\le m$, with finite supports. Put
\[
\mathcal L=\sum_w\sigma\Big(w,\sum_{j=1}^mF_j(w)\Big),\qquad E_u=\sum_w\sigma\Big(w,\sum_{j=1}^mE_{u,j}(w)\Big),
\]
with $\sigma(w,y)=2\ell_we_wy$ as in Remark~\ref{r2:placeform}. Let $W_u$ be the finite set of the places $w$ above the primes of $\mathcal P$ up to $X''$, above $T_u$, and in the supports of the $F_j$ and $E_{u,j}$, and for $w\in W_u$ with residue characteristic $p_w$ put
\[
s_u(w)=\sum_{j=1}^m\Big(\mathrm{gain}^{\mathcal D}_j(p_w)+F_j(w)-E_{u,j}(w)\Big),
\]
where $\mathrm{gain}^{\mathcal D}_j(p)=0$ for $p\notin\mathcal P$. Here the weights are the weights $a_t$ fixed in Subsection~\ref{r2:sec:loss-setting} and the bottom weights $b_t$ of the nodes.

\begin{definition}\label{r2:Nu}
Let $\Lambda_u\subset M^{\mathrm{ram}}_u$ be the lattice whose completion at $w$ is $\NP_w$ (Theorem~\ref{r2:colength}) at the places $w$ above primes $p\in\mathcal P$ with $p\le X''$, $p\ne T_u$ and $\cV_p\ne\emptyset$, is $N^T_w$ at the places above $T_u$, and is $M^{\mathrm{ram}}_{u,w}$ at every other place. Put
\[
N_u=\prod_{w\in W_u}\mathfrak p_w^{-\lfloor e_ws_u(w)\rfloor}\,\Lambda_u,\qquad \mu^N_u=\mu(N_u),
\]
with the metrics of $M_u$.
\end{definition}

The lattice $\Lambda_u$ exists because a lattice is determined by its completions, which may be prescribed at finitely many places. The $\cO_\bK$-module $\Lambda_u$ is a sublattice of $M^{\mathrm{ram}}_u$ of the same rank $d_u$.

\begin{theorem}\label{r2:A2}
Suppose the inputs (L1) to (L8) and the size conditions (Z) of Subsection~\ref{r2:sec:lossinputs}, and that $\mathrm{Sig}(p)\le\frac12$ whenever $\cV_p\ne\emptyset$. Let $\lambda_T=\frac14C_TN_Tm(a/k+b')$ be the number of (L7). Then
\begin{enumerate}
\item[(i)] $\mu^M_u+\Gamma+\mathcal L-E_u-e_2\le\mu^N_u$ with
\[
e_2=4\,\mathrm{IO}+\mathrm{LT}+\mathrm{RX}+1,\quad\mathrm{LT}=4\lambda_T\log T_u,\quad\mathrm{RX}=2\sum_{w\in W_u}\ell_w\le4X''+2\log T_u+2X'' ;
\]
\item[(ii)] $\mu_{\max}(N_u)\le\mu_{\max}(M_u)+\Gamma+\mathcal L-E_u$.
\end{enumerate}
\end{theorem}

\begin{proof}
Write $W_T$ for the places above $T_u$, $\tau(w)=\tau(p_w)$ and $\mathrm{io}(w)=\mathrm{io}(p_w)$ for $p_w\in\mathcal P$, $p_w\le X''$, and $\tau(w)=\mathrm{io}(w)=0$ otherwise. Since $T_u\ge X'$ by (L6), the prime $T_u$ is not a candidate, so $\cV_{T_u}=\emptyset$ and $\tau=\mathrm{io}=0$ on $W_T$. Let $\mathrm{len}(w)$ be the length of $M^{\mathrm{ram}}_{u,w}/\Lambda_{u,w}$ divided by $e_w$, and put $s'(w)=\lfloor e_ws_u(w)\rfloor/e_w$, so that $0\le s_u(w)-s'(w)<1/e_w$.

\emph{Step 1, the slope of $N_u$.} Lemma~\ref{r2:degree} with $\Lambda=M^{\mathrm{ram}}_u$, $\Lambda'=\Lambda_u$ and $n_w=-e_ws'(w)$ gives the normalized degree difference. Multiplication by $2/d_u$ converts it to the slope difference and gives
\[
\mu^N_u=\mu(M^{\mathrm{ram}}_u)+\sum_{w\in W_u}2\ell_we_w\Big(s'(w)-\frac{\mathrm{len}(w)}{d_u}\Big).
\]
\emph{Step 2, the colengths.} For $w\notin W_T$ with $\cV_{p_w}\ne\emptyset$, Theorem~\ref{r2:colength}(b) gives $\mathrm{len}(w)\le d_M(\tau(w)+\mathrm{io}(w))$, using $v_w(p)=e_w$. At the other places $w\notin W_T$ the colength is $0$. For $w\in W_T$, (L7) gives $\mathrm{len}(w)\le d_M\lambda_T$. Put $\mathrm{de}=1-d_u/d_M$. By (L5), $d_u=d_M(1-1/(T_u-1))^f$, so Bernoulli's inequality gives $0\le\mathrm{de}\le f/(T_u-1)\le2f/T_u\le\frac12$ by (Z). Hence $d_M/d_u=1/(1-\mathrm{de})\le2$ and
\[
\sum_{w\notin W_T}\frac{2\ell_we_w\mathrm{len}(w)}{d_u}\le\frac1{1-\mathrm{de}}\sum_w2\ell_we_w\big(\tau(w)+\mathrm{io}(w)\big)=\frac{\mathrm{Ta}+2\,\mathrm{IO}}{1-\mathrm{de}},
\]
where $\mathrm{Ta}=\sum_w2\ell_we_w\tau(w)=\sum_p2\log p\,\tau(p)$ and $\sum_w2\ell_we_w\mathrm{io}(w)=2\,\mathrm{IO}$, both by Remark~\ref{r2:placeform}. Also $\sum_{w\in W_T}2\ell_we_w\mathrm{len}(w)/d_u\le2\lambda_T\cdot2\log T_u=\mathrm{LT}$.

\emph{Step 3, rounding.} $\sum_w2\ell_we_w(s_u(w)-s'(w))\le\sum_{w\in W_u}2\ell_w=\mathrm{RX}$. For a prime $p$, $\sum_{w\mid p}\ell_w=\sum_{w\mid p}f_w\log p/[\bK:\mathbb Q]\le\log p$, so $\mathrm{RX}\le2\sum_{p\le X''}\log p+2\log T_u+2\sum_{p\in P_{FE}}\log p$, where $P_{FE}$ is the set of the residue characteristics of the supports of the $F_j$ and $E_{u,j}$. Chebyshev's bound $\sum_{p\le y}\log p\le2y$ and (L8) give the bound for $\mathrm{RX}$.

\emph{Step 4, the gain.} By Remark~\ref{r2:placeform} and the definitions,
\[
\sum_w2\ell_we_w\big(s_u(w)-\tau(w)\big)=\Gamma+\mathcal L-E_u .
\]
Steps 1 to 4 give
\[
\mu^N_u\ge\mu(M^{\mathrm{ram}}_u)+\Gamma+\mathcal L-E_u-\mathrm{RX}-\mathrm{LT}-\frac{2\,\mathrm{IO}}{1-\mathrm{de}}-\frac{\mathrm{de}}{1-\mathrm{de}}\mathrm{Ta},
\]
because $\mathrm{Ta}/(1-\mathrm{de})=\mathrm{Ta}+\mathrm{de}\,\mathrm{Ta}/(1-\mathrm{de})$. Now $2\,\mathrm{IO}/(1-\mathrm{de})\le4\,\mathrm{IO}$, and $\mathrm{de}\,\mathrm{Ta}/(1-\mathrm{de})\le2\,\mathrm{de}\,b^2\le4fb^2/T_u$ by Lemma~\ref{r2:tau} and (Z). By (L5), $\mu(M^{\mathrm{ram}}_u)\ge\mu^M_u-f\log(T_u-1)/(T_u-2)$, and (Z) gives $4fb^2/T_u+f\log(T_u-1)/(T_u-2)\le1$. This proves (i).

(ii) Every place $w$ with $\cV_{p_w}\ne\emptyset$ has $\Lambda_{u,w}=\NP_w\subset p^{\tau(w)}M^{\mathrm{ram}}_{u,w}=\mathfrak p_w^{e_w\tau(w)}M^{\mathrm{ram}}_{u,w}$ by Theorem~\ref{r2:colength}(a), and at the other places $\tau(w)=0$ and $\Lambda_{u,w}\subset M^{\mathrm{ram}}_{u,w}$. Hence $N_u\subset\mathfrak bM^{\mathrm{ram}}_u$ with $\mathfrak b=\prod_w\mathfrak p_w^{e_w\tau(w)-e_ws'(w)}$. A sublattice $N'$ of $N_u$ is a sublattice of $\mathfrak bM^{\mathrm{ram}}_u$, and $\mathfrak b^{-1}N'$ is a sublattice of $M^{\mathrm{ram}}_u\subset M_u$, so Lemma~\ref{r2:degree} gives $\mu(N')=\mu(\mathfrak b^{-1}N')+\sum_w2\ell_we_w(s'(w)-\tau(w))\le\mu_{\max}(M_u)+\sum_w2\ell_we_w(s_u(w)-\tau(w))$. Step 4 finishes the proof.
\end{proof}

\begin{lemma}\label{r2:tau}
Under the size conditions (Z), $\mathrm{Ta}=\sum_{p\in\mathcal P}2\log p\,\tau(p)\le b^2$.
\end{lemma}

\begin{proof}
$\tau(p)=0$ unless $p<X'$ is a candidate. For a candidate, $\tau(p)\le e_p+m\sum_{t_0\in\cV_p}(t_{t_0}+D_p+2)$, and $\sum_{t_0}N(t_0)\le N_T$, $\lvert\cV_p\rvert\le N_T$. By (L2), $e_p\le(C_e+1)fk^2\log_pa$. The definitions of Subsection~\ref{r2:sec:lp} and Lemma~\ref{r2:orbitsize} give $t_{t_0}+D_p+2\le\tau^+\,N(t_0)$ with the number $\tau^+$ of (Z). So $\mathrm{Ta}\le\sum_{p<X'}2\log p\big((C_e+1)fk^2\log_pa+mN_T\tau^+\big)\le4X'\big((C_e+1)fk^2\log a+mN_T\tau^+\big)$, which is at most $b^2$ by (Z).
\end{proof}

Theorem~\ref{r2:A2}(i) is hypothesis \eqref{A2} of Proposition~\ref{acct} for the index $u$, with the error $e_2$ above. The term $4\,\mathrm{IO}$ is bounded in Subsection~\ref{r2:sec:gainprime}, and Subsection~\ref{r2:sec:lossinputs} bounds $\mathrm{LT}$ and $\mathrm{RX}$ by $x^{C\varepsilon}\mathcal B$. The next subsection proves \eqref{A1} for the same number $\Gamma$.
\subsection{The gain against the sum over primes, and hypothesis \texorpdfstring{\eqref{A1}}{(A1)}}\label{r2:sec:gainprime}

This subsection compares, prime by prime, the contribution $2\log p\sum_{j,t}\mathrm{ask}_{j,t}(p)$ of $p$ to $S$ with its contribution $2\log p\big(\sum_j\mathrm{gain}^{\mathcal D}_j(p)-\tau(p)\big)$ to $\Gamma$. The bottom weights $b_t$ are those of the nodes, so that (W1) holds with them.

\paragraph{Local losses.} Put $\ell(p)=\log_p(2a)+1$ and $\mathrm{cp}(p)=5c_G\,\ell(p)$. Let $N^\sharp=2x^{2+8\varepsilon}$, which bounds $n_p$ and $N(t_0)$, let $c_0=2b'+2c_G$, and put
\begin{align*}
\Theta^{\mathrm{II}}&=r_0(2b'+c_G)+\frac{5r_0C_{31}\log x}{\delta_1}\,(w''+\eta+c_G+N_m)+K_{tII}r_0(c_0+kQ_{p_*})\\ &\quad {}+40C_{31}\log x\big(K_{DII}(\nu_1^3+r_0^5+k^2)\log a+r_0c_0\big)+4,\\
c_6'&=320C_{31}\log x\big(K_{DI}x^{26}(\nu_1^5+k^3)\log a+1\big)\\ &\quad {}+40K_{\kappa2}C_{31}k\log a\,\log x+\frac{164\,C_{31}c_G\log x}{\delta_1}+8,\\
\Theta^{\mathrm I}(p)&=\frac{b'}x+\eta+\frac{(K_{\mathrm{Pr}}+2)C_{31}\log x}{\delta_1}(w''+\eta+c_G+N_m)\\ &\quad {}+K_{\kappa2}\big(N^\sharp Q_pc_6'+kQ_p+k\log a+c_G\big)+40K_{DI}C_{31}x^{26}(\nu_1^5+k^3)\log a\log x+4 .
\end{align*}
For a prime $p\in\mathcal P$ of class I with $p<X'$ put
\[
\mathrm{lo}(p)=\begin{cases}\mathrm{cp}(p)+\Theta^{\mathrm{II}}+K_{tII}r_0\Big(\dfrac{\mathrm{La}\,a}{pk}+a\,p^{-L_1(p)}\Big),&p<p_*,\\[1.5ex]
\mathrm{cp}(p)+\Theta^{\mathrm I}(p)+K_{\kappa1}\nu_1\Big(\dfrac a{p^2}+\dfrac{\mathrm{La}\,a}{pk}\Big),&p\ge p_*,\end{cases}
\]
and $\mathrm{lo}(p)=\mathrm{cp}(p)$ for the other $p\in\mathcal P$. Put $\mathrm{fl}(p)=\sum_{t_0\in R_p}\mathrm{dl}(t_0)+e_p$.

\begin{lemma}\label{r2:shift}
Let $p\in\mathcal P$, $1\le j\le m$, $s\in T$ and $e$ odd, and put $\rho=\pi_p(s,e)$, $A=a_s/p^e$. Then
\[
\begin{aligned}
\max\big(0,A-w_j(\rho)\big)&\le\max\Big(0,\frac{\mathrm{al}(s)}{p^e}-\mathrm{cap}_j(s,e)\Big)+3c_G\,\mathbf 1[A>w_j(\rho)],\\ \max\big(0,A-w_j(\rho)\big)&\le\frac{\mathrm{al}(s)}{p^e}+\frac{2c_G}{p^e}-\min\big(A,w_j(\rho)\big).
\end{aligned}
\]
For fixed $s$ and $j$, the number of odd $e$ with $A>w_j(\rho)$ is at most $\ell(p)$.
\end{lemma}

\begin{proof}
$\mathrm{cap}_j(s,e)=w_j(\rho)-\kappa_j+c_G$, so $A-w_j(\rho)=\mathrm{al}(s)/p^e-\mathrm{cap}_j(s,e)+2c_G/p^e+c_G-\kappa_j$, and $0\le c_G-\kappa_j\le c_G$, $2c_G/p^e\le2c_G$. This gives the first inequality when $A>w_j(\rho)$, and both sides are compared with $0$ otherwise. The second is $\max(0,A-w)=A-\min(A,w)$ and $a_s=\mathrm{al}(s)+2c_G$. Finally $w_j(\rho)\ge1$ and $a_s\le2a$, so $A>w_j(\rho)$ forces $p^e<2a$.
\end{proof}

\begin{proposition}\label{r2:star}
Let $p$ be a candidate with $\mathrm{Sig}(p)\le\theta$, $t_0\in\cV_p$ and $1\le j\le m$. Suppose the weights lie in the windows of Theorem~\ref{r2:LPthm}. Then
\[
\sum_{s\in T(t_0)}\big(\mathrm{ask}_{j,s}(p)-\mathrm{lo}(p)\big)\le g^1_j(t_0)+\mathrm{dl}(t_0)-\mathrm{Pr}_j(t_0)-\lceil t_{t_0}\rceil-D_p-1 .
\]
\end{proposition}

\begin{proof}
Write $N=N(t_0)$. By \eqref{r2:g1},
\[
g^1_j(t_0)+\mathrm{dl}(t_0)=\sum_{s\ \mathrm{even}}\sum_{e\ \mathrm{odd}}\frac{\mathrm{al}(s)}{p^e}+\sum_{s\ \mathrm{odd}}\sum_{e\ \mathrm{odd}}\max\Big(0,\frac{\mathrm{al}(s)}{p^e}-\mathrm{cap}_j(s,e)\Big),
\]
since $\gamma_p\,\mathrm{al}(s)=\sum_{e\ \mathrm{odd}}\mathrm{al}(s)/p^e$ and $y-\min(y,c)=\max(0,y-c)$. For an odd top $s$, the first inequality of Lemma~\ref{r2:shift} gives $\mathrm{ask}_{j,s}(p)\le\sum_e\max(0,\mathrm{al}(s)/p^e-\mathrm{cap}_j(s,e))+3c_G\ell(p)$. For an even top $s$, the second gives
\[
\mathrm{ask}_{j,s}(p)\le\sum_e\frac{\mathrm{al}(s)}{p^e}+2c_G\gamma_p-\sum_e\min\Big(\frac{a_s}{p^e},w_j(\pi_p(s,e))\Big).
\]
Let $y$ be an odd top with $d(y)>0$ and let $s=y-d(y)h_p$, an even top. Then $\pi_p(s,d(y))$ is the top place of $y$, where $w_j=b_y+o_j(y)\ge b_y$ by (W1). The pairs $(s,d(y))$ are distinct for distinct $y$. Since $2c_G\gamma_p\le2c_G$, summing over the tops gives
\begin{equation}\label{r2:kink}
\sum_{s\in T(t_0)}\mathrm{ask}_{j,s}(p)\le g^1_j(t_0)+\mathrm{dl}(t_0)+5c_GN\ell(p)-\sum_{y\ \mathrm{odd},\ d(y)>0}\min\Big(b_y,\frac{a_{y-d(y)h_p}}{p^{d(y)}}\Big).
\end{equation}
It remains to show $\mathrm{Pr}_j+\lceil t_{t_0}\rceil+D_p+1\le N(\mathrm{lo}(p)-\mathrm{cp}(p))+\Sigma_{\mathrm{kink}}$, where $\Sigma_{\mathrm{kink}}$ is the last sum in \eqref{r2:kink}, which is nonnegative. We use $n_p\le C_{31}N\log x/\delta_1$ (Lemma~\ref{r2:orbitsize}), $N^+=40n_p\delta_1\le40C_{31}N\log x$, $\lceil y\rceil\le y+1$, $\log_pa\le\log a$, $N\ge1$, and that $Q_p=p^{G'}$ increases with $p$.

\emph{Case $p<p_*$.} By \eqref{r2:price}, $\mathrm{Pr}_j\le N\big(r_0(2b'+c_G)+r_0C_{31}\log x(w''+\eta+c_G+N_m)/\delta_1\big)$. The definitions of $t^{\mathrm{II}}$ and $D^{\mathrm{II}}$ give
\[
\begin{aligned}
\lceil t^{\mathrm{II}}\rceil+D^{\mathrm{II}}+1&\le N\Big(K_{tII}r_0\Big(\frac{\mathrm{La}\,a}{pk}+ap^{-L_1(p)}\Big)+K_{tII}r_0(c_0+kQ_{p_*})\\ &\qquad {}+40C_{31}\log x\big(K_{DII}(\nu_1^3+r_0^5+k^2)\log a+r_0c_0\big)+\frac{4r_0C_{31}c_G\log x}{\delta_1}+3\Big).
\end{aligned}
\]
The sum of the two bounds is at most $N(\mathrm{lo}(p)-\mathrm{cp}(p))$.

\emph{Case $p\ge p_*$.} By \eqref{r2:price} and (W1), $\hat w_j(0,0,s)-b_s=o_j(s)-\kappa_j\le b'/x+\eta$, so $\mathrm{Pr}_j\le N(b'/x+\eta)+K_{\mathrm{Pr}}C_{31}N\log x(w''+\eta+c_G+N_m)/\delta_1$. In $t^{\mathrm I}$ the first sum equals $\Sigma_{\mathrm{kink}}$. In the rest, $N^+\le40C_{31}N\log x$, $n_p\le C_{31}N\log x/\delta_1$, $\delta_1\le1$ and $D^{\mathrm I}\le40K_{DI}C_{31}N\log x\,x^{26}(\nu_1^5+k^3)\log a+1$ give $c_6\le Nc_6'$, so $K_{\kappa2}N\,Q_pc_6\le K_{\kappa2}N\,N^\sharp Q_pc_6'$, and $2n_pc_G\le2C_{31}Nc_G\log x/\delta_1$. Finally $D^{\mathrm I}+2\le40K_{DI}C_{31}N\log x\,x^{26}(\nu_1^5+k^3)\log a+3$. Collecting terms, $\mathrm{Pr}_j+\lceil t^{\mathrm I}\rceil+D^{\mathrm I}+1\le\Sigma_{\mathrm{kink}}+N\big(\Theta^{\mathrm I}(p)+K_{\kappa1}\nu_1(a/p^2+\mathrm{La}\,a/(pk))\big)$.
\end{proof}

\begin{lemma}\label{r2:noop}
Let $p\in\mathcal P$ with $p\le X''$ be of class II or satisfy $p\ge X'$, let $t_0\in R_p$ and $1\le j\le m$. Then $\sum_{s\in T(t_0)}\big(\mathrm{ask}_{j,s}(p)-\mathrm{cp}(p)\big)\le\max(g^0_j(t_0),0)+\mathrm{dl}(t_0)$.
\end{lemma}

\begin{proof}
If $p$ is of class II, $\pi_p(s,e)\in\Phi_1$ for odd $e$, so no pair $(s,e)$ satisfies $\mathrm{Excl}(s,e)$. If $p\ge X'$ and $\pi_p(s,e)$ is the top place of some $s'\in T$, then (W1) and \eqref{scale} give $w_j(\pi_p(s,e))\ge b'+b'/x$ while $a_s/p^e\le(a+a/(kx^{13}))/X'=b'(1+1/(kx^{13}))$, so the term of $(s,e)$ in $\mathrm{ask}_{j,s}(p)$ vanishes. In both cases the excluded pairs contribute $0$ to $\mathrm{ask}$, and the first inequality of Lemma~\ref{r2:shift}, summed over the other pairs, gives $\sum_s\mathrm{ask}_{j,s}(p)\le g^0_j(t_0)+\mathrm{dl}(t_0)+3c_GN(t_0)\ell(p)$.
\end{proof}

\paragraph{Classes of primes.} Split the primes $p\in\mathcal P$ into five sets. The \emph{high} primes are those with $p>X''$. Among the primes $p\le X''$, the \emph{served} primes are the candidates with $\mathrm{Sig}(p)\le\theta$, the \emph{second} primes are the candidates with $\mathrm{Sig}(p)>\theta$, the \emph{plain} primes are the primes $p\notin\mathcal D$ of class II or with $p\ge X'$, and the \emph{first} primes are the remaining primes of class I. A first prime $p$ is not a candidate. If $p<X'$, then $n_p<x^{\varepsilon/4}$, or $p\ge p_*$ and $\lvert p^i-X'\rvert\le\rho_BX'$ for some $i\ge1$. If $p\ge X'$, then $p\in\mathcal D$, so $h_p$ has odd order and $h_p\ne0$, which gives $n_p\le o<x^{\varepsilon/4}$ by Lemma~\ref{r2:evenorder}, or $p<X'(1+C_r\rho_B)$ and one of the two alternatives above holds. For a served prime let $T^{\mathrm{gd}}(p)$ be the set of the tops on good orbits and $\beta_p=N_T-\lvert T^{\mathrm{gd}}(p)\rvert$.

\begin{proposition}\label{r2:perprime}
Let $p\in\mathcal P$ with $p\le X''$, and suppose the weights lie in the windows. Then $0\le\mathrm{gain}_j(p)$, and
\begin{enumerate}
\item[(a)] if $p$ is served, $\sum_j\sum_{s\in T^{\mathrm{gd}}(p)}\big(\mathrm{ask}_{j,s}(p)-\mathrm{lo}(p)\big)-m\,\mathrm{fl}(p)\le\sum_j\mathrm{gain}^{\mathcal D}_j(p)-\tau(p)$;
\item[(b)] if $p$ is plain, $\sum_j\sum_{s\in T}\big(\mathrm{ask}_{j,s}(p)-\mathrm{lo}(p)\big)-m\,\mathrm{fl}(p)\le\sum_j\mathrm{gain}^{\mathcal D}_j(p)-\tau(p)$.
\end{enumerate}
\end{proposition}

\begin{proof}
Each term of $\mathrm{gain}_j(p)$ is a maximum with $\max(g^0_j,0)\ge0$. (a) A served prime has $\cV_p$ equal to the set of good orbits, $p\notin\mathcal D$ and $\mathrm{gain}^{\mathcal D}=\mathrm{gain}$. Then $\mathrm{gain}_j(p)\ge\sum_{t_0\in\cV_p}\big(g^1_j(t_0)-\mathrm{Pr}_j(t_0)\big)$ because the terms of the unserved orbits are nonnegative. Proposition~\ref{r2:star}, summed over $t_0\in\cV_p$ and $j$, and $\tau(p)=e_p+\sum_{t_0\in\cV_p}m(\lceil t_{t_0}\rceil+D_p+1)$ give the claim, since $\sum_{t_0\in\cV_p}m\,\mathrm{dl}(t_0)+e_p\le m\,\mathrm{fl}(p)$. (b) A plain prime has $\cV_p=\emptyset$, so $\tau(p)=0$ and $\mathrm{gain}_j(p)=\sum_{t_0\in R_p}\max(g^0_j(t_0),0)$, unless $p$ is of class I with $h_p=0$, in which case $\mathrm{gain}_j(p)=0$, $p\ge X'$ and every $\pi_p(s,e)$ with $e$ odd is the top place of $s$, so every $\mathrm{ask}_{j,s}(p)$ vanishes by the argument of Lemma~\ref{r2:noop}. In the other cases Lemma~\ref{r2:noop}, summed over $t_0\in R_p$ and $j$, gives the claim, since $\mathrm{lo}(p)=\mathrm{cp}(p)$ for a plain prime.
\end{proof}

\paragraph{The total.} For $p\in\mathcal P$ put $\mathrm{dev}(p)=\sum_{j,t}\mathrm{ask}_{j,t}(p)-\sum_j\mathrm{gain}^{\mathcal D}_j(p)+\tau(p)$, so that $S-\Gamma=\sum_p2\log p\,\mathrm{dev}(p)+\sum_{j,t}\mathrm{dt}_{j,t}$ by \eqref{r2:Ssplit} and \eqref{r2:Gamma}. For $p\ge3$ we have $\gamma_p\le\frac9{8p}$, and $a_t\le\frac{17}{16}a$, so
\begin{equation}\label{r2:askbound}
\mathrm{ask}_{j,t}(p)\le a_t\gamma_p\le\frac{2a}p .
\end{equation}
Define the weighted count of second primes and the correction
\[
\mathrm{lhsW}=\sum_{p\ \mathrm{candidate}}\frac{\log p}p\Big(2\cdot\mathbf 1[\mathrm{Sig}(p)>\theta]+4\,\mathrm{Sig}(p)\Big).
\]

\begin{theorem}\label{r2:A1}
Suppose the inputs (L1) to (L8), the size conditions (Z) of Subsection~\ref{r2:sec:lossinputs}, that the weights lie in the windows, and that $\mathrm{lhsW}\le x^{-9}$. Then
\[
S\le\Gamma+e_1,\qquad e_1=2\log4\;N_Tmb+B_{\mathrm{minor}}\le(2\log4+1)\,N_Tmb,\qquad\mathrm{IO}\le N_Tmb,
\]
where $B_{\mathrm{minor}}$ is the sum of the bucket bounds (B1) to (B9) below.
\end{theorem}

\begin{proof}
We bound $\mathrm{dev}(p)$ class by class, using $\mathrm{gain}^{\mathcal D}_j\ge0$ and $\tau\ge0$.

\emph{Served primes.} By Proposition~\ref{r2:perprime}(a) and \eqref{r2:askbound}, $\mathrm{dev}(p)\le mN_T\mathrm{lo}(p)+m\,\mathrm{fl}(p)+2ma\beta_p/p$.

\emph{Plain primes.} By Proposition~\ref{r2:perprime}(b), $\mathrm{dev}(p)\le mN_T\mathrm{cp}(p)+m\,\mathrm{fl}(p)$.

\emph{Second primes.} Such a prime is a candidate with $\mathrm{Sig}(p)>\theta$, hence $p<X'$, $p$ has no gain and $p\ne T_u$ because $T_u\ge X'$, so $p\in\mathcal D$ and $\cV_p=\emptyset$. Then $\mathrm{dev}(p)\le2amN_T/p$.

\emph{First and high primes.} $\tau(p)=0$, since such a prime is not a candidate. So $\mathrm{dev}(p)\le\sum_{j,t}\mathrm{ask}_{j,t}(p)$. For a first prime \eqref{r2:askbound} gives $\mathrm{dev}(p)\le2amN_T/p$. For a high prime and odd $i\ge3$, $a_t/p^i\le2a/X''^3<w''$, and for $i=1$, $a_t/p-w_j\le a(1+\vartheta)/p-w''$ with $\vartheta=1/(kx^{13})$, which is positive when $p<X''(1+\vartheta)$ and nonpositive otherwise, and then it is at most $w''\vartheta$ since $X''w''=a$.

Summing with the weights $2\log p$ gives $S-\Gamma\le\sum_{i=0}^{8}B_i$ with
\begin{align*}
B_0&=2mN_T\sum_{p\ \mathrm{served}}\log p\,\Big(\frac{b'}x\Big),&
B_1&=2mN_T\sum_{p\le X''}\log p\,\mathrm{cp}(p),\\
B_2&=2mN_T\sum_{p<X'}\log p\,\big(\mathrm{lo}(p)-\mathrm{cp}(p)-\mathbf 1[p\ge p_*]\tfrac{b'}x\big),&
B_3&=2m\sum_{p\le X''}\log p\,\mathrm{fl}(p),\\
B_4&=4am\sum_{p\ \mathrm{served}}\frac{\log p}p\beta_p,&
B_5&=4amN_T\sum_{p\ \mathrm{second}}\frac{\log p}p,\\
B_6&=4amN_T\sum_{p\ \mathrm{first}}\frac{\log p}p,&
B_7&=2mN_T\,w''\vartheta\sum_{X''<p<X''(1+\vartheta)}\log p,\\
B_8&=\sum_{j,t}\mathrm{dt}_{j,t}.
\end{align*}
where in $B_2$ the term $b'/x$ of $\Theta^{\mathrm I}(p)$ is moved to $B_0$ for $p\ge p_*$, and the term $\Theta^{\mathrm{II}}$ stays in $B_2$. We bound each.

(B0) By the Chebyshev bound $\sum_{p\le y}\log p\le y\log4$ and $X'b'/x=b$, $B_0\le2\log4\;N_Tmb$.

(B4) By (L9), $\sum_p\frac{\log p}p\beta_p\le N_Tx^{-2+\varepsilon}$, so $B_4\le4amN_Tx^{-2+\varepsilon}=4N_Tmb\,x^{-1+\varepsilon}$.

(B5) Each second prime contributes $2\log p/p$ to $\mathrm{lhsW}$, so $B_5\le2amN_T\,\mathrm{lhsW}\le2N_Tmb\,x^{-8}$.

(B6) By Lemma~\ref{r2:firstcount}, $\sum_{p\ \mathrm{first}}\log p/p\le B^{\mathrm{ord}}+B^{\mathrm{band}}$, and (Z) bounds $4amN_T(B^{\mathrm{ord}}+B^{\mathrm{band}})$.

(B7) The interval contains at most $\vartheta X''+1$ integers, so $B_7\le2mN_T\log(2X'')(a\vartheta^2+w''\vartheta)$.

(B8) $\mathrm{dt}_{j,t}\le\log d\sum_{i\ge1}a_t/d^i\le2a\log d/d$, so $B_8\le2amN_T\log d/d$.

(B1), (B2), (B3) are explicit sums over primes, bounded by (Z). Hence $S-\Gamma\le2\log4\,N_Tmb+B_{\mathrm{minor}}$ with $B_{\mathrm{minor}}=B_1+\dots+B_8$, and (Z) gives $B_{\mathrm{minor}}\le N_Tmb$.

For $\mathrm{IO}$: for a served prime and $t_0\in\cV_p$, $S_{t_0}\le\sum_j\max(0,g^1_j(t_0))+1\le m\gamma_p\sum_{s\in T(t_0)}\mathrm{al}(s)+1\le2mN_Ta/p$, using \eqref{r2:askbound} and $mN_Ta/p\ge b'\ge1$. So $\mathrm{io}(p)\le4\,\mathrm{Sig}(p)\,mN_Ta/p$, and $\mathrm{IO}\le amN_T\,\mathrm{lhsW}\le N_Tmb\,x^{-8}$.
\end{proof}

\begin{lemma}\label{r2:firstcount}
Let $J=x^{\varepsilon/4}$ and $q=2^{e+2}$ with $e$ as in Proposition~\ref{field}. Then
\[
\sum_{p\ \mathrm{first}}\frac{\log p}p\le B^{\mathrm{ord}}+B^{\mathrm{band}},\qquad B^{\mathrm{ord}}=4J\Big(\frac{\log P_0}{P_0}+\frac{\log^2X''}q\Big),\quad B^{\mathrm{band}}=I_b\big(4\rho_B+p_*^{-1}\big)\log(4X'),
\]
where $P_0=2^{(e+3)/2}J^{-1/2}$, $I_b=\lceil\log(2X')/\log p_*\rceil$, and we assume $P_0\ge3$ and $X''\ge q$.
\end{lemma}

\begin{proof}
\emph{Small order.} By Proposition~\ref{field}(iii), $h_p$ maps to the class of $p$ in $(\mathbb Z/q)^\times/\{\pm1\}$, a cyclic group of order $2^e$. If $n_p<J$, the image has order $2^i$ with $2^i<J$, so $p^{2^i}\equiv\pm1\pmod q$. For $i\ge1$, $p^{2^i}\equiv1\pmod8$ and $-1\not\equiv1\pmod8$, so $p^{2^i}\equiv1\pmod q$. The $2$-adic valuation satisfies $v_2(p^{2^i}-1)=v_2(p^2-1)+i-1$ for $i\ge1$, by induction from $p^{2^{i+1}}-1=(p^{2^i}-1)(p^{2^i}+1)$ and $p^{2^i}+1\equiv2\pmod4$. So $v_2(p^2-1)\ge e+3-i$ and $p^2>2^{e+3-i}$, which gives $p>P_0$. For $i=0$, $p\ge q-1\ge P_0$. The solutions of $y^{2^i}=1$ in $(\mathbb Z/q)^\times\cong\mathbb Z/2\times\mathbb Z/2^e$ number at most $2^{i+1}$, so the primes in question lie in at most $\sum_{2^i<J}2^{i+1}\le4J$ residue classes modulo $q$, and are at least $P_0$. For one class $c$, the function $y\mapsto\log y/y$ decreases for $y\ge3$, so the sum of $\log n/n$ over the integers $P_0\le n\le X''$ with $n\equiv c\pmod q$ is at most $\log P_0/P_0+\frac1q\int_{P_0}^{X''}\frac{\log y}y\,dy\le\log P_0/P_0+\log^2X''/(2q)$.

\emph{Band.} If $p\ge p_*$, $p\le X''$ and $\lvert p^i-X'\rvert\le\rho_BX'$ with $\rho_B\le\frac13$, then $i\le I_b$ and $p$ lies in $[y_i,y_i(1+3\rho_B)]$ with $y_i=(X'(1-\rho_B))^{1/i}\ge p_*$, since $((1+\rho_B)/(1-\rho_B))^{1/i}\le1+3\rho_B$. This interval contains at most $3\rho_By_i+1$ integers, and a prime $p$ in it with $p\ge p_*$ contributes $\log p/p\le\log(4X')/\max(y_i,p_*)$. So the band primes for one $i$ contribute at most $(3\rho_By_i+1)\log(4X')/\max(y_i,p_*)\le(3\rho_B+1/p_*)\log(4X')$.
\end{proof}

\paragraph{The choice of weights.} The weights enter $\mathrm{lhsW}$ through the shares $\kappa_{t_0}(a,b)$. Let $\mathrm{Box}$ be the set of integer weight vectors $(a_t,b_t)_{t\in T}$ in the windows $a\le a_t\le a+a/(kx^{13})$ and $b'\le b_t\le b'+2n\delta_1x^\varepsilon\mathrm{La}\,b'/k$.

\begin{proposition}\label{r2:weights}
Under the hypotheses of Corollary~\ref{r2:meankappa} and the size condition (Z9), there is $(a,b)\in\mathrm{Box}$ with $\mathrm{lhsW}\le x^{-9}$.
\end{proposition}

\begin{proof}
For $\theta=x^{-\varepsilon}$ and $\mathrm{Sig}\ge0$, $2\cdot\mathbf 1[\mathrm{Sig}>\theta]+4\,\mathrm{Sig}\le(2/\theta+4)\mathrm{Sig}$. The share $\kappa_{t_0}$ depends on the weights at the tops of the orbit of $t_0$, and $\mathrm{Box}$ is a product of windows, so the mean of $\kappa_{t_0}$ over $\mathrm{Box}$ is its mean over the windows of the orbit, which is at most $N(t_0)x^{-12}$ by Corollary~\ref{r2:meankappa}. The orbits at $p$ partition $T$, so the mean of $\mathrm{Sig}(p)$ is at most $mN_Tx^{-12}$, and by Mertens' bound $\sum_{p\le y}\log p/p\le\log y+2$ the mean of $\mathrm{lhsW}$ is at most $(2x^\varepsilon+4)mN_Tx^{-12}(\log X'+2)$. This is at most $x^{-9}$ by (Z9), and some element of $\mathrm{Box}$ is at most the mean.
\end{proof}

Proposition~\ref{r2:weights} is the reason the bottom weights are chosen together with the weights $a_t$. Input (L6) asks that the nodes of Section~\ref{sec:proofA} have the bottom weights of Proposition~\ref{r2:weights}. With these weights Theorem~\ref{r2:A1} is hypothesis \eqref{A1} of Proposition~\ref{acct}, and it bounds the term $4\,\mathrm{IO}$ of $e_2$ in Theorem~\ref{r2:A2} by $4N_Tmb$.
\subsection{What this section takes from the other steps}\label{r2:sec:lossinputs}

The band width and its constant are $\rho_B=x^{-2}$ and $C_r=1$. The proofs above use the following inputs from Sections~\ref{sec:lifts}, \ref{r5:sec:localterms} and~\ref{sec:proofA}, and the size conditions (Z) listed after them.

\begin{enumerate}
\item[(L1)] \emph{Local structure (Section~\ref{sec:lifts}).} At every place $w$ above a prime $p\in\mathcal P$ with $p\le X''$ and $p\ne T_u$ there are data (F1) to (F5) of Subsection~\ref{r2:sec:colength}, in which the coordinates indexed by $v\in\cV_p$ are the orbits of $p$, the representation $V_v$ is the representation of Subsection~\ref{r2:sec:lp} with $n'=n_p$ and the base weights of the orbit, and $y_v=p^{A_r}T_{r,v}$ with $r=r_p$.
\item[(L2)] \emph{Integral structures (Section~\ref{sec:lifts}).} The lattices of (F3) satisfy $p^{e_p}M^{\mathrm{na}}\subset M$ with $e_p\le(C_e+1)fk^2\log_pa$ for a constant $C_e$.
\item[(L3)] \emph{The orbit bounds (Section~\ref{sec:lifts}).} The numbers $g^0_j$, $g^1_j$ and $\mathrm{Pr}_j$ of Subsection~\ref{r2:sec:loss-setting} are the orbit bounds $g^0$, $g^1$, $\mathrm{PR}$ of Theorem~\ref{lowterm} for the lifts at $w$. This input is used in Section~\ref{r3:sec}, which applies Theorem~\ref{lowterm} to the generators of $\NP_w$ described after Theorem~\ref{r2:colength}. The proofs of this section use the formulas \eqref{r2:g0} to \eqref{r2:price} as definitions.
\item[(L4)] \emph{The base lattice (Sections~\ref{sec:lifts} and~\ref{r4:sec:R4}).} $M_u$ is a hermitian $\cO_\bK$-lattice of rank $d_M$ whose completion at each place $w$ of (L1) is the lattice $M$ of (F3).
\item[(L5)] \emph{The ramified part (Section~\ref{sec:proofA}).} $M^{\mathrm{ram}}_u\subset M_u$ is saturated, its completions at the places of (L1) are the lattices $M^{\mathrm{ram}}$ of (F5), $d_u=d_M\big(1-1/(T_u-1)\big)^f$, and $\mu(M^{\mathrm{ram}}_u)\ge\mu(M_u)-f\log(T_u-1)/(T_u-2)$.
\item[(L6)] \emph{The auxiliary prime and the nodes (Section~\ref{sec:proofA}).} $T_u$ is a prime with $T_u\ge X'$. The nodes satisfy (W1) to (W3), and their bottom weights $b_t$, together with the weights $a_t$, form the element of $\mathrm{Box}$ given by Proposition~\ref{r2:weights}.
\item[(L7)] \emph{The colength at $T_u$ (Section~\ref{sec:proofA}).} At each place $w\mid T_u$ there is a sublattice $N^T_w\subset M^{\mathrm{ram}}_{u,w}$ whose colength, divided by $e_w$, is at most $d_M\lambda_T$ with $\lambda_T=\frac14C_TN_Tm(a/k+b')$.
\item[(L8)] \emph{Local terms (Section~\ref{r5:sec:localterms}).} The functions $F_j$ and $E_{u,j}$ have finite supports, and the set $P_{FE}$ of their residue characteristics has $\sum_{p\in P_{FE}}\log p\le X''$.
\item[(L9)] \emph{Bad orbits (Section~\ref{sec:proofA}).} For the served primes, $\sum_p\frac{\log p}p\beta_p\le N_Tx^{-2+\varepsilon}$.
\item[(L10)] \emph{The engines.} Hypothesis~\ref{r2:engines}, which this section uses through Theorem~\ref{r2:LPthm} and Corollary~\ref{r2:meankappa}.
\end{enumerate}

Proposition~\ref{tops} gives (L9), applied with $\mathcal P$ the set of candidates and $L_{2,p}=2\max_{p'}L_1(p')$. Indeed $L_1(p')\le4\log(kx^{13})/\log2+1$, so $L_{2,p}\le400\log x$ for large $x$, and Proposition~\ref{tops} with $\lvert T\rvert\ge\frac12\delta_1n$ and Mertens' bound gives $\sum_p\frac{\log p}p\beta_p\le8C_bx^{-2}N_T(\log X'+2)$, which is at most $N_Tx^{-2+\varepsilon}$ for large $x$ because $\log X'=x^{\varepsilon/3}$.

\paragraph{Size conditions.} Write $\Pi=\max(\eta,c_G,N_m,t_r,K_\bullet,C_\bullet)$ for the largest of the auxiliary quantities and constants above. The size conditions are the following.
\begin{enumerate}
\item[(Z1)] $x\ge x_0(\varepsilon,o)$, $e\ge1$, $P_0\ge3$, $X''\ge2^{e+2}$, and $\Pi\le\exp(x^{\varepsilon/4})$.
\item[(Z2)] $2f/T_u\le\frac12$ and $4fb^2/T_u+f\log(T_u-1)/(T_u-2)\le1$.
\item[(Z3)] $t_{t_0}+D_p+2\le\tau^+N(t_0)$ for every candidate $p$ and good orbit $t_0$, with $\tau^+=2b'+\max\big(\Theta^{\mathrm{II}},\Theta^{\mathrm I}(X')\big)+(K_{tII}r_0+K_{\kappa1}\nu_1)a$, and $4X'\big((C_e+1)fk^2\log a+mN_T\tau^+\big)\le b^2$.
\item[(Z4)] $B_1+\dots+B_8\le N_Tmb$ for the bucket bounds of Theorem~\ref{r2:A1}, and $\mathrm{LT}+\mathrm{RX}+1\le x^{C\varepsilon}\mathcal B$.
\item[(Z5)] The hypotheses of the engines in Hypothesis~\ref{r2:engines}(A), with $Y=x^{13}$ and $N_s$ the maximal number of tops of a segment.
\item[(Z9)] $(2x^\varepsilon+4)\,mN_Tx^{-12}(\log X'+2)\le x^{-9}$.
\end{enumerate}

Condition (Z9) holds for large $x$ because $m\le x^{1/2+\varepsilon}$ and $N_T\le2\delta_1n\le4x^{2+8\varepsilon}$, so the left side is at most $x^{-19/2+11\varepsilon}$. Condition (Z3) holds for large $x$. The inequality $t_{t_0}+D_p+2\le\tau^+N(t_0)$ is a consequence of the bounds in the proof of Proposition~\ref{r2:star}, because the first sum in $t^{\mathrm I}$ is at most $N(t_0)\max_yb_y\le2N(t_0)b'$. The number $\Theta^{\mathrm I}(X')$ is at most $\exp(Cx^{\varepsilon/3+\varepsilon/8}\log x)$, because $Q_{X'}=X'^{G'}$ with $G'\le Cx^{\varepsilon/8}\log x$. Then $X'mN_T\tau^+\le\exp(Cx^{11\varepsilon/24}\log x)\,a$, which is smaller than $b^2=e^{2x^\varepsilon}$. In (Z4) the bucket bounds compare with $N_Tmb=N_Tm\,e^{x^\varepsilon}$ in the way we now describe, for large $x$. In $B_1$ and $B_3$, the sums over $p\le X''$ are at most $2X''\exp(Cx^{\varepsilon/4})$ times $mN_T$ or $mf$, and $X''\le2x^{3+8\varepsilon}e^{x^{\varepsilon/3}}$, while $f\le4N_T/\delta_1$. In $B_2$, the terms with $p<p_*$ sum to at most $4p_*mN_T\Theta^{\mathrm{II}}+2K_{tII}r_0amN_T\big(\mathrm{La}(\log p_*+2)/k+2p_*(kx^{13})^{-4}\big)$, using $p^{L_1(p)}\ge(kx^{13})^4$, and $\Theta^{\mathrm{II}}\le\exp(Cx^{\varepsilon/4})b'$ with $p_*b'=k^2x^{27}b/X'$. The terms with $p_*\le p<X'$ sum to at most $4X'mN_T\Theta^{\mathrm I}(X')+2K_{\kappa1}\nu_1amN_T\big(2\log p_*/p_*+\mathrm{La}(\log X'+2)/k\big)$, and $a\,\mathrm{La}\log X'/k\le b\,x^{1+3\varepsilon}/x^{9/2}$ because $k\ge x^{9/2}$ and $\mathrm{La}\le Cx^{\varepsilon}\log^2x$. In $B_6$, $4a(B^{\mathrm{ord}}+B^{\mathrm{band}})\le Cb\big(x^{1+\varepsilon/4}\log x/P_0+x^{1+\varepsilon/4+2\varepsilon/3}/2^e+x^{-1+2\varepsilon/3}\big)$ with $P_0\ge2^{e/2}x^{-\varepsilon/8}\ge x^{1+4\varepsilon-\varepsilon/8}/\sqrt o$ and $2^e\ge x^{2+8\varepsilon}/o$. Finally $B_7\le CmN_Tb\log X''/(k^2x^{25})$, and $B_8\le4mN_Tb\log x/x$ because $d\ge x^2$. The second part of (Z4) and condition (Z2) are conditions on $T_u$, $C_T$ and $\mathcal B$ that step (R6) must verify, and (Z5) is part of the engines.


\subsection{The engines of Hecke nondegeneracy}\label{w1:sec:scope}

\paragraph{Valuations.} Throughout, $p$ is an odd prime, $v$ is the valuation on $\bC_p$ with $v(p)=1$, and $v(0)=+\infty$. For an integer $m\ge1$, Legendre's formula gives
\begin{equation}\label{w1:eq:legendre}
v(m!)=\frac{m-s_p(m)}{p-1},\qquad \frac{m}{p-1}-v(m!)=\frac{s_p(m)}{p-1}\le\log_pm+1,
\end{equation}
where $s_p(m)$ is the sum of the base-$p$ digits of $m$. Kummer's theorem gives $v\binom Nk\le\log_pN$ for $0\le k\le N$, since $v\binom Nk$ is the number of carries in the addition of $k$ and $N-k$ in base $p$, and that number is at most the number of digits of $N$ beyond the first.

\paragraph{The Fock form.} Let $\sigma$ be a finite set of variables and $K$ a field of characteristic $0$. For polynomials $F=\sum_mF_mx^m$ and $G=\sum_mG_mx^m$ in $K[x_\sigma]$ the \emph{Fock form} is
\begin{equation}\label{w1:eq:fock}
\langle F,G\rangle=\sum_mF_m\,G_m\,m!,\qquad m!=\prod_{s\in\sigma}m_s!.
\end{equation}
For a variable $x_s$ and monomials $x^m$, $x^{m'}$ one has $\langle x_sx^m,x^{m'}\rangle=\langle x^m,\partial_sx^{m'}\rangle$, because both sides vanish unless $m'=m+e_s$, and then both equal $(m+e_s)!$. So multiplication by $x_s$ and the derivative $\partial_s$ are adjoint for the Fock form.

\subsection{Engine (C): the mean of the segment shares}\label{w1:sec:C}

This subsection proves the analytic core of Hypothesis~\ref{r2:engines}(C). The core is a theorem on the $p$-adic Smith exponents of a matrix of polynomials evaluated along an arithmetic progression (Theorem~\ref{w1:thm:window}), together with an identity that writes the Gram matrix of a segment form at one level as such a matrix, with leading coefficient the untwisted Fock Gram matrix (Proposition~\ref{w1:prop:tie}). Their combination is Theorem~\ref{w1:thm:level}, the bound at one level. Subsection~\ref{w1:sec:Crest} states the steps that lead from the bound at one level to the mean of the segment share.

\subsubsection{Smith exponents and good families}

Let $L$ be a free $\bZ_p$-module with basis $e_1,\dots,e_n$ and let $\Psi\colon L\times L\to\bQ_p$ be bilinear. For $D\in\bZ$ the \emph{good count} $\mathrm{gc}(\Psi,L,D)$ is the largest $N$ for which there are $x_1,\dots,x_N,y_1,\dots,y_N\in L$ with $\Psi(x_i,y_j)=p^D\delta_{ij}$, and the \emph{bad count} is $\mathrm{bad}(\Psi,L,D)=n-\mathrm{gc}(\Psi,L,D)$.

\begin{lemma}\label{w1:lem:smithbad}
Let $s\in\bZ_p\setminus\{0\}$ and suppose that the matrix $A=\big(s\,\Psi(e_a,e_b)\big)_{a,b}$ has entries in $\bZ_p$. Write $A=U\operatorname{diag}(d_1,\dots,d_n)V$ with $U,V\in\mathrm{GL}_n(\bZ_p)$. Then for every integer $D$,
\[
\mathrm{bad}(\Psi,L,D)\le\#\{i:\ v(d_i)\ge D+1\}.
\]
\end{lemma}

\begin{proof}
Let $I=\{i:v(d_i)\le D\}$. For $i\in I$ write $d_i=p^{v(d_i)}u_i$ with $u_i\in\bZ_p^\times$, and put
\[
x_i=\sum_a(U^{-1})_{ia}\,e_a,\qquad y_i=p^{D-v(d_i)}u_i^{-1}\sum_b(V^{-1})_{bi}\,e_b .
\]
Both lie in $L$, because the coefficients are in $\bZ_p$ and $D-v(d_i)\ge0$. For $i,j\in I$,
\[
s\,\Psi(x_i,y_j)=p^{D-v(d_j)}u_j^{-1}\big(U^{-1}AV^{-1}\big)_{ij}=p^{D-v(d_j)}u_j^{-1}d_j\delta_{ij}=p^D\delta_{ij}.
\]
Replace $y_j$ by $s\,y_j$, which lies in $L$ since $s\in\bZ_p$. Then $\Psi(x_i,y_j)=p^D\delta_{ij}$, so $\mathrm{gc}\ge\lvert I\rvert$ and $\mathrm{bad}\le n-\lvert I\rvert$.
\end{proof}

The scalar $s$ allows the window theorem to be applied to a multiple of the form. The Smith exponents of $A$ exceed those of the Gram matrix of $\Psi$ by $v(s)\ge0$, so the bound is in the right direction.

\begin{lemma}\label{w1:lem:trunc}
Let $A,A'\in M_n(\cO)$, where $\cO$ is the ring of integers of a finite extension of $\bQ_p$ with normalized valuation extended from $v$, and suppose that every entry of $A-A'$ has valuation at least $T>0$. Let $s_1\le\dots\le s_n$ and $s_1'\le\dots\le s_n'$ be the Smith exponents of $A$ and $A'$, with $s_i=+\infty$ for the zero diagonal entries. Then $\min(s_i,T)=\min(s'_i,T)$ for every $i$.
\end{lemma}

\begin{proof}
Let $\pi$ be a uniformizer of $\cO$ with $\pi^e\cO=p\cO$, and assume $eT\in\bZ$, as it is in the application where $T$ is an integer. If $A=U\operatorname{diag}(\pi^{es_1},\dots)V$ then $\operatorname{coker}A\cong\bigoplus_i\cO/\pi^{es_i}\cO$, and tensoring with $R_T=\cO/\pi^{eT}\cO$ gives $\bigoplus_iR_T/\pi^{e\min(s_i,T)}R_T$. On the other hand $\operatorname{coker}A\otimes R_T$ is the cokernel of the reduction of $A$ modulo $\pi^{eT}$, which equals the reduction of $A'$. So the two finitely generated $\cO$-modules $\bigoplus_i\cO/\pi^{e\min(s_i,T)}$ and $\bigoplus_i\cO/\pi^{e\min(s'_i,T)}$ are isomorphic, and the uniqueness of elementary divisors over the discrete valuation ring $\cO$ gives the claim.
\end{proof}

Smith exponents do not change under a finite extension of the base, since $s_1+\dots+s_l$ is the least valuation of an $l\times l$ minor. So a matrix over $\bZ_p$ may be treated over a larger ring of integers.

\subsubsection{The window theorem}

For $c\in\bC_p$ and an integer $i\ge0$ write $\binom ci=c(c-1)\cdots(c-i+1)/i!$.

\begin{lemma}\label{w1:lem:binshift}
Let $h\ge1$ and $\Delta<p^h$ be integers, let $E\ge h$ be an integer and let $y\in\bC_p$ with $v(y)=E$. Then for every $z\in\bZ_{\ge0}$ and $0\le k\le\Delta$,
\[
v\Big(\binom{z+y}k-\binom zk\Big)\ge E-h+1 .
\]
\end{lemma}

\begin{proof}
Vandermonde's identity $\binom{z+y}k=\sum_{l=0}^k\binom z{k-l}\binom yl$ holds as an identity of polynomials in $y$. The terms with $l\ge1$ remain. For $1\le l\le k<p^h$,
\[
\binom yl=\frac yl\prod_{i=1}^{l-1}\frac{y-i}{i},
\]
and $v(i)\le h-1<E=v(y)$ for $1\le i<l$, so $v(y-i)=v(i)$ and the product is a unit. Hence $v\binom yl=E-v(l)\ge E-h+1$. The integers $\binom z{k-l}$ have nonnegative valuation.
\end{proof}

\begin{lemma}\label{w1:lem:closest}
Let $t_0,d\in\bZ$ with $p\nmid d$, let $Y_0\ge1$, and put $z_j=t_0+dj$ for $0\le j<Y_0$. Let $\zeta\in\bC_p$, and let $j^*$ be an index for which $v(z_{j}-\zeta)$ is largest. Then $v(z_j-\zeta)\le v_p(j-j^*)$ for $j\ne j^*$, and
\[
\sum_{j\ne j^*}v_p(j-j^*)\le\frac{Y_0-1}{p-1}.
\]
\end{lemma}

\begin{proof}
If $v(z_j-\zeta)>v(z_j-z_{j^*})$ for some $j\ne j^*$, then $v(z_{j^*}-\zeta)=v(z_j-z_{j^*})<v(z_j-\zeta)$ by the ultrametric inequality, which contradicts the choice of $j^*$. Since $p\nmid d$, $v(z_j-z_{j^*})=v_p(j-j^*)$. The numbers $\lvert j-j^*\rvert$ for $j\ne j^*$ are $1,\dots,j^*$ and $1,\dots,Y_0-1-j^*$, so the sum is $v(j^*!)+v((Y_0-1-j^*)!)\le v((Y_0-1)!)\le(Y_0-1)/(p-1)$, using that $\binom{a+b}a$ is an integer.
\end{proof}

We now state the window theorem. The proof below gives the slightly stronger bound with $h-1$ in place of $h$.

\begin{theorem}\label{w1:thm:window}
Let $n,\Delta,h,Y_0\ge1$ be integers with $\Delta+1\le p^h$ and $Y_0\le p^h$, and let $E\ge1$ be an integer. Let $c_{ab,i}\in\bC_p$ for $1\le a,b\le n$ and $0\le i\le\Delta$, with $v(c_{ab,i})\ge0$, and let
\[
A(z)=\Big(\sum_{i=0}^{\Delta}c_{ab,i}\binom zi\Big)_{a,b}\in M_n(\bC_p[z]).
\]
Suppose that $\det A(z)=\ell\prod_{\zeta\in R}(z-\zeta)$ with $\ell\ne0$ and $R$ the multiset of roots in $\bC_p$, and let $\Lambda_0$ be a real number with
\begin{equation}\label{w1:eq:leadhyp}
v(\ell)+\frac{\lvert R\rvert}{p-1}\le n\Lambda_0 .
\end{equation}
Let $t_0,d\in\bZ$ with $p\nmid d$, $z_j=t_0+dj\ge0$ for $0\le j<Y_0$, and let $s_1(j)\le\dots\le s_n(j)$ be the Smith exponents of $A(z_j)$. Then
\begin{equation}\label{w1:eq:window}
\frac1{Y_0}\sum_{j=0}^{Y_0-1}\#\{i:\ s_i(j)\ge E\}\le n\Big(\frac{\Delta}{Y_0}+\frac{\Lambda_0+\Delta(h-1)/Y_0}{E}\Big).
\end{equation}
\end{theorem}

\begin{proof}
All the numbers $c_{ab,i}$, $\ell$ and the roots lie in a finite extension $K$ of $\bQ_p$, and we may enlarge $K$ to an unramified extension of any degree. The entries of $A(z_j)$ lie in the ring of integers, because $\binom{z_j}i\in\bZ$. The degree of $\det A$ is at most $n\Delta$, so $\lvert R\rvert\le n\Delta$.

Fix $j$ and put $z=z_j$ and $E'=E+h-1$. Choose $y\in K$ with $v(y)=E'$ such that $v(z+y-\zeta)=\min(v(z-\zeta),E')$ for every $\zeta\in R$. Such $y$ exists. Write $y=p^{E'}u$ with $u$ a unit. For a root with $v(z-\zeta)\ne E'$ the equality holds for every unit $u$. For a root with $v(z-\zeta)=E'$ it fails exactly when $u\equiv-(z-\zeta)p^{-E'}$ modulo the maximal ideal, which excludes one residue class. The residue field of $K$ is larger than $\lvert R\rvert+1$ after an unramified extension, so a unit $u$ avoiding these classes exists. Put $z'=z+y$.

By Lemma~\ref{w1:lem:binshift} with $\Delta<p^h$ and $E'\ge h$, every entry of $A(z')-A(z)$ has valuation at least $E'-h+1=E$. In particular $A(z')$ has integral entries. By Lemma~\ref{w1:lem:trunc} the truncated Smith exponents agree, $\min(s_i(j),E)=\min(s_i',E)$, where $s_i'\ge0$ are the Smith exponents of $A(z')$. Hence
\begin{align*}
E\cdot\#\{i:s_i(j)\ge E\}&\le\sum_i\min(s_i(j),E)=\sum_i\min(s'_i,E)\le\sum_is'_i\\
&=v(\det A(z'))=v(\ell)+\sum_{\zeta\in R}\min\big(v(z_j-\zeta),E'\big).
\end{align*}
Sum over $j$. For each $\zeta\in R$ let $j^*(\zeta)$ be as in Lemma~\ref{w1:lem:closest}. The term $j=j^*(\zeta)$ contributes at most $E'$, and the other terms contribute at most $\sum_{j\ne j^*}v_p(j-j^*)\le Y_0/(p-1)$. So
\[
E\sum_j\#\{i:s_i(j)\ge E\}\le Y_0\,v(\ell)+\lvert R\rvert\Big(E+h-1+\frac{Y_0}{p-1}\Big)\le Y_0\,n\Lambda_0+n\Delta(E+h-1),
\]
by \eqref{w1:eq:leadhyp} and $\lvert R\rvert\le n\Delta$. Divide by $EY_0$.
\end{proof}

\begin{remark}
The first term $n\Delta/Y_0$ accounts for points of the progression that lie very close to a root of the determinant, where the matrix can lose rank. The second term is Markov's inequality for the remaining valuation, which is controlled by the leading coefficient and by the spacing of a progression whose step is prime to $p$. The step $p\nmid d$ is used in Lemma~\ref{w1:lem:closest}, and the condition $\Delta<p^h$ is used in Lemma~\ref{w1:lem:binshift}.
\end{remark}

\subsubsection{The tie operators and the Gram polynomial}

Let $r\ge1$ and let $w_1,\dots,w_r,e_1,\dots,e_r$ be distinct variables in $\sigma$. Put
\[
Q=\sum_{c=1}^rx_{w_c}x_{e_c},\qquad \mathsf E(F)=QF,\qquad \mathsf F=\sum_{c=1}^r\partial_{w_c}\partial_{e_c}.
\]
A polynomial $F$ is \emph{bihomogeneous of bidegree $(o,n)$} when every monomial of $F$ has total degree $o$ in the variables $x_w$ and $n$ in the variables $x_e$. For integers $o,n,z\ge0$ put $A_z=o+n+r+z-1$ and define the \emph{tie operator}
\begin{equation}\label{w1:eq:tiephi}
\Phi_z=\sum_{j=0}^{o}\frac{(-1)^j}{(j!)^2\binom{A_z}j}\,\mathsf E^j\mathsf F^j
\end{equation}. The variable $z$ is the free parameter of the tie, which in the application is the shifted weight $a_+-o_+$ of the next top.

\begin{lemma}\label{w1:lem:adj}
For all $F,G$ and $l\ge0$, $\langle\mathsf E^lF,G\rangle=\langle F,\mathsf F^lG\rangle$. If $G$ is bihomogeneous of bidegree $(o,n)$ and $l>\min(o,n)$, then $\mathsf F^lG=0$.
\end{lemma}

\begin{proof}
By the adjointness of $x_s$ and $\partial_s$ in Section~\ref{w1:sec:scope}, $\langle x_{w_c}x_{e_c}F,G\rangle=\langle F,\partial_{w_c}\partial_{e_c}G\rangle$. Sum over $c$ and iterate. Each application of $\mathsf F$ lowers both partial degrees by one, so $\mathsf F^lG=0$ when $l$ exceeds one of them.
\end{proof}

\begin{proposition}\label{w1:prop:tie}
Let $v_0$ be bihomogeneous of bidegree $(o,n)$, let $u$ be any polynomial, and put $J=\min(o,n)$ and
\[
m_l(u,v_0)=\frac{\langle\mathsf E^l\mathsf F^lu,v_0\rangle}{(l!)^2\binom Jl}\qquad(0\le l\le J).
\]
Then for every integer $z\ge0$,
\begin{equation}\label{w1:eq:phipair}
\binom{A_z}J\langle\Phi_zu,v_0\rangle=\sum_{l=0}^J(-1)^l\binom{A_z-l}{J-l}m_l(u,v_0)=\sum_{i=0}^J\binom zi\,c_i(u,v_0),
\end{equation}
with
\[
c_i(u,v_0)=\sum_{l=0}^{J-i}(-1)^l\binom{A_0-l}{J-l-i}m_l(u,v_0),\qquad c_J(u,v_0)=\langle u,v_0\rangle .
\]
\end{proposition}

\begin{proof}
By Lemma~\ref{w1:lem:adj}, $\langle\mathsf E^j\mathsf F^ju,v_0\rangle=\langle\mathsf F^ju,\mathsf F^jv_0\rangle=0$ for $J<j\le o$. For $j\le J$ the identity $\binom{A}J\binom Jj=\binom Aj\binom{A-j}{J-j}$ with $A=A_z\ge J$ gives
\[
\binom{A_z}J\frac{(-1)^j}{(j!)^2\binom{A_z}j}\langle\mathsf E^j\mathsf F^ju,v_0\rangle=(-1)^j\binom{A_z-j}{J-j}m_j(u,v_0),
\]
which is the first equality. For the second, Vandermonde's identity gives $\binom{A_0-l+z}{J-l}=\sum_{i=0}^{J-l}\binom zi\binom{A_0-l}{J-l-i}$, valid since $A_0-l\ge0$. Exchange the sums. For $i=J$ the inner sum has the single term $l=0$, and $m_0(u,v_0)=\langle u,v_0\rangle$.
\end{proof}

So the paired value of the tie operator, normalized by $\binom{A_z}J$, is a polynomial in $z$ of degree at most $J$ in the binomial basis. Its top coefficient is the untwisted Fock form.

\subsubsection{The window bound at one level}

\begin{theorem}\label{w1:thm:level}
Let $o,n\ge0$ and $J=\min(o,n)$. Let $T$ be a linear operator on polynomials, $d_F\in\bQ_p^\times$, and $f_1,\dots,f_{n_1}$ polynomials that are $\bZ_p$-linearly independent and bihomogeneous of bidegree $(o,n)$. Let $x_0\ge0$ and $D_1$ be integers and $\Lambda$ a real number. Suppose
\begin{enumerate}[label=\textup{(\roman*)}]
\item $v\big(p^{x_0}m_l(Tf_a,f_b)/d_F\big)\ge0$ for all $a,b$ and $0\le l\le J$;
\item $\langle Tf_a,f_b\rangle=0$ for $a\ne b$, and $\langle Tf_a,f_a\rangle\ne0$ with $v\big(\langle Tf_a,f_a\rangle/d_F\big)\le D_1$;
\item $J/(p-1)-v(J!)\le\Lambda$.
\end{enumerate}
Let $W,h\ge1$ with $W\le p^h$ and $J+1\le p^h$, let $z_0\ge0$ and $k\ge0$ be integers, and let $L=\spn_{\bZ_p}(f_1,\dots,f_{n_1})$. For $0\le i<W$ let $\Psi_i(u,v_0)=\langle\Phi_{z_0+i}Tu,v_0\rangle/d_F$. Then
\[
\sum_{i=0}^{W-1}\mathrm{bad}(\Psi_i,L,k)\le W\,n_1\Big(\frac JW+\frac{x_0+D_1+\Lambda+J(h-1)/W}{k+1}\Big).
\]
\end{theorem}

\begin{proof}
Let $c_{ab,i}=p^{x_0}c_i(Tf_a,f_b)/d_F$ with $c_i$ as in Proposition~\ref{w1:prop:tie}. By (i) and the integrality of the binomial coefficients, $v(c_{ab,i})\ge0$. Let $A(z)=\big(\sum_ic_{ab,i}\binom zi\big)_{ab}$. By Proposition~\ref{w1:prop:tie},
\[
A(z_0+i)=s_i\big(\Psi_i(f_a,f_b)\big)_{ab},\qquad s_i=p^{x_0}\binom{A_{z_0+i}}J\in\bZ_p\setminus\{0\}.
\]
Lemma~\ref{w1:lem:smithbad} gives $\mathrm{bad}(\Psi_i,L,k)\le\#\{a:\ s_a(i)\ge k+1\}$, where $s_a(i)$ are the Smith exponents of $A(z_0+i)$.

The top coefficient matrix is $C_J=\big(p^{x_0}\langle Tf_a,f_b\rangle/d_F\big)$, diagonal by (ii). Since $\binom zJ=z^J/J!+(\text{lower degree})$, the coefficient of $z^{n_1J}$ in $\det A(z)$ is $\ell=\det C_J/(J!)^{n_1}\ne0$. So $\deg\det A=n_1J=\lvert R\rvert$ and
\[
v(\ell)+\frac{\lvert R\rvert}{p-1}\le n_1(x_0+D_1)-n_1v(J!)+\frac{n_1J}{p-1}\le n_1(x_0+D_1+\Lambda).
\]
Apply Theorem~\ref{w1:thm:window} with $n=n_1$, $\Delta=J$, $\Lambda_0=x_0+D_1+\Lambda$, $E=k+1$, $t_0=z_0$, $d=1$, $Y_0=W$.
\end{proof}



\begin{lemma}\label{w1:lem:jordan}
Let $p$ be odd, let $L$ be a free $\bZ_p$-module of finite rank and $B\colon L\times L\to\bQ_p$ a symmetric bilinear form. Then $L$ has a basis $f_1,\dots,f_n$ with $B(f_a,f_b)=0$ for $a\ne b$.
\end{lemma}

\begin{proof}
If $B=0$ every basis works. Otherwise let $m$ be the least valuation of a value $B(u,u')$. If $v(B(u,u))>m$ for every $u$, then $2B(u,u')=B(u+u',u+u')-B(u,u)-B(u',u')$ has valuation larger than $m$ for all $u,u'$, and since $2$ is a unit this contradicts the choice of $m$. Choose $f_1$ with $v(B(f_1,f_1))=m$. Then $B(u,f_1)/B(f_1,f_1)\in\bZ_p$ for every $u$, and $u\mapsto u-\frac{B(u,f_1)}{B(f_1,f_1)}f_1$ projects $L$ onto $L_1=\{u:B(u,f_1)=0\}$ with kernel $\bZ_pf_1$. So $L=\bZ_pf_1\oplus L_1$ with $L_1$ free, and induction on the rank completes the proof.
\end{proof}

In the application the vectors of a basis with $B(f_a,f_a)=0$ or with $v(B(f_a,f_a))>D_1$ are counted as bad, and the remaining vectors form the family of Theorem~\ref{w1:thm:level}.

\subsubsection{From one level to the segment share}\label{w1:sec:Crest}

The passage from Theorem~\ref{w1:thm:level} to Hypothesis~\ref{r2:engines}(C) goes through the following steps.

\begin{enumerate}[label=(C\arabic*),start=6,leftmargin=*]
\item \emph{Integrality of the level coefficients}. For a block datum $G$ with all letter degrees and tie parameters at most $2A$, every $u,v_0$ in the level lattice $L_{j,j_p}$ and $l\le J$ satisfy hypothesis (i) of Theorem~\ref{w1:thm:level} with $x_0=\big((t+1)(r+1)+(N+1)r+2(t+1)\big)\lambda_N$, where $\lambda_N=\lfloor\log_p(8A+r)\rfloor+1$ and $T$ is the product of the tie operators of the other ties. The mechanism is that $\mathsf F^l/l!$ acts integrally on integral polynomials, since $\frac1{l!}\mathsf F^l=\sum_{\lvert\kappa\rvert=l}\kappa!\prod_c\partial_{w_c}^{(\kappa_c)}\partial_{e_c}^{(\kappa_c)}$ with divided derivatives $\partial^{(k)}=\partial^k/k!$, and that each tie operator \eqref{w1:eq:tiephi} has denominators $\binom{A_z}j$ of valuation at most $\log_pA_z$ by Kummer's theorem. The bookkeeping of the factorials in $d_F$ against the Fock weights $m!$ on $L_{j,j_p}$ is not written here.
\item \emph{Window bound}. Under the window and prime inequalities of Proposition~\ref{w1:prop:C7}, the bracket of Theorem~\ref{w1:thm:level} is at most $3/Y$.
\item \emph{Level decomposition of a block lattice}. The block lattice of a kept block decomposes, up to a loss of $\mathrm{levelPassageLoss}=10(s+1)(r+1)(\log_p(8a+r)+1)+2N_t(r+1)\log_pa$ in the depth, into the level lattices $L_{j,j_p}$ indexed by the $\mathfrak{sl}_2$ levels of the ties, and the block form restricted to a level equals $\langle\Phi_{z}T\cdot,\cdot\rangle/d_F$ with $z=a_+-o_+$. The identity behind it is the following: for the tie with extra variables $x_{z_v},x_{s_v}$ and $P_z(y)=\sum_l\frac{(-1)^l}{(l!)^2\binom{A_z}l}\hat{\mathsf E}^l\mathsf F^l(x_{z_v}^zy)$, one has $\hat{\mathsf F}P_z(y)=0$ and $\langle P_z(y),P_z(y')\rangle=z!\,\langle y,\Phi_zy'\rangle$.
\item \emph{Averaging over the other ties}. The bad counts $\mathrm{bY}$ of the untwisted level forms, averaged over the top weights of the segment, are at most $(30+C_M)/Y$ times the averaged dimensions. The mechanism is Theorem~\ref{w1:thm:window} applied in each tie coordinate, with the arithmetic progression of tie parameters cut out by the digit classes, and $\varepsilon_i=2S_i/Y_i+(2\kappa+2S_i\lambda/Y_i)/E$ as the error of the $i$-th coordinate.
\item \emph{Segment assembly}. Summation over blocks, levels and ties, which gives the factor $(t+1)N_1$, and the dichotomy that fixes the split index $\sigma$ uniformly over the window when no power $p^d$ lies in $[(1-2\delta)a/b',(1+2\delta)a/b']$.
\item \emph{From segment shares to the share of} $\mathrm{LP}$. On every $1$-edge there is a source-test system of rank at least $(1-\kappa_P-\sum_j\kappa_T(j))\dim V/[k]_q$, with $\kappa_P\le C_sNN_1/Y$.
\end{enumerate}

\notgiven{Statements (C6) and (C8) to (C11) are not proved here.}

\subsubsection{The numerics of the window bound (C7)}\label{w1:sec:C7}

In the setting of (C6) and (C7) let $s$ be the number of tops of the segment, $t=s-\sigma\le s$, $r\ge4$ the number of frame variables of a block, $Y\ge1$, $C_M\ge1$, $\mathrm{winA}$ the window of top weights, an interval of integers with $a/(kY)\le\lvert\mathrm{winA}\rvert\le a+1$, and $b_+\le\frac52b'$ the bottom weight of the next top. Put
\begin{gather*}
\lambda_A=\log_p(8a+r)+1,\qquad\Lambda=\lfloor\log_p(8(3s+2)a+r)\rfloor+1,\qquad\kappa_S=3(s+2)r^2\lambda_A,\\
x_0=\big((t+1)(r+1)+(2s+2)r+2(t+1)\big)\Lambda,\qquad D_{1r}=t\,(Y\kappa_S+C_M\kappa_S)+C_M\kappa_S .
\end{gather*}

\begin{proposition}\label{w1:prop:C7}
Suppose $p\ge3s+2$, $p\ge25$, $k\ge5$, $k^2Y^2\le p$ and $pb'<a$. Let $J\le b_+$ be an integer and put $W=\lvert\mathrm{winA}\rvert$, $h=\lceil\log_p(a+1)\rceil$ and $K=\lfloor YD_{1r}\rfloor$. Then $W\le p^h$, $J+1\le p^h$, and
\[
\frac JW+\frac{x_0+\lfloor D_{1r}\rfloor+\Lambda+Jh/W}{K+1}\le\frac3Y\le\frac5Y .
\]
\end{proposition}

\begin{proof}
First $J\le\frac52b'<a$, so $J+1\le a+1\le p^h$ and $W\le a+1\le p^h$. From $k^2Y^2b'\le pb'<a$ we get $kYb'<a/(kY)\le W$, so
\[
\frac JW\le\frac{5b'}{2kYb'}=\frac5{2kY}\le\frac1{2Y},
\]
using $k\ge5$. Next, $8(3s+2)a+r\le(3s+2)(8a+r)$ and $3s+2\le p$ give $\Lambda\le\log_p(8(3s+2)a+r)+1\le\lambda_A+1\le2\lambda_A$, since $\lambda_A\ge1$. Put $c=(t+1)(r+1)+(2s+2)r+2(t+1)+1$. From $t\le s$ and $r\ge4$,
\[
2c\le6(s+1)(r+1)+2\le12(s+2)r\le3(s+2)r^2,
\]
so $x_0+\Lambda=c\Lambda\le2c\lambda_A\le\kappa_S$. Since $C_M\ge1$ and $t,Y\ge0$, $\kappa_S\le D_{1r}$, and $h\le\log_p(a+1)+1\le\lambda_A\le\kappa_S$. The denominator satisfies $K+1\ge YD_{1r}>0$. Hence
\[
\frac{x_0+\lfloor D_{1r}\rfloor+\Lambda+Jh/W}{K+1}\le\frac{2D_{1r}+\frac1{2Y}D_{1r}}{YD_{1r}}=\frac2Y+\frac1{2Y^2}\le\frac2Y+\frac1{2Y},
\]
and adding $J/W\le1/(2Y)$ gives $3/Y$.
\end{proof}

For the medium primes $p\ge p_*=k^2x^{26}$ take $Y=x^{13}$. Then $k^2Y^2\le p$, and $p<X'$ gives $pb'<X'b'=a$.



Items (C1) to (C5) are Lemmas~\ref{w1:lem:smithbad}, \ref{w1:lem:trunc}, Theorem~\ref{w1:thm:window}, Proposition~\ref{w1:prop:tie} with Theorem~\ref{w1:thm:level}, and Lemma~\ref{w1:lem:jordan}.

\subsection{Engine (B): detection for the twisted Fock form}\label{w1:sec:B}

The detection statement of Hypothesis~\ref{r2:engines}(B) combines three steps, which we call Steps~B, C and DEF, with five tools. This section proves the arithmetic tools of the twist, namely the factorial digit split, the decay of the unit parts, the binomial approximation and the divided-derivative congruence, and the integrality of divided powers on the balanced lattice that controls the Fock projectors.

\subsubsection{The factorial digit split}

Let $P=p^g$ with $g\ge1$. For integers $n\ge0$ and $0\le\rho<P$ put
\begin{equation}\label{w1:eq:cd}
c_P(n)=\prod_{i=0}^{n-1}\prod_{y=1}^{P-1}\Big(1+\frac{Pi}y\Big),\qquad d_P(n,\rho)=\prod_{y=1}^{\rho}\Big(1+\frac{Pn}y\Big).
\end{equation}

\begin{proposition}\label{w1:prop:digitsplit}
For $0\le\rho<P$ and $\beta\ge0$,
\[
(\rho+P\beta)!=\rho!\,(P!)^\beta\,\beta!\,c_P(\beta)\,d_P(\beta,\rho),
\]
and $c_P(\beta)$ and $d_P(\beta,\rho)$ are $1$-units, that is, $v(c_P(\beta)-1)\ge1$ and $v(d_P(\beta,\rho)-1)\ge1$.
\end{proposition}

\begin{proof}
First $(\rho+P\beta)!/(P\beta)!=\prod_{y=1}^\rho(P\beta+y)=\rho!\,d_P(\beta,\rho)$. Next
\[
(P\beta)!=\prod_{i=0}^{\beta-1}\Big(P(i+1)\prod_{y=1}^{P-1}(Pi+y)\Big)=P^\beta\beta!\prod_{i=0}^{\beta-1}\prod_{y=1}^{P-1}y\Big(1+\frac{Pi}y\Big)=(P!)^\beta\beta!\,c_P(\beta),
\]
using $P\,(P-1)!=P!$. For $1\le y<P$ one has $v(y)\le g-1$, so $v(Pi/y)\ge1$ and each factor of \eqref{w1:eq:cd} is a $1$-unit.
\end{proof}

The identity says that adding $P\beta$ to a low digit $\rho<P$ produces no carry, so the factorial splits into the factorial of the low digit, the factorial of the high digit, a power of $P!$ and a $1$-unit.

\subsubsection{Decay of the unit parts}

For a function $f\colon\bZ_{\ge0}\to\bC_p$ let $\Delta f(n)=f(n+1)-f(n)$. The Mahler coefficients of $f$ are $a_j(f)=\Delta^jf(0)$, and $f(n)=\sum_ja_j(f)\binom nj$.

\begin{definition}
For real $\alpha,\gamma\ge0$ let $\mathcal D(\alpha,\gamma)$ be the set of functions $f\colon\bZ_{\ge0}\to\bC_p$ with $v(\Delta^jf(n))\ge\gamma+\alpha j$ for all $j,n\ge0$.
\end{definition}

\begin{lemma}\label{w1:lem:Dring}
\begin{enumerate}[label=\textup{(\roman*)}]
\item $f\in\mathcal D(\alpha,\gamma)$ if and only if $v(a_j(f))\ge\gamma+\alpha j$ for every $j$.
\item $\mathcal D(\alpha,\gamma)$ is a $\bZ_p$-module, and $\mathcal D(\alpha,\gamma)\cdot\mathcal D(\alpha,\gamma')\subset\mathcal D(\alpha,\gamma+\gamma')$.
\item If $f_m\in\mathcal D(\alpha,\gamma_m)$ with $\gamma_m\ge\gamma$ and $\gamma_m\to\infty$, then $\sum_mf_m$ converges pointwise and lies in $\mathcal D(\alpha,\gamma)$.
\end{enumerate}
\end{lemma}

\begin{proof}
(i) One direction is $n=0$. Conversely $\Delta^jf(n)=\sum_{i\ge0}\binom ni a_{i+j}(f)$, and $\binom ni\in\bZ$, $\gamma+\alpha(i+j)\ge\gamma+\alpha j$. (ii) Linearity is clear. For products, the Leibniz rule for differences, $\Delta(fg)(n)=\Delta f(n)\,g(n+1)+f(n)\,\Delta g(n)$, iterates to
\[
\Delta^j(fg)(n)=\sum_{k=0}^j\binom jk\Delta^kf(n)\,\Delta^{j-k}g(n+k),
\]
and each term has valuation at least $\gamma+\alpha k+\gamma'+\alpha(j-k)$. (iii) Each $\Delta^jf_m(n)$ is a finite integral combination of values of $f_m$, which have valuation at least $\gamma_m$, so the series of differences converge, and the bound passes to the limit.
\end{proof}

\begin{proposition}\label{w1:prop:decay}
Let $p\ge5$ and $P=p^g$, $g\ge1$. The functions $n\mapsto c_P(n)^{\pm1}$ and, for each $0\le\rho<P$, $n\mapsto d_P(n,\rho)^{\pm1}$ lie in $\mathcal D(\tfrac14,0)$. In particular their Mahler coefficients satisfy $v(a_j)\ge j/4$.
\end{proposition}

\begin{proof}
Let $\log$ be the $p$-adic logarithm on $1$-units, $\log(1+u)=\sum_{s\ge0}(-1)^su^{s+1}/(s+1)$ for $v(u)\ge1$. Since $v(u)\ge1>1/(p-1)$, $\exp$ converges at $\log(1+u)$ and $\exp\log(1+u)=1+u$, and $\log$ is additive on $1$-units.

\emph{The function $\ell=\log c_P$.} By additivity and absolute convergence,
\[
\ell(n)=\sum_{s\ge0}\frac{(-1)^sP^{s+1}H_P(s+1)}{s+1}S_{s+1}(n),\qquad H_P(m)=\sum_{y=1}^{P-1}y^{-m},\quad S_m(n)=\sum_{i=0}^{n-1}i^m .
\]
Since $v(y)\le g-1$ for $1\le y<P$, $v(P^mH_P(m))\ge mg-m(g-1)=m$. The function $S_m$ satisfies $\Delta S_m(n)=n^m$, so $a_0(S_m)=0$, $a_1(S_m)=0^m=0$ for $m\ge1$, and $a_j(S_m)=\Delta^{j-1}(n^m)(0)=(j-1)!\,\mathcal S(m,j-1)$ for $j\ge2$, where $\mathcal S$ denotes the Stirling numbers of the second kind, which vanish for $j-1>m$. Hence $a_0(\ell)=a_1(\ell)=0$, and for $j\ge2$
\[
v(a_j(\ell))\ge\min_{m\ge j-1}\big(m-v_p(m)\big)\ge(j-1)-\log_p(j-1),
\]
because $x-\log_px$ is increasing for $x\ge1$ and $v_p(m)\le\log_pm$. For $j=2$ this is $1$. For $j\ge3$ and $p\ge5$ we show $(j-1)-\log_p(j-1)\ge\frac12+\frac j4$: at $j=3$ the left side is $2-\log_52>1.56\ge1.25$, and the difference of the two sides increases with $j$. For $j=1$ and $j=0$ the coefficients vanish. So $v(a_j(\ell))\ge\frac12+\frac j4$ for all $j$, that is, $\ell\in\mathcal D(\frac14,\frac12)$ by Lemma~\ref{w1:lem:Dring}(i).

\emph{Exponentiation.} By Lemma~\ref{w1:lem:Dring}(ii), $\pm\ell^m/m!\in\mathcal D(\frac14,\frac m2-v(m!))$, and $\frac m2-v(m!)\ge\frac m2-\frac m{p-1}\ge\frac m4$ for $p\ge5$. By Lemma~\ref{w1:lem:Dring}(iii), $c_P^{\pm1}=\exp(\pm\ell)=\sum_m(\pm\ell)^m/m!$ lies in $\mathcal D(\frac14,0)$.

\emph{The function $\log d_P(\cdot,\rho)$.} Here $\log d_P(n,\rho)=\sum_{y=1}^\rho\sum_s\frac{(-1)^s(P/y)^{s+1}}{s+1}n^{s+1}$, and $a_j(n^m)=j!\,\mathcal S(m,j)$, which vanishes for $j=0$ and $j>m$. So $v(a_j(\log d_P))\ge\min_{m\ge j}(m-v_p(m))\ge j-\log_pj$ for $j\ge1$, which is at least $\frac12+\frac j4$ for $p\ge5$. The same exponentiation applies.
\end{proof}



\begin{corollary}\label{w1:cor:binapprox}
Let $p\ge5$, let $f_1,\dots,f_K\in\mathcal D(\frac14,0)$, and let $c\ge1$ be an integer. There are $b_{\mathbf J}\in\bC_p$ with $v(b_{\mathbf J})\ge0$, indexed by $\mathbf J\in\bZ_{\ge0}^K$ with $\lvert\mathbf J\rvert\le4c$, such that for every $\mathbf n\in\bZ_{\ge0}^K$
\[
v\Big(\prod_{k=1}^Kf_k(n_k)-\sum_{\lvert\mathbf J\rvert\le4c}b_{\mathbf J}\prod_k\binom{n_k}{J_k}\Big)\ge c .
\]
If the $f_k$ take rational values, the $b_{\mathbf J}$ are rational.
\end{corollary}

\begin{proof}
Write $f_k(n)=\sum_ja_{k,j}\binom nj$ with $v(a_{k,j})\ge j/4$. Then $\prod_kf_k(n_k)=\sum_{\mathbf J}\big(\prod_ka_{k,J_k}\big)\prod_k\binom{n_k}{J_k}$, a pointwise convergent sum in which the term of $\mathbf J$ has valuation at least $\lvert\mathbf J\rvert/4$. Put $b_{\mathbf J}=\prod_ka_{k,J_k}$ for $\lvert\mathbf J\rvert\le4c$. The omitted terms have valuation at least $(4c+1)/4>c$.
\end{proof}



\subsubsection{Binomial polynomials as divided derivatives}

For a variable $x$ the divided derivative is $\partial^{(m)}x^\alpha=\binom\alpha mx^{\alpha-m}$, and for multi-indices $\partial^{(\mathbf m)}=\prod_s\partial_s^{(m_s)}$. It maps integral polynomials to integral polynomials.

\begin{lemma}\label{w1:lem:binop}
For a multi-index $\mathbf m$ and every monomial $x^\alpha$, $\binom\alpha{\mathbf m}x^\alpha=x^{\mathbf m}\partial^{(\mathbf m)}x^\alpha$. Hence multiplication of the coefficient of $x^\alpha$ by a polynomial $\sum_{\mathbf m}b_{\mathbf m}\binom\alpha{\mathbf m}$ in the exponents is the differential operator $\sum_{\mathbf m}b_{\mathbf m}x^{\mathbf m}\partial^{(\mathbf m)}$, whose order is the largest $\lvert\mathbf m\rvert$.
\end{lemma}

\begin{proof}
$x^{\mathbf m}\partial^{(\mathbf m)}x^\alpha=\binom\alpha{\mathbf m}x^\alpha$ coordinatewise.
\end{proof}

Corollary~\ref{w1:cor:binapprox} and Lemma~\ref{w1:lem:binop} together say that multiplication by the twisting unit $\gamma(\alpha)$ of the twisted Fock form agrees modulo $p^c$ with an integral differential operator of order at most $4c$ in the digit variables.

\begin{proposition}\label{w1:prop:der}
Let $P=p^g$, let $\boldsymbol\rho,\mathbf M\in\bZ_{\ge0}^r$ with $\rho_a<P$, and let $c\ge1$. There are integers $e_{\boldsymbol\kappa}$, indexed by $\boldsymbol\kappa\in\bZ_{\ge0}^r$ with
\[
\lvert\boldsymbol\kappa\rvert<\frac{\lvert\mathbf M\rvert}P+c,
\]
such that for every $\boldsymbol\beta\in\bZ_{\ge0}^r$
\[
\binom{\boldsymbol\rho+P\boldsymbol\beta}{\mathbf M}\equiv\sum_{\boldsymbol\kappa}e_{\boldsymbol\kappa}\binom{\boldsymbol\beta}{\boldsymbol\kappa}\pmod{p^c}.
\]
Consequently $\partial^{(\mathbf M)}x^{\boldsymbol\rho+P\boldsymbol\beta}\equiv\big(\sum_{\boldsymbol\kappa}e_{\boldsymbol\kappa}\binom{\boldsymbol\beta}{\boldsymbol\kappa}\big)x^{\boldsymbol\rho+P\boldsymbol\beta-\mathbf M}$ modulo $p^c$, with the convention that the monomial is $0$ when an exponent is negative.
\end{proposition}

\begin{proof}
Coordinatewise, Vandermonde gives $\binom{\rho+P\beta}M=\sum_{i}\binom\rho i\binom{P\beta}{M-i}$, so it suffices to treat $\binom{P\beta}m$ for each $m\le M$, jointly over the coordinates. In $\bZ[x]$ write $(1+x)^P=1+x^P+R(x)$ with $R(x)=\sum_{0<l<P}\binom Plx^l$. By Kummer's theorem $v\binom Pl=g-v(l)\ge1$ for $0<l<P$, so $R\in p\bZ[x]$. Then
\[
(1+x)^{P\beta}=\big(1+(x^P+R(x))\big)^\beta=\sum_{k\ge0}\binom\beta k\sum_{u=0}^k\binom ku R(x)^ux^{P(k-u)} .
\]
The terms with $u\ge c$ lie in $p^c\bZ[x]$. A term with $u<c$ contributes to the coefficient of $x^m$ when $P(k-u)\le m$, and not otherwise, so $k<m/P+c$. Hence $\binom{P\beta}m\equiv\sum_{k<m/P+c}\epsilon_{m,k}\binom\beta k$ modulo $p^c$ with integers $\epsilon_{m,k}$. For several coordinates take the product of the generating functions $\prod_a(1+x_a)^{P\beta_a}$ and count the factors $R(x_a)$ jointly. A term with $u$ such factors in total lies in $p^u\bZ[x]$, so the terms with total $u<c$ remain, and the coefficient of $x^{\mathbf m}$ then involves $\binom{\boldsymbol\beta}{\boldsymbol\kappa}$ with $P(\lvert\boldsymbol\kappa\rvert-u)\le\lvert\mathbf m\rvert\le\lvert\mathbf M\rvert$, so $\lvert\boldsymbol\kappa\rvert<\lvert\mathbf M\rvert/P+c$.
\end{proof}

The divided derivative of a digit-split monomial agrees modulo $p^c$ with a combination of binomials in the high digits of order less than $\lvert\mathbf M\rvert/P+c$.

\subsubsection{Divided powers and the balanced lattice}

Let $\mathcal V$ be the space of polynomials in variables $x_{i,a}$, with $i$ in a set of letters and $1\le a\le r$, over a finite extension $\mathcal K$ of $\bQ_p$ with valuation ring $\cO$. The \emph{balanced lattice} is
\[
\mathcal B=\Big\{F=\sum_\alpha F_\alpha x^\alpha:\ v(F_\alpha)+\tfrac12v(\alpha!)\ge0\ \text{for all }\alpha\Big\},
\]
the lattice of the balanced norm $v_{\mathrm{BB}}(F)=\min_\alpha\big(v(F_\alpha)+\frac12v(\alpha!)\big)$ as in Hypothesis~\ref{r2:engines}(B). For a set $S$ of letters, each marked as a vector or a covector letter, and $a\ne b$, let $E_{ab}=\sum_{i\in S}\pm x_{i,a}\partial_{i,b}$, with the sign and the exchange of $a,b$ at covector letters that make $a\mapsto E_{ab}$ the action of $\mathfrak{gl}_r$ on the vectors and on the dual.

\begin{lemma}\label{w1:lem:divpow}
For $a\ne b$ and $k\ge0$ the divided power $E_{ab}^{(k)}=E_{ab}^k/k!$ maps $\mathcal B$ into $\mathcal B$.
\end{lemma}

\begin{proof}
The operators $x_{i,a}\partial_{i,b}$ for different $i$ commute, and $\big(\sum_iT_i\big)^{(k)}=\sum_{k_1+\dots=k}\prod_iT_i^{(k_i)}$ for commuting $T_i$. So it suffices to treat one letter and one pair $(a,b)$, and the action on the two variables $x_a=x_{i,a}$, $x_b=x_{i,b}$. On a monomial $x_a^nx_b^m$,
\[
(x_a\partial_b)^{(k)}x_a^nx_b^m=\binom mk x_a^{n+k}x_b^{m-k}.
\]
The change in the balanced weight is
\[
v\binom mk+\tfrac12v\big((n+k)!\,(m-k)!\big)-\tfrac12v(n!\,m!)=\tfrac12v\Big(\frac{m!}{(m-k)!}\Big)+\tfrac12v\Big(\frac{(n+k)!}{n!}\Big)-v(k!)\ge0,
\]
because $m!/(m-k)!$ and $(n+k)!/n!$ are both divisible by $k!$.
\end{proof}

Lemma~\ref{w1:lem:divpow} says that $\mathcal B$ is stable under the Kostant $\bZ$-form of $U(\mathfrak{gl}_r)$ acting through the letters of $S$.

\begin{proposition}\label{w1:prop:extremal}
Let $V\subset\mathcal V$ be a finite-dimensional subspace stable under the action of $\mathfrak{gl}_r$ through the letters of $S$, and let $M$ bound $\lvert\langle\nu,\alpha^\vee\rangle\rvert$ for every weight $\nu$ of $V$ and every root $\alpha$. Let $\mu$ be a dominant weight and $V_\mu$ its weight space. Order the positive roots normally and let
\[
\mathsf P=\prod_{\alpha>0}\ \sum_{k\ge0}\frac{(-1)^k}{k!}\,\frac{1}{(h_\alpha+\rho(h_\alpha)+1)\cdots(h_\alpha+\rho(h_\alpha)+k)}\,f_\alpha^ke_\alpha^k
\]
be the extremal projector of Asherova, Smirnov and Tolstoy. Then $\mathsf P$ is defined on $V_\mu$, it is the projection of $V_\mu$ onto the highest weight vectors $V_\mu^{\mathfrak n^+}$ along $V_\mu\cap\mathfrak n^-V$, and it maps $\mathcal B\cap V_\mu$ into $p^{-\binom r2\lfloor\log_p(M+r)\rfloor}(\mathcal B\cap V_\mu)$.
\end{proposition}

\begin{proof}
The formula and the projection property are the theorem of R.~M.~Asherova, Yu.~F.~Smirnov and V.~N.~Tolstoy, \emph{Projection operators for simple Lie groups}, Theor. Math. Phys.~8 (1971), 813--825, in the form of D.~P.~Zhelobenko, \emph{Representations of Reductive Lie Algebras} (1994), Chapter~IV. Each operator $f_\alpha^ke_\alpha^k$ preserves weights, so $h_\alpha$ acts on $V_\mu$ through the scalar $\langle\mu,\alpha^\vee\rangle$ in every factor. Put $N_\alpha=\langle\mu+\rho,\alpha^\vee\rangle$, a positive integer since $\mu$ is dominant, so no denominator vanishes. Since $f_\alpha^ke_\alpha^k=(k!)^2f_\alpha^{(k)}e_\alpha^{(k)}$, the $k$-th term of the factor of $\alpha$ acts on $V_\mu$ as
\[
(-1)^k\frac{k!}{(N_\alpha+1)\cdots(N_\alpha+k)}f_\alpha^{(k)}e_\alpha^{(k)}=(-1)^k\binom{N_\alpha+k}{k}^{-1}f_\alpha^{(k)}e_\alpha^{(k)} .
\]
If $e_\alpha^{(k)}V_\mu\ne0$ then $\mu+k\alpha$ is a weight of $V$, so $\langle\mu,\alpha^\vee\rangle+2k\le M$ and $N_\alpha+k\le M+\langle\rho,\alpha^\vee\rangle\le M+r$. Kummer's theorem gives $v\binom{N_\alpha+k}k\le\lfloor\log_p(M+r)\rfloor$. The divided powers preserve $\mathcal B$ by Lemma~\ref{w1:lem:divpow}, and a finite sum of elements of $p^{-c}\mathcal B$ lies in $p^{-c}\mathcal B$. So each of the $\binom r2$ factors costs at most $\lfloor\log_p(M+r)\rfloor$.
\end{proof}

Both are at least the loss of Proposition~\ref{w1:prop:extremal} when $M\le2\deg_S$. The restriction to dominant $\mu$ is necessary. For $\mathfrak{sl}_2$ and the lowest weight $-3$ of the representation of highest weight $3$, the denominator $(h+2)(h+3)$ vanishes.

\begin{proposition}\label{w1:prop:projexists}
Let $S$ be a set of letters and $\theta\in\bZ$. Let $A_S(\theta)\subset\mathcal V_d$ be the sum of the $\mathfrak{gl}_r$-isotypic components of the multidegree-$d$ space $\mathcal V_d$, for the action through $S$, whose highest weight $\mu$ has $\mu_r\ge\theta$. Then there is a unique Fock-selfadjoint projection $\pi_{S,\theta}$ of $\mathcal V_d$ onto $A_S(\theta)$. For $S\subset S'$ or $S\cap S'=\emptyset$, $\pi_{S,\theta}$ commutes with the action of $\mathfrak{gl}_r$ through $S'$ and with $\pi_{S',\theta'}$.
\end{proposition}

\begin{proof}
Everything is defined over $\bQ$, so we work on the space $\mathcal V_{d,\bQ}$ of rational polynomials and extend scalars at the end. There the Fock form is positive definite. By the adjointness of $x_s$ and $\partial_s$, the operator $x_{i,a}\partial_{i,b}$ is adjoint to $x_{i,b}\partial_{i,a}$, so $\langle Xu,u'\rangle=\langle u,X^{\mathsf T}u'\rangle$ for $X\in\mathfrak{gl}_r$ acting through $S$, where $X^{\mathsf T}$ is the transpose. The covector convention preserves this.

Let $V_\lambda$ and $V_\nu$ be two isotypic components with $\lambda\ne\nu$. The form gives a map $V_\lambda\to V_\nu^*$ that intertwines the action on $V_\lambda$ with the action $X\mapsto(X^{\mathsf T})^*$ on $V_\nu^*$. This twisted dual has the same weights as $V_\nu$, since $h^{\mathsf T}=h$ for diagonal $h$, so it is isomorphic to $V_\nu$, and Schur's lemma makes the map zero. So the isotypic components are mutually orthogonal, the form is nondegenerate on each of them by positivity, and the orthogonal projection onto $A_S(\theta)$ is the sum of the isotypic projections. It is selfadjoint. A selfadjoint projection onto a subspace on which the form is nondegenerate is the orthogonal projection, which gives uniqueness.

An operator commuting with the action through $S$ maps each isotypic component into itself, by Schur's lemma, and so commutes with $\pi_{S,\theta}$. If $S\subset S'$, the action through $S'$ is the sum of the action through $S$, which preserves the $S$-isotypic components, and the action through $S'\setminus S$, which commutes with the action through $S$. So the action through $S'$ commutes with $\pi_{S,\theta}$. Applying the same remark to $S'$, the projection $\pi_{S',\theta'}$ commutes with every operator commuting with the action through $S'$, in particular with $\pi_{S,\theta}$. If $S\cap S'=\emptyset$ the two actions commute and the same argument applies.
\end{proof}



The remaining steps of engine (B) are the following.

\begin{enumerate}[label=(B\arabic*),start=7,leftmargin=*]
\item \emph{Norm bound for the isotypic projector.} $v_{\mathrm{BB}}(\pi_{S,\theta}F)\ge v_{\mathrm{BB}}(F)-\binom r2\log_p(2\deg_S+r)$ for every $F\in\mathcal V_d$. Proposition~\ref{w1:prop:extremal} bounds the extremal projector $\mathsf P$, which projects onto highest weight vectors. The passage from $\mathsf P$ to the isotypic projector $\pi_{S,\theta}$, which reassembles a full isotypic component from its highest weight vectors by divided powers $f^{(k)}$ with normalizing binomial coefficients, is not written here. The expected loss is the one of $\mathsf P$ by the same Kummer bound, and the precise accounting is the residual.
\item \emph{Steps B, C and DEF.} The derivation of the detection bound of Hypothesis~\ref{r2:engines}(B) for a source $F$ of multidegree $(d_s)$ from the tools of this subsection. Its tools are Propositions~\ref{w1:prop:digitsplit}, \ref{w1:prop:decay}, \ref{w1:prop:der}, \ref{w1:prop:projexists}, Corollary~\ref{w1:cor:binapprox}, Lemma~\ref{w1:lem:binop} and (B7). The combination of the three steps into the detection bound $v(\mathrm B(F,\mathrm{ph}\,G'))\le n_0+D_m+D_l+3\kappa_p$ is a direct addition of the three losses once the three steps hold.
\item \emph{Transfer}. The contravariant self-duality with loss $\kappa_1(\lambda)+\kappa_1(\lambda^\flat)\le C_K\lvert T\rvert k^2\log_pa$, and surjectivity of multiplication $H^0(\lambda^\flat)\otimes H^0(\lambda^{\mathrm{rest}})\to H^0(\lambda)$ on edge lattices. The surjectivity over $\bZ$ for $\mathrm{GL}_n$ is the projective normality of flag varieties in all characteristics, S.~Ramanan and A.~Ramanathan, \emph{Projective normality of flag varieties and Schubert varieties}, Invent. Math.~79 (1985), 217--224. Its transport to the edge lattices is not written here.
\end{enumerate}

\notgiven{Statements (B8) and (B9) are not proved here, apart from the published surjectivity theorem named in (B9).}

\subsubsection{The norm of the isotypic projector}\label{w1:sec:B7}

Throughout this subsection $\mathcal K$ is a finite extension of $\bQ_p$ with valuation ring $\cO$, $p$ is odd, and $\mathfrak g=\mathfrak{gl}_r(\mathcal K)$ with the standard basis $E_{ab}$. For a positive root $\alpha=\epsilon_a-\epsilon_b$ ($a<b$) write $e_\alpha=E_{ab}$, $f_\alpha=E_{ba}$, $e_\alpha^{(k)}=e_\alpha^k/k!$ and $f_\alpha^{(k)}=f_\alpha^k/k!$. For a finite-dimensional representation $U$ of $\mathfrak g$ in which the $E_{cc}$ act semisimply with integer eigenvalues we write $U_\nu$ for the weight space of $\nu\in\bZ^r$. An $\cO$-lattice $\Lambda\subset U$ is \emph{graded} when $\Lambda=\bigoplus_\nu(\Lambda\cap U_\nu)$, and \emph{admissible} when it is graded and $e_\alpha^{(k)}\Lambda\subset\Lambda$ and $f_\alpha^{(k)}\Lambda\subset\Lambda$ for every positive root $\alpha$ and every $k\ge0$. The \emph{invariant projection} $P_U$ of $U$ is the projection onto $U^{\mathfrak g}$ along the sum of the isotypic components of nontrivial type.

\begin{lemma}\label{w1:lem:invproj}
Let $U$ be as above, let $\Lambda\subset U$ be an admissible lattice, and let $L\ge0$ be an integer such that $k\le L$ whenever $k\ge1$ and $k\alpha$ is a weight of $U$ for a positive root $\alpha$. Then
\[
p^{N}P_U\Lambda\subset\Lambda,\qquad N=\binom r2\big\lfloor\log_p(L+r-1)\big\rfloor .
\]
\end{lemma}

\begin{proof}
The invariant projection kills $U_\nu$ for $\nu\ne0$, because $\mathfrak g$-invariant vectors have weight $0$. Since $\Lambda$ is graded, $P_U\Lambda=P_U(\Lambda\cap U_0)$, and it suffices to treat $x\in\Lambda\cap U_0$.

Let $\mathsf P$ be the extremal projector of Proposition~\ref{w1:prop:extremal}. The weight $0$ is dominant, so $\mathsf P$ is defined on $U_0$ and is the projection of $U_0$ onto $U_0^{\mathfrak n^+}$ along $U_0\cap\mathfrak n^-U$. A vector of $U_0^{\mathfrak n^+}$ is a highest weight vector of weight $0$ or is zero, so it generates a trivial representation and lies in $U^{\mathfrak g}$. Conversely $U^{\mathfrak g}\cap U_0\subset U_0^{\mathfrak n^+}$. In an irreducible representation of highest weight $\lambda\ne0$ every weight vector of weight different from $\lambda$ lies in $\mathfrak n^-$ applied to the representation, so the weight-$0$ part of every nontrivial isotypic component lies in $U_0\cap\mathfrak n^-U$, and $\mathfrak n^-U$ lies in the sum of the nontrivial isotypic components. Hence $\mathsf P=P_U$ on $U_0$.

On $U_0$ the factor of $\alpha=\epsilon_a-\epsilon_b$ in $\mathsf P$ acts as $\sum_{k\ge0}(-1)^k\binom{N_\alpha+k}k^{-1}f_\alpha^{(k)}e_\alpha^{(k)}$ with $N_\alpha=\langle0+\rho,\alpha^\vee\rangle=b-a\le r-1$, as in the proof of Proposition~\ref{w1:prop:extremal}. Each $f_\alpha^{(k)}e_\alpha^{(k)}$ preserves weights, so every factor maps $U_0$ to $U_0$. If $y\in U_0$ and $e_\alpha^{(k)}y\ne0$ with $k\ge1$, then $k\alpha$ is a weight of $U$, so $k\le L$, and Kummer's theorem gives $v\binom{N_\alpha+k}k\le\lfloor\log_p(N_\alpha+k)\rfloor\le\lfloor\log_p(L+r-1)\rfloor$. Since $\Lambda$ is admissible, the factor of $\alpha$ maps $p^{-c}\Lambda\cap U_0$ into $p^{-c-\lfloor\log_p(L+r-1)\rfloor}\Lambda\cap U_0$ for every $c$. There are $\binom r2$ positive roots, so $p^N\mathsf Px\in\Lambda$.
\end{proof}

\begin{lemma}\label{w1:lem:traceformula}
Let $V$ be a finite-dimensional representation of $\mathfrak g$ as above, let $\lambda$ be a dominant weight, let $e_\lambda$ be the projection of $V$ onto its isotypic component of type $\lambda$ along the other isotypic components, and suppose that $V$ contains an irreducible subrepresentation $W$ of highest weight $\lambda$. Let $b_1,\dots,b_n$ be a basis of $W$ with dual basis $b_1^*,\dots,b_n^*$, and let $P$ be the invariant projection of $\Hom(W,V)$, on which $X\in\mathfrak g$ acts by $X\cdot\varphi=X\varphi-\varphi X$. Then for every $m\in V$
\[
e_\lambda(m)=n\sum_{i=1}^nP\big(b_i^*\otimes m\big)(b_i),
\]
where $b_i^*\otimes m$ is the map $w\mapsto b_i^*(w)m$.
\end{lemma}

\begin{proof}
Both sides are linear in $m$, so we may take $m$ in an irreducible subrepresentation $U\subset V$ of highest weight $\mu$. The maps $b_i^*\otimes m$ lie in $\Hom(W,U)$, a subrepresentation of $\Hom(W,V)$, and $P$ preserves it, so $P(b_i^*\otimes m)\in\Hom(W,U)^{\mathfrak g}$. If $\mu\ne\lambda$ this space is $0$ by Schur's lemma, and both sides vanish. If $\mu=\lambda$ choose an isomorphism $\psi\colon W\to U$ of representations. Composition with $\psi^{-1}$ identifies $\Hom(W,U)$ with $\mathrm{End}(W)$ as representations, and the invariant projection of $\mathrm{End}(W)$ is $A\mapsto(\tr A/n)\,\mathrm{id}$, because this map is equivariant, idempotent, has image $\mathrm{End}(W)^{\mathfrak g}=\mathcal K\,\mathrm{id}$, and the invariant projection is the unique equivariant projection onto the invariants. Hence $P(b_i^*\otimes m)=\frac1nb_i^*(\psi^{-1}m)\,\psi$ and
\[
n\sum_iP(b_i^*\otimes m)(b_i)=\sum_ib_i^*(\psi^{-1}m)\,\psi(b_i)=\psi\psi^{-1}m=m=e_\lambda(m).\qedhere
\]
\end{proof}

We return to the setting of Proposition~\ref{w1:prop:projexists}. Let $\mathcal V_d$ be the space of polynomials of multidegree $d$, with $\mathfrak g$ acting through the letters of $S$, and put $\mathcal B_d=\mathcal B\cap\mathcal V_d$ and $\deg_S=\sum_{i\in S}d_i$.

\begin{lemma}\label{w1:lem:Badm}
$\mathcal B_d$ is an admissible lattice of $\mathcal V_d$, every weight $\nu$ of $\mathcal V_d$ satisfies $-\deg_S\le\nu_c\le\deg_S$, and any two weights $\mu,\omega$ of $\mathcal V_d$ satisfy $\lvert\mu_c-\omega_c\rvert\le\deg_S$ for every $c$.
\end{lemma}

\begin{proof}
A monomial $x^\beta$ is a weight vector, of weight $\nu_c=\sum_{i\in S\ \mathrm{vector}}\beta_{i,c}-\sum_{i\in S\ \mathrm{covector}}\beta_{i,c}$, and $\mathcal B_d$ is the direct sum of the lines $\{c\,x^\beta:v(c)+\frac12v(\beta!)\ge0\}$ over the monomials of multidegree $d$. So $\mathcal B_d$ is a graded lattice. It is stable under $E_{ab}^{(k)}$ for $a\ne b$ by Lemma~\ref{w1:lem:divpow}. For the weight bounds, $\sum_{i\in S\ \mathrm{vector}}\beta_{i,c}$ lies between $0$ and $\sum_{i\in S\ \mathrm{vector}}d_i$, and $\sum_{i\in S\ \mathrm{covector}}\beta_{i,c}$ lies between $0$ and $\sum_{i\in S\ \mathrm{covector}}d_i$, so $\nu_c$ lies in an interval of length $\deg_S$ that contains $0$.
\end{proof}

\begin{theorem}[the norm bound (B7)]\label{w1:thm:B7}
In the setting of Proposition~\ref{w1:prop:projexists}, put
\[
N_S=\binom r2\big\lfloor\log_p(\deg_S+r-1)\big\rfloor .
\]
Then $p^{N_S}\pi_{S,\theta}\,\mathcal B_d\subset\mathcal B_d$. Equivalently $v_{\mathrm{BB}}(\pi_{S,\theta}F)\ge-N_S$ for every $F\in\mathcal V_d$ with $v_{\mathrm{BB}}(F)\ge0$. If the value group $v(\mathcal K^\times)$ contains $\frac12\bZ$, which holds when $\cO$ contains an element $\varpi$ with $\varpi^2=p$, then
\[
v_{\mathrm{BB}}(\pi_{S,\theta}F)\ge v_{\mathrm{BB}}(F)-N_S\ge v_{\mathrm{BB}}(F)-\binom r2\log_p(2\deg_S+r)\qquad(F\in\mathcal V_d).
\]
\end{theorem}

\begin{proof}
By the proof of Proposition~\ref{w1:prop:projexists}, $\pi_{S,\theta}=\sum_{\lambda\in\Sigma}e_\lambda$, where $\Sigma$ is the finite set of dominant weights $\lambda$ with $\lambda_r\ge\theta$ for which $\mathcal V_d$ has a nonzero isotypic component of type $\lambda$, and $e_\lambda$ is the isotypic projection. Since $\mathcal B_d$ is an $\cO$-module it suffices to show $p^{N_S}e_\lambda\mathcal B_d\subset\mathcal B_d$ for each $\lambda\in\Sigma$.

Fix $\lambda\in\Sigma$ and an irreducible subrepresentation $W\subset\mathcal V_d$ of highest weight $\lambda$. Put $\Lambda_W=\mathcal B_d\cap W$. Since $W$ is spanned by weight vectors and $\mathcal B_d$ is graded, $\Lambda_W=\bigoplus_\omega(\mathcal B_d\cap W_\omega)$, and $\Lambda_W$ is stable under the divided powers $e_\alpha^{(k)}$, $f_\alpha^{(k)}$. It is a free $\cO$-module of rank $n=\dim W$. Choose an $\cO$-basis $b_1,\dots,b_n$ of $\Lambda_W$, which is a $\mathcal K$-basis of $W$, with dual basis $b_i^*$. Let
\[
\Lambda_H=\{\varphi\in\Hom(W,\mathcal V_d):\ \varphi(\Lambda_W)\subset\mathcal B_d\},
\]
an $\cO$-lattice in $\Hom(W,\mathcal V_d)$, isomorphic to $\Hom_\cO(\Lambda_W,\mathcal B_d)$.

$\Lambda_H$ is graded. For $\varphi\in\Lambda_H$ its component of weight $\nu$ is $\varphi_\nu=\sum_\omega\mathrm{pr}_{\omega+\nu}\circ\varphi\circ\mathrm{pr}_\omega$, where $\mathrm{pr}$ denotes the weight projections. These preserve $\mathcal B_d$ and $\Lambda_W$ because both are graded, so $\varphi_\nu\in\Lambda_H$.

$\Lambda_H$ is stable under divided powers. The operators $L_X\varphi=X\varphi$ and $R_X\varphi=\varphi X$ commute, and $X\cdot\varphi=(L_X-R_X)\varphi$, so for $X=e_\alpha$ or $f_\alpha$
\[
X^{(k)}\cdot\varphi=\sum_{i+j=k}(-1)^j\,X^{(i)}\circ\varphi\circ X^{(j)} ,
\]
and each term maps $\Lambda_W$ into $\mathcal B_d$ because $X^{(j)}\Lambda_W\subset\Lambda_W$ and $X^{(i)}\mathcal B_d\subset\mathcal B_d$ (Lemma~\ref{w1:lem:Badm}).

The weights of $\Hom(W,\mathcal V_d)\cong W^*\otimes\mathcal V_d$ are the differences $\mu-\omega$ of a weight $\mu$ of $\mathcal V_d$ and a weight $\omega$ of $W$. The weights of $W$ are weights of $\mathcal V_d$, so if $k\alpha=\mu-\omega$ with $\alpha=\epsilon_a-\epsilon_b$ then $k=\mu_a-\omega_a\le\deg_S$ by Lemma~\ref{w1:lem:Badm}. Lemma~\ref{w1:lem:invproj} with $U=\Hom(W,\mathcal V_d)$, $\Lambda=\Lambda_H$ and $L=\deg_S$ gives $p^{N_S}P\Lambda_H\subset\Lambda_H$.

Let $m\in\mathcal B_d$. Then $b_i^*\otimes m\in\Lambda_H$, because $b_i^*(\Lambda_W)\subset\cO$. Hence $p^{N_S}P(b_i^*\otimes m)\in\Lambda_H$, and since $b_i\in\Lambda_W$, the vector $p^{N_S}P(b_i^*\otimes m)(b_i)$ lies in $\mathcal B_d$. Lemma~\ref{w1:lem:traceformula} gives $p^{N_S}e_\lambda(m)=n\sum_ip^{N_S}P(b_i^*\otimes m)(b_i)\in\mathcal B_d$, since $n$ is an integer.

For the second form, $v_{\mathrm{BB}}(cF)=v(c)+v_{\mathrm{BB}}(F)$ and $\mathcal B_d=\{F:v_{\mathrm{BB}}(F)\ge0\}$. The values $v(F_\beta)+\frac12v(\beta!)$ lie in $v(\mathcal K^\times)+\frac12\bZ$, which is $v(\mathcal K^\times)$ under the assumption, so for $F\ne0$ there is $c$ with $v(c)=-v_{\mathrm{BB}}(F)$, and the first statement applied to $cF$ gives the inequality. Finally $\lfloor\log_p(\deg_S+r-1)\rfloor\le\log_p(2\deg_S+r)$.
\end{proof}

\begin{remark}\label{w1:rem:B7audit}
In (B7) the loss is $\binom r2\lfloor\log_p(\deg_S+r-1)\rfloor$. Since $\deg_S\le\sum_xd_x$, the bound $e_P=r^2(\log_p(2\sum_xd_x+r)+1)$ in (B8) applies to each projector.
\end{remark}



\subsection{Engine (A): the rank of the symbol spaces}\label{w1:sec:A}

The statement is Hypothesis~\ref{r2:engines}(A). This section proves the algebraic steps of its proof and the analytic lemmas of the concentration estimate, and reduces the standard monomial theory to the literature.

\subsubsection{Lifting from the residue field}

\begin{lemma}\label{w1:lem:nakayama}
Let $\cO$ be a discrete valuation ring with uniformizer $\varpi$ and residue field $\bF$, let $\mathcal M$ be a free $\cO$-module with basis $(x^\mu)_{\mu\in I}$, $I$ finite, and let $\mathcal G\subset\mathcal M$. Choose $g_1,\dots,g_n\in\mathcal G$ whose reductions form a basis of the $\bF$-span of the reductions of $\mathcal G$. Then $g_1,\dots,g_n$ are linearly independent over $\cO$ and span a direct summand of $\mathcal M$. In particular $\rk_\cO\spn_\cO\mathcal G\ge\dim_\bF\spn_\bF\overline{\mathcal G}$.
\end{lemma}

\begin{proof}
The reductions $\bar g_1,\dots,\bar g_n$ are independent, so the $n\times\lvert I\rvert$ coefficient matrix of the $g_i$ has an $n\times n$ minor that is a unit. Complete $g_1,\dots,g_n$ by the basis vectors $x^\mu$ outside the columns of that minor. The resulting square matrix has unit determinant, so the completed family is a basis of $\mathcal M$.
\end{proof}

\subsubsection{The Frobenius identity and the discriminant}

\begin{lemma}\label{w1:lem:frob}
Let $\bF$ be a field of characteristic $p$ and $f\in\bF_p[X_1,\dots,X_N]$. Then $f(X_1^{p^R},\dots,X_N^{p^R})=f(X_1,\dots,X_N)^{p^R}$. If $D\in\bF[X]$ is nonzero, multiplication by $D$ is injective on $\bF[X]$, and for every finite family $(u_\iota)$ in $\bF[X]$, $\dim\spn(Du_\iota)=\dim\spn(u_\iota)$.
\end{lemma}

\begin{proof}
The Frobenius $F\mapsto F^p$ is a ring endomorphism of $\bF[X]$ fixing $\bF_p$, so $f(X)^{p^R}=f^{(\sigma^R)}(X^{p^R})=f(X^{p^R})$, where $f^{(\sigma^R)}$ has coefficients raised to the power $p^R$, equal to those of $f$. The polynomial ring is a domain.
\end{proof}

In the construction of the segment generators the coarse segment polynomials have integer coefficients, so their reductions have coefficients in $\bF_p$, and Lemma~\ref{w1:lem:frob} identifies the reduction of a coarse placement with exponent $p^{\mathrm{dist}+R_j}$ with the $p^{R_j}$-th power of the true placement. The discriminant $D_\ell$ is a product of powers of $r\times r$ determinants in pairwise distinct sets of variables, so its reduction is a product of nonzero polynomials and is nonzero. Then Lemma~\ref{w1:lem:frob} gives $\dim\spn\overline{\mathrm{gen}}\ge\dim\spn(\overline{D_\ell}\cdot\overline{\mathrm{coarse}})=\dim\spn\overline{\mathrm{coarse}}$, provided each $\overline{D_\ell}\,\bar u$ for a coarse generator $u$ is the reduction of a true generator.

The passage from coarse to true generators uses the following statement.

\begin{enumerate}[label=(A\arabic*),start=2,leftmargin=*]
\item \emph{Coarse generators become true generators}. For every coarse generator $u$ there is a true generator $g$ with $D_\ell u\equiv g$ modulo $\varpi$. The mechanism is Lemma~\ref{w1:lem:frob} together with the additivity of $\mathrm{val}_S$ and $\lambda_S$ on products and the fact that a determinant of $r$ plane vectors has type $(1,\dots,1)$, so that multiplication by $\det^{\lfloor B_o/(k^2Y)\rfloor}$ raises the lowest plane type by the margin.
\end{enumerate}

\notgiven{Statement (A2) is not proved here.}

\subsubsection{Digit decomposition}

\begin{lemma}\label{w1:lem:digits}
Let $P_1<P_2<\dots<P_J$ be powers of $p$, each dividing the next. Every $\mathbf e\in\bZ_{\ge0}^N$ has at most one expression $\mathbf e=\boldsymbol\ell+\sum_{j=1}^JP_j\mathbf m_j$ with $0\le\ell_t<P_1$ and with the $t$-th coordinate of $\mathbf m_j$ less than $P_{j+1}/P_j$ for $j<J$. Consequently, if for each $j$ the polynomials $(u_{j,\iota})_\iota$ in $\bF[X]$ are linearly independent and $\bF$ has characteristic $p$, the products $X^{\boldsymbol\ell}\prod_ju_{j,\iota_j}(X^{P_j})$, over labels $\boldsymbol\ell$ and families $(\iota_j)$ for which every monomial obeys these digit bounds, are linearly independent.
\end{lemma}

\begin{proof}
Uniqueness is that of the mixed-radix expansion. For independence, the monomials of $X^{\boldsymbol\ell}\prod_ju_{j,\iota_j}(X^{P_j})$ are $X^{\boldsymbol\ell+\sum_jP_j\mathbf m_j}$ with $X^{\mathbf m_j}$ a monomial of $u_{j,\iota_j}$, and distinct tuples $(\boldsymbol\ell,\mathbf m_1,\dots,\mathbf m_J)$ give distinct exponents. So the span is the tensor product of the spans of the $(u_{j,\iota})_\iota$ over the label monomials, and a tensor product of independent families is independent.
\end{proof}



\begin{lemma}\label{w1:lem:telescope}
Let the frame tops satisfy $f_{j+J,i}=f_{j,i}+n'$ and let $R_j$ be $J$-periodic. Put $Q_j=p^{f_{j+1,r-1}-f_{j,r-1}+R_{j+1}-R_j}$. Then $\prod_{j=0}^{J-1}Q_j=p^{n'}$.
\end{lemma}

\begin{proof}
The exponent telescopes to $f_{J,r-1}-f_{0,r-1}+R_J-R_0=n'$.
\end{proof}

\begin{lemma}\label{w1:lem:gauss}
For $1\le r<k$ and a prime power $q$, the number of $r$-dimensional subspaces of $\bF_q^k$ is at least $q^{r(k-r)}$.
\end{lemma}

\begin{proof}
The subspaces that are graphs of linear maps $\bF_q^r\to\bF_q^{k-r}$ are distinct and number $q^{r(k-r)}$.
\end{proof}

\subsubsection{Standard monomials}

Engine (A) uses the following two statements on standard monomials.

\begin{enumerate}[label=(A\arabic*),start=4,leftmargin=*]
\item \emph{(Published.)} \emph{Integral standard monomial bases.} Let $\cO$ be a discrete valuation ring. The module $\bigotimes_i\mathrm{Sym}^{b_i}(\cO^r)$, with duals at covector letters, has an $\cO$-basis of products of column polynomials $\Psi(\mu,I)$ indexed by Pieri paths $0=\mu_0,\mu_1,\dots$ (each step adds a horizontal strip of size $b_i$ at a vector letter, or removes one at a covector letter) and semistandard fillings $I$ of the last shape. For vector letters alone this is the standard monomial theory of Doubilet, Rota and Stein, \emph{On the foundations of combinatorial theory IX}, Stud. Appl. Math.~53 (1974), in the characteristic-free form of De Concini, Eisenbud and Procesi, \emph{Young diagrams and determinantal varieties}, Invent. Math.~56 (1980), 129--165, and of D\'esarm\'enien, Kung and Rota, \emph{Invariant theory, Young bitableaux, and combinatorics}, Adv. Math.~27 (1978). Mixed vector and covector letters are covered by the mixed tensor version of De Concini and Procesi, \emph{A characteristic free approach to invariant theory}, Adv. Math.~21 (1976).
\item \emph{Compatibility with prefix types}. The submodule cut out by the prefix conditions $\mathrm{val}_{\{y\le x\}}\ge\theta_x$ is spanned over $\cO$ and over the residue field by the basis elements whose path satisfies the same inequalities, and the leading part of $\Psi$ for a weight separating plane and normal coordinates is $\Psi_{\mathrm{pl}}\Psi_{\mathrm{nm}}$. This is a filtration property of standard monomials along the letters, classical in substance (the straightening law respects the order of the letters), and the exact integral form required is not matched to a published statement here.
\end{enumerate}

Statement (A4) is published, with the sources named in it. \notgiven{Statement (A5) is not proved here.}

\subsubsection{The concentration estimate}

For integers $k'\ge1$ define $f_{k'}(y)=\log\Gamma(y+k'+1)-\log\Gamma(y+1)-\log k'!$ for real $y\ge0$, so that $f_{k'}(y)=\log\binom{y+k'}{k'}$ for integers $y\ge0$.

\begin{lemma}\label{w1:lem:fT}
For $y\ge0$, $f_{k'}'(y)=\sum_{i=1}^{k'}\frac1{y+i}$. Hence $f_{k'}$ is increasing and concave, $\frac{k'}{y+k'}\le f'_{k'}(y)\le\frac{k'}{y+1}$, $f'_{k'}(y)\ge\frac{k'}{2y}$ for $y\ge k'$, $f'_{k'}(y')\le\frac12f'_{k'}(y)$ for $y'\ge4y\ge4k'$, and $f_{k'}(x)-f_{k'}(y)\le f'_{k'}(y)(x-y)$ for $x,y\ge0$.
\end{lemma}

\begin{proof}
$\Gamma(y+k'+1)/\Gamma(y+1)=\prod_{i=1}^{k'}(y+i)$, so $f_{k'}(y)=\sum_i\log(y+i)-\log k'!$. Differentiate. Each term $1/(y+i)$ lies between $1/(y+k')$ and $1/(y+1)$. For $y\ge k'$, $k'/(y+k')\ge k'/(2y)$. For $y'\ge4y\ge4k'$,
\[
f'_{k'}(y')\le\frac{k'}{y'}\le\frac{k'}{4y}=\frac12\cdot\frac{k'}{2y}\le\frac12f'_{k'}(y).
\]
The derivative is decreasing, which is concavity, and concavity gives the last inequality.
\end{proof}

\begin{lemma}\label{w1:lem:deficit}
Let $k'\ge1$ and $L\ge1$, and let $\alpha_0,\dots,\alpha_{L-1}\ge0$ and $B_0,\dots,B_{L-1}$ be reals with $B_l\ge k'$ and $4B_{l+1}\le B_l$. Suppose $\sum_{l\ge u}\alpha_l\le\sum_{l\ge u}B_l$ for every $u<L$ (sums over $u\le l<L$). Put $\delta_u=f'_{k'}(B_u)-f'_{k'}(B_{u-1})$ with $f'_{k'}(B_{-1})=0$. Then $\delta_u\ge\frac12f'_{k'}(B_u)>0$ and
\[
\sum_lf_{k'}(\alpha_l)-\sum_lf_{k'}(B_l)\le\sum_u\delta_u\Big(\sum_{l\ge u}\alpha_l-\sum_{l\ge u}B_l\Big)\le0 .
\]
If moreover $\sum_{l\ge u_0}(\alpha_l-B_l)\le-X$ with $X\ge0$, then $\sum_lf_{k'}(\alpha_l)-\sum_lf_{k'}(B_l)\le-\frac{Xk'}{4B_{u_0}}$.
\end{lemma}

\begin{proof}
Since $B_{u-1}\ge4B_u\ge4k'$, Lemma~\ref{w1:lem:fT} gives $f'(B_{u-1})\le\frac12f'(B_u)$, so $\delta_u\ge\frac12f'(B_u)$. By concavity $f(\alpha_l)-f(B_l)\le f'(B_l)(\alpha_l-B_l)$. Write $f'(B_l)=\sum_{u\le l}\delta_u$ and exchange the sums (summation by parts),
\[
\sum_lf'(B_l)(\alpha_l-B_l)=\sum_u\delta_u\sum_{l\ge u}(\alpha_l-B_l).
\]
Each $\delta_u>0$ and each tail is at most $0$. Keeping the term $u_0$ alone gives at most $\delta_{u_0}(-X)\le-\frac X2f'(B_{u_0})\le-\frac{Xk'}{4B_{u_0}}$ by Lemma~\ref{w1:lem:fT}.
\end{proof}

The lemma says that a Pieri path whose plane part takes $X$ more units of degree than the typical window, measured from a tail, loses a factor $e^{-Xk'/(4B)}$ in the dimension of its normal part.

The remaining steps of engine (A) are the following.

\begin{enumerate}[label=(A\arabic*),start=6,leftmargin=*]
\item \emph{The remaining inputs of the concentration estimate}. The first is a two-sided estimate of $\log\dim\nabla(\lambda)$ for a dominant weight of $\mathrm{GL}_n$ with at most $N_s$ positive and $N_s$ negative entries, $\lvert\log\dim\nabla(\lambda)-\sum_mf_{K-2N_s}(\lambda_m^+)-\sum_mf_{K-2N_s}(\lambda_m^-)\rvert\le(2N_sr+8N_s^2)\log(1+A_0)$. It follows from Weyl's dimension formula $\prod_{i<j}\frac{\lambda_i-\lambda_j+j-i}{j-i}$, in which the product of the factors pairing a nonzero entry with the block of zero entries is $\binom{\lambda_i^\pm+k'}{k'}$ up to $O(N_s)$ factors of size at most $1+A_0$. The proof of the stated constants is not written here. The others are elementary properties of Pieri paths (entries bounded by the cumulative degrees, horizontal strips interlace, the number of paths is at most $(1+2aN_s)^{2N_s^2}$).
\item \emph{Path-weight concentration}. An atypical Pieri path has weight at most $\exp\big(9N_s^2\log(1+2aN_s)-\min(\frac{L_ak'}{16k},\frac{k'}{12})\big)\dim$, and the total weight of atypical paths is at most $a^{-N_s^2}\dim$ when $L_a,k\ge1088N_s^2\log a$. The mechanism is Lemma~\ref{w1:lem:deficit} applied to the tails of an atypical path, with (A6) giving the comparison between path weights and sums of $f_{k'}$. The assembly is not written here.
\item[(A8)] Lemmas~\ref{w1:lem:telescope} and \ref{w1:lem:gauss}: the frame deficit over a period of segments is $q^{rK}$, and $q^{rK}\le\#\mathrm{Gr}_r(\bF_q^k)$.
\item[(A9)] \emph{Assembly of the counts}. The losses $(1-5/Y)$, $(1-1/Y)$, $(1-r/(2Y))$ per top and $(1-1/(32Y))$ per segment multiply to at least $1-\lvert T\rvert r^2k/Y$. The final inequality is elementary once the per-top factors are established, and the per-top factors are not written here.
\end{enumerate}

\notgiven{The construction in (A6), (A7) and the individual share bounds in (A9) remain open.}

\subsubsection{The final count of engine (A)}\label{w1:sec:A9final}

\begin{lemma}\label{w1:lem:A9}
Let $r\ge2$, $k\ge5$, $Y\ge1$, and let $J\le N$ be positive integers with $Nr^2k<Y$. Let $A$ be a real number with $1-\frac1{32Y}\le A\le1$. Then the numbers $1-\frac5Y$, $1-\frac1Y$, $1-\frac r{2Y}$ lie in $[0,1]$ and
\[
A^J\Big(\big(1-\tfrac5Y\big)\big(1-\tfrac1Y\big)\big(1-\tfrac r{2Y}\big)\Big)^N\ge1-\frac{Nr^2k}Y .
\]
\end{lemma}

\begin{proof}
Put $u=1/Y$. From $Nr^2ku<1$ and $r^2k\ge20$ we get $u<\frac1{20}$, so $5u\le1$ and $\frac r2u\le\frac{r^2k}{2}u<1$, and $A\ge1-\frac u{32}\ge0$. For numbers $x_1,\dots,x_M\in[0,1]$ one has $\prod_i(1-x_i)\ge1-\sum_ix_i$, by induction on $M$, since $(1-x)(1-y)\ge1-x-y$ for $y\le1$ and $1-\sum$ may be negative. Applying this to $J$ factors $1-\frac u{32}$ and $N$ triples of factors gives a lower bound $1-u\big(\frac J{32}+N(6+\frac r2)\big)\ge1-uN\big(6+\frac1{32}+\frac r2\big)$, because $J\le N$ and $A\ge1-\frac u{32}$. Finally $r^2k-\frac r2\ge r(5r-\frac12)\ge19\ge6+\frac1{32}$.
\end{proof}

In (A9) the lemma applies with $N=\lvert T\rvert$ (the tops of the orbit), $J$ the number of segments over a period and $A=1-\frac1{32Y}$, once the per-top factors $(1-5/Y)$, $(1-1/Y)$, $(1-r/(2Y))$ and the per-segment factor $(1-1/(32Y))$ are established. The per-top factors are the residual of (A9).



\subsubsection{The comparison exponent of (L2)}\label{r2:sec:L2}

For a dominant weight $\lambda$ of $\GL_k$ let $V(\lambda)$ be the irreducible representation over $\mathcal K$, $v_\lambda$ a highest weight vector and $v_\lambda^*$ the linear form that is $1$ on $v_\lambda$ and $0$ on the other weight spaces. Let $U_\bZ$ be the $\bZ$-subalgebra of $U(\mathfrak{gl}_k)$ generated by the $E_{ab}^{(n)}$ ($a\ne b$) and the $\binom{E_{cc}}n$. Put
\[
\Delta_\cO(\lambda)=U_\bZ\,\cO v_\lambda,\qquad\nabla_\cO(\lambda)=\{x\in V(\lambda):\ v_\lambda^*(ux)\in\cO\ \text{for all }u\in U_\bZ\},
\]
and $w(\lambda)=\binom k2\lfloor\log_p(\lambda_1-\lambda_k+k-1)\rfloor$.

\begin{proposition}\label{r2:prop:nabladelta}
$p^{w(\lambda)}\nabla_\cO(\lambda)\subset\Delta_\cO(\lambda)\subset\nabla_\cO(\lambda)$.
\end{proposition}

\begin{proof}
Write $\Delta=\Delta_\cO(\lambda)$, $V=V(\lambda)$, $V^*$ its dual with the contragredient action, and $\Delta'=U_\bZ\,\cO v_\lambda^*\subset V^*$. The algebra $U_\bZ$ is spanned by products $f^{(K)}h\,e^{(K')}$ of ordered divided powers of the negative root vectors, binomials in the $E_{cc}$ and divided powers of the positive root vectors (Kostant). The positive ones kill $v_\lambda$, the binomials act on $v_\lambda$ by integers, and the negative ones lower the weight, so $\Delta\cap V_\lambda=\cO v_\lambda$ and $v_\lambda^*(\Delta)\subset\cO$. Since $\Delta$ is $U_\bZ$-stable this gives $\Delta\subset\nabla_\cO(\lambda)$.

By definition, $x\in\nabla_\cO(\lambda)$ exactly when $\eta(x)\in\cO$ for every $\eta\in\Delta'$, because $\eta=u^{\mathsf c}v_\lambda^*$ with $u^{\mathsf c}$ the image of $u$ under the antipode, which preserves $U_\bZ$. So $\nabla_\cO(\lambda)$ is the $\cO$-dual of $\Delta'$.

Let $U=V\otimes V^*=\mathrm{End}(V)$ and $\Lambda=\Delta\otimes_\cO\Delta'$. The lattices $\Delta$ and $\Delta'$ are spanned by weight vectors and stable under divided powers, and $X^{(n)}(y\otimes\eta)=\sum_{i+j=n}X^{(i)}y\otimes X^{(j)}\eta$, so $\Lambda$ is admissible. The weights of $U$ are differences $\mu-\mu'$ of weights of $V$, whose entries lie between $\lambda_k$ and $\lambda_1$, so $k\alpha=\mu-\mu'$ forces $k\le\lambda_1-\lambda_k$. Lemma~\ref{w1:lem:invproj} with $r=k$ and $L=\lambda_1-\lambda_k$ gives $p^{w(\lambda)}P_U\Lambda\subset\Lambda$. Since $V$ is irreducible, $P_U(A)=\frac{\tr A}{\dim V}\mathrm{id}$, and $v_\lambda\otimes v_\lambda^*\in\Lambda$ has trace $1$. So $\frac{p^{w(\lambda)}}{\dim V}\mathrm{id}\in\Lambda$.

An element $\sum_jy_j\otimes\eta_j$ of $\Lambda$ maps $x\in\nabla_\cO(\lambda)$ to $\sum_j\eta_j(x)y_j\in\Delta$, since $\eta_j(x)\in\cO$. Applied to $\frac{p^{w(\lambda)}}{\dim V}\mathrm{id}$ this gives $p^{w(\lambda)}x\in\dim V\cdot\Delta\subset\Delta$.
\end{proof}

\begin{corollary}[input (L2)]\label{r2:cor:L2}
Suppose that, in (F3), $M$ is the lattice of forms with $F(L)\in\bigotimes_v\bigotimes_s\Delta_\cO(\lambda^*_{v,s};L_v^{(s)})\otimes H'(L)$ and $M^{\mathrm{na}}$ the lattice of forms with values in $\bigotimes_v\bigotimes_s\nabla_\cO(\lambda^*_{v,s};L_v^{(s)})\otimes H'(L)$, where the lattices at $L_v^{(s)}$ are the transports of those at $\cO^k$ by any $g$ with $g\cO^k=L_v^{(s)}$. Then $M\subset M^{\mathrm{na}}$ and $p^{e}M^{\mathrm{na}}\subset M$ with $e=\sum_v\sum_{s\in\bZ/n_p}w(\lambda^*_{v,s})$. If moreover every base weight satisfies $\lambda_1-\lambda_k+k-1\le a^2$ and the positions $(v,s)$ number at most $f$, then $e\le fk^2\log_pa$, and (L2) holds with $e_p=\lceil C_efk^2\log_pa\rceil$ for any $C_e\ge1$, so that $e_p\le(C_e+1)fk^2\log_pa$ when $fk^2\log_pa\ge1$.
\end{corollary}

\begin{proof}
Proposition~\ref{r2:prop:nabladelta} transports along $g$, because both lattices are defined from $U_\bZ$ and $g$ intertwines the actions. For lattices $A_s\subset B_s$ with $p^{w_s}B_s\subset A_s$ the tensor products satisfy $\bigotimes A_s\subset\bigotimes B_s$ and $p^{\sum w_s}\bigotimes B_s\subset\bigotimes A_s$, since both are spanned by pure tensors. The conditions on forms are pointwise and equivariance is preserved by scalars, which gives the first claim. For the bound, $\binom k2\le\frac{k^2}2$ and $\lfloor\log_p(\lambda_1-\lambda_k+k-1)\rfloor\le2\log_pa$, so each $w(\lambda^*_{v,s})\le k^2\log_pa$, and there are at most $f$ terms. Since $e\le C_efk^2\log_pa\le e_p$, also $p^{e_p}M^{\mathrm{na}}\subset M$.
\end{proof}

For the base weights of Theorem~\ref{r2:LPthm} the extreme entries are $a_s\le2a$ and $-b_s\ge-2b'$ at a top, and the other entries have absolute value at most $c_G$, so $\lambda_1-\lambda_k+k-1\le4a+2c_G+k\le a^2$ once $c_G+k\le a$. The positions $(v,s)$ are the elements of $H$, so their number is $n\le f$.

\subsubsection{The bound for $\Pi$ in (Z1)}\label{r2:sec:Z1}

Put $\eta=W+mh_0+1$, with $W$ the number of roots of unity of $E$ and $h_0=2\lceil x^{402}\rceil$ as in Section~\ref{sec:proofA}.

\begin{lemma}\label{r2:lem:W123}
The nodes $w_{u_0,j}$ of Proposition~\ref{r6:R6:prop:nodes} satisfy (W1), (W2) and (W3) with this $\eta$, and the bottom weights $b_t$ of that proposition.
\end{lemma}

\begin{proof}
At a top place, $w_{u_0,j}(0,0,t)=b_t+o_j(t)$ with $o_j(t)=\lfloor b'/x\rfloor+(j-\frac12)h_0$. Since $h_0\ge2$, $o_j(t)\ge b'/x-1+\frac12h_0\ge b'/x$, and $o_j(t)\le b'/x+mh_0\le b'/x+\eta$. At a thick place, $\lceil b\rceil\le b_1(\rho)<\lceil b\rceil+W$ gives $b\le w_{u_0,j}(\rho)\le b+1+W+mh_0=b+\eta$. At a thin place, $b_1(\rho)=\lceil w''\rceil$ gives $w''\le w_{u_0,j}(\rho)\le w''+1+mh_0\le w''+\eta$.
\end{proof}

\begin{proposition}\label{r2:prop:Z1}
Suppose that $c_G\le x^{C_c}$ for a constant $C_c$. Then $\Pi\le\exp(x^{\varepsilon/4})$ for $x\ge x_0(\varepsilon,C_c)$.
\end{proposition}

\begin{proof}
By Lemma~\ref{r6:R6:lem:tr}, $t_r\le x^{401}$. By Proposition~\ref{ranks}, $k\le2x^{9/2+10\varepsilon}$ and $m\le2\sqrt x$, so $N_m=k+m(t_r+1)\le x^{402}$ for large $x$. By the remark after Proposition~\ref{nodes}, $W\le32n^2\le128x^{9/2}$, so $\eta\le128x^{9/2}+2\sqrt x(2x^{402}+2)+1\le x^{403}$ for large $x$. The constants $K_\bullet$ and $C_\bullet$ depend on $\varepsilon$ alone. Hence $\Pi\le\max(x^{403},x^{C_c},K)$ with $K$ depending on $\varepsilon$, which is at most $\exp(x^{\varepsilon/4})$ for large $x$.
\end{proof}



%%%%%%%%%%%%%%%%%%%%%%%%%%%%%%%%%%%%%%%%%%%%%%%%%%%%%%%%%%%%%%%%%%%%%%%%%%%%%%%%%%%%%%%
% BEGIN BODY
%%%%%%%%%%%%%%%%%%%%%%%%%%%%%%%%%%%%%%%%%%%%%%%%%%%%%%%%%%%%%%%%%%%%%%%%%%%%%%%%%%%%%%%

\section{The slope inequality for the family of period maps}\label{r3:sec}

Hypothesis \eqref{A3} of Proposition~\ref{acct} asks that the rank-weighted mean of $\mu(N_u)-h_u$ over the indices $u\in W$ exceed the slope of the dual $R^\vee$ of the lattice of automorphic forms by at most $e_3$. For one injective map Bost's slope inequality bounds $\mu(N_u)-h_u$ by the maximal slope of $R^\vee$. The family needs more, because $\mu_{\max}(R^\vee)$ can exceed the mean slope $\mu(R^\vee)$ by far, and because the images of the maps $B_u$ overlap. The argument of this section replaces the maximal slope by the mean slope. It cuts each source lattice along its Harder--Narasimhan filtration, moves each piece into one step of the filtration of $R^\vee$ by the slope inequality, and bounds the rank of that step from below by a count of dimensions at one complex place, where the matrices of the maps form a nearly tight frame. Section~\ref{r3:sec:slopes} collects the slope theory, Section~\ref{r3:sec:frame} proves the frame theorem, Section~\ref{r3:sec:gram} derives the frame conditions from Hilbert--Schmidt bounds for Gram operators, Section~\ref{r3:sec:A3} proves \eqref{A3} with an explicit $e_3$ and bounds $e_3$, and Section~\ref{r3:sec:S} states the near-orthogonality theorem for Bessel Poincar\'e series that the construction has to provide.

\subsection{Hermitian lattices and slopes}\label{r3:sec:slopes}

Let $K$ be a totally complex number field of degree $n_K$, let $d_K$ be its discriminant, and let $\Sigma_K$ be its set of complex places, so that $\lvert\Sigma_K\rvert=n_K/2$. For each $v\in\Sigma_K$ fix one of the two embeddings $\sigma_v\colon K\to\mathbb C$ in the class $v$.

A \emph{hermitian lattice} $\bar L$ over $\OK$ consists of a finite-dimensional $K$-vector space $V$, a finitely generated $\OK$-submodule $L\subset V$ that spans $V$, and for each $v\in\Sigma_K$ a hermitian norm $\lVert\cdot\rVert_v$ on $V_v=V\otimes_{K,\sigma_v}\mathbb C$. We write $L$ for $\bar L$ when the norms are fixed, and $\rk L=\dim_KV$. A nonzero finitely generated $\OK$-submodule $N\subset L$ is a hermitian lattice in $KN$ with the restricted norms. It is \emph{saturated} when $N=L\cap KN$. For a saturated proper submodule $N\subset L$ the quotient $L/N$ is a hermitian lattice in $V/KN$, with the quotient norms $\lVert x+KN\rVert_v=\min\{\lVert y\rVert_v:y\in x+(KN)_v\}$.

The real vector space $V_{\mathbb R}=\bigoplus_{v\in\Sigma_K}V_v$, with the euclidean structure whose square norm is $\sum_v\lVert x_v\rVert_v^2$, has dimension $n_K\rk L$, and the map $\rho\colon V\to V_{\mathbb R}$, $x\mapsto(\sigma_v(x))_v$, sends $L$ onto a lattice of full rank. Put $\deg_{\mathbb Z}L=-\log\covol\rho(L)$ and define the \emph{slope}
\begin{equation}\label{r3:mu}
\mu(L)=\frac{2\deg_{\mathbb Z}L}{n_K\rk L}-\log2+\frac{\log\lvert d_K\rvert}{n_K}.
\end{equation}
Since $\covol\rho(\OK)=2^{-n_K/2}\lvert d_K\rvert^{1/2}$ for the norms $\lvert\sigma_v(\cdot)
vert$, the slope of $\OK$ is $0$. Dividing every norm by $c>0$ raises $\mu$ by $2\log c$. So $\mu$ is twice the normalized Arakelov slope $\widehat\mu$ of Bost~\cite{bost}, and it is measured in units of $\log\lVert\cdot\rVert^2$. All slopes in Part~II are slopes in the sense of \eqref{r3:mu}.

The \emph{maximal slope} $\mu_{\max}(L)$ is the supremum of $\mu(N)$ over the nonzero submodules $N\subset L$, and the \emph{minimal slope} $\mu_{\min}(L)$ is the infimum of $\mu(L/N)$ over the saturated proper submodules $N\subset L$, with $N=0$ allowed. For a $K$-linear map $\varphi\colon V\to V'$ between spaces with hermitian norms at the places of $\Sigma_K$, the \emph{height} of $\varphi\ne0$ is
\begin{equation}\label{r3:ht}
\hgt(\varphi)=\frac2{n_K}\sum_{v\in\Sigma_K}\log\lVert\varphi_v\rVert_{\mathrm{op}}^2 ,
\end{equation}
with $\varphi_v=\varphi\otimes_{\sigma_v}\mathbb C$. It is twice the height $\frac1{n_K}\sum_\sigma\log\lVert\varphi_\sigma\rVert$ over all embeddings $\sigma$. The \emph{dual} of $\bar L$ is the lattice $L^\vee=\{f\in V^*:f(L)\subset\OK\}$ in $V^*=\operatorname{Hom}_K(V,K)$, with the dual norms $\lVert f\rVert_v=\max_{x\ne0}\lvert f_v(x)\rvert/\lVert x\rVert_v$.

We use the following classical facts.

\begin{proposition}[Harder--Narasimhan filtration]\label{r3:HN}
Let $\bar L$ be a hermitian lattice over $\OK$ of rank $\ell\ge1$. There are real numbers $x_1\ge x_2\ge\dots\ge x_\ell$, the slopes of $\bar L$, and for each real $\theta$ a saturated submodule $L_{\ge\theta}\subset L$, with the following properties.
\begin{enumerate}
\item[(HN1)] For every $\theta$ the number of $j$ with $x_j\ge\theta$ equals $\rk L_{\ge\theta}$.
\item[(HN2)] $x_1+\dots+x_\ell=\ell\,\mu(L)$.
\item[(HN3)] $\mu_{\min}(L)\le x_j\le\mu_{\max}(L)$ for every $j$.
\item[(HN4)] If $N\subset L$ is a nonzero submodule with $\mu_{\min}(N)\ge\theta$, then $N\subset L_{\ge\theta}$.
\item[(HN5)] If $L_{\ge\theta}\ne0$, then $\mu_{\min}(L_{\ge\theta})\ge\theta$.
\end{enumerate}
Moreover $\mu(L^\vee)=-\mu(L)$ and $\mu_{\min}(L^\vee)=-\mu_{\max}(L)$.
\end{proposition}

These statements are the Harder--Narasimhan theory of Stuhler~\cite{r3stuhler} and Grayson~\cite{r3grayson} over $\mathbb Z$, extended to rings of integers in Bost's Bourbaki talk~\cite{bost} and by Gaudron~\cite{r3gaudron}, together with the behaviour of the Arakelov degree under duality~\cite{r3gaudron,r3bostkunnemann}. If $0=L_0\subsetneq L_1\subsetneq\dots\subsetneq L_s=L$ is the Harder--Narasimhan filtration, with semistable quotients of slopes $\nu_1>\dots>\nu_s$, then $x_1,\dots,x_\ell$ is the list of the $\nu_i$, each repeated $\rk(L_i/L_{i-1})$ times, and $L_{\ge\theta}=L_i$ for the largest $i$ with $\nu_i\ge\theta$, or $0$ when $\nu_1<\theta$. Statement (HN2) is the additivity of $\deg_{\mathbb Z}$ in exact sequences with quotient norms. The slope \eqref{r3:mu} is an affine function of $\widehat\mu$ with positive coefficient, and every statement of Proposition~\ref{r3:HN} is invariant under such a change.

The slope inequality we need bounds the minimal slope of an image. We give its proof, because the normalizations \eqref{r3:mu} and \eqref{r3:ht} enter the error terms of Part~II.

\begin{lemma}\label{r3:covol}
Let $\bar A$ be a hermitian lattice over $\OK$ in $V$ of rank $a$, let $V'$ have hermitian norms at the places of $\Sigma_K$, and let $\psi\colon V\to V'$ be an injective $K$-linear map. Give $\psi(A)$ the norms of $V'$ restricted to $\psi(V)$. Then
\[
\mu(\psi(A))\ge\mu(A)-\hgt(\psi).
\]
\end{lemma}

\begin{proof}
The map $\psi_{\mathbb R}=\bigoplus_v\psi_v$ sends $\rho(A)$ onto $\rho(\psi(A))\subset\psi(V)_{\mathbb R}$. In orthonormal bases of $V_v$ and $\psi(V)_v$ the complex linear map $\psi_v$ has a matrix $M_v$ of size $a$, and as a real linear map of $\mathbb C^a=\mathbb R^{2a}$ it has determinant $\lvert\det M_v\rvert^2\le\lVert\psi_v\rVert_{\mathrm{op}}^{2a}$. Hence
\[
\covol\rho(\psi(A))=\prod_v\lvert\det M_v\rvert^2\,\covol\rho(A)\le\prod_v\lVert\psi_v\rVert_{\mathrm{op}}^{2a}\,\covol\rho(A),
\]
so $\deg_{\mathbb Z}\psi(A)\ge\deg_{\mathbb Z}A-a\sum_v\log\lVert\psi_v\rVert_{\mathrm{op}}^2$. Dividing by $n_Ka/2$ gives the lemma, because the two lattices have the same rank.
\end{proof}

\begin{proposition}[slope inequality]\label{r3:SI}
Let $\bar A$ be a hermitian lattice over $\OK$ in $V$, let $V'$ have hermitian norms at the places of $\Sigma_K$, and let $\varphi\colon V\to V'$ be a $K$-linear map with $\varphi(A)\ne0$. Give $\varphi(A)$ the norms of $V'$ restricted to $K\varphi(A)$. Then
\[
\mu_{\min}(\varphi(A))\ge\mu_{\min}(A)-\hgt(\varphi).
\]
\end{proposition}

\begin{proof}
Let $M\subset\varphi(A)$ be a saturated proper submodule and $Q=\varphi(A)/M$ with the quotient norms. Let $\pi\colon K\varphi(A)\to K\varphi(A)/KM$ be the projection and $\psi=\pi\circ\varphi$ on $V$. Then $A_0=A\cap\ker\psi$ is a saturated proper submodule of $A$, and $\psi$ induces an injective map $\bar\psi\colon V/KA_0\to K\varphi(A)/KM$ that sends $A/A_0$ onto $Q$. For a class $\xi\in(V/KA_0)_v$ and any representative $y$ of $\xi$,
\[
\lVert\bar\psi_v\xi\rVert_v=\lVert\pi_v\varphi_vy\rVert_v\le\lVert\varphi_vy\rVert_v\le\lVert\varphi_v\rVert_{\mathrm{op}}\lVert y\rVert_v ,
\]
and the minimum over $y$ gives $\lVert\bar\psi_v\rVert_{\mathrm{op}}\le\lVert\varphi_v\rVert_{\mathrm{op}}$, so $\hgt(\bar\psi)\le\hgt(\varphi)$. Lemma~\ref{r3:covol} applied to $A/A_0$ and $\bar\psi$ gives $\mu(Q)\ge\mu(A/A_0)-\hgt(\varphi)$. A quotient of $A$ by a saturated submodule followed by a quotient by a saturated submodule is a quotient of $A$ by a saturated submodule, with the same quotient norms, so $\mu(A/A_0)\ge\mu_{\min}(A)$. Taking the infimum over $M$ proves the proposition.
\end{proof}

\subsection{The frame theorem}\label{r3:sec:frame}

For a complex matrix $X$ write $X^*$ for its conjugate transpose and $\lVert X\rVert_\HS^2=\tr(X^*X)$ for its squared Hilbert--Schmidt norm. For a $K$-linear map $\varphi$ between spaces with fixed bases and a place $v\in\Sigma_K$, write $[\varphi]_v$ for the matrix of $\varphi$ with $\sigma_v$ applied to its entries.

\begin{lemma}[frame lemma]\label{r3:LA}
Let $I$ be a finite set, let $\widetilde D\ge0$ and $d_i\ge0$ for $i\in I$ be integers with $\Delta=\sum_id_i>0$, let $A_i$ be complex $\widetilde D\times d_i$ matrices and $\gamma_i>0$, and put $Q=\sum_i\gamma_iA_iA_i^*$. Let $e,\eta,\varepsilon$ be real numbers with $0\le e\le\frac14$ and $\eta\ge0$ such that
\[
\sum_i\lVert\gamma_iA_i^*A_i-1\rVert_\HS^2\le e\Delta,\qquad\tr(Q^2)\le(1+\eta)\frac{(\tr Q)^2}{\widetilde D}.
\]
Let $U_i\subset\mathbb C^{d_i}$ be subspaces with $\sum_i(d_i-\dim U_i)\le\varepsilon \Delta$, and let $\mathcal W$ be the orthogonal complement of $\sum_iA_i(U_i)$ in $\mathbb C^{\widetilde D}$. Then $\dim\mathcal W\le(\eta+4\varepsilon+4\sqrt e)\,\widetilde D$.
\end{lemma}

\begin{proof}
Put $E_i=\gamma_iA_i^*A_i-1$, $k_i=d_i-\dim U_i$, $w=\dim\mathcal W$ and $t=\tr Q$. The hypothesis on $\sum k_i$ and $\Delta>0$ give $\varepsilon\ge0$. If $\widetilde D=0$ then $w=0$, so assume $\widetilde D\ge1$.

For orthonormal vectors $f_1,\dots,f_m$ in $\mathbb C^{d_i}$ we have $\langle f_l,E_if_l\rangle=\gamma_i\lVert A_if_l\rVert^2-1$, so the Cauchy--Schwarz inequality gives
\begin{equation}\label{r3:basic}
\Big\lvert\gamma_i\sum_{l=1}^m\lVert A_if_l\rVert^2-m\Big\rvert\le\sum_l\lVert E_if_l\rVert\le\sqrt m\Big(\sum_l\lVert E_if_l\rVert^2\Big)^{1/2}\le\sqrt m\,\lVert E_i\rVert_\HS .
\end{equation}
In the same way $\lvert\sum_l\langle g_l,Tg_l\rangle\rvert\le\sqrt m\,\lVert T\rVert_\HS$ for every matrix $T$ of size $\widetilde D$ and orthonormal $g_1,\dots,g_m$ in $\mathbb C^{\widetilde D}$.

By cyclicity of the trace $t=\sum_i\gamma_i\lVert A_i\rVert_\HS^2$. Inequality \eqref{r3:basic} for an orthonormal basis of $\mathbb C^{d_i}$, summed over $i$, and the Cauchy--Schwarz inequality give
\[
t\ge\sum_i\big(d_i-\sqrt{d_i}\,\lVert E_i\rVert_\HS\big)\ge \Delta-\sqrt \Delta\Big(\sum_i\lVert E_i\rVert_\HS^2\Big)^{1/2}\ge \Delta\,(1-\sqrt e).
\]
Since $\sqrt e\le\frac12$, we get $\Delta\le2t$, and in particular $t>0$.

Put $\lambda=t/\widetilde D$. Since $Q$ is self-adjoint,
\[
\lVert Q-\lambda\rVert_\HS^2=\tr(Q^2)-2\lambda t+\lambda^2\widetilde D=\tr(Q^2)-\frac{t^2}{\widetilde D}\le\frac{\eta\,t^2}{\widetilde D}.
\]
Let $g_1,\dots,g_w$ be an orthonormal basis of $\mathcal W$ and $T_{\mathcal W}=\sum_l\langle g_l,Qg_l\rangle$. The estimate after \eqref{r3:basic} with $T=Q-\lambda$ gives
\[
T_{\mathcal W}\ge\lambda w-\sqrt w\,\lVert Q-\lambda\rVert_\HS\ge\frac{t\,w}{\widetilde D}-\sqrt{\frac{w\,\eta}{\widetilde D}}\;t .
\]
For $g\in\mathcal W$ and $y\in U_i$ we have $\langle A_i^*g,y\rangle=\langle g,A_iy\rangle=0$, so $A_i^*g$ lies in $U_i^\perp$. Let $f_1,\dots,f_{k_i}$ be an orthonormal basis of $U_i^\perp$. Bessel's inequality for the orthonormal vectors $g_l$ gives
\[
\sum_l\lVert A_i^*g_l\rVert^2=\sum_{m}\sum_l\lvert\langle g_l,A_if_m\rangle\rvert^2\le\sum_m\lVert A_if_m\rVert^2,
\]
and \eqref{r3:basic} gives $\gamma_i\sum_m\lVert A_if_m\rVert^2\le k_i+\sqrt{k_i}\,\lVert E_i\rVert_\HS$. Since $\langle g,Qg\rangle=\sum_i\gamma_i\lVert A_i^*g\rVert^2$, summing over $i$ and applying the Cauchy--Schwarz inequality,
\[
T_{\mathcal W}\le\sum_ik_i+\Big(\sum_ik_i\Big)^{1/2}\Big(\sum_i\lVert E_i\rVert_\HS^2\Big)^{1/2}\le\big(\varepsilon+\sqrt{\varepsilon e}\big)\Delta\le2t\big(\varepsilon+\sqrt{\varepsilon e}\big).
\]
Comparing the two bounds for $T_{\mathcal W}$ and multiplying by $\widetilde D/t$ gives $w-\sqrt{w\eta\widetilde D}\le2\widetilde D(\varepsilon+\sqrt{\varepsilon e})$. The inequality $\sqrt{w\eta\widetilde D}\le\frac12(w+\eta\widetilde D)$ between geometric and arithmetic means turns this into $w\le\big(\eta+4\varepsilon+4\sqrt{\varepsilon e}\big)\widetilde D$. If $\varepsilon\le1$ then $\sqrt{\varepsilon e}\le\sqrt e$, which proves the lemma. If $\varepsilon>1$ the bound of the lemma exceeds $\widetilde D\ge w$.
\end{proof}

\begin{lemma}[rank at one place]\label{r3:place}
Let $I$ be a finite set. For $i\in I$ let $V_i$ be a $K$-vector space of dimension $d_i$ with a basis, let $W$ be a $K$-vector space of dimension $D$ with a basis, and let $B_i\colon V_i\to W$ be $K$-linear maps. Fix $v\in\Sigma_K$, invertible complex matrices $Y_i$ of size $d_i$ and $Y'$ of size $D$, an integer $D_1\ge0$, and put $\widetilde D=D+D_1$. Let $A_i$ be complex $\widetilde D\times d_i$ matrices whose first $D$ rows form the matrix $Y'^{-1}[B_i]_vY_i$, and let $\gamma_i>0$, $e$, $\eta$ and $\delta$ be real numbers with $0\le e\le\frac14$, $\eta\ge0$, $0\le\delta<1$ such that $\Delta=\sum_id_i>0$, $(1-\delta)\widetilde D\le D$, and, with $Q=\sum_i\gamma_iA_iA_i^*$,
\[
\sum_i\lVert\gamma_iA_i^*A_i-1\rVert_\HS^2\le e\Delta,\qquad\tr(Q^2)\le(1+\eta)\frac{(\tr Q)^2}{\widetilde D}.
\] Let $\varepsilon$ be real and let $N_i'\subset V_i$ be finitely generated $\OK$-submodules with $\sum_i(d_i-\rk N_i')\le\varepsilon \Delta$, and put $M=\sum_iB_i(N_i')\subset W$. Then
\[
D-\rk M\le\frac{\eta+4\varepsilon+4\sqrt e}{1-\delta}\,D .
\]
\end{lemma}

\begin{proof}
For $x\in V_i$ let $c_i(x)=Y_i^{-1}[x]_v\in\mathbb C^{d_i}$, and for $y\in W$ let $c'(y)=Y'^{-1}[y]_v\in\mathbb C^D$, where $[x]_v$ is the coordinate vector of $x$ with $\sigma_v$ applied. Then $c'(B_ix)=Y'^{-1}[B_i]_vY_i\,c_i(x)$. Let $U_i$ be the complex span of $c_i(N_i')$ and $T$ the complex span of $c'(M)$. A $K$-subspace $U$ of a $K$-vector space and its image under $\sigma_v$ satisfy $\dim_{\mathbb C}(U\otimes_{K,\sigma_v}\mathbb C)=\dim_KU$, so $\dim U_i=\rk N_i'$ and $\dim T=\rk M$. Let $P\colon\mathbb C^{\widetilde D}\to\mathbb C^D$ be the projection to the first $D$ coordinates. Then $PA_i(U_i)$ is spanned by the vectors $c'(B_ix)$ with $x\in N_i'$, which lie in $T$. Put $S=\sum_iA_i(U_i)$. Then $S\subset P^{-1}(T)$, which has dimension $\rk M+D_1$, so $\dim S^\perp\ge\widetilde D-\rk M-D_1=D-\rk M$. Lemma~\ref{r3:LA} bounds $\dim S^\perp$ by $(\eta+4\varepsilon+4\sqrt e)\widetilde D$, and $\widetilde D\le(1-\delta)^{-1}D$.
\end{proof}

\begin{lemma}[reverse Markov inequality]\label{r3:markov}
Let $y_1,\dots,y_m$, $s$, $u$, $\tau$ be real numbers with $u\ge0$, $\tau\ge0$, $\sum_jy_j\ge ms$ and $y_j\le s+u$ for every $j$. Let $F$ be the number of $j$ with $y_j<s-\tau$. Then $F(u+\tau)\le mu$, and $F=0$ when $u=0$.
\end{lemma}

\begin{proof}
Bounding the $F$ small terms by $s-\tau$ and the others by $s+u$ gives $ms\le\sum_jy_j\le F(s-\tau)+(m-F)(s+u)=m(s+u)-F(u+\tau)$. If $u=0$ and $F\ge1$, then every $y_j\le s$ and one $y_j<s$, so $\sum_jy_j<ms$.
\end{proof}

\begin{lemma}[layer cake]\label{r3:cake}
Let $z_1,\dots,z_m$ be real numbers with $m\ge1$, and let $s$ and $H\ge0$ be real numbers with $z_k\ge s-H$ for every $k$. Let $w$ be an integrable function on $[0,H]$ such that the number of $k$ with $z_k<s-\tau$ is at most $w(\tau)\,m$ for every $\tau\in[0,H]$. Then
\[
\frac1m\sum_kz_k\ge s-\int_0^H\min\big(1,w(\tau)\big)\,d\tau .
\]
\end{lemma}

\begin{proof}
If $s-z_k\ge0$ then $s-z_k\le H$ and $s-z_k=\int_0^H\mathbf 1[z_k<s-\tau]\,d\tau$. If $s-z_k<0$ the integral is nonnegative. Summing $s-z_k\le\int_0^H\mathbf 1[z_k<s-\tau]\,d\tau$ over $k$ and bounding the number of $k$ with $z_k<s-\tau$ by $m\min(1,w(\tau))$ proves the lemma.
\end{proof}

For $u,H\ge0$ put
\[
\Lambda(u,H)=
\begin{cases}
\dfrac{uH}{u+H},&u+H>0,\\
0,&u=H=0.
\end{cases}
\]
Then $0\le\Lambda(u,H)\le\min\{u,H\}$.

\begin{theorem}[frame theorem]\label{r3:frame}
Let $K$ be a totally complex number field, let $I$ be a nonempty finite set, and for $i\in I$ let $\bar N_i$ be a hermitian lattice over $\OK$ of rank $d_i$ in a $K$-vector space $V_i$ with a basis. Let $\bar C$ be a hermitian lattice over $\OK$ of rank $D$ in a $K$-vector space $W$ with a basis, and put $\Delta=\sum_id_i$, with $\Delta>0$. Let $B_i\colon V_i\to W$ be $K$-linear maps with $B_i(N_i)\subset C$, and let $h_i$ be real numbers with $\hgt(B_i)\le h_i$ whenever $B_i\ne0$, where $\hgt$ is computed with the norms of $\bar N_i$ and $\bar C$. Fix $v\in\Sigma_K$, invertible complex matrices $Y_i$ of size $d_i$ and $Y'$ of size $D$, an integer $D_1\ge0$, and complex $(D+D_1)\times d_i$ matrices $A_i$ whose first $D$ rows form $Y'^{-1}[B_i]_vY_i$. Let $\gamma_i>0$ and $Q=\sum_i\gamma_iA_iA_i^*$. Let $s$, $u$, $H$, $e$, $\eta$, $\delta$ be real numbers with $u\ge0$, $H\ge0$, $0\le e\le\frac14$, $\eta\ge0$ and $0\le\delta<1$ such that
\begin{enumerate}
\item[(i)] $s\Delta\le\sum_id_i\big(\mu(N_i)-h_i\big)$,
\item[(ii)] $\mu_{\max}(N_i)-h_i\le s+u$ for every $i$,
\item[(iii)] $s-\mu_{\min}(C)\le H$,
\item[(iv)] $(1-\delta)(D+D_1)\le D$,
\item[(v)] $\sum_i\lVert\gamma_iA_i^*A_i-1\rVert_\HS^2\le e\,\Delta$,
\item[(vi)] $\tr(Q^2)\le(1+\eta)(\tr Q)^2/(D+D_1)$.
\end{enumerate}
Then
\[
\mu(C)\ge s-\frac{4\Lambda(u,H)+(\eta+4\sqrt e)\,H}{1-\delta}.
\]
\end{theorem}

\begin{proof}
\emph{The target is nonzero.} If $D=0$ then $D_1=0$ by (iv), every $A_i$ is zero, and (v) reads $\Delta\le e\Delta$, which contradicts $\Delta>0$ and $e\le\frac14$. So $D\ge1$.

\emph{The pieces of the sources.}
For $\tau\ge0$ and $i\in I$ let
$N_i^\tau=(N_i)_{\ge s-\tau+h_i}$ be the
Harder--Narasimhan step of Proposition~\ref{r3:HN}.
Let $x_{i1},\dots,x_{id_i}$ be the slopes of $N_i$
and put $y_{ij}=x_{ij}-h_i$.
By (HN2) and (i),
$\sum_{i,j}y_{ij}=\sum_id_i(\mu(N_i)-h_i)\ge\Delta s$.
By (HN3) and (ii), $y_{ij}\le s+u$.
Define the actual deficit proportion by
\[
\varepsilon(\tau)=\frac1\Delta
\#\{(i,j):i\in I,\ 1\le j\le d_i,\ y_{ij}<s-\tau\}.
\]
Then $0\le\varepsilon(\tau)\le1$, and
$\varepsilon$ is a finite step function.
By (HN1), the number of $j$ with
$x_{ij}<s-\tau+h_i$ is $d_i-\rk N_i^\tau$.
Consequently,
\[
\sum_i\big(d_i-\rk N_i^\tau\big)
=\varepsilon(\tau)\,\Delta.
\]

\emph{The images lie in one step of the target.} Put $C_\tau=C_{\ge s-\tau}$. Fix $i$ with $B_i(N_i^\tau)\ne0$. Then $N_i^\tau\ne0$, so (HN5) gives $\mu_{\min}(N_i^\tau)\ge s-\tau+h_i$, and $B_i\ne0$, so $\hgt(B_i)\le h_i$. Proposition~\ref{r3:SI} gives
\[
\mu_{\min}\big(B_i(N_i^\tau)\big)\ge\mu_{\min}(N_i^\tau)-\hgt(B_i)\ge s-\tau .
\]
Since $B_i(N_i^\tau)\subset C$, (HN4) gives $B_i(N_i^\tau)\subset C_\tau$. Hence $M_\tau=\sum_iB_i(N_i^\tau)$ lies in $C_\tau$, and $\rk C_\tau\ge\rk M_\tau$.

\emph{The step of the target has small corank.} Lemma~\ref{r3:place} with $N_i'=N_i^\tau$ and $\varepsilon=\varepsilon(\tau)$ gives
\[
D-\rk C_\tau\le D-\rk M_\tau\le w(\tau)\,D,\qquad w(\tau)=\frac{\eta+4\varepsilon(\tau)+4\sqrt e}{1-\delta}.
\]

\emph{Integration.}
Let $z_1,\dots,z_D$ be the slopes of $C$.
By (HN1), the number of $k$ with $z_k<s-\tau$ is
$D-\rk C_\tau\le w(\tau)D$.
By (HN3) and (iii),
$z_k\ge\mu_{\min}(C)\ge s-H$.
By (HN2), $\mu(C)$ is the mean of the $z_k$.
Lemma~\ref{r3:cake}, using $\min(1,w(\tau))\le w(\tau)$,
therefore gives
$\mu(C)\ge s-\int_0^H w(\tau)\,d\tau$.

Write $a_+=\max\{a,0\}$.
For $u,H>0$ and $y\le s+u$, one has
\[
\min\{H,(s-y)_+\}
\le\frac{H}{u+H}(s+u-y).
\]
Indeed, this follows by considering separately
$s-y\le0$, $0\le s-y\le H$, and $s-y\ge H$.
Integrating the finite sum defining $\varepsilon$ gives
\[
\begin{aligned}
\int_0^H\varepsilon(\tau)\,d\tau
&=\frac1\Delta\sum_{i,j}\min\{H,(s-y_{ij})_+\}\\
&\le\frac{H}{u+H}
\left(s+u-\frac1\Delta\sum_{i,j}y_{ij}\right)\\
&\le\frac{uH}{u+H}
=\Lambda(u,H).
\end{aligned}
\]
If $H=0$, the integral is zero.
If $u=0$, the mean bound and $y_{ij}\le s$ force
every $y_{ij}=s$, so $\varepsilon(\tau)=0$ for
$\tau\ge0$.
Thus the same integral bound holds for all $u,H\ge0$.
Consequently,
\[
\mu(C)\ge s-\int_0^H w(\tau)\,d\tau
\ge s-\frac{4\Lambda(u,H)+(\eta+4\sqrt e)H}{1-\delta}.
\qedhere
\]
\end{proof}

\begin{remark}\label{r3:rem:frame}
The proof bounds the corank of every step $C_{\ge s-\tau}$ with $0\le\tau\le H$, so it controls the lower half of the Harder--Narasimhan polygon of $C$, and the mean slope is one consequence. The hermitian structure at $v$ in which conditions (v) and (vi) are stated is arbitrary, because $Y_i$ and $Y'$ are arbitrary invertible matrices. The norms of the lattices at all places of $\Sigma_K$ enter through the slopes and the heights in Proposition~\ref{r3:SI}, and the place $v$ enters through the dimension count of Lemma~\ref{r3:place}. For a single map, $I=\{i\}$, the theorem is weaker than Proposition~\ref{r3:SI}, and the gain lies in the passage from $\mu_{\max}(C)$ to $\mu(C)$ for a redundant family. The term $\Lambda(u,H)$ bounds the integral of the actual deficit proportion $\varepsilon(\tau)$ of the source slopes over $0\le\tau\le H$.
\end{remark}

\subsection{Gram families and the frame conditions}\label{r3:sec:gram}

In the application the matrices $A_i$ come from maps $P_i$ defined on larger spaces, compressed by isometries, and the hypotheses are Hilbert--Schmidt bounds for the Gram operators of the $P_i$. This subsection derives conditions (v) and (vi) of Theorem~\ref{r3:frame} from such bounds.

\begin{lemma}\label{r3:compress}
Let $X$ be a complex $a\times b$ matrix, and let $V$ and $V'$ be complex $a\times c$ and $b\times c'$ matrices with $V^*V=1$ and $V'^*V'=1$. Then $\lVert V^*XV'\rVert_\HS\le\lVert X\rVert_\HS$. If moreover $a=b$, $V=V'$ and $\lVert X-1\rVert_\HS^2\le\epsilon a$ with $\epsilon\ge0$, then
\begin{gather*}
\tr X-\tr(V^*XV)=(a-c)+\tr\big((X-1)(1-VV^*)\big),\\
\big\lvert\tr\big((X-1)(1-VV^*)\big)\big\rvert\le\sqrt{(a-c)\,\epsilon\,a}.
\end{gather*}
\end{lemma}

\begin{proof}
The matrix $VV^*$ is an orthogonal projection, so $\lVert V^*Y\rVert_\HS^2=\tr(Y^*VV^*Y)\le\tr(Y^*Y)$ for every $Y$ with $a$ rows, and the same argument on the right gives the first statement. For the second, $\tr(V^*XV)=\tr(XVV^*)$, so $\tr X-\tr(V^*XV)=\tr(X(1-VV^*))=\tr(1-VV^*)+\tr((X-1)(1-VV^*))$, and $\tr(1-VV^*)=a-\tr(V^*V)=a-c$. The projection $1-VV^*$ has $\lVert1-VV^*\rVert_\HS^2=a-c$, and the Cauchy--Schwarz inequality for the Hilbert--Schmidt inner product gives the bound.
\end{proof}

\begin{proposition}[from Gram bounds to the frame conditions]\label{r3:gram}
Let $I$ be a nonempty finite set, let $\widetilde D\ge0$ be an integer, and for $i\in I$ let $d_i^S\ge d_i\ge0$ be integers, $P_i$ a complex $\widetilde D\times d^S_i$ matrix, $V_i$ a complex $d^S_i\times d_i$ matrix with $V_i^*V_i=1$, and $\gamma_i>0$. Put
\[
K_{ij}=\sqrt{\gamma_i\gamma_j}\,P_i^*P_j,\qquad \Delta^S=\sum_id^S_i,\qquad \Delta=\sum_id_i,\qquad A_i=P_iV_i,\qquad Q=\sum_i\gamma_iA_iA_i^* .
\]
Let $D^*>0$, $\sigma$, $\eta_0$, $\eta_1$, $e_0$ and $\delta'$ be real numbers with
\[
0<\sigma\le\tfrac1{20},\quad-1\le\eta_0\le\tfrac\sigma2,\quad\eta_1\le\tfrac\sigma2,\quad0\le e_0,\quad\sqrt{e_0}\le\tfrac\sigma2,\quad0\le\delta'\le\tfrac14,
\]
and suppose that
\begin{enumerate}
\item[(G1)] $\lVert K_{ij}\rVert_\HS^2\le(1+\eta_0)\,d^S_id^S_j/D^*$ for $i\ne j$,
\item[(G2)] $\lVert K_{ii}-1\rVert_\HS^2\le e_0\,d^S_i$ for every $i$,
\item[(G3)] $\widetilde D\le(1+\eta_1)D^*$,
\item[(G4)] $D^*\le\sigma \Delta^S$,
\item[(G5)] $\sum_i(d^S_i-d_i)\le\delta'\Delta^S$.
\end{enumerate}
Then $\Delta\ge(1-\delta')\Delta^S>0$, $\widetilde D\ge1$, and
\[
\sum_i\lVert\gamma_iA_i^*A_i-1\rVert_\HS^2\le2e_0\Delta,\qquad\tr(Q^2)\le\big(1+5\sigma+3\sqrt{\delta'}\big)\frac{(\tr Q)^2}{\widetilde D}.
\]
\end{proposition}

\begin{proof}
Put $a=\sqrt{e_0}$ and $t=\sqrt{\delta'}$, so that $0\le a\le\frac1{40}$ and $0\le t\le\frac12$. By (G4), $\Delta^S\ge D^*/\sigma>0$, and (G5) gives $\Delta\ge(1-\delta')\Delta^S\ge\frac34\Delta^S$.

\emph{Condition (v).} We have $\gamma_iA_i^*A_i=V_i^*K_{ii}V_i$ and $V_i^*K_{ii}V_i-1=V_i^*(K_{ii}-1)V_i$. Lemma~\ref{r3:compress} and (G2) give $\sum_i\lVert\gamma_iA_i^*A_i-1\rVert_\HS^2\le e_0\Delta^S\le2e_0\Delta$.

\emph{Traces of the uncompressed family.} Put $Q^S=\sum_i\gamma_iP_iP_i^*$. Then $\tr Q^S=\sum_i\tr K_{ii}$ and
\[
\tr\big((Q^S)^2\big)=\sum_{i,j}\gamma_i\gamma_j\tr\big(P_iP_i^*P_jP_j^*\big)=\sum_{i,j}\lVert K_{ij}\rVert_\HS^2 .
\]
The estimate \eqref{r3:basic} for an orthonormal basis gives $\tr K_{ii}\ge d^S_i-\sqrt{d^S_i}\lVert K_{ii}-1\rVert_\HS\ge(1-a)d^S_i$, so $\tr Q^S\ge(1-a)\Delta^S$. The triangle inequality gives $\lVert K_{ii}\rVert_\HS\le\sqrt{d^S_i}+\sqrt{e_0d^S_i}=(1+a)\sqrt{d^S_i}$. With (G1) and $1+\eta_0\ge0$,
\[
\tr\big((Q^S)^2\big)\le(1+\eta_0)\frac{(\Delta^S)^2}{D^*}+(1+a)^2\Delta^S\le\frac{(\Delta^S)^2}{D^*}\Big((1+\eta_0)+(1+a)^2\sigma\Big),
\]
where the last step uses $\Delta^S\le\sigma(\Delta^S)^2/D^*$, which is (G4).

\emph{Compression.} The Gram operators of the compressed family are $\sqrt{\gamma_i\gamma_j}A_i^*A_j=V_i^*K_{ij}V_j$, so Lemma~\ref{r3:compress} gives $\tr(Q^2)=\sum_{i,j}\lVert V_i^*K_{ij}V_j\rVert_\HS^2\le\tr((Q^S)^2)$. The second part of Lemma~\ref{r3:compress} with $X=K_{ii}$ and $\epsilon=e_0$, the Cauchy--Schwarz inequality and (G5) give
\begin{align*}
\tr Q^S-\tr Q&\le\sum_i\Big((d^S_i-d_i)+\sqrt{d^S_i-d_i}\,a\sqrt{d^S_i}\Big)\\
&\le\delta'\Delta^S+a\Big(\sum_i(d^S_i-d_i)\Big)^{1/2}(\Delta^S)^{1/2}\le(t^2+at)\Delta^S .
\end{align*}
So $\tr Q\ge(1-a-t^2-ta)\Delta^S$, and the coefficient is positive because $a\le\frac1{40}$ and $t\le\frac12$. In particular $\tr Q>0$, so $\widetilde D\ge1$.

\emph{Condition (vi).} Since $\widetilde D\ge1$, (G3) gives $1+\eta_1>0$. By (G3) and the two bounds above,
\[
\widetilde D\,\tr(Q^2)\le(1+\eta_1)\big((1+\eta_0)+(1+a)^2\sigma\big)(\Delta^S)^2 .
\]
The right side is nondecreasing in $\eta_0$, $\eta_1$ and $a$ on the given ranges, so it is at most $(1+\frac\sigma2)^2\big(1+(1+\frac\sigma2)\sigma\big)(\Delta^S)^2$. For $0<\sigma\le\frac14$,
\[
(1+5\sigma)\big(1-\tfrac\sigma2\big)^2-\big(1+\tfrac\sigma2\big)^2\big(1+\sigma+\tfrac{\sigma^2}2\big)=\sigma\Big(2-\tfrac{13}2\sigma+\tfrac12\sigma^2-\tfrac18\sigma^3\Big)\ge\sigma\Big(2-\tfrac{13}8-\tfrac1{512}\Big)>0,
\]
and $(1-\frac\sigma2)^2\le(1-a)^2$. Hence $\widetilde D\tr(Q^2)\le(1+5\sigma)(1-a)^2(\Delta^S)^2$. Put $P=1+5\sigma$, $B=1-a$ and $c=t(t+a)$. The function $3B-2(P+3t)(t+a)$ decreases in $\sigma$, $a$ and $t$, and at $\sigma=\frac1{20}$, $a=\frac1{40}$, $t=\frac12$ it equals $\frac3{80}$. So $2(P+3t)(t+a)\le3B$, and
\[
(P+3t)(B-c)^2\ge(P+3t)(B^2-2Bc)=PB^2+tB\big(3B-2(P+3t)(t+a)\big)\ge PB^2 .
\]
Since $B-c=1-a-t^2-ta>0$ and $(B-c)\Delta^S\le\tr Q$,
\[
\widetilde D\,\tr(Q^2)\le PB^2(\Delta^S)^2\le(P+3t)(B-c)^2(\Delta^S)^2\le\big(1+5\sigma+3\sqrt{\delta'}\big)(\tr Q)^2 .\qedhere
\]
\end{proof}

The ranks $d_i$ in the application are the ranks of the parts of the spaces $C_u$ of Section~\ref{r3:sec:A3} that are ramified at the auxiliary prime $T_u$, and the following lemma gives (G5).

\begin{lemma}\label{r3:ramshare}
Let $f\ge1$ be an integer and let $T\ge2$ and $d^S\ge0$ be real numbers. If $d=d^S\big(1-\frac1{T-1}\big)^f$, then $d^S-d\le\frac{2f}T\,d^S$.
\end{lemma}

\begin{proof}
Bernoulli's inequality $(1-y)^f\ge1-fy$ for $0\le y\le1$, with $y=1/(T-1)$, gives $d^S-d\le\frac f{T-1}d^S\le\frac{2f}Td^S$.
\end{proof}

\subsection{The period maps and hypothesis~\eqref{A3}}\label{r3:sec:A3}

We now state the setting of hypothesis~\eqref{A3} of Proposition~\ref{acct}. For each $u_0\in A_0$ the construction of Part~II gives the following objects. The field $K$ is a totally complex Galois number field that contains the field $E$. The index set $W$ is a finite nonempty set of words $u$, and each word has an auxiliary prime $T_u$. There is one hermitian lattice $R$ over $\OK$ of rank $D_{\mathrm{rows}}\ge1$ in a $K$-vector space $Y$, the lattice of automorphic forms, and $\mu^C=\mu(R^\vee)$. For each $u$ there are a hermitian lattice $N_u$ of rank $d_u\ge1$ in a $K$-vector space $V_u$, the period map $B_u\colon V_u\to Y^*$, and a real number $h_u$. The numbers $\mu^N_u=\mu(N_u)$ and the weights $d_u$ are those of \eqref{A3}.

The archimedean structure of the period maps at one complex place is the following.

\begin{definition}\label{r3:bessel}
Let $v\in\Sigma_K$. A \emph{Bessel realization} of the family $(B_u)_{u\in W}$ at $v$ consists of
\begin{enumerate}
\item a complex Hilbert space $\mathcal H$ of finite dimension $D_{\mathrm{all}}$ and a conjugate-linear isometric embedding $\mathbf s\colon Y_v\to\mathcal H$, where $Y_v=Y\otimes_{K,\sigma_v}\mathbb C$ has the norm of $R$ at $v$;
\item for each $u\in W$ a complex Hilbert space $\mathcal C_u$ of dimension $d^S_u$, a linear map $\mathcal P_u\colon\mathcal C_u\to\mathcal H$, a number $c_u>0$, an inner product on $V_{u,v}=V_u\otimes_{K,\sigma_v}\mathbb C$ and a linear isometric embedding $\iota_u\colon V_{u,v}\to\mathcal C_u$;
\item a unitary operator $\mathcal V$ of $Y_v$, and for each $u$ a unitary operator $\mathcal U_u$ of $V_{u,v}$ and a complex number $a_u\ne0$,
\end{enumerate}
such that for every $u$, every $\varphi\in V_{u,v}$ and every $F\in Y_v$
\begin{equation}\label{r3:besselid}
\big(B_{u,v}\varphi\big)(F)=a_u\,\big\langle\mathbf s(\mathcal VF),\mathcal P_u\,\iota_u\,\mathcal U_u\varphi\big\rangle_{\mathcal H}\,.
\end{equation}
Here $B_{u,v}\varphi\in Y^*\otimes_{K,\sigma_v}\mathbb C$ is read as a linear form on $Y_v$, and the inner product of $\mathcal H$ is conjugate-linear in the first variable, so both sides of \eqref{r3:besselid} are linear in $\varphi$ and in $F$. The \emph{Gram operators} of the realization are $K_{uu'}=(c_uc_{u'})^{-1/2}\mathcal P_u^*\mathcal P_{u'}\colon\mathcal C_{u'}\to\mathcal C_u$.
\end{definition}

In the construction $\mathcal H$ is a space of automorphic forms on a unitary group, $\mathcal C_u$ is the space of column forms of the word $u$, $\iota_u$ is the inclusion of the forms that are ramified at $T_u$, $\mathcal P_u$ is the map to Bessel Poincar\'e series and $c_u$ is its normalizing volume, and \eqref{r3:besselid} identifies the period map with a Bessel period. Section~\ref{r3:sec:S} states what the construction has to prove about these objects.

\begin{lemma}[matrices of the period maps at $v$]\label{r3:block}
Let $(B_u)$ have a Bessel realization at $v$ with $D_{\mathrm{all}}\ge D_{\mathrm{rows}}$, and fix $K$-bases of $Y^*$ and of each $V_u$. Put $D_1=D_{\mathrm{all}}-D_{\mathrm{rows}}$. There are invertible complex matrices $Y'$ of size $D_{\mathrm{rows}}$ and $Y_u$ of size $d_u$, complex $D_{\mathrm{all}}\times d^S_u$ matrices $P_u$ and complex $d^S_u\times d_u$ matrices $V_u$ with $V_u^*V_u=1$, such that for every $u$ the first $D_{\mathrm{rows}}$ rows of $P_uV_u$ form the matrix $Y'^{-1}[B_u]_vY_u$, and such that for all $u,u'$ the matrix $(c_uc_{u'})^{-1/2}P_u^*P_{u'}$ is the matrix of $K_{uu'}$ in orthonormal bases of $\mathcal C_{u'}$ and $\mathcal C_u$. In particular $\lVert(c_uc_{u'})^{-1/2}P_u^*P_{u'}\rVert_\HS=\lVert K_{uu'}\rVert_\HS$ and $\lVert c_u^{-1}P_u^*P_u-1\rVert_\HS=\lVert K_{uu}-1\rVert_\HS$.
\end{lemma}

\begin{proof}
Let $(\epsilon_k)_{k\le D_{\mathrm{rows}}}$ be an orthonormal basis of $Y_v$. The vectors $\mathbf s(\mathcal V\epsilon_k)$ are orthonormal in $\mathcal H$, because $\mathbf s$ is a conjugate-linear isometry and $\mathcal V$ is unitary. Extend them to an orthonormal basis $(f_a)_{a\le D_{\mathrm{all}}}$ of $\mathcal H$ with $f_k=\mathbf s(\mathcal V\epsilon_k)$ for $k\le D_{\mathrm{rows}}$. Fix orthonormal bases $(c_{u,l})_l$ of $\mathcal C_u$ and $(o_{u,j})_{j\le d_u}$ of $V_{u,v}$, and put
\[
(P_u)_{a,l}=\langle f_a,\mathcal P_uc_{u,l}\rangle,\qquad(V_u)_{l,j}=\langle c_{u,l},\iota_u\mathcal U_uo_{u,j}\rangle .
\]
Since $\iota_u\mathcal U_u$ is an isometry, $V_u^*V_u=1$, and $(P_uV_u)_{a,j}=\langle f_a,\mathcal P_u\iota_u\mathcal U_uo_{u,j}\rangle$. The matrix $P_u^*P_{u'}$ has entries $\langle\mathcal P_uc_{u,l},\mathcal P_{u'}c_{u',l'}\rangle$ by Parseval's identity in $\mathcal H$, which is the matrix of $\mathcal P_u^*\mathcal P_{u'}$ in the chosen bases.

For $k\le D_{\mathrm{rows}}$, \eqref{r3:besselid} gives
\[
(P_uV_u)_{k,j}=\big\langle\mathbf s(\mathcal V\epsilon_k),\mathcal P_u\iota_u\mathcal U_uo_{u,j}\big\rangle=a_u^{-1}\big(B_{u,v}o_{u,j}\big)(\epsilon_k).
\]
Let $(b_l)$ be the fixed basis of $V_u$ and $(b'_m)$ the fixed basis of $Y^*$. Let $\beta_u$ be the matrix with $o_{u,j}=\sum_l(\beta_u)_{lj}\,b_l\otimes1$, and let $\beta'$ be the matrix with entries $\beta'_{km}=(b'_m\otimes1)(\epsilon_k)$. Since $B_{u,v}(b_l\otimes1)=\sum_m([B_u]_v)_{ml}\,b'_m\otimes1$, we get $\big((B_{u,v}o_{u,j})(\epsilon_k)\big)_{k,j}=\beta'[B_u]_v\beta_u$. The matrix $\beta_u$ is invertible because $(o_{u,j})$ is a basis of $V_{u,v}$, and $\beta'$ is invertible because it is the matrix of the perfect pairing between $Y^*_v$ and $Y_v$ in the bases $(b'_m\otimes1)$ and $(\epsilon_k)$. So $Y'=\beta'^{-1}$ and $Y_u=a_u^{-1}\beta_u$ have the required property.
\end{proof}

The next theorem is hypothesis~\eqref{A3} of Proposition~\ref{acct}, with $\mu^N_u=\mu(N_u)$, $\mu^C=\mu(R^\vee)$, the weights $d_u=\rk N_u$ and the error $e_3$ of \eqref{r3:e3}. Its hypotheses (a) to (f) are the outputs of the construction that it uses.

\begin{theorem}[hypothesis~\eqref{A3}]\label{r3:A3}
Let $K$ be a totally complex number field, $W$ a finite nonempty set, $\bar R$ a hermitian lattice over $\OK$ of rank $D_{\mathrm{rows}}\ge1$ in a $K$-vector space $Y$, and for $u\in W$ let $\bar N_u$ be a hermitian lattice over $\OK$ of rank $d_u\ge1$ in a $K$-vector space $V_u$, $B_u\colon V_u\to Y^*$ a $K$-linear map and $h_u$ a real number. Put
\[
s=\frac{\sum_ud_u\big(\mu(N_u)-h_u\big)}{\sum_ud_u},\qquad u_F=\max_{u\in W}\big(\mu_{\max}(N_u)-h_u\big)-s .
\]
Let $D^*>0$, $H>0$, $\sigma$, $\eta_0$, $\eta_1$, $e_0$, $\delta'$, $\delta$ be real numbers with
\[
0<\sigma\le\tfrac1{20},\quad-1\le\eta_0\le\tfrac\sigma2,\quad\eta_1\le\tfrac\sigma2,\quad0\le e_0\le\tfrac18,\quad\sqrt{e_0}\le\tfrac\sigma2,\quad0\le\delta'\le\tfrac14,\quad0\le\delta<1,
\]
and suppose the following.
\begin{enumerate}
\item[(a)] $B_u(N_u)\subset R^\vee$ for every $u$.
\item[(b)] $\hgt(B_u)\le h_u$ for every $u$ with $B_u\ne0$, the height being computed with the norms of $\bar N_u$ and $\bar R^\vee$.
\item[(c)] $s+\mu_{\max}(R)\le H$.
\item[(d)] The family $(B_u)$ has a Bessel realization at a place $v\in\Sigma_K$ with $D_{\mathrm{all}}\ge D_{\mathrm{rows}}$ and $(1-\delta)D_{\mathrm{all}}\le D_{\mathrm{rows}}$.
\item[(e)] The Gram operators of the realization satisfy $\lVert K_{uu'}\rVert_\HS^2\le(1+\eta_0)d^S_ud^S_{u'}/D^*$ for $u\ne u'$, $\lVert K_{uu}-1\rVert_\HS^2\le e_0d^S_u$ for every $u$, and $D_{\mathrm{all}}\le(1+\eta_1)D^*$.
\item[(f)] $D^*\le\sigma\sum_ud^S_u$ and $\sum_u(d^S_u-d_u)\le\delta'\sum_ud^S_u$.
\end{enumerate}
Then the number
\begin{equation}\label{r3:e3}
e_3=\frac{4\Lambda(u_F,H)+\big(5\sigma+3\sqrt{\delta'}+4\sqrt{2e_0}\big)H}{1-\delta}
\end{equation}
is nonnegative, and
\[
\frac{\sum_ud_u\big(\mu(N_u)-h_u\big)}{\sum_ud_u}-e_3\le\mu(R^\vee).
\]
\end{theorem}

\begin{proof}
Lemma~\ref{r3:block} gives invertible matrices $Y'$, $Y_u$ and matrices $P_u$, $V_u$ with $V_u^*V_u=1$ such that the first $D_{\mathrm{rows}}$ rows of $A_u=P_uV_u$ form $Y'^{-1}[B_u]_vY_u$, and such that the matrices $\sqrt{\gamma_u\gamma_{u'}}P_u^*P_{u'}$ with $\gamma_u=c_u^{-1}$ have the Hilbert--Schmidt norms of the Gram operators. Proposition~\ref{r3:gram} with $\widetilde D=D_{\mathrm{all}}$ applies by (e) and (f), and it gives conditions (v) and (vi) of Theorem~\ref{r3:frame} with $e=2e_0\le\frac14$ and $\eta=5\sigma+3\sqrt{\delta'}$.

We apply Theorem~\ref{r3:frame} with $I=W$, the lattices $N_u$, the target $C=R^\vee$ of rank $D=D_{\mathrm{rows}}$, $D_1=D_{\mathrm{all}}-D_{\mathrm{rows}}$, and the numbers $s$, $u=u_F$ and $H$. The inclusions $B_u(N_u)\subset C$ and the height bounds are (a) and (b). Condition (i) holds with equality by the definition of $s$. Condition (ii) holds by the definition of $u_F$, and $u_F\ge0$ because $\mu(N_u)\le\mu_{\max}(N_u)$ makes $s$ at most the maximum of $\mu_{\max}(N_u)-h_u$. Condition (iii) is (c), because $\mu_{\min}(R^\vee)=-\mu_{\max}(R)$ by Proposition~\ref{r3:HN}. Condition (iv) is (d). The conclusion of Theorem~\ref{r3:frame} is the inequality of the theorem, with $\eta+4\sqrt e=5\sigma+3\sqrt{\delta'}+4\sqrt{2e_0}$.
\end{proof}

\begin{remark}\label{r3:rem:log}
Theorem~\ref{r3:A3} also gives
\[
s-\log\lvert W\rvert-e_3\le\mu(R^\vee)
\]
by subtracting $\log\lvert W\rvert\ge0$.
Condition (K3) of Section~\ref{sec:proofA} uses
a further weakening, since
\[
\Lambda(u,H)\le u\log(1+H/u)
\qquad(u>0,\ H\ge0).
\]
This follows from
$\log(1+t)\ge t/(1+t)$ for $t\ge0$.
At $u=0$, both expressions are interpreted as zero.
Thus the logarithmic accounting bounds in
Section~\ref{sec:proofA} remain valid.
\end{remark}

The size of $e_3$ is controlled by the spread $u_F$, which the next lemma bounds by the two-sided form of hypothesis \eqref{A2} and the spreads of the base slopes, the local terms and the heights.

\begin{lemma}\label{r3:uF}
Let $K$ be a totally complex number field and $W$ a finite nonempty set, and for $u\in W$ let $\bar N_u$ and $\bar M_u$ be hermitian lattices over $\OK$, with $d_u=\rk N_u\ge1$, and let $h_u$, $E_u$, $\beta_u^+$, $\beta_u^-$ be real numbers. Let $\Gamma$ and $\mathcal L$ be real numbers, put $s=\sum_ud_u(\mu(N_u)-h_u)/\sum_ud_u$ and $u_F=\max_u(\mu_{\max}(N_u)-h_u)-s$, and suppose that for every $u$
\[
\mu_{\max}(N_u)\le\mu_{\max}(M_u)+\Gamma+\mathcal L-E_u+\beta^+_u,\qquad\mu(N_u)\ge\mu(M_u)+\Gamma+\mathcal L-E_u-\beta^-_u .
\]
Then
\[
u_F\le\max_{u,u'\in W}\Big(\big(\mu_{\max}(M_u)-\mu(M_{u'})\big)+\lvert E_u-E_{u'}\rvert+\lvert h_u-h_{u'}\rvert\Big)+\max_u\beta^+_u+\frac{\sum_ud_u\beta^-_u}{\sum_ud_u}.
\]
\end{lemma}

\begin{proof}
For $u,u'\in W$ the two inequalities give
\[
\big(\mu_{\max}(N_u)-h_u\big)-\big(\mu(N_{u'})-h_{u'}\big)\le\big(\mu_{\max}(M_u)-\mu(M_{u'})\big)+(E_{u'}-E_u)+(h_{u'}-h_u)+\beta^+_u+\beta^-_{u'} .
\]
Fix $u$, multiply by $d_{u'}$, sum over $u'$ and divide by $\sum_{u'}d_{u'}$. The left side becomes $\mu_{\max}(N_u)-h_u-s$, and the right side is at most the bound of the lemma. Taking the maximum over $u$ proves the lemma.
\end{proof}

\begin{corollary}[size of $e_3$]\label{r3:size}
Let $x\ge1$, $C_H\ge0$, $b>0$, $u_F\ge0$,
$U\ge\max(u_F,b)$, $0\le\delta\le\frac12$,
$\delta'\ge0$ and $e_0\ge0$ be real numbers with
\[
(3\sqrt{\delta'}+4\sqrt{2e_0})\,x^{C_H}\le1.
\]
Put $\sigma=x^{-C_H}$, $H=x^{C_H}b$ and
\[
e_3=
\frac{4\Lambda(u_F,H)
+\big(5\sigma+3\sqrt{\delta'}+4\sqrt{2e_0}\big)H}
{1-\delta}.
\]
Then $e_3\le8\,U+12\,b$.
\end{corollary}

\begin{proof}
We have
\[
\begin{aligned}
\big(5\sigma+3\sqrt{\delta'}+4\sqrt{2e_0}\big)H
&=5b+\big(3\sqrt{\delta'}+4\sqrt{2e_0}\big)x^{C_H}b\\
&\le6b.
\end{aligned}
\]
Also $0\le\Lambda(u_F,H)\le u_F\le U$ and
$(1-\delta)^{-1}\le2$.
Consequently,
\[
e_3\le2(4U+6b)=8U+12b.
\qedhere
\]
\end{proof}

\begin{lemma}[the size conditions of Theorem~\ref{r3:A3}]\label{r3:sizecond}
Let $x\ge1$, $C_H>0$, $\varepsilon>0$, $C_0\ge1$, $\Gamma_0>0$, $b\ge3$ and an integer $f\ge1$ be given, put $\sigma=x^{-C_H}$ and $\omega=\exp(-x^{\varepsilon/8})$, and assume
\[
x^{C_H}\ge20,\qquad11\,\omega\,x^{C_H}\le1,\qquad4f\le b\,\omega^2 .
\]
Let $W$ be a finite set with $\lvert W\rvert\ge C_0\Gamma_0^2x^{C_H}$, let $D^*>0$, and for $u\in W$ let $d^S_u\ge D^*/(C_0\Gamma_0^2)$ and $T_u\ge b$ be real numbers and $d_u=d^S_u(1-\frac1{T_u-1})^f$. Let $\eta_0,\eta_1,e_0$ be real numbers with $\eta_0\ge-1$, $e_0\ge0$ and $\eta_0,\eta_1,\sqrt{e_0}\le\omega$, and put $\delta'=2f/b$. Then
\begin{gather*}
D^*\le\sigma\sum_ud^S_u,\qquad\sum_u(d^S_u-d_u)\le\delta'\sum_ud^S_u,\qquad0<\sigma\le\tfrac1{20},\qquad\eta_0,\eta_1,\sqrt{e_0}\le\tfrac\sigma2,\\
e_0\le\tfrac18,\qquad0\le\delta'\le\tfrac14,\qquad\big(3\sqrt{\delta'}+4\sqrt{2e_0}\big)x^{C_H}\le1 .
\end{gather*}
\end{lemma}

\begin{proof}
Summing $d^S_u\ge D^*/(C_0\Gamma_0^2)$ over $W$ gives $\sum_ud^S_u\ge\lvert W\rvert D^*/(C_0\Gamma_0^2)\ge x^{C_H}D^*=D^*/\sigma$. Lemma~\ref{r3:ramshare} and $T_u\ge b$ give $d^S_u-d_u\le(2f/b)d^S_u$, and summing over $u$ gives the second inequality. Now $\sqrt{\delta'}\le\sqrt{\omega^2/2}\le\omega$, $\sigma\le\frac1{20}$, and $\omega\le\frac{\sigma}{11}\le\frac\sigma2$, so $\eta_0,\eta_1,\sqrt{e_0}\le\frac\sigma2$ and $e_0\le\omega^2\le\frac18$, $\delta'\le\omega^2/2\le\frac14$. Finally $(3\sqrt{\delta'}+4\sqrt{2e_0})x^{C_H}\le(3+4\sqrt2)\,\omega\,x^{C_H}\le1$.
\end{proof}

The integrality (a) comes from the divisibility of the composite of the lifts and from the definition of $N_u$ as a rescaled sublattice. The following lemma isolates the step from local divisibility to (a).

\begin{lemma}[integrality from local divisibility]\label{r3:int}
Let $K$ be a number field, $R$ a lattice in a $K$-vector space $Y$, $V$ a $K$-vector space and $B\colon V\to Y^*$ a $K$-linear map. For each finite place $w$ of $K$ let $\Lambda_w$ be an $\mathcal O_{K,w}$-submodule of $V\otimes_KK_w$, and suppose that $\operatorname{ord}_w\big(B(x)(y)\big)\ge0$ for every $x\in\Lambda_w$ and $y\in R$, where $B$ is extended $K_w$-linearly. Let $N\subset V$ be a lattice with $N\subset\Lambda_w$ for every $w$. Then $B(N)\subset R^\vee$. In particular, if $N^0_w\subset V\otimes K_w$ and integers $s(w)\ge0$ satisfy $\operatorname{ord}_w(B(x)(y))\ge s(w)$ for $x\in N^0_w$ and $y\in R$, and $N\otimes\mathcal O_{K,w}\subset\varpi_w^{-s(w)}N^0_w$ for every $w$, then $B(N)\subset R^\vee$.
\end{lemma}

\begin{proof}
For $x\in N$ and $y\in R$ the number $B(x)(y)\in K$ has $\operatorname{ord}_w\ge0$ at every finite place $w$, so it lies in $\OK$. Hence $B(x)\in R^\vee$. For the second statement put $\Lambda_w=\varpi_w^{-s(w)}N^0_w$, so that $\operatorname{ord}_w(B(\varpi_w^{-s(w)}x)(y))\ge-s(w)+s(w)=0$ for $x\in N^0_w$.
\end{proof}

\subsection{The automorphic inputs of Theorem~\ref{r3:A3}}\label{r3:sec:S}

Theorem~\ref{r3:A3} uses three properties of the construction that are not consequences of slope theory. The first is the Bessel realization (d), that is, the identity \eqref{r3:besselid} of the period map with a Bessel period of a Poincar\'e series, together with the bound $(1-\delta)D_{\mathrm{all}}\le D_{\mathrm{rows}}$ with $0\le\delta\le\exp(-x^{\varepsilon_z/7})$, where $\varepsilon_z$ is $\varepsilon$ divided by a fixed integer $J\ge1$. The second is the near orthogonality (e). The third is the pair of counting conditions in (f) and Lemma~\ref{r3:sizecond}.

We record the objects of the second. Let $G$ be the unitary group of the hermitian space of the top step over $E$, let $\mathcal H$ be the space of automorphic forms on $G$ whose component at the archimedean places is the fixed discrete series and whose component at the finite places is fixed by the level idempotent of the construction, with the Petersson inner product, and let $D^*$ be the product of the formal degree of the discrete series, the volume of the automorphic quotient and the value at $1$ of the level idempotent. For each word $u$ the Bessel datum consists of a flag of isotropic vectors $e_1,\dots,e_r$ and a vector $z_u$ orthogonal to them with $\langle z_u,z_u\rangle=-T_u$. The seed of $u$ is built from the Bessel function $W_u$ of the discrete series for this datum and a character of a level subgroup, $\mathcal P_u\colon\mathcal C_u\to\mathcal H$ sends a column vector $T$ to the Poincar\'e series $x\mapsto\sum_{\gamma\in S_u(F)\backslash G(F)}\mathrm{Ph}_{u,T}(\gamma x)$ of the seed, where $S_u$ is the stabilizer of the datum, and $c_u=\mathrm{vol}_N(u)\mathrm{vol}_H(u)/\dim\Pi_u$ with $\dim\Pi_u$ the product over the archimedean places of the dimensions of the $K_\infty$-types of $u$. Put $\Gamma_0=\exp(C_\Gamma m^5fx^\varepsilon)$.

Section~\ref{sec:bessel} proves Theorem~\ref{w2:S} (near orthogonality), Theorem~\ref{w2:I} (the share of irreducible rows) and Theorem~\ref{w2:B} (the Bessel identity) from the statements listed there.

\section{Bessel Poincar\'e series}\label{sec:bessel}

This section proves the three automorphic inputs (d) and (e) of Theorem~\ref{r3:A3} from the statements listed in it: the near orthogonality of Bessel Poincar\'e series, the share of irreducible rows, and the Bessel identity for the period maps.



\subsection{The three statements}\label{w2:sec:statements}

We keep the notation of Section~\ref{r3:sec}. The field $K$ is a totally complex Galois number field, $v$ is a complex place of $K$, $Y$ is a $K$-vector space with a hermitian lattice $R$ of rank $D_{\mathrm{rows}}$, and for each word $u$ of the finite set $W$ the period map is a $K$-linear map $B_u\colon V_u\to Y^*$. Write $Y_v=Y\otimes_{K,\sigma_v}\C$ and $V_{u,v}=V_u\otimes_{K,\sigma_v}\C$, with the hermitian norms of $R$ and of $N_u$ at $v$.

On the automorphic side $F$ is a totally real field of degree $f$, $E/F$ is a CM extension, $V$ is a nondegenerate hermitian space over $E$, and $G=\mathrm U(V)$ is its unitary group over $F$. Write $\A$ and $\Af$ for the adeles and the finite adeles of $F$, $G_\infty=G(F\otimes_{\mathbb Q}\R)$ and $[G]=G(F)\backslash G(\A)$. The space $\mathcal H$ is the finite dimensional space of automorphic forms on $G$ whose archimedean component is a fixed discrete series and whose finite part is fixed by the level idempotent of the construction, with the Petersson inner product, which is conjugate linear in the first variable. Its dimension is $D_{\mathrm{all}}$. The number $D^*$ is the product of the formal degree of the discrete series, the volume of $[G]$ and the value at $1$ of the level idempotent.

For each word $u$ the column space $\mathcal C_u$ has dimension $d^S_u$, the Poincar\'e map $\mathcal P_u\colon\mathcal C_u\to\mathcal H$ sends a column form to its Bessel Poincar\'e series (Section~\ref{w2:sec:unfold}), $c_u>0$ is its normalizing volume, and
\[
K_{uu'}=(c_uc_{u'})^{-1/2}\mathcal P_u^*\mathcal P_{u'}\colon\mathcal C_{u'}\to\mathcal C_u .
\]
The family parameters are $x$, $\varepsilon$, $b=e^{x^\varepsilon}$, $L=\log b=x^\varepsilon$, the ranks $m,k,r$ of the datum and $\Gamma_0=\exp(C_\Gamma m^5fL)$.

\begin{theorem}[near orthogonality]\label{w2:S}
There are absolute constants $C_0\ge1$ and $C_\Gamma>0$ with the following property. For every $\varepsilon$, every $x$ large in terms of $\varepsilon$, and every family of data of Part~II at the scale $x$ whose words satisfy
\begin{equation}\label{w2:window}
C_0^{-1}\le d^S_u\Gamma_0^2/D^*\le C_0\qquad\text{for every }u\in W,
\end{equation}
the Gram operators satisfy
\[
\lVert K_{uu'}\rVert_\HS^2\le(1+\eta_0)\frac{d^S_ud^S_{u'}}{D^*}\ (u\ne u'),\qquad\lVert K_{uu}-1\rVert_\HS^2\le e_0\,d^S_u,\qquad D_{\mathrm{all}}\le(1+\eta_1)D^*,
\]
with $e_0\ge0$ and $\eta_0,\eta_1,\sqrt{e_0}\le\exp(-x^{\varepsilon/8})$.
\end{theorem}

\begin{theorem}[share of irreducible rows]\label{w2:I}
Let $J\ge1$ be the fixed integer of the construction and $\varepsilon_z=\varepsilon/J$. Let $E_{\mathrm{row}}\subset\mathcal H$ be the span of the image of $R\otimes_{K,\sigma_v}\C$ under the row embedding, and define $\delta$ by $D_{\mathrm{rows}}=\dim E_{\mathrm{row}}=(1-\delta)D_{\mathrm{all}}$. For every family of data of Part~II and every $x$ large in terms of the family constants,
\[
0\le\delta\le\exp\bigl(-x^{\varepsilon_z/7}\bigr).
\]
\end{theorem}

\begin{theorem}[Bessel identity]\label{w2:B}
For every family of data of Part~II and every $x$ large, the following holds at the complex place $v$. There are a unitary operator $\mathcal V$ of $Y_v$, and for each $u\in W$ a unitary operator $\mathcal U_u$ of $V_{u,v}$ and a number $a_u\ne0$, such that
\begin{equation}\label{w2:besselid}
(B_{u,v}\varphi)(F)=a_u\,\bigl\langle\mathbf s(\mathcal VF),\,\mathcal P_u\iota_u\mathcal U_u\varphi\bigr\rangle_{\mathcal H}\qquad(\varphi\in V_{u,v},\ F\in Y_v),
\end{equation}
where $\mathbf s\colon Y_v\to\mathcal H$ is the conjugate linear isometric row embedding and $\iota_u\colon V_{u,v}\to\mathcal C_u$ is the linear isometric column embedding of Definition~\ref{r3:bessel}. In addition $B_u$ is $K$-rational in the sense that $(B_{u,v}(\varphi\otimes1))(y\otimes1)=\sigma_v\bigl(B_u(\varphi)(y)\bigr)$ for $\varphi\in V_u$ and $y\in Y$.
\end{theorem}

Theorem~\ref{w2:B} together with Theorem~\ref{w2:I} is hypothesis (d) of Theorem~\ref{r3:A3}, and Theorem~\ref{w2:S} is hypothesis (e) there. 

\subsection{The Bessel identity}\label{w2:sec:bessel}

The identity \eqref{w2:besselid} is the end of a chain with four links. The period map is defined by a period of the chain of theta lifts on the definite side. The seesaw relation transfers every linear relation among those periods to the local Bessel periods of the theta lifts on the indefinite side. Unfolding the Poincar\'e series transfers every linear relation among local Bessel periods to the pairings with Poincar\'e series. One nonzero pairing then turns the transfer of relations into proportionality. Section~\ref{w2:sec:unfold} proves the unfolding, Section~\ref{w2:sec:prop} proves the passage from transfer of relations to proportionality, and Section~\ref{w2:sec:assembleB} assembles the theorem.

\subsubsection{Bessel Poincar\'e series and their unfolding}\label{w2:sec:unfold}

Fix a word $u$ and drop it from the notation where no confusion arises. The construction gives the following objects.

\begin{enumerate}[label=(P\arabic*),leftmargin=*]
\item\label{w2:P1} A closed $F$-subgroup $S=N\rtimes H$ of $G$, the Bessel subgroup of the flag of $u$. Here $N$ is a unipotent $F$-group and $H=\mathrm U(W)$ for a hermitian space $W$ over $E$ that is positive definite at every real place of $F$. A unitary character $\psi\colon N(\A)\to\C^\times$, trivial on $N(F)$ and invariant under conjugation by $H(\A)$.
\item\label{w2:P2} A compact open subgroup $K_2\subset G(\Af)$ with a unitary character $\chi_2$, and an element $x_u\in G(\Af)$.
\item\label{w2:P3} A finite dimensional unitary representation $(\sigma,V_\sigma)$ of the compact group $H_\infty=H(F\otimes\R)$, extended to $S_\infty=N_\infty H_\infty$ by $\sigma(nh)=\psi_\infty(n)\sigma(h)$, and a measurable function $W_u\colon G_\infty\to V_\sigma$ with $W_u(sy)=\sigma(s)W_u(y)$ for $s\in S_\infty$ and
\begin{equation}\label{w2:Wint}
\int_{S_\infty\backslash G_\infty}\lVert W_u(y)\rVert\,d\nu(y)<\infty .
\end{equation}
\item\label{w2:P4} The space $\mathcal C$ of column forms. A column form is a linear map $T$ from $V_\sigma$ to functions on $H(\A)$ with $T(\xi)(\gamma a)=T(\xi)(a)$ for $\gamma\in H(F)$, $T(\sigma(h)\xi)(a)=T(\xi)(ah)$ for $h\in H_\infty$, and $\lvert T(\xi)(a)\rvert\le C_T\lVert\xi\rVert$.
\item\label{w2:P5} For each $T\in\mathcal C$ a measurable seed $\mathrm{Ph}_T$ on $G(\A)$, supported on $G_\infty\times S(\Af)x_uK_2$, with $\mathrm{Ph}_T(gk)=\chi_2(k)\mathrm{Ph}_T(g)$ for $k\in K_2$, and
\begin{equation}\label{w2:seed}
\mathrm{Ph}_T\bigl(na\cdot(y,x_u)\bigr)=\psi(n)\,T\bigl(W_u(y)\bigr)(a)\qquad(n\in N(\A),\ a\in H(\A),\ y\in G_\infty).
\end{equation}
\end{enumerate}

Measures are fixed as follows. $G(\A)=G_\infty\times G(\Af)$ has a product Haar measure, $S(\A)$, $N(\A)$ and $H(\A)$ have Haar measures with $ds=dn\,da$, discrete groups have counting measure, and $\nu$ is the quotient measure on $S_\infty\backslash G_\infty$. Lemma~\ref{w2:unimod} below shows that these quotient measures exist.

For a bounded measurable function $\theta$ on $G(\A)$ that is left $N(F)$-invariant put
\begin{equation}\label{w2:per}
\per(\Psi,\theta)=\int_{H(F)\backslash H(\A)}\int_{N(F)\backslash N(\A)}\theta(na)\,\psi(n)\,\Psi(a)\,dn\,da
\end{equation}
for every bounded measurable left $H(F)$-invariant function $\Psi$ on $H(\A)$. This is the Bessel period of $\theta$ against $\Psi$.

\begin{lemma}\label{w2:unimod}
Under \ref{w2:P1} the groups $N(\A)$, $H(\A)$ and $S(\A)$ are unimodular, the quotients $N(F)\backslash N(\A)$ and $H(F)\backslash H(\A)$ are compact, and $S(F)\backslash S(\A)$ is compact.
\end{lemma}

\begin{proof}
The modular function of $S(\A)=N(\A)\rtimes H(\A)$ is $nh\mapsto\lvert\det\mathrm{Ad}(h)|_{\mathrm{Lie}\,N}\rvert_\A^{-1}$ times the modular functions of $N(\A)$ and $H(\A)$. The map $h\mapsto\det\mathrm{Ad}(h)|_{\mathrm{Lie}\,N}$ is an $F$-rational character of $H=\mathrm U(W)$. Every $F$-rational character of $\mathrm U(W)$ factors through the determinant $\mathrm U(W)\to\mathrm U(1)$, and $\mathrm U(1)=\ker(N_{E/F}\colon R_{E/F}\mathbb G_m\to\mathbb G_m)$ is an anisotropic torus, so it has no nontrivial $F$-rational character. Hence $\det\mathrm{Ad}(h)|_{\mathrm{Lie}\,N}=1$ identically, so at every place and on the adelic points. The same argument shows that $S_\infty$ and $S(\Af)$ are unimodular. Unipotent groups and reductive groups have unimodular adelic points, so $S(\A)$ is unimodular. The quotient $N(F)\backslash N(\A)$ is compact for every unipotent $F$-group \cite{w2borel}. Since $W$ is positive definite at a real place, it has no isotropic vector over $F$, so $H$ is $F$-anisotropic and $H(F)\backslash H(\A)$ is compact by the theorem of Borel and Harish-Chandra and of Mostow and Tamagawa \cite{w2bhc}. The image of the compact set $C_N\times C_H$ under multiplication, with $C_N$ and $C_H$ compact sets mapping onto the two quotients, maps onto $S(F)\backslash S(\A)$, because $S(F)=N(F)H(F)$ and $H(\A)$ normalizes $N(\A)$.
\end{proof}

\begin{proposition}[unfolding]\label{w2:unfold}
Assume \ref{w2:P1} to \ref{w2:P5} and let $T\in\mathcal C$. The series
\[
P_T(g)=\sum_{\gamma\in S(F)\backslash G(F)}\mathrm{Ph}_T(\gamma g)
\]
converges absolutely for almost every $g\in G(\A)$, defines an element of $L^1([G])$ that does not depend on the choice of representatives, and is right $(K_2,\chi_2)$-equivariant. There is a number $c_u>0$, depending on the Haar measures and on $u$ and not on $T$, such that for every bounded measurable function $\theta$ on $G(\A)$ that is left $G(F)$-invariant and right $(K_2,\overline{\chi_2})$-equivariant,
\begin{equation}\label{w2:perlaw}
\int_{[G]}\theta(g)\,P_T(g)\,dg=c_u\int_{S_\infty\backslash G_\infty}\per\Bigl(T\bigl(W_u(y)\bigr),\ R_{(y,x_u)}\theta\Bigr)\,d\nu(y),
\end{equation}
where $R_{(y,x_u)}\theta(g)=\theta(g\cdot(y,x_u))$, and both sides converge absolutely.
\end{proposition}

\begin{proof}
\emph{Step 1, integration in stages.} Since $G(\A)$, $S(\A)$, $S(F)$ and $G(F)$ are unimodular (Lemma~\ref{w2:unimod}), the quotient integral formula of Weil \cite[\S2.6]{w2folland} applies to $S(F)\subset S(\A)\subset G(\A)$ and to $S(F)\subset G(F)\subset G(\A)$. For every measurable $\Phi\ge0$ on $G(\A)$ that is left $S(F)$-invariant,
\begin{align}
\int_{S(F)\backslash G(\A)}\Phi&=\int_{S(\A)\backslash G(\A)}\int_{S(F)\backslash S(\A)}\Phi(s\dot g)\,ds\,d\dot g,\label{w2:stage1}\\
\int_{S(F)\backslash G(\A)}\Phi&=\int_{[G]}\sum_{\gamma\in S(F)\backslash G(F)}\Phi(\gamma g)\,dg.\label{w2:stage2}
\end{align}
The second identity is the unfolding. To see it directly, let $\mathfrak F$ be a measurable fundamental domain for $G(F)$ in $G(\A)$, which exists because $G(F)$ is a countable discrete subgroup of a second countable locally compact group, and let $D$ be a set of representatives of $S(F)\backslash G(F)$. The sets $\delta\gamma\mathfrak F$ with $\delta\in S(F)$ and $\gamma\in D$ are disjoint and cover $G(\A)$, so $\bigcup_{\gamma\in D}\gamma\mathfrak F$ is a fundamental domain for $S(F)$. Left invariance of Haar measure and Tonelli's theorem give \eqref{w2:stage2}.

\emph{Step 2, the finite part.} Put $\Omega=S(\Af)x_uK_2$, an open subset of $G(\Af)$. Let $\Phi\ge0$ be measurable on $G(\A)$ with $\Phi(sg)=\Phi(g)$ for $s\in S(\A)$, right $K_2$-invariant and supported on $G_\infty\times\Omega$. Then $\Phi$ factors through $S(\A)\backslash G(\A)=(S_\infty\backslash G_\infty)\times(S(\Af)\backslash G(\Af))$ and is a function of the first coordinate on the subset $S(\Af)\backslash\Omega$, which is the image of $K_2$ under $k\mapsto S(\Af)x_uk$. The quotient measure on a product of quotients is the product of the quotient measures, so
\begin{equation}\label{w2:finite}
\int_{S(\A)\backslash G(\A)}\Phi(\dot g)\,d\dot g=c_u\int_{S_\infty\backslash G_\infty}\Phi\bigl((y,x_u)\bigr)\,d\nu(y),
\end{equation}
where $c_u$ is the quotient measure of $S(\Af)\backslash\Omega$. It is finite because $S(\Af)\backslash\Omega$ is the continuous image of the compact group $K_2$, and positive because it is open.

\emph{Step 3, absolute convergence.} Let $\theta$ be as in the statement. The product $\Phi_0=\lvert\theta\,\mathrm{Ph}_T\rvert$ is left $S(F)$-invariant, because $\theta$ is left $G(F)$-invariant and $\mathrm{Ph}_T$ is left $S(F)$-invariant. For the latter, the support $G_\infty\times\Omega=S(\A)\cdot(G_\infty\times\{x_u\})\cdot K_2$ is stable under $S(F)$, and for $\gamma=n_0a_0\in S(F)$, \eqref{w2:seed} gives $\mathrm{Ph}_T(\gamma na(y,x_u))=\psi(n_0)\psi(a_0na_0^{-1})T(W_u(y))(a_0a)=\psi(n)T(W_u(y))(a)$, by the triviality of $\psi$ on $N(F)$, its invariance under $H(\A)$, and the $H(F)$-invariance in \ref{w2:P4}. It is right $K_2$-invariant, because the characters $\chi_2$ and $\overline{\chi_2}$ cancel, and it is supported on $G_\infty\times\Omega$. By \eqref{w2:stage1}, then \eqref{w2:finite} applied to the left $S(\A)$-invariant function $g\mapsto\int_{S(F)\backslash S(\A)}\Phi_0(sg)\,ds$, then \eqref{w2:seed} and the integration formula $ds=dn\,da$ on $S(F)\backslash S(\A)$,
\begin{align*}
\int_{S(F)\backslash G(\A)}\Phi_0
&=c_u\int_{S_\infty\backslash G_\infty}\int_{[H]}\int_{[N]}\bigl\lvert\theta(na(y,x_u))\bigr\rvert\,\bigl\lvert T(W_u(y))(a)\bigr\rvert\,dn\,da\,d\nu(y)\\
&\le c_u\lVert\theta\rVert_\infty\vol[N]\vol[H]\,C_T\int_{S_\infty\backslash G_\infty}\lVert W_u\rVert\,d\nu,
\end{align*}
which is finite by Lemma~\ref{w2:unimod} and \eqref{w2:Wint}. The same computation applies to $\Phi_0=\lvert\mathrm{Ph}_T\rvert$, which is right $K_2$-invariant because $\lvert\chi_2\rvert=1$, and with \eqref{w2:stage2} it shows that $\sum_\gamma\lvert\mathrm{Ph}_T(\gamma g)\rvert$ is integrable on $[G]$, hence finite almost everywhere, so $P_T$ converges absolutely almost everywhere and lies in $L^1([G])$. Independence of representatives is the left $S(F)$-invariance of $\mathrm{Ph}_T$, and right equivariance under $(K_2,\chi_2)$ passes to the sum.

\emph{Step 4, the identity.} Since $\theta\,\mathrm{Ph}_T$ is absolutely integrable on $S(F)\backslash G(\A)$ by Step 3, the identities \eqref{w2:stage2}, \eqref{w2:stage1} and \eqref{w2:finite} hold for it by Fubini's theorem applied to the positive and negative parts of its real and imaginary parts. The left side of \eqref{w2:stage2} is $\int_{[G]}\theta P_T$ because $\theta$ is left $G(F)$-invariant. On the right, \eqref{w2:seed} gives
\[
\begin{aligned}
\int_{S(F)\backslash S(\A)}(\theta\,\mathrm{Ph}_T)(s(y,x_u))\,ds&=\int_{[H]}\int_{[N]}\theta(na(y,x_u))\,\psi(n)\,T(W_u(y))(a)\,dn\,da\\
&=\per\bigl(T(W_u(y)),R_{(y,x_u)}\theta\bigr).
\end{aligned}
\]
This is \eqref{w2:perlaw}.
\end{proof}

\begin{remark}\label{w2:rem:perlaw}
The argument shows $c_u>0$. The identity \eqref{w2:perlaw} does not need $P_T\in L^2([G])$, nor that $P_T$ is fixed by the archimedean projector. Both facts are needed to know that $P_T\in\mathcal H$, which Definition~\ref{w2:anchor} uses.
\end{remark}

\begin{corollary}[transfer of relations from Bessel periods to Poincar\'e pairings]\label{w2:bottom}
Assume \ref{w2:P1} to \ref{w2:P5}. Let $T_1,\dots,T_\ell\in\mathcal C$, let $\theta_1,\dots,\theta_\ell$ be bounded measurable functions on $G(\A)$, left $G(F)$-invariant and right $(K_2,\overline{\chi_2})$-equivariant, and let $\alpha_1,\dots,\alpha_\ell\in\C$. If
\begin{equation}\label{w2:bothyp}
\sum_{i=1}^\ell\alpha_i\,\per\bigl(T_i(\xi),\,R_{(y,x_u)}\theta_i\bigr)=0\qquad\text{for every }\xi\in V_\sigma\text{ and every }y\in G_\infty,
\end{equation}
then $\sum_i\alpha_i\int_{[G]}\theta_iP_{T_i}=0$. If \eqref{w2:bothyp} holds for $\xi$ in a spanning set of $V_\sigma$ and every $y$, it holds for every $\xi$.
\end{corollary}

\begin{proof}
The period \eqref{w2:per} is linear in $\Psi$ and $T_i$ is linear, so the left side of \eqref{w2:bothyp} is linear in $\xi$, which proves the last sentence. Apply Proposition~\ref{w2:unfold} to each $i$ with the same $c_u$, and put $\xi=W_u(y)$ in \eqref{w2:bothyp}. The integrand of the sum of the right sides of \eqref{w2:perlaw} vanishes for every $y$.
\end{proof}

\subsubsection{From transfer of relations to proportionality}\label{w2:sec:prop}

\begin{lemma}\label{w2:prop}
Let $X$ be a set and $\lambda,\mu\colon X\to\C$ functions. Assume that for every finite family $x_1,\dots,x_\ell\in X$ and $\alpha_1,\dots,\alpha_\ell\in\C$ with $\sum_i\alpha_i\lambda(x_i)=0$ one has $\sum_i\alpha_i\mu(x_i)=0$. Assume $\mu(x_0)\ne0$ for some $x_0\in X$. Then $\lambda(x_0)\ne0$, and $\mu=\kappa\lambda$ on $X$ with $\kappa=\mu(x_0)/\lambda(x_0)\ne0$.
\end{lemma}

\begin{proof}
If $\lambda(x_0)=0$, the family consisting of $x_0$ with coefficient $1$ satisfies the hypothesis, so $\mu(x_0)=0$, a contradiction. For $x\in X$ the family $(x_0,x)$ with coefficients $(\lambda(x),-\lambda(x_0))$ satisfies $\lambda(x)\lambda(x_0)-\lambda(x_0)\lambda(x)=0$, so $\lambda(x)\mu(x_0)-\lambda(x_0)\mu(x)=0$.
\end{proof}

\begin{lemma}[central translations]\label{w2:central}
Let $z$ be an element of the center of $G(\Af)$ and $R_z\theta(g)=\theta(gz)$. Then $R_z$ is a unitary operator of $L^2([G])$ that commutes with the right action of $G(\A)$. It maps the right $(K_2,\chi_2)$-equivariant functions to themselves and preserves every $G(\A)$-isotypic subspace, in particular $\mathcal H$.
\end{lemma}

\begin{proof}
$R_z$ is well defined on $[G]$ and preserves Haar measure because $G(\A)$ is unimodular, and it is invertible with inverse $R_{z^{-1}}$. Since $z$ is central, $R_zR_h=R_hR_z$ for $h\in G(\A)$, which gives the last two assertions, and $\theta(gkz)=\theta(gzk)=\chi_2(k)\theta(gz)$ for $k\in K_2$.
\end{proof}

\subsubsection{Assembly of the Bessel identity}\label{w2:sec:assembleB}

The theorem uses four inputs of the construction. The first is a definition, the second is the seesaw relation of the theta tower, the third is a pair of dictionaries between the abstract realization of Definition~\ref{r3:bessel} and the automorphic objects of Section~\ref{w2:sec:unfold}, and the fourth is rationality. They are stated here in the form used, for one word $u$ at the place $v$.

\begin{hypothesis}[the period map and the seesaw relation]\label{w2:H-see}
There are a central element $z\in G(\Af)$, a linear map $U_u$ of $V_{u,v}$ with $\lVert U_u\varphi\rVert=r_u\lVert\varphi\rVert$ for some $r_u>0$, a set $X_u$ containing $V_{u,v}\times Y_v$, a function $\lambda_u\colon X_u\to\C$ with $\lambda_u(\varphi,F)=(B_{u,v}\varphi)(F)$ on $V_{u,v}\times Y_v$, and for each $x\in X_u$ a column form $T_x\in\mathcal C_u$ and a function $\theta_x\in\mathcal H$, with $T_{(\varphi,F)}=\iota_u(U_u\varphi)$ and $\theta_{(\varphi,F)}=\mathbf s(F)$ on $V_{u,v}\times Y_v$, such that the following holds. For every finite family $x_i\in X_u$ and $\alpha_i\in\C$ with $\sum_i\alpha_i\lambda_u(x_i)=0$,
\[
\sum_i\alpha_i\,\per\Bigl(T_{x_i}(\xi),\,R_{(y,x_u)}R_z\overline{\theta_{x_i}}\Bigr)=0\qquad\text{for every }\xi\in V_\sigma\text{ and }y\in G_\infty.
\]
For some $x_0\in X_u$ the pairing $\int_{[G]}R_z\overline{\theta_{x_0}}\,P_{T_{x_0}}$ is nonzero.
\end{hypothesis}

\begin{hypothesis}[dictionaries]\label{w2:H-dict}
The space $\mathcal H$ consists of bounded continuous functions on $G(\A)$, left $G(F)$-invariant and right $(K_2,\chi_2)$-equivariant. The column space $\mathcal C_u$ of Definition~\ref{r3:bessel} is the space $\mathcal C$ of \ref{w2:P4} with the norm for which $\mathcal P_u$ is the map $T\mapsto P_T$ of Proposition~\ref{w2:unfold}, and $P_T\in\mathcal H$ for every $T$. The central translation $R_z$ maps $\mathbf s(Y_v)$ onto itself.
\end{hypothesis}

\begin{hypothesis}[rationality]\label{w2:H-rat}
For $\varphi\in V_u$ and $y\in Y$ the number $(B_{u,v}(\varphi\otimes1))(y\otimes1)$ lies in $\sigma_v(K)$ and equals $\sigma_v(B_u(\varphi)(y))$.
\end{hypothesis}

\begin{definition}\label{w2:anchor}
For $x\in X_u$ put $\mu_u(x)=\int_{[G]}R_z\overline{\theta_x}\,P_{T_x}$.
The last condition of Hypothesis~\ref{w2:H-see} asks for a pair $x_0$
with $\mu_u(x_0)\ne0$.
\end{definition}

\begin{lemma}\label{w2:punz}
Assume $\lVert K_{uu}-1\rVert_\HS^2\le e_0\,d^S_u$ and $\dim\iota_u(V_{u,v})>e_0\,d^S_u$. Then there is $\varphi_0\in V_{u,v}$ with $\mathcal P_u\iota_u(U_u\varphi_0)\ne0$.
\end{lemma}

\begin{proof}
For $w\in\mathcal C_u$ one has $\lVert\mathcal P_uw\rVert^2=c_u\langle w,K_{uu}w\rangle$, so $\ker\mathcal P_u=\ker K_{uu}$, because $K_{uu}$ is positive semidefinite. On this kernel $K_{uu}-1$ acts as $-1$, so $\dim\ker\mathcal P_u\le\lVert K_{uu}-1\rVert_\HS^2\le e_0d^S_u$. Since $U_u$ is invertible, $\iota_uU_u(V_{u,v})=\iota_u(V_{u,v})$ has dimension larger than $\dim\ker\mathcal P_u$ and is not contained in it.
\end{proof}

Under the conclusion of Theorem~\ref{w2:S}, $e_0\le\tfrac12$ for large $x$, and $\dim\iota_u(V_{u,v})=d_u\ge(1-2f/T_u)d^S_u>\tfrac12d^S_u$ by Lemma~\ref{r3:ramshare}, since $T_u\ge b>8f$. So Lemma~\ref{w2:punz} applies.

\begin{theorem}\label{w2:Bproof}
Assume \ref{w2:P1} to \ref{w2:P5} for the word $u$, Hypotheses~\ref{w2:H-see}, \ref{w2:H-dict} and \ref{w2:H-rat}. Then \eqref{w2:besselid} holds for $u$ with $\mathcal V$ the unitary operator of $Y_v$ defined by $\mathbf s(\mathcal VF)=R_z\mathbf s(F)$, with $\mathcal U_u=r_u^{-1}U_u$ and with a number $a_u\ne0$. If the hypotheses hold for every $u\in W$ with one $z$, then $\mathcal V$ does not depend on $u$ and Theorem~\ref{w2:B} holds.
\end{theorem}

\begin{proof}
By Lemma~\ref{w2:central} and Hypothesis~\ref{w2:H-dict}, $R_z$ restricts to a unitary operator of $\mathbf s(Y_v)$, and since $\mathbf s$ is a conjugate linear isometry onto its image, $\mathcal V=\mathbf s^{-1}R_z\mathbf s$ is a linear unitary operator of $Y_v$.

Let $\mu_u$ be as in Definition~\ref{w2:anchor}. Every $\theta_x\in\mathcal H$ is bounded, and so is $R_z\overline{\theta_x}$, which is left $G(F)$-invariant and right $(K_2,\overline{\chi_2})$-equivariant by Lemma~\ref{w2:central}. Hypothesis~\ref{w2:H-see} and Corollary~\ref{w2:bottom} show that every linear relation among the values of $\lambda_u$ holds among the values of $\mu_u$. By Hypothesis~\ref{w2:H-see}, $\mu_u(x_0)\ne0$. Lemma~\ref{w2:prop} gives $\kappa_u\ne0$ with $\mu_u=\kappa_u\lambda_u$ on $X_u$. On $V_{u,v}\times Y_v$ this reads
\[
\kappa_u\,(B_{u,v}\varphi)(F)=\int_{[G]}\overline{R_z\mathbf s(F)}\;\mathcal P_u\iota_u(U_u\varphi)=\bigl\langle\mathbf s(\mathcal VF),\mathcal P_u\iota_u(U_u\varphi)\bigr\rangle_{\mathcal H}=r_u\bigl\langle\mathbf s(\mathcal VF),\mathcal P_u\iota_u\mathcal U_u\varphi\bigr\rangle_{\mathcal H}.
\]
So \eqref{w2:besselid} holds with $a_u=r_u/\kappa_u\ne0$, and $\mathcal U_u=r_u^{-1}U_u$ is unitary because it is a linear isometry of a finite dimensional space. The rationality clause is Hypothesis~\ref{w2:H-rat}.
\end{proof}

\begin{remark}[what Theorem~\ref{w2:Bproof} leaves open]\label{w2:rem:Bopen}
Hypothesis~\ref{w2:H-see} combines a definition and a relation. The period map is defined from the period $\mathrm{perCol}$ of the chain of theta lifts on the definite side. The seesaw relation transfers relations among the numbers $c(t)\,\mathrm{perCol}$ to relations among the local Bessel periods of the indefinite theta lifts translated by $z$, with $c(t)=c_{\mathrm{top}}(t)s/c_{\mathrm{bot}}(t)$. The constants of the double step come from the volumes and local scalar factors in Theorem~\ref{w3s:thm:T2}. The Siegel--Weil and doubling identities enter the surjectivity in Hypothesis~\ref{w2:see:H-links}.
\end{remark}

\subsubsection{The chain of seesaw transfers}\label{w2:sec:seesaw}

The transfer follows from the theta tower identities in Hypothesis~\ref{w2:see:H-laws} and the period identifications in Hypothesis~\ref{w2:see:H-links}.

\paragraph{The tower.} Let $m=2n+2$ be the (even) number of lifts of the chain. The construction gives hermitian spaces $V_0,\dots,V_m$ and $T$ over $E$, definite hermitian spaces $X_0,\dots,X_m$ with an orthogonal decomposition $X_m=T\oplus L$, and $h\mapsto\epsilon(h)$ the induced embedding $\mathrm U(T)\to\mathrm U(X_m)$. For $0\le j$ and $\rho\ge0$ with $j+2\rho+2=m$ there is a Bessel datum of rank $\rho$ at level $j$, and $\mathcal B^\rho_j(\phi,f)$ denotes the Bessel period \eqref{w2:per} of a function $f$ on $\mathrm U(V_j)(\A)$ against a function $\phi$ on $\mathrm U(X_j)(\A)$, with the fixed additive character $\psi$. For $e=\sum_ge(g)\,g$ a finitely supported combination of finite adelic points put $\mathrm{tr}(e)f=\sum_ge(g)R_gf$.

For $j\le m-1$ let $\mathcal S_j$ be the space of smooth vectors of the archimedean Fock model at level $j$ and $\mathcal S_T$ the one of $T$. The space $\mathcal C_j$ consists of linear maps $G\colon\mathcal S_j\to\mathcal A(\mathrm U(V_j))$, and $\mathcal A$ denotes the space of automorphic functions. Define $\mathcal C_T$ for $T$ in the same way. The theta lifts of the chain are linear maps
\[
\iota_j\colon\mathcal C_{j+1}\to\mathcal C_j\quad(j\le m-2),\qquad\iota_T\colon\mathcal C_T\to\mathcal C_{m-1},
\]
and $\Theta_j(f;P)$ denotes the theta integral of a function $f$ on $\mathrm U(V_{j+1})$ against the theta kernel attached to a polynomial Fock vector $P$, with $\Theta_T$ the same for $T$. The maps $\mathrm{It}_j\colon\mathcal P_j\otimes\mathcal S_{j+1}\to\mathcal S_j$ and $\mathrm{It}_T\colon\mathcal P_T\otimes\mathcal S_T\to\mathcal S_{m-1}$ are the tensor maps of the Fock models.

The tower is indexed by a set $\mathcal T$ of tuples. For each $s\in\mathcal T$ and each level $j$ the construction gives a finite set $\mathrm{Ch}_j$, combinations $e_{j,\alpha}$ with finite support for $\alpha\in\mathrm{Ch}_j$, a set $\mathrm{Col}_j$ of columns with a subset $\mathrm{Ram}_j$, maps $S_j\colon\mathrm{Col}_j\to\mathrm{Col}_j$ and $K_j\colon\mathrm{Col}_j\to\mathrm{Col}_{j+1}$ preserving the subsets $\mathrm{Ram}$, a set $\mathrm{Fr}_j$ of frames, reading functions $\mathrm{rd}_j(\alpha,\psi',y)$ on $\mathrm U(X_j)(\A)$ for $\psi'\in\mathrm{Col}_j$ and $y\in\mathrm{Fr}_j$, a tail $\mathrm{tail}(\psi',F)\in\C$ for $\psi'\in\mathrm{Col}_m$ and $F\in\mathcal C_T$, and complex numbers $c_I(j),c_V(j),c_K(j),c_P$. The dependence on $s$ is not written. Put
\begin{equation}\label{w2:see:per}
\mathrm{per}^\rho_j(\psi',G;y,v)=\sum_{\alpha\in\mathrm{Ch}_j}\mathcal B^\rho_j\bigl(\mathrm{rd}_j(\alpha,\psi',y),\ \mathrm{tr}(e_{j,\alpha})\,G(v)\bigr)\qquad(G\in\mathcal C_j,\ v\in\mathcal S_j),
\end{equation}
$\mathrm{up}_0\psi'=\psi'$ and $\mathrm{up}_{k+1}\psi'=K_{2k+1}K_{2k}S_{2k}\,\mathrm{up}_k\psi'$ for $\psi'\in\mathrm{Col}_0$, and define the top and bottom constants
\begin{equation}\label{w2:see:const}
W_k=\prod_{l=0}^{k-1}c_V(2l)\,c_K(2l)\,c_K(2l+1),\qquad c_{\mathrm{top}}=W_{n+1}\,c_P,\qquad c_{\mathrm{bot}}=\prod_{j=0}^{m-1}c_I(j).
\end{equation}
The lift $\iota_T$ has the factor $c_I(m-1)$, so $c_{\mathrm{bot}}$ is the product of the constants of all $m$ lifts.

\begin{hypothesis}[laws of the theta tower]\label{w2:see:H-laws}
For every tuple the following hold.
\begin{enumerate}[label=(T\arabic*),leftmargin=*]
\item\label{w2:see:T1} (Lift law.) For $G\in\mathcal C_{j+1}$, $P\in\mathcal P_j$ and $v\in\mathcal S_{j+1}$, $\Theta_j(G(v);P)=c_I(j)\,(\iota_jG)(\mathrm{It}_j(P\otimes v))$, and the same holds for $\iota_T$ with $c_I(m-1)$. All $c_I(j)$ are nonzero.
\item\label{w2:see:T2} (Double-step seesaw identity.) Let $j+2(\sigma+1)+2=m$, $P_1\in\mathcal P_j$, $P_2\in\mathcal P_{j+1}$, and let $a$ and $a'$ be archimedean points of $\mathrm U(V_{j+1})$ and $\mathrm U(V_{j+2})$. There is a bivacuum kernel $C$, depending on $(j,\sigma,P_1,P_2,a,a')$ and not on the tuple, such that for $\psi'\in\mathrm{Ram}_j$, $y$, $G\in\mathcal C_{j+2}$ and $v$, the integrand of $\mathcal B^{\sigma+1}_j$ applied to the double lift $\Theta_j(\Theta_{j+1}(G(v);P_2);P_1)$, integrated first over the Jacobi quotient at level $j+1$ translated by $a$ and then over the Bessel quotient at level $j+2$ translated by $a'$, converges absolutely and equals
\[
c_V(j)\sum_{\beta}\mathcal B^\sigma_{j+2}\bigl(\mathrm{two}_j(\mathrm{rd}_j(\cdot,S_j\psi',y),C)_\beta,\ \mathrm{tr}(e_{j+2,\beta})\,R_{a'}G(v)\bigr),
\]
where $\mathrm{two}_j$ is the bilift of the reading functions through the kernels of levels $j$ and $j+1$ composed with $C$.
\item\label{w2:see:T3} (Kernel law.) For every bivacuum kernel $C$ and every $y$ there are coefficients $a(y')$, finitely many nonzero and independent of the tuple and of $\psi'\in\mathrm{Ram}_j$, with
\[
\mathrm{two}_j(\mathrm{rd}_j(\cdot,\psi',y),C)_\beta=c_K(j)\,c_K(j+1)\sum_{y'}a(y')\,\mathrm{rd}_{j+2}\bigl(\beta,K_{j+1}K_j\psi',y'\bigr).
\]
\item\label{w2:see:T4} (Last-step identity.) For $j=m-2$, $P_1\in\mathcal P_j$, $P_2\in\mathcal P_T$ there is a bivacuum kernel $C$ independent of the tuple with
\[
\mathrm{per}^0_j\bigl(\psi',\Theta_j(\Theta_T(F(w);P_2);P_1);y\bigr)=c_V(j)\,\mathrm{perT}\bigl(\mathrm{two}_j(\mathrm{rd}_j(\cdot,S_j\psi',y),C),F(w)\bigr),
\]
where $\mathrm{perT}(\phi,f)=\sum_\beta\int_{[\mathrm U(T)]}\mathrm{tr}(e_{T,\beta})f(h)\,\phi_\beta(\epsilon(h))\,dh$.
\item\label{w2:see:T5} (Top period.) For $y$ and $w$ there is $b(y,w)\in\C$, independent of the tuple, of $\psi'\in\mathrm{Ram}_m$ and of $F$, with $\mathrm{perT}(\mathrm{rd}_m(\cdot,\psi',y),F(w))=c_P\,b(y,w)\,\mathrm{tail}(\psi',F)$.
\item\label{w2:see:T6} (Unfolding of the double step.) If for every pair $(a,a')$ of archimedean points the iterated integrals in \ref{w2:see:T2} of a finite family of tuples, with coefficients $\alpha_i$, have values $v_i$ with $\sum_i\alpha_iv_i=0$, then $\sum_i\alpha_i\,\mathcal B^{\sigma+1}_j$ of the corresponding double lifts vanishes.
\item\label{w2:see:T7} (Density.) The vectors $\mathrm{It}_j(P_1\otimes\mathrm{It}_{j+1}(P_2\otimes v))$, and the vectors $\mathrm{It}_j(P_1\otimes\mathrm{It}_T(P_2\otimes v))$, span a subspace on which a finite linear relation among the functionals $v\mapsto\mathrm{per}^\rho_j(\psi_i,G_i;y,v)$ implies the same relation on all of $\mathcal S_j$.
\item\label{w2:see:T8} (Archimedean translation.) For an archimedean point $a'$ of $\mathrm U(V_{j})$ there is a linear map $\tau_{a'}$ of $\mathcal S_j$ with $R_{a'}G(v)=G(\tau_{a'}v)$ for every $G\in\mathcal C_j$ and every tuple, and $\mathcal B^\rho_j$ is linear in each argument.
\end{enumerate}
\end{hypothesis}

\begin{proposition}[one double step]\label{w2:see:step}
Assume Hypothesis~\ref{w2:see:H-laws}. Let $j+2(\sigma+1)+2=m$, let $s_1,\dots,s_\ell$ be tuples, $\alpha_i\in\C$, $\psi_i\in\mathrm{Ram}_j$ and $G_i\in\mathcal C_{j+2}$, all for the tuple $s_i$. If
\[
\sum_i\alpha_i\,c_V(j)c_K(j)c_K(j+1)\,\mathrm{per}^\sigma_{j+2}\bigl(K_{j+1}K_jS_j\psi_i,G_i;y,v\bigr)=0\qquad\text{for all }y,v,
\]
then
\[
\sum_i\alpha_i\,c_I(j)c_I(j+1)\,\mathrm{per}^{\sigma+1}_j\bigl(\psi_i,\iota_j\iota_{j+1}G_i;y,v\bigr)=0\qquad\text{for all }y,v.
\]
\end{proposition}

\begin{proof}
Fix $y$. By \ref{w2:see:T7} it suffices to prove the second identity at $v=\mathrm{It}_j(P_1\otimes\mathrm{It}_{j+1}(P_2\otimes w))$. Two applications of \ref{w2:see:T1} give
\[
c_I(j)c_I(j+1)\,(\iota_j\iota_{j+1}G_i)(v)=\Theta_j\bigl(\Theta_{j+1}(G_i(w);P_2);P_1\bigr),
\]
so the left side of the second identity is $\sum_i\alpha_i\,\mathrm{per}^{\sigma+1}_j$ of the double lifts. By \ref{w2:see:T6} it suffices to show, for fixed archimedean points $(a,a')$, that the iterated integrals $v_i$ satisfy $\sum_i\alpha_iv_i=0$. Let $C$ be the kernel of \ref{w2:see:T2} for $(j,\sigma,P_1,P_2,a,a')$. Then \ref{w2:see:T2}, \ref{w2:see:T3} applied to $S_j\psi_i\in\mathrm{Ram}_j$, and \ref{w2:see:T8} give
\[
v_i=c_V(j)c_K(j)c_K(j+1)\sum_{y'}a(y')\,\mathrm{per}^\sigma_{j+2}\bigl(K_{j+1}K_jS_j\psi_i,G_i;y',\tau_{a'}w\bigr),
\]
with $a(y')$ and $\tau_{a'}$ independent of $i$. Hence $\sum_i\alpha_iv_i=\sum_{y'}a(y')\cdot0=0$ by the hypothesis at $(y',\tau_{a'}w)$.
\end{proof}

\begin{proposition}[last step]\label{w2:see:last}
Assume Hypothesis~\ref{w2:see:H-laws} and let $j=m-2$. Let $s_i$, $\alpha_i$, $\psi_i\in\mathrm{Ram}_j$ and $F_i\in\mathcal C_T$ be finitely many data. If
\[
\sum_i\alpha_i\,c_V(j)c_K(j)c_K(j+1)\,c_P\ \mathrm{tail}\bigl(K_{j+1}K_jS_j\psi_i,F_i\bigr)=0,
\]
then $\sum_i\alpha_i\,c_I(j)c_I(j+1)\,\mathrm{per}^0_j(\psi_i,\iota_j\iota_TF_i;y,v)=0$ for all $y,v$.
\end{proposition}

\begin{proof}
By \ref{w2:see:T7} take $v=\mathrm{It}_j(P_1\otimes\mathrm{It}_T(P_2\otimes w))$. By \ref{w2:see:T1} the $i$-th term is
\[
 \alpha_i\,\mathrm{per}^0_j(\psi_i,\Theta_j(\Theta_T(F_i(w);P_2);P_1);y).
\]
By \ref{w2:see:T4}, \ref{w2:see:T3} and \ref{w2:see:T5} this equals
\[
\sum_{y'}a(y')\,b(y',w)\,\alpha_i\,c_V(j)c_K(j)c_K(j+1)\,c_P\ \mathrm{tail}\bigl(K_{j+1}K_jS_j\psi_i,F_i\bigr).
\]
The sum over $i$ vanishes for each $y'$ by hypothesis.
\end{proof}

\begin{theorem}[transfer along the tower]\label{w2:see:main}
Assume Hypothesis~\ref{w2:see:H-laws}. Let $s_i$, $\alpha_i$, $\psi_i\in\mathrm{Ram}_0$ and $F_i\in\mathcal C_T$ be finitely many data. If
\[
\sum_i\alpha_i\,c_{\mathrm{top}}(s_i)\,\mathrm{tail}\bigl(\mathrm{up}_{n+1}\psi_i,F_i\bigr)=0,
\]
then for all $y\in\mathrm{Fr}_0$ and $v\in\mathcal S_0$
\[
\sum_i\alpha_i\,c_{\mathrm{bot}}(s_i)\,\mathrm{per}^n_0\bigl(\psi_i,\ \iota_0\iota_1\cdots\iota_{m-2}\iota_TF_i;y,v\bigr)=0 .
\]
\end{theorem}

\begin{proof}
Since $S$ and $K$ preserve the subsets $\mathrm{Ram}$, $\mathrm{up}_k\psi_i\in\mathrm{Ram}_{2k}$ for $2k\le m$. Since $\mathrm{up}_{n+1}\psi_i=K_{2n+1}K_{2n}S_{2n}\,\mathrm{up}_n\psi_i$ and $c_{\mathrm{top}}=W_n\,c_V(2n)c_K(2n)c_K(2n+1)\,c_P$, Proposition~\ref{w2:see:last} with $j=2n$, columns $\mathrm{up}_n\psi_i$ and coefficients $\alpha_iW_n$ gives
\begin{equation}\label{w2:see:start}
\sum_i\alpha_iW_n\,c_I(2n)c_I(2n+1)\,\mathrm{per}^0_{2n}\bigl(\mathrm{up}_n\psi_i,\iota_{2n}\iota_TF_i;y,v\bigr)=0 .
\end{equation}
We show by induction on $k$ the following statement $(\mathrm A_k)$. If $2k+2\rho+2=m$, $G_i\in\mathcal C_{2k}$ and $\sum_i\beta_iW_k\,\mathrm{per}^\rho_{2k}(\mathrm{up}_k\psi_i,G_i;y,v)=0$ for all $y,v$, then $\sum_i\beta_i\prod_{j<2k}c_I(j)\,\mathrm{per}^{k+\rho}_0(\psi_i,\iota_0\cdots\iota_{2k-1}G_i;y,v)=0$ for all $y,v$. For $k=0$ there is nothing to prove. If $(\mathrm A_k)$ holds and the hypothesis of $(\mathrm A_{k+1})$ holds with $\rho$, then $W_{k+1}=W_k\,c_V(2k)c_K(2k)c_K(2k+1)$, and Proposition~\ref{w2:see:step} at $j=2k$, $\sigma=\rho$ with coefficients $\beta_iW_k$ gives $\sum_i\beta_iW_k\,c_I(2k)c_I(2k+1)\,\mathrm{per}^{\rho+1}_{2k}(\mathrm{up}_k\psi_i,\iota_{2k}\iota_{2k+1}G_i;y,v)=0$. Then $(\mathrm A_k)$ with $\rho+1$ and coefficients $\beta_ic_I(2k)c_I(2k+1)$ gives $(\mathrm A_{k+1})$. Finally $(\mathrm A_n)$ with $\rho=0$, $\beta_i=\alpha_ic_I(2n)c_I(2n+1)$ and $G_i=\iota_{2n}\iota_TF_i$ applies to \eqref{w2:see:start}, and $\prod_{j<2n}c_I(j)\cdot c_I(2n)c_I(2n+1)=c_{\mathrm{bot}}$ because $m=2n+2$.
\end{proof}

\paragraph{The word $u$.} Fix a word $u$. The tower attached to $u$ has its tuples indexed by the set of tuples $t$ of the chain datum, with the base tuple $\mu_0$, and $\mathrm{perCol}_t(\varphi,F)$ denotes the period of the chain on the definite side at the tuple $t$, for $\varphi\in V_{u,v}$ and a $K$-rational vector $F$ of the good part at $t$. By the definition of the period map, $(B_{u,v}\varphi)(\iota_0(q))=\mathrm{perCol}_{\mu_0}(\varphi,q)$ for $q$ in the good part at $\mu_0$, and these vectors $\iota_0(q)$ span $Y_v$. Let $\mathcal R_0$ be the space of row forms at level $0$, the span of the theta lifts of level $0$ that the construction uses, and for $F$ as above let $\mathcal R_F\in\mathcal R_0$ be the row form of $F$. Let $b\in\mathrm U(V_0)(\A)$ be the translation of the Bessel datum of the word, with $z$ and $x_u$ as in Hypothesis~\ref{w2:H-see} related to it by the pinning relation of the construction, and write
\[
\mathfrak b_u(\varphi,\mathcal R;y,v)=\mathcal B^n_0\bigl(\mathrm{rd}^u_0(\varphi,y),\ R_b\,\mathcal R(v)\bigr)
\]
for the local Bessel period at level $0$ of the translated row form $\mathcal R$ against the column $\varphi$.

\begin{hypothesis}[dictionaries of the tower]\label{w2:see:H-links}
There are maps $\mathrm{col}_0(t,\cdot)\colon V_{u,v}\to\mathrm{Col}_0$ and $\mathrm{col}_m\colon\mathrm{Col}_0\to\mathrm{Col}_m$ such that the following hold for every tuple $t$.
\begin{enumerate}[label=(D\arabic*),leftmargin=*]
\item\label{w2:see:D1} $\mathrm{col}_0(t,\varphi)\in\mathrm{Ram}_0$ and $\mathrm{per}^n_0(\mathrm{col}_0(t,\varphi),G;y,v)=\mathfrak b_u(\varphi,G;y,v)$ for every $G\in\mathcal C_0$.
\item\label{w2:see:D2} If $G\in\mathcal C_T$ is the leading map of $F$, then $\mathrm{tail}(\mathrm{col}_m\mathrm{col}_0(t,\varphi),G)=\mathrm{perCol}_t(\varphi,F)$.
\item\label{w2:see:D3} For every $F$ there is $G\in\mathcal C_T$, the leading map of $F$, with $\iota_0\cdots\iota_{m-2}\iota_TG=\mathcal R_F$.
\item\label{w2:see:D4} Every $\mathcal R\in\mathcal R_0$ equals $\iota_0\cdots\iota_{m-2}\iota_TG$ for some $G\in\mathcal C_T$ at the base tuple $\mu_0$.
\item\label{w2:see:D5} There are a linear bijection $S$ of $V_{u,v}$, $r>0$ and $\varsigma\ne0$ with $\lVert S\varphi\rVert=r\lVert\varphi\rVert$ and $\mathrm{tail}(\mathrm{up}_{n+1}\mathrm{col}_0(t,\varphi),G)=\varsigma\,\mathrm{tail}(\mathrm{col}_m\mathrm{col}_0(t,S\varphi),G)$ for all $t$, $\varphi$ and $G$.
\item\label{w2:see:D6} $c_{\mathrm{top}}(t)\ne0$.
\end{enumerate}
\end{hypothesis}

\begin{corollary}[the seesaw relation]\label{w2:see:cor}
Assume Hypotheses~\ref{w2:see:H-laws} and~\ref{w2:see:H-links}. Put $U_u=S^{-1}$, $r_u=r^{-1}$ and
\begin{equation}\label{w2:see:ct}
c(t)=\frac{c_{\mathrm{top}}(t)\,\varsigma}{c_{\mathrm{bot}}(t)}\ne0,\qquad p_{\mathcal R}(\varphi_0)=\frac{c_{\mathrm{top}}(\mu_0)}{c_{\mathrm{bot}}(\mu_0)}\,\mathrm{tail}\bigl(\mathrm{up}_{n+1}\mathrm{col}_0(\mu_0,S^{-1}\varphi_0),G_{\mathcal R}\bigr)
\end{equation}
for $\mathcal R\in\mathcal R_0$, with $G_{\mathcal R}$ chosen by \ref{w2:see:D4}. Then $\lVert U_u\varphi\rVert=r_u\lVert\varphi\rVert$, and for finitely many data $(\varphi_{0,i},\mathcal R_i)$, $(t_k,\varphi_k,F_k)$ and coefficients $\alpha_i,\alpha'_k$ with
\[
\sum_i\alpha_i\,p_{\mathcal R_i}(\varphi_{0,i})+\sum_k\alpha'_k\,c(t_k)\,\mathrm{perCol}_{t_k}(\varphi_k,F_k)=0
\]
one has, for all $y$ and $v$,
\[
\sum_i\alpha_i\,\mathfrak b_u(U_u\varphi_{0,i},\mathcal R_i;y,v)+\sum_k\alpha'_k\,\mathfrak b_u(U_u\varphi_k,\mathcal R_{F_k};y,v)=0 .
\]
\end{corollary}

\begin{proof}
The norm statement follows from $\lVert S\psi\rVert=r\lVert\psi\rVert$ with $\psi=S^{-1}\varphi$, and $c(t)\ne0$ follows from \ref{w2:see:D6}, $\varsigma\ne0$ and \ref{w2:see:T1}. For each $k$ choose $G_k$ by \ref{w2:see:D3}. Apply Theorem~\ref{w2:see:main} to the family consisting of the tuples $\mu_0$ (indices $i$) and $t_k$ (indices $k$), the columns $\mathrm{col}_0(\mu_0,S^{-1}\varphi_{0,i})$ and $\mathrm{col}_0(t_k,S^{-1}\varphi_k)$, the elements $G_{\mathcal R_i}$ and $G_k$, and the coefficients $\alpha_i/c_{\mathrm{bot}}(\mu_0)$ and $\alpha'_k/c_{\mathrm{bot}}(t_k)$. Its hypothesis is the displayed relation, because by \ref{w2:see:D5} and \ref{w2:see:D2}
\[
\frac{c_{\mathrm{top}}(t_k)}{c_{\mathrm{bot}}(t_k)}\,\mathrm{tail}\bigl(\mathrm{up}_{n+1}\mathrm{col}_0(t_k,S^{-1}\varphi_k),G_k\bigr)=\frac{c_{\mathrm{top}}(t_k)\,\varsigma}{c_{\mathrm{bot}}(t_k)}\,\mathrm{perCol}_{t_k}(\varphi_k,F_k)=c(t_k)\,\mathrm{perCol}_{t_k}(\varphi_k,F_k),
\]
and the $i$-terms are $\alpha_ip_{\mathcal R_i}(\varphi_{0,i})$ by definition. Its conclusion, after the factors $c_{\mathrm{bot}}$ cancel, is $\sum_i\alpha_i\,\mathrm{per}^n_0(\mathrm{col}_0(\mu_0,S^{-1}\varphi_{0,i}),\mathcal R_i;y,v)+\sum_k\alpha'_k\,\mathrm{per}^n_0(\mathrm{col}_0(t_k,S^{-1}\varphi_k),\mathcal R_{F_k};y,v)=0$, using \ref{w2:see:D3} and \ref{w2:see:D4}, and \ref{w2:see:D1} turns this into the assertion.
\end{proof}

\begin{remark}[passage to Hypothesis~\ref{w2:H-see}]\label{w2:see:rem}
Take $X_u$ to consist of the pairs $(\varphi,F)\in V_{u,v}\times Y_v$ and the pairs $(\varphi_0,\mathcal R)$ with $\mathcal R\in\mathcal R_0$, with $\lambda_u(\varphi,F)=(B_{u,v}\varphi)(F)$ and $\lambda_u(\varphi_0,\mathcal R)=p_{\mathcal R}(\varphi_0)/c(\mu_0)$. Only the base tuple $t=\mu_0$ is needed. A relation among the values of $\lambda_u$ with complex $F$ is a relation with $K$-rational $F$ after expanding $F$ in a $K$-basis of the good part, because both sides are linear in $F$, and multiplying the relation by $c(\mu_0)$ turns it into the hypothesis of Corollary~\ref{w2:see:cor} with the same coefficients. The conclusion of the corollary is the conclusion of Hypothesis~\ref{w2:H-see}, once the local Bessel period $\mathfrak b_u(U_u\varphi,\mathcal R;y,v)$ is identified with $\mathrm{per}(T_x(\xi),R_{(y,x_u)}R_z\overline{\theta_x})$ through the pinning relation between $b$, $z$ and $x_u$ and the row dictionary $\mathcal R_F\leftrightarrow\mathbf s(F)$ of Hypothesis~\ref{w2:H-dict}. Lemma~\ref{w2:prop} uses only relations with one or two terms, so Hypothesis~\ref{w2:H-see} may be weakened to such relations without changing Theorem~\ref{w2:Bproof}. The anchor of Lemma~\ref{w2:anchor} may be any $x_0=(\varphi_0,\mathcal R)$ with $\mu_u(x_0)\ne0$. The exact values of $c_{\mathrm{top}}$, $c_{\mathrm{bot}}$ and $\varsigma$ do not enter Theorem~\ref{w2:B}, because the number $a_u=r_u/\kappa_u$ absorbs them. Only $c(\mu_0)\ne0$ is used there, while the unit property of $c(t)/c(\mu_0)$ at the primes of the chain is used for the integral structure of the period maps.
\end{remark}

\begin{lemma}[rationality]\label{w2:see:rat}
Assume that $B_u(\varphi)(y)=\mathrm{Pt}_{\mu_0}(\varphi,q_y)$ for $\varphi\in V_u$ and $y\in Y$, where $\mathrm{Pt}_{\mu_0}$ is the $K$-bilinear period of the chain on the definite side at the base tuple and $q_y$ is the $K$-rational vector of the good part with $\iota_0(q_y)=y$. Assume also $\mathrm{perCol}_{\mu_0}(\varphi\otimes1,q\otimes1)=\sigma_v(\mathrm{Pt}_{\mu_0}(\varphi,q))$ and $(B_{u,v}\varphi)(\iota_0q)=\mathrm{perCol}_{\mu_0}(\varphi,q)$. Then Hypothesis~\ref{w2:H-rat} holds.
\end{lemma}

\begin{proof}
For $\varphi\in V_u$ and $y\in Y$, $(B_{u,v}(\varphi\otimes1))(y\otimes1)=\mathrm{perCol}_{\mu_0}(\varphi\otimes1,q_y\otimes1)=\sigma_v(\mathrm{Pt}_{\mu_0}(\varphi,q_y))=\sigma_v(B_u(\varphi)(y))$, which lies in $\sigma_v(K)$.
\end{proof}

\paragraph{Rationality on the selected record.}\label{rat:selected}
We give the descent producing this identity in the original argument.
Fix one selected record, its coefficient field $K$, its embedding
$\sigma:K\hookrightarrow\mathbb C$, its base tuple $\mu_0$ and a word $u$.
Write $N_{u,K}$ and $R_K$ for its rational column and row spaces and
$j_N,j_R$ for their $\sigma$-semilinear maps to the complex spaces.
The images of the chosen bases are complex bases. Let $B^K_u$
be its rational bilinear pairing and $\mathcal B_u$ its complex pairing.
For these data the assertion is
\begin{equation}\label{rat:base}
 \mathcal B_u(j_N\varphi,j_R y)=\sigma\big(B^K_u(\varphi,y)\big)
 \qquad(\varphi\in N_{u,K},\ y\in R_K).
\end{equation}

Choose the rational column coordinates of $\varphi$. On the row side
let $e_R:K^d\longrightarrow R_K$ be the coordinate isomorphism and let
$p:K^n\to K^d$ be the composite of the good-part projection and the
row-coordinate map. It is onto by surjectivity of the first projection
and bijectivity on the good part. Let $A(F)$ be the corresponding
complex level row. Its values span the degree-zero space $V^0$,
since the level operator has precisely this image and fixes its image.
Put $\beta(y')=B^K_u(\varphi,e_Ry')$. The row comparison on this
realization gives
\[
 \iota A(F)=j_R e_R(pF),\qquad
 \mathrm{perCol}_{\mu_0}(j_N\varphi,A(F))=\sigma(\beta(pF)).
\]
For the second equality, the comparison of a rational definite period
with the complex column gives $\sigma(\mathrm{Pt}_{\mu_0})$ on the left.
The definite period factors as $\zeta(u,\mu_0)$ times the pairing on the
projected row, and $\zeta(u,\mu_0)=1$. Transport by the rational column
and row isomorphisms gives $\beta(pF)$. At other tuples the source
specifies $\zeta(u,t)\ne0$ and valuation zero at the designated places.
Identifying this scalar with $c(t)/c(\mu_0)$ requires a separate
comparison of period normalizations.

Extend $\beta e_R^{-1}$ to the complex linear functional $\Lambda$
by prescribing its $\sigma$-embedded values on the chosen row basis.
The last two identities show that $\mathrm{perCol}_{\mu_0}(j_N\varphi,-)$
and $\Lambda\iota$ agree on every $A(F)$, hence on $V^0$. Put
$U=V^0\cap\ker\iota$ and $Q=V^0\cap U^\perp$. On $Q$, $\iota$ is
injective. The section at $\mu_0$ first projects onto $\iota(Q)$ and
then inverts $\iota|_Q$. Given $y=e_R y'$, take $F$ with $pF=y'$.
Now $q=A(F)-\operatorname{proj}_U A(F)$ belongs to $Q$ and
$\iota q=j_R y$. Thus the section takes $j_Ry$ to $q$.
The defining tie for $\mathcal B_u$ evaluates the period on this
section. Its value is $\Lambda(j_R y)=\sigma(B^K_u(\varphi,y))$.
This proves \eqref{rat:base}. It also proves vanishing on $U$ for the
columns just treated. For a general column, extending the good-part
identity to all of $V^0$ requires vanishing on $U$, by the same
orthogonal decomposition.

In a selected family this argument holds at every index $i$ once $x$
exceeds a common threshold. Take the maximum of the threshold for
the row facts, the selection threshold, and the initial scale plus one.
The realization supplies the assertion that all tuples are used, which
allows the definite-period comparison to be applied to these rows.
For condition (d) of Theorem~\ref{r3:A3} one place suffices for the
whole family. Take the place of $\sigma$ and the same chosen record for
all words. Use its rational pairing for $B_u$. For bases $(b_i)$ of
$N_{u,K}$ and $(c_l)$ of $R_K$, equation~\eqref{rat:base} gives the
matrix of the analytic pairing as
$(\sigma(B_u(b_i)(c_l)))_{il}$. Bilinearity now identifies the two
maps on the full complex spaces and gives Hypothesis~\ref{w2:H-rat}
for this realization at the place required by (d). If the seesaw
calculation uses different coordinates, transport its columns and rows
by their respective basis isomorphisms. In common coordinates the
remaining comparison is equality of the two resulting matrices, or
equivalently equality on every pair $(b_i,c_l)$. The height comparison
in Section~\ref{r4:sec:R4} still uses the metrics at all infinite places.



\subsubsection{The double-step seesaw identity}\label{w3s:sec}

This subsection proves the double-step identity (T2) of Hypothesis~\ref{w2:see:H-laws}. The identity is an iterated integral over two quotients of finite adelic groups. We show that the integral collapses to an integral over the finitely many places of the tuple, with a volume factor $\kappa_j$, and we identify the result with the bilift of the reading functions through one kernel $C$ that does not depend on the tuple. The analytic unfolding over the archimedean and automorphic variables is not part of (T2). It is item (T6), where rapid decay of cusp forms enters.

\paragraph{Quotient integrals.} Let $\mathbf G$ be a second countable locally compact group with a left Haar measure $\mu$, and let $R\subset\mathbf G$ be a closed subgroup with a left Haar measure $\rho$. We say that a function $f\colon\mathbf G\to\C$ has the \emph{absolutely convergent quotient integral} $z$ over $R\backslash\mathbf G$, and write $\int^{\mathrm{abs}}_{R\backslash\mathbf G}f=z$, if there is $F\in L^1(\mathbf G,\mu)$ such that for every $g\in\mathbf G$ the function $r\mapsto F(rg)$ is $\rho$-integrable, $f(g)=\int_RF(rg)\,d\rho(r)$, and $z=\int_{\mathbf G}F\,d\mu$. For two such pairs $(\mathbf G_1,R_1)$, $(\mathbf G_2,R_2)$ and a function $I$ on $\mathbf G_1\times\mathbf G_2$, we say that $I$ has the \emph{absolutely convergent iterated quotient integral} $z$ if there is $f_1\colon\mathbf G_1\to\C$ with $\int^{\mathrm{abs}}_{R_1\backslash\mathbf G_1}f_1=z$ and $\int^{\mathrm{abs}}_{R_2\backslash\mathbf G_2}I(q_1,\cdot)=f_1(q_1)$ for every $q_1\in\mathbf G_1$. These are the formal notions \texttt{IsQuotIntL1} and \texttt{Iter2}.

\begin{lemma}[uniqueness]\label{w3s:lem:uniq}
Assume that $\mathbf G$ and $R$ are unimodular. If $\int^{\mathrm{abs}}_{R\backslash\mathbf G}f=z$ and $\int^{\mathrm{abs}}_{R\backslash\mathbf G}f=z'$, then $z=z'$. The value is linear in $f$, it vanishes for $f=0$, and $\int^{\mathrm{abs}}_{R\backslash\mathbf G}R_bf=\int^{\mathrm{abs}}_{R\backslash\mathbf G}f$ for every $b\in\mathbf G$, where $R_bf(g)=f(gb)$.
\end{lemma}

\begin{proof}
Since the modular functions of $\mathbf G$ and $R$ are both trivial, there is a right $\mathbf G$-invariant Radon measure $\dot\mu$ on $R\backslash\mathbf G$ with $\int_{\mathbf G}F\,d\mu=\int_{R\backslash\mathbf G}\int_RF(rg)\,d\rho(r)\,d\dot\mu(g)$ for every $F\in L^1(\mathbf G)$ \cite[\S2.6]{w2folland}. Hence $z=\int_{R\backslash\mathbf G}f\,d\dot\mu$ depends only on $f$, and linearity follows. If $F$ is a witness for $f$, then $R_bF$ is a witness for $R_bf$, and $\int R_bF\,d\mu=\int F\,d\mu$ because $\mu$ is right invariant.
\end{proof}

\begin{lemma}[collapse to a factor]\label{w3s:lem:collapse}
Let $\mathbf G$ be a second countable locally compact unimodular group, and let $G_1,G_2\subset\mathbf G$ be closed unimodular subgroups whose elements commute with each other and for which $(a,b)\mapsto ab$ is a homeomorphism $G_1\times G_2\to\mathbf G$. Let $R\subset\mathbf G$ be a closed unimodular subgroup such that $ab\in R$ with $a\in G_1$, $b\in G_2$ implies $a,b\in R$. Put $R_i=R\cap G_i$. Let $K_i\subset G_i$ be compact open subgroups and $K=K_1K_2$. Let $\mu,\rho,\mu_1,\rho_1$ be Haar measures on $\mathbf G,R,G_1,R_1$ with $\mu_1(K_1)=1$ and $\rho_1(R_1\cap K_1)=1$. Let $f\colon\mathbf G\to\C$ satisfy
\begin{enumerate}[label=(\roman*),leftmargin=*]
\item $f(rg)=f(g)$ for $r\in R$ and $g\in\mathbf G$;
\item $f(gk)=f(g)$ for $k\in K_2$ and $g\in\mathbf G$;
\item $f(g)\ne0$ implies $g\in RG_1K_2$.
\end{enumerate}
If $\int^{\mathrm{abs}}_{R_1\backslash G_1}f|_{G_1}=z$, then
\[
\int^{\mathrm{abs}}_{R\backslash\mathbf G}f=\frac{\mu(K)}{\rho(R\cap K)}\,z .
\]
\end{lemma}

\begin{proof}
The product map is a continuous homomorphism because $G_1$ and $G_2$ commute, so it is an isomorphism of topological groups $G_1\times G_2\cong\mathbf G$. By the hypothesis on $R$, it restricts to an isomorphism $R_1\times R_2\cong R$, and $R_1,R_2$ are closed. The modular function of $R_1\times R_2$ is the product of those of $R_1$ and $R_2$, so $R_1$ and $R_2$ are unimodular. By uniqueness of Haar measure there are Haar measures $\mu_2$ on $G_2$ and $\rho_2$ on $R_2$ with $\mu=\mu_1\otimes\mu_2$ and $\rho=\rho_1\otimes\rho_2$ under these isomorphisms. The group $R_2\cap K_2$ is compact and open in $R_2$, so $0<\rho_2(R_2\cap K_2)<\infty$, and likewise $0<\mu_2(K_2)<\infty$.

\emph{Step 1.} We show $f(ab)=f(a)\,\mathbf 1_{R_2K_2}(b)$ for $a\in G_1$ and $b\in G_2$. If $b=r_2k$ with $r_2\in R_2$ and $k\in K_2$, then $ab=r_2ak$, so (ii) and (i) give $f(ab)=f(r_2a)=f(a)$. If $f(ab)\ne0$, then (iii) gives $ab=r_1r_2xk$ with $r_i\in R_i$, $x\in G_1$, $k\in K_2$, that is $ab=(r_1x)(r_2k)$, and uniqueness of the decomposition gives $b=r_2k\in R_2K_2$.

\emph{Step 2.} Let $F_1\in L^1(G_1)$ be a witness for $f|_{G_1}$, and put $F(ab)=F_1(a)\,\mathbf 1_{K_2}(b)/\rho_2(R_2\cap K_2)$. Then $F\in L^1(\mathbf G)$ and $\int_{\mathbf G}F\,d\mu=z\,\mu_2(K_2)/\rho_2(R_2\cap K_2)$. For $g=ab$ and $r=r_1r_2$ we have $F(rg)=F_1(r_1a)\mathbf 1_{K_2}(r_2b)/\rho_2(R_2\cap K_2)$. The second factor is the indicator of the set $R_2\cap K_2b^{-1}$, of measure at most $\rho_2(R_2\cap K_2)$, so $r\mapsto F(rg)$ is $\rho$-integrable and Fubini's theorem gives
\[
\int_RF(rg)\,d\rho(r)=f(a)\,\frac{\rho_2(R_2\cap K_2b^{-1})}{\rho_2(R_2\cap K_2)} .
\]
If $b=r_0k\in R_2K_2$, then $R_2\cap K_2b^{-1}=(R_2\cap K_2)r_0^{-1}$, which has measure $\rho_2(R_2\cap K_2)$ because $\rho_2$ is right invariant. If $R_2\cap K_2b^{-1}$ is nonempty, then $b\in R_2K_2$. So the right side equals $f(a)\mathbf 1_{R_2K_2}(b)$, which is $f(ab)$ by Step 1. Hence $\int^{\mathrm{abs}}_{R\backslash\mathbf G}f=z\,\mu_2(K_2)/\rho_2(R_2\cap K_2)$.

\emph{Step 3.} By the hypothesis on $R$, $R\cap K=(R_1\cap K_1)(R_2\cap K_2)$. Hence $\mu(K)=\mu_1(K_1)\mu_2(K_2)=\mu_2(K_2)$ and $\rho(R\cap K)=\rho_1(R_1\cap K_1)\rho_2(R_2\cap K_2)=\rho_2(R_2\cap K_2)$.
\end{proof}

\paragraph{The double-step integrand.} We use the notation of Subsection~\ref{w2:sec:seesaw}. For $0\le i<m$ the pair $(\mathrm U(V_i),\mathrm U(V_{i+1}))$ acts on functions on $W_i(\A)$, $W_i=\Hom(V_{i+1},V_i)$, by the Weil representation $\omega_i$ of the tuple attached to $\psi$ and to the splitting characters of the tuple. The pairs $(\mathrm U(X_i),\mathrm U(X_{i+1}))$ of definite spaces act on functions on $W'_i(\A)$ in the same way, by $\omega'_i$. For a finitely supported function $e$ on $W_i(\Af)$ put $\mathrm{comb}(e)\Phi=\sum_ue(u)\,\mathrm{heis}(u)\Phi$, where $\mathrm{heis}(u)$ is the Heisenberg translation by $u$ with the character $\psi$. We say that $e$ is \emph{supported at} a set $P$ of places if every $u$ in its support has trivial components outside $P$. For self-dual lattices $N_i\subset V_i$ and a set $P$ of finite places, a theta function $\Phi$ on $W_i(\A)$ is \emph{vacuum at} $P$ if it is invariant under the Heisenberg action of the lattice $\Hom(N_{i+1},N_i)\otimes\prod_{p\in P}\cO_p$, so that its component at $p\in P$ is a multiple of the characteristic function of that lattice. For finitely supported functions $e$ on $W'_i(\Af)$ and $e'$ on $W'_{i+1}(\Af)$ and a function $C$ on $W'_i(\A)\times W'_{i+1}(\A)$ write $\mathrm{biComb}(e,e')C(u,v)=\mathrm{comb}(e)\bigl(u'\mapsto\mathrm{comb}(e')C(u',\cdot)(v)\bigr)(u)$.

Fix $j$ and $\sigma$ with $j+2(\sigma+1)+2=m$. The Bessel datum of rank $\sigma+1$ at level $j$, the Fourier--Jacobi datum of rank $\sigma+1$ at level $j+1$ and the Bessel datum of rank $\sigma$ at level $j+2$ give closed subgroups of $\mathrm U(V_{j+1})(\A)$ and $\mathrm U(V_{j+2})(\A)$. Their intersections with the finite adelic groups $\mathbf G_{j+1}=\mathrm U(V_{j+1})(\Af)$ and $\mathbf G_{j+2}=\mathrm U(V_{j+2})(\Af)$ are denoted $R_J$ and $R_B$, with Haar measures $\rho_J$, $\rho_B$, and $\mu_{j+1},\mu_{j+2}$ are the Haar measures of $\mathbf G_{j+1},\mathbf G_{j+2}$ fixed by the construction. For theta functions $\Phi_1$ on $W_j(\A)$ and $\Phi_2$ on $W_{j+1}(\A)$, a function $\phi$ on $\mathrm U(X_j)(\A)$ and points $g,h_1,h_2$ of $\mathrm U(V_j)(\A)$, $\mathrm U(V_{j+1})(\A)$, $\mathrm U(V_{j+2})(\A)$, let
\[
\mathbf H(\Phi_1,\Phi_2;\phi;g,h_1,h_2)=\mathrm{hand}_{\mathrm{FB}}\Bigl(\mathrm{hand}_{\mathrm{BF}}\bigl(\phi,\mathrm{flat}_{\mathrm{BF}}(\omega_j(g,h_1)\Phi_1)\bigr),\ \mathrm{flat}_{\mathrm{FB}}(\omega_{j+1}(h_1,h_2)\Phi_2)\Bigr)
\]
be the function on $\mathrm U(X_{j+2})(\A)$ obtained by the Bessel-to-Fourier--Jacobi step at level $j$ followed by the Fourier--Jacobi-to-Bessel step at level $j+1$ (formal \texttt{Chain.hand}). Here $\mathrm{flat}_{\mathrm{BF}}$ and $\mathrm{flat}_{\mathrm{FB}}$ are the partial integrations of a theta function along the unipotent directions of the two flags, taken with Haar measures of covolume one, $\mathrm{hand}_{\mathrm{BF}}$ integrates $\phi$ over $[\mathrm U(X_j)]$ against the restricted kernel, and $\mathrm{hand}_{\mathrm{FB}}$ contracts the result with the second kernel. The map $\mathbf H$ is bilinear in $(\Phi_1,\Phi_2)$.

The contracted kernel of two theta functions $\Psi_1,\Psi_2$ is the function $\mathrm{ctr}(\Psi_1,\Psi_2)$ on $W'_j(\A)\times W'_{j+1}(\A)$ given by the integral over the link variable of the product $\mathrm{flat}_{\mathrm{FB}}\Psi_2\cdot\mathrm{flat}_{\mathrm{BF}}\Psi_1$ (formal \texttt{Comm.ctr}). For a function $C$ on $W'_j(\A)\times W'_{j+1}(\A)$ and a function $\phi$ on $\mathrm U(X_j)(\A)$ the bilift is
\[
\mathrm{bl}_j(\phi,C)(h_2)=\int_{[\mathrm U(X_{j+1})]}\int_{[\mathrm U(X_j)]}\phi(h_0)\,\Bigl(\omega'_{j+1}(h_1,h_2)\bigl(v\mapsto\omega'_j(h_0,h_1)C(\cdot,v)(0)\bigr)\Bigr)(0)\,dh_0\,dh_1,
\]
and for a family $\phi=(\phi_\alpha)_{\alpha\in\mathrm{Ch}_j}$ the double kernel lift of Hypothesis~\ref{w2:see:H-laws} is $\mathrm{two}_j(\phi,C)_\beta=\sum_{\alpha,\gamma}\mathrm{bl}_j\bigl(\phi_\alpha,\mathrm{biComb}(\mathrm{dd}_j(\alpha,\gamma),\mathrm{dd}_{j+1}(\gamma,\beta))C\bigr)$ for $\beta\in\mathrm{Ch}_{j+2}$, where $\mathrm{dd}_i$ are the kernel combinations of the tuple. A \emph{bivacuum kernel} is a continuous $C$ that is a theta function vacuum at every finite place in each variable, and it is \emph{real} if it is moreover smooth in the archimedean directions of both variables.

For a tuple $s$, archimedean points $a\in\mathrm U(V_{j+1})_\infty$, $a'\in\mathrm U(V_{j+2})_\infty$, polynomial Fock vectors $P_1,P_2$, a column $\psi'\in\mathrm{Ram}_j$, a frame $y$, $G\in\mathcal C_{j+2}$ and $v\in\mathcal S_{j+2}$, the \emph{double-step integrand} is the function on $\mathbf G_{j+1}\times\mathbf G_{j+2}$
\begin{equation}\label{w3s:integrand}
I_s(q_1,q_2)=\sum_{\alpha\in\mathrm{Ch}_j}\sum_ge_{j,\alpha}(g)\,\mathcal B^\sigma_{j+2}\Bigl(\mathbf H\bigl(\mathrm{pt}_j(P_1),\mathrm{pt}_{j+1}(P_2);\mathrm{rd}_j(\alpha,\psi',y);g,aq_1,a'q_2\bigr),\ R_{a'q_2}G(v)\Bigr),
\end{equation}
where $\mathrm{pt}_i(P)$ is the theta function of the tuple at level $i$ attached to $P$ (formal \texttt{Tup.integrand}). It is the integrand of $\mathcal B^{\sigma+1}_j$ on the double lift after the two unfoldings, as a function of the two finite adelic variables.

\begin{hypothesis}[general laws of the Weil representation]\label{w3s:hyp:gen}
\begin{enumerate}[label=(W\arabic*),leftmargin=*]
\item\label{w3s:W1} (Disjoint places.) Heisenberg combinations supported at disjoint sets of places commute. A combination supported at $P$ maps theta functions vacuum at a set $P'$ disjoint from $P$ to theta functions vacuum at $P'$, and it commutes with $\omega_i(g,h)$ for $g,h$ with trivial components at $P$. A theta function is vacuum at $P\cup P'$ if and only if it is vacuum at $P$ and at $P'$.
\item\label{w3s:W2} (Normalized steps.) The maps $\mathrm{flat}_{\mathrm{BF}}$ and $\mathrm{flat}_{\mathrm{FB}}$ on theta functions depend only on the Bessel and Fourier--Jacobi data and not on the remaining data of the step of the tuple. Consequently $\mathrm{ctr}(\Psi_1,\Psi_2)$ does not depend on the tuple.
\item\label{w3s:W3} (Representation.) $\omega_i$ is a representation of $\mathrm U(V_i)(\A)\times\mathrm U(V_{i+1})(\A)$, and archimedean and finite adelic points commute.
\item\label{w3s:W4} (Bi-action.) For theta data on $(X_j,X_{j+1})$ and $(X_{j+1},X_{j+2})$ and a bikernel $C$ that is a theta function in each variable, the two actions $\omega'_j$, $\omega'_{j+1}$ on the two variables of $C$ commute and preserve this property (formal \texttt{BiAct}, Lean theorem \texttt{seeSawWeilBiAction\_commutes}).
\item\label{w3s:W5} (Flat maps.) The flat maps send theta functions to theta functions and their defining integrals converge absolutely (formal \texttt{FlatLaw}, Lean theorem \texttt{flatLaw}).
\item\label{w3s:W6} (Seesaw exchange.) For $\phi$ continuous and left $\mathrm U(X_j)(F)$-invariant and theta functions $\Xi_1,\Xi_2$ on the reduced spaces, $\mathrm{hand}_{\mathrm{FB}}\bigl(\mathrm{hand}_{\mathrm{BF}}(\phi,\Xi_1),\Xi_2\bigr)=\mathrm{bl}_j\bigl(\phi,\mathrm{contr}(\Xi_1,\Xi_2)\bigr)$, where $\mathrm{contr}$ is the contraction over the link variable, so that $\mathrm{ctr}(\Psi_1,\Psi_2)=\mathrm{contr}(\mathrm{flat}_{\mathrm{BF}}\Psi_1,\mathrm{flat}_{\mathrm{FB}}\Psi_2)$
(formal \texttt{ContrLaw}, Lean theorem \texttt{contrLaw}). Its proof is Fubini's theorem for a bounded continuous integrand on the product of a compact fundamental domain of $[\mathrm U(X_j)]$ and a fundamental domain of covolume one of the link lattice, followed by the factorization of the Weil representation of an orthogonal sum $V^0=A\oplus L$ as the tensor product of the Weil representations of $A$ and $L$, with splitting characters multiplying (\texttt{weil\_datum\_product\_compatibility}).
\item\label{w3s:W7} (Linearity.) $\mathcal B^\sigma_{j+2}(c\phi',c'f)=cc'\,\mathcal B^\sigma_{j+2}(\phi',f)$ and $\mathbf H(c\Phi_1,c'\Phi_2;\dots)=cc'\,\mathbf H(\Phi_1,\Phi_2;\dots)$.
\end{enumerate}
\end{hypothesis}

\begin{hypothesis}[local laws of the tuple]\label{w3s:hyp:loc}
There are data common to all tuples, namely finite disjoint sets $g\mathrm P$ and $t\mathrm P$ of finite places, self-dual lattices $N_i\subset V_i$, base points $b_i\in\mathbf G_i$, compact open subgroups $J_i$ of $\mathrm U(V_i)(F_{g\mathrm P})$, theta functions $\Psi_i(P;a,a')$ on $W_i(\A)$ vacuum at every finite place, combinations $e^G_i$, $d^G_i$ supported at $g\mathrm P$, combinations $d^T_i$ and an operator $\mathrm{op}_j=\sum_k\mathrm{biComb}(u_k,u'_k)$ supported at $t\mathrm P$, a number $c_G(j)$, a set $\mathrm{RamAd}_j$ of continuous automorphic functions on $\mathrm U(X_j)(\A)$ with a map $\mathrm{Sad}_j$ of it into itself. Put $K_i=\{\kappa k:\kappa\text{ stabilizes }N_i\text{ away from }g\mathrm P,\ k\in J_i\}$ and
\[
\kappa_j=\frac{\mu_{j+1}(K_{j+1})}{\rho_J(R_J\cap K_{j+1})}\cdot\frac{\mu_{j+2}(K_{j+2})}{\rho_B(R_B\cap K_{j+2})} .
\]
For each tuple $s$ there are a finite set $Q=Q_s$ of finite places disjoint from $g\mathrm P\cup t\mathrm P$, numbers $\mathrm{sc}_i$, $c_Q(j)$, $c_T(j)$, combinations $e^Q_i$, $d^Q_i(\alpha,\gamma)$ supported at $Q$, and finitely supported functions $\varepsilon_{i,\alpha}$ on $\mathrm U(V_i)(F_Q)$, such that the following hold.
\begin{enumerate}[label=(L\arabic*),leftmargin=*]
\item\label{w3s:L1} (Vacuum form.) $\omega_i(ab_i,a'b_{i+1})\,\mathrm{pt}_i(P)=\mathrm{sc}_i\,\mathrm{comb}(e^G_i)\,\mathrm{comb}(e^Q_i)\,\Psi_i(P;a,a')$ for archimedean $a,a'$. Moreover $e_{i,\alpha}(g)=\varepsilon_{i,\alpha}(gb_i^{-1})$, the combination $\mathrm{dd}_i(\alpha,\gamma)$ acts on theta functions as $\mathrm{comb}(d^Q_i(\alpha,\gamma))\,\mathrm{comb}(d^G_i)\,\mathrm{comb}(d^T_i)$, for $\psi'\in\mathrm{Ram}_j$ one has $\mathrm{rd}_j(\alpha,\psi',y)\in\mathrm{RamAd}_j\cap\mathrm{FX}_{j,\alpha}$ and $\mathrm{rd}_j(\alpha,S_j\psi',y)=\mathrm{Sad}_j(\mathrm{rd}_j(\alpha,\psi',y))$, and for $G\in\mathcal C_i$ and archimedean $a$ the function $R_{ab_i}G(v)$ is automorphic, lies in $\mathrm{FY}_i$ and is invariant under the stabilizer of $N_i$ at the places outside $g\mathrm P\cup Q$. Here $\mathrm{FX}_{j,\alpha}$ and $\mathrm{FY}_i$ are sets of functions given with the tuple (formal \texttt{VTup.FX}, \texttt{VTup.FY}).
\item\label{w3s:L2} (Cells.) Let $\Psi_1,\Psi_2$ be theta functions vacuum at every finite place, put $\Phi^0_i=\mathrm{comb}(e^G_{j+i-1})\mathrm{comb}(e^Q_{j+i-1})\Psi_i$, let $\phi_\alpha\in\mathrm{RamAd}_j$, let $f=R_{a'b_{j+2}}G(v)$ and
\[
I^0(q_1,q_2)=\sum_\alpha\sum_g\varepsilon_{j,\alpha}(g)\,\mathcal B^\sigma_{j+2}\bigl(\mathbf H(\Phi^0_1,\Phi^0_2;\phi_\alpha;g,q_1,q_2),R_{q_2}f\bigr).
\]
Write $\mathbf G_{i,Q}$ for the points of $\mathrm U(V_i)$ at $Q$ and $K_i=K_{i,Q}K'_i$ with $K_{i,Q}$ the stabilizer of $N_i$ at $Q$ and $K'_i$ the part of $K_i$ away from $Q$. Then $I^0$ is left invariant under $R_J$ in $q_1$ and under $R_B$ in $q_2$, right invariant under $K'_{j+1}$ in $q_1$ and under $K'_{j+2}$ in $q_2$, and $I^0(q_1,q_2)\ne0$ implies $q_1\in R_J\mathbf G_{j+1,Q}K'_{j+1}$, and also $q_2\in R_B\mathbf G_{j+2,Q}K'_{j+2}$ when $q_1\in\mathbf G_{j+1,Q}$.
\item\label{w3s:L3} (Twist places.) Let $\rho_{J,Q},\mu_{j+1,Q},\rho_{B,Q},\mu_{j+2,Q}$ be Haar measures on $R_J\cap\mathbf G_{j+1,Q}$, $\mathbf G_{j+1,Q}$, $R_B\cap\mathbf G_{j+2,Q}$, $\mathbf G_{j+2,Q}$ giving $K_{j+1,Q}$, $K_{j+2,Q}$ and their intersections with $R_J$, $R_B$ the volume one. For theta functions $\Psi'_1,\Psi'_2$ vacuum at $Q$, continuous automorphic $\phi_\alpha\in\mathrm{FX}_{j,\alpha}$ and automorphic $f\in\mathrm{FY}_{j+2}$, the function
\[
(q_1,q_2)\mapsto\sum_\alpha\sum_g\varepsilon_{j,\alpha}(g)\,\mathcal B^\sigma_{j+2}\bigl(\mathbf H(\mathrm{comb}(e^Q_j)\Psi'_1,\mathrm{comb}(e^Q_{j+1})\Psi'_2;\phi_\alpha;g,q_1,q_2),R_{q_2}f\bigr)
\]
on $\mathbf G_{j+1,Q}\times\mathbf G_{j+2,Q}$ has the absolutely convergent iterated quotient integral
\[
c_Q(j)\sum_\beta\mathcal B^\sigma_{j+2}\Bigl(\sum_{\alpha,\gamma}\mathrm{bl}_j\bigl(\phi_\alpha,\mathrm{biComb}(d^Q_j(\alpha,\gamma),d^Q_{j+1}(\gamma,\beta))\,\mathrm{ctr}(\Psi'_1,\Psi'_2)\bigr),\ \mathrm{tr}(\varepsilon_{j+2,\beta})f\Bigr).
\]
\item\label{w3s:L4} (Level places.) For theta functions $\Psi_1,\Psi_2$ vacuum at $g\mathrm P$,
$\mathrm{ctr}(\mathrm{comb}(e^G_j)\Psi_1,\mathrm{comb}(e^G_{j+1})\Psi_2)=c_G(j)\,\mathrm{biComb}(d^G_j,d^G_{j+1})\,\mathrm{ctr}(\Psi_1,\Psi_2)$.
\item\label{w3s:L5} (Places of $t\mathrm P$.) For theta functions $\Psi_1,\Psi_2$ vacuum at $t\mathrm P$, combinations $d,d'$ supported away from $t\mathrm P$ and $\phi\in\mathrm{RamAd}_j$,
\begin{multline*}
\mathrm{bl}_j\bigl(\phi,\mathrm{biComb}(d,d')\,\mathrm{ctr}(\Psi_1,\Psi_2)\bigr)\\
=c_T(j)\,\mathrm{bl}_j\bigl(\mathrm{Sad}_j\phi,\ \mathrm{biComb}(d,d')\,\mathrm{biComb}(d^T_j,d^T_{j+1})\,\mathrm{op}_j\,\mathrm{ctr}(\Psi_1,\Psi_2)\bigr),
\end{multline*}
and $\mathrm{op}_j\,\mathrm{ctr}(\Psi_1,\Psi_2)$ is a real bivacuum kernel when $\Psi_1,\Psi_2$ are vacuum at every finite place.
\item\label{w3s:L6} (Normalization.) $c_V(j)=\kappa_j\,\mathrm{sc}_j\,\mathrm{sc}_{j+1}\,c_Q(j)\,c_G(j)\,c_T(j)$.
\end{enumerate}
\end{hypothesis}

\begin{theorem}[the double-step identity (T2)]\label{w3s:thm:T2}
Assume Hypotheses~\ref{w3s:hyp:gen} and~\ref{w3s:hyp:loc}. Let $j+2(\sigma+1)+2=m$, let $P_1,P_2$ be polynomial Fock vectors and $a\in\mathrm U(V_{j+1})_\infty$, $a'\in\mathrm U(V_{j+2})_\infty$. Put $\Psi_1=\Psi_j(P_1;1,a)$, $\Psi_2=\Psi_{j+1}(P_2;a,a')$ and $C=\mathrm{op}_j\,\mathrm{ctr}(\Psi_1,\Psi_2)$. Then $C$ is a real bivacuum kernel that depends only on $(j,\sigma,P_1,P_2,a,a')$ and not on the tuple. For every tuple $s$, every $\psi'\in\mathrm{Ram}_j$, $y$, $G\in\mathcal C_{j+2}$ and $v$, the integrand \eqref{w3s:integrand} has an absolutely convergent iterated quotient integral over $R_J\backslash\mathbf G_{j+1}$ and then $R_B\backslash\mathbf G_{j+2}$, equal to
\begin{equation}\label{w3s:value}
c_V(j)\sum_{\beta\in\mathrm{Ch}_{j+2}}\mathcal B^\sigma_{j+2}\Bigl(\mathrm{two}_j\bigl(\mathrm{rd}_j(\cdot,S_j\psi',y),C\bigr)_\beta,\ \mathrm{tr}(e_{j+2,\beta})\,R_{a'}G(v)\Bigr),
\end{equation}
with $c_V(j)=c_V(j)(s)=\kappa_j\,\mathrm{sc}_j(s)\,\mathrm{sc}_{j+1}(s)\,c_Q(j)(s)\,c_G(j)\,c_T(j)(s)$ and $\kappa_j>0$ independent of $s$.
\end{theorem}

\begin{proof}
\emph{Step 1, the kernel.} The functions $\Psi_1,\Psi_2$ belong to the common data, and by \ref{w3s:W2} so does $\mathrm{ctr}(\Psi_1,\Psi_2)$. Hence $C$ depends only on $(j,\sigma,P_1,P_2,a,a')$. It is a real bivacuum kernel by \ref{w3s:L5}. The groups $K_{j+1}$ and $K_{j+2}$ are compact and open, and their intersections with $R_J$ and $R_B$ are compact open subgroups of $R_J$ and $R_B$, so the four volumes in $\kappa_j$ are positive and finite and $\kappa_j>0$. It does not involve the tuple.

\emph{Step 2, the twist places.} Fix $s$, $\psi'$, $y$, $G$, $v$, and put $\phi_\alpha=\mathrm{rd}_j(\alpha,\psi',y)$, which lies in $\mathrm{RamAd}_j\cap\mathrm{FX}_{j,\alpha}$ by \ref{w3s:L1}, and $f=R_{a'b_{j+2}}G(v)$, which is automorphic and lies in $\mathrm{FY}_{j+2}$ by \ref{w3s:L1}. The functions $\Psi_i$ are vacuum at every finite place, so by \ref{w3s:W1} the functions $\Psi'_i=\mathrm{comb}(e^G_{j+i-1})\Psi_i$ are vacuum at $Q$, because $Q$ and $g\mathrm P$ are disjoint, and $\mathrm{comb}(e^Q)\Psi'_i=\Phi^0_i$ in the notation of \ref{w3s:L2}, because combinations at disjoint places commute. By \ref{w3s:L3} the restriction of the function $I^0$ of \ref{w3s:L2} to $\mathbf G_{j+1,Q}\times\mathbf G_{j+2,Q}$ has an absolutely convergent iterated quotient integral $z_Q$ for the measures normalized at $Q$, with $z_Q$ the value displayed in \ref{w3s:L3} for $\mathrm{ctr}(\Psi'_1,\Psi'_2)$. Let $f_Q$ be the corresponding inner function on $\mathbf G_{j+1,Q}$.

\emph{Step 3, collapse of the inner integral.} Apply Lemma~\ref{w3s:lem:collapse} to $\mathbf G=\mathbf G_{j+2}$, $G_1=\mathbf G_{j+2,Q}$, $G_2$ the points away from $Q$, $R=R_B$, $K_1=K_{j+2,Q}$, $K_2=K'_{j+2}$. The groups involved are the finite adelic points of unitary groups, which are unimodular, and the Bessel and Fourier--Jacobi subgroups, which are semidirect products of a unipotent group by a unitary group and are unimodular by the argument of Lemma~\ref{w2:unimod}. The finite adelic unitary group is the restricted product of its local groups, which gives the product decomposition and the splitting of $R_B$, since $R_B$ is the group of finite adelic points of an $F$-subgroup. For $q_1\in\mathbf G_{j+1,Q}$ the function $I^0(q_1,\cdot)$ satisfies (i) to (iii) by \ref{w3s:L2}, so its quotient integral over $R_B\backslash\mathbf G_{j+2}$ converges absolutely with value $\nu_{j+2}f_Q(q_1)$, where $\nu_{j+2}=\mu_{j+2}(K_{j+2})/\rho_B(R_B\cap K_{j+2})$. For an arbitrary $q_1\in\mathbf G_{j+1}$, either $I^0(q_1,\cdot)=0$, and the inner integral is $0$, or by \ref{w3s:L2} $q_1=r\,x\,k$ with $r\in R_J$, $x\in\mathbf G_{j+1,Q}$, $k\in K'_{j+1}$, and $I^0(q_1,\cdot)=I^0(x,\cdot)$ by the invariances of \ref{w3s:L2}. Let $f_1(q_1)$ be the value of the inner integral. By Lemma~\ref{w3s:lem:uniq} it is well defined, and the invariances and support of $I^0$ in $q_1$ pass to $f_1$. On $\mathbf G_{j+1,Q}$ one has $f_1=\nu_{j+2}f_Q$.

\emph{Step 4, collapse of the outer integral.} Lemma~\ref{w3s:lem:collapse}, applied to $\mathbf G_{j+1}$, $R_J$, $K_{j+1}=K_{j+1,Q}K'_{j+1}$ and the function $f_1$, shows that $f_1$ has the absolutely convergent quotient integral $\nu_{j+1}\nu_{j+2}z_Q=\kappa_jz_Q$ over $R_J\backslash\mathbf G_{j+1}$, with $\nu_{j+1}=\mu_{j+1}(K_{j+1})/\rho_J(R_J\cap K_{j+1})$. So $I^0$ has the absolutely convergent iterated quotient integral $\kappa_jz_Q$.

\emph{Step 5, base points.} For $p_1\in\mathbf G_{j+1}$, $p_2\in\mathbf G_{j+2}$ and $g'$ in the support of $\varepsilon_{j,\alpha}$, \ref{w3s:W3} gives $(g'b_j,a\,p_1b_{j+1})=(g',p_1)(b_j,ab_{j+1})$ and $(a\,p_1b_{j+1},a'\,p_2b_{j+2})=(p_1,p_2)(ab_{j+1},a'b_{j+2})$. By \ref{w3s:L1},
\begin{align*}
\omega_j(g'b_j,a\,p_1b_{j+1})\,\mathrm{pt}_j(P_1)&=\mathrm{sc}_j\,\omega_j(g',p_1)\Phi^0_1,\\
\omega_{j+1}(a\,p_1b_{j+1},a'\,p_2b_{j+2})\,\mathrm{pt}_{j+1}(P_2)&=\mathrm{sc}_{j+1}\,\omega_{j+1}(p_1,p_2)\Phi^0_2,
\end{align*}
and $R_{a'p_2b_{j+2}}G(v)=R_{p_2}f$. Since $e_{j,\alpha}(g)=\varepsilon_{j,\alpha}(gb_j^{-1})$, \ref{w3s:W7} gives $I_s(p_1b_{j+1},p_2b_{j+2})=\mathrm{sc}_j\mathrm{sc}_{j+1}\,I^0(p_1,p_2)$. The right translations by $b_{j+1}^{-1}$ and $b_{j+2}^{-1}$ preserve absolutely convergent quotient integrals with their values (Lemma~\ref{w3s:lem:uniq}). Hence $I_s$ has the absolutely convergent iterated quotient integral $\mathrm{sc}_j\mathrm{sc}_{j+1}\kappa_jz_Q$.

\emph{Step 6, the value.} By \ref{w3s:L4} and \ref{w3s:W1}, \ref{w3s:W4},
\begin{multline*}
\mathrm{biComb}(d^Q_j(\alpha,\gamma),d^Q_{j+1}(\gamma,\beta))\,\mathrm{ctr}(\Psi'_1,\Psi'_2)\\
=c_G(j)\,\mathrm{biComb}\bigl(d^Q_j(\alpha,\gamma)d^G_j,\ d^Q_{j+1}(\gamma,\beta)d^G_{j+1}\bigr)\,\mathrm{ctr}(\Psi_1,\Psi_2),
\end{multline*}
where products of combinations denote composition of the operators. These combinations are supported away from $t\mathrm P$. By \ref{w3s:L5}, then \ref{w3s:L1} for $\mathrm{dd}$ and for $\mathrm{Sad}_j$,
\[
\sum_{\alpha,\gamma}\mathrm{bl}_j\bigl(\phi_\alpha,\mathrm{biComb}(d^Q_j(\alpha,\gamma),d^Q_{j+1}(\gamma,\beta))\,\mathrm{ctr}(\Psi'_1,\Psi'_2)\bigr)=c_G(j)\,c_T(j)\,\mathrm{two}_j\bigl(\mathrm{rd}_j(\cdot,S_j\psi',y),C\bigr)_\beta .
\]
Finally $\mathrm{tr}(e_{j+2,\beta})R_{a'}G(v)=\mathrm{tr}(\varepsilon_{j+2,\beta})R_{a'b_{j+2}}G(v)=\mathrm{tr}(\varepsilon_{j+2,\beta})f$, because $e_{j+2,\beta}(g)=\varepsilon_{j+2,\beta}(gb_{j+2}^{-1})$ and $a'$ commutes with $b_{j+2}$. By \ref{w3s:W7} the iterated integral of Step 5 equals $\kappa_j\mathrm{sc}_j\mathrm{sc}_{j+1}c_Q(j)c_G(j)c_T(j)$ times the sum in \eqref{w3s:value}, and \ref{w3s:L6} identifies the constant with $c_V(j)$.
\end{proof}

\begin{proposition}[nonvanishing of $c_V(j)$]\label{w3s:prop:cV}
Under \ref{w3s:L6}, $c_V(j)(s)\ne0$ if and only if $\mathrm{sc}_j(s)$, $\mathrm{sc}_{j+1}(s)$, $c_Q(j)(s)$, $c_G(j)$ and $c_T(j)(s)$ are all nonzero.
\end{proposition}

\begin{proof}
This is $\kappa_j>0$ (Step 1 of the proof of Theorem~\ref{w3s:thm:T2}) together with \ref{w3s:L6}.
\end{proof}

\begin{remark}[normalization]\label{w3s:rem:corr}
The constant $c_V(j)$ depends on the tuple through $\mathrm{sc}_j$, $\mathrm{sc}_{j+1}$, $c_Q(j)$ and $c_T(j)$. The object that does not depend on the tuple is the kernel $C$, together with the volume factor $\kappa_j$. In Proposition~\ref{w2:see:step} the constants $c_V(j)c_K(j)c_K(j+1)$ stand inside the sum over the tuples $s_i$, and the proof needs only that one kernel $C$ serves all tuples, so that the coefficients $a(y')$ of the kernel law (T3) can be pulled out of the sum. Nonvanishing of $c_V(j)(s)$ is not used in Proposition~\ref{w2:see:step} or Theorem~\ref{w2:see:main}. It enters only through $c_{\mathrm{top}}(t)\ne0$ in (D6), and Proposition~\ref{w3s:prop:cV} reduces it to the nonvanishing of the five scalars of the construction.
\end{remark}

\begin{remark}[what remains]\label{w3s:rem:open}
Theorem~\ref{w3s:thm:T2} reduces (T2) to Hypotheses~\ref{w3s:hyp:gen} and~\ref{w3s:hyp:loc}. Items \ref{w3s:W1} to \ref{w3s:W7} are general properties of the Weil representation, the flat maps and the Bessel period, and \ref{w3s:W6} is Fubini's theorem combined with the tensor factorization of the Weil representation of an orthogonal sum. The proof above uses \ref{w3s:W1} to \ref{w3s:W4} and \ref{w3s:W7} directly. Items \ref{w3s:W5} and \ref{w3s:W6} are used in the formal derivation of the cell properties \ref{w3s:L2} and of the twist-place law \ref{w3s:L3}. Items \ref{w3s:L1} to \ref{w3s:L6} are local computations of the construction at the twist places $Q$, at the level places $g\mathrm P$, at the places $t\mathrm P$ and at the places where every datum is unramified. Of these, \ref{w3s:L3} is the local double-step identity at the twist places. The Siegel--Weil formula and the doubling method do not enter (T2). They enter Hypothesis~\ref{w2:see:H-links} only through the surjectivity (D4).
\end{remark}

\subsection{The share of irreducible rows}\label{w2:sec:irr}

\subsubsection{Labels and weights}

Let $g$ be a prime and $\Gamma_N=\mathrm{GL}_N(\F_g)$. A \emph{label} is a monic polynomial $c\in\F_g[X]$ of degree $N$ with $c(0)\ne0$. Every irreducible representation $\pi$ of $\Gamma_N$ has a label, the characteristic polynomial of the semisimple class $s_\pi$ of its Deligne--Lusztig series. For a label $c$ let $E_c\in\C[\Gamma_N]$ be the sum of the central idempotents of the irreducible representations with label $c$. The $E_c$ are orthogonal central idempotents with $\sum_cE_c=1$. For $e\in\C[\Gamma_N]$ with $e(1)\ne0$ put
\[
\mathrm{wt}(e,c)=\frac{\operatorname{Re}(eE_c)(1)}{\operatorname{Re}e(1)},
\]
where $a(1)$ is the coefficient of the identity. For a squarefree label put $\sq(c)=\prod_{P\mid c}(g^{\deg P}-1)^{-1}$, the product over the monic irreducible factors.

\begin{lemma}\label{w2:wtprob}
If $e=e^*=e^2$ and $e(1)>0$, then $\mathrm{wt}(e,c)\ge0$ for every label and $\sum_c\mathrm{wt}(e,c)=1$.
\end{lemma}

\begin{proof}
In the left regular representation $\lambda$ on $\ell^2(\Gamma_N)$ one has $a(1)=\langle\lambda(a)\delta_1,\delta_1\rangle$. Central idempotents of $\C[\Gamma_N]$ are self-adjoint, so $eE_c=E_ceE_c$ and $(eE_c)(1)=\lVert\lambda(e)\lambda(E_c)\delta_1\rVert^2\ge0$. Summing over $c$ gives $e(1)$.
\end{proof}

Let $U\subset\Gamma_N$ be the upper unitriangular group, $\psi_0$ a nontrivial additive character of $\F_g$, and $e_U=|U|^{-1}\sum_{u\in U}\overline{\psi_0(\sum_iu_{i,i+1})}\,u$ the Gelfand--Graev idempotent. It is a self-adjoint idempotent with $e_U(1)=|U|^{-1}$, because $u\mapsto\psi_0(\sum_iu_{i,i+1})$ is a character of $U$.

\begin{proposition}[Gelfand--Graev weights]\label{w2:GG}
For every squarefree label $c$, $\mathrm{wt}(e_U,c)=\sq(c)$.
\end{proposition}

\begin{proof}
By Frobenius reciprocity $(e_UE_\pi)(1)=|\Gamma_N|^{-1}\dim\pi\cdot m_\pi$, where $m_\pi$ is the multiplicity of $\pi$ in the Gelfand--Graev representation $\mathrm{Ind}_U^{\Gamma_N}\psi$, so $\mathrm{wt}(e_U,c)=[\Gamma_N:U]^{-1}\sum_{\pi\text{ of label }c}m_\pi\dim\pi$. By Gelfand--Graev multiplicity one and the theory of Deligne and Lusztig \cite[\S10]{w2dl}, \cite[Ch.~8]{w2carter}, each series contains exactly one constituent of the Gelfand--Graev representation, with multiplicity one, the regular character $\chi_{s,\mathrm{reg}}$, and $\chi_{s,\mathrm{reg}}(1)=[\Gamma_N:C(s)]_{g'}$, the prime-to-$g$ part of the index of the centralizer. For $\Gamma_N=\mathrm{GL}_N(\F_g)$ this is also Green's theorem \cite{w2green}. If $c$ is squarefree with irreducible factors of degrees $d_1,\dots,d_j$, the centralizer of $s$ is the torus $\prod_i\F_{g^{d_i}}^\times$, of order $\prod_i(g^{d_i}-1)$, prime to $g$. Since $[\Gamma_N:U]=|\Gamma_N|_{g'}$, $\mathrm{wt}(e_U,c)=|\Gamma_N|_{g'}/(\prod_i(g^{d_i}-1)\,|\Gamma_N|_{g'})=\sq(c)$.
\end{proof}

\subsubsection{Good tuples}

Fix a prime $g$, an integer $N\ge100$, a finite set $\Pi$ of $f$ places, a finite set $J$ of $m$ indices and elements $\chi_{w,j}\in\F_g^\times$ for $w\in\Pi$, $j\in J$. Let $\bar\F_g$ be an algebraic closure. For a prime $\ell\ne g$ and $x\in\bar\F_g^\times$ let $x_{\ell'}$ be the prime-to-$\ell$ component of $x$ in the torsion group $\bar\F_g^\times$, and let
\[
BB(g,N)=\{g\}\cup\bigl\{\ell\le g^N\text{ prime}:\ g^d<\ell^{4v_\ell(g^d-1)}\text{ for some }1\le d\le N\bigr\}.
\]
A tuple $(c_w)_{w\in\Pi}$ of labels is \emph{good} if the following hold.
\begin{enumerate}[label=(G\arabic*),leftmargin=*]
\item\label{w2:G1} Every $c_w$ is squarefree.
\item\label{w2:G2} For every prime $\ell\notin BB(g,N)$ and every $w$, the map $x\mapsto x_{\ell'}$ is injective on the roots of $c_w$.
\item\label{w2:G3} For every prime $\ell\ne g$, every $w$, $j$, every sign and every root $x$ of $c_w$, $x_{\ell'}\ne(\chi_{w,j}^{\pm1})_{\ell'}$.
\item\label{w2:G4} For every $1\le a\le N/2$ there is $w$ such that $c_w$ has no divisor of degree $a$.
\item\label{w2:G5} There is $w$ such that $c_w$ has an irreducible factor of prime degree $p$ with $N<2p$ and $p+3\le N$.
\end{enumerate}


\subsubsection{Counting lemmas}

\begin{lemma}\label{w2:counts}
Let $N\ge2$ and $g\ge2$. (i) The number $I_d$ of monic irreducible polynomials of degree $d$ in $\F_g[X]$ other than $X$ satisfies $dI_d\le g^d-1$. (ii) At least $g^N-2g^{N-1}$ labels of degree $N$ are squarefree. (iii) For every squarefree label $c$ of degree $N$, $\sq(c)\ge g^{-N}$.
\end{lemma}

\begin{proof}
(i) The roots of the counted polynomials are distinct elements of $\F_{g^d}^\times$, so $dI_d\le g^d-1$. (ii) The number of squarefree monic polynomials of degree $N\ge2$ is $g^N-g^{N-1}$ \cite[Prop.~2.3]{w2rosen}, and those with $c(0)=0$ are $X$ times a squarefree monic polynomial of degree $N-1$, at most $g^{N-1}$ of them. (iii) $\prod_P(g^{\deg P}-1)^{-1}\ge\prod_Pg^{-\deg P}=g^{-N}$.
\end{proof}

\begin{lemma}[types versus cycle types]\label{w2:types}
Let $\mathcal A$ be a property of partitions of $N$. Then
\[
\sum_{\substack{c\text{ squarefree label}\\ \mathcal A(\text{degrees of the factors of }c)}}\sq(c)\le\frac{\#\{\tau\in S_N:\ \mathcal A(\text{cycle type of }\tau)\}}{N!}.
\]
In particular $\sum_{c\text{ squarefree}}\sq(c)\le1$.
\end{lemma}

\begin{proof}
A squarefree label with $n_d$ irreducible factors of degree $d$ is a choice of $n_d$ distinct irreducibles of each degree $d$, none equal to $X$. The total weight of the labels of this type is $\prod_d\binom{I_d}{n_d}(g^d-1)^{-n_d}\le\prod_d\frac{1}{n_d!}\bigl(\frac{I_d}{g^d-1}\bigr)^{n_d}\le\prod_d\frac{1}{d^{n_d}n_d!}$ by Lemma~\ref{w2:counts}(i). By Cauchy's formula the last product is the proportion of permutations of cycle type $(n_d)$. Sum over the types with property $\mathcal A$.
\end{proof}

\begin{lemma}[multiples]\label{w2:mult}
Let $h$ be a squarefree monic polynomial of degree at most $N$ with $h(0)\ne0$. Then $\sum\sq(c)\le\sq(h)$, the sum over squarefree labels $c$ of degree $N$ divisible by $h$.
\end{lemma}

\begin{proof}
Such $c$ are $c=hk$ with $k$ squarefree, prime to $h$, of degree $N-\deg h$ and $k(0)\ne0$, and $\sq(hk)=\sq(h)\sq(k)$. By Lemma~\ref{w2:types} in degree $N-\deg h$ (or trivially if the degree is $0$), $\sum_k\sq(k)\le1$.
\end{proof}

\begin{lemma}[invariant sets of permutations]\label{w2:78}
For $1\le a<n$, the proportion of $\tau\in S_n$ that leave some $a$-element subset invariant is at most $7/8$.
\end{lemma}

\begin{proof}
Let $X(\tau)$ be the number of $\tau$-invariant $a$-subsets. Counting pairs $(\tau,A)$ with $A$ invariant gives $\sum_\tau X(\tau)=\binom na\,a!\,(n-a)!=n!$. Among these pairs consider those for which $\tau|_A$ and $\tau|_{A^c}$ both have a fixed point. With $D_j$ the number of derangements of $j$ points, $2D_j\le j!$ for $j\ge1$, so the proportion of such pairs is $(1-D_a/a!)(1-D_{n-a}/(n-a)!)\ge1/4$. If $(\tau,A)$ is such a pair, with fixed points $p\in A$ and $q\notin A$, then $A\setminus\{p\}\cup\{q\}$ is another invariant $a$-subset, so $X(\tau)\ge2$. Let $B=\{\tau:X(\tau)\ge2\}$ and $n_1=\#\{\tau:X(\tau)=1\}$. Then $\sum_{\tau\in B}X(\tau)\ge n!/4$, so $n_1=n!-\sum_{\tau\in B}X(\tau)\le\tfrac34n!$, and $\#B\le\tfrac12\sum_{\tau\in B}X(\tau)=\tfrac12(n!-n_1)$. Hence $\#\{X\ge1\}=n_1+\#B\le\tfrac12n!+\tfrac12n_1\le\tfrac78n!$.
\end{proof}

\begin{lemma}[long prime cycles]\label{w2:primecycles}
For $N\ge100$ the proportion of $\tau\in S_N$ that have a cycle of prime length $p$ with $N/2<p\le N-3$ is $\sum_{N/2<p\le N-3}1/p\ge1/(5\log N)$.
\end{lemma}

\begin{proof}
A permutation has at most one cycle of length greater than $N/2$, and the number of permutations with a cycle of a given length $L>N/2$ is $N!/L$, so the events are disjoint and the proportion is the sum. For $100\le N\le2\cdot10^6$ the inequality was checked by direct computation (the smallest value of $5\log N\sum1/p$ is $2.94$, at $N=226$). For $N>2\cdot10^6$, the bounds $x/\log x<\pi(x)<1.25506\,x/\log x$ of Rosser and Schoenfeld \cite[Cor.~1]{w2rs}, valid for $x\ge17$, give $\sum_{N/2<p\le N-3}1/p\ge(\pi(N-3)-\pi(N/2))/(N-3)$ with $\pi(N-3)-\pi(N/2)\ge\frac{N-3}{\log N}-0.6276\frac{N}{\log(N/2)}\ge\frac{0.339N-3}{\log N}$, because $\log(N/2)\ge0.95\log N$ for $N>2\cdot10^6$, and $(0.339N-3)/(N-3)\ge1/5$.
\end{proof}

\begin{lemma}[prime power parts of $g-1$]\label{w2:share}
Let $g$ be a prime with $16\mid g-1$ and let $q_1\ne d$ be odd primes dividing $g-1$ with $\max(q_1,d)\ge16$. Then $s(g):=\sum_{\ell\mid g-1}\ell^{v_\ell(g-1)}/(g-1)\le\frac{49}{256\min(q_1,d)}$.
\end{lemma}

\begin{proof}
Write $g-1=\prod_ia_i$ with $a_i=\ell_i^{v_{\ell_i}(g-1)}$, so $s(g)=\sum_i1/\prod_{j\ne i}a_j$. Here $a_2\ge16$, $a_{q_1}\ge q_1$, $a_d\ge d$, and the remaining factors $b_1,\dots,b_t\ge3$ have product $B$ with $\sum_ib_i\le B$ when $t\ge1$. The terms are at most $1/(q_1d)$, $1/(16d)$, $1/(16q_1)$, and $\sum_ib_i/(16q_1dB)\le1/(16q_1d)$. With $m_0=\min(q_1,d)$ and $\max(q_1,d)\ge16$ the total is at most $(1/16+1/8+1/256)/m_0=49/(256m_0)$.
\end{proof}

\subsubsection{The bad weight at one prime}

\begin{theorem}\label{w2:badweight}
Let $g\ge2^{12}$ be a prime, $N\ge100$, and let $q_1,d$ be as in Lemma~\ref{w2:share} with $m_0=\min(q_1,d)$. Let $P$ be a nonnegative function on labels of degree $N$ with total mass $1$ that equals $\sq$ on squarefree labels. Then the $P^{\otimes\Pi}$-mass of the tuples that are not good is at most
\[
\begin{aligned}
W(g)={}&\frac{2f}{g}+\frac{f\log_2g}{\sqrt g}\Bigl(1.31+\frac{2.02\,(1+\log N)}{g^{1/4}}\Bigr)\\
&+2fm\Bigl(\frac1{4m_0}+\frac{1+\log N}{g-1}\Bigr)+\frac N2\Bigl(\frac78\Bigr)^f+e^{-f/(5\log N)}.
\end{aligned}
\]
\end{theorem}

\begin{proof}
A tuple that is not good satisfies one of the following: (I) some $c_w$ is not squarefree; (II) some $c_w$ is squarefree and fails \ref{w2:G2}; (III) for some $(w,j,\pm)$, $c_w$ is squarefree and fails \ref{w2:G3} for that index; (IV) all $c_w$ are squarefree and \ref{w2:G4} fails; (V) all $c_w$ are squarefree and \ref{w2:G5} fails. The mass is at most the sum of the masses of the five events. For an event that concerns one place $w$, the mass is the $P$-mass of the corresponding set of labels at $w$, because $P$ has mass $1$ at the other places.

(I) By Lemma~\ref{w2:counts}(ii),(iii), the squarefree labels have $\sq$-mass at least $(g^N-2g^{N-1})g^{-N}=1-2/g$. So the non-squarefree labels have $P$-mass at most $2/g$, and the union over $f$ places at most $2f/g$.

(II) Fix $w$ and a prime $\ell\notin BB(g,N)$. For $1\le\delta\le N$ the Sylow $\ell$-subgroup of $\F_{g^\delta}^\times$ has order $\ell^{v_\ell(g^\delta-1)}\le g^{\delta/4}$: for $\ell\le g^N$ this is the definition of $BB$, and for $\ell>g^N$ the group is trivial. A squarefree label fails \ref{w2:G2} if it has two distinct roots with equal $\ell'$-components. Either both are roots of one irreducible factor, or they are roots of two distinct factors. By Lemma~\ref{w2:mult} the mass is at most the sum of $\sq(Q)$ over the irreducible $Q$ with a self-collision plus the sum of $\sq(QQ')$ over the ordered pairs $Q\ne Q'$ with a cross collision.

Self-collisions. Let $Q$ have degree $\delta\ge2$ and roots $x,x^{g^i}$ with $1\le i<\delta$ and $(x^{g^i})_{\ell'}=x_{\ell'}$, so $x^{g^i-1}$ lies in the Sylow $\ell$-subgroup of $\F_{g^\delta}^\times$, which is then nontrivial, so $\ell\mid g^\delta-1$. The kernel of $y\mapsto y^{g^i-1}$ on $\F_{g^\delta}^\times$ is $\F_{g^{\gcd(i,\delta)}}^\times$, of order less than $g^{\delta/2}$. So for fixed $i$ and $\ell$ there are at most $g^{3\delta/4}$ such $x$. There are at most $\delta\log_2g$ primes dividing $g^\delta-1$ and $\delta-1$ values of $i$, so at most $(\delta-1)\log_2g\,g^{3\delta/4}$ such $Q$ of degree $\delta$, summed over all $\ell\notin BB$. With $\rho=g^{-1/4}\le1/8$ their mass is at most
\[
\log_2g\sum_{\delta\ge2}(\delta-1)\frac{g^{3\delta/4}}{g^\delta-1}\le\bigl(1+2^{-23}\bigr)\log_2g\,\frac{\rho^2}{(1-\rho)^2}\le1.31\,\frac{\log_2g}{\sqrt g}.
\]
Cross collisions. Let $Q\ne Q'$ have degrees $\delta,\delta'$ and roots $x,y$ with $x_{\ell'}=y_{\ell'}=:z$. Write $x=zx_\ell$, $y=zy_\ell$ with $\ell$-power components $x_\ell\in\F_{g^\delta}^\times$, $y_\ell\in\F_{g^{\delta'}}^\times$. Since $x\ne y$, one of $x_\ell,y_\ell$ is not $1$. Let $\mathcal A$ be the set of $(Q,Q')$ that admit such $(x,y,\ell)$ with $y_\ell\ne1$, and $\mathcal B$ the set with $x_\ell\ne1$. The Frobenius $t\mapsto t^g$ maps such a triple $(x,y,\ell)$ to another one with the same property, and moves $x$ through all $\delta$ roots of $Q$. So every $(Q,Q')\in\mathcal A$ yields at least $\delta$ triples with distinct $x$. Given $x$, the prime $\ell$ divides $g^{\delta'}-1$ (at most $\delta'\log_2g$ choices) and $y=x_{\ell'}y_\ell$ is determined by $y_\ell$ (at most $g^{\delta'/4}$ choices). Hence $\#\mathcal A\le g^\delta\delta'\log_2g\,g^{\delta'/4}/\delta$, and symmetrically $\#\mathcal B\le g^{\delta'}\delta\log_2g\,g^{\delta/4}/\delta'$. The mass is at most
\[
\sum_{\delta,\delta'=1}^N\frac{\#\mathcal A+\#\mathcal B}{(g^\delta-1)(g^{\delta'}-1)}\le2\Bigl(1+\frac1{g-1}\Bigr)^2\log_2g\Bigl(\sum_{\delta\le N}\frac1\delta\Bigr)\sum_{k\ge1}kg^{-3k/4}\le2.02\,\frac{(1+\log N)\log_2g}{g^{3/4}}.
\]
Adding and summing over the $f$ places gives the second term of $W(g)$.

(III) Fix $(w,j,\pm)$ and put $\chi=\chi_{w,j}^{\pm1}\in\F_g^\times$. For $\ell\ne g$, $x_{\ell'}=\chi_{\ell'}$ holds exactly when $x/\chi$ has $\ell$-power order. Let $Q$ be irreducible of degree $\delta$ with such a root $x$. Then $x\in\chi\cdot\Syl_\ell(\F_{g^\delta}^\times)$, and every root of $Q$ has the same property because $t\mapsto t^g$ fixes $\chi$. Hence the number of such $Q$ is at most $\delta^{-1}\sum_{\ell\mid g^\delta-1}\ell^{v_\ell(g^\delta-1)}$; the primes $\ell\nmid g^\delta-1$ contribute the single element $x=\chi$, which lies in every coset counted. For $\ell\mid g-1$, the lifting of the exponent gives $v_\ell(g^\delta-1)=v_\ell(g-1)+v_\ell(\delta)$, valid for odd $\ell$ and for $\ell=2$ because $4\mid g-1$, so $\ell^{v_\ell(g^\delta-1)}\le\delta\,\ell^{v_\ell(g-1)}$. For $\ell\nmid g-1$ the prime powers $\ell^{v_\ell(g^\delta-1)}$ are pairwise coprime divisors of $(g^\delta-1)/(g-1)$, so their sum is at most $(g^\delta-1)/(g-1)$. With $s(g)$ of Lemma~\ref{w2:share}, the mass of the labels with a bad factor is at most, by Lemma~\ref{w2:mult},
\[
\sum_{\delta=1}^N\frac{\delta(g-1)s(g)+(g^\delta-1)/(g-1)}{\delta\,(g^\delta-1)}\le s(g)\Bigl(1+\frac1{g-1}\Bigr)+\frac{1+\log N}{g-1}\le\frac1{4m_0}+\frac{1+\log N}{g-1}.
\]
There are $2fm$ indices $(w,j,\pm)$.

(IV) A squarefree label has a divisor of degree $a$ exactly when the degrees of its irreducible factors have a sub-multiset with sum $a$, a property of the type. By Lemmas~\ref{w2:types} and~\ref{w2:78} the $\sq$-mass of the squarefree labels with a divisor of degree $a$ is at most $7/8$. For fixed $a$ the event that this holds at every place has mass at most $(7/8)^f$, and there are at most $N/2$ values of $a$.

(V) By Lemmas~\ref{w2:types} and~\ref{w2:primecycles} the $\sq$-mass of the squarefree labels without an irreducible factor of prime degree in $(N/2,N-3]$ is at most $1-1/(5\log N)$. At all $f$ places the mass is at most $(1-1/(5\log N))^f\le e^{-f/(5\log N)}$.
\end{proof}

\subsubsection{From weights to dimensions}

The construction attaches to each element $\tau$ of an orthonormal eigenbasis of $\mathcal H$ a tuple of labels $\mathrm{lab}_k(\tau)=(\mathrm{lab}_{k,w}(\tau))_{w\in\Pi_k}$ at each of two auxiliary primes $g_0,g_1$, where $\Pi_k$ is the set of $f$ places of $F$ above $g_k$ and the level at $w$ contains $\mathrm{GL}_N(\F_{g_k})$ as a quotient, with $N=r$. For each $(k,w)$ it fixes an element $e_{k,w}\in\C[\mathrm{GL}_N(\F_{g_k})]$, the local factor of the level idempotent, and characters $\chi_{w,j}\in\F_{g_k}^\times$, $j\in J$, $|J|=m$. For a set $S$ of label tuples at $k$ put
\[
w_k(S)=\sum_{t\in S}\prod_{w\in\Pi_k}\mathrm{wt}(e_{k,w},t_w).
\]

\begin{hypothesis}[local weights]\label{w2:H-wt}
Each $e_{k,w}$ is a self-adjoint idempotent with $e_{k,w}(1)>0$, and $(1-1/g_k)\,\mathrm{wt}(e_U,c)\le\mathrm{wt}(e_{k,w},c)$ for every label $c$.
\end{hypothesis}

\begin{hypothesis}[multiplicity formula]\label{w2:H-cnt}
There is $\varepsilon_C\in[0,\tfrac12]$ with $\lvert D_{\mathrm{all}}-D^*\rvert\le\varepsilon_CD^*$ and, for $k\in\{0,1\}$ and every set $S$ of label tuples at $k$,
\[
\Bigl\lvert\#\{\tau:\mathrm{lab}_k(\tau)\in S\}-D^*w_k(S)\Bigr\rvert\le\varepsilon_C\,D^*w_k(S).
\]
\end{hypothesis}

\begin{hypothesis}[rows are the good forms]\label{w2:H-rows}
$D_{\mathrm{rows}}$ equals the number of $\tau$ whose label tuples at $g_0$ and at $g_1$ are both good for the characters $\chi_{w,j}$.
\end{hypothesis}

\begin{lemma}[comparison of weights]\label{w2:transfer}
Let $w_U,w_J$ be probability weights on a finite set at each of $f$ places with $(1-1/g)w_U\le w_J$. For every set $A$ of tuples, $w_J^{\otimes f}(A)\le w_U^{\otimes f}(A)+f/g$.
\end{lemma}

\begin{proof}
Write $w_J=(1-1/g)w_U+\tfrac1g\nu$ with $\nu$ a probability weight. Expand the product. The term $(1-1/g)^fw_U^{\otimes f}(A)$ is at most $w_U^{\otimes f}(A)$, and the other terms have total mass $1-(1-1/g)^f\le f/g$.
\end{proof}

\begin{theorem}\label{w2:Iproof}
Assume Hypotheses~\ref{w2:H-wt}, \ref{w2:H-cnt} and~\ref{w2:H-rows}, that $16\mid g_k-1$ and $q_1,d$ divide $g_k-1$ for both $k$, with $q_1\ne d$ odd primes and $\max(q_1,d)\ge16$, and the sizes
\begin{equation}\label{w2:sizes}
g_k\ge e^{8x^{\varepsilon'}},\quad x^2\le f\le x^{C_s},\quad m\le x^{C_s},\quad100\le N\le x^{C_s},\quad\min(q_1,d)\ge e^{x^{\varepsilon'/6}},
\end{equation}
with $0<\varepsilon'\le1$ and $C_s\ge1$. Then $0\le\delta\le\exp(-x^{\varepsilon'/7})$ for $x\ge x_0(\varepsilon',C_s)$. With $\varepsilon'=\varepsilon_z$ this is Theorem~\ref{w2:I}.
\end{theorem}

\begin{proof}
Let $B_k$ be the set of label tuples at $g_k$ that are not good. By Hypothesis~\ref{w2:H-rows} every $\tau$ outside the rows has $\mathrm{lab}_0(\tau)\in B_0$ or $\mathrm{lab}_1(\tau)\in B_1$, so $0\le D_{\mathrm{all}}-D_{\mathrm{rows}}\le\sum_k\#\{\tau:\mathrm{lab}_k(\tau)\in B_k\}$, which gives $\delta\ge0$. By Hypothesis~\ref{w2:H-cnt},
\[
\delta=\frac{D_{\mathrm{all}}-D_{\mathrm{rows}}}{D_{\mathrm{all}}}\le\frac{(1+\varepsilon_C)D^*(w_0(B_0)+w_1(B_1))}{(1-\varepsilon_C)D^*}\le3\bigl(w_0(B_0)+w_1(B_1)\bigr).
\]
By Lemma~\ref{w2:wtprob} and Hypothesis~\ref{w2:H-wt} the weights $\mathrm{wt}(e_{k,w},\cdot)$ and $\mathrm{wt}(e_U,\cdot)$ are probability weights with $(1-1/g_k)\mathrm{wt}(e_U,\cdot)\le\mathrm{wt}(e_{k,w},\cdot)$, so Lemma~\ref{w2:transfer} gives $w_k(B_k)\le\mathrm{wt}(e_U,\cdot)^{\otimes f}(B_k)+f/g_k$. By Proposition~\ref{w2:GG} the weight $P=\mathrm{wt}(e_U,\cdot)$ satisfies the hypothesis of Theorem~\ref{w2:badweight}, and the sizes give $g_k\ge2^{12}$ for large $x$. Hence $w_k(B_k)\le W(g_k)+f/g_k$.

Under \eqref{w2:sizes} the terms of $W(g_k)+f/g_k$ are bounded as follows for large $x$. The terms $3f/g_k$ and $2fm(1+\log N)/(g_k-1)$ are at most $4x^{2C_s}(1+C_s\log x)e^{-8x^{\varepsilon'}}$. Since $t\mapsto\log t/\sqrt t$ decreases for $t\ge e^2$, the second term is at most $3x^{C_s}\cdot12x^{\varepsilon'}e^{-4x^{\varepsilon'}}$. The term $2fm/(4m_0)$ is at most $\tfrac12x^{2C_s}e^{-x^{\varepsilon'/6}}$. Finally $\tfrac N2(7/8)^f\le x^{C_s}e^{-x^2/8}$ and $e^{-f/(5\log N)}\le e^{-x^2/(5C_s\log x)}$. Since $\varepsilon'\le1$, every term is at most $x^{2C_s}e^{-x^{\varepsilon'/6}}$ for large $x$, so $\delta\le36\,x^{2C_s}e^{-x^{\varepsilon'/6}}\le e^{-x^{\varepsilon'/7}}$ once $x^{\varepsilon'/6}-x^{\varepsilon'/7}\ge2C_s\log x+\log36$.
\end{proof}

\begin{remark}\label{w2:rem:irr}
The irreducibility certificates \ref{w2:G2} to \ref{w2:G5} are used later, where Jordan\'s theorem enters. Theorem~\ref{w2:I} itself is a count and does not use them.
\end{remark}

\subsubsection{The level idempotents at the auxiliary primes}\label{w2x:sec:wt}
Let $g$ be a prime, $\F=\F_g$, and let $\psi_0\colon\F\to\C^\times$ be a nontrivial additive character. Put $G_n=\GL_n(\F)$, let $U_n\subset G_n$ be the group of upper unitriangular matrices, and index rows and columns by $0,\dots,n-1$. For a subgroup $H$ of a finite group $G$ and a character $\lambda\colon H\to\C^\times$ put
\[
e(H,\lambda)=\lvert H\rvert^{-1}\sum_{h\in H}\overline{\lambda(h)}\,h\in\C[G].
\]
This is a self-adjoint idempotent with $e(H,\lambda)(1)=\lvert H\rvert^{-1}$. For a representation $(\pi,V)$ of $G$ the operator $\pi(e(H,\lambda))$ is the orthogonal projection onto $V^{(H,\lambda)}=\{v\in V:\pi(h)v=\lambda(h)v\ (h\in H)\}$. For $a\in\C[G]$ and a class function $\chi$ on $G$ put $\tr(a,\chi)=\sum_{x\in G}a(x)\chi(x)$. If $\chi=\chi_\pi$ is the character of $\pi$, then $\tr(a,\chi_\pi)=\Tr\pi(a)$, so
\begin{equation}\label{w39:trmult}
\tr\bigl(e(H,\lambda),\chi_\pi\bigr)=\dim V^{(H,\lambda)}=\dim\Hom_H(\lambda,\pi|_H).
\end{equation}
Since the class functions are the linear combinations of the irreducible characters and $\tr(a,\cdot)$ is linear, an identity between the numbers \eqref{w39:trmult} for all irreducible $\pi$ is an identity of $\tr(a,\cdot)$ on all class functions.

\paragraph{Pattern groups.}

Let $I=\{0,\dots,N-1\}$ and let $\prec$ be a strict partial order on $I$, that is, an irreflexive transitive relation. Write $a\preceq b$ for $a=b$ or $a\prec b$. The \emph{pattern group} of $\prec$ is
\[
U_\prec=\{u\in\GL_N(\F):u_{aa}=1\text{ for all }a,\ u_{ab}=0\text{ unless }a\prec b\text{ or }a=b\}.
\]

\begin{lemma}\label{w39:pattern}
(i) $U_\prec$ is a subgroup of $\GL_N(\F)$ of order $g^{\#\{(a,b):a\prec b\}}$.
(ii) Let $E$ be a set of pairs $(a,b)$ with $a\prec b$, and let $c\colon E\to\F$. Suppose that for every $(a,b)\in E$ there is no $m$ with $a\prec m\prec b$. Then $\lambda_{E,c}(u)=\psi_0\bigl(\sum_{(a,b)\in E}c_{ab}u_{ab}\bigr)$ is a character of $U_\prec$.
\end{lemma}

\begin{proof}
(i) For $u,v\in U_\prec$, $(uv)_{ab}=\sum_mu_{am}v_{mb}$, and a nonzero term needs $a\preceq m\preceq b$. By transitivity $a\preceq b$, and if $a=b$ the only term is $m=a$, because $a\prec m\prec a$ would give $a\prec a$. So $uv\in U_\prec$ with unit diagonal. Write $u=1+x$. Every product $x^j$ is supported on pairs $a\prec b$, and $x^N=0$ because a strictly increasing chain in $I$ has at most $N$ elements. Hence $u^{-1}=\sum_{j<N}(-x)^j\in U_\prec$. The entries at the pairs $a\prec b$ are free, which gives the order.
(ii) For $u,v\in U_\prec$ and $(a,b)\in E$ we have $(uv)_{ab}=u_{ab}+v_{ab}+\sum_{a\prec m\prec b}u_{am}v_{mb}=u_{ab}+v_{ab}$.
\end{proof}

\begin{lemma}[one double coset]\label{w39:mackey}
Let $H$ be a finite group with subgroups $H_1,H_2$ and characters $\lambda_i\colon H_i\to\C^\times$ that agree on $H_1\cap H_2$. Suppose that for $i\in\{1,2\}$ and every $h\in H\setminus H_i$ there is $a\in H_i$ with $h^{-1}ah\in H_i$ and $\lambda_i(h^{-1}ah)\ne\lambda_i(a)$. Then $\operatorname{Ind}_{H_1}^H\lambda_1\cong\operatorname{Ind}_{H_2}^H\lambda_2$. Consequently, for every finite group $G\supseteq H$ and every representation $\pi$ of $G$, $\dim\pi^{(H_1,\lambda_1)}=\dim\pi^{(H_2,\lambda_2)}$, and $\tr(e(H_1,\lambda_1),\chi)=\tr(e(H_2,\lambda_2),\chi)$ for every class function $\chi$ on $G$.
\end{lemma}

\begin{proof}
For subgroups $A,B\subseteq H$ and characters $\alpha,\beta$, Mackey's formula (J.-P. Serre, \emph{Linear representations of finite groups}, \S7.4) gives
\[
\bigl\langle\operatorname{Ind}_A^H\alpha,\operatorname{Ind}_B^H\beta\bigr\rangle_H=\sum_{h\in A\backslash H/B}\bigl\langle\alpha,\beta^h\bigr\rangle_{A\cap hBh^{-1}},\qquad\beta^h(y)=\beta(h^{-1}yh).
\]
Each term is $1$ if $\alpha=\beta^h$ on $A\cap hBh^{-1}$ and $0$ otherwise. Take $A=B=H_i$ and $\alpha=\beta=\lambda_i$. For $h\notin H_i$ the element $a$ of the hypothesis lies in $H_i\cap hH_ih^{-1}$ and satisfies $\lambda_i^h(a)=\lambda_i(h^{-1}ah)\ne\lambda_i(a)$, so the term of $h$ vanishes. The double coset of $1$ contributes $1$. Hence $\operatorname{Ind}_{H_i}^H\lambda_i$ is irreducible. Next take $A=H_1$, $B=H_2$. The term of $h=1$ is $1$ because $\lambda_1=\lambda_2$ on $H_1\cap H_2$, and the other terms are nonnegative, so the two irreducible induced representations have a common constituent and are isomorphic. By Frobenius reciprocity and transitivity of induction, $\dim\pi^{(H_i,\lambda_i)}=\dim\Hom_G(\operatorname{Ind}_H^G\operatorname{Ind}_{H_i}^H\lambda_i,\pi)$, which does not depend on $i$. The last claim follows from \eqref{w39:trmult} and linearity.
\end{proof}

\begin{lemma}[exchange]\label{w39:exchange}
Let $\prec_1,\prec_2$ be strict partial orders on $I$ whose union $\prec$ is again a strict partial order. Let $p,q\in I$, $W\subseteq I\setminus\{p,q\}$, and let $E$ be a set of pairs with a function $c\colon E\to\F$. Assume the following.
\begin{enumerate}[label=(X\arabic*),leftmargin=*]
\item\label{w39:X1} The pairs with $a\prec_1b$ and not $a\prec_2b$ are the pairs $(p,w)$, $w\in W$. The pairs with $a\prec_2b$ and not $a\prec_1b$ are the pairs $(w,q)$, $w\in W$.
\item\label{w39:X2} $(p,q)\in E$ and $c_{pq}\ne0$.
\item\label{w39:X3} If $w\in W$, $x\preceq p$ and $w\preceq y$, then $x\prec_1y$, and $(x,y)\in E$ only for $(x,y)=(p,q)$.
\item\label{w39:X4} If $w\in W$, $x\preceq w$ and $q\preceq y$, then $x\prec_2y$, and $(x,y)\in E$ only for $(x,y)=(p,q)$.
\item\label{w39:X5} For $(a,b)\in E$ and $j\in\{1,2\}$ there is no $m$ with $a\prec_jm\prec_jb$.
\end{enumerate}
Then $\lambda=\lambda_{E,c}$ is a character of $U_{\prec_1}$ and of $U_{\prec_2}$, and $\dim\pi^{(U_{\prec_1},\lambda)}=\dim\pi^{(U_{\prec_2},\lambda)}$ for every representation $\pi$ of every finite group containing $U_\prec$.
\end{lemma}

\begin{proof}
By Lemma~\ref{w39:pattern} the groups $H=U_\prec$, $H_1=U_{\prec_1}$, $H_2=U_{\prec_2}$ satisfy $H_1,H_2\subseteq H$, and by \ref{w39:X5} $\lambda$ is a character of $H_1$ and of $H_2$. The two characters are given by one formula, so they agree on $H_1\cap H_2$. We verify the hypothesis of Lemma~\ref{w39:mackey}.

Let $h\in H\setminus H_1$. By \ref{w39:X1} the pairs of $\prec$ outside $\prec_1$ are the pairs $(w,q)$, so $h_{wq}\ne0$ for some $w\in W$. Since $p\prec_1w$, the element $a=1+tE_{pw}$ lies in $H_1$ for every $t\in\F$. Then $(h^{-1}ah)_{xy}=\delta_{xy}+t\,(h^{-1})_{xp}h_{wy}$. A nonzero off-diagonal entry needs $(h^{-1})_{xp}\ne0$ and $h_{wy}\ne0$, so $x\preceq p$ and $w\preceq y$ because $h,h^{-1}\in U_\prec$, and \ref{w39:X3} gives $x\prec_1y$. Hence $h^{-1}ah\in H_1$. By \ref{w39:X3} the only pair of $E$ among these is $(p,q)$, and $(h^{-1})_{pp}=1$, so $\lambda(h^{-1}ah)=\psi_0(c_{pq}t\,h_{wq})$. By \ref{w39:X3} with $(x,y)=(p,w)$ the pair $(p,w)$ lies in $E$ only if $w=q$, which is excluded, so $\lambda(a)=1$. Since $c_{pq}h_{wq}\ne0$ and $\psi_0\ne1$ there is $t$ with $\psi_0(c_{pq}t\,h_{wq})\ne1$.

Let $h\in H\setminus H_2$. Then $h^{-1}\notin H_2$, and by \ref{w39:X1} $(h^{-1})_{pw}\ne0$ for some $w\in W$. Put $a=1+tE_{wq}\in H_2$. The off-diagonal entries of $h^{-1}ah$ are $t\,(h^{-1})_{xw}h_{qy}$ with $x\preceq w$ and $q\preceq y$, so \ref{w39:X4} gives $h^{-1}ah\in H_2$ and $\lambda(h^{-1}ah)=\psi_0(c_{pq}t\,(h^{-1})_{pw})$, while $\lambda(a)=1$ because $(w,q)\in E$ would force $w=p$. Choose $t$ as before.
\end{proof}

\paragraph{Exchange of roots.}

Let $k\ge1$ and $s\ge0$ be integers and $N=k+2s+1$. Put $W=\{s,\dots,s+k-1\}$, $c_0=s+k$, and $O=I\setminus W=\{0,\dots,s-1\}\cup\{c_0,\dots,N-1\}$ with its natural order. Define
\[
U_J=\{u\in U_N:u_{a,c_0}=0\ (a\in W)\},\qquad U'=\{u\in U_N:u_{a,k}=0\ (0\le a<k)\},
\]
\[
\psi_J(u)=\psi_0\Bigl(\sum_{(a,b)\in E}u_{ab}\Bigr),\qquad\psi'(u)=\psi_0\Bigl(\sum_{l=0}^{N-2}u_{l,l+1}\Bigr),
\]
where $E$ is the set of pairs $(a,b)$ that are consecutive in $O$ (the pairs $(l,l+1)$ with $l+1<s$, the pair $(s-1,c_0)$ when $s\ge1$, and the pairs $(l,l+1)$ with $l\ge c_0$) together with the pairs $(l,l+1)$ with $l,l+1\in W$. Put $e_J=e(U_J,\psi_J)$, $e'=e(U',\psi')$ and $e_U=e(U_N,\psi')$. In the construction $(N,k,s)=(r,k,r_S)$ with $r_S=m/2-1$, and $e_J$ is the local factor of the level idempotent at every place above $g_0$ and $g_1$.

\begin{theorem}[exchange of roots]\label{w39:thmA}\label{w2x:exch}
For every class function $\chi$ on $G_N$, $\tr(e_J,\chi)=\tr(e',\chi)$. Equivalently $\dim\pi^{(U_J,\psi_J)}=\dim\pi^{(U',\psi')}$ for every representation $\pi$ of $G_N$.
\end{theorem}

\begin{proof}
For $0\le i\le s$ let $T_i=\{b\in O:b>i\}$ if $i<s$ and $T_s=\{c_0+1,\dots,N-1\}$. Define $a\prec_ib$ if one of the following holds: $a,b\in W$ and $a<b$; $a,b\in O$ and $a<b$; $a\in O$, $a<i$ and $b\in W$; $a\in W$ and $b\in T_i$.

\emph{$\prec_i$ is a strict partial order.} It is contained in no cycle because a relation from $W$ to $O$ ends in $T_i\subseteq\{b>i\}$ and a relation from $O$ to $W$ starts in $\{a<i\}$, so no chain passes from $W$ to $O$ and back to $W$. Transitivity is checked on the four types. For $a\prec_iw\prec_ib$ with $a,b\in O$ and $w\in W$ we have $a<i<b$ when $b<s$, and $a<s\le c_0\le b$ otherwise, so $a<b$. For $w\prec_ib\prec_ib'$ with $b,b'\in O$ we have $b'>b$, and $b\in T_i$ implies $b'\in T_i$ since $T_i$ is an upper set of $O$ (for $i<s$ it is $\{b\in O:b>i\}$, and $T_s$ is the top of $O$). The remaining cases are immediate.

\emph{The character.} Each pair of $E$ is a covering pair of $\prec_i$. For a pair $(a,b)$ consecutive in $O$ no element of $O$ lies between, and $a\prec_iw\prec_ib$ with $w\in W$ needs $a<i$ and $b\in T_i$. If $b<s$ then $b=a+1\le i$ contradicts $b>i$. If $b=c_0$ then $a=s-1<i$ forces $i=s$, and $c_0\notin T_s$. If $b>c_0$ then $a\ge c_0>i$. For a pair $(l,l+1)$ in $W$ an intermediate element of $O$ would need $l\prec_im\prec_il+1$, that is $m\in T_i$ and $m<i$, which is impossible. By Lemma~\ref{w39:pattern} $\psi_J=\lambda_{E,1}$ is a character of every $U_{\prec_i}$.

\emph{The end points.} For $i=s$ the relation $\prec_s$ consists of the pairs $a<b$ other than $(w,c_0)$, $w\in W$, so $U_{\prec_s}=U_J$. For $i=0$ let $\sigma$ be the permutation of $I$ with $\sigma(a)=s+a$ for $a<k$, $\sigma(a)=a-k$ for $k\le a<k+s$, and $\sigma(a)=a$ for $a\ge k+s$. It maps $\{0,\dots,k-1\}$ increasingly onto $W$ and $\{k,\dots,N-1\}$ increasingly onto $O$. A direct check of the four types gives $\sigma(a)\prec_0\sigma(b)$ if and only if $a<b$ and not ($a<k$ and $b=k$), and $(\sigma(a),\sigma(b))\in E$ if and only if $b=a+1\ne k$. With the permutation matrix $P$, $Pe_a=e_{\sigma(a)}$, we get $(P^{-1}uP)_{ab}=u_{\sigma(a)\sigma(b)}$, hence $P^{-1}U_{\prec_0}P=U'$ and $\psi_J(PuP^{-1})=\psi'(u)$ for $u\in U'$ (the entry $u_{k-1,k}$ vanishes on $U'$). Since $\chi$ is a class function, $\tr(e(U_{\prec_0},\psi_J),\chi)=\tr(e',\chi)$.

\emph{One step.} Let $1\le i\le s$. We apply Lemma~\ref{w39:exchange} to $\prec_1=\prec_i$, $\prec_2=\prec_{i-1}$, $p=i-1$, $q=i$ if $i<s$ and $q=c_0$ if $i=s$, the set $W$ and $E$ with $c\equiv1$. Note that $q$ is the successor of $p$ in $O$. The union $\prec$ is the relation of the same four types with $\{a<i\}$ in the third type and $T'=\{b\in O:b\ge i\}$ in the fourth, which is a strict partial order by the argument above. Condition \ref{w39:X1} holds because $\{a\in O:a<i\}\setminus\{a<i-1\}=\{i-1\}$ and $T_{i-1}\setminus T_i=\{q\}$. Condition \ref{w39:X2} holds because $(p,q)$ is consecutive in $O$. For \ref{w39:X3} let $x\preceq p$ and $w\preceq y$ with $w\in W$. Since $p\in O$ and no element of $W$ precedes $p$ (that would need $p\ge i$), $x\in O$ and $x\le i-1$. Either $y\in W$, and then $x\prec_iy$ by the third type, or $y\in O$ with $y\ge i>x$, and then $x\prec_iy$ by the second type. A pair of $E$ of this form joins two consecutive elements $x<y$ of $O$ with $x\le i-1<i\le y$, so it is $(p,q)$. For \ref{w39:X4} let $x\preceq w$ and $q\preceq y$. Then $y\in O$ with $y\ge q\ge i$, because nothing in $W$ follows $q\ge i$ in $\prec$. If $x\in W$ then $y\in T_{i-1}$, and if $x\in O$ then $x<i\le y$. In both cases $x\prec_{i-1}y$, and a pair of $E$ of this form has $x\in O$, $x<i\le q\le y$, $x,y$ consecutive in $O$, so it is $(p,q)$. Condition \ref{w39:X5} was verified above. Lemma~\ref{w39:exchange} gives $\tr(e(U_{\prec_i},\psi_J),\chi)=\tr(e(U_{\prec_{i-1}},\psi_J),\chi)$.

Chaining the steps $i=s,s-1,\dots,1$ with the two end points proves the theorem. For $s=0$ there is no step and $U_{\prec_0}=U_{\prec_s}$.
\end{proof}

\paragraph{Column expansion.}

For $n\ge1$ let $P_n\subseteq G_n$ be the mirabolic subgroup of matrices whose last row is $(0,\dots,0,1)$. Embed $G_{n-1}\subseteq P_n$ by $h\mapsto\operatorname{diag}(h,1)$ and let $Y_n=\{y(x):x\in\F^{n-1}\}$, where $y(x)$ is the unipotent matrix with last column $(x,1)$. Then $P_n=G_{n-1}\ltimes Y_n$, $hy(x)h^{-1}=y(hx)$, and $U_n=U_{n-1}\ltimes Y_n$. For $0\le i\le n-1$ let
\[
\psi^{(i)}_n(u)=\psi_0\Bigl(\sum_{0\le l\le n-2,\ l\ne i-1}u_{l,l+1}\Bigr)\qquad(u\in U_n),
\]
so that $\psi^{(0)}_n=\psi_n$ is the generic character. For a representation $\tau$ of $P_n$ put
\[
A(\tau)=\dim\tau^{(U_{n-1},\psi_{n-1})},\qquad B_i(\tau)=\dim\tau^{(U_n,\psi^{(i)}_n)}\quad(0\le i\le n-1).
\]
For a representation $\rho$ of $G_n$ put $\mathcal W(\rho)=\dim\rho^{(U_n,\psi_n)}$.

\begin{lemma}[orbits of the unipotent group]\label{w39:orbits}
Let $m\ge1$ and let $\sigma$ be a representation of $P_m$. For $0\le j\le m-2$ let $S_j=\{u\in U_{m-1}:u_{jl}=0\ (l>j)\}$, and for $c\in\F^\times$ let $\chi'_{j,c}(s\,y(x))=\psi_{m-1}(s)\psi_0(cx_j)$ on $S_jY_m$. Then $\chi'_{j,c}$ is a character and
\[
A(\sigma)=B_{m-1}(\sigma)+\sum_{j=0}^{m-2}\sum_{c\in\F^\times}X_{j,c}(\sigma),\qquad X_{j,c}(\sigma)=\dim\sigma^{(S_jY_m,\chi'_{j,c})}.
\]
\end{lemma}

\begin{proof}
For a row vector $\mu\in\F^{m-1}$ let $\sigma_\mu=\{v:\sigma(y(x))v=\psi_0(\mu x)v\}$. Since $Y_m$ is abelian, $\sigma=\bigoplus_\mu\sigma_\mu$, and $\sigma(h)\sigma_\mu=\sigma_{\mu h^{-1}}$ for $h\in G_{m-1}$. The group $U_{m-1}$ acts on row vectors by $\mu\mapsto\mu u^{-1}$. Its orbits are $\{0\}$ and, for $0\le j\le m-2$ and $c\in\F^\times$, the sets $O_{j,c}$ of $\mu$ with $\mu_l=0$ for $l<j$ and $\mu_j=c$. Indeed $\mu u^{-1}$ has the same first nonzero coordinate as $\mu$, and $c\,e_j^{T}(1+\sum_{l>j}c^{-1}\mu_lE_{jl})=\mu$ for $\mu\in O_{j,c}$. The stabilizer of $c\,e_j^T$ is $S_j$, which therefore fixes the character $\psi_0(cx_j)$ of $Y_m$, so $\chi'_{j,c}$ is a character of $S_j\ltimes Y_m$. The sum $\sigma_{O}=\bigoplus_{\mu\in O}\sigma_\mu$ over one orbit is $U_{m-1}$-stable and isomorphic to $\operatorname{Ind}_{S_j}^{U_{m-1}}\sigma_{ce_j}$, so by Frobenius reciprocity $\dim\sigma_O^{(U_{m-1},\psi_{m-1})}=\dim\sigma_{ce_j}^{(S_j,\psi_{m-1})}=X_{j,c}(\sigma)$. The orbit $\{0\}$ gives $\sigma^{Y_m}$, and $\dim(\sigma^{Y_m})^{(U_{m-1},\psi_{m-1})}=B_{m-1}(\sigma)$ because $\psi^{(m-1)}_m(uy)=\psi_{m-1}(u)$.
\end{proof}

\begin{lemma}[induction from the mirabolic subgroup]\label{w39:indP}
Let $m\ge1$ and let $\sigma$ be a representation of $P_m$. For $0\le j\le m-1$ let $Y_m^{<j}=\{y(x):x_l=0\ (l\ge j)\}$ and $H_j=U_{m-1}Y_m^{<j}$, and for $c\in\F^\times$ let
\[
\chi_{j,c}(u\,y(x))=\psi_0\Bigl(\sum_{0\le l\le m-3,\ l\ne j-1}u_{l,l+1}+c\,x_{j-1}\Bigr),
\]
with $x_{-1}=0$. Then $\chi_{j,c}$ is a character of the group $H_j$ and
\[
\mathcal W\bigl(\operatorname{Ind}_{P_m}^{G_m}\sigma\bigr)=\sum_{j=0}^{m-1}\sum_{c\in\F^\times}D_{j,c}(\sigma),\qquad D_{j,c}(\sigma)=\dim\sigma^{(H_j,\chi_{j,c})}.
\]
\end{lemma}

\begin{proof}
By Mackey's formula and Frobenius reciprocity,
\[
\mathcal W(\operatorname{Ind}_{P_m}^{G_m}\sigma)=\sum_{x\in U_m\backslash G_m/P_m}\dim\sigma^{(x^{-1}(U_m\cap xP_mx^{-1})x,\ \psi_m(x\,\cdot\,x^{-1}))}.
\]
The map $gP_m\mapsto e_{m-1}^Tg^{-1}$ identifies $G_m/P_m$ with the nonzero row vectors, and $U_m$ acts by $\mu\mapsto\mu u^{-1}$. As in Lemma~\ref{w39:orbits} the orbits are indexed by $(j,c)$ with $0\le j\le m-1$ and $c\in\F^\times$, with representative $c\,e_j^T$ and stabilizer $U_{(j)}=\{u\in U_m:u_{jl}=0\ (l>j)\}$. Let $\beta$ be the bijection of $I$ with $\beta(l)=l$ for $l<j$, $\beta(j)=m-1$ and $\beta(l)=l-1$ for $l>j$, let $w$ be the permutation matrix with $we_l=e_{\beta(l)}$, and let $t=\operatorname{diag}(1,\dots,1,c)$. Then $x$ with $x^{-1}=tw$ satisfies $e_{m-1}^Tx^{-1}=c\,e_j^T$, so $U_m\cap xP_mx^{-1}=U_{(j)}$. Since $(wzw^{-1})_{\beta(a)\beta(b)}=z_{ab}$ and $\beta$ is increasing on $I\setminus\{j\}$, the group $wU_{(j)}w^{-1}$ consists of the matrices of $U_m$ whose block on $\{0,\dots,m-2\}$ is arbitrary in $U_{m-1}$ and whose last column is supported on the rows $l<j$ (these come from the entries $(l,j)$). Conjugation by $t$ rescales the last column, so $x^{-1}U_{(j)}x=H_j$. For $z\in H_j$ put $u=xzx^{-1}=w^{-1}(t^{-1}zt)w$, so $u_{ab}=(t^{-1}zt)_{\beta(a)\beta(b)}$ and $(t^{-1}zt)_{a,m-1}=c\,z_{a,m-1}$ for $a<m-1$. The superdiagonal entries of $u$ are $u_{l,l+1}=z_{l,l+1}$ for $l+1<j$, $u_{j-1,j}=c\,z_{j-1,m-1}$, $u_{j,j+1}=0$, and $u_{l,l+1}=z_{l-1,l}$ for $l>j$. Hence $\psi_m(xzx^{-1})=\chi_{j,c}(z)$.
\end{proof}

\begin{lemma}\label{w39:DX}
In the notation of Lemmas~\ref{w39:orbits} and~\ref{w39:indP}, $D_{j+1,c}(\sigma)=X_{j,c}(\sigma)$ for $0\le j\le m-2$ and $c\in\F^\times$.
\end{lemma}

\begin{proof}
On $I=\{0,\dots,m-1\}$ let $a\prec_1b$ mean $a<b\le m-2$, or $b=m-1$ and $a\le j$, and let $a\prec_2b$ mean $a<b\le m-2$ and $a\ne j$, or $b=m-1>a$. Then $U_{\prec_1}=H_{j+1}$ and $U_{\prec_2}=S_jY_m$. Let $E$ consist of the pairs $(l,l+1)$ with $0\le l\le m-3$, $l\ne j$, with coefficient $1$, and of $(j,m-1)$ with coefficient $c$. On $H_{j+1}$ the character $\lambda_{E,c}$ is $\chi_{j+1,c}$. On $S_jY_m$ it is $\chi'_{j,c}$, because $(s\,y(x))_{j,m-1}=(sx)_j=x_j$ when the row $j$ of $s$ is $e_j^T$. If $j=m-2$ the two relations coincide. Otherwise we apply Lemma~\ref{w39:exchange} with $p=j$, $q=m-1$ and $W=\{j+1,\dots,m-2\}$. The union of the two relations is the natural order $<$. \ref{w39:X1} holds by definition. \ref{w39:X3}: if $x\le j<w\le y$ then $x\prec_1y$, and a pair $(x,x+1)$ with $x\le j<x+1$ has $x=j$ and is not in $E$, so only $(j,m-1)$ occurs. \ref{w39:X4}: if $x\le w$ and $y\ge m-1$ then $y=m-1$ and $x\prec_2y$, and $(x,m-1)\in E$ only for $x=j$. \ref{w39:X5}: the pairs $(l,l+1)$ are covering pairs of $<$. An element $b$ with $j\prec_1b\prec_1m-1$ would satisfy $b>j$ and $b\le j$. An element $b$ with $j\prec_2b$ equals $m-1$. Lemma~\ref{w39:exchange} applies since $P_m\supseteq U_m=U_\prec$.
\end{proof}

\begin{proposition}\label{w39:propM}
For every $n\ge1$ and every representation $\tau$ of $P_n$,
\[
A(\tau)=B_{n-1}(\tau)+\sum_{i=0}^{n-2}(g-1)g^{\,n-2-i}B_i(\tau).
\]
\end{proposition}

\begin{proof}
For $n=1$ both sides equal $\dim\tau$. Let $n\ge2$ and assume the claim for $n-1$. Put $m=n-1$ and let $\theta(y(x))=\psi_0(x_{n-2})$, a character of $Y_n$ whose stabilizer in $G_{m}$ for the action $\mu\mapsto\mu h^{-1}$ on row vectors is $P_m$. Clifford theory as in the proof of Lemma~\ref{w39:orbits}, now with the transitive action of $G_m$ on nonzero row vectors, gives $\tau|_{G_m}\cong\tau^{Y_n}\oplus\operatorname{Ind}_{P_m}^{G_m}\sigma$ with $\sigma=\tau^{(Y_n,\theta)}$, a representation of $P_m$. Since $\psi^{(n-1)}_n$ is trivial on $Y_n$, $\mathcal W(\tau^{Y_n})=B_{n-1}(\tau)$. For $0\le i\le n-2$ the character $\psi^{(i)}_n$ restricts to $\theta$ on $Y_n$ and to $\psi^{(i)}_m$ on $U_m$, so $B_i(\tau)=B_i(\sigma)$, the right side computed in $P_m$.

By Lemmas~\ref{w39:indP}, \ref{w39:DX} and~\ref{w39:orbits},
\[
\mathcal W(\operatorname{Ind}_{P_m}^{G_m}\sigma)=\sum_cD_{0,c}(\sigma)+\sum_{j=0}^{m-2}\sum_cX_{j,c}(\sigma)=(g-1)A(\sigma)+\bigl(A(\sigma)-B_{m-1}(\sigma)\bigr),
\]
because $H_0=U_{m-1}$ and $\chi_{0,c}=\psi_{m-1}$. By the induction hypothesis $gA(\sigma)-B_{m-1}(\sigma)=(g-1)B_{m-1}(\sigma)+\sum_{i\le m-2}(g-1)g^{m-1-i}B_i(\sigma)=\sum_{i=0}^{m-1}(g-1)g^{m-1-i}B_i(\sigma)$. Adding $B_{n-1}(\tau)$ and using $B_i(\sigma)=B_i(\tau)$ and $m-1-i=n-2-i$ gives the claim.
\end{proof}

For a representation $\pi$ of $G_N$ and $0\le i\le N-1$ put $W_i(\pi)=\dim\pi^{(U_N,\psi^{(i)}_N)}$, so that $W_0(\pi)=\Tr\pi(e_U)$.

\begin{theorem}[column expansion]\label{w39:thmB}\label{w2x:der}
Let $1\le k<N$. For every representation $\pi$ of $G_N$,
\[
\Tr\pi(e')=\sum_{i=0}^{k-1}\bigl(g^{k-i}-g^{k-i-1}\bigr)W_i(\pi)+W_k(\pi).
\]
\end{theorem}

\begin{proof}
Put $n=k+1\le N$ and embed $G_n\subseteq G_N$ in the upper left corner. Let $Q\subseteq U_N$ be the group of $u\in U_N$ whose upper left $n\times n$ block is $1$, and $\psi_Q(z)=\psi_0(\sum_{l=n-1}^{N-2}z_{l,l+1})$, a character of $Q$. For $p\in P_n$ the conjugate $pzp^{-1}$ replaces the upper right block $X$ of $z$ by $pX$ and keeps the lower right block, and the row $n-1$ of $pX$ equals that of $X$, so $P_n$ normalizes $Q$ and fixes $\psi_Q$. Hence $\tau=\pi^{(Q,\psi_Q)}$ is a representation of $P_n$. Every $u\in U_N$ factors uniquely as $u=u_0z$ with $u_0\in U_n$ and $z\in Q$, and $u_0$ is the upper left block of $u$. Thus $U'=U_kQ$ with $U_k\subseteq G_k\subseteq P_n$, and $\psi'(u_0z)=\psi_k(u_0)\psi_Q(z)$, because the row $k=n-1$ of $u_0$ is $e_k^T$. Hence $\Tr\pi(e')=\dim\tau^{(U_k,\psi_k)}=A(\tau)$. In the same way $\psi^{(i)}_N(u_0z)=\psi^{(i)}_n(u_0)\psi_Q(z)$ for $0\le i\le k$, since the omitted entry $(i-1,i)$ lies in the upper left block, so $W_i(\pi)=B_i(\tau)$. Proposition~\ref{w39:propM} with $n=k+1$ and $(g-1)g^{k-1-i}=g^{k-i}-g^{k-i-1}$ ends the proof.
\end{proof}

\begin{corollary}[Hypothesis~\ref{w2:H-wt}]\label{w39:Hwt}\label{w2x:idem}\label{w2x:cmp}
Let $(N,k,s)=(r,k,r_S)$ with $k\ge1$. The element $e_J$ is a self-adjoint idempotent with $e_J(1)=g^{k-N(N-1)/2}>0$, and for every family $(E_c)$ of orthogonal central idempotents of $\C[G_N]$ with $\sum_cE_c=1$, in particular for the label idempotents of Section~\ref{w2:sec:irr},
\[
(1-1/g)\,\mathrm{wt}(e_U,c)\le\mathrm{wt}(e_J,c)\qquad\text{for every }c.
\]
\end{corollary}

\begin{proof}
$U_J$ is a subgroup (Lemma~\ref{w39:pattern}, it is $U_{\prec_s}$) defined in $U_N$ by $k$ vanishing entries above the diagonal, so $e_J(1)=\lvert U_J\rvert^{-1}=g^ke_U(1)$. By Theorems~\ref{w39:thmA} and~\ref{w39:thmB} and $W_i\ge0$, $\Tr\pi(e_J)\ge(g^k-g^{k-1})W_0(\pi)=g^k(1-1/g)\Tr\pi(e_U)$ for every irreducible $\pi$. For $a\in\C[G_N]$, $(aE_\pi)(1)=\lvert G_N\rvert^{-1}\dim\pi\,\Tr\pi(a)$. Summing over the $\pi$ with $E_\pi E_c=E_\pi$ gives $(e_JE_c)(1)\ge g^k(1-1/g)(e_UE_c)(1)$, and dividing by $e_J(1)=g^ke_U(1)$ gives the claim.
\end{proof}

\begin{lemma}\label{w2x:trace}
Let $a\in\C[G_N]$ be an idempotent and $c$ a label. Then
\[
(aE_c)(1)=\frac1{\lvert G_N\rvert}\sum_{\pi\text{ of label }c}\dim\pi\cdot\operatorname{rk}\pi(a)\ \ge0,\qquad\sum_c(aE_c)(1)=a(1).
\]
In particular Lemma~\ref{w2:wtprob} holds for every idempotent $a$ with $\operatorname{Re}a(1)>0$, self-adjoint or not.
\end{lemma}

\begin{proof}
The central idempotent of $\pi$ is $E_\pi=\frac{\dim\pi}{\lvert G_N\rvert}\sum_x\overline{\chi_\pi(x)}\,x$, and $\overline{\chi_\pi(x^{-1})}=\chi_\pi(x)$, so $(aE_\pi)(1)=\sum_xa(x)E_\pi(x^{-1})=\frac{\dim\pi}{\lvert G_N\rvert}\Tr\pi(a)$. The trace of an idempotent operator is its rank. Summing over the $\pi$ of label $c$ gives the first identity. Summing over all $c$ gives $\sum_\pi\frac{\dim\pi}{\lvert G_N\rvert}\Tr\pi(a)$, which is the trace of $a$ in the left regular representation divided by $\lvert G_N\rvert$, that is $a(1)$.
\end{proof}

\subsubsection{Rows, labels and the cylinder idempotents}\label{w2x:sec:rows}

Let $(\tau)$ be the orthonormal eigenbasis of $\mathcal H$ of Section~\ref{w2:sec:irr}. For $k\in\{0,1\}$ let $\mathcal G_k=\prod_{w\in\Pi_k}\GL_N(\F_{g_k})$ be the quotient of the level at $g_k$, and for a set $S$ of label tuples at $k$ let $\Pi_S=\sum_{t\in S}\bigotimes_{w\in\Pi_k}E_{t_w}$, a central idempotent of $\C[\mathcal G_k]$. The level idempotent factors as $e=e^{(k)}\otimes\bigotimes_{w\in\Pi_k}e_{k,w}$, with $e^{(k)}$ supported away from $g_k$, and we put $e_S=e\,\Pi_S$.

\begin{hypothesis}[label dictionary]\label{w2x:H-lab}
(i) For $k\in\{0,1\}$ and every set $S$ of label tuples at $k$, the space of automorphic forms of Section~\ref{w2:sec:statements} defined with $e_S$ in place of $e$ has dimension $\#\{\tau:\mathrm{lab}_k(\tau)\in S\}$. (ii) There is an injective map $\iota$ from an index set of the rows of $R$ to the eigenbasis whose image is the set of $\tau$ whose label tuples at $g_0$ and at $g_1$ are both good, and $\mathrm{lab}_k(\iota(\sigma))$ is the label tuple of the row $\sigma$.
\end{hypothesis}

Indeed the mutually orthogonal central idempotents $E_t$ indexed by pairs of label tuples sum to $1$ and commute with $e$. Hence $\mathcal H$ is the orthogonal direct sum of the spaces $R(E_t)\mathcal H$. An orthonormal basis adapted to this sum gives (i) by taking dimensions. For (ii), suppose each row basis vector lies in a good summand and every good summand lies in the row space. Then the row space is precisely the sum over good $t$. Choosing a bijection of the two bases separately in each summand gives the required row indices.

\begin{lemma}\label{w2x:DS}
The element $e_S$ is a self-adjoint idempotent, and $D^*(e_S)=D^*\,w_k(S)$, where $D^*(e')$ denotes the product of the formal degree, the volume of $[G]$ and $\operatorname{Re}e'(1)$.
\end{lemma}

\begin{proof}
$\Pi_S$ is a self-adjoint central idempotent of $\C[\mathcal G_k]$, so it commutes with $\bigotimes_we_{k,w}$, and the product of two commuting self-adjoint idempotents is a self-adjoint idempotent. The coefficient at $1$ of a tensor product is the product of the coefficients at $1$, so
\[
e_S(1)=e^{(k)}(1)\sum_{t\in S}\prod_{w\in\Pi_k}(e_{k,w}E_{t_w})(1)=e(1)\sum_{t\in S}\prod_{w\in\Pi_k}\frac{(e_{k,w}E_{t_w})(1)}{e_{k,w}(1)}.
\]
By Lemma~\ref{w2x:trace} every factor is real, and $e_{k,w}(1)>0$ by Corollary~\ref{w39:Hwt}, so the right side is $e(1)\,w_k(S)$ and $D^*(e_S)=D^*w_k(S)$.
\end{proof}

\subsubsection{The multiplicity formula}\label{w2x:sec:count}

We prove Hypotheses~\ref{w2:H-cnt} and~\ref{w2:H-DIM} from three analytic statements about the pre-trace formula, which are not proved here. Let $\varphi_\infty(y)=d_\infty\prod_{v}\overline{\langle\pi_v(y_v)w_v,w_v\rangle}$, where $w_v$ is a unit lowest $K$-type vector of the discrete series $\pi_v$ at the real place $v$ and $d_\infty$ is the product of the formal degrees. For a self-adjoint idempotent $e'$ of the Hecke algebra of $G(\Af)$ put $\varphi=\varphi_\infty\otimes e'$, let $R(\varphi)$ be its right convolution operator on $L^2([G])$, let $D(e')$ be the dimension of the space of Section~\ref{w2:sec:statements} defined with $e'$, and let $D^*(e')=d_\infty\vol([G])\operatorname{Re}e'(1)$. Let $\mathfrak F$ be a fundamental domain for $G(F)$ in $G(\A)$ and $n=\dim V=r$.

\begin{hypothesis}[pre-trace inputs, not proved here]\label{w2x:H-pre}
There are absolute constants $c_1>0$, $c_b>0$, and for every $A\ge1$ constants $n_0$, $C_A$ and $A_1$, such that the following holds for every datum with $n\ge n_0$ whose level and Harish-Chandra parameters satisfy the bounds of the construction with the constant $A$, whose idempotent $e'$ is of positive type and supported on a compact open subgroup that is unipotent modulo an auxiliary prime split in $E$ and larger than $4n$, and whose Harish-Chandra parameters have gap at least $h$ at every real place, with
\begin{equation}\label{w2x:gap}
h\ge C_A\,n^{19}f^5\log(b\,\mathrm{rd}_E),\qquad\log f\le\sqrt n,\qquad\log\log b\le c_b\,nf\log n.
\end{equation}
\begin{enumerate}[label=(\roman*),leftmargin=*]
\item $R(\varphi)$ is an orthogonal projection of rank $D(e')$.
\item For almost every $x\in\mathfrak F$ the series $\sum_{\gamma\in G(F)}\varphi(x^{-1}\gamma x)$ converges absolutely, and its sum is $\sum_j\lvert F_j(x)\rvert^2$ for every orthonormal basis $(F_j)$ of the range of $R(\varphi)$.
\item There is a measurable $M\subset\mathfrak F$ such that, with $\eta=\exp(-c_1nf\log n)$,
\[
\begin{gathered}
\vol(\mathfrak F\setminus M)\le\eta\vol(\mathfrak F),\\
\Bigl\lvert\int_M\sum_{\gamma\ne1}\varphi(x^{-1}\gamma x)\,dx\Bigr\rvert\le A_1\eta D^*(e'),\\
\Bigl\lvert\int_{\mathfrak F\setminus M}\sum_{\gamma}\varphi(x^{-1}\gamma x)\,dx\Bigr\rvert\le A_1e^{-n^2f}D(e').
\end{gathered}
\]
\end{enumerate}
\end{hypothesis}

\begin{theorem}\label{w2x:count}
Under Hypothesis~\ref{w2x:H-pre} put $c=c_1/2$. There is $n_1(A)$ such that every datum as in Hypothesis~\ref{w2x:H-pre} with $n\ge n_1(A)$ satisfies $\lvert D(e')-D^*(e')\rvert\le\exp(-c\,nf\log n)\,D^*(e')$.
\end{theorem}

\begin{proof}
Write $D=D(e')$, $D^*=D^*(e')$ and $K(x)=\sum_\gamma\varphi(x^{-1}\gamma x)$. By (i) and (ii), integrating $\sum_j\lvert F_j\rvert^2$ over $\mathfrak F$ gives $D=\int_{\mathfrak F}K$. The term $\gamma=1$ is $\varphi(1)=d_\infty e'(1)$, because $\lVert w_v\rVert=1$. Since $e'$ is a self-adjoint idempotent, $e'(1)=\lVert e'\delta_1\rVert^2$ is real and nonnegative, so $\varphi(1)\vol(\mathfrak F)=D^*$. Splitting $\mathfrak F=M\cup(\mathfrak F\setminus M)$ gives
\[
D-D^*=-\varphi(1)\vol(\mathfrak F\setminus M)+\int_M\sum_{\gamma\ne1}\varphi(x^{-1}\gamma x)\,dx+\int_{\mathfrak F\setminus M}K,
\]
and $\varphi(1)\vol(\mathfrak F\setminus M)\le\eta D^*$, so by (iii) $\lvert D-D^*\rvert\le(1+A_1)\eta D^*+A_1\theta D$ with $\theta=e^{-n^2f}$. Take $n$ so large that $A_1\theta\le\frac12$ and $n\ge c_1\log n$. Then $D\le2(1+(1+A_1)\eta)D^*\le2(2+A_1)D^*$, and $\theta\le\eta$, so $\lvert D-D^*\rvert\le B\eta D^*$ with $B=1+A_1+2A_1(2+A_1)$. Finally $B\eta\le\eta^{1/2}=\exp(-c\,nf\log n)$ once $c\,nf\log n\ge\log B$.
\end{proof}

\begin{corollary}\label{w2x:cnt}
Assume Hypotheses~\ref{w2x:H-lab} and~\ref{w2x:H-pre}, and that the construction satisfies \eqref{w2x:gap} with $n=r$ for $e$ and for every $e_S$, the unipotent condition for $e_S$ being imposed at the prime $g_{1-k}$. Then for $x$ large Hypothesis~\ref{w2:H-cnt} holds with $\varepsilon_C=\exp(-c\,rf\log r)\le\frac12$, and Hypothesis~\ref{w2:H-DIM} holds with the same $c$.
\end{corollary}

\begin{proof}
The set $S$ of all tuples has $e_S=e$ and $D(e)=D_{\mathrm{all}}$, so Theorem~\ref{w2x:count} gives Hypothesis~\ref{w2:H-DIM} and the first clause of Hypothesis~\ref{w2:H-cnt}. For general $S$, Lemma~\ref{w2x:DS} shows that $e_S$ is a self-adjoint idempotent with $D^*(e_S)=D^*w_k(S)$, and Hypothesis~\ref{w2x:H-lab}(i) gives $D(e_S)=\#\{\tau:\mathrm{lab}_k(\tau)\in S\}$. Theorem~\ref{w2x:count} applied to $e_S$ is the second clause. Since $r\ge100$ and $f\ge x^2$, $\exp(-c\,rf\log r)\le\frac12$ for $x$ large.
\end{proof}

\subsubsection{Sizes}\label{w2x:sec:sizes}

\begin{lemma}\label{w2x:sizes}
Let $0<\varepsilon<\frac1{52}$, $J\ge1$ and $\varepsilon_z=\varepsilon/J$. Suppose $x^{2+8\varepsilon}\le n$, $x^{9/2+10\varepsilon}\le k\le2x^{9/2+10\varepsilon}$, $m\le2n\le k$, $g_0,g_1\ge(kb)^8$ and $\min(q,d)\ge\exp(x^{\varepsilon_z/6})$. Then \eqref{w2:sizes} holds with $\varepsilon'=\varepsilon_z$, $C_s=6$, $N=r=k+m-1$ and $f=2n$, for $x\ge4$ large in terms of $\varepsilon$.
\end{lemma}

\begin{proof}
$g_k\ge(kb)^8\ge b^8=e^{8x^\varepsilon}\ge e^{8x^{\varepsilon_z}}$. Next $f=2n\ge x^2$, and $f\le k\le2x^{9/2+10\varepsilon}\le x^5$ since $\varepsilon<\frac1{52}$ and $x$ is large, so $f\le x^6$ and $m\le f\le x^6$. The rank satisfies $k\le r\le2k\le4x^{9/2+10\varepsilon}\le x^6$, and $r\ge k\ge x^{9/2}\ge100$. The last condition is the hypothesis on $q$ and $d$.
\end{proof}



\subsection{Near orthogonality of Bessel Poincar\'e series}\label{w2:sec:S}

\subsubsection{Assembly}

\begin{proposition}\label{w2:assembly}
Let $C_0\ge1$, $\Gamma_0>1$, $D^*>0$, $\eta_0\ge0$, $\xi\ge0$, $\eta_1\ge0$, and assume \eqref{w2:window}. Let $K_{uu'}=K^0_{uu'}+K^{\deg}_{uu'}$ with
\begin{enumerate}[label=(\roman*),leftmargin=*]
\item $K^{\deg}_{uu'}=0$ for $u\ne u'$,
\item $\lVert K^{\deg}_{uu}\rVert_\HS^2\le\xi d^S_u$,
\item $\lVert K^0_{uu'}\rVert_\HS^2\le(1+\eta_0)d^S_ud^S_{u'}/D^*$ for $u\ne u'$,
\item $\lVert K^0_{uu}-1\rVert_\HS^2\le\Gamma_0(d^S_u)^2/D^*$,
\item $\lvert D_{\mathrm{all}}-D^*\rvert\le\eta_1D^*$.
\end{enumerate}
Then the conclusion of Theorem~\ref{w2:S} holds with these $\eta_0,\eta_1$ and $e_0=2C_0/\Gamma_0+2\xi$.
\end{proposition}

\begin{proof}
Off the diagonal use (i) and (iii). On the diagonal, $\lVert A+B\rVert^2\le2\lVert A\rVert^2+2\lVert B\rVert^2$ with $A=K^0_{uu}-1$, $B=K^{\deg}_{uu}$ and (iv), (ii) give $\lVert K_{uu}-1\rVert_\HS^2\le2\Gamma_0(d^S_u)^2/D^*+2\xi d^S_u$, and \eqref{w2:window} gives $\Gamma_0(d^S_u)^2/D^*\le(C_0/\Gamma_0)d^S_u$.
\end{proof}

\subsubsection{The pair bound}

Unfolding $\langle\mathcal P_uT,\mathcal P_{u'}T'\rangle$ as in Proposition~\ref{w2:unfold}, with $\theta=\overline{P_{T}}$, writes $K_{uu'}$ as a sum over $S_u(F)\backslash G(F)/S_{u'}(F)$ of archimedean orbital integrals of $W_u\otimes\overline{W_{u'}}$. We organize the terms as follows, and the next theorem is the Hilbert space lemma that closes the off-diagonal clause.

\begin{definition}[term families]\label{w2:terms}
A term family consists of a countable set of terms $t$, each with an operator $\mathrm{op}(t)$ in a Hilbert space of operators, a vector $\mathrm{vec}(t)$ in a Hilbert space $H_1$, a row $\mathrm{row}(t)$, a Gram value $G(t)$ and two flags, $\mathrm{isId}(t)$ and $\mathrm{gen}(t)$; a finite set $I$ with a map $\iota$ from rows to $I$; a hermitian matrix $(\Gamma_{ij})_{i,j\in I}$ and a number $\Gamma_1>0$; for pairs of terms in one row a corank $s_0(t,t')\in\mathbb Z_{\ge0}$ and a flag $\mathrm{scal}(t,t')$. Put $\mathrm{id}=\{\mathrm{isId}\}$, $\mathrm{deg}=\{G=0,\neg\mathrm{isId}\}$, $\mathrm{ng}=\{G\ne0,\neg\mathrm{gen}\}$, $\mathrm{gen}=\{G\ne0,\mathrm{gen}\}$, $\mathrm{part}(S)=\sum_{t\in S}\mathrm{op}(t)$, $v_i(S)=\sum_{t\in S,\ \iota(\mathrm{row}\,t)=i}\mathrm{vec}(t)$. For a set $Q$ of pairs in one row put $\Sigma(Q)=\Gamma_1\sum_Q\operatorname{Re}\langle\mathrm{vec}(t),\mathrm{vec}(t')\rangle$ and $\lvert\Sigma\rvert(Q)=\Gamma_1\sum_Q\lvert\langle\mathrm{vec}(t),\mathrm{vec}(t')\rangle\rvert$. Let $\mathrm{big}(S)$ be the pairs in $S\times S$ in one row with $s_0\ge2$, $\mathrm{one}(S)$ those with $s_0=1$, split into $\mathrm{sc}(S)$ (with $\mathrm{scal}$) and $\mathrm{nsc}(S)$.
\end{definition}

\begin{theorem}[pair bound]\label{w2:pair}
Assume $\sum_t\lVert\mathrm{op}(t)\rVert<\infty$, $\sum_t\lVert\mathrm{vec}(t)\rVert<\infty$, and
\begin{enumerate}[label=(\alph*),leftmargin=*]
\item $\lVert\mathrm{part}(S)\rVert^2=\operatorname{Re}\sum_{i,j}\Gamma_{ij}\langle v_i(S),v_j(S)\rangle$ for $S=\mathrm{gen}$ and $S=\mathrm{ng}$,
\item $P_g\Gamma_1\cdot1-(\Gamma_{ij})$ is positive semidefinite for some $P_g\ge0$,
\item $\langle\mathrm{vec}(t),\mathrm{vec}(t')\rangle=0$ if $\iota(\mathrm{row}\,t)=\iota(\mathrm{row}\,t')$ and $\mathrm{row}\,t\ne\mathrm{row}\,t'$,
\item $s_0(t,t')\ge1$ if $t,t'$ lie in one row and $G(t),G(t')\ne0$,
\item $\mathrm{part}(\mathrm{id})=0$ and $\mathrm{part}(\mathrm{deg})=0$,
\item with $dd\ge0$: $\lvert\Sigma\rvert(\mathrm{big}(\mathrm{gen}\cup\mathrm{ng}))\le\varepsilon\,dd$, $\lvert\Sigma\rvert(\mathrm{nsc}(\mathrm{gen}))\le\varepsilon\,dd$, $\Sigma(\mathrm{sc}(\mathrm{gen}))\le A\,dd$ and $\Sigma(\mathrm{one}(\mathrm{ng}))\le\varepsilon'\,dd$.
\end{enumerate}
Then $\lVert\sum_t\mathrm{op}(t)\rVert^2\le\bigl(\sqrt{P_g(A+2\varepsilon)}+\sqrt{P_g(\varepsilon+\varepsilon')}\bigr)^2dd$.
\end{theorem}

\begin{proof}
By (b), $\operatorname{Re}\sum_{ij}\Gamma_{ij}\langle a_i,a_j\rangle\le P_g\Gamma_1\sum_i\lVert a_i\rVert^2$ for all vectors $a_i$. With (a), $\lVert\mathrm{part}(S)\rVert^2\le P_g\Phi(S)$ where $\Phi(S)=\Gamma_1\sum_i\lVert v_i(S)\rVert^2$. Expanding $\lVert v_i(S)\rVert^2$ and using (c), only pairs in one row remain, so $\Phi(S)=\Sigma(\text{pairs of }S\text{ in one row})$. By (d), $\Phi(\mathrm{gen})=\Sigma(\mathrm{big}(\mathrm{gen}))+\Sigma(\mathrm{nsc}(\mathrm{gen}))+\Sigma(\mathrm{sc}(\mathrm{gen}))\le(A+2\varepsilon)dd$ and $\Phi(\mathrm{ng})=\Sigma(\mathrm{big}(\mathrm{ng}))+\Sigma(\mathrm{one}(\mathrm{ng}))\le(\varepsilon+\varepsilon')dd$, using $\Sigma\le\lvert\Sigma\rvert$ and monotonicity of $\lvert\Sigma\rvert$. By (e), $\sum_t\mathrm{op}(t)=\mathrm{part}(\mathrm{gen})+\mathrm{part}(\mathrm{ng})$, and the triangle inequality ends the proof.
\end{proof}

\begin{lemma}\label{w2:partition}
For $g\ge3$ and every $k$, $\prod_{i=1}^k(1-g^{-i})^{-1}\le1+3/g$.
\end{lemma}

\begin{proof}
With $y=1/g$, $\prod_{i\ge1}(1-y^i)^{-1}=\sum_{n\ge0}p(n)y^n$, and $p(n)\le2^{n-1}$ for $n\ge1$ because sorting the parts of a composition gives every partition. So the product is at most $1+y/(1-2y)\le1+3y$ for $y\le1/3$.
\end{proof}

Lemma~\ref{w2:partition} is the estimate behind the value $P_g=(1+3/g_1)^f(1+3/g_2)^f$ in Hypothesis~\ref{w2:H-KE} below: the matrix $(\Gamma_{ij})$ contains the finite group factors at the typed places, and its diagonal dominance costs a $p$-adic volume factor of the form $\prod_i(1-g^{-i})^{-1}$ at each of them.

\begin{lemma}[closing numerics]\label{w2:close}
Let $b\ge16$, $0\le P_g\le e^{6b^{-7}}$, $A\le(1+b^{-7})e^{b^{-14}}$ and $0\le\varepsilon,\varepsilon'\le b^{-14}$. Then $\bigl(\sqrt{P_g(A+2\varepsilon)}+\sqrt{P_g(\varepsilon+\varepsilon')}\bigr)^2\le1+b^{-6}$.
\end{lemma}

\begin{proof}
Put $t=b^{-7}\le16^{-7}$. Then $P_g(A+2\varepsilon)\le e^{6t}(1+t)e^{t^2}+2e^{6t}t^2\le1+8t$ and $P_g(\varepsilon+\varepsilon')\le2e^{6t}t^2\le3t^2$. The square is at most $1+8t+3t^2+2\sqrt{3t^2(1+8t)}\le1+12t$, and $12t=12b^{-1}b^{-6}\le b^{-6}$.
\end{proof}

\begin{corollary}[off-diagonal and diagonal clauses from the row estimates]\label{w2:offdiag}
Let $b\ge16$. Suppose that for each pair of words the terms of $K_{uu'}$ form a term family with (a) to (d) of Theorem~\ref{w2:pair}, with $P_g\le e^{6b^{-7}}$, and with $\mathrm{part}(\mathrm{deg})=K^{\deg}_{uu'}$, $\mathrm{part}(\mathrm{id})=0$ for $u\ne u'$ and $\mathrm{part}(\mathrm{id})=1$ for $u=u'$. Suppose the row estimates (f) hold with $\varepsilon,\varepsilon'\le b^{-14}$ and
\[
dd=d^S_ud^S_{u'}/D^*,\qquad A\le(1+b^{-7})e^{b^{-14}}\ (u\ne u'),\qquad A\le\gamma_B\ (u=u'),
\]
with $\gamma_B\ge1$ and $\Gamma_0\ge6\gamma_B$. Then (iii) of Proposition~\ref{w2:assembly} holds with $\eta_0=b^{-6}$, and (iv) holds.
\end{corollary}

\begin{proof}
For $u\ne u'$ apply Theorem~\ref{w2:pair} to the terms other than the degenerate ones, whose sum is $K^0_{uu'}$, and Lemma~\ref{w2:close}. For $u=u'$ apply Theorem~\ref{w2:pair} to the terms other than the identity and degenerate ones, whose sum is $K^0_{uu}-1$. With $(\sqrt a+\sqrt b)^2\le2a+2b$ the bound is $2P_g(\gamma_B+4b^{-14})(d^S_u)^2/D^*\le6\gamma_B(d^S_u)^2/D^*\le\Gamma_0(d^S_u)^2/D^*$.
\end{proof}

\subsubsection{The inputs from the construction}\label{w2:sec:Sinputs}

Put $r_S=m/2-1$, the rank parameter of a row. The ambient rank remains $r=k+m-1$. Proposition~\ref{w2:assembly} and Corollary~\ref{w2:offdiag} reduce Theorem~\ref{w2:S} to the following statements about the concrete Gram operators.

\begin{hypothesis}[S-KE, kernel expansion]\label{w2:H-KE}
For every pair of words the unfolding of $\langle\mathcal P_uT,\mathcal P_{u'}T'\rangle$ converges absolutely and produces a term family with (a) to (d) of Theorem~\ref{w2:pair}, $P_g=(1+3/g_1)^f(1+3/g_2)^f$, $\mathrm{part}(\mathrm{id})=0$ and $\mathrm{part}(\mathrm{deg})=0$ for $u\ne u'$, $\mathrm{part}(\mathrm{id})=1$ for $u=u'$.
\end{hypothesis}

\begin{hypothesis}[S-MR, main row]\label{w2:H-MR}
$\Sigma(\mathrm{sc}(\mathrm{gen}))\le(1+\delta)e_1\,d^S_ud^S_{u'}/D^*$ for $u\ne u'$, with $0\le\delta\le b^{-7}$ and $1\le e_1\le\exp(12f2^{-(k-2r_S-2)})$; and $\Sigma(\mathrm{sc}(\mathrm{gen}))\le\gamma_B(d^S_u)^2/D^*$ for $u=u'$, with $\gamma_B=\exp(C_Bm^5f\log b)$.
\end{hypothesis}

\begin{hypothesis}[S-LR, loops]\label{w2:H-LR}
Each of the sums $\lvert\Sigma\rvert(\mathrm{big}(\mathrm{gen}\cup\mathrm{ng}))$ and $\lvert\Sigma\rvert(\mathrm{nsc}(\mathrm{gen}))$ is at most $\exp(-\tfrac13kf\log k)\,dd$.
\end{hypothesis}

\begin{hypothesis}[S-NG, nongeneric terms]\label{w2:H-NG}
$\Sigma(\mathrm{one}(\mathrm{ng}))\le2^{-(k-2r_S-5)}\exp(C_4m^5f\log b)\,dd$.
\end{hypothesis}

\begin{hypothesis}[S-DG, degenerate terms]\label{w2:H-DG}
$\lVert K^{\deg}_{uu}\rVert_\HS^2\le\exp(-kf\log k)\,d^S_u$.
\end{hypothesis}

\begin{hypothesis}[S-DIM, multiplicity formula]\label{w2:H-DIM}
$\lvert D_{\mathrm{all}}-D^*\rvert\le\exp(-c\,rf\log r)\,D^*$ for an absolute $c>0$.
\end{hypothesis}

\begin{hypothesis}[S-SZ, sizes]\label{w2:H-SZ}
$b=e^{x^\varepsilon}\ge16$, $(kb)^8\le g_1,g_2$, $f\le k^8b$, $70\log b\le k$, $k\ge8\log\Gamma_0$, $C_\Gamma\ge\max(C_B,C_4)+1$, the quantities $\exp(-\tfrac13kf\log k)$ and $2^{-(k-2r_S-5)}\exp(C_4m^5f\log b)$ are at most $b^{-14}$, $12f2^{-(k-2r_S-2)}\le b^{-14}$, and $c\,rf\log r\ge x^{\varepsilon/8}$.
We also require $32\le k$, $4r_S+12\le k$ and $8C_4m^5f\log b\le k$.
\end{hypothesis}

\begin{lemma}[sizes]\label{w39:sizes}
Let $b=e^{x^\varepsilon}$, and let $m,k,n$ be as in Proposition~\ref{ranks} with $f=2n$, $r_S=m/2-1$. For $x$ large in terms of $\varepsilon$ and $C_4$: $32\le k$, $4r_S+12\le k$, $f\le k$, $70\log b\le k$, $8C_4m^5f\log b\le k$, and consequently
\[
12f2^{-(k-2r_S-2)}\le b^{-14},\qquad2^{-(k-2r_S-5)}\exp(C_4m^5f\log b)\le b^{-14}.
\]
\end{lemma}

\begin{proof}
By Proposition~\ref{ranks}, $m\le2\sqrt x$, $f=2n\le4x^{2+8\varepsilon}$ and $k\ge x^{9/2+10\varepsilon}$. So $4r_S+12=2m+8\le k$, $f\le k$, $70x^\varepsilon\le k$, and $8C_4m^5f\log b\le1024\,C_4x^{9/2+9\varepsilon}\le k$ for large $x$. From $4r_S+12\le k$ we get $k-2r_S-2\ge k/2+4$ and $k-2r_S-5\ge k/2+1$. Then $12f2^{-(k-2r_S-2)}\le f2^{-k/2}\le e^{-k/5}$ since $f\le k\le e^{(\frac{\log2}2-\frac15)k}$ for $k\ge32$, and $2^{-(k-2r_S-5)}e^{C_4m^5f\log b}\le2^{-k/2}e^{k/8}\le e^{-k/5}$. Finally $e^{-k/5}\le b^{-14}$ because $70\log b\le k$.
\end{proof}

\begin{theorem}\label{w2:Sproof}
Hypotheses~\ref{w2:H-KE} to~\ref{w2:H-SZ} imply Theorem~\ref{w2:S} for $x$ large in terms of $\varepsilon$ and $C_0$.
\end{theorem}

\begin{proof}
By \ref{w2:H-SZ}, $P_g\le\exp(3f(g_1^{-1}+g_2^{-1}))\le\exp(6f(kb)^{-8})\le e^{6b^{-7}}$, $(1+\delta)e_1\le(1+b^{-7})e^{b^{-14}}$, and $\varepsilon=\exp(-\tfrac13kf\log k)$, $\varepsilon'$ of \ref{w2:H-NG} are at most $b^{-14}$. Also $\Gamma_0/\gamma_B=\exp((C_\Gamma-C_B)m^5fL)\ge e^{L}\ge6$ since $L=\log b$. Corollary~\ref{w2:offdiag} gives (iii) with $\eta_0=b^{-6}$ and (iv) of Proposition~\ref{w2:assembly}; \ref{w2:H-KE} gives (i); \ref{w2:H-DG} gives (ii) with $\xi=\exp(-kf\log k)$; \ref{w2:H-DIM} gives (v) with $\eta_1=\exp(-crf\log r)$. Proposition~\ref{w2:assembly} gives the three clauses with $e_0=2C_0/\Gamma_0+2\xi$. With $y=x^{\varepsilon/8}$: $\eta_0=e^{-6y^8}\le e^{-y}$; $\eta_1\le e^{-y}$ by \ref{w2:H-SZ}; $\log\Gamma_0\ge C_\Gamma y^8$ and $k\ge8\log\Gamma_0$ give $2C_0/\Gamma_0+2\xi\le e^{-2y}$ for large $y$, so $\sqrt{e_0}\le e^{-y}$.
\end{proof}

\begin{remark}[the window]\label{w2:rem:window}
Theorem~\ref{w2:S} assumes \eqref{w2:window}. Theorem~\ref{r3:A3} uses its conclusion, so the construction of the words must establish \eqref{w2:window}. The lower half of \eqref{w2:window} is the hypothesis $d^S_u\ge D^*/(C_0\Gamma_0^2)$ of Lemma~\ref{r3:sizecond}.
\end{remark}

\begin{remark}[no circularity]\label{w2:rem:circ}
Lemma~\ref{w2:punz} derives the nonvanishing used in the proof of Theorem~\ref{w2:B} from the diagonal clause of Theorem~\ref{w2:S}. Theorem~\ref{w2:S} is a statement about Poincar\'e series alone and its inputs \ref{w2:H-KE} to \ref{w2:H-SZ} do not involve the period maps, so there is no cycle. Theorem~\ref{r3:A3} uses all three theorems.
\end{remark}

\subsubsection{The error in the multiplicity formula}\label{count:error}

Let $n$ denote the rank in the counting theorem, $f$ the degree of the
totally real field, $D$ the multiplicity and $M$ the identity contribution.
With the Haar measures used for the formal degrees $d_v$, one has
$M=(\prod_{v\mid\infty}d_v)\operatorname{vol}(\mathfrak F)
\operatorname{Re}e(1)$ for the level idempotent $e$. In particular $M\ge0$.
The trace estimate also requires finite dimensionality of the multiplicity
space. In the varying-rank estimate the gap in the weights is at least
\[
  8\cdot10^{14}(C_g+4)n^{19}f^5\log(b\,\mathrm{rd}_E),
\]
where the level parameter $g$ satisfies $1\le g\le b^{C_g}$.
The hypotheses include a common signature $(p,q)$ with $q^2\le n$,
a prime $g_0>4n$ split completely in $E$, and the bounds
$\log f\le\sqrt n$, $\log\log\mathrm{rd}_E\le n$,
$\log\log b\le nf\log n/16$, and $|\lambda_{v,i}|\le b^2$.



\subsubsection{Compactly supported Gram-fibre identities}

Let $E/F$ be a quadratic CM extension with involution $\sigma$. Let $V=E^M$ be equipped with a nonsingular Hermitian form $h_J$ and unitary group $U_J$, and let $\mathcal H_n$ denote Hermitian $n\times n$ matrices. For $X=(X_1,\ldots,X_n)\in V^n$, define $\operatorname{Gram}_J(X)_{ab}=h_J(X_a,X_b)$ and $\mathcal H_n^\times=\{S:\det S\ {\rm is\ a\ unit}\}$. Fix an $F$-basis of $E$. At every place $v$ of $F$, choose additive Haar measure with integral covolume one and transport it to matrix and Hermitian coordinates as follows. The diagonal Hermitian coordinates are $\tfrac12\operatorname{Tr}_{E/F}(S_{ii})$, and the off-diagonal coordinates are the coordinates of $S_{ij}$ in the chosen basis. Write $dX_v^\circ,dS_v^\circ$ for these measures and $L_v$ for the standard local lattice of matrices.

On $\mathcal H_n^\times$, disintegrate $dX_v^\circ$ over $dS_v^\circ$ along $\operatorname{Gram}_J$. Let $\omega_{v,S}^\circ$ be the fibre measures and define
\[
 D_v^\circ(\varphi,S)=\int_{\operatorname{Gram}_J^{-1}(S)}\varphi\,d\omega_{v,S}^\circ,
 \qquad \ell_v^\circ(T)=\omega_{v,T}^\circ(L_v).
\]
Thus
\[
 \int_{\operatorname{Gram}_J^{-1}(K)}\varphi\,dX_v^\circ
 =\int_KD_v^\circ(\varphi,S)\,dS_v^\circ
\]
for every compact $K\subset\mathcal H_n^\times$. In particular, $D_v^\circ(\varphi,\cdot)$ is continuous on that locus for continuous compactly supported $\varphi$. Empty fibres have measure zero.

Let $\psi$ be a nontrivial additive character of $\mathbb A_F$. If $dX_v,dS_v$ are self-dual for the trace pairings used in the Schrödinger model, write $dX_v=c_vdX_v^\circ$ and $dS_v=e_vdS_v^\circ$. Fourier inversion gives
\[
 c_v=e_v=1\quad\hbox{for almost all }v,\qquad
 c_\infty\prod_vc_v=e_\infty\prod_ve_v=1.
\]
Uniqueness of fibre disintegration implies
\[
 D_v^{\rm sd}={c_v\over e_v}D_v^\circ
\]
for the densities taken with $dX_v,dS_v$.

\begin{proposition}\label{prop:normalized-unitary-sw}
Suppose $2n<M$ and $T\in\mathcal H_n(E)$ is nonsingular. Let $\Phi_\infty$ be continuous, rapidly decreasing, smooth and compactly supported, and let $\Phi_v$ be Schwartz--Bruhat, with $\Phi_v=1_{L_v}$ outside a finite set. Put $\Phi=\Phi_\infty\prod_v\Phi_v$. For every Haar measure $\mu$ on $U_J(\mathbb A)$ and every fundamental domain $\mathcal F$ for $U_J(F)$, one has $0<\mu(\mathcal F)<\infty$ and
\[
 \mu(\mathcal F)^{-1}\int_{\mathcal F}
 \sum_{X\in V^n:\ \operatorname{Gram}_J(X)=T}\Phi(h^{-1}X)\,d\mu(h)
 =D_\infty^{\rm sd}(\Phi_\infty,T)\prod_vD_v^{\rm sd}(\Phi_v,T).
\]
The finite product multiplies. Moreover, the fibre measures form an Euler product: $\ell_v^\circ(T)<\infty$ for every finite $v$, $\sum_v|\ell_v^\circ(T)-1|<\infty$, and, for every sufficiently large finite set $S$,
\[
 \int\Phi_\infty\prod_{v\in S}\Phi_v\,d\mathcal M_T
 =D_\infty^\circ(\Phi_\infty,T)\prod_{v\in S}D_v^\circ(\Phi_v,T)
 \prod_{v\notin S}\ell_v^\circ(T),
\]
where $\mathcal M_T$ is the normalized pushforward of $\mu|_{\mathcal F}$ along the rational fibres, and the functions are evaluated on the same orbit variables on both sides.
\end{proposition}

The proof uses only the global convergent-range Siegel--Weil formula as an external analytic theorem. All local normalizations and tails are derived below.

\subsubsection{The normalized global formula}

For the adelic Weil representation $\omega$, write
$\Theta(\Phi,g,h)=\sum_{X\in V^n}\omega(g,h)\Phi(X)$, and let $E(g,s,\Phi)$ be the associated Siegel Eisenstein series. Weil's convergent-range formula \cite[Theorem 5]{Weil1965}, in the unitary formulation also recalled in \cite[Theorem 2.3.1]{LiTheta}, asserts for $M>2n$ that
\[
 {1\over2}\int_{[U_J]}\Theta(\Phi,g,h)\,dh=E(g,(M-n)/2,\Phi),
\]
where $dh$ is Tamagawa measure. Any admissible auxiliary splitting character may be used for the dual pair; the big-cell comparison below removes that choice from the identity at the identity. No signature restriction is needed. If the Witt index $r$ is at most $n$, then $M-r>n$; otherwise $M\ge2r$ and again $M-r>n$. The full unitary Tamagawa covolume is two: $\tau(SU_J)=1$ \cite{Mars1969}, Ono's Main Theorem gives covolume two for the quadratic norm-one torus \cite{Ono1963}, and the relative theorem transports this through the determinant sequence \cite{Ono1965,WeilAdeles}. Hence division by the covolume removes the displayed factor $1/2$. Taking a nonsingular $T$-Fourier coefficient at the identity gives
\[
 \overline I_T(\Phi)=\prod_vW_{v,T}(1,\Phi_v).
\]
Borel's theorem \cite[Theorem 5.6(i)]{borel63} applies because the abelian quotient of $U_J$ is the norm-one torus, which has no nontrivial $F$-rational character. This gives finite covolume. Comparison of fundamental domains and Haar scaling give it for every $\mu,\mathcal F$ in the proposition, and covolume-normalized means are independent of those choices.

In the large Bruhat cell, the Schrödinger action and Fubini give
\begin{equation}\label{eq:whittaker-coarea}
 W_{v,T}(1,\varphi)=D_v^\circ(\varphi,T).
\end{equation}
Indeed, after integrating over the unipotent Hermitian variable, one obtains the Fourier transform at $T$ of $\operatorname{Gram}_*(\varphi\,dX_v^\circ)$; the coarea definition identifies this with $D_v^\circ(\varphi,T)$. The same equality holds with any consistent scaling of input and target measures. For $n=1$, the Hermitian scalar is identified with an element of $F$ by half the trace, and \eqref{eq:whittaker-coarea} is the rank-one Gram-fibre identity.

Varying the test functions determines an identity of locally finite Radon measures. Compact finite-place cylinders are approximated by locally constant functions; compactly supported Borel functions at infinity by smooth functions; and general nonnegative measurable functions by increasing compact truncations and monotone convergence. The scalar identity therefore includes the measurable functions in the eventual Euler-product statement.

\subsubsection{The cofinite error estimate}

It remains to bound $\ell_v^\circ(T)-1$. Begin with $n=1$. At almost every finite place, the trace form in the chosen coordinates is an integral unimodular quadratic form $Q$ in $2M$ variables, while $2$ and the fixed nonzero scalar $t$ are both units. Every residue-field solution is then smooth: otherwise nondegeneracy forces $x=0$, contradicting $Q(0)=0\ne t$. Hensel lifting identifies the local mass with the first-level count:
\[
 \ell_v^\circ(t)=q_v^{1-2M}\#\{x\in k_v^{2M}:Q(x)=t\}.
\]
The character-sum evaluation of this count gives
\[
 N_1=q_v^{2M-1}+q_v^{-1}\sum_{a\in k_v^\times}G(aQ)e(-at).
\]
The Gauss sum of a nondegenerate quadratic form in $2M$ variables has modulus $q_v^M$, so
\[
 |\ell_v^\circ(t)-1|\le q_v^{1-M}\le2q_v^{1-M}.
\]
For $M\ge3$, the error is at most $2q_v^{-2}$.

At good inert places, every unimodular Hermitian lattice has an orthonormal basis. A unit-norm vector exists: otherwise polarization, with $1/2$, would make every Hermitian pairing nonunit, contrary to unimodularity. Such a vector splits off after scaling, using surjectivity of the norm on units, and induction applies to its complement. The split case follows from a basis and its dual basis.

Now split $T\simeq(t)\perp T_1$. Two-stage coarea over the first represented vector gives
\[
 \ell_v^\circ((t)\perp T_1)
 =\delta_v(t)^n\ell_v^\circ(T_1)\ell_v^\circ(t),
\]
where $\delta_v(t)$ is the coordinate Jacobian of multiplication by $t$. At almost every place $\delta_v(t)=1$, and all congruence and Gram determinants are units. The identity
$xy-1=(x-1)(y-1)+(x-1)+(y-1)$ shows that an error $2q_v^{-M}$ for the first factor and $C_1q_v^{-2}$ for the second gives an error $(3|C_1|+2)q_v^{-2}$ for their product. Induction yields a finite set $S_0$ and a constant $C$ such that
\[
 |\ell_v^\circ(T)-1|\le Cq_v^{-2}\qquad(v\notin S_0).
\]
All lattice masses are finite by compactness. Hence $\sum_v|\ell_v^\circ(T)-1|<\infty$. Since $D_v^\circ(\Phi_v,T)=\ell_v^\circ(T)$ whenever $\Phi_v=1_{L_v}$, the finite local densities form a multipliable family for every test function in the proposition. Empty fibres cause no convergence problem.

Combining the normalized global formula, \eqref{eq:whittaker-coarea}, and this tail estimate proves the coordinate version of the proposition in every rank. Finally,
\[
 \biggl({c_\infty\over e_\infty}\biggr)\prod_v{c_v\over e_v}=1,
\]
so replacing $D^\circ$ by $D^{\rm sd}$ leaves the completed product unchanged. The scaled densities have the required continuity and compact-set identities by definition of fibre disintegration. Rank zero is elementary: all relevant spaces are singletons, all tail masses are one, both sides equal $\Phi(0)$, and a self-dual measure on a singleton has mass one.

\subsubsection{The Witt-block form}

The same normalization is compatible with induction on the target rank. Let $J'=(t)\perp J$ and $T'=(t)\perp T$. Witt extension \cite[Theorem 42:17]{OMeara1973} and measurable selection on the locally closed orbit give a section $s(y)$ from the distinguished vector $e$ to $y$ on $q(y)=t$. Define
\[
 \Psi(y)=\int_X\varphi(s(y),X)\,d\omega_T^\circ(X).
\]
Stabilizer invariance makes $\Psi$ independent of the section, and two-stage coarea gives
\[
 \int\varphi\,d\omega_{T'}^\circ
 =\delta_v(t)^n\int\Psi(y)\,d\omega_{(t),1}^\circ(y).
\]
Here $\delta_v(t)=1$ for almost all $v$ and $\delta_\infty(t)\prod_v\delta_v(t)=1$. If $\mathcal M_{J,T}$ is the averaged orbit measure and $\mathcal M^{(1)}_{J',(t)}$ is its rank-one analogue, decomposition of global frames by their distinguished vector gives
\[
 \int\varphi\,d\mathcal M_{J',T'}
 =\Big(\prod_{v\notin S}\ell_v^\circ(T)\Big)
 \int\Psi\,d\mathcal M^{(1)}_{J',(t)}.
\]
For each fixed first vector, the inner integral is the lower-rank orbit integral; stabilizer and fundamental-domain invariance identify it with the lower Euler formula. Nonnegative truncation followed by real/imaginary decomposition permits the interchange. The product formula then cancels $\delta_\infty^n\prod_{v\in S}\delta_v^n$ before passage to the limit, so the block Euler product has constant one.

Congruence of $J$ or $T$ preserves the statement: the induced group conjugacy identifies fundamental domains, while the source and target Jacobians enter in the same ratio and have adelic product one. If $q_J$ does not represent $t$ rationally, the scalar and block global measures vanish. The scalar Euler identity forces an archimedean or a finite local scalar fibre measure to vanish; the block fibre at the same place also vanishes by projecting to the first vector. The absolutely convergent zero-Euler-product statement follows. Peeling off an anisotropic scalar and applying these three operations proves the block-induction form without excluding unrepresented targets.

\subsubsection{Measure of an adapted flag class}

Let $\mathfrak A=(I,I_0,e,x_f,S,\Lambda,x_0,T,T',T_0)$ be an adapted
model. Thus $e$ is compatible with the involutions, $n=n_0+2J$ where
$J=\sum_i j_i$, each $\Lambda_i$ is a full lattice, and the blocks
satisfy
\[
 T_i^*JT_k=(T_i')^*JT_k'=0,\qquad T_i^*JT_k'=0\quad(i\ne k),
\]
together with $T_i^*JT_i'=1$, $T_i^*JT_0=(T_i')^*JT_0=0$, and
$T_0^*JT_0=J_0$. If $L,L_0$ are the two self-dual lattices and
$\Lambda_i^\vee$ is the trace dual, the adapted-lattice condition is
\[
 z\in L\quad\Longleftrightarrow\quad
 z=\sum_iT_iu_i+\sum_iT_i'u_i'+T_0e(y),\qquad
 u_i\in\Lambda_i,\quad u_i'\in\Lambda_i^\vee,\quad y\in L_0,
\]
and $S_k$ is the span of the initial $T$-blocks. Assume $j_i\ge1$.
Fix the common signature $(p,q)$ of $I$. The whole-group and subgroup
gauges are normalized by their Hilbert--Schmidt small-ball asymptotics:
in dimension $d$, the measure of the radius-$r$ ball, divided by $r^d$,
tends to the volume of the Euclidean unit ball. These limits use
$n^2,d_N,(n-2J)^2+2\sum_ij_i^2$, respectively; the compact subgroup
measure is a probability. Fix gauge constants $c=(c_v)$ and
$c_0=(c_{0,v})$ (necessarily positive). Isotropy gives $J\le\min(p,q)$.

Write $f=[F:\mathbf Q]$, $v=v_{\mathrm{inv}}$, and
$d_N=4\sum_{i<k}j_ij_k+2J(n-2J)+\sum_i j_i^2$. Let
$\mathcal H_a$ consist of those $\tau$ with
$0<\tau_1<\cdots<\tau_a<1/2$. Define
\[
 \mathcal L_f(\tau)=
 \prod_{i<k}(\tau_i/\tau_k)^{2j_ij_kf}
                (\tau_i\tau_k)^{2j_ij_kf}
 \prod_i\tau_i^{2j_i(n-2J)f}(\tau_i^2)^{j_i^2f},
 \qquad
 D(\tau)=\mathcal L_f(\tau)\prod_i\frac{2fj_i}{\tau_i}.
\]
Let $\ell_v(m)$ be the same expression with $f=1$, evaluated at the
local determinant ratios
$(d_{v,i+1}(m)/d_{v,i}(m))^{-1/j_i}$.

For a group $\Gamma$ acting on the left, put
\[
 \operatorname{CV}_\Gamma(\mu;A)=
 \inf_D\mu(D),
\]
where $D$ runs through measurable sets meeting every left
$\Gamma$-orbit in $A$. Let $\Gamma,\Gamma_0$ be the flag and residual
lattice stabilizers, and let $A_B$ be the points whose successive flag
sizes lie in $B$. Let $\mathsf V_i$ be the volume of the determinant
slab $1\le d_i<\mathrm e$ modulo $\operatorname{Aut}(\Lambda_i)$, and let
$\mathsf V_0$ be $(\prod_vc_{0,v})$ times the covering covolume of
$\Gamma_0$. Denote by $\iota(x)$ the value of a nonnegative real $x$ in
$[0,\infty]$. Finally, with $U_r$ normalized by
$U_0=1$ and $U_r/U_{r-1}=2^{r-1}(2\pi^r)/(r-1)!$, put
\[
 c_{\mathrm{Iw}}=2^{-d_N/2}\frac{U_pU_q}
 {\bigl(\prod_iU_{j_i}2^{j_i^2/2}\bigr)U_{p-J}U_{q-J}}.
\]

\begin{theorem}
For every adapted model, every $a$ and $(j_i)$ with $j_i\ge1$, every
choice of the above gauges, and every measurable
$B\subseteq\mathcal H_a$,
\begin{equation}\label{eq:adapted-class-bound}
\begin{split}
&\iota\!\left(\prod_vc_v\right)\operatorname{CV}_\Gamma(A_B)\\
&\quad\le \iota\!\left(c_{\mathrm{Iw}}^f(2v)^{fd_N/2}\right)
 \left(\prod_i\iota(2^{fj_i^2})\mathsf V_i\right)\mathsf V_0
 \int_B\iota(D(\tau))\,d\tau .
\end{split}
\end{equation}
\end{theorem}

\begin{proof}[Proof]
The index $a=0$ and the case $A_B=\varnothing$ are included, the latter
by taking the empty covering domain.

\textbf{1. The real-place measure identity.} Let $N,M,K$ be the
unipotent radical, the flag Levi subgroup, and the standard maximal
compact subgroup of the real unitary group. Orthogonal projections
adapted to the flag give $G=NMK$ and show that $um\in K$ with
$(u,m)\in N\times M$ forces $u=1$. Since the ambient group is
reductive, $N$ is nilpotent, and $K$ is compact, the closed subgroups
satisfy the hypotheses of the Haar-decomposition theorem
\cite[Theorem 8.32 and Corollary 8.31]{Knapp}. The modular function is
computed by the determinant of
the adjoint action \cite[Proposition 8.27]{Knapp}.

To evaluate it here, double the flag by duality. If $D_s$ is the
determinant on the $s$-th step and $r_s$ its complex dimension, the
adjoint determinant is
$\prod_sD_s^{r_{s+1}}D_{s+1}^{-r_s}$. Self-duality identifies opposite
steps and yields $\dim F_{2a+1-s}+\dim F_s=n$. Consequently the
pushforward of unipotent Haar measure by $u\mapsto mum^{-1}$ is
$\iota(\ell(m))$ times itself. Hence there is a finite nonzero
$\kappa$ with
\[
 \int\phi\,dg=\kappa\int_N\int_M\iota(\ell(m))
 \int_K\phi(umk)\,dk\,dm\,dn.
\]

The value of $c\kappa$ is obtained from small balls, not from a generic
measure comparison. Put $d_M=(n-2J)^2+2\sum_ij_i^2$ and
$d_L=(p-J)^2+(q-J)^2+\sum_ij_i^2$. Then
$n^2=2d_N+d_M$ and $p^2+q^2=d_N+d_L$. The volume of $K$ is $U_pU_q$,
and the volume of $K\cap M$ is
$U_L=(\prod_iU_{j_i}2^{j_i^2/2})U_{p-J}U_{q-J}$; the powers of $2$
come from the two copies of each isotropic block in the
Hilbert--Schmidt metric. Orthogonal diagonalization of the Lie-algebra
Jacobian gives a unit body of volume
$2^{d_N/2}\operatorname{ballVol}(2d_N+d_M)
 \operatorname{ballVol}(d_L)$. If $\mu_K$ is the probability measure
and $\mu_0$ is the Riemannian gauge on $K$, then
$\mu_0=(U_pU_q)\mu_K$ by Haar uniqueness. Comparing the two small-ball
asymptotics, with a slab on $K\cap M$ squeezed between radii
$r/(1+\eta)$ and $(1+\eta)r$, gives
$c\kappa U_L2^{d_N/2}=U_pU_q$. Thus $c\kappa=c_{\mathrm{Iw}}$. Taking
the finite product over the $f$ places and applying Tonelli gives
\begin{equation}\label{eq:iwasawa-product}
\begin{split}
&\iota\!\left(\prod_vc_v\right)\int\phi\,dg\\
&\quad=\iota(c_{\mathrm{Iw}}^f)\int_{N_G}\int_{M_G}
 \prod_v\iota(\ell_v(m_v))\int_{K_G}\phi(umk)\,dk\,dm\,dn .
\end{split}
\end{equation}

\textbf{2. Integral radical covolume.} Complete the flag by duality and
let $r_b$ and $h_b$ be its graded ranks and heights. Setting
$\tau_i$ equal to the size of the $i$-th graded piece at the identity,
we have $h_i=\tau_i^{2fj_i}$, $h_{2a-i}=\tau_i^{-2fj_i}$, $h_a=1$,
$r_i=r_{2a-i}=j_i$, and $r_a=n-2J$. For $0\le c<2a+1$, define
\[
 A_c=\sum_br_br_{b+c+1},\qquad
 Q_c=\prod_bh_b^{r_{b+c+1}}h_{b+c+1}^{-r_b}.
\]
The $c$-th nilpotent layer consists of homomorphism blocks. If a
biorthogonal block frame changes $\Lambda_b$ to $G_b\Lambda_b$, then
$\operatorname{cov}(\Lambda_b)\prod_\nu|\det G_{\nu,b}|^2
 =v^{fr_b}h_b$. Evaluating the source and target homomorphism
covolumes therefore gives $\operatorname{cov}(L_c)=v^{fA_c}Q_c$.

Now separate real and imaginary coordinates. Conjugation symmetry gives
$\operatorname{cov}(A)\operatorname{cov}(B)
 \le2^{fA_c}\operatorname{cov}(L_c)$. The two covolumes are equal. To
see this, for each rational prime $p$ choose a skew algebraic integer
$\delta$, purely imaginary at every selected embedding, whose real
squared determinant $N$ is prime to $p$. Such a $\delta$ is obtained by
the Chinese remainder theorem away from the finitely many maximal
ideals above $p$. Multiplication by $\delta$ sends $A$ to $B$ and $B$
to $A$. If $T$ denotes this map, integral bases yield
$|\det T|\operatorname{cov}(A)=n_1\operatorname{cov}(B)$ and
$|\det T|\operatorname{cov}(B)=n_2\operatorname{cov}(A)$ with
$n_1n_2=N^{A_c}$. If $\rho$ is the covolume ratio, then for every $p$
there are natural $a,b$ such that $\rho^2b=a$ and $p\nmid ab$. Varying
$p$ proves $\rho=1$. A real fundamental parallelepiped $D_c$ therefore
satisfies $|D_c|\le\sqrt{(2v)^{fA_c}Q_c}$.

The exponential chart has unipotent triangular coordinate changes
between successive layers. Correcting one layer at a time therefore
gives a covering domain with
$\operatorname{CV}_{\Gamma_N}\le\prod_c|D_c|$. Reindexing by ordered
pairs gives $\sum_cA_c=d_N$ and $\prod_cQ_c=\mathcal L_f(\tau)^2$.
It follows that
\begin{equation}\label{eq:radical-covolume}
 \operatorname{CV}_{\Gamma_N}
 \le\iota\!\left((2v)^{fd_N/2}\mathcal L_f(\tau)\right).
\end{equation}

\textbf{3. Levi integration.} The adapted basis gives a measurable
Levi chart $\Theta$ for which
\[
 \operatorname{size}_i(\Theta(z,g_0))
 =(d_i(z_i)/\kappa_i)^{-1/(2fj_i)},\qquad
 \kappa_i=\tau_i^{2fj_i},
\]
and $d\mu_M=(\prod_i2^{fj_i^2})(\prod_vc_{0,v})$ times the pushforward
of the product of the matrix Haar measures with $d\mu_{G_0}$. The
singular locus has measure zero, because $j_i\ge1$ makes the matrix
Haar density vanish there, so restriction to $d_i>0$ loses no mass.
Choose a measurable scale-stable fundamental domain $\mathcal D_i$ for
$\operatorname{Aut}(\Lambda_i)$, using a shortest orbit representative
and a real-coordinate lexicographic tie-break. If $N_i=2fj_i$, then
\begin{equation}\label{eq:slab}
 \int_{\mathcal D_i}H((d_i/\kappa_i)^{-1/N_i})\,dz_i
 =\mathsf V_i\int_0^\infty H(t)\frac{N_i}{t}\,dt .
\end{equation}
Indeed, translation in $\log d_i$ reduces the left side to one slab
with the weight $dt/t$, and $t=\kappa_i\tau^{-N_i}$ has
$|dt/t|=N_i\,d\tau/\tau$. Multiplication over $i$ and covering the
residual factor by a $\Gamma_0$-domain shows that the weighted Levi
covolume is at most
\[
 \left(\prod_i\iota(2^{fj_i^2})\mathsf V_i\right)\mathsf V_0
 \int_{\tau_i>0}F(\tau)\iota\!\left(\prod_i\frac{2fj_i}{\tau_i}\right)
 d\tau .
\]
The residual covolume is nonzero because $\Gamma_0$ is discrete, so
passing to the infimum over residual domains is legitimate. Put
$F=\mathbf1_B\iota(\mathcal L_f)$. On $B$, every factor is positive,
and the last integral is $\int_B\iota(D(\tau))\,d\tau$.

\textbf{4. Assembly.} If $A_B\ne\varnothing$, the size-movement formula
\[
 \operatorname{size}_i(umk)^f
 =\tau_i^f\prod_v\operatorname{locSize}_v(m)_i
\]
and membership of one point in $A_B$ prove $\tau_i>0$. Also
$\mathcal L_f(\operatorname{size}(m))
 =\mathcal L_f(\tau)\prod_v\ell_v(m_v)$. Take
$Z=\{m:\operatorname{size}(m)\in B\}$. A flag lattice stabilizer
fixes $L$ and every $S_k$, hence leaves every graded size unchanged;
in particular, $A_B$ is invariant under its action. The radical estimate
and the size relation allow the base $\mathcal L_f(\tau)$ to be taken outside
the $N$-integral and cancelled against the indicator-weighted Levi
integrand.

More explicitly, fix measurable $N$- and $M$-covering domains and a
measurable section of $K/(K\cap M)$. Form the corresponding measurable
$NMK$ slice. Its orbit meets $A_B$, and \eqref{eq:iwasawa-product}
bounds the flag covering covolume by
$c_{\mathrm{Iw}}^f\mu_N(D_N)\int_{D_M\cap Z}\prod_v\ell_v\,dm$.
Taking the infimum over $D_N$, then using
\eqref{eq:radical-covolume} and the slab computation, and finally
taking the infimum over $D_M$, gives \eqref{eq:adapted-class-bound}.
The zero and infinite $M$-integral cases are handled directly; all
other scalar--infimum interchanges involve finite nonzero constants.
Every real quantity embedded by $\iota$ is nonnegative on the indicated
domain.
\end{proof}

\subsubsection{Definite unimodular Hermitian planes}\label{sec:rank-two-planes}

\paragraph{Notation.} Let $E$ be a totally complex CM field with maximal real
subfield $F$, and let $x\mapsto\bar x$ be complex conjugation. Put
\[
\mathcal O=\mathcal O_E,\qquad f=[F:\mathbb Q],\qquad D=|\Delta_E|,
\qquad \rho=D^{1/(2f)},\qquad Y=\rho^f,\qquad L=1+f\log\rho .
\]
Thus $[E:\mathbb Q]=2f$, $\sqrt D=Y$, $\rho\geq1$, $Y\geq1$, and $L\geq1$.
For a number field $K$, $h_K$ is its ideal class number and
$N\mathfrak a=[\mathcal O_K:\mathfrak a]$ is the absolute norm of a nonzero
integral ideal. A definite unimodular plane is a full $\mathcal O$-lattice
$\Lambda$ in a two-dimensional $E$-vector space, equipped with a Hermitian
form $h$ which is positive definite at every complex embedding and satisfies
\[
\Lambda=\{x:h(y,x)\in\mathcal O\quad(y\in\Lambda)\}.
\]
Let $\operatorname{defClassNum}(E,2)$ be the number of its isometry classes.
We use Minkowski's convex-body theorem, the Minkowski ideal bound, and
Dirichlet's unit theorem in the precise forms of Marcus, \emph{Number Fields},
2nd ed., Chapter 5: the lemma on pp.~96--97, Theorem 37, Corollary~2, and
Theorem~38, respectively.

\begin{lemma}[Ideal counting]\label{lem:hyperbola-ideals}
For every number field $K$ of degree $d$ and every $X\geq1$,
\[
\#\{\mathfrak a\subseteq\mathcal O_K:\mathfrak a\neq0,\,
N\mathfrak a\leq X\}\leq X(1+\log X)^{d-1}.
\]
Consequently $h_K\leq\sqrt{|\Delta_K|}(1+\log\sqrt{|\Delta_K|})^d$; in
particular $h_E\leq YL^{2f}$.
\end{lemma}
\begin{proof}
Unique factorization of ideals and $\sum_{\mathfrak p\mid p}e_{\mathfrak p}
f_{\mathfrak p}=d$ are Marcus, Chapter~3, Theorems~16 and~21. Colour the primes
over each rational prime $p$ by $d$ colours, injectively for fixed $p$. If
$\mathfrak a=\prod_{\mathfrak p}\mathfrak p^{e_{\mathfrak p}}$, form
$A_i=\prod_{\operatorname{colour}(\mathfrak p)=i}(N\mathfrak p)^{e_{\mathfrak p}}$.
The tuple $(A_1,\ldots,A_d)$ determines $\mathfrak a$. Hence at most $d_d(m)$
ideals have norm $m$, where $d_d(m)$ counts ordered $d$-tuples of product $m$.

Set $D_1(N)=N$ and
$D_{r+1}(N)=\sum_{m=1}^{N}D_r(\lfloor N/m\rfloor)$. Induction on $r$, using
$\sum_{m\leq N}m^{-1}\leq1+\log N$, gives
$D_r(N)\leq N(1+\log X)^{r-1}$ for $N\leq X$. Since
$D_d(N)=\sum_{m\leq N}d_d(m)$, this proves the first bound. Every ideal class
has an integral ideal of norm at most
$(4/\pi)^{r_2}(d!/d^d)\sqrt{|\Delta_K|}\leq\sqrt{|\Delta_K|}$ by the
Minkowski bound, so selecting one such ideal for each class proves the second
bound. Finally $\sqrt D=Y$ and $[E:\mathbb Q]=2f$.
\end{proof}

\begin{lemma}[Gram determinant and index]\label{lem:gram-index}
Let $X\subseteq E^r$ be an integral self-dual Hermitian lattice containing
$\mathcal O^r$. If $a$ is the determinant of its standard-basis Gram matrix
and $n=[X:\mathcal O^r]$, then $|N_{E/\mathbb Q}(a)|=n^2$.
\end{lemma}
\begin{proof}
Successive coordinate projections and splitting by invertible fractional ideals
give a pseudobasis $X=\bigoplus_iC_i z_i$, where $\mathcal O\subseteq C_i$
and the rows $z_i$ are unipotent triangular in the standard coordinates. Thus
$n=\prod_i[C_i:\mathcal O]$. Let $C=\prod_iC_i$ and
$\mathfrak J=\bar C(a)C$. If $\bar s_i\in C_i$ and $t_i\in C_i$, the matrix
$(H(\bar s_i z_i,t_jz_j))_{ij}$ is integral; its determinant is
$(\prod_i s_i)a(\prod_it_i)$. Hence $\mathfrak J\subseteq\mathcal O$.
Conversely, for $u_i\in\bar C_i^{-1}$, form the standard dual vectors with
coordinates $Q^{-1}(u_ie_i)$, where $Q$ is the conjugate triangular matrix
times the Gram matrix. Self-duality puts these vectors in $X$. Comparing their
pseudobasis coefficient determinants gives $\prod_i u_i\in(a)C$. Therefore
$\bar C^{-1}\subseteq(a)C$, so $\mathfrak J=\mathcal O$. The identities
$N(C_i)[C_i:\mathcal O]=N(\bar C_i)[C_i:\mathcal O]=1$ now give
$N(\mathfrak J)n^2=|N_{E/\mathbb Q}(a)|$.
\end{proof}

\begin{lemma}[A shortest vector at every embedding]\label{lem:short-all}
Every definite unimodular plane contains a nonzero $u\in\Lambda$ such that
$h_\varphi(u,u)\leq\rho$ for every $\varphi:E\hookrightarrow\mathbb C$.
\end{lemma}
\begin{proof}
Choose an $E$-basis $v_1,v_2$ lying in $\Lambda$ and set
$\Gamma=\mathcal O v_1\oplus\mathcal O v_2$ and $n=[\Lambda:\Gamma]$.
For each complex place $w$, let $G_w$ be the Gram matrix in this basis and
$\delta_w=\det G_w$. Positive definiteness on $E$-rational vectors, density of
$E$ in $\mathbb C$, and nonvanishing of the algebraic Gram determinant show
that $G_w$ is positive definite over $\mathbb C$.

Let $a$ be the algebraic Gram determinant and $P=\prod_w\delta_w$.
Lemma~\ref{lem:gram-index} and the principal-ideal norm formula give
$P^2=|N_{E/\mathbb Q}(a)|=n^2$, hence $P\leq n$. The mixed embedding of
$\Gamma$ has covolume $(2^{-f}\sqrt D)^2=D/4^f$, so the image of $\Lambda$
in $\prod_w\mathbb C^2$ has covolume $D/(4^fn)$. For $R^2=\rho$, define
\[
\mathcal E_R=\{(z_w)_w:z_w^*G_wz_w\leq R^2\quad(w)\}.
\]
The volume of a complex Hermitian ellipsoid follows from its Cholesky
factorization:
$\operatorname{vol}\{z:z^*G_wz\leq R^2\}=\pi^2R^4/(2\delta_w)$.
Since $R^{4f}=\rho^{2f}=D$,
\[
\operatorname{vol}(\mathcal E_R)=\frac{(\pi^2/2)^fD}{P},
\qquad 2^{4f}\operatorname{covol}(\Lambda)=\frac{4^fD}{n}.
\]
Minkowski's theorem therefore gives
\[
\#(\mathcal E_R\cap\Lambda)\geq
\frac{\operatorname{vol}(\mathcal E_R)}{2^{4f}\operatorname{covol}(\Lambda)}
=\frac nP\Bigl(\frac{\pi^2}{8}\Bigr)^f>1.
\]
This supplies the required nonzero vector at the place representatives.
Conjugate embeddings have the same real Hermitian value, proving the bound
at every embedding.
\end{proof}

\paragraph{Diagonal splitting.} Choose $u$ as in
Lemma~\ref{lem:short-all}, and complete it to an orthogonal basis $(u,w)$.
Writing $\gamma=(u,w)$ gives
$\gamma^*J\gamma=\operatorname{diag}(\beta_1,\beta_2)$, where
$\beta_1=h(u,u)$, $\beta_2=h(w,w)$ are totally positive. Put
$L=\gamma^{-1}\Lambda$,
\[
A=\{a:(a,0)\in L\},\qquad B=\{b:(0,b)\in L\},
\]
and let $P,Q$ be the two coordinate projections of $L$. Then $1\in A$ and
\[
t\in A\iff\bar p\beta_1t\in\mathcal O\quad(p\in P),\qquad
t\in B\iff\bar q\beta_2t\in\mathcal O\quad(q\in Q).
\tag{1}
\]
Both projections are nonzero because $L$ is full. Using inverse fractional
ideals in (1) shows that $A$ and $B$ are nonzero full rank-one lattices, just
as in the dual-colon calculation below. Define
\[
\mathfrak a_1=\operatorname{span}_{\mathcal O}
\{\bar x\beta_1y:x,y\in A\},\qquad
\mathfrak a_2=\operatorname{span}_{\mathcal O}
\{\bar x\beta_2y:x,y\in B\}.
\]
For $c\in E$, the identity
\[
\bar b\beta_1((\bar x\beta_1y)a)
=(\bar b\beta_1y)\overline{(\bar a\beta_1x)}
\]
and inversion of the fractional ideal $P$ give
$c\in\mathfrak a_1\iff cP\subseteq A$. Similarly
$c\in\mathfrak a_2\iff cQ\subseteq B$. If $cP\subseteq A$ and $(p,q)\in L$,
subtracting $(cp,0)$ from $c(p,q)$ proves $(0,cq)\in L$. Hence
$\mathfrak n:=\mathfrak a_1=\mathfrak a_2$.

The ideal $\mathfrak n$ is generated by its conjugation-fixed elements. Indeed,
let $\mathfrak b$ be the ideal generated by those elements. If
$(\mathfrak b:\mathfrak n)\neq\mathcal O$, choose a maximal $\mathfrak p$
containing it. By (1) and $P^{-1}P=\mathcal O$, there are $x,y\in A$ and
$a,b\in P$ with $\bar x\beta_1a$ and $\bar b\beta_1y$ outside $\mathfrak p$.
Examining the other two cross products shows that one of $x$, $y$, or $x+y$
can be used so that both
$s=\bar x\beta_1a$ and $t=\bar b\beta_1x$ are outside $\mathfrak p$. For every
generator $\bar y\beta_1z$ of $\mathfrak n$,
\[
st(\bar y\beta_1z)=(\bar y\beta_1a)(\bar b\beta_1z)(\bar x\beta_1x)
\]
belongs to $\mathfrak b$. Thus $st\in\mathfrak p$, a contradiction.
Consequently $\mathfrak n=(\mathfrak n\cap\mathcal O_F)\mathcal O$.
Finally $1\in A$ gives $\beta_1\in\mathfrak n$, and the shortest-vector
estimate at all $2f$ embeddings gives
\[
N\mathfrak n\leq|N_{E/\mathbb Q}(\beta_1)|
=\prod_\varphi|h_\varphi(u,u)|\leq\rho^{2f}=Y^2.
\tag{2}
\]

\begin{lemma}[Real ideals and lines]\label{lem:real-lines}
For $Y\geq1$, there are at most $YL^f$ nonzero fixed-generated ideals with
$N\mathfrak n\leq Y^2$. For each such $\mathfrak n$, there are at most
$h_E2^f$ equivalence classes of pairs $X=(\beta,N)$ with $\beta$ totally
positive, $N$ a full $\mathcal O$-lattice in $E$, and
$\operatorname{span}_{\mathcal O}\{\bar x\beta y:x,y\in N\}
=\mathfrak n$. Here $(\beta,N)\sim(\beta',N')$ if
$\bar c\beta'c=\beta$ and $cN=N'$ for some $c\in E^\times$.
\end{lemma}
\begin{proof}
Contraction $\mathfrak n\mapsto\mathfrak a=\mathfrak n\cap\mathcal O_F$ is
injective by fixed-generation, and
$N_{E/\mathbb Q}(\mathfrak a\mathcal O)=(N_{F/\mathbb Q}\mathfrak a)^2$.
The first claim follows from Lemma~\ref{lem:hyperbola-ideals} applied to $F$.

For lines, choose a reference $(\beta^0,N^0)$ in each occurring fractional
ideal class. For $X$ in that class choose $aN=N^0$. Equality of norm ideals
gives $u_X\in\mathcal O^\times$ with $\bar a\beta^0a\,u_X=\beta$: the multiplier
$t=\bar a\beta^0a/\beta$ and its inverse stabilize the nonzero finitely
generated ideal $\mathfrak n$, hence are integral. Record the parity of the
exponents of $u_X$ in a fundamental unit basis. Dirichlet's theorem gives
$f-1$ free generators, hence at most $2^{f-1}$ parity vectors.

If $X,Y$ have the same ideal class and parity, then $u_X=u_Y\zeta v^2$ with
$\zeta$ torsion and $v$ a unit. Put $c=a_Xv/a_Y$. This gives $cN_X=N_Y$ and
\[
\bar c\beta_Yc=\bar a_X\beta^0a_Xu_Y\bar v v.
\]
At any embedding, $|\zeta|=1$ and $|v^2|=|\bar v v|$, so this positive fixed
value equals $\beta_X=\bar a_X\beta^0a_Xu_X$. Thus $X\sim Y$. There are at
most $h_E2^{f-1}\leq h_E2^f$ line classes.
\end{proof}

\begin{lemma}[Unimodular gluings]\label{lem:glue-count}
Fix line representatives $(\beta_1,N_1)$ and $(\beta_2,N_2)$, with the norm
ideal of the first equal to $\mathfrak n$. The lattices $M\subseteq E^2$ which
are self-dual for $\operatorname{diag}(\beta_1,\beta_2)$ and have prescribed
axis intersections $N_1,N_2$ number at most $N\mathfrak n$.
\end{lemma}
\begin{proof}
Write $N^{\vee_i}=\{a:\bar s\beta_i a\in\mathcal O\ (s\in N)\}$. Choose
$0\neq s\in N_1$ and $d=\bar s\beta_1$. For the integral ideals
$I=\{c:d^{-1}c\in N_1\}$ and $J=\{c:d^{-1}c\in N_1^{\vee_1}\}$,
prime factorization gives $y$ with $J=I+(y)$. Taking $u=d^{-1}y$ gives
$N_1^{\vee_1}=N_1+\mathcal O u$. Also
$\mathfrak nN_1^{\vee_1}\subseteq N_1$, since
$(\bar x\beta_1y)u'=(\bar x\beta_1u')y$.

If $P$ is the first projection of an admissible $M$ and
$a\in N_1^{\vee_1}$, then $a\in P$. Indeed, for $q\in P^{-1}$,
$\bar q\beta_1^{-1}\in N_1$; hence $qa$ is integral and
$a(pq)=(qa)p\in P$. Summing over $PP^{-1}=\mathcal O$ proves $a\in P$.
Thus there is $b$ with $(a,b)\in M$, and similarly in the second coordinate.

Fix $M_0$ if the family is nonempty, and choose $(u,v_M)\in M$ for every $M$.
For each $M$, lift $v_M$ to $(a,v_M)\in M_0$ and choose $c_M\in\mathcal O$
with $a-c_Mu\in N_1$. Then $v_M-c_Mv_0\in N_2$. If
$c_M-c_{M'}\in\mathfrak n$, it follows that $v_M-v_{M'}\in N_2$. For any
$(a,b)\in M$, choose $d$ with $a-du\in N_1$; then $b-dv_M\in N_2$ and
\[
(a,b)=(a-du,0)+d(u,v_{M'})
 +(0,b-dv_M+d(v_M-v_{M'}))\in M'.
\]
Symmetry gives $M=M'$. Therefore $M\mapsto c_M\bmod\mathfrak n$ is an
injection into $\mathcal O/\mathfrak n$.
\end{proof}

\begin{theorem}\label{thm:rank-two-class-bound}
With the notation above,
\[
\operatorname{defClassNum}(E,2)\leq4^f\rho^{5f}(1+f\log\rho)^{5f}.
\]
\end{theorem}
\begin{proof}
By Lemma~\ref{lem:real-lines} and $h_E\leq YL^{2f}$, there are at most $YL^f$
ideals and at most $2^fYL^{2f}$ classes for each line. By
Lemma~\ref{lem:glue-count} and (2), each pair of line representatives has at
most $\lfloor Y^2\rfloor$ gluings. Thus the number of tuples
$(\mathfrak n,[X_1],[X_2],M)$ is at most
\[
(YL^f)(h_E2^f)^2\lfloor Y^2\rfloor
\leq4^fY^5L^{5f}.
\]
Every class has such a tuple. Namely, after splitting, choose for each
resulting line $X_i=(\beta_i,N_i)$ its reference representative
$X_i^0=(\beta_i^0,N_i^0)$, and scalars with
$\bar c_i\beta_i c_i=\beta_i^0$ and $c_iN_i^0=N_i$.
With $\delta=\gamma\operatorname{diag}(c_1,c_2)$ and $M=\delta^{-1}\Lambda$,
the diagonal form is $\operatorname{diag}(\beta_1^0,\beta_2^0)$ and $M$ has
exactly the required axis intersections. Fixed tuple data determine this
diagonal pair, and the original plane is isometric to it; hence the map from
classes to tuples is injective. Since $Y=\rho^f$ and $L=1+f\log\rho$, this is
the claimed bound.
\end{proof}

\paragraph{Absolute-volume substitution.} Let $\ell_E(M)$ denote the local
$L$-factor in the rank recurrence and set
$\operatorname{rankFac}(E,M)=\rho^{f(M-1/2)}\ell_E(M)$. The local comparison
$\ell_E(M)\leq\zeta(M)^f$ and the explicit estimate
$\prod_{M=3}^{n}\zeta(M)<1.4$ give
\[
\prod_{2<M\leq n}\operatorname{rankFac}(E,M)
\leq\rho^{f(n^2-4)/2}1.4^f,
\qquad\sum_{2<M\leq n}(M-\tfrac12)=\frac{n^2-4}{2}.
\tag{3}
\]
For the zeta product, truncate the $\zeta(3)$ and $\zeta(5)$ series and use
$\zeta(4)=\pi^4/90$, $\zeta(6)=\pi^6/945$; for $M\geq6$, comparison term by
term gives $\zeta(M)-1\leq2^{6-M}(\zeta(6)-1)$, hence
$\sum_{M\geq6}(\zeta(M)-1)<0.0347$. Together with $\zeta(3)<1.2021$,
$\zeta(4)<1.0824$, $\zeta(5)<1.0385$, and $e^{0.0347}<1.0354$, this gives the
product bound $1.3991<1.4$.

Assume now that the rank recurrence has reduced the volume $V_{p,q}$ to the
definite signature $(2,0)$ and that $n=p+q\geq8$, $q^2\leq n$, $X\geq2$,
$\rho\leq X$, and $2f\rho\leq X^2$. (The first two imply $p\geq2$.) Since $\operatorname{volU}(2)=8\pi^3$,
Theorem~\ref{thm:rank-two-class-bound} and (3) yield
\[
V_{p,q}\leq(8\pi^3)^f\operatorname{defClassNum}(E,2)
\prod_{2<M\leq n}\operatorname{rankFac}(E,M)
\leq(8\pi^3\cdot4\cdot1.4)^f\rho^{5f}L^{5f}\rho^{f(n^2-4)/2}.
\]
The remaining comparisons are explicit:
\[
8\pi^3\cdot4\cdot1.4\leq8(3.15)^3(4)(1.4)<2^{11}\leq X^{11},
\]
and, using $\log\rho\leq\rho-1$ and $f\geq1$,
\[
L=1+f\log\rho\leq1+f(\rho-1)\leq f\rho\leq X^2/2.
\]
Consequently
\[
V_{p,q}\leq X^{11f}X^{5f}X^{10f}X^{f(n^2-4)/2}
=X^{f(26+(n^2-4)/2)}\leq X^{n^2f},
\]
where the last inequality uses $n^2\geq64$, hence $48\leq n^2$.

\subsubsection{Decay and the far orbits}

Write $r=\min(p,q)\ge1$. Let $h$ be a lower bound for the gaps between opposite-sign weights and suppose $h\ge120nr^{3/2}$. In polar coordinates $k_1e^Xk_2$ put $t=\|X\|_{\rm HS}$. For unit vectors in the lowest $K$-type the absolute value of the matrix coefficient is at most
\[
 \exp\{-\tfrac h8\min(t^2,t/\sqrt{2r})\}.
\]
Indeed the compression to the lowest $K$-type is the compact integral of
$\prod_{j<n}\Delta_j(k^{-1}e^Xk)^{-(\lambda_j-\lambda_{j+1}+1)}$
against the two highest-vector matrix coefficients, with factor $\dim S$ for the lowest type $S$. The normalization follows by Schur orthogonality at the identity. The integral and the discrete-series coefficient satisfy the same lowering equations: commuting a compact derivative past a noncompact one gives the Lie bracket, and an ordered lowering derivative kills the highest vector. The gap $h>2n$ excludes the complementary types in the successive jets. Induction on the noncompact degree gives equality of the Taylor expansions, hence of the analytic coefficients along $\exp(tX)$.

For $Y=k^{-1}Xk$ the leading minor $\Delta_j(e^Y)$ is a convex sum of exponentials of sums of $j$ eigenvalues, by Cauchy--Binet. These sums have mean zero, absolute value at most $\sqrt{2r}\,t$, and variance $B_j=\sum_{a>j,\,b\le j}|Y_{ab}|^2$. The centered exponential estimate at $\min(1,1/(\sqrt{2r}\,t))$ and convexity give
$\log\Delta_j(e^Y)\ge B_j\min(1,1/(\sqrt{2r}\,t))/4$.
Summation over the cuts gives
$\sum_j(\lambda_j-\lambda_{j+1}+1)B_j\ge ht^2/2$.
Cauchy--Schwarz and Schur orthogonality in the compact integral give the asserted bound. In the operator radius $s=\|X\|_{\rm op}$ the paired nonzero eigenvalues give $t\ge\sqrt2s$, hence the exponent is at most $-(h/8)\min(2s^2,s/\sqrt r)$.

Multiply the local estimates by the product $d_\infty$ of the formal degrees and the idempotent bound $a=e(1)$. For least lattice scale $\tau$ put $A(x)=(3\sqrt n/\min(1/2,\tau)^2)^{2n^2f}$. The number of elements with radii at most $s_v$ is bounded by $A(x)\exp(2n^2\sum_v s_v)$, by packing lattice points in the product of balls. Put $\kappa=2n^2$, $T=h\varepsilon^2/16$, and take $\varepsilon\le1/\sqrt{2r}$, $\kappa\varepsilon\le\log2$, $32\kappa\le h\varepsilon$, $T\ge1$. In a dyadic shell with upper radius $2^j\varepsilon$, $j\ge1$, decay times the counting factor is at most $\exp(-T2^{j-1})$; the bottom shell contributes at most $2$. The sum of the positive-index terms is at most $2e^{-T}\le1$. A far orbit has at least one positive index, so the product over $f$ places gives an absolutely convergent sum bounded by $d_\infty a A(x)\,2f3^fe^{-T}$.

For the deep region put $N=n^2f$, $P=8qN$, and write $\tau$ for the least lattice scale, $\rho=n^2\sqrt f\,g\,\mathrm{rd}_E(2\tau)^{1/q}/4$. Let $\beta$ bound the one-step derivative constants through order $P+N$. The local Sobolev estimate on jets has squared evaluation bound
\[
 (6\beta)^N e^f(17\sqrt n/\tau^2)^{2N}
 D\,\prod_v c_v/\mu_f(K_g).
\]
Here $c_v$ are the local gauge constants. The lattice scale on the radius-$\rho$ neighbourhood is at least $\tau/2$. At each of the $P$ steps the constant term along the isotropic radical vanishes by cuspidality; comparison with its average costs at most $2\rho\sqrt f\,\beta$. Squaring gives the additional factor $(2\rho\sqrt f\,\beta)^{2P}$. Applying the frame estimate first to finite subspaces and integrating with the absolute volume bound gives finite dimensionality. An orthonormal basis then gives the diagonal kernel as the squared evaluation norm.

\subsubsection{Near orbits and exceptional flags}\label{fr:sec:near-orbits}

Throughout this subsection let $E/F$ be the totally imaginary quadratic extension of the totally real field in the automorphic datum, and put $f=[F:\mathbb Q]$. Thus $E$ has $f$ complex places. Let $n$ be the hermitian rank and write $J=\sum_i j_i$ for a type $(j,\tau)$ of length $a$. Define
\[
 \mathcal P_\delta(j,\tau)=
 \prod_{i<k}\!\left[G_\delta\!\left(\frac{\tau_i}{\tau_k}\right)^{2j_ij_kf}
 G_\delta(\tau_i\tau_k)^{2j_ij_kf}\right]
 \prod_i\!\left[G_\delta(\tau_i)^{2j_i(n-2J)f}
 G_\delta(\tau_i^2)^{j_i^2f}\right],
 \qquad G_\delta(u)=1+\frac{12\delta}{u}.
\]
Write $\mathcal F$ for the fundamental domain with measure $\mu$, and $\varphi$ for the product of the archimedean matrix coefficients, their formal degrees and the finite level idempotent. Write $s_{\min}(x)$ for the least lattice scale. For the lattice at $x$, let $S_{\le}(t)$ and $S_{<}(t)$ be the spans of sublattices of largest size at most $t$ and strictly less than $t$, respectively. Let $\mathcal N(x,\delta)$ consist of the rational elements whose conjugates by $x$ lie in the finite level and within Frobenius distance $\delta$ of the identity at each infinite place. Write $\mathcal N^*(x)$ for the nonidentity near orbits at the original Cartan cutoff.

\begin{proposition}[Near unipotence]\label{fr:prop:near-unipotent}
Assume $n\ge1$, $4n<g_0$, and that $(g_0)\mathcal O_E$ is radical. For $0<c\le10^{-5}$ put
\[
 \varepsilon=\frac{c}{n^8f^2},\qquad \delta=3\sqrt{n\varepsilon}.
\]
Then $\delta<1/4$, and $\delta qnf\le10^{-2}$ whenever $q^2\le n$. Moreover every automorphic $\gamma$ for which $\varphi(x^{-1}\gamma x)\ne0$ and whose archimedean Cartan distances after conjugation by $x$ are at most $\varepsilon$ has
\[
 (\gamma_f-1)^n=0,\qquad (\gamma_v-1)^n=0\quad(v\mid\infty),
\]
and belongs to $\mathcal N(x,\delta)$.
\end{proposition}

\begin{proof}
The scalar assertions are uniform in $n,f\ge1$. With $a=n\varepsilon$ one has
\[
 a(n^3f^2)=c/n^4\le c,\qquad a\le c\le10^{-5}.
\]
Hence $3\sqrt a<1/4$. If $q^2\le n$, then
\[
 (3\sqrt a\,qnf)^2=9a(q^2n^2f^2)\le9a(n^3f^2)\le9c<10^{-4},
\]
which gives the last radius bound. Also $n\varepsilon\le1/4$.

For an invertible complex matrix $A$, write $A=ke^X$ with $k$ unitary and $X$ hermitian. This is the unitary polar decomposition: $A^*A=e^{2X}$ and $k=Ae^{-X}$. The Cartan distance in the statement is $\|X\|_F$. Suppose $\|X\|_F\le\varepsilon$ and put $E=e^X$. Then
\[
 \|E-1\|_F\le e^\varepsilon-1,\qquad
 \|(kE)^r\|_F\le e^{r\varepsilon}\sqrt n,
\]
by expanding the exponential and inducting on $r$. Cauchy--Schwarz on diagonal entries gives $|\operatorname{tr}M|\le\sqrt n\,\|M\|_F$. Since $r\varepsilon\le n\varepsilon\le1/4$ and $e^{1/4}<3/2$, it follows for $0\le r\le n$ that
\begin{equation}\label{fr:eq:archtrace}
 |\operatorname{tr}((kE)^r)-n|\le ne^{r\varepsilon}+n\le3n.
\end{equation}
Now suppose in addition $\operatorname{tr}(kE)=n$. Using $\varepsilon\le1$ gives $\|E-1\|_F\le2\varepsilon$. The identities
\[
 \|k-1\|_F^2=2n-2\Re\operatorname{tr}k,
 \qquad \operatorname{tr}k=n-\operatorname{tr}(k(E-1))
\]
imply $\|k-1\|_F\le2\sqrt{n\varepsilon}$. Since $\sqrt\varepsilon\le\sqrt n$ and $\sqrt{n\varepsilon}\le1/2$,
\[
 \|kE-1\|_F\le2\sqrt{n\varepsilon}+2\varepsilon\le3\sqrt{n\varepsilon}=\delta.
\tag{1}
\]

Nonvanishing of the level idempotent at the conjugated finite component puts that component in the lattice stabilizer and makes it unipotent modulo $g_0$, by the two support clauses of admissibility. Let $G$ be the rational matrix supplied by rationality of the automorphic orbit, and let $\Lambda_x$ be its self-dual $\mathcal O_E$-lattice. Thus $G\Lambda_x\subseteq\Lambda_x$ and
\[
 (G-1)^n\Lambda_x\subseteq g_0\Lambda_x.
\]
For fixed $r$, set $S_r=\sum_{s=0}^{r-1}G^s$. Since $(G-1)S_r=G^r-1$ and $S_r$ commutes with $G-1$, the same divisibility holds for $(G^r-1)^n$.

We use the following trace congruence. If $A\Lambda\subseteq\Lambda$ and $A^m\Lambda\subseteq g_0\Lambda$, then $\operatorname{tr}A\in g_0\mathcal O_E$. Indeed, for an eigenvalue $z$ of $A$ in an algebraic closure, a nonzero left eigenvector and a common denominator in $\Lambda$ show that $z$ is integral over $\mathcal O_E$, while $z^m/g_0$ is integral. Thus $z$ lies in the radical of $g_0$ in the integral closure. Summing eigenvalues gives $(\operatorname{tr}A)^N\in g_0\mathcal O_E$ for some $N$; radicality of $(g_0)$ gives the assertion. Consequently
\[
 \operatorname{tr}(G^r)-n=g_0t_r\qquad(t_r\in\mathcal O_E).
\]

The archimedean realizations of $G^r$ are conjugate to the matrices in \eqref{fr:eq:archtrace}, so $|\sigma(g_0t_r)|\le3n$ at every complex embedding $\sigma$. If $t_r\ne0$, the product formula for its nonzero algebraic norm gives some $|\sigma(t_r)|\ge1$; hence $g_0\le3n$, contrary to $4n<g_0$. Thus $\operatorname{tr}(G^r)=n$ for $0\le r\le n$.

For $N=G-1$, the binomial expansion now yields $\operatorname{tr}(N^k)=n(1-1)^k=0$ for $1\le k\le n$. Newton's identities over the characteristic-zero field $E$ force the characteristic polynomial of $N$ to be $X^n$. Cayley--Hamilton gives $N^n=0$. Ring homomorphisms and conjugation preserve this identity, giving unipotence at every finite and infinite component. In particular the trace at $r=1$ is $n$, so (1) supplies exactly the archimedean radius in $\mathcal N(x,\delta)$. The other two membership clauses are rationality and the lattice stabilizer already established.
\end{proof}

At the operative hypotheses, $g_0$ is prime and the number of primes of $\mathcal O_E$ over $g_0$ equals $[E:\mathbb Q]$. The identity $\sum_{\mathfrak p\mid g_0}e_{\mathfrak p}f_{\mathfrak p}=[E:\mathbb Q]$ therefore forces every ramification index and residue degree to be $1$. In particular $(g_0)$ is squarefree, hence radical. This is the total-splitting input used above.

\begin{lemma}[Geometric-mean packing]\label{fr:lem:gm-packing}
Let $U_v$, $v=1,\dots,f$, be real vector subspaces of a finite-dimensional normed space, with $\dim_{\mathbb R}U_v\le N$. Suppose $y_t(v)\in U_v$ has norm at most $X>0$, and distinct $t,t'$ satisfy
\[
 \left(\prod_v\|y_t(v)-y_{t'}(v)\|\right)^{1/f}\ge\lambda>0.
\]
Then the number of indices $t$ is at most $(1+2X/\lambda)^{Nf}$.
\end{lemma}

\begin{proof}
Put $G=(\prod_vX)^{1/f}=X$ and, in the general version with radii $X_v$, rescale coordinate $v$ by $G/X_v$. These scalars have product $1$, so the geometric mean of every difference is unchanged. After rescaling all coordinates have norm at most $G$, while the geometric mean of a difference is at most its sup norm. Thus the points are $\lambda$-separated in sup norm in a ball of radius $G$. Their balls of radius $\lambda/2$ are disjoint and lie in the ball of radius $G+\lambda/2$. Haar-volume comparison in dimension at most $Nf$ proves the bound; padding each $U_v$ to dimension $N$ removes the word ``at most''.
\end{proof}

We next isolate the slope estimate behind the block count. For an $\mathcal O_E$-lattice $L$ of rank $r$, write
\[
 H(L)=2^{fr}|d_E|^{-r/2}\operatorname{covol}(L),\qquad
 \operatorname{size}(L)=H(L)^{1/(2fr)},
\]
and use the analogous normalized covolume $H(A/A_0)$ and $\operatorname{size}(A/A_0)$ for a saturated $A_0\subsetneq A$. Let $\operatorname{leastSize}(B)$ be the infimum of sizes of nonzero sublattices of $B$, and $\operatorname{largestSize}(A)$ the supremum of quotient sizes over proper saturated $A_0$.

If $\phi:A\to B$ has nonzero image and local operator norms $c_v$, use the singular-value decomposition of the induced map from $A/(A\cap\ker\phi)$ to its image at each complex place. Its $r=\operatorname{rk}\phi(A)$ singular values are at most $c_v$, so the determinant of that induced map has squared absolute value at most $c_v^{2r}$. Applied to a free full-rank sublattice of the quotient, this gives, with $P=\prod_vc_v$,
\[
 H(\phi A)\le P^{2r}H(A/(A\cap\ker\phi)).
\]
Taking the $(2fr)^{-1}$-st root gives
\[
 \operatorname{size}(\phi A)\le P^{1/f}\operatorname{size}(A/(A\cap\ker\phi)).
\]
The image itself is a test lattice for $\operatorname{leastSize}(B)$, and $\ker\phi\cap A$ is a test quotient for $\operatorname{largestSize}(A)$. Therefore
\begin{equation}\label{fr:eq:size-inequality}
 \frac{\operatorname{leastSize}(B)}{\operatorname{largestSize}(A)}
 \le\left(\prod_vc_v\right)^{1/f}.
\end{equation}
Quotienting by a saturated sublattice and replacing operator norms by distance-to-subspace bounds gives the same estimate; this is the form used below.

\begin{proposition}[Filtration-aware near count]\label{fr:prop:near-count}
For every point $x$ of type $(j,\tau)$ and every $0<\delta<1/4$,
\[
 \mathcal N(x,\delta)\text{ is finite},\qquad
 \#\mathcal N(x,\delta)\le\mathcal P_\delta(j,\tau).
\]
No additional numerical hypothesis on the type is needed.
\end{proposition}

\begin{proof}
The type supplies a saturated flag $0=S_0\subsetneq\cdots\subsetneq S_a$ with rank increments $j_i$ and
\[
 S_{i+1}=S_{\le}(\tau_i),\qquad S_i=S_{<}(\tau_i),\qquad S_a=S_{<}(1/2).
\]
Complete it by hermitian duality: $G_p=S_p$ for $p\le a$ and $G_p=G_{2a+1-p}^{\perp}$ for $p>a$. Put $t_i=\tau_i$ for $i<a$, $t_a=2$, and $t_p=\tau_{2a-p}^{-1}$ for $p>a$. Self-duality gives $G_{2a+1}=\Lambda_x$, $\operatorname{rk}G_p+\operatorname{rk}G_{2a+1-p}=n$, and the strict-step descriptions of every $G_p$. If $r_p=\operatorname{rk}G_{p+1}-\operatorname{rk}G_p$, then
\[
 r_i=r_{2a-i}=j_i\ (i<a),\qquad r_a=n-2J.
\]

The condition $\delta<1/4$ is used at the central step. A near matrix lowers $S_{\le}(t)$ to $S_{\le}(\delta t)$. If it did not, applying \eqref{fr:eq:size-inequality} to the induced nonzero map on the quotient would give a ratio at most $\delta$, whereas the strict-step lower size and the largest-size upper bound make that ratio greater than $\delta$. At $p=a$, $2\delta<1/2$ gives $S_{\le}(2\delta)\subseteq S_{<}(1/2)=G_a$. Consequently a rational near matrix lies in
\[
 \mathcal L=\{\gamma:\gamma\text{ isometric},\ ( \gamma-1)G_{p+1}\subseteq G_p\ (p\le2a),\
 \|\gamma_v-1\|\le\delta\ \forall v\}.
\]
The rational representation is injective: its archimedean and finite matrices determine the adelic element. It therefore remains to count $\mathcal L$.

Index blocks by $\mathcal B=\{(q,p):0\le q<a,\ q<p,\ q+p\le2a\}$ and order them by decreasing height $p-q$. The relation that $D=\gamma-\gamma'$ maps $G_{p+1}$ to $G_q$ is an equivalence. Unitarity gives the exact identity
\[
 h(Du,\gamma v)+h(\gamma'u,Dv)=0.
\]
A strong induction on $p$, with a second induction on the dual index, propagates the relation from smaller heights to all pairs $(q,p)$ of height $p-q<c$ and proves $D(G_p)\subseteq G_{p-c}$. It also proves
\[
 h(Du,v)+h(u,Dv)=0\qquad
 (u\in G_{p+1},\ v\in G_{p'+1},\ p+p'\le2a+c).
\]
At height $2a+1$, this makes $D$ zero on $G_{2a+1}=\Lambda_x$, hence zero because $\Lambda_x$ spans.

For a fixed block $(q,p)$, subtract a basepoint $\gamma_0$. The corresponding archimedean block is the projection of $(\gamma-\gamma_0)_v$ from the graded piece of rank $r_p$ to the dual graded piece of rank $r_{2a-q}$. It has real dimension $2r_pr_{2a-q}$, except when $q+p=2a$; in that diagonal case the skew identity restricts it to a skew-hermitian space of real dimension $r_p^2$. These dimensions are exactly
\[
 d_{q,p}=\begin{cases}r_p^2,&q+p=2a,\\2r_pr_{2a-q},&q+p<2a.\end{cases}
\]
All block norms are at most $2\delta$.

If two elements differ, the relation at level $(q,p)$ fails for a quotient-separating subset unless they agree. By the previous paragraph, nontriviality of the induced map supplies a nonzero quotient image. The quotient form of \eqref{fr:eq:size-inequality}, the strict-step lower bound $t_q$ and the largest-step upper bound $t_p$ give geometric-mean separation at least $t_q/t_p$. Lemma~\ref{fr:lem:gm-packing} therefore bounds this fibre by
\[
 \left(1+\frac{2(2\delta)}{t_q/t_p}\right)^{d_{q,p}f}.
\]
Taking fibre products in decreasing height proves finiteness and
\begin{equation}\label{fr:eq:raw-block-product}
 \#\mathcal L\le\prod_{(q,p)\in\mathcal B}
 \left(1+\frac{4\delta}{t_q/t_p}\right)^{d_{q,p}f}.
\end{equation}

It remains to evaluate the product. For $i<k<a$ the two pairs $(i,k)$ and $(i,2a-k)$ have dimensions $2j_ij_k$ and ratios $\tau_i/\tau_k$ and $\tau_i\tau_k$, respectively. The central pairs $(i,a)$ and $(i,2a-i)$ have dimensions $2j_i(n-2J)$ and $j_i^2$ and ratios $\tau_i/2$ and $\tau_i^2$. All these identifications use $t_a=2$. Finally, for positive $\ell,\ell'$ with $\ell\le3\ell'$,
\[
 \left(1+\frac{4\delta}{\ell'}\right)^m
 \le\left(1+\frac{12\delta}{\ell}\right)^m,
\]
because $4\delta/\ell'\le12\delta/\ell$. It applies with equality of ratios away from $(i,a)$, and with $\ell=\tau_i$, $\ell'=\tau_i/2$ at $(i,a)$. Substitution in \eqref{fr:eq:raw-block-product} gives exactly $\mathcal P_\delta(j,\tau)$, including every exponent and both factors for each off-diagonal pair.
\end{proof}

For the integrated estimate write $L=nf\log n$, $X=nfrd\,g\,b$, $\ell_*=X^{-11q}$, and
\[
 Y=2fq(11q\log X).
\]
Let $\mathrm M=(\prod_vd_v)\mu(\mathcal F)\operatorname{Re}e(1)$ be the main term, and let $\kappa_{p,q,j}$ be the flag-class constant. The pointwise orbit input gives $|\varphi(y)|\le M_0=(\prod_vd_v)\operatorname{Re}e(1)$ and $\mathrm M=M_0\mu(\mathcal F)\ge0$.

\begin{proposition}[Near terms on the moderate region]\label{fr:prop:near-moderate}
Use the flag-measure bound, its constant bound, near unipotence, and the pointwise bound with their stated hypotheses. For each $C_g\ge0$ there is an $n_0$, which can be taken to be
\[
 n_0=\max\!\left\{3^{15},\left\lceil e^{16(1+22(|C_g|+4))}\right\rceil\right\},
\]
such that, under the full counting hypotheses with $q\ge1$,
\[
 \int_{\mathcal F\cap\{\ell_*\le s_{\min}\}}
 \sum_{\gamma\in\mathcal N^*(x)}|\varphi(x^{-1}\gamma x)|\,d\mu(x)
 \le e^{-L/4}\mathrm M.
\]
The left side is interpreted as an extended-nonnegative-real integral and the right side through $\mathbb R_{\ge0}\cup\{\infty\}$-valued realification.
\end{proposition}

\begin{proof}
Set $\delta=3\sqrt{n\,10^{-6}/(n^8f^2)}$. Near unipotence and the split-prime argument supply $0<\delta<1/4$ and $\delta qnf\le10^{-2}$. If $s_{\min}(x)\ge1/2$, the same lowering argument as in the preceding proof forces $\mathcal N(x,\delta)=\{1\}$, so the nonidentity near sum is zero.

Fix a composition type $j$ with $J\le q$. On its moderate part, Proposition~\ref{fr:prop:near-count}, finiteness of the near set and the pointwise bound give the pointwise majorant $M_0\mathcal P_\delta(j,\tau)$. The flag-measure bound turns integration over the type into
\[
 M_0\mu(\mathcal F)\kappa_{p,q,j}
 \int_{\tau\in\mathcal S_a,\ \ell_*\le\tau_i}
 \mathcal P_\delta(j,\tau)\operatorname{dens}_{n,f,j}(\tau)\,d\tau.
\]
Here $\operatorname{dens}=\mathcal P_0^{\mathrm{id}}\prod_i(2fj_i/\tau_i)$, where $\mathcal P_0^{\mathrm{id}}$ is the product defining $\mathcal P_\delta$ with $G_\delta(u)$ replaced by $u$. Every type parameter satisfies $0<\tau_i<1/2$ and is increasing. For any argument $\lambda$ occurring in this product,
\[
 (1+12\delta/\lambda)\lambda=\lambda+12\delta\le e^{12\delta}.
\]
The sum of all exponents is at most $3Jnf\le3qnf$, so
\[
 \mathcal P_\delta(j,\tau)\mathcal P_0^{\mathrm{id}}(j,\tau)
 \le e^{36\delta qnf}\le e^{36/100}\le2.
\]
The last inequality uses $36/100<\log2$. Consequently the one-type integral is at most
\[
 \mathrm M\,2\kappa_{p,q,j}\prod_i\int_{\ell_*}^{1/2}\frac{2fj_i}{\tau_i}\,d\tau_i
 =\mathrm M\,2\kappa_{p,q,j}\prod_i\left(2fj_i\log\frac1{2\ell_*}\right).
\]
Since $\ell_*=X^{-11q}$, $\log(1/(2\ell_*))\le11q\log X$, and $j_i\le q$, this is at most
\begin{equation}\label{fr:eq:one-type}
 \mathrm M\,2\kappa_{p,q,j}Y^a.
\end{equation}
There is no lost volume factor: $\mathrm M=M_0\mu(\mathcal F)$. If $\kappa_{p,q,j}\le0$, the realification is zero and the inequality is immediate; otherwise the same real calculation applies.

The constant bound, with $p+q=n$, $f_{\mathrm{deg}}(E)=f$ and $2v_{\mathrm{inv}}(E)=rd$, gives
\[
 \kappa_{p,q,j}\le e^{-JL/2}.
\]
Also $a\le J$, since every block is positive. The next paragraph proves $Y\le e^{L/8}$. Thus
\[
 2\kappa_{p,q,j}Y^a
 \le2\exp(-JL/2+aL/8)\le2e^{-3L/8}.
\]

To prove $Y\le e^{L/8}$, put $K=22(|C_g|+4)$ and $E_0=e^{L/16}$. The chosen threshold gives $\log n\ge16(1+K)$. Hence $nf+K\le L/16$ and $(nf)K\le E_0$. Using $q^2\le n$, $\log rd\le e^n\le E_0$, $\log f\le\sqrt n\le E_0$, $\log b\le E_0$, and $\log g\le C_g\log b\le|C_g|E_0$ gives
\[
 \log X\le(|C_g|+4)E_0,
 \qquad Y=22fq^2\log X\le(nfK)E_0\le E_0^2=e^{L/8}.
\]
Moreover $2^q\le e^q$ and $1+q\le L/8$, so
\begin{equation}\label{fr:eq:final-absorption}
 2\,2^q e^{-3L/8}\le e^{1+q-3L/8}\le e^{-L/4}.
\end{equation}

Finally, points outside the finite type union contribute zero. Canonically filtering a lattice with $s_{\min}<1/2$ (Grayson's canonical filtration \cite[Theorem 1.18 and Definition 1.19, p. 609]{r3grayson}) gives a type with $1\le J\le q$; isotropy gives the upper bound. There are $2^{J-1}$ compositions of $J$. Subadditivity and \eqref{fr:eq:one-type} therefore leave the factor
\[
 \sum_{J=1}^q2^{J-1}=2^q-1<2^q.
\]
Applying \eqref{fr:eq:final-absorption} completes the proof with the original exponent $-L/4$.
\end{proof}

\begin{proposition}[Rare flags]\label{fr:prop:rare-flags}
The nonzero-flag bound implies the rare-flag bound with exponent $-L/4$, with threshold $\lceil e^{15}\rceil$.
\end{proposition}

\begin{proof}
If $n\ge\lceil e^{15}\rceil$, then $\log n\ge15$. Apply the nonzero-flag bound with the unchanged signature, $q^2\le n$, $\log rd\le e^n$, and $\log f\le\sqrt n$ hypotheses. The step characterization gives
\[
 \{s_{\min}<1/2\}=\{S_{<}(1/2)\ne0\}.
\]
Thus its measure divided by $\mu(\mathcal F)$ is at most $e^{-L/3}$. Since $n,f\ge1$ and $\log n\ge15$, one has $L>0$ and
\[
 e^{-L/3}\le e^{-L/4}.
\]
No strengthening of the degree or discriminant hypotheses is required.
\end{proof}

\paragraph{The total error.}
Put $L=nf\log n$, $\delta=\exp(-n^2f)$, and $X=|D-M|$.
The required orbit estimates give
\[
 X\le (2e^{-L/4}+\delta)M+\delta D .
\]
For $n\ge e^{15}$ put $\eta=e^{-L/16}$. Then $240\le L\le n^2f$,
so $e^{-L/4},\delta\le\eta/8$ and $\delta\le1/2$. Since $D\le M+X$,
\[
 X\le (2e^{-L/4}+2\delta)M+\delta X
   \le \tfrac12\eta M+\tfrac12X .
\]
Thus $|D-M|\le e^{-nf\log n/16}M$.



\section{Heights, the CM period, and the slopes of the target and base lattices}\label{r4:sec:R4}

This section proves hypotheses \eqref{A4}, \eqref{A6} and \eqref{A7} of Proposition~\ref{acct}. Subsection~\ref{r4:sec:R4conv} fixes the normalization of slopes and heights. Subsection~\ref{r4:sec:R4period} proves that the constant $\log\Omega$ of Proposition~\ref{acct} is the normalized arithmetic degree of the line of periods of a CM elliptic curve, which is the Chowla--Selberg formula in the form needed here. Subsection~\ref{r4:sec:R4height} bounds the height of the period map $B_u$ by the product of the norms of the $m$ lifts and of the final pairing. Subsection~\ref{r4:sec:R4target} states the bound for the slope of the target lattice $R^\vee$ that the base side gives. Subsection~\ref{r4:sec:R4A4} combines the two bounds into \eqref{A4}. Subsection~\ref{r4:sec:R4A6} proves \eqref{A6} from the Weyl dimension formula and the Siegel mass formula, and Subsection~\ref{r4:sec:R4A7} proves \eqref{A7} from the integral theory of Weyl modules. Subsection~\ref{r4:sec:R4sizes} compares the three errors with the quantity $\mathcal B$ of \eqref{calB}.

Throughout, $x$, $\varepsilon$, $L_\varepsilon=x^\varepsilon$, $b=e^{L_\varepsilon}$, $a=xb$, $X'$ and $b'=a/X'$ are as in \eqref{scale}, and the field $E$, the CM type $\Phi$, the integer $f=2n=\lvert\Phi\rvert$, the number of lifts $m$, the ranks $k$, $t_r$ and $r=k+m-1$, the set of tops $T$ with $N_T=\lvert T\rvert$, the weights $a_t$ and the nodes $w_{j\rho}$ are those of Sections~\ref{sec:auxiliary} and~\ref{sec:reduction}. We write $\mathrm{rd}_E=\lvert d_E\rvert^{1/(2f)}$ for the root discriminant of $E$ and put
\begin{equation}\label{r4:R4:v}
v=\tfrac12\,\mathrm{rd}_E,\qquad \Lambda_a=\log(a\,\mathrm{rd}_E),\qquad g=\log(g_0g_1),
\end{equation}
where $g_0,g_1$ are the primes of (D5). The field $K=\mathbb Q(\sqrt{-d})$ is the imaginary quadratic subfield of $E$ fixed by $H$, with its embedding $\varphi_0$, and $\chi_d$ is the quadratic character of $K$. Since $d\equiv3\pmod4$ is prime, $\chi_d(n)=\big(\frac{-d}n\big)=\big(\frac nd\big)$, the discriminant of $K$ is $-d$, and $\chi_d$ is the character $\chi_1$ of Proposition~\ref{excess}. We put
\[
\lambda_1=\frac{L'(1,\chi_d)}{L(1,\chi_d)},
\]
which is real.

\subsection{Slopes and heights}\label{r4:sec:R4conv}

Let $F$ be a totally complex number field with complex conjugation, and let $n_F=[F:\mathbb Q]/2$ be the number of its archimedean places. A \emph{hermitian lattice} over $\mathcal O_F$ is a finitely generated projective $\mathcal O_F$-module $\Lambda$ with a positive definite hermitian form $h_v$ on $\Lambda\otimes_{\mathcal O_F,\sigma_v}\mathbb C$ for each archimedean place $v$, where $\sigma_v$ is one of the two embeddings of $F$ at $v$. For a basis $e_1,\dots,e_s$ of $\Lambda\otimes F$ contained in $\Lambda$ put
\[
\degh(\Lambda)=\log\#\big(\Lambda/(\mathcal O_Fe_1+\dots+\mathcal O_Fe_s)\big)-\sum_v\log\det\big(h_v(e_i,e_j)\big)_{i,j}.
\]
The product formula shows that $\degh(\Lambda)$ does not depend on the basis, and it is the arithmetic degree of Bost~\cite{bost}, since the two embeddings at $v$ give equal determinants. We use the normalized slope and height
\begin{equation}\label{r4:R4:slope}
\mu(\Lambda)=\frac{\degh(\Lambda)}{n_F\,\rk\Lambda},\qquad \mathrm{ht}(\phi)=\frac1{n_F}\sum_v\log\lVert\phi\rVert_v^2
\end{equation}
for a nonzero $F$-linear map $\phi$ between hermitian lattices, where $\lVert\phi\rVert_v$ is the operator norm at $v$. With $\mu_{\max}$ and $\mu_{\min}$ the largest and smallest slopes of the Harder--Narasimhan filtration, the slope inequality of Bost~\cite{bost} in this normalization reads $\mu(\Lambda_1)\le\mu_{\max}(\Lambda_2)+\mathrm{ht}(\phi)$ for an injective $\phi\colon\Lambda_1\to\Lambda_2$. The dual lattice $\Lambda^\vee=\operatorname{Hom}_{\mathcal O_F}(\Lambda,\mathcal O_F)$ has the dual forms, and $\degh(\Lambda^\vee)=-\degh(\Lambda)$, so that
\begin{equation}\label{r4:R4:dual}
\mu(\Lambda^\vee)=-\mu(\Lambda),\qquad \mu_{\min}(\Lambda^\vee)=-\mu_{\max}(\Lambda).
\end{equation}
For positive real numbers $\kappa_v$ we write $\langle\log\kappa\rangle=n_F^{-1}\sum_v\log\kappa_v$. Multiplying the forms $h_v$ by $\kappa_v$ changes $\mu$ by $-\langle\log\kappa\rangle$ and $\mathrm{ht}$ of a map into $\Lambda$ by $+\langle\log\kappa\rangle$.

\begin{remark}\label{r4:R4:normrem}
These are twice the slopes and heights of~\cite{bost}, and they coincide with the slope~\eqref{r3:mu} and the height~\eqref{r3:ht} of Section~\ref{r3:sec}: both vanish on the trivial lattice and both equal twice the normalized slope and height of Bost. Section~\ref{r2:sec:loss} uses only differences of slopes of lattices over one field. The doubling is the normalization of Proposition~\ref{acct}, where the sum $S$ over primes has the factor $2$ of \eqref{excessid}, and every hypothesis \eqref{A1} to \eqref{A7} is stated in it.
\end{remark}

We write $\ph(y)=y\log y-y$ for $y>0$ and $\ph(0)=0$.

\subsection{The CM period}\label{r4:sec:R4period}

For a lattice $\Lambda\subset\mathbb C$ let $g_2(\Lambda)=60\sum'_{\omega\in\Lambda}\omega^{-4}$, $g_3(\Lambda)=140\sum'_{\omega\in\Lambda}\omega^{-6}$, $\Delta(\Lambda)=g_2^3-27g_3^2$ and let $A(\Lambda)$ be the area of a fundamental parallelogram. Since $g_2(c\Lambda)=c^{-4}g_2(\Lambda)$, $g_3(c\Lambda)=c^{-6}g_3(\Lambda)$ and $A(c\Lambda)=\lvert c\rvert^2A(\Lambda)$, the number
\[
\mathcal A(\Lambda)=A(\Lambda)\,\lvert\Delta(\Lambda)\rvert^{1/6}
\]
depends on the homothety class of $\Lambda$ alone.

\begin{lemma}\label{r4:R4:areaeta}
Let $\Lambda=\mathbb Z\omega_1+\mathbb Z\omega_2$ with $\tau=\omega_2/\omega_1$ in the upper half plane, and let $\eta$ be the Dedekind eta function. Then $\mathcal A(\Lambda)=4\pi^2\operatorname{Im}(\tau)\lvert\eta(\tau)\rvert^4$.
\end{lemma}

\begin{proof}
The product formula of Jacobi gives $\Delta(\mathbb Z+\mathbb Z\tau)=(2\pi)^{12}\eta(\tau)^{24}$, so $\Delta(\Lambda)=\omega_1^{-12}(2\pi)^{12}\eta(\tau)^{24}$ by homogeneity. Hence $\lvert\Delta(\Lambda)\rvert^{1/6}=4\pi^2\lvert\eta(\tau)\rvert^4\lvert\omega_1\rvert^{-2}$, and $A(\Lambda)=\lvert\omega_1\rvert^2\operatorname{Im}\tau$.
\end{proof}

\begin{proposition}\label{r4:R4:kronecker}
Let $d>3$ be a prime with $d\equiv3\pmod4$, let $K=\mathbb Q(\sqrt{-d})$ with class number $h$, and for each ideal class $C$ of $K$ let $\Lambda_C\subset\mathbb C$ be the image under $\varphi_0$ of a fractional ideal in $C$. Then
\[
\frac1h\sum_C\log\frac{\mathcal A(\Lambda_C)}\pi=\gamma+\log\frac{2\pi}{\sqrt d}-\lambda_1 .
\]
\end{proposition}

\begin{proof}
For $\tau$ in the upper half plane and $\operatorname{Re}s>1$ put $Z(\tau,s)=\sum'_{(m,n)}(\operatorname{Im}\tau)^s\lvert m+n\tau\rvert^{-2s}$. Kronecker's first limit formula states that
\begin{equation}\label{r4:R4:KLF}
Z(\tau,s)=\frac\pi{s-1}+2\pi\Big(\gamma-\log2-\log\big(\sqrt{\operatorname{Im}\tau}\,\lvert\eta(\tau)\rvert^2\big)\Big)+O(s-1).
\end{equation}
Let $C$ be a class, let $\mathfrak b$ be a fractional ideal in $C^{-1}$ with an oriented basis $\mathfrak b=\mathbb Z\beta_1+\mathbb Z\beta_2$, and put $\tau_C=\beta_2/\beta_1$. The integral ideals of $C$ are the ideals $x\mathfrak b^{-1}$ with $x\in\mathfrak b\setminus\{0\}$, and $x$ and $x'$ give the same ideal exactly when $x'/x$ is one of the $w=2$ units. The covolume of $\mathfrak b$ is $N(\mathfrak b)\sqrt d/2$, so for $x=m\beta_1+n\beta_2$
\[
N(x\mathfrak b^{-1})=\frac{\lvert\beta_1\rvert^2\lvert m+n\tau_C\rvert^2}{N(\mathfrak b)}=\frac{\sqrt d}{2}\,\frac{\lvert m+n\tau_C\rvert^2}{\operatorname{Im}\tau_C}.
\]
Hence the partial zeta function of $C$ is $\zeta(s,C)=\frac1w\big(\frac{\sqrt d}2\big)^{-s}Z(\tau_C,s)$. Summing over $C$ and using $\big(\frac{\sqrt d}2\big)^{-s}=\frac2{\sqrt d}\big(1-(s-1)\log\frac{\sqrt d}2+O((s-1)^2)\big)$ we obtain
\[
\zeta_K(s)=\frac2{w\sqrt d}\Big(\frac{h\pi}{s-1}+2\pi\sum_C\big(\gamma-\log2-\log(\sqrt{\operatorname{Im}\tau_C}\lvert\eta(\tau_C)\rvert^2)\big)-h\pi\log\frac{\sqrt d}2\Big)+O(s-1).
\]
On the other hand $\zeta_K(s)=\zeta(s)L(s,\chi_d)=\frac{L(1,\chi_d)}{s-1}+\gamma L(1,\chi_d)+L'(1,\chi_d)+O(s-1)$. Comparing residues gives $L(1,\chi_d)=2\pi h/(w\sqrt d)$, which is Dirichlet's class number formula. Dividing the constant terms by this value gives
\[
\gamma+\lambda_1=\frac2h\sum_C\Big(\gamma-\log2-\log\big(\sqrt{\operatorname{Im}\tau_C}\lvert\eta(\tau_C)\rvert^2\big)\Big)-\log\frac{\sqrt d}2,
\]
that is,
\[
\frac1h\sum_C\log\big(\operatorname{Im}\tau_C\,\lvert\eta(\tau_C)\rvert^4\big)=\gamma-\log2-\tfrac12\log d-\lambda_1 .
\]
The lattice $\Lambda_{C^{-1}}$ is homothetic to $\mathbb Z+\mathbb Z\tau_C$, and $C\mapsto C^{-1}$ permutes the classes. Lemma~\ref{r4:R4:areaeta} gives $\log(\mathcal A(\Lambda_{C^{-1}})/\pi)=\log(4\pi)+\log(\operatorname{Im}\tau_C\lvert\eta(\tau_C)\rvert^4)$, and the average over $C$ is $\log(4\pi)+\gamma-\log2-\frac12\log d-\lambda_1=\gamma+\log(2\pi/\sqrt d)-\lambda_1$.
\end{proof}

Fix a nonzero ideal $\mathfrak c$ of $\mathcal O_K$ prime to $2d$ and a complex number $\delta$ with $\delta^{12}=\Delta(\varphi_0(\mathfrak c))$. The lattice $\Lambda_0=\delta\,\varphi_0(\mathfrak c)$ satisfies $\Delta(\Lambda_0)=1$. Its invariants $g_2(\Lambda_0)$ and $g_3(\Lambda_0)$ are algebraic, because $g_2(\Lambda_0)^3=j(\mathfrak c)/1728$ and $27g_3(\Lambda_0)^2=g_2(\Lambda_0)^3-1$ with $j(\mathfrak c)$ algebraic. Let $M\subset\mathbb C$ be a number field that contains $g_2(\Lambda_0)$, $g_3(\Lambda_0)$ and $\tau(E)$ for every $\tau\in\Phi$. For each embedding $s\colon M\to\mathbb C$ let $\Lambda_0^s$ be the lattice with $g_2(\Lambda_0^s)=s(g_2(\Lambda_0))$ and $g_3(\Lambda_0^s)=s(g_3(\Lambda_0))$. It exists and is unique by the uniformization theorem, since $s(g_2)^3-27s(g_3)^2=s(1)=1\ne0$.

We identify $\Phi$ with the set of $K$-embeddings $\tau\colon E\to M$ that restrict to $\varphi_0$ on $K$. Let $\mathfrak d=\mathfrak d_{E/K}$ be the relative different, let $\mathfrak b$ be the ideal of $\mathcal O_E$ with $\mathfrak d=\mathfrak b\bar{\mathfrak b}$ of Lemma~\ref{hostideal}(ii), and put $\mathfrak A=\mathfrak b^{-1}$. Let $R_M\subset M^\Phi$ be the $\mathcal O_M$-module generated by the vectors $(\tau(\alpha))_{\tau\in\Phi}$ with $\alpha\in\mathfrak A$, and for $\tau\in\Phi$ let
\[
\ell_\tau=\{s\in M:\ s\,e_\tau\in R_M\},
\]
where $e_\tau$ is the coordinate vector of $\tau$. On $\ell_\tau$ we put the metrics $\lVert s\rVert_\rho^2=\pi\lvert\rho(s)\rvert^2/A(\Lambda_0^\rho)$ for the embeddings $\rho\colon M\to\mathbb C$.

\begin{lemma}\label{r4:R4:coord}
In this situation $\ell_\tau=\tau(\mathfrak A\mathfrak d)\mathcal O_M$, and $N_{M/\mathbb Q}(\ell_\tau)^{[E:\mathbb Q]}=N_{E/\mathbb Q}(\mathfrak b)^{[M:\mathbb Q]}$.
\end{lemma}

\begin{proof}
The map $E\otimes_KM\to M^\Phi$, $\alpha\otimes c\mapsto(\tau(\alpha)c)_\tau$, is an isomorphism of $M$-algebras because $E/K$ is separable and $\lvert\Phi\rvert=[E:K]$, and it takes the trace $\operatorname{Tr}_{E\otimes M/M}$ to $y\mapsto\sum_\tau y_\tau$. Under it $R_M$ is the image of $\mathfrak A\otimes_{\mathcal O_K}\mathcal O_M$, where $\mathcal O_M$ is an $\mathcal O_K$-algebra through $\varphi_0$. The trace pairing $\mathfrak A\times\mathfrak A^*\to\mathcal O_K$, $(\alpha,\beta)\mapsto\operatorname{Tr}_{E/K}(\alpha\beta)$, with $\mathfrak A^*=\mathfrak A^{-1}\mathfrak d^{-1}$, is perfect, and $\mathfrak A$ is a projective $\mathcal O_K$-module, so the base change $\mathfrak A\otimes\mathcal O_M\times\mathfrak A^*\otimes\mathcal O_M\to\mathcal O_M$ of this pairing is perfect as well. In $M^\Phi$ the base-changed pairing is $\langle x,y\rangle=\sum_\tau x_\tau y_\tau$, and $R_M$ is the set of $x$ with $\langle x,y\rangle\in\mathcal O_M$ for all $y$ in the image $R_M^*$ of $\mathfrak A^*\otimes\mathcal O_M$. For $s\in M$ and $y\in R_M^*$ we have $\langle se_\tau,y\rangle=sy_\tau$, and the coordinates $y_\tau$ of $R_M^*$ generate $\tau(\mathfrak A^*)\mathcal O_M$. So $se_\tau\in R_M$ exactly when $s\,\tau(\mathfrak A^*)\mathcal O_M\subset\mathcal O_M$, that is, when $s\in\tau(\mathfrak A\mathfrak d)\mathcal O_M$. The norm of the extension to $M$ of an ideal $\mathfrak e$ of $\tau(E)$ is $N(\mathfrak e)^{[M:\tau(E)]}$, and $N_{E/\mathbb Q}(\mathfrak A\mathfrak d)=N_{E/\mathbb Q}(\bar{\mathfrak b})=N_{E/\mathbb Q}(\mathfrak b)$.
\end{proof}

\begin{definition}\label{r4:R4:Omega}
For $\tau\in\Phi$ and $s\in\ell_\tau\setminus\{0\}$ put
\[
\log\Omega_{\tau}=\frac1{[M:\mathbb Q]}\Big(2\log[\ell_\tau:s\mathcal O_M]-\sum_{\rho\colon M\to\mathbb C}\log\lVert s\rVert_\rho^2\Big).
\]
\end{definition}

The product formula $\sum_\rho\log\lvert\rho(s)\rvert^2=2\log\lvert N_{M/\mathbb Q}(s)\rvert$ shows that $\log\Omega_\tau$ does not depend on $s$. It is $2/[M:\mathbb Q]$ times the arithmetic degree of the hermitian line $\ell_\tau$, and Proposition~\ref{r4:R4:period} shows that it depends neither on $M$, nor on $\tau$, nor on $\mathfrak c$.

\begin{proposition}\label{r4:R4:period}
Let $d\equiv7\pmod{16}$ be a prime, $K=\mathbb Q(\sqrt{-d})$, let $E\supset K$ be a CM field with $[E:K]=f$ and with an ideal $\mathfrak b$ of $\mathcal O_E$ such that $\mathfrak d_{E/K}=\mathfrak b\bar{\mathfrak b}$, and let $\mathfrak c$, $\delta$, $M$ and $\tau\in\Phi$ be as above. Then
\[
\log\Omega_\tau=\gamma+\log\frac\pi v-\lambda_1,\qquad v=\tfrac12\mathrm{rd}_E .
\]
\end{proposition}

\begin{proof}
By Lemma~\ref{r4:R4:coord}, $[\ell_\tau:s\mathcal O_M]=\lvert N_{M/\mathbb Q}(s)\rvert/N_{M/\mathbb Q}(\ell_\tau)$, and
\[
\frac2{[M:\mathbb Q]}\log N_{M/\mathbb Q}(\ell_\tau)=\frac2{[E:\mathbb Q]}\log N_{E/\mathbb Q}(\mathfrak b)=\frac1{2f}\log N_{E/\mathbb Q}(\mathfrak d_{E/K}).
\]
The discriminant of the tower $E/K/\mathbb Q$ is $\lvert d_E\rvert=d^{\,f}N_{E/\mathbb Q}(\mathfrak d_{E/K})$, so this number is $\log\mathrm{rd}_E-\frac12\log d$. With the product formula,
\[
\log\Omega_\tau=-\Big(\log\mathrm{rd}_E-\tfrac12\log d\Big)+\frac1{[M:\mathbb Q]}\sum_\rho\log\frac{A(\Lambda_0^\rho)}\pi .
\]
For each $\rho$ the lattice $\Lambda_0^\rho$ has $\Delta=1$ and $j(\Lambda_0^\rho)=\rho(j(\mathfrak c))$. By the theory of complex multiplication the numbers $j(\mathfrak a)$, for $\mathfrak a$ running over representatives of the $h$ ideal classes, are distinct and form one orbit under $\operatorname{Gal}(\bar{\mathbb Q}/\mathbb Q)$, so $j(\mathfrak c)$ has degree $h$ over $\mathbb Q$ and, as $\rho$ runs over the $[M:\mathbb Q]$ embeddings, $\rho(j(\mathfrak c))$ takes each value $j(\mathfrak a)$ exactly $[M:\mathbb Q]/h$ times. Two lattices with the same $j$-invariant are homothetic, so $\Lambda_0^\rho$ is homothetic to $\varphi_0(\mathfrak a)$ for the class of $\mathfrak a$ with $j(\mathfrak a)=\rho(j(\mathfrak c))$, and since $\Delta(\Lambda_0^\rho)=1$ we get $A(\Lambda_0^\rho)=\mathcal A(\varphi_0(\mathfrak a))$. Proposition~\ref{r4:R4:kronecker} gives
\[
\frac1{[M:\mathbb Q]}\sum_\rho\log\frac{A(\Lambda_0^\rho)}\pi=\frac1h\sum_C\log\frac{\mathcal A(\Lambda_C)}\pi=\gamma+\log\frac{2\pi}{\sqrt d}-\lambda_1 .
\]
Adding the two contributions gives $\gamma+\log(2\pi/\mathrm{rd}_E)-\lambda_1=\gamma+\log(\pi/v)-\lambda_1$.
\end{proof}

We write $\log\Omega=\gamma+\log(\pi/v)-\lambda_1$ for the common value. It is the constant of Proposition~\ref{acct} with this choice of $v$ and $\lambda_1=\ell_{\chi_1}$. The lattice $\Lambda_0$ is the period lattice of the CM curve of Section~\ref{sec:inert} in the model of discriminant one, and $\log\Omega_\tau$ compares the archimedean norm of its invariant differential with the integral structure that the model defines at the finite places. This is the comparison that the archimedean norms of the theta kernels need, because the kernel law of Section~\ref{sec:inert} expresses the kernels in the algebraic parameter of the model, while their archimedean norms are computed in the complex coordinate.

\begin{lemma}\label{r4:R4:Omegasize}
In the situation of Proposition~\ref{r4:R4:period}, suppose that $v\ge\pi$ and $\lvert\lambda_1\rvert\le L_\varepsilon-1$. Then $\lvert\log\Omega\rvert\le\Lambda_a$ for $x\ge1$.
\end{lemma}

\begin{proof}
Since $v\ge\pi$, the number $\log(v/\pi)$ lies between $0$ and $\log\mathrm{rd}_E$, so $\lvert\gamma-\log(v/\pi)\rvert\le\max(\gamma,\log\mathrm{rd}_E)\le\gamma+\log\mathrm{rd}_E$. Hence $\lvert\log\Omega\rvert\le\gamma+\log\mathrm{rd}_E+L_\varepsilon-1\le\log x+L_\varepsilon+\log\mathrm{rd}_E=\Lambda_a$, because $\gamma<1$ and $\log x\ge0$.
\end{proof}

The hypothesis $\lvert\lambda_1\rvert\le L_\varepsilon-1$ holds for large $x$, because Theorem~\ref{stage2}(ii) gives $\lvert\lambda_1\rvert\le602c_0\log x$ and $L_\varepsilon=x^\varepsilon$.

\subsection{The heights of the period maps}\label{r4:sec:R4height}

Let $F$ be the totally complex field over which the lattices of the chain are defined (the field $\mathbb L$ of Section~\ref{sec:chain}), and let $v$ run over its archimedean places. For each index $u\in W$ the period map has the form of a chain. There are hermitian spaces $V_0,\dots,V_m$ over $F$ with frames $X_i$ at the archimedean places, the source $V_0$ of $N_u$, $F$-linear maps $K_j\colon V_{j-1}\to V_j$ for $1\le j\le m$ (the lifts), and an $F$-linear map $P_u\colon V_m\to R_F^{\vee}$ (the final pairing with forms on the base side) such that
\begin{equation}\label{r4:R4:chain}
B_u=P_u\circ K_m\circ\dots\circ K_1 .
\end{equation}
The target $R_F$ has a frame $X_R$ and real numbers $\kappa_v>0$ and $\lambda_j(p,v)>0$, indexed by the lifts $j$, by the base-type subspaces $p$ of $R_F$ and by the places $v$, such that the base-type inclusion multiplies the hermitian form on the base-type subspace $p$ by $\kappa_v\prod_j\lambda_j(p,v)$ at $v$. These scalars are part of the data of the base side in Subsection~\ref{r4:sec:R4target}.

\begin{lemma}\label{r4:R4:composite}
Let $\phi_1\colon U_0\to U_1,\dots,\phi_s\colon U_{s-1}\to U_s$ be linear maps between finite-dimensional complex hermitian spaces. Then $\lVert\phi_s\circ\dots\circ\phi_1\rVert\le\prod_i\lVert\phi_i\rVert$. Let $U$ be a complex vector space with two hermitian forms $h$ and $h'$ and an orthogonal decomposition $U=\bigoplus_pU^p$ for both, with $h=c_p\,h'$ on $U^p$ for real numbers $c_p>0$. Then for every linear map $\phi$ from a hermitian space into the dual space $U^\vee$, the operator norm of $\phi$ for the dual form $h^\vee$ is at most $(\min_pc_p)^{-1/2}$ times its operator norm for $h'^\vee$.
\end{lemma}

\begin{proof}
The first statement is submultiplicativity of operator norms. For the second, a linear form $\xi=\sum_p\xi_p$ on $U$ has $h^\vee(\xi,\xi)=\sum_pc_p^{-1}h'^\vee(\xi_p,\xi_p)\le(\min_pc_p)^{-1}h'^\vee(\xi,\xi)$.
\end{proof}

\begin{proposition}\label{r4:R4:ht}
Suppose that $B_u\ne0$, that for every archimedean place $v$ of $F$
\[
\log\lVert P_u\rVert_v^2\le h_P,\qquad \log\lVert K_j\rVert_v^2\le h_{K,j}(v)\quad(1\le j\le m),
\]
and that $\sum_j\log\lambda_j(p,v)\ge\ell_\lambda$ for every nonzero base-type summand $p$ and every $v$, where the operator norm of $P_u$ is taken with respect to the frame $X_m$ and the dual of the form $h'$ that equals $(\kappa_v\prod_j\lambda_j(p,v))^{-1}$ times the form of $X_R$ on the base-type subspace $p$, and the operator norm of $K_j$ with respect to $X_{j-1}$ and $X_j$. Then
\[
\mathrm{ht}(B_u)\le h_P+\sum_{j=1}^m\frac1{n_F}\sum_vh_{K,j}(v)-\ell_\lambda-\langle\log\kappa\rangle,
\]
where $\mathrm{ht}$ is computed with the frame $X_R$.
\end{proposition}

\begin{proof}
At a place $v$, Lemma~\ref{r4:R4:composite} applied to \eqref{r4:R4:chain}, with $U=R_F\otimes_v\mathbb C$, its decomposition into base-type subspaces, $h$ the form of the frame $X_R$, and $c_p=\kappa_v\prod_j\lambda_j(p,v)$, gives
\[
\log\lVert B_u\rVert_v^2\le h_P+\sum_jh_{K,j}(v)-\log\kappa_v-\min_p\sum_j\log\lambda_j(p,v)\le h_P+\sum_jh_{K,j}(v)-\log\kappa_v-\ell_\lambda .
\]
Taking the mean over the $n_F$ places gives the statement.
The minimum is over the nonzero summands. Zero summands have zero
projection in the dual-norm calculation. Since $B_u\ne0$, at least
one summand is nonzero.
\end{proof}

The three hypotheses of Proposition~\ref{r4:R4:ht} are the archimedean content of the period maps, and we now state them in the form in which they enter \eqref{A4}. For a lift $j$ and an embedding $\rho\in\Phi$ the kernel of $K_j$ at $\rho$ is a polynomial in the matrix coordinates, and its row weights form a vector $\nu_{j\rho}=(\nu_{j\rho1},\dots,\nu_{j\rho N})$ of nonnegative reals with $N=k+m(t_r+1)$. Let $h_{j\rho}$ be the logarithm of the normalizing constant of the frame on the $\rho$-component of $V_j$, so that $\sum_\rho h_{j\rho}$ is the quantity $\ell_\lambda$ contributed by the lift $j$ (the hook constants of the base side). Let $C_0\ge1$ be a constant.

\begin{hypothesis}[kernel norms]\label{r4:R4:HK}
For a fixed word $u$, write $s_j=s_j(u)$ for the following bound. For every lift $j$, $\frac1{n_F}\sum_vh_{K,j}(v)\le s_j$ with
\[
s_j=\sum_{\rho\in\Phi}\Big(\sum_{i=1}^N\ph(\nu_{j\rho i})-\Big(\sum_{i=1}^N\nu_{j\rho i}\Big)\log\Omega\Big)+C_0\big(fN^2\Lambda_a+fm^3N^2L_\varepsilon\big).
\]
\end{hypothesis}

\begin{hypothesis}[final pairing]\label{r4:R4:HP}
$h_P\le-\sum_{t\in T}a_t\log T_u+e_P$ with $e_P\le C_0\big(N_Tb'N+f(m\eta+S)N+fN^2\big)\Lambda_a$, where $\eta=W+mh_0$ is the spread of Section~\ref{sec:auxiliary} and $S\ge0$ bounds the small weights.
\end{hypothesis}

\begin{hypothesis}[frame constants]\label{r4:R4:HL}
$\ell_\lambda\ge\sum_j\sum_\rho h_{j\rho}-e_\lambda$ with $e_\lambda\le C_0mfN^2\Lambda_a$.
\end{hypothesis}

\begin{remark}\label{r4:R4:HKmech}
Hypothesis~\ref{r4:R4:HK} is the place where the factor $m$ in front of $\ph(a_t)$ in \eqref{A4} arises. Each lift is a finite sum over classes of local theta kernels, as in (D1). At an archimedean place the local kernel restricted to a $K$-type is a polynomial whose monomials of row degree $\nu$ have Fock norm $\prod_i\nu_i!$, and $\log\nu!=\ph(\nu)+O(\log(\nu+2))$. The coordinates in which the kernel law of Section~\ref{sec:inert} makes the kernel integral differ from the complex coordinates by the period of the CM curve, and Proposition~\ref{r4:R4:period} computes the mean over the embeddings of the logarithm of this period. The class sum costs the logarithm of the number of classes, which is $O(fm^2N^2(L_\varepsilon+\log T_u))$, and this is the second term of $s_j$. Hypothesis~\ref{r4:R4:HP} bounds the pairing with a form of level $T_u$, whose branching to the smaller group divides the norm by $T_u$ to the power of the degree, and the degree at a top $t$ is $a_t$. Subsection~\ref{r4:sec:R4open} lists the status of these three bounds in this section.
\end{remark}

\begin{definition}\label{r4:R4:hu}
Put $e_P^{\max}=C_0(N_Tb'N+f(m\eta+S)N+fN^2)\Lambda_a$ and $e^{\max}_{\lambda,j}=C_0fN^2\Lambda_a$, and
\[
h^0_u=-\sum_{t\in T}a_t\log T_u+e_P^{\max}+\sum_{j=1}^m\Big(s_j-\sum_\rho h_{j\rho}+e^{\max}_{\lambda,j}\Big),\qquad h_u=h^0_u-\langle\log\kappa\rangle .
\]
\end{definition}

\begin{corollary}\label{r4:R4:htcor}
Under Hypotheses~\ref{r4:R4:HK}, \ref{r4:R4:HP} and~\ref{r4:R4:HL}, $\mathrm{ht}(B_u)\le h_u$ for every $u\in W$ with $B_u\ne0$.
\end{corollary}

\begin{proof}
Proposition~\ref{r4:R4:ht} with the three hypotheses gives $\mathrm{ht}(B_u)\le-\sum_ta_t\log T_u+e_P+\sum_js_j-\sum_{j,\rho}h_{j\rho}+e_\lambda-\langle\log\kappa\rangle$, and $e_P\le e_P^{\max}$, $e_\lambda\le\sum_je^{\max}_{\lambda,j}$.
\end{proof}

The number $h_u$ is the height of Proposition~\ref{acct}. It enters \eqref{A3} through Corollary~\ref{r4:R4:htcor} and \eqref{A4} through Theorem~\ref{r4:R4:A4} below.

\subsection{The slope of the target lattice}\label{r4:sec:R4target}

The target lattice is a lattice $R$ in an $F$-vector space $R_F$ of dimension $D$, defined place by place by the local types of the forms on the top group. Its dual $R^\vee$ is the target of the maps $B_u$. We put
\[
\mu^C=\mu(R^\vee)\quad\text{for the frame }X_R .
\]
The height $h_u$ of Definition~\ref{r4:R4:hu} bounds $\mathrm{ht}(B_u)$ for the same frame by Corollary~\ref{r4:R4:htcor}, so the slope inequality \eqref{A3} for the maps $B_u$ and the inequality \eqref{A4} use the same pair $(\mu^C,h_u)$. Inequality \eqref{A4} involves $\mu^C$ and $h_u$ through the sum $\mu^C+h_u=(\mu^C-\langle\log\kappa\rangle)+h^0_u$.

The base side consists of lattices $M_j\subset W_j$ for $1\le j\le m$, where $W_j$ is a space of automorphic forms on a unitary group of signature $(p_j,q_j)$ at the places of $\Phi$, of weight $w_{j\rho}$ at $\rho$, and of linear maps $\theta_j$ from $R_F$ to $W_j$ that multiply the hermitian forms of the base-type subspaces by $\lambda_j(p,v)$. The local terms at the twisted places are the numbers
\begin{equation}\label{r4:R4:calL}
\mathcal L=\sum_{j=1}^m\sum_{y\in\mathrm{Tw}}2\,\ell_y\,e_y\,F_j(y),
\end{equation}
where $\mathrm{Tw}$ is the finite set of places where the local type of the forms is twisted, $\ell_y=\log N(y)/[F:\mathbb Q]$, $e_y$ is the ramification index and $F_j(y)\ge0$ is the local length of the lift $j$ at $y$. This is the quantity $\mathcal L$ of Proposition~\ref{acct}, and it is the same quantity in \eqref{A2}.

\begin{hypothesis}[target slope]\label{r4:R4:HT}
There are real numbers $\mu_0$, $P_{\mathrm{pass}}$ and $E_b$ with
\begin{align*}
\mu^C-\langle\log\kappa\rangle&\le\mu_0+P_{\mathrm{pass}},& P_{\mathrm{pass}}&\le1+C_0\lambda_0,\\
\mu_0&\le\sum_{j=1}^m\sum_{\rho\in\Phi}(p_j-q_j)\,w_{j\rho}\log\Omega+\mathcal L+E_b,& E_b&\le C_0\,m\,x^{\varepsilon}(B_0+N_c),
\end{align*}
where $\lambda_0\le b$, $\lvert p_j-q_j-k\rvert\le m$, and $B_0=fmb\,x^\varepsilon$ and $N_c=N_2kb\,x^{2\varepsilon}/n^2$ are the two terms of $\mathcal B$ in \eqref{calB}.
\end{hypothesis}

\begin{remark}\label{r4:R4:HTmech}
Hypothesis~\ref{r4:R4:HT} is the composite of three statements. The first, the passage, compares the slope of $R^\vee$ at the top level of the chain with the slope of the corresponding lattice at the level of the base side, at a cost $3\delta(\mu+H_0+28\lambda_0)+28\lambda_0$ with $\delta\le1/3$. The second bounds the slope of the dual of the base-side target, minus $\langle\log\kappa\rangle$, by $\sum_j\operatorname{lcv}(M_j)/D+\mathcal L+E$, where $\operatorname{lcv}(M_j)$ is the normalized logarithmic covolume of $M_j$ and $E$ collects the local indices at the twisted places. The third is the covolume bound $\sum_j\operatorname{lcv}(M_j)/D\le\sum_{j,\rho}(p_j-q_j)w_{j\rho}\log\Omega+E_{l3}$, which is the arithmetic degree of the automorphic line bundle of weight $w_j$ on the unitary Shimura variety of signature $(p_j,q_j)$, evaluated through CM points where Proposition~\ref{r4:R4:period} applies, with the flat part of degree $\frac12(p_j-q_j)\log\Omega$ per unit of weight. The local indices at the twisted places are bounded by the Galois-theoretic bounds of the twisted places, which is the input $E_b\le C_0mx^\varepsilon(B_0+N_c)$. Subsection~\ref{r4:sec:R4open} lists which of these statements are proved here.
\end{remark}

\subsection{Hypothesis (A4)}\label{r4:sec:R4A4}

The weights of the kernels satisfy two relations, which express the degree of the kernel of the lift $j$ at the embedding $\rho$ in terms of the nodes and the tops.

\begin{hypothesis}[weights of the kernels]\label{r4:R4:HW}
For all $j$ and $\rho$,
\begin{align*}
\sum_{i=1}^N\nu_{j\rho i}&=\mathbf 1_T(\rho)\,a_\rho+k\,w_{j\rho}+\delta_{j\rho},\\
\lvert\delta_{j\rho}\rvert&\le C_0\big(w_{j\rho}+k(N+S)+m(t_r+1)\eta+\mathbf 1_T(\rho)kb'\big),\\
\sum_{i=1}^N\ph(\nu_{j\rho i})-h_{j\rho}&\le\mathbf 1_T(\rho)\ph(a_\rho)+\epsilon_{j\rho},\\
\epsilon_{j\rho}&\le C_0\big(w_{j\rho}+S+N^2+m(t_r+1)\eta\big)\Lambda_a,
\end{align*}
where $\mathbf 1_T(\rho)=1$ if $\rho$ is a top and $0$ otherwise, and $a_\rho=a_t$ at the top $\rho=t$.
\end{hypothesis}

\begin{theorem}\label{r4:R4:A4}
Let $x\ge3$, $0<\varepsilon<1$, $L_\varepsilon=x^\varepsilon$, $b=e^{L_\varepsilon}$, $a=xb$, and suppose that $\log x\le L_\varepsilon$, $1\le\mathrm{rd}_E$, $\log\mathrm{rd}_E\le L_\varepsilon$. Let $m,f,k,N\ge1$, $S,t_r,\eta,b',W''\ge0$ and $C_0\ge1$ satisfy
\begin{gather}
0\le w_{j\rho}\le2b,\quad \lvert\log\Omega\rvert\le\Lambda_a,\quad N^2\le b,\quad m^3N^2L_\varepsilon\le b,\quad m\eta NL_\varepsilon\le b,\quad kN\le b,\label{r4:R4:sz1}\\
kS\le W'',\quad m(t_r+1)\eta\le W'',\quad NW''\le b,\quad b'N\le a,\quad mkb'\le a.\label{r4:R4:sz2}
\end{gather}
Assume Hypotheses~\ref{r4:R4:HK}, \ref{r4:R4:HP}, \ref{r4:R4:HL}, \ref{r4:R4:HT} and~\ref{r4:R4:HW}, and let $h_u$ be as in Definition~\ref{r4:R4:hu}. Put $e_4=(50C_0+10)\big(N_Ta\Lambda_a+m\,x^\varepsilon(B_0+N_c)\big)$. Then for every $u\in W$
\[
\mu^C+h_u-\mathcal L\le\sum_{t\in T}\big(m\ph(a_t)-ma_t\log\Omega-a_t\log T_u\big)+e_4 .
\]
\end{theorem}

\begin{proof}
Write $\sum_{j,\rho}$ for the sum over $1\le j\le m$ and $\rho\in\Phi$. Since $\mu^C+h_u=(\mu^C-\langle\log\kappa\rangle)+h^0_u$, Hypothesis~\ref{r4:R4:HT} and Definition~\ref{r4:R4:hu} give
\[
\mu^C+h_u-\mathcal L\le\sum_{j,\rho}(p_j-q_j)w_{j\rho}\log\Omega+E_b+P_{\mathrm{pass}}-\sum_ta_t\log T_u+e_P^{\max}+\sum_j s_j-\sum_{j,\rho}h_{j\rho}+\sum_je^{\max}_{\lambda,j}.
\] Insert the definition of $s_j$ and the two relations of Hypothesis~\ref{r4:R4:HW}. For each $j$ and $\rho$,
\begin{align*}
&\sum_i\ph(\nu_{j\rho i})-h_{j\rho}-\Big(\sum_i\nu_{j\rho i}\Big)\log\Omega+(p_j-q_j)w_{j\rho}\log\Omega\\
&\qquad\le\mathbf 1_T(\rho)\big(\ph(a_\rho)-a_\rho\log\Omega\big)+\epsilon_{j\rho}+(p_j-q_j-k)w_{j\rho}\log\Omega-\delta_{j\rho}\log\Omega .
\end{align*}
Summing over $j$ and $\rho$, the first term gives $m\sum_{t\in T}(\ph(a_t)-a_t\log\Omega)$, and the inequality to prove reduces to
\[
\Sigma_1+\Sigma_2\le e_4,\qquad \Sigma_1=\sum_{j,\rho}\big(\epsilon_{j\rho}+m\,w_{j\rho}\Lambda_a+\lvert\delta_{j\rho}\rvert\Lambda_a\big),
\]
where $\Sigma_2=E_b+P_{\mathrm{pass}}+e_P^{\max}+\sum_jC_0(fN^2\Lambda_a+fm^3N^2L_\varepsilon)+\sum_je^{\max}_{\lambda,j}$, and where we used $\lvert p_j-q_j-k\rvert\le m$ and $\lvert\log\Omega\rvert\le\Lambda_a$.

The size conditions give $S\le kS\le W''\le b/N\le b$, $N^2\le b$, $m(t_r+1)\eta\le W''\le b$ and $k(N+S)\le kN+W''\le2b$. Hence $w_{j\rho}+S+N^2+m(t_r+1)\eta\le5b$ and $w_{j\rho}+k(N+S)+m(t_r+1)\eta\le5b$, so for each pair $(j,\rho)$
\[
\epsilon_{j\rho}+mw_{j\rho}\Lambda_a+\lvert\delta_{j\rho}\rvert\Lambda_a\le\big(10C_0+2m\big)b\,\Lambda_a+\mathbf 1_T(\rho)C_0kb'\Lambda_a .
\]
There are $m$ lifts and $f$ embeddings, and $mkb'\le a$, so $\Sigma_1\le(10C_0+2)m^2fb\Lambda_a+C_0N_Ta\Lambda_a$. Since $\log x\le L_\varepsilon$ and $\log\mathrm{rd}_E\le L_\varepsilon$ we have $\Lambda_a=\log x+L_\varepsilon+\log\mathrm{rd}_E\le3L_\varepsilon=3x^\varepsilon$, and $m^2fbx^\varepsilon=mB_0$. Therefore $\Sigma_1\le3(10C_0+2)mB_0+C_0N_Ta\Lambda_a$.

For $\Sigma_2$: $E_b\le C_0mx^\varepsilon(B_0+N_c)$ and $P_{\mathrm{pass}}\le1+C_0b\le(1+C_0)B_0$ because $B_0\ge b\ge1$. Next $N_Tb'N\Lambda_a\le N_Ta\Lambda_a$; $fm\eta N\Lambda_a\le fb\Lambda_a$ because $m\eta N\le b/L_\varepsilon\le b$; $fSN\Lambda_a\le fb\Lambda_a$ because $SN\le W''N\le b$; and $fN^2\Lambda_a\le fb\Lambda_a$. So $e_P^{\max}\le C_0N_Ta\Lambda_a+3C_0fb\Lambda_a\le C_0N_Ta\Lambda_a+9C_0B_0$. Finally $mC_0(fN^2\Lambda_a+fm^3N^2L_\varepsilon)\le C_0(3mfbx^\varepsilon+mfb)\le4C_0B_0$ and $\sum_je^{\max}_{\lambda,j}=C_0mfN^2\Lambda_a\le3C_0B_0$. Collecting terms and using $B_0\le mx^\varepsilon B_0$,
\[
\Sigma_1+\Sigma_2\le2C_0N_Ta\Lambda_a+(48C_0+7)\,m\,x^\varepsilon(B_0+N_c)\le e_4 .\qedhere
\]
\end{proof}

\begin{corollary}\label{r4:R4:A4cor}
Under the hypotheses of Theorem~\ref{r4:R4:A4}, inequality \eqref{A4} of Proposition~\ref{acct} holds for every $u\in W$ with $\mu^C=\mu(R^\vee)$, the heights $h_u$ of Definition~\ref{r4:R4:hu}, the local terms $\mathcal L$ of \eqref{r4:R4:calL}, $\log\Omega=\gamma+\log(\pi/v)-\lambda_1$ with $v=\frac12\mathrm{rd}_E$ and $\lambda_1=\ell_{\chi_1}$, and
\[
e_4=(50C_0+10)\Big(N_Ta\Lambda_a+m\,x^\varepsilon(B_0+N_c)\Big).
\]
\end{corollary}

\begin{proof}
The right side of \eqref{A4} is $\sum_{t=1}^{N_T}(m\ph(a_t)-ma_t\log\Omega-a_t\log T_u)+e_4$, and Theorem~\ref{r4:R4:A4} is this inequality. Proposition~\ref{r4:R4:period} identifies $\log\Omega$.
\end{proof}

\subsection{Hypothesis (A6)}\label{r4:sec:R4A6}

Hypothesis \eqref{A6} is a lower bound for the auxiliary prime $T_u$. The prime $T_u$ is chosen so that the rank of the lattice $N_u$, which grows like $T_u^{kf}$, matches the dimension budget of the target lattice. We first define the two numbers that are matched.

For a strictly decreasing vector $y=(y_0,\dots,y_{s-1})$ of reals put
\[
\operatorname{vand}(y)=\prod_{0\le i<i'<s}(y_i-y_{i'}),\qquad \dim(y)=\frac{\operatorname{vand}(y)}{\prod_{i<s}i!} .
\]
For an integral dominant weight $y$ of $GL_s$, Weyl's dimension formula says that $\dim(y+\rho_s)$, with $\rho_s=(s-1,\dots,0)$, is the dimension of the irreducible representation of highest weight $y$. For $M\ge1$ let $\operatorname{mass}(M)$ be the mass of the standard unimodular hermitian lattice $\mathcal O_E^M$, that is, the sum of $1/\lvert\operatorname{Aut}\rvert$ over the classes in its genus.

For a prime $g\ge3$ and integers $k\ge2$, $m\ge2$ even, with $r=k+m-1$ and $r_S=m/2-1$, let $I_\Pi(g)=\lvert GL_k(\mathbb F_g)\rvert/\lvert U_k(\mathbb F_g)\rvert$, where $U_k$ is the group of upper unitriangular matrices, and let $I_\Sigma(g)=\lvert GL_r(\mathbb F_g)\rvert/\lvert P(\mathbb F_g)\rvert$, where $P$ is the group of upper unitriangular $r\times r$ matrices whose entries in the column $r_S+k$ vanish in the rows $r_S,\dots,r_S+k-1$. These are the indices of the level structures at $g$ of the column groups and of the row group. Put $I_\Pi=(I_\Pi(g_0)I_\Pi(g_1))^f$ and $I_\Sigma=(I_\Sigma(g_0)I_\Sigma(g_1))^f$.

\begin{lemma}\label{r4:R4:index}
For a prime $g\ge3$, $\frac12g^{k(k+1)/2}\le I_\Pi(g)\le g^{k(k+1)/2}$ and $\frac12g^{e}\le I_\Sigma(g)\le g^{e}$ with $e=\frac12(r^2+3k+1+2r_S)$. Moreover $e-\frac12k(k+1)=\frac12m(2k+m-1)$, and
\[
\frac{I_\Sigma}{I_\Pi}\ge2^{-2f}(g_0g_1)^{fm(2k+m-1)/2}.
\]
\end{lemma}

\begin{proof}
We have $\lvert GL_s(\mathbb F_g)\rvert=g^{s(s-1)/2}\prod_{i=1}^s(g^i-1)$ and $1\ge\prod_i(1-g^{-i})\ge1-\sum_{i\ge1}g^{-i}=1-\frac1{g-1}\ge\frac12$. The group $U_k$ has order $g^{k(k-1)/2}$, so $I_\Pi(g)=\prod_{i=1}^k(g^i-1)$ lies between $\frac12g^{k(k+1)/2}$ and $g^{k(k+1)/2}$. The group $P$ consists of the unitriangular matrices with $k$ entries above the diagonal fixed to zero, so $\lvert P\rvert=g^{r(r-1)/2-k}$ and $I_\Sigma(g)=g^k\prod_{i=1}^r(g^i-1)$ lies between $\frac12g^{r(r+1)/2+k}$ and $g^{r(r+1)/2+k}$. Since $r=k+1+2r_S$ we have $r(r+1)/2+k=\frac12(r^2+r+2k)=e$. Next $2e-k(k+1)=r^2+3k+1+2r_S-k^2-k=(k+m-1)^2-k^2+k+m-1=m(2k+m-1)$. Finally $I_\Sigma/I_\Pi\ge\prod_{g\in\{g_0,g_1\}}\big(\tfrac12g^e/g^{k(k+1)/2}\big)^f$.
\end{proof}

For each $\rho\in\Phi$ let the row parameter of the target group be
\[
\lambda^{\mathrm{row}}_\rho=(v_{\rho,0},\dots,v_{\rho,k-2},-b_{\rho,0},\dots,-b_{\rho,m-1})\in\mathbb R^r,
\]
and for each $u\in W$ let $\lambda^{\mathrm{col}}_{u,\rho}=(u_{\rho,0},\dots,u_{\rho,k-1})$ be the column parameter of the source group at level $T_u$. The data of Section~\ref{sec:auxiliary} and Proposition~\ref{nodes} give real numbers $\Delta\ge1$ and these parameters with the following properties, for every $\rho$ (we drop $\rho$ and $u$ from the notation).
\begin{enumerate}
\item[(P1)] $0\le v_i\le4k\Delta$ for $i\le k-2$, and $v_i-v_{i+1}=4\Delta$ for $i\le k-3$.
\item[(P2)] $v_i<u_i$ and $u_{i+1}<v_i$ for $i\le k-2$, and $-b_0<u_{k-1}$.
\item[(P3)] Either $\rho$ is a top, $u_0\le2a$ and $-a\le u_{k-1}$, or $\rho$ is not a top, $4k\Delta\le u_0-v_0\le8k\Delta$ and $u_{k-1}=v_{k-2}-\frac12$.
\item[(P4)] $b_j<w_{j\rho}$ and $\lvert w_{j\rho}-b_j\rvert\le b_j/x$ for $j<m$, $w_{j\rho}<b_{j+1}$ and $b_{j+1}-b_j\ge1$ for $j<m-1$, $4k\Delta\le b_0/x$, and $w_{m-1,\rho}\le2b$.
\end{enumerate}
Here the nodes $w_{j\rho}$ are indexed by $0\le j<m$, and $w_{j\rho}$ is the node of the lift $j+1$ in the notation of \eqref{A6}.

\begin{definition}\label{r4:R4:Xu}
Put
\[
d^*_u=\operatorname{mass}(k)\,I_\Pi\prod_{\rho\in\Phi}\dim(\lambda^{\mathrm{col}}_{u,\rho}),\qquad \log X_u=\frac1{fk}\log\frac{D^*}{d^*_u},
\]
where $D^*$ is the dimension budget of the target lattice.
\end{definition}

\begin{hypothesis}[dimension budget]\label{r4:R4:HD}
$D^*\ge\big(\tfrac23\big)^f\operatorname{mass}(r)\,I_\Sigma\prod_{\rho\in\Phi}\dim(\lambda^{\mathrm{row}}_\rho)$.
\end{hypothesis}

The number $d^*_u(T^k-1)^f$ is the dimension of the space of forms of level $T$ on the source group when the stabilizers of the classes are trivial, by the class number formula $\#\mathrm{Cl}=\operatorname{mass}(k)I_\Pi(T^k-1)^f$ and Weyl's dimension formula. The budget $D^*$ is the product of the volume of the target group, the index $I_\Sigma$ and the formal degrees of the discrete series with parameters $\lambda^{\mathrm{row}}_\rho$. For the Haar measure of the compact dual, Harish-Chandra's formula gives the formal degree $\dim(\lambda^{\mathrm{row}}_\rho)$ exactly, and the volume equals $\operatorname{mass}(r)$ when the hermitian space of the target is locally isomorphic at every finite place to the standard unimodular space. These two published inputs give Hypothesis~\ref{r4:R4:HD} with the factor $1$ in place of $(2/3)^f$.

\begin{lemma}\label{r4:R4:vand}
Let $k\ge2$, $m\ge1$, $\Delta\ge1$, $x\ge1$, $a\ge1$, and let $v,u,b,w$ satisfy (P1) to (P4) with $b\le a$. Then $\operatorname{vand}(u)>0$, $\operatorname{vand}(\lambda^{\mathrm{row}})>0$ and
\[
\log\operatorname{vand}(\lambda^{\mathrm{row}})-\log\operatorname{vand}(u)\ge(k-1)\sum_{j<m}\log b_j-2k\log(8a)-k(1+\log k).
\]
\end{lemma}

\begin{proof}
The entries of $\lambda^{\mathrm{row}}$ are strictly decreasing by (P1), (P4) and $v_{k-2}\ge0>-b_0$, and those of $u$ are strictly decreasing by (P2), so both Vandermonde products are positive. The differences of entries of $\lambda^{\mathrm{row}}$ are $v_i-v_{i'}$, $v_i+b_j$ and $b_{j'}-b_j$ for $i<i'$ and $j<j'$, and $b_{j'}-b_j\ge1$ and $v_i+b_j\ge b_j$. Hence
\[
\operatorname{vand}(\lambda^{\mathrm{row}})\ge\operatorname{vand}(v)\prod_{j<m}b_j^{\,k-1},\qquad v=(v_0,\dots,v_{k-2}).
\]
For $1\le i<i'\le k-2$, (P2) gives $u_i-u_{i'}<v_{i-1}-v_{i'}=(i'-i+1)4\Delta$, while $v_{i-1}-v_{i'-1}=(i'-i)4\Delta$. The pairs $(i-1,i'-1)$ run over the pairs $0\le p<q\le k-3$, and the remaining factors $v_p-v_{k-2}$ of $\operatorname{vand}(v)$ are at least $4\Delta\ge1$. Therefore
\[
\prod_{1\le i<i'\le k-2}(u_i-u_{i'})\le\operatorname{vand}(v)\prod_{d=1}^{k-3}\Big(\frac{d+1}d\Big)^{k-2-d}\le\operatorname{vand}(v)\,e^{k(1+\log k)},
\]
since $\sum_{d\le k}(k-2-d)\log(1+1/d)\le k\sum_{d\le k}1/d\le k(1+\log k)$. The other factors of $\operatorname{vand}(u)$ are $u_0-u_{i'}$ for $1\le i'\le k-1$ and $u_i-u_{k-1}$ for $1\le i\le k-2$, at most $2k$ factors. Each lies in $(0,8a]$. In the top case all $u_i$ lie in $[-a,2a]$. In the other case $u_0\le v_0+8k\Delta\le12k\Delta\le3b_0/x\le6b\le6a$ by (P3) and (P4), and $u_{k-1}\ge-\frac12$, so all $u_i$ lie in $[-\frac12,6a]$. Hence $\operatorname{vand}(u)\le\operatorname{vand}(v)e^{k(1+\log k)}(8a)^{2k}$, and taking logarithms gives the statement.
\end{proof}

\begin{proposition}\label{r4:R4:weyl}
Let $0<\varepsilon<1$ and $C_1\ge0$. There is a constant $C_6$, depending on $\varepsilon$ and $C_1$ alone, with the following property. Let $x\ge1$, $L_\varepsilon=x^\varepsilon$, $b=e^{L_\varepsilon}$, $a=xb$, let $f\ge1$, $m\ge2$ even with $m\le x$, $2\le k\le x^{C_1}$, let $\pi\le v$ with $\log v\le L_\varepsilon$, and let $g_0,g_1\ge3$ be primes. Suppose that the parameters of every $\rho\in\Phi$ satisfy (P1) to (P4), that
\[
\frac{\operatorname{mass}(M)}{\operatorname{mass}(M-1)}=\frac{\big((M-1)!\,(v/\pi)^M\big)^f}{(2v)^{f/2}}\,\mathbf L(M),\qquad\lvert\log\mathbf L(M)\rvert\le f\log\zeta(M)\qquad(k<M\le r),
\]
and that $D^*$ satisfies Hypothesis~\ref{r4:R4:HD}. Then for every $u$
\[
\log X_u\ge\sum_{j<m}\frac1f\sum_{\rho\in\Phi}\log\frac{w_{j\rho}\,v}\pi+m\log(g_0g_1)-C_6\Big(L_\varepsilon+\frac{m^2L_\varepsilon}k\Big).
\]
\end{proposition}

\begin{proof}
By Hypothesis~\ref{r4:R4:HD} and the definition of $d^*_u$,
\[
\log\frac{D^*}{d^*_u}\ge f\log\tfrac23+\log\frac{\operatorname{mass}(r)}{\operatorname{mass}(k)}+\log\frac{I_\Sigma}{I_\Pi}+\sum_\rho\big(\log\dim(\lambda^{\mathrm{row}}_\rho)-\log\dim(\lambda^{\mathrm{col}}_{u,\rho})\big).
\]
We treat the four terms in turn. The mass ratio telescopes over $M=k+1,\dots,r$, and there are $r-k=m-1$ such $M$:
\[
\log\frac{\operatorname{mass}(r)}{\operatorname{mass}(k)}=\sum_{M=k+1}^r\Big(f\log(M-1)!+fM\log\frac v\pi-\frac f2\log(2v)+\log\mathbf L(M)\Big).
\]
Since $\log(v/\pi)\ge0$, $\sum_MfM\log(v/\pi)\ge f(m-1)k\log(v/\pi)$, and $\log\mathbf L(M)\ge-f\log\zeta(2)\ge-f$. For the dimensions, $\log\dim(\lambda^{\mathrm{row}})-\log\dim(\lambda^{\mathrm{col}})=\log\operatorname{vand}(\lambda^{\mathrm{row}})-\log\operatorname{vand}(\lambda^{\mathrm{col}})-\sum_{M=k}^{r-1}\log M!$, and summed over the $f$ embeddings the factorials cancel the terms $f\log(M-1)!$ of the mass ratio exactly. Lemma~\ref{r4:R4:index} bounds the index term below by $-2f\log2+fmk\,g$, with $g=\log(g_0g_1)\ge0$, because $m(2k+m-1)/2\ge mk$. Lemma~\ref{r4:R4:vand} bounds the Vandermonde terms. Dividing by $fk$ we obtain
\[
\log X_u\ge\frac{k-1}{k}\cdot\frac1f\sum_\rho\sum_{j<m}\log b_{\rho,j}+(m-1)\log\frac v\pi+mg-\mathcal E,
\]
with
\[
\mathcal E=2\log(8a)+1+\log k+\frac{(m-1)\log(2v)}{2k}+\frac{m-1}k+\frac{\log(3/2)+2\log2}k .
\]
By (P4), $w_{j\rho}\le b_{\rho,j}(1+1/x)$, so $\log b_{\rho,j}\ge\log w_{j\rho}-1/x$, and $1\le b_{\rho,j}\le w_{j\rho}\le2b$, so $\frac1k\sum_{j<m}\log b_{\rho,j}\le\frac mk(L_\varepsilon+1)$. Hence
\[
\frac{k-1}k\cdot\frac1f\sum_{\rho,j}\log b_{\rho,j}\ge\sum_{j<m}\frac1f\sum_\rho\log w_{j\rho}-\frac mx-\frac mk(L_\varepsilon+1),
\]
and $(m-1)\log(v/\pi)\ge m\log(v/\pi)-L_\varepsilon$. It remains to bound $\mathcal E+\frac mx+\frac mk(L_\varepsilon+1)+L_\varepsilon$ by $C_6(L_\varepsilon+m^2L_\varepsilon/k)$. We have $\log(8a)=\log8+\log x+L_\varepsilon$, $\log x\le L_\varepsilon/(e\varepsilon)$ for $x\ge1$, $\log k\le C_1\log x$, $\log(2v)\le1+L_\varepsilon$ and $m/x\le1$. Since $L_\varepsilon\ge1$ and $m\le m^2$, each term is at most a constant depending on $\varepsilon$ and $C_1$ times $L_\varepsilon+m^2L_\varepsilon/k$.
\end{proof}

\begin{theorem}\label{r4:R4:A6}
Assume the hypotheses of Proposition~\ref{r4:R4:weyl}, let $\Gamma\ge1$ with $\log\Gamma\le fk$, and suppose that for every $u\in W$ the prime $T_u$ satisfies $T_u\ge X_u\Gamma^{-2/(fk)}$. Then \eqref{A6} holds for every $u\in W$ with $g=\log(g_0g_1)$, with $v=\frac12\mathrm{rd}_E$, with the nodes $w_{j\rho}$ of the lifts, and with
\[
e_6=C_6\Big(L_\varepsilon+\frac{m^2L_\varepsilon}k\Big)+2 .
\]
\end{theorem}

\begin{proof}
$\log T_u\ge\log X_u-\frac2{fk}\log\Gamma\ge\log X_u-2$, and Proposition~\ref{r4:R4:weyl} bounds $\log X_u$ below by the right side of \eqref{A6} plus $e_6-2$.
\end{proof}

\begin{remark}\label{r4:R4:Tu}
The prime $T_u$ is the auxiliary prime of the index $u$. It must also satisfy $T_u\le(1+\frac1{fk})X_u\Gamma^{-2/(fk)}$, so that the rank of $N_u$ matches $D^*\Gamma^{-2}$ up to a bounded factor, as the frame condition of \eqref{A3} requires, and it must split in $K$ as $T_u=\pi\bar\pi$ with $\pi\equiv1$ modulo $Q_Sqd(g_0g_1)^2$. The existence of such a prime in the interval $[X_u\Gamma^{-2/(fk)},(1+\frac1{fk})X_u\Gamma^{-2/(fk)}]$ is a statement on primes in a ray class of $K$ in a short interval, and the choice belongs to the construction of the data. The upper bound $\log T_u\le C\,mL_\varepsilon$ follows from the same comparison of dimensions and is used in Subsection~\ref{r4:sec:R4A7}. Here $\Gamma=\exp(C_\Gamma m^5fL_\varepsilon)$ is the factor of the frame condition, and $\log\Gamma\le fk$ is the size condition $C_\Gamma m^5L_\varepsilon\le k$.
\end{remark}

\subsection{Hypothesis (A7)}\label{r4:sec:R4A7}

The base lattice $M_u$ is the lattice of forms of level zero of the chain for the index $u$. We first recall the integral theory of polynomial representations that describes it.

Let $s\ge1$ and let $\mathbb Z[x]=\mathbb Z[x_{ab}:1\le a,b\le s]$ with $GL_s$ acting on the rows. On $\mathbb Q[x]$ the Fock form
\[
\{p,q\}=\sum_\alpha p_\alpha q_\alpha\,\alpha!,\qquad\alpha!=\prod_{a,b}\alpha_{ab}!,
\]
where $p=\sum_\alpha p_\alpha x^\alpha$, is positive definite, and over $\mathbb C$ its hermitian version is invariant under the unitary group acting on rows. For a partition $\nu=(\nu_1\ge\dots\ge\nu_s\ge0)$ let $hw_\nu=\prod_{i=1}^s\Delta_i^{\nu_i-\nu_{i+1}}$, with $\Delta_i$ the leading principal $i\times i$ minor and $\nu_{s+1}=0$. The \emph{Weyl module} $V(\nu)_{\mathbb Z}$ is the $\mathbb Z$-span of the images of $hw_\nu$ under the divided powers of the row operators $x_{b\cdot}\partial/\partial x_{a\cdot}$, and the \emph{dual Weyl module} $\nabla(\nu)_{\mathbb Z}$ is the set of $p\in\mathbb Z[x]$ of column weight $\nu$ that are invariant under the column substitutions $x_{\cdot j}\mapsto x_{\cdot j}+t\,x_{\cdot l}$ for $l<j$. Both are lattices in the same $\mathbb Q$-space, the irreducible representation of highest weight $\nu$, which is spanned over $\mathbb Q$ by $hw_\nu$ and its translates.

\begin{definition}\label{r4:R4:std}
For a partition $\nu$ with $s$ parts let $\mu_{\mathrm{std}}(\nu)=-\frac1{r_\nu}\log\det\big(\{e_i,e_{i'}\}/\{hw_\nu,hw_\nu\}\big)$, where $e_1,\dots,e_{r_\nu}$ is a $\mathbb Z$-basis of $V(\nu)_{\mathbb Z}$, and put $c_\nu=\{hw_\nu,hw_\nu\}/\prod_i\nu_i!$.
\end{definition}

\begin{lemma}\label{r4:R4:hook}
$\{hw_\nu,hw_\nu\}=\prod_{i=1}^s(\nu_i+s-i)!\big/\prod_{i<i'}(\nu_i-\nu_{i'}+i'-i)$, and $1\le c_\nu\le(\nu_1+s)^{s(s-1)/2}$.
\end{lemma}

\begin{proof}
The first formula is the Capelli identity applied to the highest weight vector, which computes $hw_\nu(\partial)hw_\nu$ (see the published input listed in Subsection~\ref{r4:sec:R4open}). Then $c_\nu=\prod_i\frac{(\nu_i+s-i)!}{\nu_i!}\big/\prod_{i<i'}(\nu_i-\nu_{i'}+i'-i)$. Since $\nu_{i'}\ge0$, $\prod_{i'>i}(\nu_i-\nu_{i'}+i'-i)\le\prod_{i'>i}(\nu_i+i'-i)=(\nu_i+s-i)!/\nu_i!$, which gives $c_\nu\ge1$. Since each factor of the denominator is at least $1$ and $(\nu_i+s-i)!/\nu_i!\le(\nu_i+s)^{s-i}\le(\nu_1+s)^{s-i}$, we get $c_\nu\le(\nu_1+s)^{\sum_i(s-i)}=(\nu_1+s)^{s(s-1)/2}$.
\end{proof}

\begin{proposition}\label{r4:R4:stdslope}
$\mu_{\mathrm{std}}(\nu)\ge-\big(\lvert\nu\rvert\log s+\log c_\nu\big)$, where $\lvert\nu\rvert=\sum_i\nu_i$.
\end{proposition}

\begin{proof}
Let $G$ and $G'$ be the Gram matrices of $\{\cdot,\cdot\}/\{hw_\nu,hw_\nu\}$ on $\mathbb Z$-bases of $V(\nu)_{\mathbb Z}$ and of $\nabla(\nu)_{\mathbb Z}$. Under the form $\{\cdot,\cdot\}/\{hw_\nu,hw_\nu\}$ the lattices $V(\nu)_{\mathbb Z}$ and $\nabla(\nu)_{\mathbb Z}$ are dual to each other, so $\det G\cdot\det G'=1$; this is the duality of Weyl modules and induced modules under the contravariant form normalized at the highest weight vector, listed in Subsection~\ref{r4:sec:R4open}. It suffices to prove $\det G'\ge A^{r_\nu}$ with $A=\prod_i\nu_i!/(s^{\lvert\nu\rvert}\{hw_\nu,hw_\nu\})$, because then $\mu_{\mathrm{std}}(\nu)=-\frac1{r_\nu}\log\det G=\frac1{r_\nu}\log\det G'\ge\log A=-(\lvert\nu\rvert\log s+\log c_\nu)$.

Let $f_1,\dots,f_{r_\nu}$ be a $\mathbb Z$-basis of $\nabla(\nu)_{\mathbb Z}$ and let $F$ be the integral matrix of their coefficients in the monomial basis. The unnormalized Gram matrix is $F\,\mathrm{diag}(\alpha!)\,F^{\mathsf T}$, and the Cauchy--Binet formula gives
\[
\det\big(F\,\mathrm{diag}(\alpha!)F^{\mathsf T}\big)=\sum_{\lvert S\rvert=r_\nu}\det(F_S)^2\prod_{\alpha\in S}\alpha!,
\]
where $S$ runs over the sets of $r_\nu$ monomials and $F_S$ is the corresponding square minor. Every monomial $x^\alpha$ that occurs has column weight $\nu$, that is $\sum_a\alpha_{aj}=\nu_j$ for each column $j$, and the multinomial bound $\nu_j!\le s^{\nu_j}\prod_a\alpha_{aj}!$ gives $\alpha!\ge\prod_j\nu_j!/s^{\lvert\nu\rvert}$. Since $F$ has rank $r_\nu$, some $\det F_S$ is a nonzero integer, and therefore $\det(F\mathrm{diag}(\alpha!)F^{\mathsf T})\ge(\prod_j\nu_j!/s^{\lvert\nu\rvert})^{r_\nu}$. Dividing by $\{hw_\nu,hw_\nu\}^{r_\nu}$ gives $\det G'\ge A^{r_\nu}$.
\end{proof}

We now describe $M_u$. For each class $c$ of the genus of level zero and each $\tau\in\Phi$ there is a hermitian lattice $Q_{c,\tau}$ of rank $k$ over the field $F$, obtained from the class representative $L_c$ by the embedding $\tau$, with $L_c\subset L_c^\#$ of index $T_u^{2f}$. The lattice $M_u$ is the direct sum over the classes $c$ of the tensor products over $\tau\in\Phi$ of the Schur functors of highest weight $\lambda^{\mathrm{col}}_{u,\tau}-\rho_k$ applied to $Q_{c,\tau}$, with the integral structure of the Weyl modules at the finite places and with the Fock metrics, normalized by the norm of the highest weight vector, in orthonormal frames at the archimedean places. We write $\lambda_\tau=(\lambda_{\tau,0}\ge\dots\ge\lambda_{\tau,k-1})$ for this highest weight, $\lvert\lambda_\tau\rvert=\sum_i\lambda_{\tau,i}$, and $\nu_\tau=\lambda_\tau-\lambda_{\tau,k-1}(1,\dots,1)$, which is a partition.

\begin{lemma}\label{r4:R4:schur}
Let $Q$ be a hermitian lattice of rank $s$ over $\mathcal O_F$, let $\lambda$ be a dominant weight of $GL_s$ with $\nu=\lambda-\lambda_{s-1}(1,\dots,1)$, and let $\mathbb S_\lambda Q$ be the lattice of the Schur functor of $\lambda$ with the integral structure given at each finite place by the Weyl module $V(\nu)_{\mathbb Z}$ twisted by $\det^{\lambda_{s-1}}$ and with the Fock metrics normalized at the highest weight vector in orthonormal frames. Then
\[
\mu(\mathbb S_\lambda Q)=\lvert\lambda\rvert\,\mu(Q)+\mu_{\mathrm{std}}(\nu).
\]
\end{lemma}

\begin{proof}
Let $r_\lambda$ be the dimension of the representation of highest weight $\lambda$. For $g\in GL_s$ the determinant of $\mathbb S_\lambda(g)$ is $\det(g)^{\lvert\lambda\rvert r_\lambda/s}$, because both sides are characters of $GL_s$ that agree on scalar matrices. Every hermitian lattice of rank $s$ is the translate $g\cdot\mathcal O_F^s$ of the standard lattice with the standard metric by an element $g=(g_w,g_v)$ of $GL_s$ of the adeles of $F$, with $g_w$ in $GL_s(\mathcal O_{F_w})$ for almost all finite places $w$. Translating by $g$ changes $\degh$ of a lattice of rank $s'$ by $-\sum_w\log\lvert\det g_w\rvert_w-\sum_v\log\lvert\det g_v\rvert_v^2$, with the normalized absolute values, when the determinant is taken on the lattice of rank $s'$. The integral structures of $\mathbb S_\lambda$ are $GL_s(\mathcal O_{F_w})$-stable and the normalized Fock metrics are invariant under the unitary group, so $\mathbb S_\lambda(g\cdot\mathcal O_F^s)=\mathbb S_\lambda(g)\cdot\mathbb S_\lambda(\mathcal O_F^s)$, and the determinant formula gives
\[
\degh(\mathbb S_\lambda Q)-\degh(\mathbb S_\lambda\mathcal O_F^s)=\frac{\lvert\lambda\rvert r_\lambda}{s}\big(\degh(Q)-\degh(\mathcal O_F^s)\big)=\frac{\lvert\lambda\rvert r_\lambda}{s}\degh(Q).
\]
The twist by $\det^{\lambda_{s-1}}$ does not change the standard lattice, which has degree zero, so $\degh(\mathbb S_\lambda\mathcal O_F^s)=\degh(V(\nu)_{\mathbb Z}\otimes\mathcal O_F)=n_Fr_\lambda\,\mu_{\mathrm{std}}(\nu)$. Dividing by $n_Fr_\lambda$ gives the formula, since $\mu(Q)=\degh(Q)/(n_Fs)$.
\end{proof}

\begin{proposition}\label{r4:R4:Mslope}
$\mu(M_u)=-\sum_{\tau\in\Phi}\frac{\lvert\lambda_\tau\rvert}k\log T_u+\sum_{\tau\in\Phi}\mu_{\mathrm{std}}(\nu_\tau)$, and therefore
\[
\mu(M_u)\ge-\sum_{\tau\in\Phi}\frac{\lvert\lambda_\tau\rvert}k\log T_u-\sum_{\tau\in\Phi}\Big(\lvert\nu_\tau\rvert\log k+\frac{k^2}2\log\big(\lambda_{\tau,0}-\lambda_{\tau,k-1}+k\big)\Big).
\]
\end{proposition}

\begin{proof}
The hermitian form $h$ of $L_c$ takes values in $\mathcal O_E$ on $L_c$, and $L_c^\#=\{y:h(L_c,y)\subset\mathcal O_E\}$ is isometric to the dual lattice $L_c^\vee$. Hence $\degh(L_c^\#)=-\degh(L_c)$ and $\degh(L_c^\#)-\degh(L_c)=\log[L_c^\#:L_c]=2f\log T_u$, so $\degh(L_c)=-f\log T_u$ and $\mu(L_c)=-f\log T_u/(fk)=-\log T_u/k$. The normalized slope does not change under the base change to $F$ by $\tau$, so $\mu(Q_{c,\tau})=-\log T_u/k$, and Lemma~\ref{r4:R4:schur} gives $\mu(\mathbb S_{\lambda_\tau}Q_{c,\tau})=-\frac{\lvert\lambda_\tau\rvert}k\log T_u+\mu_{\mathrm{std}}(\nu_\tau)$. The slope of a tensor product of hermitian lattices is the sum of the slopes, and a direct sum of lattices of equal slope has that slope, which gives the formula for $\mu(M_u)$. Proposition~\ref{r4:R4:stdslope} and Lemma~\ref{r4:R4:hook} with $s=k$ give the inequality.
\end{proof}

\begin{theorem}\label{r4:R4:A7}
Let $C_n\ge1$ and suppose that the weights of level zero satisfy, for every $u$ and $\tau$, $\lambda_{\tau,i}\in\mathbb Z$, and
\begin{enumerate}
\item[(W1)] if $\tau$ is a top, $xb\le\lambda_{\tau,0}\le2xb$, $b'\le-\lambda_{\tau,k-1}<2b'$, and $0\le\lambda_{\tau,i}\le S$ for $0<i<k-1$;
\item[(W2)] if $\tau$ is not a top, $0\le\lambda_{\tau,i}\le C_nS$ for all $i$,
\end{enumerate}
with $1\le S\le b'$. Suppose also that $\log T_u\le C_TmL_\varepsilon$, $\log k\le5\log x$, $fk(b'+S)\le N_Ta$ and $fk^2\le N_Ta$, and that $\log x\le L_\varepsilon$ and $\log\mathrm{rd}_E\le L_\varepsilon$. Then \eqref{A7} holds for every $u\in W$ with
\[
e_7=C_7\,N_Ta\Lambda_a\Big(1+\frac mk\Big),
\]
where $C_7$ depends on $C_n$ and $C_T$ alone.
\end{theorem}

\begin{proof}
By (W1) and (W2), $\lvert\lambda_\tau\rvert\le2a+kS$ at a top and $\lvert\lambda_\tau\rvert\le kC_nS$ elsewhere, where we drop the negative entry $\lambda_{\tau,k-1}$ at a top. Further $\lvert\nu_\tau\rvert=\lambda_{\tau,0}+\sum_{0<i<k-1}\lambda_{\tau,i}-(k-1)\lambda_{\tau,k-1}\le2a+kS+2kb'$ at a top and $\lvert\nu_\tau\rvert\le kC_nS$ elsewhere. Proposition~\ref{r4:R4:Mslope} gives
\begin{multline*}
-\mu(M_u)\le N_T\frac{2a}k\log T_u+f(1+C_n)S\log T_u+\big(2aN_T+fk(3+C_n)(S+b')\big)\log k\\
+\sum_\tau\frac{k^2}2\log\big(\lambda_{\tau,0}-\lambda_{\tau,k-1}+k\big).
\end{multline*}
We bound each term by a multiple of $N_Ta\Lambda_a(1+m/k)$. First $N_T\frac{2a}k\log T_u\le2C_TN_Ta\Lambda_a\frac mk$. Next $fS\log T_u\le fSC_TmL_\varepsilon\le C_TfkS\Lambda_a\le C_TN_Ta\Lambda_a$ because $m\le k$. Next $\log k\le5\log x\le5\Lambda_a$ and $fk(S+b')\le N_Ta$, so the third term is at most $5(5+C_n)N_Ta\Lambda_a$. Finally $\lambda_{\tau,0}-\lambda_{\tau,k-1}+k\le(C_n+5)a$ at every $\tau$, so the last term is at most $\frac12fk^2\big(\log(C_n+5)+\Lambda_a\big)\le\big(1+\log(C_n+5)\big)N_Ta\Lambda_a$, because $fk^2\le N_Ta$ and $\Lambda_a\ge1$. So $\mu(M_u)\ge-C_7N_Ta\Lambda_a(1+m/k)$ with $C_7=2C_T+C_T(1+C_n)+5(5+C_n)+1+\log(C_n+5)$.
\end{proof}

\begin{remark}\label{r4:R4:Mmax}
The same description of $M_u$ bounds $\mu_{\max}(M_u)$ above by $C\,N_Ta\Lambda_a$, through the duality of Proposition~\ref{r4:R4:stdslope} and the bound $c_\nu\le(\nu_1+s)^{s(s-1)/2}$. That bound enters \eqref{A2} and is not needed for \eqref{A7}.
\end{remark}

\subsection{The sizes of the errors}\label{r4:sec:R4sizes}

\begin{proposition}\label{r4:R4:sizes}
Let $\varepsilon\in(0,\frac1{52})$ and suppose that the size conditions \eqref{sz1} and \eqref{sz2} of Proposition~\ref{sizes} hold, that $A=\sum_ta_t\le2N_Ta$, and that $\log x\le L_\varepsilon$ and $\log\mathrm{rd}_E\le L_\varepsilon$. Then
\[
e_4+Ae_6+e_7\le C_8\,x^{\varepsilon}\,\mathcal B,
\]
with $C_8$ depending on $C_0$, $C_6$ and $C_7$ alone, where $\mathcal B=N_Ta(1+m^2/k)+m(B_0+N_c)$ is the number \eqref{calB}.
\end{proposition}

\begin{proof}
Since $\Lambda_a\le3x^\varepsilon$, $e_4\le(50C_0+10)\cdot3x^\varepsilon\big(N_Ta+m(B_0+N_c)\big)\le3(50C_0+10)x^\varepsilon\mathcal B$. Next $Ae_6\le2N_Ta\big(C_6x^\varepsilon(1+m^2/k)+2\big)\le2(C_6+2)x^\varepsilon N_Ta(1+m^2/k)$. Finally $e_7\le3C_7x^\varepsilon N_Ta(1+m/k)\le3C_7x^\varepsilon N_Ta(1+m^2/k)$.
\end{proof}

The constant $C_8$ depends on $\varepsilon$ through $C_6$, which contains the factor $1/(e\varepsilon)$ of the bound $\log x\le L_\varepsilon/(e\varepsilon)$. For $x\ge x_0(\varepsilon)$ this factor is absorbed in $x^{\varepsilon}$, so the three errors are at most $x^{2\varepsilon}\mathcal B$ with an absolute implied constant once $x$ is large in terms of $\varepsilon$.

\subsection{Inputs of this section}\label{r4:sec:R4open}

This section proves, from the hypotheses named in them, the mean period formula (Proposition~\ref{r4:R4:period}), the height bound (Proposition~\ref{r4:R4:ht} and Corollary~\ref{r4:R4:htcor}), the deduction of \eqref{A4} (Theorem~\ref{r4:R4:A4}), the deduction of \eqref{A6} (Proposition~\ref{r4:R4:weyl} and Theorem~\ref{r4:R4:A6}), the deduction of \eqref{A7} (Theorem~\ref{r4:R4:A7}), and the comparison with $\mathcal B$ (Proposition~\ref{r4:R4:sizes}). Subsections~\ref{w3:sec:HK} to~\ref{w3:sec:HT} reduce Hypotheses~\ref{r4:R4:HK}, \ref{r4:R4:HP}, \ref{r4:R4:HL} and~\ref{r4:R4:HT} to the statements stated there. \notgiven{Hypotheses~\ref{r4:R4:HK}, \ref{r4:R4:HP}, \ref{r4:R4:HL}, \ref{r4:R4:HT}, \ref{r4:R4:HW} and~\ref{r4:R4:HD} are not proved here.} The published theorems used are Kronecker's first limit formula, the product formula of Jacobi for $\Delta$, Dirichlet's class number formula, the main theorem of complex multiplication for the conjugates of $j(\mathfrak c)$, the discriminant formula for a tower of fields, Siegel's mass formula for unimodular hermitian lattices, Harish-Chandra's formula for formal degrees, Weyl's dimension formula, the Capelli identity, and the duality of Weyl modules and dual Weyl modules under the contravariant form.

\subsection{The archimedean hypotheses: conventions}\label{w3:sec:conv}

We use the conventions of Subsection~\ref{r4:sec:R4conv}. $F$ is a totally complex field with $n_F$ archimedean places, $\degh$ is the arithmetic degree, $\mu=\degh/(n_F\rk)$ is twice the slope of Bost, and $\mathrm{ht}(\phi)=n_F^{-1}\sum_v\log\lVert\phi\rVert_v^2$. For positive reals $\kappa_v$ we write $\langle\log\kappa\rangle=n_F^{-1}\sum_v\log\kappa_v$. We write $\ph(y)=y\log y-y$ for $y>0$ and $\ph(0)=0$. For a hermitian lattice $\Lambda$ of rank $D$ we put $\lcv(\Lambda)=-\degh(\Lambda)/n_F$, so that $\mu(\Lambda^\vee)=\lcv(\Lambda)/D$.

For a finite set $\kappa$ we write $\mathbb C^\kappa$ with the standard norm $\lvert u\rvert^2=\sum_x\lvert u_x\rvert^2$, Lebesgue measure $du$, and for $e\in\mathbb N^\kappa$ we put $u^e=\prod_xu_x^{e_x}$, $e!=\prod_xe_x!$ and $\lvert e\rvert=\sum_xe_x$.

\subsection{The archimedean norm of each lift}\label{w3:sec:HK}

\subsubsection{Three elementary estimates}

\begin{lemma}\label{w3:W3:shift}
Let $y\ge0$, $0<\epsilon<1$ and $s\ge0$. Then $s^y\exp(-(1-\epsilon)s^2/2)\le(1-\epsilon)^{-y/2}\exp(\ph(y)/2)$. For $n\ge0$ one has $(1-\frac1{n+1})^{-n/2}\le e^{1/2}$.
\end{lemma}

\begin{proof}
For $y=0$ both sides are at most $1$. For $y>0$ the function $g(s)=y\log s-(1-\epsilon)s^2/2$ has its maximum at $s^2=y/(1-\epsilon)$, where $g=\frac y2\log y-\frac y2-\frac y2\log(1-\epsilon)=\frac12\ph(y)-\frac y2\log(1-\epsilon)$. For the second statement, $(1-\frac1{n+1})^{-n/2}=(1+\frac1n)^{n/2}\le e^{1/2}$ for $n\ge1$, and the case $n=0$ is trivial.
\end{proof}

\begin{lemma}\label{w3:W3:gauss}
Let $\iota$ be a finite set of cardinality $q$, let $\nu_1,\dots,\nu_r\ge0$ with $\lvert\nu\rvert=\sum_i\nu_i$, and let $s_1,\dots,s_r$ be nonnegative measurable functions on $\mathbb C^\iota$ with $\sum_is_i(u)^2=\lvert u\rvert^2$. Then
\[
\pi^{-q}\int_{\mathbb C^\iota}\prod_is_i(u)^{\nu_i}\,e^{-\lvert u\rvert^2/2}\,du\le e^{1/2}\,(2\lvert\nu\rvert+2)^{q}\,\exp\Big(\frac12\sum_i\ph(\nu_i)\Big).
\]
\end{lemma}

\begin{proof}
Put $\epsilon=1/(\lvert\nu\rvert+1)$. Write $e^{-\lvert u\rvert^2/2}=\prod_ie^{-(1-\epsilon)s_i^2/2}\cdot e^{-\epsilon\lvert u\rvert^2/2}$. Lemma~\ref{w3:W3:shift} applied to each factor $s_i^{\nu_i}e^{-(1-\epsilon)s_i^2/2}$ bounds the integrand by $(1-\epsilon)^{-\lvert\nu\rvert/2}e^{\frac12\sum\ph(\nu_i)}e^{-\epsilon\lvert u\rvert^2/2}$. The Gaussian integral is $\pi^{-q}\int e^{-\epsilon\lvert u\rvert^2/2}du=(2/\epsilon)^q=(2\lvert\nu\rvert+2)^q$, and $(1-\epsilon)^{-\lvert\nu\rvert/2}\le e^{1/2}$ by the second part of Lemma~\ref{w3:W3:shift}.
\end{proof}

\begin{lemma}[coefficients in the Fock space]\label{w3:W3:jet}
Let $G$ be an entire function on $\mathbb C^\kappa$ with Taylor expansion $G(u)=\sum_ec_eu^e$ and $\lvert G(u)\rvert\le Se^{\lvert u\rvert^2/2}$. Then for every $e\in\mathbb N^\kappa$ the function $G(u)\bar u^{\,e}e^{-\lvert u\rvert^2}$ is integrable and
\[
\int_{\mathbb C^\kappa}G(u)\,\bar u^{\,e}\,e^{-\lvert u\rvert^2}\,du=\pi^{\lvert\kappa\rvert}\,e!\,c_e .
\]
\end{lemma}

\begin{proof}
Integrability follows from $\lvert G(u)\bar u^{e}\rvert e^{-\lvert u\rvert^2}\le S\lvert u\rvert^{\lvert e\rvert}e^{-\lvert u\rvert^2/2}$. On the polydisc $P_R=\{\lvert u_x\rvert\le R\}$ the Taylor series converges uniformly, so $\int_{P_R}G\bar u^ee^{-\lvert u\rvert^2}=\sum_{e'}c_{e'}\int_{P_R}u^{e'}\bar u^{e}e^{-\lvert u\rvert^2}$. In polar coordinates in each variable the integral vanishes unless $e'=e$, and for $e'=e$ it equals $\prod_x2\pi\int_0^Rr^{2e_x+1}e^{-r^2}dr$. Letting $R\to\infty$ and using dominated convergence on the left gives $\pi^{\lvert\kappa\rvert}e!c_e$.
\end{proof}

\subsubsection{The block bound}

Let $\Phi$ be a finite set of places (the embeddings $\tau$), let $\iota$ be a finite set with $q=\lvert\iota\rvert$, and put $\kappa=\Phi\times\iota$. Fix $n_\tau\in\mathbb N$ and weights $\nu_\tau=(\nu_{\tau1},\dots,\nu_{\tau r})$ with $\sum_i\nu_{\tau i}=n_\tau$.

\begin{proposition}[block bound]\label{w3:W3:block}
Let $W$ be a Banach space and $R(u)=\sum_eu^eR_e$ a $W$-valued polynomial on $\mathbb C^\kappa$ whose monomials satisfy $\sum_{x\in\iota}e(\tau,x)=n_\tau$ for every $\tau$. Suppose there are nonnegative functions $s_{\tau i}$ on $\mathbb C^\iota$ with $\sum_is_{\tau i}(w)^2=\lvert w\rvert^2$ and
\[
\lVert R(\bar u)\rVert\le\prod_\tau\prod_is_{\tau i}(u_\tau)^{\nu_{\tau i}}\qquad(u=(u_\tau)_\tau\in\mathbb C^\kappa).
\]
Let $G=\sum c_eu^e$ be entire on $\mathbb C^\kappa$ with $\lvert G(u)\rvert\le Se^{\lvert u\rvert^2/2}$, and let $\lambda_\tau\in\mathbb C$ with $\lvert\lambda_\tau\rvert^2=\Omega_\tau^{-1}$, $\Omega_\tau>0$. Then the vector $J=\sum_ec_e\,e!\,\prod_{(\tau,x)}\lambda_\tau^{e(\tau,x)}R_e$ satisfies
\[
\lVert J\rVert^2\le S^2\prod_\tau\Omega_\tau^{-n_\tau}\,e\,(2n_\tau+2)^{2q}\exp\Big(\sum_i\ph(\nu_{\tau i})\Big).
\]
\end{proposition}

\begin{proof}
By homogeneity, $\prod_{(\tau,x)}\lambda_\tau^{e(\tau,x)}=\prod_\tau\lambda_\tau^{n_\tau}$ for every $e$ in the support of $R$. Lemma~\ref{w3:W3:jet} gives $c_ee!=\pi^{-\lvert\kappa\rvert}\int G(u)\bar u^ee^{-\lvert u\rvert^2}du$, hence $J=\prod_\tau\lambda_\tau^{n_\tau}\cdot\pi^{-\lvert\kappa\rvert}\int G(u)R(\bar u)e^{-\lvert u\rvert^2}du$, the sum being finite. Therefore
\[
\lVert J\rVert\le\prod_\tau\Omega_\tau^{-n_\tau/2}\,S\,\prod_\tau\Big(\pi^{-q}\int_{\mathbb C^\iota}\prod_is_{\tau i}(u_\tau)^{\nu_{\tau i}}e^{-\lvert u_\tau\rvert^2/2}du_\tau\Big),
\]
and Lemma~\ref{w3:W3:gauss} bounds each factor. Squaring gives the statement.
\end{proof}

In the application $R$ is the $K$-type projection of the matrix coefficient, $R_e$ are the operators that send a highest weight vector of row degree $e$ to its image, and $s_{\tau i}(w)$ are the singular values of the matrix $w\in M_{a\times b}(\mathbb C)$. The hypothesis $\lVert R(\bar u)\rVert\le\prod s^{\nu}$ is the statement that the polynomial representation of $U(a)\times U(b)$ on monomials of degree $\nu_\tau$ in the singular values is contracting, which holds for the standard Fock-space realization because the operator norm of a matrix power representation is the product of the singular values to the power of the weight (Cartan decomposition $w=k_1\,\mathrm{diag}(s)\,k_2$ with unitary $k_1,k_2$, and the highest weight vector of weight $\nu$ is multiplied by $\prod s_i^{\nu_i}$). Proposition~\ref{w3:W3:block} applies with $s$ the singular values and $n_\tau=\lvert\nu_\tau\rvert$.

\subsubsection{Class sums}

\begin{lemma}\label{w3:W3:schur}
Let $\mathcal C$ and $\mathcal C'$ be finite sets partitioned into blocks $\mathcal C=\bigsqcup_{a\in A}\mathcal C_a$ and $\mathcal C'=\bigsqcup_{b\in A'}\mathcal C'_b$ with $\lvert\mathcal C_a\rvert\le h$ and $\lvert\mathcal C'_b\rvert\le h'$. Let $0<\gamma_L\le1$ and $0<\gamma'_{L'}\le1$, let $\epsilon_{LL'}$ be complex numbers of modulus one, and let $T_{LL'}\colon W\to W'$ be bounded operators with $\lVert T_{LL'}\rVert^2\le t$ which vanish unless $L$ and $L'$ lie in blocks $a$, $b$ with $(a,b)$ in a set $E\subset A\times A'$. Define $\Theta(G)_{L'}=\sum_L\gamma_L\epsilon_{LL'}T_{LL'}G_L$. Then
\[
\sum_{L'}\gamma'_{L'}\lVert\Theta(G)_{L'}\rVert^2\le\lvert A\rvert\,\lvert A'\rvert\,h\,h'\,t\sum_L\gamma_L\lVert G_L\rVert^2 .
\]
\end{lemma}

\begin{proof}
By Cauchy--Schwarz, $\lVert\Theta(G)_{L'}\rVert^2\le\big(\sum_{L}\gamma_L\big)\big(\sum_L\gamma_L\lVert T_{LL'}G_L\rVert^2\big)\le\lvert A\rvert h\cdot t\sum_L\gamma_L\lVert G_L\rVert^2$, using $\gamma_L\le1$. Summing over $L'$ with $\gamma'_{L'}\le1$ multiplies by $\lvert\mathcal C'\rvert\le\lvert A'\rvert h'$.
\end{proof}

For the coupled last two steps the same argument applied twice, with an intermediate block set $Y$ of cardinality bounded by the coset count, gives the factor $\lvert Y\rvert^2$ times the two single-step factors. We do not repeat it.

\subsubsection{The theta kernel}

The local step kernel of the lift $j$ at an archimedean place is a theta function on the abelian variety $E^{n}$, $n=fN_{j-1}N_j$, where $E$ is the CM elliptic curve of the host and $f=\lvert\Phi\rvert$. We need its sup norm.

\begin{lemma}[sup norm of a theta function]\label{w3:W3:supL2}
Let $\Lambda\subset\mathbb C^n$ be a lattice of covolume $V$ for the standard volume, and let $F$ be holomorphic on $\mathbb C^n$ with $\lvert F(z+\ell)\rvert^2e^{-\lvert z+\ell\rvert^2}=\lvert F(z)\rvert^2e^{-\lvert z\rvert^2}$ for $\ell\in\Lambda$. Put $\lVert F\rVert^2=V^{-1}\int_{\mathbb C^n/\Lambda}\lvert F\rvert^2e^{-\lvert z\rvert^2}dz$. Let $\rho>0$ and let $N_\Lambda(\rho)$ be the largest number of translates $\ell+P$ of a fixed fundamental parallelotope $P$ that meet a ball of radius $\rho$. Then
\[
\sup_z\lvert F(z)\rvert^2e^{-\lvert z\rvert^2}\le\frac{n!\,e^{\rho^2}}{\rho^{2n}\pi^n}\,N_\Lambda(\rho)\,V\,\lVert F\rVert^2 .
\]
\end{lemma}

\begin{proof}
Fix $z$ and put $g(w)=F(z+w)e^{-\langle w,z\rangle-\lvert z\rvert^2/2}$, which is holomorphic in $w$, with $\lvert g(w)\rvert^2=\lvert F(z+w)\rvert^2e^{-\lvert z+w\rvert^2}e^{\lvert w\rvert^2}$. The sub-mean value property of $\lvert g\rvert^2$ on the ball $B_\rho$ of volume $\pi^n\rho^{2n}/n!$ gives
\[
\lvert F(z)\rvert^2e^{-\lvert z\rvert^2}=\lvert g(0)\rvert^2\le\frac{n!}{\pi^n\rho^{2n}}\int_{B_\rho}\lvert g\rvert^2\le\frac{n!\,e^{\rho^2}}{\pi^n\rho^{2n}}\int_{B_\rho(z)}\lvert F\rvert^2e^{-\lvert\zeta\rvert^2}d\zeta .
\]
The ball $B_\rho(z)$ is covered by at most $N_\Lambda(\rho)$ translates of $P$, and on each translate the integral of the periodic function $\lvert F\rvert^2e^{-\lvert\zeta\rvert^2}$ equals $V\lVert F\rVert^2$.
\end{proof}

For $\Lambda=\mathfrak a^n$ with $\mathfrak a\subset\mathbb C$ a lattice of area $A$ whose reduced basis has lengths $\ell_1\le\ell_2$, a fundamental parallelotope is the product of $n$ reduced parallelograms of diameter at most $\ell_1+\ell_2$, and with $\rho=1$ one has $N_\Lambda(1)\le\big((1+2(1+\ell_1+\ell_2)/\ell_1)(1+2(1+\ell_1+\ell_2)\ell_1/A)\big)^n$. For the CM lattice of the host, after the normalization of Section~8 of the manuscript, $\ell_1\ge c$ and $\ell_2\le C\,v$ with $v=\frac12\mathrm{rd}_E$, so that the logarithm of the constant of Lemma~\ref{w3:W3:supL2} is at most $n\log(C'nv)$. Below we use only the bound $c_5\,n\,\Lambda_a$ for this logarithm.

\begin{hypothesis}[normalized kernel]\label{w3:W3:K3}
For each lift $j$ and each embedding $s$ of the coefficient field $M$ of the kernels, the local kernel of $K_j$ at $s$ is of the form $\Theta_j=\sum_{L,L'}\gamma_L\epsilon_{LL'}J_{LL'}$ as in Lemma~\ref{w3:W3:schur}, where each $J_{LL'}$ is the vector $J$ of Proposition~\ref{w3:W3:block} built from a theta function $G_{LL'}$ which is a combination $\sum_kr_kG_0(z+y_k)$ of translates of one normalized theta function $G_0$ on $E^{fN_{j-1}N_j}$, with $\sum_k\lvert r_k^{s}\rvert\le\exp(\mathrm{MN}_{j,s})$, and with $\lVert G_0\rVert^2=\lVert\theta_0^{s}\rVert^{2fN_{j-1}N_j}$, where $\theta_0$ is the genus-one theta function of the base vector and $\lVert\cdot\rVert$ is the norm of Lemma~\ref{w3:W3:supL2}. Moreover $\mathrm{MN}_{j,s}\le C_2fm^2N^2(\log(g_0g_1)+\log T_u)$ and $\lvert c_{D,\tau}r_\tau\rvert^2=\Omega_{\tau,s}^{-1}$ with $\Omega_{\tau,s}$ the conjugate periods of Proposition~\ref{r4:R4:period}.
\end{hypothesis}

The tensor power identity in Hypothesis~\ref{w3:W3:K3} is the statement that a theta function for the product polarization on $E^n$ whose space is one-dimensional is the product of the genus-one theta functions in each coordinate, and its norm is the $n$-th power. The coefficient bound is a count of the translates $y_k$ needed to refine the coset indicators of the level structure, of size at most $(g_0g_1T_u)^{C_2fm^2N^2}$. Both depend on the normalization of $\theta_0$ and of the constant $c$ of Theorem~\ref{kernellaw}, which is Hypothesis~\ref{r1:hyp:H3}. \notgiven{Hypothesis~\ref{w3:W3:K3} is not proved here.}

\begin{hypothesis}[theta-null mean]\label{w3:W3:K4}
There is an absolute constant $C_{\theta}$ such that for every admissible host and every number field $M$ containing the data,
\[
\frac1{[M:\mathbb Q]}\sum_{s\colon M\to\mathbb C}\log\lVert\theta_0^{s}\rVert^2\le C_\theta\big(\log v+\lvert\lambda_1\rvert\big),\qquad\lambda_1=L'/L(1,\chi_{-d}).
\]
\end{hypothesis}

A sharper bound,
\[
  \frac1{[M:\mathbb Q]}\sum_{\rho:M\to\mathbb C}
    \log\lVert\theta_0^\rho\rVert^2
  \leq \frac12\bigl(\log\sqrt d-\gamma+\lambda_1\bigr),
\]
follows from a mean-logarithm argument, with no pointwise estimate on Weber
values.  Put $n=[M:\mathbb Q]$.  For each $\rho$ choose compatible period
coordinates, and write $y_\rho=\operatorname{Im}\tau_\rho$,
$\Theta_\rho$ for the norm at zero of the even theta function whose zero is the marked half-period, $E_\rho=|\eta(\tau_\rho)|$,
$R_\rho=\lVert\theta_0^\rho\rVert^2$, and $A_\rho$ for the area of the
corresponding lattice.  The theta-null and period-product identities give
\[
  \log R_\rho=-\tfrac12\log2-\tfrac12\log y_\rho-2\log\Theta_\rho,
  \qquad
  \log A_\rho=\log4+2\log\pi+\log y_\rho+4\log E_\rho .
\]
The Weber construction gives an element $X\in M$, the same element for
all embeddings, such that $X\ne0$, $X$ is integral over $\mathbb Z$, and
\[
  |\rho X|=(\Theta_\rho/E_\rho)^{12}
\]
for every $\rho$ and every compatible choice of the vanishing even
characteristic.  Indeed, Jacobi's product formula and the half-period
identities identify these twelfth powers, up to sign, with the roots of
$(T-16)^3=jT$; integrality of $j$ gives integrality.  The algebraic expression
in the torsion coordinates makes the formula compatible with $\rho$
(Weber, \emph{Lehrbuch der Algebra} III, \S\S125--127; Yui and Zagier,
\emph{Math.\ Comp.}\ 66 (1997), \S1).  Since $X$ is a nonzero algebraic
integer, the norm of $X$ is a nonzero rational integer, and consequently
\[
  0\leq\sum_\rho\log|\rho X|=12(S_\Theta-S_E),\qquad
  S_\Theta=\sum_\rho\log\Theta_\rho,\quad S_E=\sum_\rho\log E_\rho .
\]
Set $S_y=\sum_\rho\log y_\rho$ and $S_R=\sum_\rho\log R_\rho$.  Averaging the
two displayed identities and using the Chowla--Selberg formula as in
Proposition~\ref{r4:R4:period} give
\[
  S_R=-\tfrac n2\log2-\tfrac12S_y-2S_\Theta,
  \qquad
  S_y+4S_E=n(\gamma-\log2-\log\sqrt d-\lambda_1).
\]
Since $S_\Theta\geq S_E$, substitution yields
\[
  \frac{S_R}{n}
  \leq-\tfrac12\log2-\tfrac1{2n}S_y-\tfrac2nS_E
  =\tfrac12(\log\sqrt d-\gamma+\lambda_1).
\]
\notgiven{The transformation/integrality inputs and the compatible period
normalizations used above are not reproved in this subsection.}

\begin{hypothesis}[class and index counts]\label{w3:W3:K5}
With $U=fm^2N^2\log(g_0g_1T_u)$ there is an absolute $C_c$ with: the number of classes in each block of the class decomposition at level $i$ is at most $e^{C_cU}$, the number of blocks is at most $e^{C_cU}$, and the coset count of the coupled step is at most $e^{C_cU}$.
\end{hypothesis}

Classes are orbits of an arithmetic group on decorated lattices of a fixed genus with level structure at $g_0g_1T_u$. Their number is at most the class number of the genus times the index of the principal congruence subgroup, and both are bounded by $N(\mathfrak n)^{C\,\mathrm{rank}^2}$ with $N(\mathfrak n)\le(g_0g_1T_u)^{Cf}$. \notgiven{The count is not written out here, so Hypothesis~\ref{w3:W3:K5} is not proved here.}

\subsubsection{The kernel norm theorem}

Put $n_j=fN_{j-1}N_j$, $\nu_{j\tau}$ the row weights of the kernel of the lift $j$ at $\tau$, and for an embedding $s$
\[
\begin{aligned}
\mathcal S_j(s)&=\sum_{\tau}\Big(\sum_i\ph(\nu_{j\tau i})-\lvert\nu_{j\tau}\rvert\log\Omega_{\tau,s}\Big)+n_j\log\lVert\theta_0^s\rVert^2+n_j\log c_\Lambda\\ &\quad{}+2\,\mathrm{MN}_{j,s}+\sum_\tau\big(1+2N_{j-1}N_j\log(2\lvert\nu_{j\tau}\rvert+2)\big),
\end{aligned}
\]
where $c_\Lambda$ is the constant of Lemma~\ref{w3:W3:supL2} per coordinate.

\begin{proposition}[one embedding]\label{w3:W3:onesigma}
Under Hypotheses~\ref{w3:W3:K3} and~\ref{w3:W3:K5}, for every lift $j$ and every embedding $s$,
\[
\log\lVert K_j\rVert_s^2\le\mathcal S_j(s)+4C_cU .
\]
\end{proposition}

\begin{proof}
By Lemma~\ref{w3:W3:supL2} and the triangle inequality, $\lvert G_{LL'}(z)\rvert\le S\,e^{\lvert z\rvert^2/2}$ with $S^2=c_\Lambda^{n_j}\lVert\theta_0^s\rVert^{2n_j}e^{2\mathrm{MN}_{j,s}}$, after the change of variable $z\mapsto r\,z$ that makes the kernel exponent $\lvert z\rvert^2/2$. Proposition~\ref{w3:W3:block} with $\Omega_\tau=\Omega_{\tau,s}$ and $q=N_{j-1}N_j$ gives $\lVert J_{LL'}\rVert^2\le\exp(\mathcal S_j(s))$, the term $\sum_\tau1$ coming from the factors $e$. Lemma~\ref{w3:W3:schur} multiplies by $\lvert A\rvert\lvert A'\rvert hh'\le e^{4C_cU}$ by Hypothesis~\ref{w3:W3:K5}. The frames of the chain are the class decompositions (the isometry $\sum_L\gamma_L\lVert\Phi(x)_L\rVert^2=\lVert x\rVert_{X}^2$ of the chain record), so the operator norm of $K_j$ in the frames equals the norm of $\Theta_j$ on the weighted sums.
\end{proof}

\begin{theorem}[Hypothesis~\ref*{r4:R4:HK}]\label{w3:W3:HK}
Assume Hypotheses~\ref{w3:W3:K3}, \ref{w3:W3:K4}, \ref{w3:W3:K5}, Proposition~\ref{r4:R4:period}, and the size conditions $\log T_u\le C_TmL_\varepsilon$, $\log(g_0g_1)\le C_8L_\varepsilon$, $\log(8+32n_jv)\le c_5\Lambda_a$, $\log(2\lvert\nu_{j\tau}\rvert+2)\le c_6\Lambda_a$, $N_{j-1}N_j\le N^2$, $L_\varepsilon\le\Lambda_a$, $1\le\Lambda_a$, $m,N\ge1$. Then Hypothesis~\ref{r4:R4:HK} holds with
\[
C_0=C_\theta'+c_5+2c_6+1+2C_2(C_T+C_8)+2C_c(1+C_T+C_8)+2C_8,
\]
where $C_\theta'$ is the constant with $C_\theta(\log v+\lvert\lambda_1\rvert)\le C'_\theta L_\varepsilon$ under $\log\mathrm{rd}_E\le L_\varepsilon$ and $\lvert\lambda_1\rvert\le L_\varepsilon$.
\end{theorem}

\begin{proof}
Average Proposition~\ref{w3:W3:onesigma} over the embeddings $s$, which is the mean over archimedean places of $F$ since each place of $F$ has the same number of embeddings of $M$ above it. Proposition~\ref{r4:R4:period} gives $\mathrm{mean}_s\log\Omega_{\tau,s}=\log\Omega$ for every $\tau$, so the first term becomes $\sum_\rho(\sum_i\ph(\nu_{j\rho i})-\lvert\nu_{j\rho}\rvert\log\Omega)$, which is the main term of $s_j$. The remaining terms are bounded as follows. Hypothesis~\ref{w3:W3:K4} gives $n_j\,\mathrm{mean}\log\lVert\theta_0\rVert^2\le C'_\theta fN^2L_\varepsilon$. Next $n_j\log c_\Lambda\le c_5fN^2\Lambda_a$, $2\,\mathrm{MN}\le2C_2fm^2N^2(C_8+C_T)L_\varepsilon$ because $\log T_u\le C_TmL_\varepsilon$ and $\log(g_0g_1)\le C_8L_\varepsilon$, $\sum_\tau(1+2N_{j-1}N_j\log(2\lvert\nu\rvert+2))\le f(1+2c_6N^2\Lambda_a)$, and $4C_cU\le4C_cfm^2N^2(C_8+C_T)mL_\varepsilon$. Each error term is at most the corresponding multiple of $fN^2\Lambda_a$ or of $fm^3N^2L_\varepsilon$, and adding the coefficients gives $C_0$.
\end{proof}

\begin{remark}
Theorem~\ref{w3:W3:HK} has the shape $C_0(fN^2\Lambda_a+fm^3N^2L_\varepsilon)$ of Hypothesis~\ref{r4:R4:HK}. It does not use the archimedean decay of the kernel beyond the Gaussian bound. The factor $m$ in front of $\ph(a_t)$ in \eqref{A4} arises because each of the $m$ lifts contributes $\sum_\rho\sum_i\ph(\nu_{j\rho i})$, and Hypothesis~\ref{r4:R4:HW} then extracts $\ph(a_\rho)$ at the tops for each $j$.
\end{remark}

\subsection{The final pairing}\label{w3:sec:HP}

\subsubsection{The branching map}

Let $N=n+1$, let $\lambda$ be a dominant weight of $GL_N$ and $\sigma$ a dominant weight of $GL_n$ which interlace, $\lambda_{i+1}\le\sigma_i\le\lambda_i$, with the same determinant shift. Realize $V_\lambda$ in polynomials on $M_N(\mathbb C)$ with the Fock form $\langle f,g\rangle=\sum_ef_e\bar g_e\,e!$, which is $U(N)$-invariant. Put $n_{\mathrm{Br}}=\lvert\lambda\rvert-\lvert\sigma\rvert$. By the Gelfand--Tsetlin branching rule, $V_\lambda$ restricted to $GL_n\times GL_1$ contains $V_\sigma\otimes\chi^{n_{\mathrm{Br}}}$ exactly once, where $\chi$ is the identity character of $GL_1$. Let $\mathrm{br}\colon V_\lambda\to V_\sigma$ be the $GL_n$-equivariant map normalized by $\mathrm{br}(\xi_{\mathrm{Br}})=c_{\mathrm{Br}}\kappa_{\mathrm{Br}}\,\mathrm{hw}_\sigma$, where $\xi_{\mathrm{Br}}$ is the product of leading minors of degree $\lambda$ that defines the branching map (\notgiven{its explicit form is not written out here}), $c_{\mathrm{Br}}$ is the integer of Lemma~\ref{w3:W3:cbr}, and $\kappa_{\mathrm{Br}}=\lVert\xi_{\mathrm{Br}}\rVert^2/\lVert\mathrm{hw}_\lambda\rVert^2$.

Let $G'$ be a positive definite hermitian $n\times n$ matrix, $g>0$, and give $V_\lambda$ the invariant form for the Gram matrix $\mathrm{diag}(G',g)$ and $V_\sigma$ the one for $G'$. Let $t\in\mathbb C$ with $\bar tgt=T>0$ and put $\mathrm{br}_t=t^{-n_{\mathrm{Br}}}\mathrm{br}$.

\begin{proposition}[scaled coisometry]\label{w3:W3:br}
For all $y\in V_\sigma$ and $x\in V_\lambda$ one has $\lVert\mathrm{br}_t\,x\rVert^2\le T^{-n_{\mathrm{Br}}}c_{\mathrm{Br}}^2\kappa_{\mathrm{Br}}\lVert x\rVert^2$, with equality on the $\sigma$-isotypic component. In particular the operator norm of $\mathrm{br}_t$ is $T^{-n_{\mathrm{Br}}/2}\lvert c_{\mathrm{Br}}\rvert\kappa_{\mathrm{Br}}^{1/2}$.
\end{proposition}

\begin{proof}
First let $G'=1$ and $g=1$. The $\sigma$-isotypic component $V^\sigma\subset V_\lambda$ is irreducible, and the other components are orthogonal to it by invariance of the form under the compact group $U(n)\times U(1)$. So $\mathrm{br}$ vanishes on $(V^\sigma)^\perp$, and on $V^\sigma$ it is a $U(n)$-equivariant map between irreducible unitary representations, hence a scalar multiple $\beta$ of an isometry by Schur's lemma. The inclusion $\mathrm{tr}\colon V_\sigma\to V_\lambda$ normalized by $\mathrm{tr}(\mathrm{hw}_\sigma)=c_{\mathrm{Br}}\xi_{\mathrm{Br}}$ is adjoint to $\mathrm{br}$ up to the same scalar. Evaluating at $\xi_{\mathrm{Br}}$, which lies in $V^\sigma$ and is a highest weight vector for $U(n)$ there, gives $\lvert\beta\rvert^2\lVert\xi_{\mathrm{Br}}\rVert^2=c_{\mathrm{Br}}^2\kappa_{\mathrm{Br}}^2\lVert\mathrm{hw}_\sigma\rVert^2$. With the extremal normalization $\lVert\mathrm{hw}_\sigma\rVert=\lVert\mathrm{hw}_\lambda\rVert$ of the Fock model this is $\lvert\beta\rvert^2=c_{\mathrm{Br}}^2\kappa_{\mathrm{Br}}$.

For general $G'$ and $g$, write $\mathrm{diag}(G',g)=h^*h$ with $h=\mathrm{diag}(h',g^{1/2})$. The invariant form for $\mathrm{diag}(G',g)$ is the pull back of the standard form by the action of $h$, and $\mathrm{br}$ is $GL_n$-equivariant, so the computation reduces to $G'=1$ except for the $GL_1$-factor. On $V^\sigma$ the element $\mathrm{diag}(1,u)$ acts by $u^{n_{\mathrm{Br}}}$. Hence the form for $(G',g)$ on $V^\sigma$ is $g^{n_{\mathrm{Br}}}$ times the form for $(G',1)$, and replacing $\mathrm{br}$ by $t^{-n_{\mathrm{Br}}}\mathrm{br}$ multiplies squared norms by $\lvert t\rvert^{-2n_{\mathrm{Br}}}$. Since $\lvert t\rvert^2g=T$, the squared norm ratio on $V^\sigma$ is $T^{-n_{\mathrm{Br}}}c_{\mathrm{Br}}^2\kappa_{\mathrm{Br}}$.
\end{proof}

\begin{lemma}\label{w3:W3:cbr}
Let $\mathrm{sp}(\lambda)=\lambda_1-\lambda_N+N$ and $d_1=\sum_{i<N}\min(\lambda_i-\lambda_{i+1},d_0)$ with $d_0=\sum_i(\sigma_i-\lambda_{i+1})$. If $\mathrm{sp}(\lambda)\le A$ then $0<\lvert c_{\mathrm{Br}}\rvert\le A^{n\,d_1}$.
\end{lemma}

\begin{proof}
By definition $c_{\mathrm{Br}}=\prod_{i<N}\prod_{v\in r_i\setminus\{\sigma_i\}}\prod_{j<N}\big((v-i)-(\sigma_j-j)\big)$ with $r_i=[\lambda_{i+1},\min(\lambda_i,\lambda_{i+1}+d_0)]$. The number of $v$ is $\lvert r_i\rvert-1\le\min(\lambda_i-\lambda_{i+1},d_0)$, so there are at most $n\,d_1$ factors. Each factor is an integer of absolute value at most $\lambda_1-\lambda_N+n-1\le A$. No factor vanishes, because a factor vanishes only if $v-i=\sigma_j-j$, which interlacing excludes for $v\ne\sigma_i$.
\end{proof}

\begin{hypothesis}[Fock norm of $\xi_{\mathrm{Br}}$]\label{w3:W3:xi}
$\lVert\xi_{\mathrm{Br}}\rVert^2\le\prod_{i<j\le N}(\lambda_i-\lambda_j+j-i)\,\lVert\mathrm{hw}_\lambda\rVert^2$. Consequently $\kappa_{\mathrm{Br}}\le\mathrm{sp}(\lambda)^{N(N-1)/2}$.
\end{hypothesis}

The second sentence follows from the first because each factor of the Weyl product is at most $\mathrm{sp}(\lambda)$. The first is a comparison of Fock norms of two products of minors with the same total degree. The Fock norm of a product of leading minors is a ratio of factorials of shifted weights by the Capelli identity. The required comparison includes the normalization of both highest weight vectors. \notgiven{Hypothesis~\ref{w3:W3:xi} is not proved here.}

\subsubsection{Hypothesis~\ref{r4:R4:HP}}

\begin{theorem}[Hypothesis~\ref*{r4:R4:HP}]\label{w3:W3:HP}
Assume that the final pairing $P_u$ is, at each archimedean place, the composite of $\mathrm{br}_t$ (with $\bar tgt=T_u$) on each embedding $\tau$ with an evaluation of operator norm squared at most $e^{\mathrm{ev}}$, $\mathrm{ev}\le C_{\mathrm{ev}}fN^2\Lambda_a$, and assume Hypothesis~\ref{w3:W3:xi} and the weight relations
\[
n_\tau=\lambda_{1,\tau}-D_{0,\tau},\quad\lambda_{1,t}=a_t-\tfrac{N-k}2\ (t\in T),\quad\lambda_{1,\tau}\ge S-\tfrac{N-1}2\ (\tau\notin T),
\]
\[
D_{0,\tau}\le S+m\eta+\mathbf 1_T(\tau)b_{T,\tau},\quad D_{1,\tau}\le3D_{0,\tau}+S+\eta,\quad b_{T,t}\le2b',\quad\mathrm{sp}\le2a,\quad\log T_u\le C_TmL_\varepsilon,
\]
together with $mL_\varepsilon\le N\Lambda_a$ and $\log2a\le2\Lambda_a$. Then $h_P\le-\sum_{t\in T}a_t\log T_u+e_P$ with $e_P\le C_0(N_Tb'N+f(m\eta+S)N+fN^2)\Lambda_a$ for $C_0=2(16+C_T+C_{\mathrm{ev}})$.
\end{theorem}

\begin{proof}
Proposition~\ref{w3:W3:br}, Lemma~\ref{w3:W3:cbr} and Hypothesis~\ref{w3:W3:xi} give
\[
h_P\le-\sum_\tau n_\tau\log T_u+\sum_\tau\big(2(N-1)D_{1,\tau}+\tfrac{N(N-1)}2\big)\log2a+\mathrm{ev}.
\]
At a top, $-n_t\log T_u=-a_t\log T_u+(\frac{N-k}2+D_{0,t})\log T_u$. Away from the tops $-n_\tau\log T_u\le(\frac{N-1}2+D_{0,\tau})\log T_u$ since $S\ge0$. Summing, the coefficient of $\log T_u$ besides $-\sum_ta_t$ is at most $f\frac N2+f(S+m\eta)+2N_Tb'$, and $\log T_u\le C_TmL_\varepsilon\le C_TN\Lambda_a$. The $c$- and $\kappa$-terms are at most $(2N(4S+4m\eta+6\mathbf 1_Tb')f+fN^2)\cdot2\Lambda_a$ after inserting $D_1\le3D_0+S+\eta$. Collecting coefficients gives the bound with the stated $C_0$, with room to spare.
\end{proof}

The assumption that $P_u$ factors through $\mathrm{br}_t$ is part of the construction of the final pairing, which is not written in this paper.

\subsection{The frame constants}\label{w3:sec:HL}

The frame record uses two spaces at the base tuple. The inclusion
$\iota$ maps the row space to $R_F$, and the $j$th lift $\theta_j$
maps its reduced source to $W_j$. The inclusion has scalar
$\kappa_v\prod_j\lambda_j(p,v)$ on the row part of type $p$.
The lift has scalar $\lambda_j(p,v)$ on the corresponding reduced
source part, for the source form $X_{T,j}$ and target form $X_{W_j}$.
In the selected record the two parts have the same dimension.
Indeed their defining preimages under $\iota$ and $\theta_j$
have the dimensions of the target parts because both coordinate
maps are bijections, and the record identifies these two dimensions
for each label. Thus an active row label has an active lift part
for every $j$. The height estimate uses the common scalar attached
to this label. Hypothesis~\ref{r4:R4:HL} asks for
$\sum_j\log\lambda_j(p,v)\ge\sum_{j,\rho}h_{j\rho}-C_0mfN^2\Lambda_a$.

\begin{proposition}[structure of Hypothesis~\ref{r4:R4:HL}]\label{w3:W3:HLprop}
Assume the preceding rank and scalar identities, with $\theta_j$
mapping the reduced source part of type $p$ into $W_{j,p}$. Suppose
that for each $j$, each active type $p$ and each $v$, for the adjoint
in these same forms,
\begin{enumerate}
\item[(i)] $\theta_j\theta_j^*$ acts on $W_{j,p}$ by a scalar $a_j\,z_j$ with $z_j=L^S(\frac{1+t_0}2,c\,\chi\mu_{t,j+1}^{-1})$ a partial $L$-value and $a_j$ a product of local zeta integrals (the Rallis inner product formula);
\item[(ii)] $\log(a_jz_j)\ge\sum_\rho\log\mathrm{hook}_{j\rho}-C_{\mathrm{sz}}fN^2L_p$, where $\log\mathrm{hook}_{j\rho}$ is the archimedean formal degree term written as a hook product, and $L_p$ is the size of the local and $L$-value corrections.
\end{enumerate}
If $L_p\le C_L\Lambda_a$ for an absolute $C_L>0$ and $\sum_\rho h_{j\rho}\le\sum_\rho\log\mathrm{hook}_{j\rho}$ for each $j$, the frame bound
in Hypothesis~\ref{r4:R4:HL} holds on the active parts with
$e_\lambda\le C_L C_{\mathrm{sz}}mfN^2\Lambda_a$.
\end{proposition}

\begin{proof}
Fix a nonzero $z$ in the reduced source part and put $y=\theta_jz$.
Then $\|y\|^2=\lambda_j(p,v)\|z\|^2>0$. The scalar identity on the
source gives $\theta_j\theta_j^* y=\theta_j(\theta_j^*\theta_jz)
=\lambda_j(p,v)y$. By (i) the same vector is $a_jz_j y$.
Thus $\lambda_j(p,v)=a_jz_j$. Sum (ii) over $j$ and apply the
bound for $L_p$. The height minimum only uses these active parts.
\end{proof}

\paragraph{The forms in the selected realization.}\label{frame:forms}
At the distinguished embedding take the reduced source to be the
rational preimage of the range of the adjoint of the analytic lift.
In finite dimension that range is the orthogonal complement of its
kernel. In reduced coordinates the complex lift sends a basis vector
to the analytic lift of the corresponding rational form. Its matrix
$T_j$ is the matrix of these images in the chosen target basis.
Rationality of its entries and the common image condition give
$w_{j,0}(T_j z)=\theta'_j(j_L z)$ in these coordinates, with
$j_L$ including the reduced source. If $T_j z=0$, the included vector
lies in the kernel and its orthogonal complement, so it vanishes;
the target basis spans the prescribed image. This gives the
bijection on the reduced spaces.

Three normalizations enter here. The Petersson pairing on multiplicity
maps is the integral of their matrix coefficients at unit representation
vectors over the chosen fundamental domain. The archimedean frame
multiplies this pairing by $\prod_\rho\dim K_\rho$, where $K_\rho$
is the lowest $K$-type at $\rho$. For a finite level representation
with group $H$ the adelic adjoint is $|H|^{-1}T^*$; at the trivial
level this agrees with the ordinary adjoint. In the scalar formula
the local index $I$ satisfies $\operatorname{vol}([V])=I\,\mathrm{Mass}(V)$
and $I\|f\|^2=\prod_{v\in S}d_v^{-1}$ with the specified test vector
and measures. Substitution cancels $I$ from this product, giving
$\operatorname{vol}([V])\|f\|^2/(\prod_{v\notin S}d_v)
=\mathrm{Mass}(V)/(\prod_v d_v)$.

At another embedding $\sigma$ representing $v$, the source uses a
lower bound on the reduced quadratic form. Write $s>0$ for the size of
the finite quotient of the chosen level, $C_\sigma$ for the sum of
the conjugate source pairings, and $J_{W,\sigma}$ for the realization
of the target coordinates. For a rational reduced vector $z$, the two
forms in that estimate are
\begin{align*}
 h_{T,j,v}(z,z)
   &=s^{-1}\operatorname{Re} C_\sigma(\mathrm{toF}\,z,\mathrm{toF}\,z),\\
 h_{W,j,v}(T_jz,T_jz)
   &=s^{-1}\|J_{W,\sigma}(\operatorname{coord}(T_jz))\|^2 .
\end{align*}
In the second line the realization scales the basis vectors by
$s^{-1/2}$. In the first line the majorant on the rational reduced
space is expressed in that same quotient normalization. The conjugate
lift estimate therefore gives $h_{W,j,v}(T_jz,T_jz)\ge
\exp(b_j)h_{T,j,v}(z,z)$, where
\[
 b_j=\mathrm{hookPlus}_j-8fn_V^2 L_p(B),\qquad
 L_p(B)=\log\!\bigl(B\,\operatorname{rd}_E\,g_0g_1 n_V\bigr).
\]
Explicitly, if $\lambda^+_{j\rho}$ is the Young diagram obtained by
sorting the $p_{j\rho}$ positive-side row lengths of the record, put
$H^+_{j\rho}=\prod_{\square\in\lambda^+_{j\rho}}h(\square)$, where
$h(\square)$ is one plus the number of boxes to its right in the same row
and below it in the same column.
Then $\mathrm{hookPlus}_j=\sum_\rho\log H^+_{j\rho}$.
At the base embedding the analytic main term contains the product
$H^+_{j\rho}H^-_{j\rho}$ for the two signature parts. The negative-side
hook product is a positive integer. Dropping its logarithm gives the
one-sided bound used here, including for an empty diagram.
For an active rational part, insert one nonzero vector into the scalar
identity in these forms and divide by its positive source norm.
This gives $\log\lambda_j(p,v)\ge b_j$. The selected realization uses
the lower bound at every off-base place and the positive eigenvalue
estimate at the base. The row scalar for the same $p$ uses precisely
this $\lambda_j$ in its product; both scalar identities belong to
one record. Thus transferring the scalar between the two spaces in
that record only requires the shared label and the rank equality.

For the rows (P1)--(P4) and the grid of Proposition~\ref{r6:R6:prop:nodes}, take $B=b$. The nodes $e_j=w_{j\rho}$ are positive and at most $2b$ for large $x$. The small rows satisfy $0\le v_l\le4k\Delta\le b_0/x\le2b/x\le b$. Within one embedding the grid gives $|b_l-e_j|\le mh_0\le b$, and (D5) gives $(kb)^8\le g_i$. These verify the size conditions for the indicated row parameters. The ranks give
$n_V=k+m-1+t_r\le N$, since $N-n_V=(m-1)t_r+1$.
Lemma~\ref{r5:lem:primesizes}(i),(iv) gives $N^2\le b$ and
$\log(g_0g_1)\le18L_\varepsilon$ at sufficiently large scale. Then
$\log n_V\le L_\varepsilon/2$ and
\[
 L_p(b)\le\tfrac{39}{2}L_\varepsilon+\log\operatorname{rd}_E
     \le20\Lambda_a .
\]
Here $L_\varepsilon=\log b$ and $\Lambda_a=\log(xb\operatorname{rd}_E)$.
Consequently the error summed over the lifts is bounded by
$160mfN^2\Lambda_a$. To apply this estimate to the earlier period
chain one must identify its row inclusion, type decomposition and
level forms with the selected record, and compare $\mathrm{hookPlus}_j$ with $\sum_\rho h_{j\rho}$.
For example, if the height uses another source frame, equality of its
Gram matrix in the specified rational basis with $X_{T,j}$ is needed
before applying the displayed quadratic inequality to that frame.


Item (i) is the Rallis inner product formula in the convergent range, which follows from the Siegel--Weil formula and the doubling method. Item (ii) combines Harish-Chandra's formula for the formal degree of a holomorphic discrete series, written as a hook product, with lower bounds for the partial $L$-value at $s=\frac{1+t_0}2$ in the range of absolute convergence ($\log z_j\ge-\sum_{\mathfrak p}\log(1-N\mathfrak p^{-\sigma})^{-1}$). \notgiven{Items (i) and (ii) for the theta lifts of the chain, and the identification of $\theta_j$ with the period map of Theorem~\ref{w2:B}, are not proved here.}

\subsection{The slope of the target lattice}\label{w3:sec:HT}

Recall from R4 that Hypothesis~\ref{r4:R4:HT} has three parts. The passage compares two levels of the chain, the base-side bound compares $\mu(R^\vee)$ with the covolumes of the base lattices, and the covolume bound evaluates those covolumes.

\subsubsection{The base-side bound}

Let $R\subset F^D$ and $M_1,\dots,M_m\subset F^D$ be lattices with hermitian forms $X_E$ on $R$ and $X_{W_j}$ on $M_j$. Suppose there are invertible matrices $\iota_t$ and $\theta_{j,t}$ indexed by local types $t$, lattices $A_t$ and $T_{j,t}$, and forms $X_{Q,t}$, $X_{T,j,t}$, such that:
\begin{enumerate}
\item[(a)] at every finite place $y$, $R_y=\iota_{\mu(y)}A_{\mu(y),y}$;
\item[(b)] at every archimedean $v$ and for the base type $t_0$,
\[
\det(\iota_{t_0}^*X_E\iota_{t_0})/\det X_{Q,t_0}=\kappa_v^D\prod_j\det(\theta_{j,t_0}^*X_{W_j}\theta_{j,t_0})/\det X_{T,j,t_0}
\]
\item[(c)] at a twisted type $t$ the same identity holds at the distinguished place $v_0$ with $\kappa_{v_0}\rho_t$ in place of $\kappa_{v_0}$, $\rho_t$ the archimedean value of an element of $F$ which is a unit at all finite places.
\end{enumerate}

\begin{proposition}[base side]\label{w3:W3:base}
Under (a)--(c),
\[
\lcv(R)\le D\langle\log\kappa\rangle+\sum_j\lcv(M_j)+\sum_y\ell_y\Big(v_y(G_{\mu(y),y})+\sum_j\big(i_y(M_j)+i_{\bar y}(M_j)+j_y(T_j)\big)\Big),
\]
where $\ell_y=\log Ny/[F:\mathbb Q]$, $G_{t,y}$ is the Gram determinant of a local basis of $A_{t,y}$ for $X_{Q,t}$, $i_y(M_j)$ is the length of $M_{j,y}$ relative to $\theta_{j,\mu(y)}T_{j,\mu(y),y}$, and $j_y$ is the local discriminant exponent of $T_j$ for the hermitian form. Consequently $\mu(R^\vee)-\langle\log\kappa\rangle\le\sum_j\lcv(M_j)/D+\frac1D\sum_y\ell_y(\cdots)$.
\end{proposition}

\begin{proof}
Choose a basis of $F^D$ and compute $\degh$ of each lattice as $\log$ of the index of the lattice relative to the standard lattice minus the archimedean log-determinants. The archimedean part of $\lcv(R)-\sum_j\lcv(M_j)$ is $\frac1{n_F}\sum_v\big(\log\det(\iota^*X_E\iota)-\sum_j\log\det(\theta_j^*X_{W_j}\theta_j)\big)$ plus the archimedean contributions of $X_{Q}$ and $X_{T,j}$, and (b), (c) turn it into $D\langle\log\kappa\rangle$ plus $\frac1{n_F}\sum_v(\log\det X_Q-\sum_j\log\det X_{T,j})+\log\lvert\rho\rvert$-terms. The product formula applied to $\det\iota_t$, $\det\theta_{j,t}$ and $\rho_t$ (a unit at every finite place) converts the archimedean $\log\lvert\det\rvert$-terms into finite valuations of the same determinants, and (a) together with the definitions of $G_{t,y}$, $i_y$ and $j_y$ identifies the resulting finite part with the sum over $y$. At untwisted places all local terms vanish by the hypotheses of the record. The second statement uses $\mu(R^\vee)=\lcv(R)/D$.
\end{proof}

The determinant identity (b) follows from the frame condition (the base-type subspaces are orthogonal, $\iota$ multiplies the form on $p$ by $\kappa_v\prod_j\lambda_j(p,v)$, $\theta_j$ multiplies it by $\lambda_j(p,v)$, and $\dim W_{j,p}=\dim E_p$), since the determinant of a map that is a scalar $c_p$ on orthogonal pieces of dimension $d_p$ is $\prod_pc_p^{d_p}$. The local terms on the right are bounded by Theorem~\ref{r5:thm:DIVinN} and enter $\mathcal L+E_b$.

\subsubsection{The passage}

\begin{lemma}[averaging]\label{w3:W3:avg}
Let $r_s>0$, $d_s\in\mathbb R$ and $n_s>0$ for $s$ in a finite set, with $d_s\le Hr_s$ for all $s$, and with $(1-\delta)\bar n\le n_s\le(1+\delta)\bar n$ for some $\bar n>0$ and $0\le\delta\le\frac13$. Put $\mu_0=\sum d_s/\sum r_s$ and $\mu_1=\sum n_sd_s/\sum n_sr_s$. Then $\mu_1\ge\mu_0-3\delta(H-\mu_0)$.
\end{lemma}

\begin{proof}
Write $d_s=\mu_0r_s+e_s$ with $\sum e_s=0$ and $e_s\le(H-\mu_0)r_s$. Then $\sum\lvert e_s\rvert=2\sum e_s^+\le2(H-\mu_0)\sum r_s$, and $\mu_1-\mu_0=\sum(n_s-\bar n)e_s/\sum n_sr_s\ge-\delta\bar n\sum\lvert e_s\rvert/((1-\delta)\bar n\sum r_s)\ge-\frac{2\delta}{1-\delta}(H-\mu_0)\ge-3\delta(H-\mu_0)$.
\end{proof}

\begin{lemma}[projectors]\label{w3:W3:proj}
Let a finite group $G$ act on $F^D$ preserving a lattice $R$ and the forms at every archimedean place, let $P\le G$ and $\psi\colon P\to F^\times$ a character with values in $\mathcal O_F^\times$, and let $e=\lvert P\rvert^{-1}\sum_{p}\psi(p)^{-1}p$. Then $e$ is an orthogonal projector, $\lvert P\rvert\,eR\subset R\cap eF^D\subset eR$, and every saturated sublattice $Y\subset eF^D$ satisfies $\lvert\mu(Y\cap eR)-\mu(Y\cap R)\rvert\le2\log\lvert P\rvert$.
\end{lemma}

\begin{proof}
$e^2=e$ because $\psi$ is a character, and $e^*=e$ because each $p$ is unitary and $\psi(p)^{-1}=\overline{\psi(p)}$. If $x\in R$ then $\lvert P\rvert ex=\sum\psi(p)^{-1}px\in R$. If $x\in R\cap eF^D$ then $x=ex\in eR$. Hence $Y\cap R$ has index at most $\lvert P\rvert^{\rk Y}$ in $Y\cap eR$ (the quotient is killed by $\lvert P\rvert$), and the slopes differ by at most $2\log\lvert P\rvert$ in our normalization.
\end{proof}

\begin{hypothesis}[near-constant multiplicities]\label{w3:W3:mult}
At the base level the lattice $R^{(0)}$ decomposes up to index $\lvert G\rvert^{\rk}$ as $\bigoplus_sR_s$ over the irreducible constituents $\tau_s$, and at the sharp level $R^{(1)}$ decomposes as $\bigoplus_sR_s^{\oplus n_s}$ with $n_s$ the multiplicity of $\psi_1$ in $\tau_s|_{P_1}$. These multiplicities satisfy $(1-\delta)\bar n\le n_s\le(1+\delta)\bar n$ with $\bar n=\prod_{P\in\mathcal P_g}\lvert k_P\rvert^k$ and $\delta=8fr/\min(g_0,g_1)$.
\end{hypothesis}

The near-constancy follows from Frobenius reciprocity and a character sum count $q^k(1-1/q)\le n\le q^k(1+(2r-1)/q)$. \notgiven{Two of the counts behind Hypothesis~\ref{w3:W3:mult}, the formula for the multiplicity of $\psi_1$ in a parabolically induced constituent and the goodness of the cuspidal products that occur, are not proved here.}

\begin{proposition}[passage]\label{w3:W3:passage}
Under Hypothesis~\ref{w3:W3:mult}, with $\lambda_0=\log\lvert G\rvert$ and $H_0\ge\mu_{\max}$ at the base level,
\[
\mu^{(1)}(R^\vee)\le\mu^{(0)}(R^\vee)+3\delta\big(\mu^{(0)}(R^\vee)+H_0+28\lambda_0\big)+28\lambda_0 .
\]
\end{proposition}

\begin{proof}
Apply Lemma~\ref{w3:W3:proj} at each level to replace $R^{(i)}$ by the lattice cut out by the projector, at a cost of $O(\lambda_0)$ in slope and in maximal slope. Apply Lemma~\ref{w3:W3:avg} to $r_s=\rk R_s$, $d_s=\degh(R_s)/n_F$ and the multiplicities $n_s$, with $H=H_0+O(\lambda_0)$, to get $\mu^{(1)}(R)\ge\mu^{(0)}(R)-3\delta(H-\mu^{(0)}(R))-O(\lambda_0)$. Dualizing with $\mu(R^\vee)=-\mu(R)$ gives the statement. \notgiven{The constant in the $O(\lambda_0)$ terms is not computed here.}
\end{proof}

With $\delta\le\frac13$, $\lambda_0\le b$ and $\mu^{(0)}+H_0$ bounded by the record, this gives $P_{\mathrm{pass}}\le1+C_0\lambda_0$ of Hypothesis~\ref{r4:R4:HT}.

\subsubsection{The covolume bound}

\begin{hypothesis}[covolume bound]\label{w3:W3:L3}
There is an absolute constant $C$ such that for every admissible datum and every level $1\le j\le m$ of the chain,
\[
\frac{\lcv(M_j)}D\le\sum_{\rho\in\Phi}(p_j-q_j)\,w_{j\rho}\log\Omega+C\,U_j,
\]
where $U_j$ is bounded by the terms $C_0mx^\varepsilon(B_0+N_c)$ of Hypothesis~\ref{r4:R4:HT}.
\end{hypothesis}

\paragraph{Two normalizations.}
Let $F$ be a totally complex number field with $n_F=[F:\mathbb Q]/2$ archimedean places. For a hermitian lattice $\Lambda\neq0$ over $\mathcal O_F$ we use the slope $\mu(\Lambda)=\degh(\Lambda)/(n_F\rk\Lambda)$ of \eqref{r4:R4:slope} and we put
\[
\delta(\Lambda)=\frac{\degh(\Lambda)}{[F:\mathbb Q]}=\frac12\,\rk(\Lambda)\,\mu(\Lambda).
\]
We call $\delta$ the degree in the $[F:\mathbb Q]$-normalization. For an $\mathcal O_F$-sublattice we write $\mu_{\max}(\Lambda)$ for the supremum of $\mu(\Lambda')$ over the nonzero $\mathcal O_F$-submodules $\Lambda'\subseteq\Lambda$. We use four elementary facts. If $\Lambda'\subseteq\Lambda$ have the same rank and the same forms, then $\degh(\Lambda')\le\degh(\Lambda)$. For an integer $g\ge1$, $\mu(g^{-1}\Lambda)=\mu(\Lambda)+2\log g$, because $[g^{-1}\Lambda:\Lambda]=g^{2n_F\rk\Lambda}$. If $\Lambda_1\subseteq\Lambda$ is saturated and we give the quotient its induced forms, then $\degh(\Lambda)=\degh(\Lambda_1)+\degh(\Lambda/\Lambda_1)$. If $y_1,\dots,y_r\in\Lambda$ are linearly independent over $F$ and span $\Lambda\otimes F$, then Hadamard's inequality gives $\degh(\Lambda)\ge-\sum_v\sum_i\log h_v(y_i,y_i)$.

The formal proof measures sections on a cover by a volume slope. Let $Z$ be a component of a cover over a field $F$, let $u$ be a weight with $\deg_Z(u)>0$, where $\deg_Z(u)$ is the top geometric degree of the automorphic line bundle of weight $u$ on $Z$, and let $\Lambda_B(Z;ku)$ be the lattice of integral sections of weight $ku$ for the integral structure of the cutoff $B$ described in Hypothesis~\ref{w3:W3:unif}. With $D_{\mathrm{fl}}$ the dimension of the flag variety put
\begin{gather*}
V_B(Z;u;t)=\limsup_{k\to\infty}\frac{D_{\mathrm{fl}}!\,\log\#S_B(Z;8ku;8kt)}{[F:\mathbb Q]\,(8k)^{D_{\mathrm{fl}}+1}},\\
\sigma_B(Z;u)=\inf_{t\in\mathbb R}\Big(\frac{V_B(Z;u;t)}{\deg_Z(u)}-t\Big),
\end{gather*}
where $S_B(Z;v;t)$ is the set of $f\in\Lambda_B(Z;v)$ with $\lVert f\rVert\le e^{t}$ at every frame.
The number $\sigma_B(Z;u)$ divides by $[F:\mathbb Q]$. On the lattices of the chain and of the covers the formal proof uses $\mu$, which divides by $n_F$. A lower bound without main term for $\sigma_B$ turns into a lower bound for $\mu$ through the factor $2$ of input (I17) below, and that factor multiplies a constant. The main term is compared only in Lemma~\ref{cov:twist}, which is stated for $\mu$ and in which the flat degree $\frac12(p-q)\log\Omega$ per unit weight is a degree $\delta$. Remark~\ref{cov:factor2} records the bookkeeping.

\paragraph{The covolume inequality.}
Write $\Lambda_B=\log(B\,\mathrm{rd}_E)+|\lambda_1|$ and let $C_{\mathrm{unif}}$ denote the absolute constant in Hypothesis~\ref{w3:W3:unif}.

For the chain we write $p_j=k+j-1$ and $q_j=m-j$, so that $\lvert p_j-q_j-k\rvert\le m$, and we write $w_{j\rho}=e_{j}(\rho)$ for the exponent of index $j$ of the archimedean parameter at $\rho$ (the weights $w_{j\rho}$ of Hypothesis~\ref{r4:R4:HW}). Put
\[
\mathrm{main}_j=\sum_{\rho\in\Phi}(p_j-q_j)\,w_{j\rho}\log\Omega,\qquad U=B_0\log\log x,\qquad B_0=fmb\,x^\varepsilon .
\]
Since $\log\log x\le x^\varepsilon$ for large $x$ and $m\ge1$, $U\le m\,x^\varepsilon(B_0+N_c)$, so $U$ is admissible as the $U_j$ of Hypothesis~\ref{w3:W3:L3}.

\begin{hypothesis}[covolumes on the chain]\label{cov:H58}
There is $\varepsilon_1>0$ such that for every $0<\varepsilon\le\varepsilon_1$ and every admissible family there are $C_0\ge0$ and $x_0$ with the following property. For every $x\ge x_0$, every point of the family at $x$, every $1\le j\le m$ and every $C\ge C_0$,
\[
\frac{\lcv(M_j)}D\le\mathrm{main}_j+C\,U .
\]
\end{hypothesis}

\paragraph{Four lattice lemmas.}

\begin{lemma}[truncation]\label{cov:trunc}
Let $L\subseteq M$ be hermitian $\mathcal O_F$-lattices with $L\ne0$, let $e$ be an $F$-linear idempotent of $M\otimes F$ with $h_v(ex+(y-ey),ex+(y-ey))=h_v(ex,ex)+h_v(y-ey,y-ey)$ for all $x,y$ and all $v$, and let $g\ge1$ be an integer with $g\,e(L)\subseteq L$. Put $L_1=L\cap\operatorname{im}e$ and $L_0=L\cap\ker e$. Let $0\le\delta_0\le\frac1{12}$ and $Q\ge0$ satisfy $\rk L_0\le3\delta_0\rk L$, $\mu_{\max}(M)\le Q$, $-Q\le\mu(L)$ and $2\log g\le Q$. Then $\mu(L_1)\ge\mu(L)-12\delta_0Q$.
\end{lemma}

\begin{proof}
Write $r=\rk L$, $r_1=\rk L_1$ and $r_0=\rk L_0$. For $x\in L$ we have $g\,ex\in L_1$, so $(L_1)\otimes F=e(L\otimes F)$. Since $L\otimes F=e(L\otimes F)\oplus(1-e)(L\otimes F)$ orthogonally, the map $1-e$ identifies $L/L_1$ isometrically with $(1-e)L$, and $r=r_1+r_0$. For $x\in L$ the element $g(1-e)x=gx-g\,ex$ lies in $L\cap\ker e=L_0$, so $L_0\subseteq(1-e)L\subseteq g^{-1}L_0$. If $r_0>0$, then $\mu(L_0)\le\mu_{\max}(M)\le Q$, hence
\[
\degh(L/L_1)=\degh((1-e)L)\le\degh(g^{-1}L_0)=n_Fr_0\big(\mu(L_0)+2\log g\big)\le2n_Fr_0Q,
\]
and the same bound holds trivially when $r_0=0$. Since $L_1$ is saturated in $L$, $\degh(L_1)=\degh(L)-\degh(L/L_1)\ge n_F(r\mu(L)-2r_0Q)$. From $r_0\le3\delta_0r\le\frac r4$ we get $r_1\ge\frac34r>0$ and $r_0/r_1\le3\delta_0/(1-3\delta_0)\le4\delta_0$. Hence
\[
\mu(L_1)-\mu(L)\ge\frac{r\mu(L)-2r_0Q}{r_1}-\mu(L)=\frac{r_0(\mu(L)-2Q)}{r_1}\ge-\frac{3r_0Q}{r_1}\ge-12\delta_0Q,
\]
where we used $-Q\le\mu(L)\le\mu_{\max}(M)\le Q$.
\end{proof}

\begin{lemma}[descent]\label{cov:descent}
Let $A\ne0$ and $B$ be hermitian $\mathcal O_F$-lattices, let $\varphi_1,\dots,\varphi_n\colon A\otimes F\to B\otimes F$ be $F$-linear maps with $\varphi_h(A)\subseteq B$, with $\bigcap_h\ker\varphi_h=0$ and with $\sum_h\lVert\varphi_hf\rVert_v^2\le e^{2P}\lVert f\rVert_v^2$ for all $f$ and $v$. Suppose that $B$ is spanned over $F$ by elements of $B$ of norm at most $e^\Lambda$ at every place, with $\Lambda\ge0$, and that $n\rk B\le(1+\epsilon)\rk A$ with $\epsilon\ge0$. If $t\le\mu(A)-2P$, then $\mu(B)\ge t-\epsilon(\lvert t\rvert+2\Lambda)$.
\end{lemma}

\begin{proof}
Give $B^n$ the orthogonal sum of the forms of $B$ and put $\Phi=(\varphi_h)_h\colon A\to B^n$. The map $\Phi$ is injective with $\lVert\Phi f\rVert_v\le e^P\lVert f\rVert_v$. Put $r=\rk A$ and $N=n\rk B$, so $r\le N\le(1+\epsilon)r$. For a basis $e_1,\dots,e_r$ of $A\otimes F$ in $A$ the index of $\sum\mathcal O_F\Phi e_i$ in $\Phi(A)$ equals that of $\sum\mathcal O_Fe_i$ in $A$, and the Gram determinants satisfy $\det h_v(\Phi e_i,\Phi e_j)\le e^{2Pr}\det h_v(e_i,e_j)$. Hence $\mu(\Phi A)\ge\mu(A)-2P\ge t$. Let $S$ be the saturation of $\Phi(A)$ in $B^n$, so $\degh(S)\ge\degh(\Phi A)\ge n_Frt$. The quotient $B^n/S$ has rank $N-r$ and is spanned over $F$ by images of the elements $(0,\dots,y,\dots,0)$ with $y$ in the given spanning set, whose quotient norms are at most $e^\Lambda$. Choosing $N-r$ of these images linearly independent and applying Hadamard's inequality gives $\degh(B^n/S)\ge-2n_F(N-r)\Lambda$. Therefore $n\degh(B)=\degh(B^n)\ge n_F\big(rt-2\Lambda(N-r)\big)$ and
\[
\mu(B)=\frac{n\degh(B)}{n_FN}\ge\frac{rt}N-2\Lambda\frac{N-r}N\ge t-\lvert t\rvert\frac{N-r}N-2\Lambda\frac{N-r}N\ge t-\epsilon(\lvert t\rvert+2\Lambda),
\]
because $(N-r)/N\le(N-r)/r\le\epsilon$.
\end{proof}

\begin{lemma}[from component heights to the slope]\label{cov:mean}
Let $\mathcal Z$ be a finite nonempty set with $n=\#\mathcal Z$ elements. Let $T,D_{\mathrm{fl}},A,R,c_s,c_c,X,\rho,b_n,\delta'$ and $E$ be real numbers with $D_{\mathrm{fl}}\ge0$, $T\ge3D_{\mathrm{fl}}+5$, $A>0$, $R\ge0$, $c_s,c_c,X\ge0$, $0\le\rho<1$, $b_n>0$, $\delta'>0$ and $E\ge1$, and for $Z\in\mathcal Z$ let $h_Z,s_Z,\delta_Z,r_Z,d_Z$ be real numbers. Assume
\begin{enumerate}
\item[(a)] $\delta'\le\delta_Z\le\delta'/(1-\rho)$ and $b_n\delta_Z\le r_Z$ for every $Z$,
\item[(b)] $h_Z+c_s\ge0$ and $(1-\rho)\max(0,s_Z+c_s)\le h_Z+c_s+c_c$ for every $Z$, and $\frac1n\sum_Zs_Z\ge-X$,
\item[(c)] $d_Z\ge b_n\delta_Z\big(Th_Z-3(D_{\mathrm{fl}}+1)\lvert h_Z\rvert-A\log(2+\lvert h_Z\rvert)-R\big)$ for every $Z$,
\item[(d)] $TX\le C_1E$, $Tc_c\le C_1E$, $c_s\big(6D_{\mathrm{fl}}+10+2T\rho\big)\le C_1E$ and $A\log\max(A,1)+R\le C_2E$.
\end{enumerate}
Then $2\sum_Zd_Z/\sum_Zr_Z\ge-2(3C_1+C_2+2)E$.
\end{lemma}

\begin{proof}
Put $H_Z=h_Z+c_s\ge0$. If $h_Z\ge0$ then $\lvert h_Z\rvert=h_Z$, and if $h_Z<0$ then $\lvert h_Z\rvert=-h_Z\le c_s$, so in both cases $\lvert h_Z\rvert\le h_Z+2c_s$, and $y_Z=h_Z+2c_s\ge0$. For $y\ge0$ the function $A\log(2+y)-y$ is at most $A\log A-A+2$ when $A\ge2$ and at most $A\log2<2$ when $A<2$, so $A\log(2+y)\le y+A\log\max(A,1)+2$. Hence the bracket in (c) satisfies
\[
\beta_Z\ge(T-3D_{\mathrm{fl}}-3)h_Z-6(D_{\mathrm{fl}}+1)c_s-(h_Z+2c_s)-A\log\max(A,1)-2-R=a\,h_Z-K_0,
\]
with $a=T-3D_{\mathrm{fl}}-4\in[1,T]$ and $K_0=(6D_{\mathrm{fl}}+8)c_s+A\log\max(A,1)+2+R$. By (b), $H_Z\ge(1-\rho)(s_Z+c_s)-c_c$, so $\sum_ZH_Z\ge n\big((1-\rho)(c_s-X)-c_c\big)$. Since $H_Z\ge0$, (a) gives $\sum_Z\delta_ZH_Z\ge\delta'\sum_ZH_Z$ and $\sum_Z\delta_Z\le n\delta'/(1-\rho)$, so
\[
\frac{\sum_Z\delta_ZH_Z}{\sum_Z\delta_Z}\ge(1-\rho)\max\Big(0,\frac1n\sum_ZH_Z\Big)\ge(1-\rho)^2(c_s-X)-(1-\rho)c_c .
\]
With $h_Z=H_Z-c_s$ and $1-(1-\rho)^2\le2\rho$ this gives $m_\delta:=\sum_Z\delta_Zh_Z/\sum_Z\delta_Z\ge-X-c_c-2\rho c_s$. Hence
\[
\frac{\sum_Z\delta_Z(ah_Z-K_0)}{\sum_Z\delta_Z}=a\,m_\delta-K_0\ge-T(X+c_c+2\rho c_s)-K_0=:-K_1 ,
\]
using $1\le a\le T$ when $m_\delta<0$. Put $N=\sum_Zd_Z$, so $N\ge b_n\sum_Z\delta_Z(ah_Z-K_0)$ by (c). If $N\ge0$ then $N/\sum_Zr_Z\ge0\ge-K_1$. If $N<0$ then, since $\sum_Zr_Z\ge b_n\sum_Z\delta_Z>0$ by (a), $N/\sum_Zr_Z\ge N/(b_n\sum_Z\delta_Z)\ge-K_1$. By (d) and $E\ge1$, $K_1\le TX+Tc_c+c_s(6D_{\mathrm{fl}}+8+2T\rho)+A\log\max(A,1)+R+2\le(3C_1+C_2+2)E$.
\end{proof}

\begin{lemma}[flat twist]\label{cov:twist}
Let $\Psi$ be a finite set and let $a_\tau\in\mathbb Z$ and $c_\tau\in\mathbb R$ for $\tau\in\Psi$. Let $p\ge q\ge0$ be integers and let $\log\Omega,\ L\ge10,\ b,\ U,\ \mathrm{main},\ C_A\ge0,\ l_r,\ m_0\ge0,\ f_r\ge0,\ n_V,\ l_S$ be real numbers with
\begin{enumerate}
\item[(a)] $\lvert c_\tau+\frac12(p-q)\log\Omega\rvert\le3qL$ for every $\tau$,
\item[(b)] $\big\lvert\mathrm{main}-\sum_\tau(p-q)a_\tau\log\Omega\big\rvert\le C_AU$ and $q\,L\sum_\tau\lvert a_\tau\rvert\le C_AU$,
\item[(c)] $l_r\le2L$, $2m_0\le n_V$, $f_rn_VL\le b$, $b\le U$, $L\le U$ and $l_S\le40L$.
\end{enumerate}
Let $\Lambda$ and $\Lambda'$ be hermitian lattices of the same rank with
\[
\delta(\Lambda')\ge\delta(\Lambda)+\rk(\Lambda)\Big(\sum_\tau a_\tau c_\tau-q\,l_r\sum_\tau\lvert a_\tau\rvert-m_0f_r\big(\tfrac12\log(2\pi)+l_r\big)-l_S\Big).
\]
Then $\mu(\Lambda')\ge\mu(\Lambda)-\mathrm{main}-(11C_A+83)U$.
\end{lemma}

\begin{proof}
Since $\mu=2\delta/\rk$, the hypothesis gives $\mu(\Lambda')\ge\mu(\Lambda)+2\Theta$ with $\Theta$ the bracket. By (a), $a_\tau c_\tau\ge-\frac12(p-q)a_\tau\log\Omega-3qL\lvert a_\tau\rvert$, so by (b)
\[
2\sum_\tau a_\tau c_\tau\ge-\sum_\tau(p-q)a_\tau\log\Omega-6qL\sum_\tau\lvert a_\tau\rvert\ge-\mathrm{main}-7C_AU .
\]
By (c) and (b), $2q\,l_r\sum_\tau\lvert a_\tau\rvert\le4qL\sum_\tau\lvert a_\tau\rvert\le4C_AU$. Since $\log(2\pi)\le2\pi-1\le7$ and $L\ge10$, $\frac12\log(2\pi)+l_r\le3L$, so $2m_0f_r(\frac12\log(2\pi)+l_r)\le3n_Vf_rL\le3b\le3U$. Finally $2l_S\le80L\le80U$. Adding, $2\Theta\ge-\mathrm{main}-(11C_A+83)U$.
\end{proof}

\begin{remark}[the factor $2$]\label{cov:factor2}
Lemma~\ref{cov:twist} is the only place where the main term is compared. Its input (a) is the flat degree of input (I13), which is a degree $\delta$ in the $[F:\mathbb Q]$-normalization, namely $\frac12(p-q)\log\Omega$ per unit of the anchor weight up to $3qL$. The conversion $\mu=2\delta/\rk$ doubles it, and the doubled value $\sum_\tau(p-q)a_\tau\log\Omega$ is the main term in the normalization $\mu$ of \eqref{r4:R4:slope}, the normalization of $\lcv(M_j)/D$. This agrees with Definition~\ref{r4:R4:Omega}, where $\log\Omega_\tau$ is $2/[M:\mathbb Q]$ times an arithmetic degree, that is, a slope $\mu$ of a line. The lower bound of Hypothesis~\ref{w3:W3:unif} is a bound for $\sigma_B$, which divides by $[F:\mathbb Q]$, and it enters only through Lemma~\ref{cov:mean} and input (I17), whose factor $2$ appears in the constant $2(3C_1+C_2+2)$.
\end{remark}

\begin{lemma}[size of the error of Hypothesis~\ref{w3:W3:unif}]\label{cov:size}
Let $T,f,m,b,L,x,\ell,q$ and $\mathrm{rd}_E$ be positive reals with $L\ge10$, $b=e^L$, $\log\log x\ge1$, $\ell\le9L$ and $\log\mathrm{rd}_E\le2L$, put $U=fmbL\log\log x$, and let $u'$ be a weight and $B$ a real number with $T\lVert u'\rVert_1\le834fmb$, $1\le B\le132b$ and $q\,\mathrm{spread}(u')\le6\lVert u'\rVert_1$. Let $c_B$ be a real number with $0\le c_B\le\log B+20+44\log\mathrm{rd}_E$. Then
\[
T\big((\log B+\ell)\lVert u'\rVert_1+q\,\mathrm{spread}(u')\big)\le83400\,U,\qquad T\,c_B\lVert u'\rVert_1\le83400\,U.
\]
\end{lemma}

\begin{proof}
From $B\le132b$ we get $\log B\le\log132+L\le L+22$. Hence $\log B+\ell+6\le10L+28\le100L$ and $c_B\le L+42+88L\le100L$, because $L\ge10$. So both left sides are at most $100L\cdot T\lVert u'\rVert_1\le83400fmbL\le83400U$.
\end{proof}

\paragraph{The objects of the chain and the inputs.}
Fix $\varepsilon$, an admissible family, $x$, a point of the family at $x$ and a level $1\le j\le m$. At this level the formal proof attaches to the point hermitian $\mathcal O_F$-lattices $\Lambda^{\mathrm{type}}_j\subseteq\Lambda^{\mathrm{all}}_j$ (forms of the local types inside all forms of the level, with the archimedean norms of the level), an idempotent $e_j$ of $\Lambda^{\mathrm{all}}_j\otimes F$ with a finite group $G_j$, the good part $\Lambda^{\mathrm{good}}_j=\Lambda^{\mathrm{type}}_j\cap\operatorname{im}e_j$ and the bad part $\Lambda^{\mathrm{bad}}_j=\Lambda^{\mathrm{type}}_j\cap\ker e_j$, and a core lattice $\Lambda^{\mathrm{core}}_j$. For $j<m$ it attaches deep covers $Y$ of the unitary Shimura variety of the datum of level $j$, of signature $(p,q)$, with lattices $\Lambda_Y(G)$ of cusp forms of weight $G$ and slopes $\mu_Y(G)=\mu(\Lambda_Y(G))$, a sublattice $\Lambda^{\mathrm{deep}}_Y\subseteq\Lambda_Y(G_k)$ cut out by the forms that come from the chain, the shifted weights $G_s(\pm c_0)$, the weight $u^-$ and the anchor weight $a=(a_\tau)$ with $G_s(-c_0)$ equal, up to an equivalence that preserves lattices, ranks and slopes, to $u^-+a$ shifted by $m_0=n_V/2$. We write $K_j=PP\cdot\kappa_j$ for the size parameter of the level, $T=\mathrm{len}_T\ge1$ for the length of the filtration and $D_{\mathrm{fl}}$ for the dimension of the flag variety. Assume the following bounds with constants independent of $x$, of the point and of $j$, valid for $x\ge x_0$, where $x_0$ and the constants depend only on $\varepsilon$ and the family. We write $C_K,C_D,C_A,C_s,C_m,C_{HS}$ for these constants.
\begin{enumerate}
\item[(I1)] (recursion) $\lcv(M_j)/D\le-\mu(\Lambda^{\mathrm{good}}_j)$.
\item[(I2)] (cut) $e_j$ is idempotent and orthogonal at every embedding, $\lvert G_j\rvert\,e_j(\Lambda^{\mathrm{type}}_j)\subseteq\Lambda^{\mathrm{type}}_j$, $2\log\lvert G_j\rvert\le b$ and $\Lambda^{\mathrm{good}}_j\ne0$.
\item[(I3)] (bad locus) $\rk\Lambda^{\mathrm{bad}}_j\le3e^{-x^{\varepsilon_z/7}}\rk\Lambda^{\mathrm{type}}_j$, with $\varepsilon_z>0$ a parameter of the family.
\item[(I4)] (maximal slope) $\mu_{\max}(\Lambda^{\mathrm{all}}_j)\le x^Ab$.
\item[(I5)] (sizes) $1\le b\le U$, $\lvert\mathrm{main}_j\rvert\le x^{A'}b$, $\lvert\mathrm{main}_j\rvert\le K_j$, $b\le K_j$, and for $j<m$ the set of quotient weights $u_1$ is nonempty and each has $3D_{\mathrm{fl}}+5\le T$, $0\le\rho(u_1)<1$ and $0\le z(u_1)<1$, with $\rho,z$ the numbers of the filtration.
\item[(I6)] (core comparison) $\mu(\Lambda^{\mathrm{type}}_j)\ge\mu(\Lambda^{\mathrm{core}}_j)-C_KU$.
\item[(I7)] (top level, CM evaluation) for $j=m$, $\mu(\Lambda^{\mathrm{core}}_m)\ge-\mathrm{main}_m-C_DU$.
\item[(I8)] (existence) for $j<m$ there is a good deep cover $Y$ at level $j$, and it satisfies the hypotheses of Hypothesis~\ref{w3:W3:unif}.
\item[(I9)] (ring) $\mu(\Lambda^{\mathrm{deep}}_Y)-2b\le\mu(\Lambda^{\mathrm{core}}_j)$.
\item[(I10)] (maps) there are $n$ maps $\Lambda_Y(G_s(-c_0))\to\Lambda^{\mathrm{deep}}_Y$ satisfying the hypotheses of Lemma~\ref{cov:descent} with $P=b$, and $n\rk\Lambda^{\mathrm{deep}}_Y\le\rk\Lambda_Y(G_s(c_0))$.
\item[(I11)] (rank) $\rk\Lambda_Y(G_s(c_0))\le(1+b/K_j)\rk\Lambda_Y(G_s(-c_0))$ and $\Lambda_Y(G_s(-c_0))\ne0$.
\item[(I12)] (generators) $\Lambda^{\mathrm{deep}}_Y$ is spanned over $F$ by elements of norm at most $e^{K_j}$ at every embedding.
\item[(I13)] (flat twist) the hypotheses of Lemma~\ref{cov:twist} hold for $\Lambda=\Lambda_Y(u^-)$, $\Lambda'=\Lambda_Y(G_s(-c_0))$, the anchor weight $a$, $c_\tau$ the mean over the components of minus the degree $\delta$ of the flat section of type $\tau$, $\mathrm{main}=\mathrm{main}_j$, $l_r=\log\mathrm{rd}_E$, $f_r=f(p+q)$, $l_S=\sum_{\ell\in S_{\mathrm{bad}}}\log\ell$, with the constant $C_A$.
\item[(I14)] (heights of components) for $j<m$, every quotient weight $u_1$ and every component $Z$ of $Y$ there are the numbers of Lemma~\ref{cov:mean}, namely the height $h_Z$, the degree $\delta_Z$, the rank $r_Z$ and the degree $d_Z$ of the filtered pieces of weight $u^-$, $b_n=\binom{T-1}{D_{\mathrm{fl}}}$, $A,R$, $c_s$ and a multiplicity cost $c_m\ge0$, satisfying (a), (c), $h_Z+c_s\ge0$, $c_s(6D_{\mathrm{fl}}+10+T(z+2\rho))\le C_sU$, $Tc_m\le C_mU$ and $A\log\max(A,1)+R\le C_{HS}U$.
\item[(I15)] (comparison of the volume slope with the height) with $u'=u'(u_1)$ the weight attached to $u_1$, $g$ its gap and $B$ as in Hypothesis~\ref{w3:W3:unif}, every real $s\le\sigma_B(Z;u')$ satisfies $(1-\rho)\max(0,s+c_s)\le h_Z+c_s+c_m+c_B\lVert u'\rVert_1$, where $c_B$ is the cost of the Hodge--Tate integral structure up to $B$.
\item[(I16)] (numerics of the weight) the hypotheses of Lemma~\ref{cov:size} hold for $u'$, $B$, $c_B$ and $T$, and $u'$ satisfies the hypotheses on $u$ and $g$ of Hypothesis~\ref{w3:W3:unif}.
\item[(I17)] (sum) $\mu_Y(u^-)\ge2\sum_Zd_Z/\sum_Zr_Z$.
\end{enumerate}
\notgiven{Inputs (I1)--(I17) are not proved here.}

\begin{theorem}\label{cov:main}
Assume Hypothesis~\ref{w3:W3:unif} and inputs (I1)--(I17). Put $C_F=11C_A+83$, $C_H=C_s+C_m+83400(1+C_{\mathrm{unif}})$, $C_U=2(3C_H+C_{HS}+2)$, $C_{\mathrm{low}}=\max\big(C_D,2(C_F+C_U)+9\big)$ and $C_0=C_K+C_{\mathrm{low}}+1$. Then Hypothesis~\ref{cov:H58} holds with this $C_0$, for $x$ large in terms of $\varepsilon$, the family, $A$, $A'$ and $C_0$.
\end{theorem}

\begin{proof}
Step 1 (slope of the weight $u^-$ on the cover). Let $j<m$, let $Y$ be the cover of (I8) and let $u_1$ be a quotient weight, which exists by (I5). By (I16) and (I8), Hypothesis~\ref{w3:W3:unif} applies to $u'=u'(u_1)$ and its gap $g$, and gives numbers $s_Z\le\sigma_B(Z;u')$ with mean at least $-X$, $X=C_{\mathrm{unif}}\mathrm{err}(u',g)$. By (I15), condition (b) of Lemma~\ref{cov:mean} holds with $c_c=c_m+c_B\lVert u'\rVert_1$. By Lemma~\ref{cov:size}, $TX\le83400\,C_{\mathrm{unif}}U$ and $Tc_c\le(C_m+83400)U$. With (I14) and $2T\rho\le T(z+2\rho)$, condition (d) of Lemma~\ref{cov:mean} holds with $C_1=C_H$, $C_2=C_{HS}$ and $E=U\ge1$. Lemma~\ref{cov:mean} and (I17) give $\mu_Y(u^-)\ge-C_UU$.

Step 2 (flat twist). By (I13) and Lemma~\ref{cov:twist}, $\mu_Y(G_s(-c_0))\ge\mu_Y(u^-)-\mathrm{main}_j-C_FU\ge-\mathrm{main}_j-(C_F+C_U)U$.

Step 3 (descent to the level of the chain). Apply Lemma~\ref{cov:descent} to $A=\Lambda_Y(G_s(-c_0))$, $B=\Lambda^{\mathrm{deep}}_Y$, the maps of (I10), $P=b$, $\Lambda=K_j$, $\epsilon=b/K_j\in[0,1]$ and $t=-\mathrm{main}_j-(C_F+C_U)U-2b$. Its hypotheses hold by (I10), (I11), (I12), (I5) and Step 2. It gives $\mu(\Lambda^{\mathrm{deep}}_Y)\ge t-\epsilon(\lvert t\rvert+2K_j)$. Since $\epsilon K_j=b$, $\epsilon\lvert\mathrm{main}_j\rvert\le b$ and $\epsilon\le1$, we get $\epsilon(\lvert t\rvert+2K_j)\le b+(C_F+C_U)U+2b+2b$, so $\mu(\Lambda^{\mathrm{deep}}_Y)\ge-\mathrm{main}_j-2(C_F+C_U)U-7b$. By (I9) and $b\le U$, $\mu(\Lambda^{\mathrm{core}}_j)\ge-\mathrm{main}_j-(2(C_F+C_U)+9)U$. For $j=m$, (I7) gives $\mu(\Lambda^{\mathrm{core}}_m)\ge-\mathrm{main}_m-C_DU$. In both cases $\mu(\Lambda^{\mathrm{core}}_j)\ge-\mathrm{main}_j-C_{\mathrm{low}}U$.

Step 4 (truncation). Put $\delta_0=e^{-x^{\varepsilon_z/7}}$ and $Q=(x^A+x^{A'}+C_K+C_{\mathrm{low}}+1)U$. Take $x$ so large that $12\delta_0(x^A+x^{A'}+C_K+C_{\mathrm{low}}+1)\le1$, which is possible since $\delta_0$ decays faster than any power of $x$. Then $\delta_0\le\frac1{12}$ and $12\delta_0Q\le U$. By (I6) and Step 3, $\mu(\Lambda^{\mathrm{type}}_j)\ge-\mathrm{main}_j-(C_K+C_{\mathrm{low}})U\ge-x^{A'}b-(C_K+C_{\mathrm{low}})U\ge-Q$ by (I5). By (I4) and (I5), $\mu_{\max}(\Lambda^{\mathrm{all}}_j)\le x^Ab\le Q$, and by (I2), $2\log\lvert G_j\rvert\le b\le Q$. Lemma~\ref{cov:trunc}, with (I2) and (I3), gives $\mu(\Lambda^{\mathrm{good}}_j)\ge\mu(\Lambda^{\mathrm{type}}_j)-U$. By (I1),
\[
\frac{\lcv(M_j)}D\le-\mu(\Lambda^{\mathrm{good}}_j)\le\mathrm{main}_j+(C_K+C_{\mathrm{low}}+1)U=\mathrm{main}_j+C_0U\le\mathrm{main}_j+CU
\]
for every $C\ge C_0$, since $U\ge0$.
\end{proof}

The constants $C_K,C_D,C_A,C_s,C_m,C_{HS}$ of the inputs and the constant $C_{\mathrm{unif}}$ of Hypothesis~\ref{w3:W3:unif} are the only constants in $C_0$, and none of them depends on $x$, on the point, on $j$ or on the level of the cover. The constants of the arithmetic Hilbert--Samuel theorem on a single model, which depend on the level, enter only inside the proof of Hypothesis~\ref{w3:W3:unif}, where the limit in the cutoff removes them.

\begin{remark}[the sign of the main term]\label{cov:sign}
The main term enters Hypothesis~\ref{cov:H58} with the sign $+\mathrm{main}_j$ on the upper bound for $\lcv(M_j)/D$, through $\mu(\Lambda^{\mathrm{core}}_j)\ge-\mathrm{main}_j-C_{\mathrm{low}}U$. The weights are $w_{j\rho}=e_j(\rho)$, the exponents of the archimedean parameters, and the sign of $\mathrm{main}_j$ is that of $\sum_\rho(p_j-q_j)w_{j\rho}\log\Omega$. On the cover the same term is $\sum_\tau(p-q)a_\tau\log\Omega$ for the anchor weight $a$, by (I13) and Lemma~\ref{cov:twist}, where a positive $a_\tau$ lowers $\mu$ by $(p-q)a_\tau\log\Omega$. That the anchor weight is the weight $w_{j+1}(v,s)=\varepsilon_{\tau_{v,s}}-c_j$ of the step character of Proposition~\ref{covariance}(iii) with this sign is condition (N1) of Section~\ref{sec:chain}, which this subsection does not prove.
\end{remark}

\paragraph{The local lattice relations at good primes.}
The exact identity at the good primes in the proof of Hypothesis~\ref{w3:W3:unif} uses, at every prime $p$ outside the bad set and every chart of the model over a discrete valuation ring $O$ containing $\zeta_p$ with $w=\zeta_p-1$, $w^{p-1}=u\,p$ for a unit $u$, the following relations. For lattices $E\subseteq E'$ with the same generic fibre let $\mathrm{ck}_k(E,E')$ be the degree on the special fibre of the class of $p_k(E')-p_k(E)$ against the fixed classes $L,Q$ with $Q$ homogeneous of degree $d+1-k$, and for an element $c\in O$ let $\rho_c(E)$ be the degree of $\mathrm{div}(c)\cdot p_{k-1}(E)\cdot Q$. Let $\mathrm{Lie}_\tau\subseteq W_\tau\subseteq w^{-1}\mathrm{Lie}_\tau$ and $\omega^{\mathrm{HT}}_\tau\subseteq\omega_\tau\subseteq w^{-1}\omega^{\mathrm{HT}}_\tau$ be the de Rham and Hodge--Tate lattices, let $X_{\tau,1},\dots,X_{\tau,n}$ be the character ideals of the Hodge--Tate map at the label $\tau$, $T$ the trivial line, $H_\tau$ the lattice dual to the annihilator lattice of the Hodge--Tate map and $T_n$ the standard lattice, and let $\varphi$ be the Frobenius permutation of the labels.

\begin{proposition}\label{cov:ck}
Let $p\ge2$ and $1\le k\le d+1$. Assume, for every label $\tau$,
\begin{enumerate}
\item[(G1)] $\mathrm{ck}_k(\mathrm{Lie}_\tau,W_\tau)+\mathrm{ck}_k(\omega_\tau,\omega^{\mathrm{HT}}_\tau)=\mathrm{ck}_k(T_n,H_\tau)-k\,\rho_w(\mathrm{Lie}_\tau)$,
\item[(G2)] $(p-1)\rho_w(\mathrm{Lie}_\tau)=\rho_p(\mathrm{Lie}_\tau)$,
\item[(G3)] $\mathrm{ck}_k(T_n,H_\tau)=(-1)^k\sum_{i=1}^n\mathrm{ck}_k(T,X_{\tau,i})$,
\item[(G4)] $k\,\rho_p(\mathrm{Lie}_{\varphi\tau})=(-1)^{k-1}\sum_{i=1}^n\big(\mathrm{ck}_k(T,X_{\varphi\tau,i})-p^k\,\mathrm{ck}_k(T,X_{\tau,i})\big)$.
\end{enumerate}
Put $h_\tau=(-1)^k\sum_i\mathrm{ck}_k(T,X_{\tau,i})$ and $\lambda_\sigma=k\,\rho_p(\mathrm{Lie}_\sigma)$. Then for every $\tau$
\[
\mathrm{ck}_k(\mathrm{Lie}_\tau,W_\tau)+\mathrm{ck}_k(\omega_\tau,\omega^{\mathrm{HT}}_\tau)=h_\tau-\frac{\lambda_\tau}{p-1},\qquad p^kh_\tau=h_{\varphi\tau}+\lambda_{\varphi\tau}.
\]
\end{proposition}

\begin{proof}
By (G1) and (G3) the left side of the first relation is $h_\tau-k\rho_w(\mathrm{Lie}_\tau)$, and by (G2) $k\rho_w(\mathrm{Lie}_\tau)=\lambda_\tau/(p-1)$. For the second relation put $A_i=\mathrm{ck}_k(T,X_{\varphi\tau,i})$ and $B_i=\mathrm{ck}_k(T,X_{\tau,i})$. Then (G4) reads $\lambda_{\varphi\tau}=(-1)^{k-1}\sum_i(A_i-p^kB_i)$, so $h_{\varphi\tau}+\lambda_{\varphi\tau}=(-1)^k\sum_iA_i-(-1)^k\sum_i(A_i-p^kB_i)=p^k(-1)^k\sum_iB_i=p^kh_\tau$.
\end{proof}

Proposition~\ref{cov:ck} identifies the numbers $h_\tau$ of the relations as the Chern characters of the character ideals of the Hodge--Tate map, so the relations are two exact computations. The first uses the exact sequence that compares $w^{-1}\mathrm{Lie}_\tau/W_\tau$, $H_\tau/T_n$ and $\omega_\tau/\omega^{\mathrm{HT}}_\tau$ together with $\mathrm{div}(w)=\frac1{p-1}\mathrm{div}(p)$. The second uses the fact that the step graph of the Hodge--Tate map at $\varphi\tau$ restricted to the first factor is the span of the $p$-th powers of the character ideal at $\tau$. The Frobenius recursion (Lemma \texttt{w3:U:frob} of the insertion that proves Hypothesis~\ref{w3:W3:unif}) then converts the two relations into (L2). \notgiven{Statements (G1)--(G4), the passage from test points to charts, and the refinement of the models in which they are computed are not proved here.}



\subsubsection{Uniformity in the level}\label{w3:sec:unif}

The mean slope bound on the covers must hold with a constant independent of the level of the cover, and the level grows with $x$. A uniform arithmetic Hilbert--Samuel theorem for the family of covers would give such a bound. No such theorem is published, and the argument does not use one.

\begin{hypothesis}[uniform mean slope on deep covers]\label{w3:W3:unif}
For an integer weight $u=(u_{\tau i})$ on a Shimura datum of signature
$(p,q)$, regarded in real coordinates, and a real number $g>0$, put
\[
  B(u,g)=4\,|\Phi_H|(p+q)^2
    \Bigl(1+\frac{\max_{\tau,i}|u_{\tau i}|}{g}\Bigr),\qquad
  \lVert u\rVert_1=\sum_{\tau,i}|u_{\tau i}|,
\]
\[
  \operatorname{spread}(u)=\sum_{\tau}
    \Bigl(\max_i u_{\tau i}-\min_i u_{\tau i}\Bigr).
\]
Say that $(u,g)$ has a gap when every $u_\tau$ is strictly decreasing and
\[
  g\leq |u_{\tau a}-u_{\tau b}|
  \qquad\text{for all $\tau$ and all $a\ne b$ with $p\leq a$}.
\]
Also put
\[
  \Lambda_{B(u,g)}=\log\bigl(B(u,g)\,\operatorname{rd}_E\bigr)
      +|\lambda_1|,
  \qquad
  \operatorname{err}(u,g)=\Lambda_{B(u,g)}\lVert u\rVert_1
      +q\,\operatorname{spread}(u).
\]
There is an absolute constant $C\geq0$ with the following property.  For
every admissible host, every chain datum satisfying its laws with $m^2\leq k$,
every $1\leq j<m$, and every good deep cover at level $j$ whose level has at
most fifteen prime factors, and for every pair $(u,g)$ with a gap, there are
auxiliary real numbers $\operatorname{sl}_Z$, indexed by the components of
the cover, such that
\[
  \operatorname{sl}_Z\leq \operatorname{sl}^{\,B(u,g)}_Z(u)
  \qquad\text{for every component $Z$},
\]
and
\[
  \frac1{\#\operatorname{Comp}}
    \sum_{Z\in\operatorname{Comp}}\operatorname{sl}_Z
  \geq -C\,\operatorname{err}(u,g).
\]
Here $\operatorname{Comp}$ denotes the set of components and
$\operatorname{sl}^{\,B}_Z(u)$ is the volume slope $\sigma_B(Z;u)$ defined in the covolume calculation, with cutoff $B$, weight $8ku$ and the common archimedean normalization $\xi_R(D)$.
\end{hypothesis}

The argument for Hypothesis~\ref{w3:W3:unif} uses a tower of models indexed by a prime cutoff $B$. For each $B$ one constructs a normal projective integral model $\mathcal Y_B$ of the toroidal compactification, smooth with good reduction away from the primes $\le B$ and a fixed bad set, with hermitian Hodge line bundles of every weight. On each model the arithmetic self-intersection $P_B(u)$ of the Hodge bundle of weight $u$ is a polynomial in $u$, and the arithmetic Hilbert--Samuel theorem of Gillet--Soul\'e, Zhang and Moriwaki, applied to the single model $\mathcal Y_B$, identifies the leading growth of the lattice of sections with $P_B(u)$. Zhang's inequality for successive minima forces $P_B$ to vanish on the Weyl walls up to $O(\log B/B)$, and an antisymmetry under the simple reflections puts it in the form $\mathrm{vol}\cdot(h_{\mathrm{Fl}}(u)-\alpha\,\mathrm{Sh}(u)-\beta\,\mathrm{Tot}(u))$ up to $O(\log B/B)$. An exact identity at the good primes relates $P_B$ and $P_{B'}$ prime by prime. The constants of the asymptotic steps depend on the model, hence on the level, but the exact identity allows $B\to\infty$, and in the limit these constants disappear. What remains is a point-height constant (a slope bound of Bost's type and the formula of Chowla and Selberg for $\mathbb Q(\sqrt{-d})$), the constant $4$ of a bound for cusp weights, the constant $4$ of the first Chern class, the constant $2$ of the tail over the primes above $B$, and the constant $60$ of a bound of Mertens type, none of which depends on the level. The published inputs are the arithmetic Hilbert--Samuel theorem for each fixed model, Zhang's theorem on successive minima, Hirzebruch--Mumford proportionality, and the existence of toroidal compactifications with their integral models. \notgiven{The construction of the tower, the exact identity at the good primes, and the limit $B\to\infty$ are not written out here, and Hypothesis~\ref{w3:W3:unif} is not proved here.} The uniformity in the level rests on the exactness of the identity at the good primes.

\subsubsection{Assembly of Hypothesis~\ref{r4:R4:HT}}

\begin{theorem}\label{w3:W3:HT}
Assume Proposition~\ref{w3:W3:passage}, Proposition~\ref{w3:W3:base}, Hypothesis~\ref{w3:W3:L3} for each $j$, and the bounds of Theorems~\ref{r5:thm:Ebound} and~\ref{r5:thm:DIVinN} for the local terms. Then Hypothesis~\ref{r4:R4:HT} holds with $\mu_0=\mu^{(0)}(R^\vee)$, $P_{\mathrm{pass}}=\mu^{(1)}-\mu^{(0)}$ and $E_b=E_{l3}+LG+\sum_j(E_{1,j}+E_{2,j})+E_3$.
\end{theorem}

\begin{proof}
Proposition~\ref{w3:W3:passage} bounds $P_{\mathrm{pass}}$. Proposition~\ref{w3:W3:base} bounds $\mu_0-\langle\log\kappa\rangle$ by $\sum_j\lcv(M_j)/D$ plus the local terms, which Theorem~\ref{r5:thm:DIVinN} and Theorem~\ref{r5:thm:Ebound} bound by $\mathcal L+LG+\sum_j(E_{1,j}+E_{2,j})+E_3$. Hypothesis~\ref{w3:W3:L3} bounds $\sum_j\lcv(M_j)/D$ by $\sum_{j,\rho}(p_j-q_j)w_{j\rho}\log\Omega+E_{l3}$ with $E_{l3}=C\sum_jU_j$.
\end{proof}



\section{The local terms and hypothesis~\texorpdfstring{\eqref{A5}}{(A5)}}\label{r5:sec:localterms}

This section defines the numbers $E_u$ of Proposition~\ref{acct} and proves hypothesis \eqref{A5} with
\[
g=\log(g_0g_1),\qquad e_5=C_E\,m\,B_0,\qquad B_0=f\,m\,b\,L_\varepsilon=2\,n\,m\,b\,x^{\varepsilon},
\]
where $g_0,g_1$ are the two primes of (D5) and $C_E$ is a constant that depends on $\varepsilon$ alone. The number $E_u$ is a sum over the finite places of the field of the lattices, one term for each lift, and each term measures the failure of integrality of one lift at one place. Each lift loses the amount $A\log g_i$ at the places above the level prime $g_i$, so $E_u$ has the main term $mA\log(g_0g_1)$, and Proposition~\ref{acct} cancels it against the term $mg$ of \eqref{A6}. All other contributions to $E_u$ are bounded by a multiple of $mB_0$, which is the second term of $\mathcal B$ in \eqref{calB}.

Section~\ref{sec:twisted} treats the places at which the local type of the lifts is twisted. There a second local quantity, the congruence defect $\eta_j(w)$ of the lattice of integral forms, is bounded by the lengths of Tate--Shafarevich groups of $E$ through a congruence-module estimate, and the lengths are bounded by $\sqrt b$. That bound enters the heights of the period maps in \eqref{A4} and the slope inequality \eqref{A3}. It does not enter \eqref{A5}.

\subsection{Places and normalized lengths}\label{r5:sec:places}

Let $L$ be the number field over which the hermitian lattices $N_u$, $M_u$ and $R$ of Part~II are defined. It is a Galois extension of $\mathbb Q$ that contains $E$, and in particular it contains $K=\mathbb Q(\sqrt{-d})$. For a finite place $w$ of $L$ write $p_w$ for its residue characteristic, $e_w$ and $f_w$ for its ramification index and residue degree over $\mathbb Q$, $v_w$ for the normalized valuation with $v_w(L^\times)=\mathbb Z$, and
\[
\ell_w=\frac{f_w\log p_w}{[L:\mathbb Q]},\qquad\sigma(w,y)=2\,\ell_w\,e_w\,y\qquad(y\in\mathbb R).
\]
Here $y$ is in units with $v_p(p)=1$. If $e_wy$ is integral, rescaling the lattice at $w$ by $\varpi_w^{-e_wy}$ changes its doubled normalized slope by $\sigma(w,y)$. For arbitrary $y$ the actual lattice uses $\lfloor e_wy\rfloor$ and Section~\ref{r2:sec:Nu} bounds the rounding error. Each finite place in the sum occurs once.

\begin{lemma}\label{r5:lem:prodformula}
Let $p$ be a prime number.
\begin{enumerate}
\item[(i)] $\sum_{w\mid p}\ell_we_w=\log p$.
\item[(ii)] For $y\in L^\times$, $\sum_{w\mid p}\ell_w\,v_w(y)=\dfrac{v_p\big(N_{L/\mathbb Q}(y)\big)\log p}{[L:\mathbb Q]}$.
\end{enumerate}
\end{lemma}

\begin{proof}
(i) The places above $p$ satisfy $\sum_{w\mid p}e_wf_w=[L:\mathbb Q]$, so $\sum_{w\mid p}\ell_we_w=\log p\sum_{w\mid p}e_wf_w/[L:\mathbb Q]=\log p$.

(ii) The fractional ideal $y\mathcal O_L$ factors as $\prod_w\mathfrak p_w^{v_w(y)}$, and the absolute norm of $\mathfrak p_w$ is $p_w^{f_w}$. Hence $\lvert N_{L/\mathbb Q}(y)\rvert=\prod_wp_w^{f_wv_w(y)}$, and the exponent of $p$ in this product is $v_p(N_{L/\mathbb Q}(y))=\sum_{w\mid p}f_wv_w(y)$. Multiplying by $\log p/[L:\mathbb Q]$ gives (ii).
\end{proof}

\begin{lemma}\label{r5:lem:halfnorm}
Let $g$ be a prime number and $\pi\in\mathcal O_K$ with $\pi\bar\pi=g$. For every automorphism $\tau$ of $L$,
\[
\sum_{w\mid g}\ell_w\,v_w\big(\tau(\bar\pi)\big)=\frac12\log g .
\]
\end{lemma}

\begin{proof}
Since $L/\mathbb Q$ is Galois, $N_{L/\mathbb Q}(\tau(\bar\pi))=N_{L/\mathbb Q}(\bar\pi)$, and since $\bar\pi\in K$, transitivity of the norm gives $N_{L/\mathbb Q}(\bar\pi)=N_{K/\mathbb Q}(\bar\pi)^{[L:K]}=(\bar\pi\pi)^{[L:K]}=g^{[L:K]}$. Lemma~\ref{r5:lem:prodformula}(ii) with $p=g$ gives $[L:K]\log g/[L:\mathbb Q]=\frac12\log g$.
\end{proof}

\begin{lemma}\label{r5:lem:vallog}
Let $n\ge1$ be an integer and $P$ a finite set of primes. Then $\sum_{p\in P}v_p(n)\log p\le\log n$.
\end{lemma}

\begin{proof}
The integer $\prod_{p\in P}p^{v_p(n)}$ divides $n$, so it is at most $n$, and its logarithm is the sum on the left.
\end{proof}

\begin{lemma}\label{r5:lem:logsum}
$\displaystyle\sum_{p}\frac{\log p}{p^2-1}<1$, the sum running over all primes.
\end{lemma}

\begin{proof}
The terms with $p\le13$ are $0.2310$, $0.1374$, $0.0671$, $0.0405$, $0.0200$ and $0.0152$, each rounded up in the fourth decimal, with sum less than $0.512$. For $t\ge e^{1/2}$ the function $\log t/t^2$ is decreasing, and for an integer $n\ge17$ we have $n^2-1\ge\frac{255}{256}n^2$. Hence
\[
\sum_{p\ge17}\frac{\log p}{p^2-1}\le\frac{256}{255}\sum_{n\ge17}\frac{\log n}{n^2}\le\frac{256}{255}\int_{16}^\infty\frac{\log t}{t^2}\,dt=\frac{256}{255}\cdot\frac{\log16+1}{16}<0.237 .
\]
The total is less than $0.75$.
\end{proof}

\subsection{The local exponent data}\label{r5:sec:localdata}

Fix admissible data at the scale $x$ (Definition~\ref{def:admissible}) and $u_0\in A_0$. Part~II attaches to them the following objects. Each item names the section in which the object is constructed.
\begin{enumerate}
\item[(L1)] The finite set $W$ of indices $u$ of Proposition~\ref{acct}, and for each $u\in W$ a prime $T_u\notin\{g_0,g_1,d\}$ with $T_u=\varpi_u\bar\varpi_u$ for some $\varpi_u\in\mathcal O_K$ (Section~\ref{sec:proofA}).
\item[(L2)] The elements $\pi_0,\pi_1\in\mathcal O_K$ of (D5), with $\pi_i\bar\pi_i=g_i$.
\item[(L3)] A positive integer $D_\Gamma$, a common multiple of the orders of the stabilizers in the class sums that define the lifts $K_1,\dots,K_m$ (Section~\ref{sec:lifts}, condition (D1)).
\item[(L4)] For $1\le j\le m$ and each element $\rho$ of the CM type $\Phi$, the node $w_j(\rho)=w_{u_0,j}(\rho)>0$ of Theorem~A (Section~\ref{sec:proofA}). For $u\in W$, $1\le j\le m$ and $\rho\in\Phi$ a real number $n_{u,j}(\rho)$, the exponent of $\bar\pi_i$ in the $j$-th lift at $\rho$, which is the same for $i=0$ and $i=1$. Here $\Phi\subset\mathcal G=\mathrm{Gal}(E/\mathbb Q)$ as in Section~\ref{sec:excess}, and we write $\tau_\rho$ for an automorphism of $L$ whose restriction to $E$ is $\rho$. The valuations $v_w(\tau_\rho(y))$ for $y\in E$ below do not depend on the choice of $\tau_\rho$. A real number $\Delta_j\ge0$, the Fock degree of the polynomial model of the weight of level $j$, and a real number $\kappa_j\ge0$ (Section~\ref{sec:lifts}).
\item[(L5)] Real numbers $\varepsilon^T_u\ge0$ and $\varepsilon^d_j$, and for each prime $p$ in a finite set $P^{\rm in}$ of primes $p\le X''$ with $X''=x\,n\,X'$ a real number $\varepsilon^{\rm in}_{j,p}$ (Sections~\ref{r2:sec:loss} and~\ref{sec:proofA}). The set $P^{\rm in}$ consists of the primes at which Section~\ref{sec:gain} and Section~\ref{r2:sec:loss} attach a gain or a loss to the lifts.
\end{enumerate}
For $u\in W$ let $S_u$ be the set of places $w\mid T_u$ of $L$ such that $v_w(\tau_\rho(\varpi_u))>0$ for every $\rho\in\Phi$.

\begin{definition}\label{r5:def:Eterms}
Let $u\in W$, $1\le j\le m$, and let $w$ be a finite place of $L$ with residue characteristic $p=p_w$. Define $E_{u,j}(w)$ by the first applicable line.
\[
E_{u,j}(w)=\begin{cases}
\varepsilon^T_u&\text{if }p=T_u\text{ and }w\in S_u,\\
0&\text{if }p=T_u\text{ and }w\notin S_u,\\
v_p(D_\Gamma)-\dfrac{\Delta_j}2+\kappa_j+\dfrac1{e_w}\displaystyle\sum_{\rho\in\Phi}n_{u,j}(\rho)\,v_w\big(\tau_\rho(\bar\pi_i)\big)&\text{if }p=g_i,\ i\in\{0,1\},\\
v_p(D_\Gamma)+\varepsilon^d_j&\text{if }p=d,\\
v_p(D_\Gamma)+\varepsilon^{\rm in}_{j,p}&\text{if }p\in P^{\rm in},\\
v_p(D_\Gamma)&\text{otherwise}.
\end{cases}
\]
Put
\[
E_{u,j}=\sum_w\sigma\big(w,E_{u,j}(w)\big),\qquad E_u=\sum_{j=1}^mE_{u,j}.
\]
\end{definition}

The sum over $w$ is finite, because $E_{u,j}(w)=0$ unless $p_w$ lies in the finite set $\{T_u,g_0,g_1,d\}\cup P^{\rm in}\cup\{p:p\mid D_\Gamma\}$. The same numbers $E_{u,j}(w)$ enter the lattice $N_u$ in hypothesis \eqref{A2} and the integrality of the period maps in \eqref{A3}. Section~\ref{r5:sec:R5interface} states the form in which those hypotheses use them.

For the main term put $n_{u,j}=\sum_{\rho\in\Phi}n_{u,j}(\rho)$ and $A=\sum_{t\in T}a_t$. The bound for $E_u$ uses the following six inequalities, in which $C_g,C_T,C_\Gamma,C_d,C_{\rm in}\ge0$ are constants and $\delta_j\ge0$.
\begin{enumerate}
\item[(E1)] $\lvert n_{u,j}-\Delta_j-A\rvert\le\delta_j$.
\item[(E2)] $(\delta_j+2\kappa_j)\log(g_0g_1)\le C_g\,B_0$.
\item[(E3)] $2\,\varepsilon^T_u\log T_u\le C_T\,B_0$.
\item[(E4)] $2\log D_\Gamma\le C_\Gamma\,B_0$.
\item[(E5)] $2\log d\,\lvert\varepsilon^d_j\rvert\le C_d\,B_0$.
\item[(E6)] $\sum_{p\in P^{\rm in}}2\log p\,\lvert\varepsilon^{\rm in}_{j,p}\rvert\le C_{\rm in}\,B_0$.
\end{enumerate}
Inequality (E1) is the degree identity of the lifts at the level primes. Subsection~\ref{r5:sec:Hsizes} derives (E2) to (E5) from size bounds for the data and proves (E6) for the part of $P^{\rm in}$ whose loss term has the closed form \eqref{r5:eq:EatC}.

\subsection{The bound for the local terms}\label{r5:sec:EAbound}

\begin{theorem}\label{r5:thm:Ebound}
Let $u\in W$ and $1\le j\le m$, let the data (L1) to (L5) be given, and suppose that $g_0\ne g_1$ and that the inequalities (E1) to (E6) hold for $u$ and $j$. Then
\[
\big\lvert E_{u,j}-A\log(g_0g_1)\big\rvert\le C_E\,B_0,\qquad C_E=C_g+C_T+C_\Gamma+C_d+C_{\rm in}.
\]
\end{theorem}

\begin{proof}
We split the sum $E_{u,j}=\sum_w\sigma(w,E_{u,j}(w))$ into the places above $g_0$, above $g_1$, above $T_u$, and the remaining places. The four sets of places are disjoint because $T_u$, $g_0$ and $g_1$ are distinct primes.

\emph{The places above $g_i$.} Fix $i$ and put $c_i=v_{g_i}(D_\Gamma)-\frac{\Delta_j}2+\kappa_j$. For $w\mid g_i$ the definition gives
\[
\sigma\big(w,E_{u,j}(w)\big)=2\ell_we_wc_i+2\ell_w\sum_{\rho\in\Phi}n_{u,j}(\rho)\,v_w\big(\tau_\rho(\bar\pi_i)\big).
\]
Summing over $w\mid g_i$, Lemma~\ref{r5:lem:prodformula}(i) turns the first term into $2c_i\log g_i$, and Lemma~\ref{r5:lem:halfnorm} turns the second into $\sum_\rho n_{u,j}(\rho)\log g_i=n_{u,j}\log g_i$. Hence
\begin{align*}
\sum_{w\mid g_i}\sigma\big(w,E_{u,j}(w)\big)&=\big(2v_{g_i}(D_\Gamma)-\Delta_j+2\kappa_j+n_{u,j}\big)\log g_i\\
&=A\log g_i+\big(n_{u,j}-\Delta_j-A+2\kappa_j+2v_{g_i}(D_\Gamma)\big)\log g_i .
\end{align*}
By (E1) the absolute value of the bracket minus $2v_{g_i}(D_\Gamma)$ is at most $\delta_j+2\kappa_j$, so the sum over the places above $g_0$ and $g_1$ differs from $A\log(g_0g_1)$ by at most
\[
(\delta_j+2\kappa_j)\log(g_0g_1)+2v_{g_0}(D_\Gamma)\log g_0+2v_{g_1}(D_\Gamma)\log g_1 .
\]

\emph{The places above $T_u$.} Here $E_{u,j}(w)=\varepsilon^T_u\ge0$ for $w\in S_u$ and $E_{u,j}(w)=0$ otherwise. By Lemma~\ref{r5:lem:prodformula}(i),
\[
0\le\sum_{w\mid T_u}\sigma\big(w,E_{u,j}(w)\big)=2\varepsilon^T_u\sum_{w\in S_u}\ell_we_w\le2\varepsilon^T_u\log T_u .
\]

\emph{The remaining places.} Let $p\notin\{T_u,g_0,g_1\}$. For $w\mid p$ the definition gives $E_{u,j}(w)=v_p(D_\Gamma)+\lambda_p$, where $\lambda_d=\varepsilon^d_j$, $\lambda_p=\varepsilon^{\rm in}_{j,p}$ for $p\in P^{\rm in}\setminus\{d\}$, and $\lambda_p=0$ otherwise. The value does not depend on $w\mid p$, so Lemma~\ref{r5:lem:prodformula}(i) gives $\sum_{w\mid p}\sigma(w,E_{u,j}(w))=2(v_p(D_\Gamma)+\lambda_p)\log p$. Summing over these primes,
\[
\Big\lvert\sum_{p_w\notin\{T_u,g_0,g_1\}}\sigma\big(w,E_{u,j}(w)\big)\Big\rvert\le2\sum_{p\notin\{T_u,g_0,g_1\}}v_p(D_\Gamma)\log p+2\log d\,\lvert\varepsilon^d_j\rvert+\sum_{p\in P^{\rm in}}2\log p\,\lvert\varepsilon^{\rm in}_{j,p}\rvert .
\]

\emph{Conclusion.} The terms $2v_p(D_\Gamma)\log p$ occur above for $p=g_0$, $p=g_1$ and for the primes outside $\{T_u,g_0,g_1\}$, each prime once, so by Lemma~\ref{r5:lem:vallog} their total is at most $2\log D_\Gamma$. Adding the three estimates,
\begin{multline*}
\big\lvert E_{u,j}-A\log(g_0g_1)\big\rvert\le(\delta_j+2\kappa_j)\log(g_0g_1)+2\varepsilon^T_u\log T_u+2\log D_\Gamma\\
+2\log d\,\lvert\varepsilon^d_j\rvert+\sum_{p\in P^{\rm in}}2\log p\,\lvert\varepsilon^{\rm in}_{j,p}\rvert,
\end{multline*}
and (E2) to (E6) bound the five terms by $C_gB_0$, $C_TB_0$, $C_\Gamma B_0$, $C_dB_0$ and $C_{\rm in}B_0$.
\end{proof}

\begin{corollary}\label{r5:cor:A5}
Suppose that the hypotheses of Theorem~\ref{r5:thm:Ebound} hold for every $u\in W$ and every $1\le j\le m$. Put $g=\log(g_0g_1)$ and $e_5=C_E\,m\,B_0$. Then for every $u\in W$
\[
mAg-e_5\le E_u\le mAg+e_5 ,
\]
so hypothesis \eqref{A5} of Proposition~\ref{acct} holds. Moreover $e_5\le C_E\,\mathcal B$, with $\mathcal B$ the number \eqref{calB}.
\end{corollary}

\begin{proof}
Summing Theorem~\ref{r5:thm:Ebound} over $1\le j\le m$ gives $\lvert E_u-mA\log(g_0g_1)\rvert\le mC_EB_0=e_5$. The second term of \eqref{calB} is $m\cdot2nmbx^\varepsilon=mB_0$, and the other terms of \eqref{calB} are nonnegative, so $mB_0\le\mathcal B$.
\end{proof}

\begin{remark}\label{r5:rem:gcancel}
The number $g$ of Proposition~\ref{acct} occurs in \eqref{A5} and in \eqref{A6}, and the proof of Proposition~\ref{acct} cancels the two occurrences. Corollary~\ref{r5:cor:A5} fixes $g=\log(g_0g_1)$. Hypothesis \eqref{A6} must be proved with the same number, that is, in the form
\[
\log T_u\ge\sum_{j=1}^m\frac1f\sum_{\rho\in\Phi}\log\frac{w_j(\rho)v}\pi+m\log(g_0g_1)-e_6 .
\]
This is consistent with the choice of $T_u$ as a prime of size about $(g_0g_1)^m$ times the product of the nodes, and with the congruence $\varpi_u\equiv1$ modulo $Q_Sqd\,(g_0g_1)^2$ of Section~\ref{sec:proofA}.
\end{remark}

\begin{remark}\label{r5:rem:Cuniform}
The constant $C_E$ depends on $\varepsilon$ through $C_{\rm in}$ (Subsection~\ref{r5:sec:Hsizes}). Theorem~A asks for an absolute constant $C$ in $x^{C\varepsilon}\mathcal B$, and Corollary~\ref{r5:cor:A5} gives $e_5\le C_E\mathcal B\le x^{\varepsilon}\mathcal B$ for every $x$ with $x^\varepsilon\ge C_E$. The threshold $Y_1$ of Theorem~A may depend on $\varepsilon$, so this contributes $1$ to the absolute constant $C$.
\end{remark}

\subsection{The inequalities (E2) to (E6) from the sizes}\label{r5:sec:Hsizes}

We use the scale \eqref{scale}, $L_\varepsilon=x^\varepsilon$, $b=e^{L_\varepsilon}$, $a=xb$, $X'=\exp(x^{\varepsilon/3})$, $b'=a/X'$, the ranks of Proposition~\ref{ranks}, and
\[
N_m=k+m(t_r+1),
\]
the dimension of the polynomial model of the weight at the last level of the chain.

\begin{lemma}\label{r5:lem:primesizes}
Let $0<\varepsilon<\frac1{52}$. There is $x_0$ such that for admissible data at a scale $x\ge x_0$
\begin{enumerate}
\item[(i)] $\log g_i\le9L_\varepsilon$ for $i=0,1$, so $\log(g_0g_1)\le18L_\varepsilon$;
\item[(ii)] $\log d\le L_\varepsilon$;
\item[(iii)] $t_r\le x^{401}$ and $N_m\le x^{402}$;
\item[(iv)] $N_m^2\,X''\log a\le b$, where $X''=x\,n\,X'$.
\end{enumerate}
\end{lemma}

\begin{proof}
(i) By (D5), $g_i\le2(kb)^8$, so $\log g_i\le\log2+8\log k+8L_\varepsilon$. Proposition~\ref{ranks} gives $k\le2x^{9/2+10\varepsilon}\le2x^5$, so $\log g_i\le9+40\log x+8L_\varepsilon\le9L_\varepsilon$ once $x^\varepsilon\ge9+40\log x$.

(ii) By (D1) the prime $d$ satisfies the conclusion of Theorem~\ref{stage2} with $K=\frac32$ and $z'=z\log^2Y$, $z=Y^{\varepsilon/5}$, so $\log d\le z'=Y^{\varepsilon/5}\log^2Y$. Since $Y\le x$, this is at most $x^{\varepsilon/5}\log^2x\le x^\varepsilon$ for large $x$.

(iii) Since $\varepsilon<\frac1{52}$, Proposition~\ref{ranks} and (D1) give $f=2n\le4x^{2+8\varepsilon}\le4x^{2.2}$, $1\le m\le2\sqrt x$, $k\le2x^{9/2+10\varepsilon}\le2x^{4.7}$, $r=k+m-1\le3x^{4.7}$, and $L_\varepsilon^2=x^{2\varepsilon}\le x^{0.04}$. The definition of $n_W$ in Section~\ref{sec:data} and $m\ge1$ give
\[
n_W\le10\,r\big(f+f^2k^3L_\varepsilon^2\big)+8\le10\cdot3x^{4.7}\cdot2\cdot16x^{4.4}\cdot8x^{14.1}\cdot x^{0.04}+8\le x^{23.5}
\]
for large $x$. With $T_O=\lceil x^{377}\rceil\le2x^{377}$ and $t_r=8fr+4+T_On_W-r$ we get $t_r\le8fr+4+2x^{400.5}\le x^{401}$ for large $x$. Then $N_m=k+m(t_r+1)\le2x^{4.7}+2\sqrt x\,(x^{401}+1)\le x^{402}$ for large $x$.

(iv) We have $n\le2x^{3}$, $\log a=\log x+L_\varepsilon\le2L_\varepsilon$, and by (iii)
\[
N_m^2X''\log a\le x^{804}\cdot x\cdot2x^3\cdot\exp(x^{\varepsilon/3})\cdot2x^\varepsilon\le\exp(x^\varepsilon)=b
\]
for large $x$, because $\log$ of the left side is $x^{\varepsilon/3}+O(\log x)$.
\end{proof}

\begin{lemma}\label{r5:lem:H2toH5}
Let $0<\varepsilon<\frac1{52}$ and let $x_0$ be as in Lemma~\ref{r5:lem:primesizes}. Let $c_1,\dots,c_5\ge0$ be constants, and suppose that for admissible data at a scale $x\ge x_0$, for $u\in W$ and for $1\le j\le m$,
\begin{gather}
\delta_j\le c_1\,f\,b,\qquad\kappa_j\le c_2\,f\,m\,b,\label{r5:Z1}\\
\varepsilon^T_u\le c_3\,f\,b,\qquad\log T_u\le c_4\,m\,L_\varepsilon,\label{r5:Z2}\\
\log D_\Gamma\le c_5\,f\,N_m\log(fN_m),\qquad\lvert\varepsilon^d_j\rvert\le2\,f\,b.\label{r5:Z3}
\end{gather}
Then (E2) holds with $C_g=18(c_1+2c_2)$, (E3) with $C_T=2c_3c_4$, (E4) with $C_\Gamma=2c_5$, and (E5) with $C_d=4$, for $x\ge x_0$ enlarged if necessary.
\end{lemma}

\begin{proof}
Recall $B_0=fmbL_\varepsilon$ and $m\ge2$.

(E2) By Lemma~\ref{r5:lem:primesizes}(i), $(\delta_j+2\kappa_j)\log(g_0g_1)\le(c_1fb+2c_2fmb)\cdot18L_\varepsilon\le18(c_1+2c_2)B_0$.

(E3) $2\varepsilon^T_u\log T_u\le2c_3fb\cdot c_4mL_\varepsilon=2c_3c_4B_0$.

(E4) By Lemma~\ref{r5:lem:primesizes}(iii) and $f\le4x^3$, $\log(fN_m)\le\log(4x^{405})\le L_\varepsilon$ and $N_m\le x^{402}\le b$ for large $x$. Hence $2\log D_\Gamma\le2c_5fbL_\varepsilon\le2c_5B_0$.

(E5) By Lemma~\ref{r5:lem:primesizes}(ii), $2\log d\,\lvert\varepsilon^d_j\rvert\le2L_\varepsilon\cdot2fb\le4B_0$.
\end{proof}

The loss term at a prime $p\in P^{\rm in}$ has one of two forms. At the primes at which the lifts gain by the class of the element of Section~\ref{sec:gain} that is trivial, it has the closed form
\begin{equation}\label{r5:eq:EatC}
\varepsilon^{\rm in}_{j,p}=c_{\rm dig}\,f\,N_m^2\,\log_pa+\frac1{p^2-1}\sum_{\rho\in\Phi}w_j(\rho)
\end{equation}
with an absolute constant $c_{\rm dig}$. At the other primes of $P^{\rm in}$ it is the orbit loss of Section~\ref{r2:sec:loss}, built from the solution of the linear programme of Section~\ref{sec:gain} along the Frobenius orbits. Write $P^{\rm in}=P^{\rm C}\sqcup P^{\rm I}$ accordingly.

\begin{lemma}\label{r5:lem:H6C}
Suppose that $0\le w_j(\rho)\le3b$ for every $\rho\in\Phi$, and that $x\ge x_0$ with $x_0$ as in Lemma~\ref{r5:lem:primesizes}. If $\varepsilon^{\rm in}_{j,p}$ has the form \eqref{r5:eq:EatC} for $p\in P^{\rm C}$, then
\[
\sum_{p\in P^{\rm C}}2\log p\,\lvert\varepsilon^{\rm in}_{j,p}\rvert\le(2c_{\rm dig}+6)\,f\,b\le(c_{\rm dig}+3)\,B_0 .
\]
\end{lemma}

\begin{proof}
Both terms of \eqref{r5:eq:EatC} are nonnegative. Since $\log p\cdot\log_pa=\log a$ and $P^{\rm C}$ consists of primes $p\le X''$, the first terms contribute at most $2c_{\rm dig}fN_m^2\log a\cdot X''\le2c_{\rm dig}fb$ by Lemma~\ref{r5:lem:primesizes}(iv). Since $\sum_\rho w_j(\rho)\le3fb$, Lemma~\ref{r5:lem:logsum} bounds the second terms by $2\cdot3fb\sum_p\log p/(p^2-1)\le6fb$. Finally $fb\le\frac12fmbL_\varepsilon=\frac12B_0$.
\end{proof}

\begin{hypothesis}[inert orbit losses]\label{r5:prop:H6I}
Let $0<\varepsilon<\frac1{52}$. There are $C'_{\rm in}$ and $x_0$ such that for admissible data at a scale $x\ge x_0$ and every $1\le j\le m$, $\sum_{p\in P^{\rm I}}2\log p\,\lvert\varepsilon^{\rm in}_{j,p}\rvert\le C'_{\rm in}B_0$.
\end{hypothesis}

The orbit losses $\varepsilon^{\rm in}_{j,p}$ for $p\in P^{\rm I}$ come from the construction of Section~\ref{r2:sec:loss}. \notgiven{Hypothesis~\ref{r5:prop:H6I} is not proved here.}

With Lemma~\ref{r5:lem:H6C} and Hypothesis~\ref{r5:prop:H6I}, (E6) holds with $C_{\rm in}=c_{\rm dig}+3+C'_{\rm in}$.

\begin{hypothesis}[exponents and sizes of the lifts]\label{r5:prop:H1}
Let $0<\varepsilon<\frac1{52}$. There are constants $c_1,\dots,c_5$ and $x_0$ such that for admissible data at a scale $x\ge x_0$, the exponents and sizes of the lifts satisfy (E1) and \eqref{r5:Z1} to \eqref{r5:Z3} for every $u\in W$ and $1\le j\le m$, and $0\le w_j(\rho)\le3b$ for every $\rho\in\Phi$.
\end{hypothesis}

The degree identity (E1) is a computation with the highest weights of the lifts of Section~\ref{sec:lifts}, and \eqref{r5:Z1} to \eqref{r5:Z3} are size bounds for the data of Section~\ref{sec:proofA}. \notgiven{Hypothesis~\ref{r5:prop:H1} is not proved here.}

\subsection{How hypotheses \texorpdfstring{\eqref{A2}}{(A2)} and \texorpdfstring{\eqref{A3}}{(A3)} use the local terms}\label{r5:sec:R5interface}

The numbers $E_{u,j}(w)$ of Definition~\ref{r5:def:Eterms} are bounds for the denominators of the lifts. With the gain $\mathrm{gain}_j(w)$ and the local terms $F_j(w)$ of Sections~\ref{sec:lifts} and~\ref{r2:sec:loss}, the lattice $N_u$ of \eqref{A2} is the rescaling at each place $w$ with $p_w\ne T_u$ of a sublattice of the base lattice by the exponent
\[
s_u(w)=\sum_{j=1}^m\big(\mathrm{gain}_j(w)+F_j(w)-E_{u,j}(w)\big),
\]
so that $\mu^N_u$ differs from the slope of the sublattice by $\sum_w\sigma(w,s_u(w))$, which equals $\Gamma+\mathcal L-E_u$ up to the terms $\tau(w)$ of the loss and the rounding to integers. This is why $E_u$ enters \eqref{A2} with a minus sign. Hypothesis \eqref{A3} needs, at each place $w$ with $p_w\ne T_u$, that the period of a lifted form on the rescaled lattice is integral at $w$. The argument uses this in the following form. For $\varphi$ in the base lattice and $F$ an integral form of the local type $\mu_w$ of the lifts at $w$, the period $y\in L$ of the lift of $\varphi$ against $F$ is $0$ or satisfies
\[
v_w(y)\ge e_w\sum_{j=1}^m\big(F_j(w)-E_{u,j}(w)\big).
\]
The proof of this integrality belongs to Sections~\ref{r2:sec:loss} and~\ref{r3:sec}, and they use Definition~\ref{r5:def:Eterms} for $E_{u,j}(w)$.


\section{The twisted places}\label{sec:twisted}

At the places where the local type of the lifts is twisted, the lattice of integral lifted forms differs from the lattice that the period pairing sees, and a congruence module measures the difference. This section defines the congruence defects $\eta_j(w)$, states the two bounds that control them, and proves the parts of those bounds that reduce to commutative algebra and Galois cohomology. Theorem~\ref{r5:thm:DIVinN} combines the two bounds. Its conclusion enters the heights of the period maps in \eqref{A4} and the frame parameters in \eqref{A3}. It does not enter \eqref{A5}, which Corollary~\ref{r5:cor:A5} proves without it.

The first bound, Hypothesis~\ref{w4:thm:vwin}, bounds the lengths of Tate--Shafarevich groups of the host field on windows of scales. Subsection~\ref{tw:sec:vandiver} shows that it implies a bound in the direction of Vandiver's conjecture beyond the bounds published by Soul\'e~\cite{soule99} and by de~la~Torre~Dur\'an~\cite{delatorre}, and traces the argument for it to an automorphy statement for residually reducible Galois representations that is not a published theorem.

\subsection{Congruence defects}\label{tw:sec:eta}

For each finite place $w$ of $L$ the lifts have a local type $\mu_w$, which is either the standard type $\mu_0$ or a twisted type $\mu^{\rm tw}_w$, and $\mu_{\bar w}=\mu_w$. We call $w$ \emph{twisted} when $\mu_w\ne\mu_0$. At a twisted place the integral structure of the twisted lifts at level $j$ differs from the integral structure of the eigen-components that the period pairing sees, and the difference is measured by a congruence module.

\begin{definition}\label{r5:def:eta}
Let $w$ be a twisted place, $l=p_w$, and $1\le j\le m$. Let $\mathcal O$ be the ring of integers of the completion of $L$ at $w$, let $\Lambda$ be the $\mathcal O$-lattice spanned by the integral forms of level $j$ and local type $\mu^{\rm tw}_w$, and let $\mathbb T$ be the $\mathcal O$-algebra of Hecke operators acting on $\Lambda$. For a maximal ideal $\mathfrak m$ of $\mathbb T$ split $\Lambda_{\mathfrak m}\otimes\mathbb Q$ into the sum $V_{\rm good}\oplus V_{\rm bad}$ of the eigen-components whose eigensystems occur in the lifts of level $j$ and of the others. Put
\[
\eta_j(w)=\sum_{\mathfrak m}\frac{\operatorname{length}_{\mathcal O}\Big(\big(V_{\rm good}\cap(\Lambda_{\mathfrak m}+V_{\rm bad})\big)\big/\big(\Lambda_{\mathfrak m}\cap V_{\rm good}\big)\Big)}{\operatorname{length}_{\mathcal O}(\mathcal O/l)} ,
\]
the normalized length of the congruence module of $V_{\rm good}$ in $\Lambda$. For a non-twisted place put $\eta_j(w)=0$.
\end{definition}

The module $(V_{\rm good}\cap(\Lambda+V_{\rm bad}))/(\Lambda\cap V_{\rm good})$ is the projection of $\Lambda$ to $V_{\rm good}$ modulo the forms of $\Lambda$ that lie in $V_{\rm good}$. It vanishes when no good eigensystem is congruent to a bad one modulo $\mathfrak m$. For $1\le j\le m$ we write $d^0_j$ for the rank of the space of lifted forms of level $j$ and standard type $\mu_0$, and $N_c=\lvert T_2\rvert\,k\,b\,x^{2\varepsilon}/n^2$, so that $\mathcal B=N_Ta(1+m^2/k)+m(B_0+N_c)$.

\subsection{Tate--Shafarevich groups}\label{tw:sec:sha}

Let $F$ be a number field, $\overline F$ an algebraic closure, $G_F=\Gal(\overline F/F)$ with the Krull topology, and $l$ a prime. Let $\chi\colon G_F\to\Zl^\times$ be the $l$-adic cyclotomic character, so that $\sigma(\zeta)=\zeta^{\chi(\sigma)}$ for every $l$-power root of unity $\zeta$. For an integer $e$ let $M_e$ be the discrete group $\Ql/\Zl$ on which $\sigma$ acts by multiplication by $\chi(\sigma)^{-e}$. For a finite place $v$ of $F$ fix a place of $\overline F$ above it, with decomposition group $D_v$ and inertia group $I_v$. Cohomology is continuous cohomology with discrete coefficients.

\begin{definition}\label{w4:def:sha}
$\Sha_l(F,e)$ is the subgroup of $H^1(G_F,M_e)$ of classes $c$ with $c|_{D_v}=0$ for every $v\mid l$ and $c|_{I_v}=0$ for every finite $v$. For a finite free $\Zl$-algebra $O$ put $M_{e,O}=O\otimes_{\Zl}M_e$, with $G_F$ acting on the second factor, and define $\Sha_l(F,e;O)\subset H^1(G_F,M_{e,O})$ by the same conditions. For a finite abelian $l$-group $X$ write $\lambda(X)=\log|X|=\len_{\Zl}(X)\log l$.
\end{definition}

The conditions do not depend on the chosen places above $v$, because restriction to conjugate subgroups differs by an automorphism of cohomology. There is no condition at archimedean places. For odd $l$ none is needed, because $H^1$ of a group of order $2$ with coefficients in an $l$-primary group vanishes.

\begin{lemma}\label{w4:lem:unram}
Let $v\nmid l$ be a finite place of $F$ and $e\ge1$. An element of $H^1(D_v,M_e)$ that vanishes on $I_v$ is zero.
\end{lemma}

\begin{proof}
The group $I_v$ acts trivially on $M_e$, because $\chi$ is unramified at $v$. Inflation--restriction gives an exact sequence $0\to H^1(D_v/I_v,M_e)\to H^1(D_v,M_e)\to H^1(I_v,M_e)$, so a class that vanishes on $I_v$ comes from $H^1(D_v/I_v,M_e)$. The group $D_v/I_v$ is procyclic, topologically generated by a Frobenius element $\varphi$, and $M_e$ is a discrete torsion module, so $H^1(D_v/I_v,M_e)\cong M_e/(\varphi-1)M_e$. The element $\varphi$ acts on $M_e$ as multiplication by $a=Nv^{-e}$. Since $Nv>1$ and $e\ge1$, $a-1$ is a nonzero element of $\Ql$, say $a-1=l^ku$ with $u\in\Zl^\times$. Multiplication by $l^ku$ is surjective on the divisible group $\Ql/\Zl$, so $M_e/(\varphi-1)M_e=0$.
\end{proof}

\begin{lemma}\label{w4:lem:indepS}
Let $e\ge1$ and let $S$ be a finite set of finite places of $F$ containing the places above $l$. Let $F_S\subset\overline F$ be the maximal extension of $F$ unramified outside $S$ and the archimedean places, and $G_{F,S}=\Gal(F_S/F)$. Inflation identifies $\Sha_l(F,e)$ with
\[
\ker\Bigl(H^1(G_{F,S},M_e)\longrightarrow\prod_{v\in S}H^1(D_v,M_e)\Bigr).
\]
In particular $\Sha_l(F,e)$ does not depend on $S$.
\end{lemma}

\begin{proof}
The module $M_e$ is unramified outside $l$, so $H=\Gal(\overline F/F_S)$ acts trivially on it. Inflation--restriction gives $0\to H^1(G_{F,S},M_e)\to H^1(G_F,M_e)\to\Hom(H,M_e)$. The group $H$ is the closed normal subgroup of $G_F$ generated by the inertia groups at places outside $S$. A class $c\in H^1(G_F,M_e)$ that vanishes on every $I_v$ restricts to a homomorphism on $H$ that is trivial on these generators and is invariant under conjugation by $G_F$ (as an element of $\Hom(H,M_e)^{G_F}$), hence trivial on $H$. So $c$ is inflated from $G_{F,S}$. Conversely an inflated class vanishes on $I_v$ for $v\notin S$. At $v\notin S$ an unramified class is zero by Lemma~\ref{w4:lem:unram}. At $v\in S$ with $v\nmid l$, vanishing on $D_v$ implies vanishing on $I_v$, and vanishing on $I_v$ implies vanishing on $D_v$ by Lemma~\ref{w4:lem:unram}. At $v\mid l$ both definitions impose vanishing on $D_v$, which implies vanishing on $I_v$. This proves the identification, and it shows that the group does not depend on $S$.
\end{proof}

\begin{lemma}\label{w4:lem:coef}
Let $O$ be a finite free $\Zl$-algebra which is a discrete valuation ring with ramification index $e(O)$ over $\Zl$. Then $\Sha_l(F,e;O)\cong O\otimes_{\Zl}\Sha_l(F,e)$ as $O$-modules, and if $\Sha_l(F,e)$ is finite,
\[
\len_O\Sha_l(F,e;O)=e(O)\,\len_{\Zl}\Sha_l(F,e),\qquad \len_O(O/lO)=e(O).
\]
\end{lemma}

\begin{proof}
Choose a $\Zl$-basis $o_1,\dots,o_r$ of $O$. Then $M_{e,O}=\bigoplus_i o_i\otimes M_e$ as $G_F$-modules, so $H^1(G_F,M_{e,O})=\bigoplus_iH^1(G_F,M_e)$, compatibly with restriction to every closed subgroup. Hence $\Sha_l(F,e;O)=\bigoplus_io_i\otimes\Sha_l(F,e)=O\otimes_{\Zl}\Sha_l(F,e)$, and the $O$-module structures agree because $O$ acts on the first tensor factor in both. The functor $X\mapsto O\otimes_{\Zl}X$ is exact, and it sends $\Z/l$ to $O/lO$, whose $O$-length is $e(O)$. By additivity of length along a composition series of $X$, $\len_O(O\otimes X)=e(O)\len_{\Zl}X$ for every $\Zl$-module $X$ of finite length.
\end{proof}

\begin{lemma}\label{w4:lem:pm}
Let $l$ be odd, $E$ a CM field, $c\in\Gal(E/\Q)$ complex conjugation, $e\ge1$. The natural action of $c$ on $H^1(G_E,M_e)$ preserves $\Sha_l(E,e)$ and squares to the identity. If $\Sha_l(E,e)$ has finite length then $\len\Sha_l(E,e)=\len\Sha_l(E,e)^{c=(-1)^e}+\len\Sha_l(E,e)^{c=(-1)^{e+1}}$.
\end{lemma}

\begin{proof}
The natural action of $\sigma\in G_\Q$ on $H^1(G_E,M_e)$ is $(\sigma\cdot f)(g)=\sigma f(\sigma^{-1}g\sigma)$ on cocycles. It factors through $\Gal(E/\Q)$, permutes the decomposition and inertia groups, and so preserves $\Sha_l(E,e)$. The action of $c$ is an involution. Since $2$ is a unit in $\Zl$, $x=\frac12(x+cx)+\frac12(x-cx)$ splits $\Sha_l(E,e)$ as the direct sum of the $c$-eigenspaces for $\pm1$, and length is additive.
\end{proof}

\begin{proposition}[Restriction to abelian fields]\label{w4:prop:res}
Let $N\ge1$, let $E\subset\Q(\mu_N)$ be a subfield, let $l\nmid N$ be an odd prime and let $e\ge1$ with $(l-1)\nmid e$. Restriction $H^1(G_\Q,M_e)\to H^1(G_E,M_e)$ maps $\Sha_l(\Q,e)$ injectively into $\Sha_l(E,e)$, onto classes fixed by $\Gal(E/\Q)$. If $E$ is a CM field and $e$ is even, the image lies in $\Sha_l(E,e)^{c=(-1)^e}$. In particular $\len\Sha_l(\Q,e)\le\len\Sha_l(E,e)$.
\end{proposition}

\begin{proof}
Restriction of a class that vanishes on $D_v$ (respectively $I_v$) vanishes on $D_w=D_v\cap G_E$ (respectively $I_w$) for every place $w$ of $E$ above $v$, so the image lies in $\Sha_l(E,e)$. The kernel of restriction is $H^1(\Gal(E/\Q),M_e^{G_E})$. Since $l\nmid N$, the fields $\Q(\mu_N)$ and $\Q(\mu_{l^\infty})$ are linearly disjoint over $\Q$, because a subfield of both is unramified at every prime. Hence $\chi(G_E)\supseteq\chi(G_{\Q(\mu_N)})=\Zl^\times$. An element $x\in M_e[l]$ fixed by $G_E$ satisfies $u^{-e}x=x$ for all $u\in\Zl^\times$, and since $(l-1)\nmid e$ there is $u$ with $u^{-e}\not\equiv1\pmod l$, so $x=0$. A nonzero $G_E$-fixed element of $M_e$ of order $l^k$ would give the nonzero fixed element $l^{k-1}x$ of $M_e[l]$. So $M_e^{G_E}=0$ and restriction is injective. A restricted class $f=c|_{G_E}$ satisfies $(\sigma\cdot f)(g)=\sigma c(\sigma^{-1}g\sigma)=c(g)$ modulo coboundaries, by the cocycle relation for $c$, so it is fixed by $\Gal(E/\Q)$. For $E$ CM and $e$ even, $c$ acts trivially on the image, and $1=(-1)^e$.
\end{proof}

\begin{proposition}[Class groups]\label{w4:prop:classgroup}
Let $l$ be an odd prime, $e\ge1$ with $(l-1)\nmid e$, $F_l=\Q(\mu_l)$, $\Delta=\Gal(F_l/\Q)$, $\omega\colon\Delta\to\F_l^\times$ the Teichm\"uller character, and $A=\Cl(F_l)\otimes\Zl$. For a character $\theta$ of $\Delta$ let $A^{(\theta)}=\{a\in A:\delta a=\theta(\delta)a\text{ for all }\delta\}$. Then
\[
\dim_{\F_l}\Sha_l(\Q,e)[l]=\dim_{\F_l}\bigl(A/lA\bigr)^{(\omega^{-e})}.
\]
In particular $\Sha_l(\Q,e)=0$ if and only if $A^{(\omega^{-e})}=0$.
\end{proposition}

\begin{proof}
\emph{Step 1.} $H^0(G_\Q,M_e)=0$ and $H^0(D_l,M_e)=0$. The group $D_l$ maps onto $\Delta$ (the prime $l$ is totally ramified in $F_l$) and acts on $M_e[l]$ through $\omega^{-e}$, which is nontrivial because $(l-1)\nmid e$. So $M_e[l]^{D_l}=0$, and as in the proof of Proposition~\ref{w4:prop:res} this gives $M_e^{D_l}=0$, and a fortiori $M_e^{G_\Q}=0$.

\emph{Step 2.} From $0\to M_e[l]\to M_e\xrightarrow{l}M_e\to0$ and Step 1, the map $H^1(G_\Q,M_e[l])\to H^1(G_\Q,M_e)[l]$ is an isomorphism. For $v\ne l$ the group $I_v$ acts trivially on $M_e$, so the kernel of $H^1(I_v,M_e[l])\to H^1(I_v,M_e)$ is $M_e/lM_e=0$. At $l$ the kernel of $H^1(D_l,M_e[l])\to H^1(D_l,M_e)$ is $H^0(D_l,M_e)/l=0$. Hence $\Sha_l(\Q,e)[l]$ is identified with the group $\mathcal S$ of $c\in H^1(G_\Q,\F_l(\omega^{-e}))$ that vanish on $I_v$ for $v\ne l$ and on $D_l$.

\emph{Step 3.} Since $|\Delta|=l-1$ is prime to $l$, restriction gives an isomorphism $H^1(G_\Q,\F_l(\omega^{-e}))\cong H^1(G_{F_l},\F_l(\omega^{-e}))^\Delta=\Hom(G_{F_l},\F_l)^{\Delta}$, where $G_{F_l}$ acts trivially on the coefficients and $\Delta$-invariance of $f$ means $f(\tilde\delta g\tilde\delta^{-1})=\omega^{-e}(\delta)f(g)$. For $v\ne l$ the extension $F_l/\Q$ is unramified at $v$, so $I_v\subset G_{F_l}$, and $c$ vanishes on $I_v$ iff $f$ does. Restriction from $D_l$ to $D_l\cap G_{F_l}=D_\lambda$ is injective on $H^1(-,\F_l(\omega^{-e}))$ (index prime to $l$), so $c|_{D_l}=0$ iff $f|_{D_\lambda}=0$, where $\lambda=(1-\zeta_l)$ is the prime above $l$.

\emph{Step 4.} The field $F_l$ is totally imaginary. By class field theory a homomorphism $G_{F_l}\to\F_l$ that vanishes on every inertia group factors through the Artin isomorphism $\Art\colon\Cl(F_l)\to\Gal(H/F_l)$ of the Hilbert class field $H$. Its restriction to $D_\lambda$ is determined by its value on $\Art([\lambda])$, which is $0$ because $\lambda$ is principal. The Artin map satisfies $\tilde\delta\Art(a)\tilde\delta^{-1}=\Art(\delta a)$. So $\mathcal S\cong\Hom_\Delta(A/lA,\F_l(\omega^{-e}))$, homomorphisms with $f(\delta a)=\omega^{-e}(\delta)f(a)$. Since $|\Delta|$ is prime to $l$, $A/lA=\bigoplus_\theta(A/lA)^{(\theta)}$, and such an $f$ vanishes on $(A/lA)^{(\theta)}$ for $\theta\ne\omega^{-e}$. Hence $\mathcal S\cong\Hom_{\F_l}((A/lA)^{(\omega^{-e})},\F_l)$. Finally $\Sha_l(\Q,e)$ is an $l$-primary torsion group, so it vanishes iff its $l$-torsion does, and $(A/lA)^{(\theta)}=A^{(\theta)}/lA^{(\theta)}$ vanishes iff $A^{(\theta)}$ does.
\end{proof}

For even $e$ the character $\omega^{-e}$ is even. Vandiver's conjecture for $l$ is the vanishing of $A^{(\theta)}$ for all even $\theta$. By Proposition~\ref{w4:prop:classgroup} it is equivalent to $\Sha_l(\mathbb Q,e)=0$ for all even $e$ with $2\le e\le l-3$, since the trivial character contributes $0$.

For comparison we record the following published identification. Let $l$ be odd and $e\ge1$ with $(l-1)\nmid e$. By Lemma~\ref{w4:lem:indepS} with $S=\{l\}$ and Poitou--Tate duality~\cite[Theorem~I.4.10]{milneADT}, $\Sha_l(\mathbb Q,e)$ is Pontryagin dual to $H^2(\mathbb Z[1/l],\Zl(1+e))$, and this group is $K_{2e}(\mathbb Z)\otimes\Zl$ by the theorem of Rost and Voevodsky on the Quillen--Lichtenbaum conjecture for $\mathbb Z$ (see the survey of Weibel~\cite{weibelK}). Nothing below depends on this identification.

\subsection{The Selmer bound on windows}\label{tw:sec:vwin}

\begin{definition}[Window pairs]\label{w4:def:winpair}
Let $0<\varepsilon<\frac1{52}$ and $Y\ge2$. A \emph{window pair at scale $Y$} is a pair of primes $(q,d)$ with $q\equiv3\pmod8$, $\frac12Y^{\varepsilon/5}<\log q\le Y^{\varepsilon/5}$, $d\equiv7\pmod{16}$, $\frac12Y^{\varepsilon/5}\log^2Y<\log d\le Y^{\varepsilon/5}\log^2Y$, $q\mid d+1$, and $p\mid d+1$ for every odd prime $p\le Y^{\varepsilon/5}$.
\end{definition}

\begin{definition}[$\chi$-fields]\label{w4:def:chifield}
Let $\psi_0$ be a primitive Dirichlet character of conductor $N_0$, $e\ge0$, and $q,d$ odd primes. Put $N=\operatorname{lcm}(2^{e+2},N_0,q,d)$. The field $E(\psi_0,e,q,d)$ is the subfield of $\Q(\mu_N)$ fixed by the automorphisms $\zeta\mapsto\zeta^a$ with $a\equiv\pm1\pmod{2^{e+2}}$, $\psi_0(a)=1$, $\left(\frac aq\right)=\left(\frac ad\right)=1$. For $q\equiv3\pmod4$ it is a CM field.
\end{definition}

\begin{hypothesis}[Selmer bound on windows]\label{w4:thm:vwin}
For every primitive Dirichlet character $\psi_0$, every $0<\varepsilon<\frac1{52}$ and every function $e\colon[1,\infty)\to\mathbb N$ with $x^{2+8\varepsilon}\le2^{e(x)}\le2x^{2+8\varepsilon}$, there are $J\ge1$ and $X_0$ with the following property. For every $X\ge X_0$ and every window pair $(q,d)$ at scale $X^{1/J}$ there is $Y$ with $X^{1/J}\le Y\le X^{1/2}$ such that for every $x\in[Y,Y^2]$, with $t=\lfloor x^{401}\rfloor$ and $E=E(\psi_0,e(x),q,d)$, every $\Sha_l(E,e')$ with $0<e'<t$ is finite and
\[
\sum_{l}\ \sum_{1\le e'<t}\lambda\bigl(\Sha_l(E,e')\bigr)\le\exp\bigl(\tfrac12x^{\varepsilon}\bigr),
\]
the sum over all primes $l$ (including $l=2$).
\end{hypothesis}

Proposition~\ref{r6:R6:thm:selmer} deduces (D6) by rescaling the window.

For a CM field $E$, $t\ge1$ and $B\ge0$ write
\begin{itemize}[leftmargin=*,itemsep=1pt]
\item $(\mathrm P_2)(E,t,B)$: $\sum_{1\le e'<t}\len_{\mathbb Z_2}\Sha_2(E,e')\le B$, all finite;
\item $(\mathrm P_\pm)(E,t,B)$: for odd $l$, the $c=\pm(-1)^{e'}$ parts are finite and
\[
\sum_{l\ \mathrm{odd}}\log l\sum_{e'<t}\len\Sha_l(E,e')^{c=\pm(-1)^{e'}}\le B.
\]
\end{itemize}

\begin{theorem}[decomposition]\label{w4:thm:decomp}
Hypothesis~\ref{w4:thm:vwin} follows from the following two statements, with the notation of Hypothesis~\ref{w4:thm:vwin}.
\begin{enumerate}[leftmargin=*,label=(\Roman*)]
\item \emph{(Plus part and $l=2$ on windows.)} For every $\psi_0$, $\varepsilon$ and every pair of functions $e(x),t(x)$ with $2\le e(x)$, $2^{e(x)}\le x^4$, $2\le t(x)\le x^{401}$ for $x\ge2$, there are $J,C,P,X_1>0$ such that for $X\ge X_1$ and primes $q\equiv3\ (8)$, $d\equiv7\ (16)$ prime to $N_0$, with $q,d>\max(P,2X^{401}+1)$ and $\log(qd)\le X^{\varepsilon/(4J)}$, there is $Y$ with $X^{1/J}\le Y\le X^{1/2}$ such that for all $x\in[Y,Y^2]$ and $E=E(\psi_0,e(x),q,d)$: $(\mathrm P_2)(E,t(x),e^{x^{4\varepsilon/5}})$ and $(\mathrm P_+)(E,t(x),(\cond(E)t(x))^C+Ce^{x^{4\varepsilon/5}})$.
\item \emph{(Minus part.)} There is $C$ such that for all such $E$ of degree $n$ and conductor $f$, and $t\ge2$: $(\mathrm P_-)(E,t,\beta)$ with $\beta=3nt^2\log(tf)+C\sum_{l\mid o,\ l\ \mathrm{odd}}l\,n^2t^2\log^2(tlf)$, where $o$ is the order of $\psi_0$.
\end{enumerate}
\end{theorem}

\begin{proof}
Fix $\psi_0$, $\varepsilon$ and $e(x)$ as in Hypothesis~\ref{w4:thm:vwin}, and put $t(x)=\lfloor x^{401}\rfloor$. For $x\ge2$, $2^{e(x)}\le2x^{2+8\varepsilon}\le x^4$ and $2^{e(x)}\ge x^2\ge4$, so (I) applies; let $J_{\rm I},C,P,X_1$ be its constants and put $J=\max(J_{\rm I},1)$, and let $C_-$ be the constant of (II).

\emph{Sizes.} Let $(q,d)$ be a window pair at scale $Y_0=X^{1/J}$ and $x\ge Y_0$. Then $\log q+\log d\le Y_0^{\varepsilon/5}(1+\log^2Y_0)\le x^{\varepsilon/4}$ once $Y_0$ is large. Also $q,d\ge\exp(\frac12Y_0^{\varepsilon/5})>\max(P,2X^{401}+1)$ and $\log(qd)\le X^{\varepsilon/(4J)}$ for $X$ large, so the hypotheses of (I) hold. The field $E$ lies in $\Q(\mu_N)$ with $N=\operatorname{lcm}(2^{e(x)+2},N_0,q,d)$, so $f=\cond(E)\le N$ and $n=[E:\Q]\le\varphi(N)\le N$. Since $2^{e(x)}\le2x^{3}$,
\[
\log(ft)\le(e(x)+2)\log2+\log N_0+x^{\varepsilon/4}+401\log x\le2x^{\varepsilon/4}
\]
for $x$ large in terms of $\varepsilon$ and $N_0$.

\emph{Assembly.} Let $x\in[Y,Y^2]$ with $Y$ from (I). For odd $l$, Lemma~\ref{w4:lem:pm} splits each length into the two parts. Summing over $l$,
\[
\sum_l\sum_{e'<t}\lambda(\Sha_l(E,e'))\le e^{x^{4\varepsilon/5}}\log2+\beta+(ft)^C+Ce^{x^{4\varepsilon/5}}.
\]
Since $n\le f$, $\ell\le o$ for $\ell\mid o$, and $\log^2y\le y$ for $y\ge1$, each summand of $\beta$ is at most $o^2(ft)^3$, so $\beta\le(3+C_-o^3)(ft)^3\le(3+C_-o^3)e^{6x^{\varepsilon/4}}$; also $(ft)^C\le e^{2Cx^{\varepsilon/4}}$. Each of the four terms is at most $\frac14\exp(\frac12x^\varepsilon)$ once $x$ is large in terms of $\varepsilon$, $C$, $C_-$, $o$, because $\frac45\varepsilon<\varepsilon$ and $\frac14\varepsilon<\varepsilon$. Choose $X_0\ge X_1$ so large that $Y_0=X_0^{1/J}$ exceeds all thresholds above. Then for $X\ge X_0$ every $x\ge Y\ge X^{1/J}$ is large enough.
\end{proof}

\notgiven{Neither of the statements (I) and (II) is proved here.} The plus part of (I) splits, by the residual position of $l$ relative to the host, into five counts of the same type. The generic one is a bound for the groups $\Sha_l(\mathbb Q,e';\psi)$ attached to the Dirichlet characters $\psi$ of the host, and it follows from statement (CorGK) of Subsection~\ref{tw:sec:vandiver} for $K=\mathbb Q(\sqrt{-p_K})$. Cases (II) and (III) use $\mathrm{GSp}_4$. Case (V) uses a deformation problem over a frame of the host field. The part of (I) at $l=2$ and at the bad primes uses Poitou--Tate duality, local duality, a surjection onto $H^2(G_{L,S},\mathbb Z_l(e+1))$ from a module that maps to $\mathbb Z_l\otimes H_{2e}(\mathrm{GL}_{4e+4}(\mathcal O_L),\mathbb Z)$ with kernel killed by $l^{c(2e)}$, torsion bounds of the type of Soul\'e and of Dwyer and Friedlander, and Borel's computation of ranks. The bound is available only on windows of scales, and the constant $J=2^{R+2}(2C_K\cdot401/\varepsilon)^R$ of the formal proof, with $R=R(\psi_0)$, depends on $\psi_0$ and $\varepsilon$. Statement (II) is of the type of the Main Conjecture for abelian fields (Mazur and Wiles, Wiles, Rubin) together with the theorem of Ferrero and Washington, which bound the minus parts by generalized Bernoulli numbers, whose logarithms are polynomial in $n$, $t$ and $\log f$. The formal proof derives (II) from Stickelberger's theorem and a Kolyvagin system of Gauss sums.

\begin{proposition}[plus parts through characters]\label{w4:prop:pluschar}
Let $N\ge1$, let $H\subset(\Z/N)^\times$ be a subgroup with $-1\notin H$, and let $E\subset\Q(\mu_N)$ be the fixed field of $H$, so that $E$ is a CM field with $\Gal(E/\Q)\cong(\Z/N)^\times/H$. Let $l$ be an odd prime with $l\nmid[(\Z/N)^\times:H]$, let $e\ge1$, and let $O$ be a finite extension of $\Zl$ which is a discrete valuation ring and contains the values of all characters of $(\Z/N)^\times/H$. Let $X_e$ be the set of Dirichlet characters $\psi$ modulo $N$ with values in $O$ that are trivial on $H$ and satisfy $\psi(-1)=(-1)^e$. Then $\Sha_l(E,e)^{c=(-1)^e}$ is finite if and only if every $\Sha_l(\Q,e;\psi)$ with $\psi\in X_e$ is finite, and in that case
\[
\len_{\Zl}\Sha_l(E,e)^{c=(-1)^e}=\sum_{\psi\in X_e}\frac{\len_O\Sha_l(\Q,e;\psi)}{e(O)} .
\]
\end{proposition}

\begin{proof}
Put $\Gamma=\Gal(E/\Q)$, $M=M_{e,O}$ and $A=O[\Gamma]$, on which $G_\Q$ acts through $\Gamma$ by left multiplication. Since $M$ is a $G_\Q$-module, the induced module $\operatorname{Ind}_{G_E}^{G_\Q}M$ is isomorphic to $A\otimes_OM$ with the diagonal action, and Shapiro's lemma gives an isomorphism $H^1(G_E,M)\cong H^1(G_\Q,A\otimes_OM)$. The order of $\Gamma$ is a unit in $O$ and $O$ contains the values of the characters of $\Gamma$, so the elements $\varepsilon_\psi=\lvert\Gamma\rvert^{-1}\sum_{\gamma\in\Gamma}\psi(\gamma)^{-1}\gamma$ are orthogonal idempotents with $\sum_\psi\varepsilon_\psi=1$ and $\gamma\varepsilon_\psi=\psi(\gamma)\varepsilon_\psi$. Hence $A\otimes_OM=\bigoplus_\psi\varepsilon_\psi A\otimes_OM\cong\bigoplus_\psi M(\psi)$ as $G_\Q$-modules, the sum over all characters $\psi$ of $\Gamma$, where $M(\psi)$ is $M$ with the action twisted by $\psi$. This gives $H^1(G_E,M)\cong\bigoplus_\psi H^1(G_\Q,M(\psi))$.

Let $v$ be a finite prime and let $D$ be $D_v$ or $I_v$. By Mackey's formula the restriction of $\operatorname{Ind}_{G_E}^{G_\Q}M$ to $D$ is the sum over the double cosets $D\gamma G_E$ of the modules induced from $D\cap\gamma G_E\gamma^{-1}$, and Shapiro's lemma for each summand identifies the restriction map to $H^1(D,A\otimes_OM)$ with the product of the restriction maps of $H^1(G_E,M)$ to the decomposition groups, respectively the inertia groups, of the places of $E$ above $v$. So a class of $H^1(G_E,M)$ lies in $\Sha_l(E,e;O)$ if and only if its image in $H^1(G_\Q,A\otimes_OM)$ vanishes on $D_l$ and on every $I_v$. The decomposition $A\otimes_OM=\bigoplus_\psi M(\psi)$ is stable under $G_\Q$, hence under $D$, so these conditions hold if and only if they hold in each component. Therefore $\Sha_l(E,e;O)\cong\bigoplus_\psi\Sha(\psi)$, where $\Sha(\psi)\subset H^1(G_\Q,M(\psi))$ is the subgroup of classes vanishing on $D_l$ and on every $I_v$. Since $\psi\mapsto\psi^{-1}$ is a bijection of the characters of $\Gamma$, $\Sha(\psi)$ is the group $\Sha_l(\Q,e;\psi^{-1})$ of the text.

The natural action of $\gamma\in\Gamma$ on $H^1(G_E,M)$ of Lemma~\ref{w4:lem:pm} corresponds under Shapiro's isomorphism to the action of $A\otimes_OM$ induced by multiplication by $\gamma$ or by $\gamma^{-1}$ on $A$, according to the convention, which commutes with $G_\Q$ because $\Gamma$ is abelian. On $\varepsilon_\psi A$ it is multiplication by $\psi(\gamma)^{\pm1}$. Complex conjugation $c$ is the class of $-1$, and $\psi(-1)^{\pm1}=\psi(-1)$. Hence $\Sha_l(E,e;O)^{c=(-1)^e}\cong\bigoplus_{\psi(-1)=(-1)^e}\Sha(\psi)$. The set $X_e$ is stable under $\psi\mapsto\psi^{-1}$, so the right side is $\bigoplus_{\psi\in X_e}\Sha_l(\Q,e;\psi)$. Since $2$ is a unit in $O$, the eigenspace decomposition commutes with $O\otimes_{\Zl}-$, and Lemma~\ref{w4:lem:coef} gives $\len_O\Sha_l(E,e;O)^{c=(-1)^e}=e(O)\len_{\Zl}\Sha_l(E,e)^{c=(-1)^e}$. Length is additive on finite direct sums, which gives both assertions.
\end{proof}

\begin{proposition}[the generic plus part from (CorGK)]\label{w4:prop:genericplus}
Assume statement (CorGK) of Proposition~\ref{w4:prop:corToBound} for an imaginary quadratic field $K$, with constants $c_5,A$. Let $N\ge16$ be an integer, $e\ge1$ and $t\ge e+1$ integers, $r_1\ne r_2$ primes not dividing $30N\lvert d_K\rvert$, and $k\ge0$ an integer with $2c_5(\log r_1+\log r_2)\le16^k$. Let $L$ be a finite set of primes and, for $l\in L$, let $O_l$ be a coefficient ring as in Definition~\ref{w4:def:generic} and $X_l$ a set of Dirichlet characters modulo $N$ with values in $O_l$, each in generic position for $(K,l,e)$. Then every $\Sha_l(\Q,e;\psi)$ with $l\in L$ and $\psi\in X_l$ is finite, and
\[
\sum_{l\in L}\log l\sum_{\psi\in X_l}\frac{\len_{O_l}\Sha_l(\Q,e;\psi)}{e(O_l)}\le2c_5(\log r_1+\log r_2)\,N\,\bigl(N(e+1)\bigr)^A\le(Nt)^{k+1+A}.
\]
\end{proposition}

\begin{proof}
For a subgroup $H'\subset(\Z/N)^\times$ let $\mathcal O_l(H')$ be the set of characters with values in $O_l$ whose kernel is $H'$. Let $T$ be the set of kernels of the characters in $\bigcup_lX_l$. Each $X_l$ is contained in $\bigcup_{H'\in T}\mathcal O_l(H')$, and for $H'\in T$ let $L(H')$ be the set of $l\in L$ such that $X_l$ meets $\mathcal O_l(H')$. For $l\in L(H')$ some character of $\mathcal O_l(H')$ is in generic position, so (CorGK) with $H=H'$, $p=r_1$, $p'=r_2$ and the set $L(H')$ gives finiteness and
\[
\sum_{l\in L(H')}\log l\sum_{\psi\in\mathcal O_l(H')}\frac{\len\Sha_l(\Q,e;\psi)}{e(O_l)}\le2c_5\,\varphi([(\Z/N)^\times:H'])\,\bigl(N(e+1)\bigr)^A(\log r_1+\log r_2).
\]
All terms are nonnegative, so summing over $H'\in T$ bounds the left side of the proposition. The quotient $(\Z/N)^\times/H'$ is cyclic for $H'\in T$, so exactly $\varphi([(\Z/N)^\times:H'])$ complex characters of $(\Z/N)^\times$ have kernel $H'$. Distinct $H'$ give disjoint sets of characters, hence $\sum_{H'\in T}\varphi([(\Z/N)^\times:H'])\le\varphi(N)\le N$. This gives the first inequality. For the second, $2c_5(\log r_1+\log r_2)\le16^k\le N^k$, $N\le Nt$ and $N(e+1)\le Nt$.
\end{proof}

The hypotheses $r_1,r_2\nmid30N\lvert d_K\rvert$ and the generic position of the characters are the points at which the host enters. \notgiven{The verification that the characters of the host attached to a twist $e$ and a prime $l$ outside the finite set of bad primes with $(l-1)\nmid2e$ are in generic position for $K=\Q(\sqrt{-p_K})$, and that the host level $N$ is polynomial in the conductor of $E$, are not written out here.}

\subsection{A consequence in the direction of Vandiver's conjecture}\label{tw:sec:vandiver}

Consider the following statement.
\begin{enumerate}
\item[(V)] There are constants $C,A$ such that for every even $e\ge2$ and every finite set $L$ of primes $l\ge5$ with $(l-1)\nmid2e$, each $\Sha_l(\mathbb Q,e)$ is finite and $\sum_{l\in L}\lambda\bigl(\Sha_l(\mathbb Q,e)\bigr)\le C(e+1)^A$.
\end{enumerate}
The formal proof contains (V) as a theorem, derived from the statement that also gives the generic plus part of (I).

\begin{corollary}[Vanishing of eigenspaces]\label{w4:cor:large}
If (V) holds with constants $C,A$, then for every even $e\ge2$ and every prime $l\ge5$ with $(l-1)\nmid2e$ and $\log l>C(e+1)^A$, the $\omega^{-e}$ eigenspace of the $l$-class group of $\Q(\mu_l)$ vanishes.
\end{corollary}

\begin{proof}
If it did not vanish, $\Sha_l(\Q,e)\ne0$ by Proposition~\ref{w4:prop:classgroup}, so $\lambda(\Sha_l(\Q,e))\ge\log l$, and (V) with $L=\{l\}$ gives $\log l\le C(e+1)^A$.
\end{proof}

The exclusion $(l-1)\mid2e$ removes nothing of substance. If $(l-1)\mid2e$ and $(l-1)\nmid e$, then $\omega^{-e}$ is the quadratic character of $\Delta$, whose field is $\mathbb Q(\sqrt l)$, and the corresponding eigenspace is the $l$-part of a class group whose order is smaller than $l$.

\paragraph{Comparison with published results.} Kurihara~\cite{kurihara92} proved $A^{(\omega^{-2})}=0$ for every $l$, through $K_4(\mathbb Z)=0$, which is the case $e=2$ of Corollary~\ref{w4:cor:large} with no threshold. With $n=e+1$ odd, Soul\'e~\cite{soule99} proved $A^{(\omega^{-e})}=0$ once $\log l>n^{224n^4}$, a threshold of size $\exp(O(e^4\log e))$. De~la~Torre~Dur\'an~\cite{delatorre} improved the exponential factor for $F=\mathbb Q$ from $e^{O(n^4\log n)}$ to $e^{O(n^2\log n)}$, a threshold of size $\exp(O(e^2\log e))$. Corollary~\ref{w4:cor:large} has the threshold $C(e+1)^A$, one exponential smaller. In terms of $T(e)=\sum_l\lambda(\Sha_l(\mathbb Q,e))$, the published methods bound $\log T(e)$ by a constant times $e^2\log e$ at best, and (V) asserts $\log T(e)\le A\log(e+1)+\log C$. The methods of Soul\'e and of de~la~Torre~Dur\'an count cells of the Voronoi complex of $\mathrm{GL}_n(\mathbb Z)$, and the number of perfect forms grows like $e^{\Theta(n^2\log n)}$~\cite{vanwoerden}, so cell counting alone does not reach a polynomial bound. No finite computation can contradict (V), because its constants are existential. The computed groups $K_4(\mathbb Z)=0$ and $K_8(\mathbb Z)=0$ and the verification of Vandiver's conjecture for $l<2^{31}$ are consistent with every upper bound, and they give no information about (V).

\begin{remark}[odd twists]\label{w4:rem:odd}
For odd $e$ the analogous sum is classically polynomial. By the identification with $K_{2e}(\mathbb Z)\otimes\Zl$ in Subsection~\ref{tw:sec:sha} and the computation of $K_{4k-2}(\Z)$ from the Main Conjecture of Mazur and Wiles (Weibel's survey), $|\Sha_l(\Q,e)|$ is the $l$-part of the numerator of $B_{e+1}/(e+1)$ for $e=2k-1$ and $(l-1)\nmid e$. Hence $\sum_l\lambda(\Sha_l(\Q,e))\le\log|B_{e+1}|+o(e)=(1+o(1))\,e\log e$, using that the denominator of $B_{e+1}/(e+1)$ has logarithm $o(e)$. So the shape of (V) is the natural even analogue of a known odd statement. This supports plausibility and proves nothing.
\end{remark}

\paragraph{The chain of statements.} Let $K$ be an imaginary quadratic field. For a Dirichlet character $\psi$ write $\Sha_l(\mathbb Q,e;\psi)$ for the group of Definition~\ref{w4:def:sha} for the character $\psi^{-1}\chi^{-e}$ in place of $\chi^{-e}$, with coefficients in a finite extension $O$ of $\Zl$ that contains the values of $\psi$.

\begin{definition}[generic position]\label{w4:def:generic}
Let $K$ be a number field, $l$ a prime, $e\ge0$ and $N\ge1$ integers, $O$ a finite extension of $\Zl$ which is a discrete valuation ring, and $\psi$ a Dirichlet character modulo $N$ with values in $O$. For a prime $r$ let $\psi_{(r)}$ be the restriction of $\psi$ to the factor $(\Z/r^{v_r(N)})^\times$ of $(\Z/N)^\times$, viewed again as a character modulo $N$, and put $\psi^\circ=\psi\,\psi_{(l)}^{-1}$. Let $\omega\colon G_\Q\to\mathbb F_l^\times$ be the reduction of the cyclotomic character modulo $l$. We say that $\psi$ is in generic position for $(K,l,e)$ when the following five conditions hold.
\begin{enumerate}[leftmargin=*,label=(\alph*)]
\item $l\ne2$, $e\ge1$ and $N$ is even.
\item $l$ does not divide the order of $\psi$, $\psi_{(l)}^2=1$, and $\psi(-1)=(-1)^e$.
\item $K$ is imaginary quadratic, $l\nmid\lvert d_K\rvert h_K$, and $\gcd(d_K,N)=1$.
\item The subgroup $\ker\omega\cap\ker\psi^\circ$ of $G_\Q$ is not contained in $G_K$.
\item $(l-1)\nmid2e$.
\end{enumerate}
\end{definition}

For $\psi=1$ one has $\psi_{(l)}=\psi^\circ=1$, and conditions (b) and (d) become the conditions written in the proof of Proposition~\ref{w4:prop:corToBound}. For a general $\psi$ the condition in (b) concerns the $l$-component $\psi_{(l)}$ and not $\psi$ itself, and the condition in (d) concerns $\psi^\circ$.

For $\psi$ as above and $e\ge1$ we write $M_{e,O}(\psi^{-1})$ for $O\otimes_{\Zl}\Ql/\Zl$ with $\sigma\in G_\Q$ acting by multiplication by $\psi^{-1}(\sigma)\chi(\sigma)^{-e}$, where $\psi$ is viewed as a character of $G_\Q$ through $\Gal(\Q(\mu_N)/\Q)\cong(\Z/N)^\times$. Then $\Sha_l(\Q,e;\psi)$ is the subgroup of $H^1(G_\Q,M_{e,O}(\psi^{-1}))$ of the classes that vanish on $D_l$ and on $I_v$ for every finite $v$.

\begin{lemma}\label{w4:lem:unramchi}
Let $v\ne l$ be a prime, $e\ge1$, and let $\psi$ have order prime to $l$. A class in $H^1(D_v,M_{e,O}(\psi^{-1}))$ that vanishes on $I_v$ is zero. In particular $\Sha_l(\Q,e;\psi)$ is also the group of classes that vanish on $D_v$ for every $v\mid lN$ and on $I_v$ for every finite $v$.
\end{lemma}

\begin{proof}
Write $M=M_{e,O}(\psi^{-1})$ and let $\mathfrak m$ be the maximal ideal of $O$. Inflation and restriction give the exact sequence $0\to H^1(D_v/I_v,M^{I_v})\to H^1(D_v,M)\to H^1(I_v,M)$. Suppose first that $\psi$ is ramified at $v$, and choose $g\in I_v$ with $\psi(g)\ne1$. The element $\psi(g)$ is a root of unity of order $m>1$ prime to $l$, so $1-\psi(g)$ divides $m$ in $O[\psi(g)]$ and $1-\psi(g)\in O^\times$. The group $I_v$ acts trivially through $\chi$ because $v\ne l$, so $g$ acts on $M$ by $\psi(g)^{-1}$, and $(\psi(g)^{-1}-1)m'=0$ forces $m'=0$. Hence $M^{I_v}=0$. Suppose next that $\psi$ is unramified at $v$. Then $M^{I_v}=M$, and a Frobenius element $\varphi$ acts by $a=\psi(v)^{-1}v^{-e}$. If $a=1$, then $v^e$ would be a root of unity in the field of fractions of $O$, which is impossible because $v^e>1$. So $a-1\ne0$, multiplication by $a-1$ is surjective on the divisible $O$-module $M$, and $H^1(D_v/I_v,M)\cong M/(\varphi-1)M=0$. In both cases the inflated group vanishes. For the last assertion, a class that vanishes on $I_v$ for $v\mid N$, $v\ne l$, then vanishes on $D_v$, and conversely vanishing on $D_v$ implies vanishing on $I_v$.
\end{proof}

\begin{proposition}[(V) from (CorGK)]\label{w4:prop:corToBound}
Assume the following statement (CorGK) for $K=\Q(\sqrt{-3})$. There are $c_5,A$ such that for every even $N\ge2$, every subgroup $H\subset(\Z/N)^\times$, all distinct primes $p,p'\nmid30N\lvert d_K\rvert$, every finite set $L$ of primes such that for each $l\in L$ some character $\psi$ in the orbit of characters with kernel $H$ is in generic position in the sense of Definition~\ref{w4:def:generic}, one has finiteness and
\[
\sum_{l\in L}\log l\sum_{\psi}\frac{\len\Sha_l(\Q,e;\psi)}{e(O)}\le2c_5\,\varphi([(\Z/N)^\times:H])\,\bigl(N(e+1)\bigr)^A\bigl(\log p+\log p'\bigr),
\]
where $\Sha_l(\Q,e;\psi)$ is the group of Definition~\ref{w4:def:sha} for the character $\psi^{-1}\chi^{-e}$. Then (V) holds with $C=36c_52^A$ and the same $A$.
\end{proposition}

\begin{proof}
Take $N=2$, $H=(\Z/2)^\times$ (the trivial group), so the only character is $\psi=1$ and $\varphi(1)=1$, and take $p=7$, $p'=11$, which do not divide $30\cdot2\cdot3$. For $\psi=1$, $\Sha_l(\Q,e;1)=\Sha_l(\Q,e)$. The generic-position conditions for $K=\Q(\sqrt{-3})$, $\psi=1$, $l\ge5$, $e\ge2$ even and $(l-1)\nmid2e$ hold, as follows. (a) $l$ odd, $e\ge1$, $N=2$ even. (b) $l$ does not divide the order $1$ of $\psi$, $\psi^2=1$, and $\psi(-1)=1=(-1)^e$. (c) $K$ is imaginary quadratic with discriminant $-3$ and class number $1$, so $l\nmid3$; and $\gcd(3,2)=1$. (d) The last condition asks that $\ker(\chi\bmod l)\cap\ker\psi$ is not contained in $G_K$. Here $\ker\psi=G_\Q$, so it asks that $K\not\subset\Q(\mu_l)$, which holds because $3$ ramifies in $K$ and not in $\Q(\mu_l)$ for $l\ne3$. (e) $(l-1)\nmid2e$. The right side becomes $2c_5(2(e+1))^A(\log7+\log11)\le2c_5\,2^A(e+1)^A\cdot 18$. Lemma~\ref{w4:lem:coef} converts the left side into the left side of (V).
\end{proof}

\begin{proposition}[(CorGK) from (PropGK)]\label{w4:prop:prop2cor}
Assume the following statement (PropGK). For a number field $K$ there are $c_5,A$ such that for $e,N,H$ and every prime $p\nmid30N\lvert d_K\rvert$ there is an integer $NN_p\ne0$ with $\log\lvert NN_p\rvert\le c_5\varphi(\mathrm{index}\,H)(N(e+1))^A\log p$, and for every prime $l\ne p$ and every coefficient ring at which some character of the orbit is in generic position, the lengths are finite and $\sum_\psi\len\Sha_l(\Q,e;\psi)/e(O)\le2v_l(NN_p)$. Then (CorGK) holds with the same $c_5,A$.
\end{proposition}

\begin{proof}
Fix distinct admissible $p,p'$ and a finite set $L$ of primes in generic position. For $l\ne p$ use $NN_p$, and for $l=p$ use $NN_{p'}$. Then
\[
\sum_{l\in L}\log l\sum_\psi\frac{\len\Sha_l(\Q,e;\psi)}{e(O)}\le2\sum_l\log l\,\bigl(v_l(NN_p)+v_l(NN_{p'})\bigr)\le2\bigl(\log\lvert NN_p\rvert+\log\lvert NN_{p'}\rvert\bigr),
\]
because $\sum_lv_l(n)\log l\le\log|n|$ for a nonzero integer $n$. Insert the bounds for $\log|NN_p|$ and $\log|NN_{p'}|$.
\end{proof}

\begin{proposition}[(PropGK) from (Count) and (OneChar)]\label{w4:prop:onechar2prop}
Let $K$ be an imaginary quadratic field, $p\ge7$ a prime with $p\nmid N\,d_K$, $w\mid p$ a place of $K$ and $\sigma$ a geometric Frobenius element at $w$. Let $\mathcal O$ be a finite set of roots of unity $z$ stable under $\operatorname{Aut}(\mathbb C)$, and put $t(z)=Nw^{-e}z+1+Nw^{e}z^{-1}$, the trace of $\sigma$ on the Eisenstein triple $\Psi\oplus1\oplus\Psi^{-1}$ suitably normalized. Assume the following two statements.
\begin{enumerate}[leftmargin=*,label=(\alph*)]
\item \emph{(Count.)} There is $NN\in\Z\setminus\{0\}$ with $\log|NN|\le c\,|\mathcal O|\,(N(e+1))^A\log p$, such that for every prime $l\ne2,p$, every $\iota\colon\overline{\Q}_l\cong\mathbb C$ and every choice of finite sets $X_z$ of pairwise non-isomorphic admissible $\rho\colon G_K\to GL_3(\overline{\Q}_l)$ (level $N$, weight $e$), one has $\operatorname{tr}\rho(\sigma)\ne t(z)$ for $\rho\in X_z$ and $\lvert NN\rvert_l\le\prod_z\prod_{\rho\in X_z}\lvert\operatorname{tr}\rho(\sigma)-t(z)\rvert_l$.
\item \emph{(OneChar.)} For each character $\psi$ in generic position, if no admissible $\rho$ has $\operatorname{tr}\rho(\sigma)=t$, then $\Sha_l(\Q,e;\psi)$ is finite and there is a finite set $X_\psi$ of pairwise non-isomorphic admissible $\rho$ with $\len\Sha_l(\Q,e;\psi)/e(O)\le2\sum_{\rho\in X_\psi}v_l(\operatorname{tr}\rho(\sigma)-t_\psi)$.
\end{enumerate}
Then the conclusion of (PropGK) holds at $p$ with $NN_p=NN$.
\end{proposition}

\begin{proof}
The orbit of characters $\psi$ with kernel $H$ is in bijection with an $\operatorname{Aut}(\mathbb C)$-stable set $\mathcal O$ of roots of unity, $z_\psi=\psi(Nw)^{-1}$ under $\iota$, and $t_\psi=t(z_\psi)$; the number of elements is $\varphi(\mathrm{index}\,H)$. By~(a) no admissible $\rho$ has trace $t_\psi$, so~(b) applies to each $\psi$. Summing~(b) over $\psi$ and applying~(a) to the sets $X_\psi$,
\[
\sum_\psi\frac{\len\Sha_l(\Q,e;\psi)}{e(O)}\le2\sum_\psi\sum_{\rho\in X_\psi}v_l\bigl(\operatorname{tr}\rho(\sigma)-t_\psi\bigr)\le2v_l(NN).\qedhere
\]
\end{proof}

The content of (Count) is classical in type. The admissible $\rho$ correspond to Frobenius eigenvalues at $w$ of algebraic automorphic forms on a definite unitary group in three variables of level $N$ and weight of height $e$, whose number is polynomial in $N(e+1)$. The number $NN$ is the product of the scaled differences, rational by Galois stability, and nonzero because purity gives $|\operatorname{tr}\rho(\sigma)|\le3$ while $|t(z)|\ge p^{e}-2>3$. \notgiven{Statement (Count) is not proved here.}

\begin{proposition}[a consequence of the Selmer bound]\label{w4:prop:need}
Assume Hypothesis~\ref{w4:thm:vwin} for one primitive character $\psi_0$ of conductor $N_0$, and assume that for every $\varepsilon\in(0,\frac1{52})$ and $J\ge1$ there are arbitrarily large $X$ for which a window pair $(q,d)$ at scale $X^{1/J}$ exists (Definition~\ref{w4:def:winpair}). Then for every $\varepsilon\in(0,\frac1{52})$ there are arbitrarily large real $x$ and primes $q_x,d_x$ such that for every finite set $L$ of odd primes $l\nmid N_0q_xd_x$,
\[
\sum_{l\in L}\ \sum_{\substack{1\le e<x^{401}\\(l-1)\nmid e}}\lambda\bigl(\Sha_l(\Q,e)\bigr)\le\exp\bigl(\tfrac12x^{\varepsilon}\bigr).
\]
Consequently the product $P(T)=\prod_{e<T}\prod_{l}\lvert\Sha_l(\Q,e)\rvert$ over odd $l\nmid N_0$ with $(l-1)\nmid e$, outside two primes that depend on $T$, satisfies $\log\log P(T)\le\frac12T^{\varepsilon/401}$ for arbitrarily large $T$, for every $\varepsilon>0$.
\end{proposition}

\begin{proof}
Choose any function $e(x)$ with $x^{2+8\varepsilon}\le2^{e(x)}\le2x^{2+8\varepsilon}$. Let $J,X_0$ be given by Hypothesis~\ref{w4:thm:vwin}, take $X\ge X_0$ with a window pair $(q,d)$ at scale $X^{1/J}$, and take the $Y$ that Hypothesis~\ref{w4:thm:vwin} provides. Put $x=Y$, which tends to infinity with $X$. Let $N=\operatorname{lcm}(2^{e(x)+2},N_0,q,d)$ and let $E$ be the fixed field in $\Q(\mu_N)$ of the subgroup of Definition~\ref{w4:def:chifield}. Hypothesis~\ref{w4:thm:vwin} gives $\sum_{l\in L}\sum_{1\le e<t}\lambda(\Sha_l(E,e))\le\exp(\frac12x^\varepsilon)$ with $t=\lfloor x^{401}\rfloor$. For odd $l\nmid N_0qd$ we have $l\nmid N$, and Proposition~\ref{w4:prop:res} gives $\lambda(\Sha_l(\Q,e))\le\lambda(\Sha_l(E,e))$ whenever $(l-1)\nmid e$. The terms are nonnegative, so the inequality follows. For the consequence take $T=x^{401}$ and logarithms.
\end{proof}

By the identification with $K_{2e}(\mathbb Z)\otimes\mathbb Z_l$ and the comparison above, the published estimates give at best $\log\log P(T)\ll T^{2}\log T$, and nothing better along any sequence of $T$. Proposition~\ref{w4:prop:need} needs $\log\log P(T)\le T^{\delta}$ along a sequence of $T$, for every $\delta>0$. So the restriction to windows does not reduce what Hypothesis~\ref{w4:thm:vwin} requires to a published bound. The Selmer bound on windows implies a bound in the direction of Vandiver's conjecture that is stronger than the published bounds of Soul\'e and of de~la~Torre~Dur\'an. Alp\"oge's announcement~\cite{note} describes this bound as a weakening of Vandiver's conjecture.

\paragraph{The statement (OneChar).} Fix $K$, $l\ge5$, $e$ and a character $\psi$ in generic position, put $\Psi=\chi^e\psi$, and consider the ordinary deformation problem for representations of $G_{\mathbb Q}$ into the group of unitary type in three variables with residual representation $\Psi\oplus1\oplus\Psi^{-1}$. The formal proof derives (OneChar) from four inputs.
\begin{enumerate}[leftmargin=*,label=(G\arabic*)]
\item \emph{Tangent injection.} A class of $\Sha_l(\mathbb Q,e;\psi)$ gives a tangent vector of weight $-1$ at the Eisenstein point of a graded model of the ordinary deformation ring, with an action of a torus (the argument of Ribet).
\item \emph{An intersection inequality.} On a trigraded Rees algebra, under hypotheses (M1) to (M5), the cotangent module of weight $-1$ modulo the differentials of the Selmer equations has length at most $2\sum_\xi v(\xi(a))$, the sum over representatives of the torus orbits of the two-signed points $\xi$, with $a=\operatorname{tr}r(\sigma)-t$. This is commutative algebra with no arithmetic input.
\item \emph{Finiteness and automorphy.} The ordinary deformation problem is finite over $O$ (hypothesis (M2)), and every ordinary point $r$ with values in a finite extension of $O$ that is irreducible on every open subgroup of the kernel of the similitude character restricts on $G_K$ to a representation $\rho$ such that $\rho\otimes\chi^{-1}$ is automorphic and conjugate self-dual.
\item \emph{Rigidity.} At each two-signed point the tangent space of the Selmer locus is one-dimensional (hypothesis (M4)), which is the vanishing of the adjoint Bloch--Kato Selmer group. At the reducible fixed points $W\oplus1$ (hypothesis (M3)) the Bloch--Kato Selmer groups of the twist $W$ of a newform of weight $2e+1$ and of $W\otimes\theta$ vanish.
\end{enumerate}
\notgiven{None of (G1) to (G4) is proved here, and (OneChar) is not proved here.}

The closest published theorem to (G3) is Theorem~1.1 of Allen, Newton and Thorne~\cite{ant20}. It allows any number of residual constituents, but its hypothesis~(4) requires each constituent to be conjugate self-dual, and its hypotheses (5) and (6c) require a place at which the seed is an unramified twist of the Steinberg representation. In the configuration of (G3), on $G_K$ the constituents $\Psi$ and $\Psi^{-1}$ are exchanged by conjugate duality, so they are not individually conjugate self-dual when $e\ne0$, and no Steinberg place is assumed. The earlier theorem of Thorne~\cite{thorne15} has stronger hypotheses. So (G3) is not a published theorem. In characteristic zero at irreducible points, (G4) has the shape of Theorem~A of Newton and Thorne~\cite{nt23}, which needs the image to be enormous and does not need residual irreducibility. For the cyclotomic image, the trace criterion in Subsection~\ref{tw:rankfour} applies once the image spans the matrix algebra and its regular semisimple elements are dense.

At a reducible point let $f$ have weight $k=2e+1$. The Frobenius polynomials identify $W$ with $V(f)(e)$ by Chebotarev density and Brauer--Nesbitt. Take a cyclotomic field $F$ on which $\theta$ is trivial. Since $1\le e\le k-1$ and $e\ne k/2$, Kato's Theorem~14.2(1)~\cite{kato04} gives $H^1_f(F,V(f)(e))=0$. Restriction to $F$ is injective in characteristic zero. It gives the two vanishings over $\mathbb Q$ in (G4).

For the rank-three crossing in (G3), write the odd relations of a minimal graded presentation as $My=0$. Multiplication by the adjugate gives $(\det M)y=0$. Thus on the open set where $\det M$ is a unit all the odd coordinates vanish. In the complementary dimension argument the section bound gives dimension at least $3d+2$, whereas the triangular locus has dimension less than $3d$, with $d$ the degree parameter of the presentation. Propagation still requires finiteness over the weight algebra, a point known to be automorphic on the initial component, and a crossing theorem along the successive thickenings at reducible primes.

\begin{remark}[normalizations in (Count) and (OneChar)]\label{w4:rem:normalization}
Let $\sigma_w$ be an arithmetic Frobenius element at $w$, so that $\chi(\sigma_w)=Nw$ and $\psi(\sigma_w)=\psi(Nw)$, and let $\sigma=\sigma_w^{-1}$ be the geometric Frobenius. With $\Psi=\chi^e\psi$ the trace of $\sigma$ on $\Psi\oplus1\oplus\Psi^{-1}$ is $Nw^{-e}\psi(Nw)^{-1}+1+Nw^{e}\psi(Nw)$. With $z_\psi=\psi(Nw)^{-1}$ this is
\[
t(z_\psi),\qquad t(z)=Nw^{-e}z+1+Nw^{e}z^{-1}.
\]
The involution $z\mapsto z^{-1}$ permutes $\mathcal O$. The admissible $\rho\colon G_K\to\mathrm{GL}_3(\overline{\Q}_l)$ in both statements are of two kinds. Either $\rho$ is irreducible, $\rho\otimes\chi^{-1}$ is automorphic and conjugate self-dual, and $\rho$ has level $N$ and Hodge--Tate weights $(e,0,-e)$, or $\rho$ is conjugate to $W|_{G_K}\oplus1$ for an irreducible $W\colon G_\Q\to\mathrm{GL}_2(\overline{\Q}_l)$ with $W\otimes\chi^{-e}$ attached to a newform of weight $2e+1$, $\det W$ of finite order and level $N\lvert d_K\rvert$. In both cases the Frobenius eigenvalues at $w$ have absolute value $1$ under $\iota$, so $\lvert\operatorname{tr}\rho(\sigma)\rvert\le3$, while $\lvert t(z)\rvert\ge Nw^{e}-1-Nw^{-e}>3$ for $p\ge7$ and $e\ge1$. The explicit bound $4c_5\varphi(\mathrm{index})\lvert d_K\rvert^5N^{18}e^7(\log p+\log p')$ quoted in Subsection~\ref{tw:sec:vwin} belongs to a variant of (CorGK) that is not used. The argument uses (CorGK) with an unspecified exponent $A$, through Proposition~\ref{w4:prop:genericplus}.
\end{remark}


\subsubsection{The rank-four counts}\label{tw:rankfour}

Cases (II) and (III) of the plus part use symplectic representations of
rank four. Both use automorphy on every irreducible component through the
residual point. If $A$ is the deformation ring, $J$ the automorphic ideal
and $\mathfrak p$ the prime of the point, the starting condition is
$J\subseteq\mathfrak q$ for one minimal prime $\mathfrak q\subseteq
\mathfrak p$. To obtain the condition for every such $\mathfrak q$ one
joins the components by a chain. At each adjacent pair one needs the
local crossing assertion. Induction along the chain then gives the desired
containment for the last component. Finiteness of the set of systems of
Hecke eigenvalues alone gives no such containment.

At a locally generic automorphic point the following observation verifies
the trace condition for enormity in the adjoint vanishing theorem
of Newton and Thorne~\cite{nt23}. Enlarge the coefficient field so that
the eigenvalues in question split. Let $H$ be the image on the cyclotomic
subgroup. Assume that $H$ spans $M_n$ and that its regular semisimple
elements are dense in its Zariski closure. Let $V\subset M_n$ be a nonzero
simple $H$-submodule under conjugation. For regular $h$ write
$h=\sum_\alpha\alpha e_{h,\alpha}$, with $e_{h,\alpha}$ the spectral
projectors. If $\operatorname{tr}(e_{h,\alpha}v)=0$ for every such $h$,
every $\alpha$ and every $v\in V$, then $\operatorname{tr}(hv)=0$.
By density this holds on $H$ and by linearity on its span. Nondegeneracy
of the trace pairing gives $V=0$, a contradiction. In the symplectic
case the projection from endomorphisms to the symplectic Lie algebra is
$X\mapsto(X-X^\dagger)/2$, with $X^\dagger$ the adjoint for the alternating
form. The finite local condition is the inertia coboundary condition away
from $l$ and the crystalline condition above $l$.

The two counts use the same rigidity bound. In case (III) the prime
$l$ lies outside the bad set, which contains $2$ and $3$, and $l-1\mid 2e$.
Thus $e\ge2$. In the trace separation estimate put $y=p^e$ at an auxiliary
prime $p\ge5$. The required strict inequality is
$7y/2+y+p+1<py$: indeed $(p-9/2)y\ge p^2/2>p+1$.
In the product over conjugate roots of unity of order $k$, one unit
modulo $k$ acts on all three roots at once. The orbit therefore contributes
at most $\varphi(k)$ choices. Repeated roots retain their multiplicities.

\paragraph{The host count.}
In case (V), $l$ belongs to the exceptional set $B_\chi$ and $e=1$. On the branch $R_1\le 2^{c_2(\psi)}$, let $O$ be the coefficient ring with ramification index $E_O$, and normalize $v_l(l)=1$. At distinct two-signed points $r_i$ write
\[
 S=\sum_i v_l\bigl(\operatorname{tr}r_i(\varphi)-\operatorname{tr}r_z(\varphi)\bigr),
\]
where $r_z$ is the Eisenstein representation and $\varphi$ is geometric Frobenius. The host deformation and intersection bounds give $\operatorname{length}_O V\le 2E_OS$ for the Selmer module under consideration. If all trace differences are nonzero and $S\log l\le (lNe)^C$, division by $E_O$ gives
\[
 E_O^{-1}\operatorname{length}_O V\log l\le 2(lNe)^C\le(lNe)^{C+1}.
\]
This deduction requires the host deformation bound and the polynomial estimate for both kinds of representations in Remark~\ref{w4:rem:normalization}.

\subsubsection{Automorphic points and base change}\label{tw:levelpoints}


Let $E/F$ be the CM extension in the level construction, unramified at every finite place. Then $[F:\mathbb Q]\ge2$, since a quadratic field over $\mathbb Q$ ramifies at a finite prime. The definite unitary quotient is compact. Its irreducible summand $\pi$ therefore occurs discretely, with finite-dimensional archimedean type. Labesse's base change (Corollaire~5.3 of \emph{Changement de base CM et s\'eries discr\`etes}) gives one isobaric sum of discrete constituents $\Pi_i$ for $\pi$. At split places it gives the split local component and at spherical places it gives stable Satake transfer. At a nonsplit unramified place with residue cardinality $q$ the upstairs cardinality is $q^2$.

By the residual spectrum theorem of M{\oe}glin and Waldspurger, each $\Pi_i$ is a Speh representation of length $d_i$ on a normalized cuspidal representation of rank $a_i$, with $\sum_i a_id_i=n$. Conjugate self-duality follows for its seed from uniqueness of cuspidal support for the conjugate self-dual constituent. In the normalized induction the exponents are $(d_i-1)/2-j$, $0\le j<d_i$. Reversing their order gives the subrepresentation realization; in the unnormalized constant term the modulus contributes $a_i((d_i-1)/2-j)$ to the reversed exponent. Fix the local factors of these same global constituents at every place. Each prescribed local component belongs to the corresponding induced product as a subquotient: restrict the tensor product to a finite set containing that place and every nonspherical factor, and take its Jordan--Hölder factors. At the remaining places where only a subquotient is prescribed, finite length gives one. The infinitesimal multisets add under induction. For highest weight $\lambda$ their union at infinity consists of $\lambda_j+(n-1)/2-j$, $0\le j<n$.

For each Speh block put $b_{ij}=d_i-2j-n$. Apply the Galois representation theorem for regular algebraic polarized cuspidal representations, in the local compatibility hypotheses of this section, to the cuspidal seeds twisted by $(b_{ij}+a_i-1)/2$. Take the direct sum of the resulting representations. It is continuous, semisimple and of rank $n$. Under normalized local Langlands, parabolic products give sums and absolute-value twists give the corresponding twists of the parameter. The Speh exponents followed by the common twist $(1-n)/2$ thus give precisely the parameter of this sum at split places. This calculation applies to every correspondence satisfying the product and twist laws. On a seed infinitesimal weight $z$ the Hodge--Tate weight is $-z-b_{ij}/2$. The local conditions above $l$ pass to the sum under the spherical or tame hypotheses of the attachment theorem. If the sum is irreducible it consists of a single summand.

Now fix a common Taylor--Wiles base, a permissible auxiliary level, the free module $M$ and its Hecke annihilator $J$. For $m\ne0$ choose a nonzero coordinate in a free basis. Density of admissible weights gives a specialization where that coordinate survives. Decompose into irreducible automorphic summands and project onto a nonzero isotypic component. Take the local factors and the preceding Galois representation on this same summand. The spherical Hecke algebra acts there by its Satake character, so the Hecke polynomial is the characteristic polynomial at inverse Frobenius. The determinant and local types follow on the same summand by the product and twist laws. Projection detects the specialized vector. Varying $m$ gives separating classical points at the chosen level, with common base and annihilator $J$.

In rank two over a totally real field, take a regular algebraic polarized
Hilbert modular representation in the common-parity normalization $(k,w)$.
Skinner~\cite[eq.~(1), Thm.~1]{skinner09} gives the Frobenius-semisimplified
Weil--Deligne parameter with the twist $|\det|^{-1/2}$ and Hodge degrees
$(w-k_\tau)/2$ and $(w+k_\tau-2)/2$. Their sum is $w-1$.
At a place above $l$ an Iwahori-fixed vector gives trivial inertia of the
parameter and hence semistability. In the spherical case the monodromy
also vanishes, giving crystallinity. Away from $l$, purity includes the
monodromy filtration~\cite[Cor.~7]{blasius}. Use inverse Frobenius for the
geometric eigenvalues. If an algebraic character $\eta$ was used to put
$\pi$ in this normalization, the representation for $\pi$ is
$r_{\pi\otimes\eta}\otimes r_\eta^{-1}$. These assertions concern the
rank-two Hilbert modular input.

\subsubsection{The effective cohomology bound}\label{tw:effectiveK}

Let $p$ be a prime, $L$ a totally imaginary number field of degree $d$ and
absolute discriminant $D$, and $S$ the set of places above $p$.
Suppose $L$ contains a primitive $p$th root of unity if $p$ is odd
and a primitive fourth root if $p=2$. Write
\[
 B_C(e,d,D)=\exp((2ed)^C\log(2D)).
\]
The bound required by the arithmetic window argument is uniform: for one
real $C_K$, all such $p,L$ and all integers $e\ge1$, the group
$A=H^2(G_{L,S},\mathbb Z_p(e+1))$ is finite and
$\log |A|\le B_{C_K}(e,d,D)$. In forming the cochains, fix an algebraic
closure, a prolongation of each $v\in S$ to it, and the cyclotomic
character on the quotient unramified outside $S$. The estimate holds for
each such datum.

\paragraph{The localization kernel.}
For the discrete module $M=(\mathbb Q_p/\mathbb Z_p)(-e)$ let
$s_L$ denote the length over $\mathbb Z_p$ of the kernel of
\[
 H^1(G_{L,S},M)\longrightarrow \prod_{v\in S} H^1(G_{L_v},M).
\]
Lengths can initially be infinite. The source identifies this length with
the invariant defined using an absolute Galois group and a fixed algebraic
closure. At the end of the argument one returns from this invariant to
the same local conditions. The two twists involved here are $-e$ in degree
one and $e+1$ in degree two.

\paragraph{Counting in cyclotomic layers.}
Put $k=e-1$. The auxiliary field construction supplies a totally imaginary
field $L'$ with $d'=[L':\mathbb Q]=m d$, maximal cyclotomic depth $c$
($\mu_{p^c}\subset L'$, $\mu_{p^{c+1}}\not\subset L'$), and
\[
 2e+1<\min(p^c,r_2(L')), \quad
 m\le(d+1)(4e+4), \quad p^c\le 2d',\quad
 D'\le D^m(p^c)^{d'},\quad s_L\le s_{L'}.
\]
Here $D'=|\Delta_{L'}|$, $r_2$ counts pairs of complex embeddings, and
$c\ge1$, with $c\ge2$ at $p=2$. The strict inequalities are used before
applying the odd-degree homology estimate.

Indeed, if $c_0$ is the cyclotomic depth of $L$, take the layer
$L'=L(\mu_{p^{c_0+h}})$ with $h$ the least integer for which
$p^h\ge4e+4$. The cyclotomic layer theorem gives $m=p^h$.
Thus $m\le p(4e+4)$, $c\ge c_0+h$, and $r_2(L')=md/2>2e+1$.
The cyclotomic degree over $\mathbb Q$ gives
$p\le d+1$ and $p^c\le 2d'$. To bound the different, let $z$ generate
the layer as a primitive $N$th root, $N=p^{c_0+h}$, and let $P$ be
its relative minimal polynomial. Writing $X^N-1=PQ$ and differentiating
at $z$ gives $N=zP'(z)Q(z)$. The different contains $P'(z)$ and hence
$N$. Its norm divides $N^{d'}$, giving the displayed discriminant bound.
For the length comparison use the tower descent description at levels
$0$ and $h$. These lengths are those of the cokernels of $T-1$ and
$T^{p^h}-1$ for the same twisted operator $T$. The second image lies in
the first by the geometric-sum identity, giving a surjection of
cokernels in the required direction.

In a cyclic Galois layer $F/L'$ let $\zeta$ have order $q=p^\nu$,
$\nu\ge c$, and fix a generator $\gamma$ of order $p^{\nu-c}$
with $\gamma\zeta=\zeta^{1+p^c}$. Set $a=(1+p^c)^e$. Count the
functions $\varphi$ from nonzero prime ideals of $\mathcal O_F$ to
$\mathbb Z/q$ satisfying
\[
 \sum_P\operatorname{ord}_P(\alpha)\varphi(P)=0\quad(\alpha\in F^\times),
 \qquad \varphi(\gamma^{-1}P)=a\varphi(P).
\]
The first condition makes these characters of the ideal class group.
Write $h_\nu$ for their number, $T_j$ for the order of the torsion
subgroup of $H_j(\mathrm{GL}_{4k+10}(\mathcal O_{L'}),\mathbb Z)$, and
\begin{align*}
 b&=2k+2+(2k+3)(c+2k+3), & f&=v_p(e),\\
 g&=2d'+\lfloor\log_2 h_{L'}\rfloor+2, & X&=(b+f)g+c+f .
\end{align*}

Two counts enter. For the first count, consider a finite subset of
$F^\times$ whose elements have
$\operatorname{ord}_P(\alpha)=\mathfrak d(\gamma P)-a\mathfrak d(P)$
for a divisor $\mathfrak d$ depending on $\alpha$. The chosen elements
are distinguished modulo $q$th powers by the $p^b$th powers of their
twisted norms. In the archive convention the norm is
$\prod_{j=0}^{|\gamma|-1}\gamma^j(\alpha)^{a^j}$, with the field
automorphism acting on $\alpha$. Each element gives a cycle
in degree $2k+3$. One chooses the residue tests simultaneously away
from $p$ and the finitely many denominators of these cycles, at primes
where $\mathcal O_{L'}$ surjects onto the residue field. If two cycles
have the same scaled evaluations at every such prime, the ratio of their
powered twisted norms is a $q$th power in every tested residue field.
The power-detection lemma gives a $q$th power in $F^\times$.
The separation condition then identifies the two elements. The image
bound for the resulting primitive family of evaluations bounds the first
count by $q^{r_2(L')}T_{2k+3}T_{2k+2}$.

\paragraph{The cycle and its determinant.}\label{kone:cycle}
Write $A_q=\mathbb Z/q$, $C=\mathbb Z/q$ as an additive group,
$M=2k+3$, $\kappa=1+p^c$ and $a=\kappa^e$, so $e=k+1$.
In the bar complex of $C$ with coefficients in $A_q$ use
\[
 z=\sum_{j_1,\ldots,j_e\in C}[j_1|1|\cdots|j_e|1].
\]
This has degree $2e$. In its boundary the first face sums to zero
because its multiplicity is $q$. The remaining faces cancel in pairs
after shifting a summation variable by one. Put $d_t(x)=(x,tx)$,
$0\le t\le M$. The product-group homology in degree $2e$ with these
coefficients has a generating set of size at most $2e+1=M$ (use the
first projection onto $C$). More explicitly, constants, cyclic
shift minus identity and summation give an exact sequence
from the trivial module through two copies of the cyclic permutation
module to the trivial module. The two connecting sequences and
induction from the kernel bound the generator number in degree $d+2$
by that in degree $d$ plus the kernel-group bounds in degrees $d+1,d+2$.
For a cyclic kernel both bounds are one; degree zero needs one
generator. Hence degree $2e$ needs at most $2e+1$ generators.
The $M+1$ classes $(d_t)_*z$
therefore satisfy a relation $\sum_t n_t(d_t)_*z=\partial B$ with
some $n_{t_0}$ a unit. Indeed their subgroup is a finite abelian
$p$-group with at most $M$ cyclic factors. A redundant generator
modulo $p$ is redundant in the subgroup, by successive lifting
through the powers of $p$. This gives the unit coefficient.

Index the coordinates of $F^{C^2}$ by $j=(j_1,j_2)$.
On coordinate $j$ the diagonal representation $\rho$ of $C^2$
has character $(x,y)\mapsto\zeta^{j_1x+j_2y}$. For $t$ let
$s_t(j)=j_1+tj_2$. The lattice construction uses a finitely supported
integer function $\eta$ on $C^2$. Its sums on the fibers of $s_t$
vanish for $t\ne t_0$. At $t_0$ its pushforward is
\[
 \sum_{i=0}^{M}[X^i]P(X)\,[\kappa^i],\qquad
 P(X)=\prod_{j=0}^{M-1}(X-\kappa^j).
\]
Here $[x]$ denotes a basis vector in $\mathbb Z[C]$. To see the
interpolation, use the product of $[(0,1)]-[(t,0)]$ over $t\ne t_0$
in $\mathbb Z[C^2]$. Each unwanted pushforward kills a factor. At
$t_0$ the product equals a translation times
$\prod_{t\ne t_0}([t_0-t]-1)$. Every $t_0-t$ generates a subgroup
containing $p^c$ in $C$, since $0<|t_0-t|\le M<p^c$. Thus each factor
divides $[p^c]-1$. The desired pushforward is divisible by
$([p^c]-1)^M$: in polynomial coordinates apply
$\prod_{j<M}(U-\kappa^j)$ to $X$, where $Uf(X)=f(X^\kappa)$.
The operator $U-\kappa^j$ carries $(X^{p^c}-1)^j$ times any polynomial
into $(X^{p^c}-1)^{j+1}$ times a polynomial, by factoring
$X^{\kappa p^c}-1$ and reducing the quotient modulo $X^{p^c}-1$.
Lift the quotient in the group ring by $x\mapsto(x,0)$ and undo the
translation. This constructs $\eta$ with both specified properties.

Modify the standard lattice coordinate $j$ by the fractional ideal
with divisor $\eta(j)\mathfrak d$. For each $t$ construct an
$L'$-linear $g_t$ taking the standard lattice to this modified lattice
and commuting with $\rho d_t$. The determinant-one identifications on
blocks use the Steinitz relation for a direct sum of fractional ideals:
one can move the sum of two divisors to one coordinate, and iterate.
Scaling one distinguished coordinate realizes a prescribed determinant
when its divisor is the difference of the two total divisors.
For $t\ne t_0$ all block determinants can be taken as one. At $t_0$
the construction groups coordinates in Galois orbits and leaves the
determinant
\[
 \beta=\prod_i\gamma^{-i}(\alpha)^{[X^i]P_e(X)},\qquad
 P_e(X)=\prod_{0\le j<M,\ j\ne e}(X-\kappa^j)
\]
at character $1$. Indeed the divisor of $\gamma^{-i}\alpha$ at $P'$
is $\mathfrak d(\gamma^{i+1}P')-a\mathfrak d(\gamma^iP')$.
Multiplying this difference by the coefficients of $P_e$ gives
precisely the divisor prescribed by $P=(X-a)P_e$. This explains
both the inverse action in $\beta$ and the missing factor in $P_e$.

Let $H$ be the stabilizer of the standard lattice inside the localized
group containing all $g_t$ and $\rho(C^2)$. One base-field denominator
excludes the support of the divisors and the bad primes of the splitting
map. Set $\rho'=g_{t_0}^{-1}\rho g_{t_0}$ and $h_t=g_{t_0}^{-1}g_t$.
Both representations take values in $H$, and $h_t\in H$ since both
maps $g_t,g_{t_0}$ take the standard lattice to the same lattice.
For the bar prism $P_h$ one has
$\partial P_h+P_h\partial=(\mathrm{conj}_{h^{-1}})_*-1$.
Consequently the chain
\[
 Y=\rho'_*B-\rho_*B+\sum_t n_t P_{h_t}((\rho'd_t)_*z)
\]
is closed in $H$. In fact $h_t^{-1}\rho'd_th_t=\rho d_t$, so the
boundary of the sum is $\rho_*\partial B-\rho'_*\partial B$.
In the localized group $Y$ is homologous to
$\sum_t n_t P_{g_t}((\rho d_t)_*z)$. This follows by applying the
prism identity to $\rho_*B$ and the prism product identity to
$g_t=g_{t_0}h_t$ on each cycle. Multiply $Y$ by $n_{t_0}^{-1}$.
The standard lattice is finite projective over $\mathcal O_{L'}$.
Adding a projective complement with trivial action realizes this cycle
in some $\mathrm{GL}_N(\mathcal O_{L'})$.

At an allowed residue field $k_0$, restriction of scalars splits into
the conjugate blocks indexed by $\sigma\in\mathrm{Gal}(F/L')$.
For a character $v:k_0^\times\to A_q$, the finite-field functional
on a scalar-character block of weight $x$ evaluates a prism for $g$
as $x^e v(\det g)$. On the scalar line this is the product of $e$
cyclic carry cocycles with $v$: each consecutive pair in $z$ sums to
$x$, and insertion in the prism leaves $v(g)$. Transfer over projective
lines gives the additive functional on vector spaces. The line count
for direct sums uses $|k_0|=1$ in $A_q$, since the reduced $\zeta$ has
order $q$. Additivity kills the trivial complement. Blocks of determinant
one give zero. In the surviving block let $\sigma\zeta=\zeta^{\chi(\sigma)}$.
Summing the conjugate blocks gives
\[
 \sum_\sigma \chi(\sigma)^e v(\overline{\sigma\beta})
   =P_e(a)\,v\big(\overline{\operatorname{nrm}(\alpha)}\big).
\]
Indeed replacing $\alpha$ by $\gamma^{-i}\alpha$ multiplies the
weighted sum by $a^i$ after reindexing, and the sum for $\alpha$ is
the evaluation of its twisted norm. The unit coefficient cancels the
normalization of $Y$. For a cyclic residue unit group take $v$ with
kernel the $q$th powers. Each factor $\kappa^e-\kappa^j$ of $P_e(a)$
has valuation $c+v_p(|e-j|)\le c+M$, by lifting the exponent
(with $c\ge2$ at $p=2$). In $A_q$, $P_e(a)$ therefore divides
$p^{M(c+M)}$. Equality of two evaluations scaled by $p^{2e}$
implies that the ratio of their twisted norms raised to
$p^{2e+M(c+M)}=p^b$ is a $q$th power at this prime.

\paragraph{Detecting powers and bounding the image.}\label{kone:rank}
If $\xi\in F^\times$ passes all these residue tests, choose
$x$ with $x^q=\xi$ in an algebraic closure. For each automorphism
$g$ over $F$, the condition $\sigma x=gx$ describes an open coset.
Degree-one Chebotarev gives a Frobenius in this coset outside the finite
exclusion set and the primes dividing $q$. Its residue field has prime
order, so reduction from $\mathcal O_{L'}$ is surjective. In the residue
field $\xi$ has a $q$th root and all the $q$th roots of unity are
present. Hence the residue of $x$ is fixed by Frobenius. The quotient
$\sigma x/x$ is a $q$th root of unity congruent to one. If it were
different from one, its geometric sum of length $q$ would be zero,
giving $q=0$ in the residue field. Thus $gx=x$ for every $g$ and
$x\in F$. Applying this argument to the ratio proves injectivity
for the simultaneous evaluations of the finite set.

For the image bound put $m_1=2e$, and stabilize the finitely many
cycles to rank $n\ge2(m_1+1)+4$. An integral lift $Z$ of a
mod-$q$ cycle has $\partial Z=qV$. The homology class of $V$ is
$q$-torsion in degree $m_1$, independent of the lift. When this class
vanishes one can subtract $q$ times an integral chain from $Z$ to
make it a cycle. There are at most $T_{m_1}$ fibers of this obstruction
by homological stability. Differences in a fiber therefore reduce the
count to evaluations on integral cycles.

Write $\psi_p$ for the $p$-fold block sum on integral homology and
$i$ for stabilization to the same rank. Starting with $z_0=x$ put
$z_{t+1}=(\psi_p)_*z_t-p^{t+2}i_*z_t$. A primitive evaluation
$\phi$ is stable and additive on block sums. Hence
$\phi(z_{t+1})=p(1-p^{t+1})\phi(z_t)$. The factors $1-p^{t+1}$
are units modulo $q$. The scaled evaluations $p^{m_1}\phi(x)$ thus
factor through the integral span of the $z_{m_1}$. Its rank is at
most $r_2(L')$. To obtain this last inequality the source passes to
real cycles and pairs with one odd-degree functional at each infinite
place. A class in the common kernel, if nonzero, survives stabilization
and is separated by a closed cochain. The surjectivity theorem represents
that cochain on cycles by an invariant form. Its weight decomposition
consists of the place functionals and terms of weights $j=2,\ldots,m_1+1$.
The latter vanish on the displayed span: the factor at $t=j-2$ is
$p^j-p^j=0$. The former vanish by the kernel condition. This contradicts
separation. The real image thus has dimension at most the number of
complex places, giving the integral rank bound.

For any finite set of evaluations use the finitely generated subgroup
spanned by lifts of its elements. Its free part contributes at most
$q^{r_2(L')}$ values. Its torsion contributes at most $T_{m_1+1}$
by stability, also at the enlarged rank for $z_{m_1}$. Multiplication
by the number of obstruction fibers gives the claimed bound.

\input{sec17_line_transfer}


For the second count write $Y_F=F^\times$ additively, let $I_F$ be the
divisor group, and let $\partial:Y_F\to I_F$ be the divisor map.
On elements put $\delta=\gamma^{-1}-a$; on divisors the corresponding
operator is $(\delta_I d)(P)=d(\gamma P)-a d(P)$. Put
\[
 U=\partial^{-1}(\delta_I I_F),\quad
 W=U/(\delta Y_F+Y_{F,\mathrm{tors}}),\quad
 C'=I_F/(\partial Y_F+\delta_I I_F).
\]
Let $E'$ be $\ker\partial$ modulo $\delta(\ker\partial)$ and torsion.
The divisor and unit counts give finite groups with $|W|=|E'||C'|$,
$h_\nu\le p^{v_p(|C'|)}$, and
\[
 v_p(|E'|)+c+f=r_2(L')(\nu+f).
\]
The last equality uses $v_p(a-1)=c+f$ and
$v_p(a^{p^{\nu-c}}-1)=\nu+f$.

Here are the two order computations. On $I_F$ the operator $\delta_I$
is injective, since an element in its kernel satisfies
$(a^{|\gamma|}-1)d=0$ and the divisor group is torsion free.
For $y\in U$ the equation $\partial y=\delta_I d$ now determines $d$.
Sending $y$ to the ideal class of $d$ maps $W$ onto the kernel of
$\delta_I$ on the class group. Indeed a class in that kernel has
$\delta_I d$ principal, and altering $y$ by $\delta z$ alters $d$ by
$\partial z$. If $d$ is principal one can therefore take $y$ to be
a unit. Two such units agree in $W$ precisely modulo the twisted
differences of units and torsion, again by injectivity on divisors.
Thus the kernel is $E'$. The kernel and cokernel of an endomorphism
of a finite group have the same cardinality, and here the cokernel
is $C'$. Equivariant class characters factor through $C'$; on each
cyclic $p$-summand of order $p^i$ there are at most $p^i$ choices.

The logarithmic unit lattice spans the real hyperplane of coordinates
at infinite places with sum zero. On all coordinates $\gamma^{-1}$
permutes $r_2(L')$ cycles of length $n=p^{\nu-c}$, since each complex
place downstairs has $n$ lifts. The absolute determinant of
$\gamma^{-1}-a$ on the permutation space is $(a^n-1)^{r_2(L')}$.
Its absolute determinant on the one-dimensional quotient is $a-1$.
The determinant on the lattice modulo roots of unity therefore gives
$|E'|(a-1)=(a^n-1)^{r_2(L')}$, as required.

The powered twisted norm induces a map
$W\to Y_F/(qY_F+Y_{F,\mathrm{tors}})$. Its image has size at most
$B=q^{r_2(L')}T_{2k+3}T_{2k+2}$ by choosing one representative for
each image element and applying the first count. Write
$a^{p^{\nu-c}}-1=p^{\nu+f}m_0$ with $p\nmid m_0$.
The norm identity kills the kernel by $p^{b+f}m_0$.
Indeed, write $g=\gamma^{-1}$, $n=|\gamma|$, $N_a=\sum_{j<n}a^jg^{n-j}$.
Then $(g-a)N_a=N_a(g-a)=(1-a^n)g$. If $p^bN_a y=qz+t$
with $t$ torsion, applying $g-a$ and dividing by $q$ modulo torsion
shows that $p^{b+f}m_0g y$ lies in $(g-a)Y_F+Y_{F,\mathrm{tors}}$.
This subgroup is invariant under $g$ and $g^{-1}$, proving the assertion.
On its $p$-primary subgroup this gives an exponent bound $p^{b+f}$.
The Selmer injection and the base-field Selmer count bound $|W[p]|$
by $p^g$. A finite abelian $p$-group killed by $p^{b+f}$, with
at most $p^g$ elements killed by $p$, has order at most $p^{(b+f)g}$.
Indeed multiplication by $p$ takes the subgroup killed by $p^{j+1}$
into the subgroup killed by $p^j$ with kernel of order at most $p^g$.
Induction from $j=0$ gives the bound. Applying it to the kernel gives
$p^{v_p(|W|)}\le Bp^{(b+f)g}$. In particular
\[
 h_\nu q^{r_2(L')}\le
 p^{v_p(|C'|)+v_p(|E'|)+c+f}
 =p^{v_p(|W|)+c+f}\le Bp^X.
\]
For the Selmer bound just used, lift $w\in W[p]$ to $u\in U$ and write
$pu=\delta y+t$ with $t$ torsion. If $\partial u=\delta_I d$, then
$\partial y=pd$ by injectivity of $\delta_I$. Also $gy-y$ vanishes
modulo $pY_F+$ torsion, since $p\mid a-1$. The class of the chosen
$y$ in that quotient distinguishes the $w$: equality of two such
classes gives $y-y'=pz+t'$, whence
$p(u-u'-\delta z)$ is torsion, as is $u-u'-\delta z$ itself.
To descend these classes, fix a generator $\rho$ of the roots of unity
of $F$. In multiplicative notation a representative $\alpha$ has
$\gamma\alpha/\alpha=v^p\rho^r$ for $0\le r<p$ after absorbing a
$p$th power. The pair $(r,N_{F/L'}v)$ takes at most $p^2$ values
because $(N_{F/L'}v)^p=(N_{F/L'}\rho)^{-r}$. Within a fiber the ratio
of two such $v$ has norm one. Hilbert 90 supplies a coboundary, so
after removal of a $p$th power the ratio of the two $\alpha$ descends
to $L'$. Outside $p$ the cyclotomic extension is unramified, so the
valuations of the descended element are divisible by $p$ there.
The resulting base-field classes modulo powers have size at most
$p^{2d'+\lfloor\log_2 h_{L'}\rfloor}$. Indeed valuations modulo $p$
at places over $p$ cost at most $p^{d'}$. With these fixed, taking
the $p$th root of the principal divisor costs at most
$|\mathrm{Cl}(L')[p]|\le p^{\lfloor\log_2 h_{L'}\rfloor}$.
The remaining units modulo $p$th powers cost at most $p^{d'}$
by the unit theorem. This gives the stated $p^g$.

Substituting the value of $B$ gives
\[
 h_\nu q^{r_2(L')}
 \ \le\ q^{r_2(L')}T_{2k+3}T_{2k+2}\,p^X.
\]
Cancel the positive common factor. All quantities on the right of
$h_\nu\le T_{2k+3}T_{2k+2}p^X$ are independent of $\nu$.
This uniformity allows passage back to the localization kernel:
restriction to the infinite cyclotomic tower is injective on that kernel,
and each finite set in its image is detected at a common finite level.
At that level the corresponding torsion class characters inject into
the functions counted above. Thus every such finite set has size at most
$T_{2k+3}T_{2k+2}p^X$. It follows that the kernel is finite and
\[
 s_L\log p\le s_{L'}\log p
 \le\log T_{2k+3}+\log T_{2k+2}+X\log p.
\]

For clarity on the constants, write $u=2ed'$ and $Y=\log(2D')$.
For the torsion estimate set $N=4e+6$ and $d_* = d'N$.
The integral face complex for $\mathrm{GL}_N(\mathcal O_{L'})$ is
exact in the range used to compute $H_j$. In dimensions up to $j$
its number of orbits is bounded by
$A_*=(2D'd_*)^{12(j+1)d_*^3N}$ and the orders of the stabilizers by
$s_*=3^{d_*N}$. The boundary from dimension $j+1$ is the alternating
sum of $j+2$ faces, hence has column $\ell^1$-norm at most $j+2$.
The equivariant chain estimate with these bounds gives
\[
 T_j\le (j+2)^{A_*}\,s_*^{\,jA_*s_*^j}.
\]
This estimate uses exactness of the augmented complex, orbit counting
for integral lattices, and the bounds for homology of finite stabilizers.
Taking logarithms gives
$A_*\log(j+2)+jA_*s_*^j\log s_*$. The numerical estimate for this
expression at $j\le2e+1$ gives
$\log T_j\le \exp(u^{C_0}Y)$ with $C_0=24$ in the two degrees above. The numerical
comparison for the layer gives $X\log p\le u^9Y$. Absorbing the three
terms and descending using the degree and discriminant bounds gives
$s_L\log p\le B_C(e,d,D)$ with $C=5(\max(C_0,9)+2)+11$.

\paragraph{From the localization kernel to the required order.}
Poitou--Tate duality and the local finiteness and duality inputs give
\[
 \operatorname{length}_{\mathbb Z_p} A
 \le s_L+\sum_{v\in S}\operatorname{length}_{\mathbb Z_p}H^0(G_{L_v},M).
\]
For each $v$ the cyclotomic image $U_v\subset\mathbb Z_p^\times$ has
index $I_v\le d$: it contains the inertia image, whose index is at most
$d$. The element $(1+p^2)^{I_v}$ lies in $U_v$.
An invariant of order $p^t$ forces
$p^t\mid (1+p^2)^{I_v e}-1$. The valuation of the latter is
$2+v_p(I_v)+v_p(e)$, also for $p=2$ since the base is $1\bmod 4$.
Consequently the invariants form a subgroup of a cyclic group of order
at most $p^{2+\lfloor\log_p d\rfloor+e}$. There are at most $d$ places
above $p$, so
\[
 \operatorname{length}_{\mathbb Z_p} A
 \le s_L+d(2+\lfloor\log_p d\rfloor+e).
\]
A finite-length $\mathbb Z_p$-module of length $t$ has order $p^t$,
by a composition series with successive quotients $\mathbb F_p$.
This proves finiteness. The root-of-unity hypothesis gives $p\le d+1$,
and $\lfloor\log_p d\rfloor\log p\le\log(d+1)$. Hence the extra term
in $\log|A|$ is at most $d(3+e)\log(d+1)$. For $e,d,D\ge1$ this is
at most $4d\log(2D)(2e)(2e+3)\log(d+1)$, which the final numerical
comparison absorbs into $B_{\max(C,0)+9}(e,d,D)$ together with $B_C$.
Following the exhibited witnesses gives $5(24+2)+11=141$ at the
localization-kernel step and $141+9=150$ here. The interface used for
the window only asks for the existence of one absolute exponent.

\paragraph{Scope of this reconstruction.}
This deduction uses the arithmetic and analytic inputs to the odd-degree
rank bound, the integral torsion estimate, the Selmer and class-group
comparisons, and the tower and duality theorems in their indicated
normalizations. Paragraph~\ref{kone:lines} gives the line transfer used in the cycle
argument. A standalone derivation still requires the integral lattice
and splitting results and the arithmetic and analytic inputs above. In the window argument
the same uniform bound is supplied to the case estimates together with the
separate rank-two Selmer input.

\paragraph{Use of the bound on the selected family.}\label{kone:consumer}
The window argument takes the existential pair $(C_K,\text{cohomology
bound at }C_K)$ once. It uses that same pair in the count of type I,
in the joint count of types II and III, and in the count of type V.
Type I also uses the rank-two Selmer estimate and the unramified
automorphy input. Types II and III use the symplectic input, while
type V uses the reduced unitary lifting input. The four counts and
the common pair yield the window estimate of
Hypothesis~\ref{w4:thm:vwin} in the source normalization.
At each selected index the window supplies finiteness of the odd
localization kernels and bounds the sum of their lengths, weighted by
$\log p$, by $\sqrt b$. Indeed the record identifies its finite
lengths with the cohomological ones at every $1\le e<t$; nonprime
indices contribute zero, and summation gives the required inequality.
With the local lattice comparison and equality of the small-level
ranks at that index this gives the twisted-place divisibility used
in the division estimate of Subsection~\ref{sec:frontier}.
These rank and automorphy hypotheses enter independently of the
cohomological order bound.



\subsection{The congruence bound}\label{tw:sec:linkn}

\begin{hypothesis}[congruence bound]\label{w4:thm:linkn}
There are $C$ and $x_0$ such that for every admissible datum at scale $x\ge x_0$ for which every $\Sha_l(E,e)$ with $0<e<t_r$ is finite, and for every lift index $j$,
\[
\sum_{w\ \mathrm{twisted}}\ell_we_w\,\eta_j(w)\ \le\ d^0_j\Bigl(C\bigl(f\,m\,b\,L_\varepsilon^2+(\log x)^2N_c\bigr)+\sum_l\log l\sum_{1\le e<t_r}\len\Sha_l(E,e)\Bigr),
\]
provided the last sum converges.
\end{hypothesis}

The argument for Hypothesis~\ref{w4:thm:linkn} has three steps. A lattice defect is bounded by congruence numbers (step one), congruence numbers by Fitting lengths of a cocycle module (step two), and Fitting lengths by Selmer lengths plus local terms (step three).

Let $O$ be a complete discrete valuation ring with fraction field $K$, $X$ a finite set, $V=\bigoplus_{x\in X}V_x$ a finite-dimensional $K$-vector space, and $\Lambda\subset V$ an $O$-lattice. Let $T\subset\prod_{x\in X}O$ be an $O$-subalgebra acting on $V$ by $t\cdot v=(t_xv_x)_x$ with $T\Lambda\subset\Lambda$. For $x\in X$ define the \emph{congruence ideal} $\mathfrak c_x=\{t_x: t\in T,\ t_y=0\text{ for all }y\ne x\}$ and $\delta_x=\min\{v(a):a\in\mathfrak c_x\}\in\mathbb N\cup\{\infty\}$. For $Z\subset X$ write $V_Z=\bigoplus_{x\in Z}V_x$ and $\pi_Z\colon V\to V_Z$ for the projection.

\begin{proposition}\label{w4:prop:cong}
For $Z\subset X$, put $M^{\rm up}=\pi_Z(\Lambda)$ and $M^{\rm low}=\Lambda\cap V_Z$. Then
\begin{enumerate}[leftmargin=*,label=(\roman*)]
\item $M^{\rm up}/M^{\rm low}\cong\Lambda/\bigl((\Lambda\cap V_Z)\oplus(\Lambda\cap V_{X\setminus Z})\bigr)$;
\item $\len_O(M^{\rm up}/M^{\rm low})\le\sum_{x\in Z}\dim_KV_x\cdot\delta_x$.
\end{enumerate}
\end{proposition}

\begin{proof}
(i) The map $\Lambda\to M^{\rm up}/M^{\rm low}$, $v\mapsto\pi_Z(v)$, is surjective. Its kernel is the set of $v\in\Lambda$ with $\pi_Z(v)\in\Lambda$. For such $v$, $\pi_{X\setminus Z}(v)=v-\pi_Z(v)\in\Lambda$, so $v\in(\Lambda\cap V_Z)\oplus(\Lambda\cap V_{X\setminus Z})$, and the converse inclusion is clear.

(ii) If some $\delta_x=\infty$ there is nothing to prove. Otherwise choose $a_x\in\mathfrak c_x$ with $v(a_x)=\delta_x$ and $t^{(x)}\in T$ with $t^{(x)}_x=a_x$ and $t^{(x)}_y=0$ for $y\ne x$. For $v\in\Lambda$, $t^{(x)}v=a_x\pi_x(v)$ lies in $\Lambda\cap V_x$. Hence
\[
\bigoplus_{x\in Z}a_x\pi_x(\Lambda)\ \subset\ \bigoplus_{x\in Z}(\Lambda\cap V_x)\ \subset\ M^{\rm low}\ \subset\ M^{\rm up}\ \subset\ \bigoplus_{x\in Z}\pi_x(\Lambda).
\]
Each $\pi_x(\Lambda)$ is a lattice in $V_x$, free of rank $\dim V_x$, so $\len(\pi_x(\Lambda)/a_x\pi_x(\Lambda))=\dim V_x\cdot\delta_x$. Therefore $\len(M^{\rm up}/M^{\rm low})\le\len\bigl(\bigoplus_{x\in Z}\pi_x(\Lambda)/\bigoplus_{x\in Z}a_x\pi_x(\Lambda)\bigr)=\sum_{x\in Z}\dim V_x\cdot\delta_x$.
\end{proof}

Proposition~\ref{lk:prop:ident} identifies this length with the summand in Definition~\ref{r5:def:eta} after division by $\len(O/l)$.

\begin{lemma}\label{w4:lem:det}
Let $Q\in M_n(O)$. Then $v(\det Q)=\len_O\operatorname{coker}Q$, and if $\det Q\ne0$, $\len_O\operatorname{coker}Q=\len_O\ker\bigl(Q\colon(K/O)^n\to(K/O)^n\bigr)$. More generally, if $P\colon O^m\to O^k$ is injective after $\otimes K$, then $\ker(P\otimes K/O)\cong\operatorname{coker}(P)_{\rm tors}$.
\end{lemma}

\begin{proof}
By the elementary divisor theorem $Q=U\operatorname{diag}(\varpi^{a_1},\dots,\varpi^{a_n})U'$ with $U,U'$ invertible; both sides of the first two equalities equal $\sum a_i$ (with $\infty$ allowed). For the last statement, tensor $0\to O^m\to O^k\to\operatorname{coker}P\to0$ with $K/O$: since $O^k$ is flat, $\ker(P\otimes K/O)=\operatorname{Tor}_1^O(\operatorname{coker}P,K/O)$, which is the torsion submodule of $\operatorname{coker}P$.
\end{proof}

Let $R'$ be an $O$-algebra, not necessarily commutative, and $P\in M_{a\times b}(R')$. For an $O$-algebra map $\rho\colon R'\to M_d(O)$ write $\rho(P)\in M_{ad\times bd}(O)$ for the block matrix acting $O^{bd}=(O^d)^b\to(O^d)^a$. Then $\Hom_{R'}(\operatorname{coker}(R'^a\xrightarrow{\cdot P}R'^b),V)=\ker(\rho(P)\colon V^b\to V^a)$ for $V=K^d$ or $(K/O)^d$ with $R'$ acting through $\rho$, since an $R'$-linear map from $R'^b$ is given by the images $u_1,\dots,u_b$ of the basis, and it kills the image of $\cdot P$ exactly when $\sum_jP_{ij}u_j=0$ for all $i$.

\begin{proposition}[Fitting inequality]\label{w4:prop:fitting}
Let $T\subset\prod_{x\in X}O$ and $\mathfrak c_x$, $\delta_x$ be as above. Let $R'$ be an $O$-algebra with $O$-algebra maps $\rho_x\colon R'\to M_{d_x}(O)$ for $x\in X$, and $P\in M_{a\times b}(R')$ with $a\ge b$. Assume:
\begin{enumerate}[leftmargin=*,label=(F\arabic*)]
\item for every $Q\in M_b(R')$ the function $x\mapsto\det\rho_x(Q)$ lies in $T$;
\item for every $x\ne x_0$, $\ker(\rho_x(P)\otimes K)\ne0$, that is, there is a nonzero $R'$-linear map $\operatorname{coker}P\to K^{d_x}$;
\item in a suitable $O$-basis, $\rho_{x_0}$ is block upper triangular with diagonal blocks $\rho_1,\dots,\rho_g$, $\rho_k\colon R'\to M_{n_k}(O)$, and the map $R'\to\prod_kM_{n_k}(k_O)$, $r\mapsto(\bar\rho_k(r))_k$, is surjective, where $k_O$ is the residue field;
\item $\lvert k_O\rvert>b\sum_kn_k$.
\end{enumerate}
Then
\[
\delta_{x_0}\ \le\ \sum_{k=1}^g\len_O\Hom_{R'}\bigl(\operatorname{coker}P,(K/O)^{n_k}\bigr)=\sum_{k=1}^g\len_O\ker\bigl(\rho_k(P)\ \text{on}\ ((K/O)^{n_k})^b\bigr).
\]
\end{proposition}

\begin{proof}
If some term on the right is infinite there is nothing to prove, so assume all are finite. Then each $\rho_k(P)\otimes K$ is injective: a nonzero $K$-kernel vector, scaled, gives an injection of $K/O$ into the kernel on $(K/O)^{bn_k}$, of infinite length. Let $W_k\subset O^{an_k}$ be the saturation of the image of $\rho_k(P)\colon O^{bn_k}\to O^{an_k}$, a direct summand of rank $bn_k$.

\emph{Choice of $C$.} For $C\in M_{b\times a}(R')$ put $Q=CP\in M_b(R')$. Then $\rho_k(Q)=\rho_k(C)\rho_k(P)$. There is $C$ such that for every $k$ the map $\rho_k(C)\colon O^{an_k}\to O^{bn_k}$ restricts to an isomorphism $W_k\to O^{bn_k}$. By Nakayama it suffices that $\bar\rho_k(C)$ restricts to an isomorphism $\overline W_k\to k_O^{bn_k}$, where $\overline W_k=W_k/\varpi W_k\subset k_O^{an_k}$ has dimension $bn_k$ ($W_k$ is a direct summand). By (F3), as $C$ ranges over $M_{b\times a}(R')$ the tuple $(\bar\rho_k(C))_k$ ranges over all of $\prod_kM_{bn_k\times an_k}(k_O)$, choosing each entry of $C$ independently. Fix bases of the $\overline W_k$. The function $\Phi(Y_1,\dots,Y_g)=\prod_k\det(Y_k|_{\overline W_k})$ is a polynomial on $\prod_kM_{bn_k\times an_k}(k_O)$ of total degree $\sum_kbn_k$. It is not the zero polynomial, because for each $k$ a matrix $Y_k$ that maps $\overline W_k$ isomorphically onto $k_O^{bn_k}$ exists over $k_O$. A nonzero polynomial of degree $D$ over a field with more than $D$ elements has a non-zero value at a point with coordinates in that field (induction on the number of variables). By (F4) there is such a point, and by (F3) it is $(\bar\rho_k(C))_k$ for some $C$.

\emph{The function.} Let $f(x)=\det\rho_x(Q)$. By (F1), $f\in T$. For $x\ne x_0$, $\ker\rho_x(Q)\otimes K\supseteq\ker\rho_x(P)\otimes K\ne0$ by (F2), so $f(x)=0$. Hence $f(x_0)\in\mathfrak c_{x_0}$ and $\delta_{x_0}\le v(f(x_0))$.

\emph{The value at $x_0$.} After reordering the basis, $\rho_{x_0}(Q)$ is block upper triangular with diagonal blocks $\rho_k(Q)$, so $f(x_0)=\prod_k\det\rho_k(Q)$. For each $k$, $\rho_k(C)$ maps $\operatorname{im}\rho_k(P)\subset W_k$ into $O^{bn_k}$ with $O^{bn_k}/\rho_k(C)\operatorname{im}\rho_k(P)\cong W_k/\operatorname{im}\rho_k(P)=\operatorname{coker}(\rho_k(P))_{\rm tors}$. By Lemma~\ref{w4:lem:det}, $v(\det\rho_k(Q))=\len\operatorname{coker}\rho_k(Q)=\len\operatorname{coker}(\rho_k(P))_{\rm tors}=\len\ker(\rho_k(P)\otimes K/O)$. Summing over $k$ gives the claim.
\end{proof}

In the application $R'$ is the algebra generated by the Galois group acting on an integral model of the representation attached to an eigensystem, $P$ presents the universal cocycle module of the reducible representation $\rho_\pi\oplus\bigoplus_c\chi^{c}$ with its local relations, and (F2) is a nonzero map of Ribet type at every eigensystem outside the congruence class. The application of Proposition~\ref{w4:prop:fitting} also needs a descent along a map of Hecke algebras and the joint residual surjectivity (F3) for the blocks $\rho_\pi,\chi^1,\dots,\chi^{t_r}$. \notgiven{The descent, the verification of (F2) and of (F3) for the blocks that occur, and step three are not proved here.} Step three consists of the following statements.
\begin{enumerate}[leftmargin=*,itemsep=2pt]
\item \emph{Block bounds.} For a rank one block with character $\chi^c$, the length of the space of homomorphisms from the cocycle module to $K/O$ is at most $\len\Sha_l(E,c-1)$ plus local defects at the places above $l$, at the auxiliary places and at the column places. For the block $\chi^1$ the Selmer term is $v_l(h_E)$.
\item \emph{The defect at $l$.} A lifting theorem for crystalline classes along surjections of crystalline lattices with Hodge--Tate weights in $[-h,h]$, proved with Wach modules, whose loss is $O(h/(l-1)^{j+1}+g\ell_{\max}+\log_lh)$ with main term $hl/(l-1)^2$.
\item \emph{The twist.} A ray class character of auxiliary conductor that removes the block $\rho_\pi$.
\item \emph{Twisted theta series.} Rationality of twisted theta series, their $l$-integrality with bounded denominators, the action of Hecke operators on integral models, and a reciprocity law of Shimura type.
\item \emph{The sums} over the eigensystems $\pi$, the places $w$ and the primes $l$, including the rows at $2$ and $3$.
\end{enumerate}
The published inputs of these statements include the Galois representations attached to regular algebraic polarized cuspidal automorphic representations of $\mathrm{GL}_n$ over CM fields with local--global compatibility at all places~\cite{blggt14}, the local Langlands correspondence for $\mathrm{GL}_n$, Poitou--Tate duality and the Euler characteristic formulas~\cite{nsw}, the theory of Bloch and Kato for $H^1_f$, Fontaine's period rings, Wach modules and Berger's description of Wach lattices, the Herr complex, Berger's theorem that de Rham representations are potentially semistable~\cite{berger02}, class field theory and the Chebotarev density theorem.

\subsubsection*{The local Hecke setting}

Throughout this subsection $O$ is a complete discrete valuation ring with fraction field $K$, uniformizer $\varpi$, normalized valuation $v$ and residue field $k_O$, and $l$ is the residue characteristic, so that $v(l)=\len_O(O/lO)$. We write $e_O=v(l)$. Let $X$ be a finite set, let $\mathbb T\subset\prod_{x\in X}O$ be an $O$-subalgebra which separates the points of $X$, let $M$ be a free $O$-module of finite rank with an $O$-algebra map $\mathbb T\to\operatorname{End}_O(M)$, and put $V=M\otimes_OK$, in which $M$ is a lattice. For $x\in X$ put $V_x=\{u\in V:tu=t(x)u\ \text{for all }t\in\mathbb T\}$ and $\operatorname{mult}(x)=\dim_KV_x$. For $Y\subset X$ write $V_Y=\sum_{x\in Y}V_x$. For $Y\subset Z\subset X$ put
\[
\mathfrak c_Z=\{t\in\mathbb T:t(y)=0\text{ for all }y\notin Z\},\qquad \kappa_Z(x)=\min\{v(t(x)):t\in\mathfrak c_Z\}\in\mathbb N\cup\{\infty\},
\]
\[
M^{\rm up}_Y=V_Y\cap(M+V_{X\setminus Y}),\qquad M^{\rm low}_Y=M\cap V_Y,\qquad \eta(M,Y)=\frac{\len_O(M^{\rm up}_Y/M^{\rm low}_Y)}{e_O}.
\]
The ideal $\mathfrak c_{\{x\}}$ is the congruence ideal $\mathfrak c_x$ of Proposition~\ref{w4:prop:cong}, and $\kappa_Z(x)\le\kappa_{\{x\}}(x)=\delta_x$ for $x\in Z$.

\begin{lemma}\label{lk:lem:split}
There are $C\in O\setminus\{0\}$ and $e_x\in\mathbb T$ for $x\in X$ with $e_x(y)=C$ if $y=x$ and $e_x(y)=0$ otherwise. The operators $P_x=C^{-1}e_x$ on $V$ are orthogonal idempotents with $\sum_xP_x=1$ and $P_xV=V_x$. In particular $V=\bigoplus_{x\in X}V_x$, and for $Y\subset X$ the projection $\pi_Y=\sum_{x\in Y}P_x$ of $V$ onto $V_Y$ along $V_{X\setminus Y}$ commutes with $\mathbb T$.
\end{lemma}

\begin{proof}
For $x\ne y$ choose $t\in\mathbb T$ with $t(x)\ne t(y)$ and put $s_{xy}=t-t(y)\cdot1\in\mathbb T$, so that $s_{xy}(y)=0\ne s_{xy}(x)$. The product $f_x=\prod_{y\ne x}s_{xy}\in\mathbb T$ has $f_x(x)=c_x\ne0$ and vanishes off $x$. With $C=\prod_xc_x$ the elements $e_x=(C/c_x)f_x$ lie in $\mathbb T$ and have the stated values. As functions on $X$, $e_xe_y=\delta_{xy}Ce_x$, $\sum_xe_x=C\cdot1$ and $te_x=t(x)e_x$ for $t\in\mathbb T$. Applying the algebra map to $\operatorname{End}_K(V)$ gives $P_xP_y=\delta_{xy}P_x$, $\sum_xP_x=1$, and $tP_x=t(x)P_x$, so $P_xV\subset V_x$. Conversely $P_x$ acts on $V_x$ as $C^{-1}e_x(x)=1$, so $V_x\subset P_xV$.
\end{proof}

\begin{lemma}\label{lk:lem:proj}
For $Y\subset X$, $M^{\rm up}_Y=\pi_Y(M)$. Hence $M^{\rm up}_Y/M^{\rm low}_Y$ is the module $M^{\rm up}/M^{\rm low}$ of Proposition~\ref{w4:prop:cong} for $Z=Y$ and the decomposition of Lemma~\ref{lk:lem:split}.
\end{lemma}

\begin{proof}
If $u\in V_Y$ and $u=\lambda+b$ with $\lambda\in M$ and $b\in V_{X\setminus Y}$, then $u=\pi_Y(u)=\pi_Y(\lambda)$. Conversely $\pi_Y(\lambda)=\lambda-\pi_{X\setminus Y}(\lambda)\in(M+V_{X\setminus Y})\cap V_Y$.
\end{proof}

\subsubsection*{Identification of the congruence defects}

Let $w$ be a twisted place, $l=p_w$, $1\le j\le m$, and let $\Lambda$ and $\mathbb T$ be as in Definition~\ref{r5:def:eta}, with $\mathcal O$ enlarged, if necessary, by a finite unramified extension so that every eigensystem of $\mathbb T$ on $\Lambda\otimes\mathbb Q$ has values in $\mathcal O$. This changes no length in Definition~\ref{r5:def:eta}, because for a finite unramified extension $\mathcal O'$ of $\mathcal O$ and a finite length $\mathcal O$-module $N$ one has $\len_{\mathcal O'}(N\otimes_{\mathcal O}\mathcal O')=\len_{\mathcal O}N$, and intersections and sums of submodules commute with the flat base change $\otimes_{\mathcal O}\mathcal O'$.

\begin{proposition}\label{lk:prop:ident}
The algebra $\mathbb T$ is the product $\prod_{\mathfrak m}\mathbb T_{\mathfrak m}$ of its localizations at its maximal ideals, with idempotents $\epsilon_{\mathfrak m}\in\mathbb T$, and $\Lambda=\bigoplus_{\mathfrak m}\Lambda_{\mathfrak m}$ with $\Lambda_{\mathfrak m}=\epsilon_{\mathfrak m}\Lambda$. For each $\mathfrak m$ let $X_{\mathfrak m}$ be the set of eigensystems of $\mathbb T$ on $\epsilon_{\mathfrak m}(\Lambda\otimes\mathbb Q)$, so that $\mathbb T_{\mathfrak m}$ is an $\mathcal O$-subalgebra of $\prod_{x\in X_{\mathfrak m}}\mathcal O$ separating points, and let $G_{\mathfrak m}\subset X_{\mathfrak m}$ be the eigensystems which occur in the lifts of level $j$. Then
\[
\Lambda_{\mathfrak m}=\epsilon_{\mathfrak m}(\Lambda\otimes\mathbb Q)\cap\big(\Lambda+(1-\epsilon_{\mathfrak m})(\Lambda\otimes\mathbb Q)\big),\qquad
\eta_j(w)=\sum_{\mathfrak m}\eta(\Lambda_{\mathfrak m},G_{\mathfrak m}),
\]
and each summand is the normalized length of the module $M^{\rm up}/M^{\rm low}$ of Proposition~\ref{w4:prop:cong} for $(\mathcal O,X_{\mathfrak m},\mathbb T_{\mathfrak m},\Lambda_{\mathfrak m})$ and $Z=G_{\mathfrak m}$.
\end{proposition}

\begin{proof}
The algebra $\mathbb T$ is finite over the complete local ring $\mathcal O$, so it is the product of its localizations at its finitely many maximal ideals, and the corresponding idempotents $\epsilon_{\mathfrak m}$ lie in $\mathbb T$ (idempotents lift along $\mathbb T\to\mathbb T/\varpi\mathbb T$ by completeness). Since $\epsilon_{\mathfrak m}\in\mathbb T$ preserves $\Lambda$, $\Lambda=\bigoplus_{\mathfrak m}\epsilon_{\mathfrak m}\Lambda$. If $u=\epsilon_{\mathfrak m}u=\lambda+(1-\epsilon_{\mathfrak m})u'$ with $\lambda\in\Lambda$, applying $\epsilon_{\mathfrak m}$ gives $u=\epsilon_{\mathfrak m}\lambda\in\Lambda_{\mathfrak m}$, and conversely $\epsilon_{\mathfrak m}\lambda=\lambda-(1-\epsilon_{\mathfrak m})\lambda$. This is the first formula. In Definition~\ref{r5:def:eta} the summand at $\mathfrak m$ is the normalized length of $(V_{\rm good}\cap(\Lambda_{\mathfrak m}+V_{\rm bad}))/(\Lambda_{\mathfrak m}\cap V_{\rm good})$ with $V_{\rm good}=V_{G_{\mathfrak m}}$ and $V_{\rm bad}=V_{X_{\mathfrak m}\setminus G_{\mathfrak m}}$ inside $\Lambda_{\mathfrak m}\otimes\mathbb Q$, which is $\eta(\Lambda_{\mathfrak m},G_{\mathfrak m})$ by definition, with $e_{\mathcal O}=\len_{\mathcal O}(\mathcal O/l)$. Lemma~\ref{lk:lem:proj} identifies the module with $M^{\rm up}/M^{\rm low}$.
\end{proof}

\begin{proposition}[the congruence bound at one maximal ideal]\label{lk:prop:hlb}
In the local Hecke setting let $Z'\subset Z\subset X$, $A\in\mathbb Z\setminus\{0\}$, and let $\alpha\in\operatorname{End}_O(M)$ act on $V_x$ as multiplication by $A$ for $x\in Z'$ and as $0$ for $x\in Z\setminus Z'$. Then
\[
\eta(M,Z')\ \le\ \sum_{x\in Z'}\operatorname{mult}(x)\Big(\frac{\kappa_Z(x)}{e_O}+v_l(A)\Big).
\]
\end{proposition}

\begin{proof}
If $\kappa_Z(x)=\infty$ for some $x\in Z'$ there is nothing to prove. Extending scalars to a finite unramified extension of $O$ changes neither side, as explained before Proposition~\ref{lk:prop:ident}: lengths are preserved, $M^{\rm up}_{Z'}$ and $M^{\rm low}_{Z'}$ commute with the flat base change, $\mathfrak c_Z$ is the kernel of the restriction $\mathbb T\to\prod_{y\notin Z}O$ and so commutes with it too, and the minimum of the valuations of $O'$-linear combinations of the values $t(x)$, $t\in\mathfrak c_Z$, is $\kappa_Z(x)$. So we may assume $\lvert k_O\rvert>\lvert Z'\rvert$.

\emph{A uniform element.} For $x\in Z'$ choose $t_x\in\mathfrak c_Z$ with $v(t_x(x))=\kappa_Z(x)$, and write $t_y(x)=\varpi^{\kappa_Z(x)}u_{yx}$ with $u_{yx}\in O$, so that $u_{xx}\in O^\times$. For $\lambda\in O^{Z'}$ the element $t=\sum_y\lambda_yt_y\in\mathfrak c_Z$ has $t(x)=\varpi^{\kappa_Z(x)}\sum_y\lambda_yu_{yx}$. For each $x$ the reduction of $\lambda\mapsto\sum_y\lambda_yu_{yx}$ is a nonzero linear form on $k_O^{Z'}$, and a vector space over $k_O$ is not the union of $\lvert Z'\rvert<\lvert k_O\rvert$ proper hyperplanes. So there is $\lambda$ with $v(t(x))=\kappa_Z(x)$ for every $x\in Z'$.

\emph{The operator.} Put $\beta=\alpha\circ t$ on $V$. For $\lambda\in M$, $t\lambda\in M$ and $t\lambda=\sum_{x\in Z}t(x)P_x\lambda\in V_Z$ by Lemma~\ref{lk:lem:split}, so $\alpha t\lambda=A\sum_{x\in Z'}t(x)P_x\lambda$ lies in $M\cap V_{Z'}$. The same formula gives $\beta(\pi_{Z'}\lambda)=A\sum_{x\in Z'}t(x)P_x\lambda$. With Lemma~\ref{lk:lem:proj},
\[
\beta\,M^{\rm up}_{Z'}\subset M^{\rm low}_{Z'}\subset M^{\rm up}_{Z'} .
\]
The module $M^{\rm up}_{Z'}=\pi_{Z'}(M)$ is a lattice in $V_{Z'}$, and $\beta$ acts on $V_{Z'}$ diagonally, as $At(x)$ on $V_x$, with determinant $\prod_{x\in Z'}(At(x))^{\operatorname{mult}(x)}\ne0$. Lemma~\ref{w4:lem:det} gives
\[
\len_O\big(M^{\rm up}_{Z'}/M^{\rm low}_{Z'}\big)\le\len_O\big(M^{\rm up}_{Z'}/\beta M^{\rm up}_{Z'}\big)=\sum_{x\in Z'}\operatorname{mult}(x)\big(v(A)+\kappa_Z(x)\big).
\]
Divide by $e_O$ and use $v(A)=e_O\,v_l(A)$ for $A\in\mathbb Z$.
\end{proof}

With $\alpha=1$, $A=1$ and $Z'=Z$ Proposition~\ref{lk:prop:hlb} gives $\eta(M,Z)\le e_O^{-1}\sum_{x\in Z}\operatorname{mult}(x)\kappa_Z(x)$, which refines Proposition~\ref{w4:prop:cong}(ii) because $\kappa_Z(x)\le\delta_x$.

\begin{proposition}[removing eigensystems]\label{lk:prop:cut}
In the local Hecke setting let $G\subset S\subset X$, $A\in\mathbb Z\setminus\{0\}$, and let $\alpha\in\operatorname{End}_O(M)$ act on $V_x$ as $A$ for $x\in G$ and as $0$ for $x\notin S$. Let $\bar M=M/(M\cap V_{X\setminus S})$, identified by $\pi_S$ with the lattice $\pi_S(M)$ of $V_S$, with the action of the restriction of $\mathbb T$ to $S$. Then $M\cap V_{X\setminus S}=\{m\in M:tm=0\text{ for all }t\in\mathfrak c_S\}$, and
\[
\eta(M,G)\le\eta(\bar M,G)+v_l(A)\sum_{x\in G}\operatorname{mult}(x).
\]
\end{proposition}

\begin{proof}
An element $t\in\mathfrak c_S$ kills $V_{X\setminus S}$, and $t_S=\sum_{x\in S}e_x\in\mathfrak c_S$ acts as $C\pi_S$, so $m\in M$ is killed by $\mathfrak c_S$ exactly when $\pi_Sm=0$. Since $\pi_G=\pi_G\pi_S$, $M^{\rm up}_G=\pi_G(M)=\pi_G(\bar M)=\bar M^{\rm up}_G$, and $M\cap V_G\subset\bar M\cap V_G$. If $u\in\bar M\cap V_G$, say $u=\pi_S\lambda$ with $\lambda\in M$, then $\lambda=u+\pi_{X\setminus S}\lambda$ and $\alpha\lambda=Au\in M\cap V_G$. So $A(\bar M\cap V_G)\subset M\cap V_G\subset\bar M\cap V_G$, and since $\bar M\cap V_G$ is a lattice in $V_G$,
\[
\len_O\frac{M^{\rm up}_G}{M\cap V_G}=\len_O\frac{\bar M^{\rm up}_G}{\bar M\cap V_G}+\len_O\frac{\bar M\cap V_G}{M\cap V_G}\le\len_O\frac{\bar M^{\rm up}_G}{\bar M\cap V_G}+v(A)\dim_KV_G .
\]
Divide by $e_O$.
\end{proof}

\begin{lemma}[descent along a map of Hecke algebras]\label{lk:lem:descent}
Let $\mathbb T\subset\prod_{x\in X}O$ and $\mathbb T'\subset\prod_{x'\in X'}O$ be $O$-subalgebras, $p\colon X'\to X$ a map, $Z\subset X$, and $b_1,\dots,b_q\colon X'\to O$. Assume that every $f'\in\mathbb T'$ can be written $f'=\sum_k(a_k\circ p)\,b_k$ with $a_k\in\mathbb T$, and that for $x\notin Z$ the restrictions of $b_1,\dots,b_q$ to $p^{-1}(x)$ are linearly independent over $O$. If $f'\in\mathbb T'$ vanishes on $p^{-1}(X\setminus Z)$, then for every $\pi'\in X'$, $f'(\pi')$ lies in the ideal of $O$ generated by $\{c(p(\pi')):c\in\mathfrak c_Z\}$, and $v(f'(\pi'))\ge\kappa_Z(p(\pi'))$.
\end{lemma}

\begin{proof}
Write $f'=\sum_k(a_k\circ p)b_k$. For $x\notin Z$ and $x'\in p^{-1}(x)$, $\sum_ka_k(x)b_k(x')=f'(x')=0$, so $a_k(x)=0$ for all $k$ by the independence. Hence $a_k\in\mathfrak c_Z$ and $f'(\pi')=\sum_ka_k(p(\pi'))b_k(\pi')$. For the valuation, $v(\sum_ir_iy_i)\ge\min_iv(y_i)$ in a discrete valuation ring.
\end{proof}

\begin{proposition}[from the trimmed data to Hypothesis~\ref{w4:thm:linkn}]\label{lk:prop:gd}
Fix an admissible datum and $1\le j\le m$. For each twisted place $w$ let the local components $i$ at $(j,w)$ be the maximal ideals of Proposition~\ref{lk:prop:ident}, with good sets $G_i$, and suppose that at each $i$ with $G_i\ne\emptyset$ there are $S_i\supset G_i$ and $\alpha_i$ as in Proposition~\ref{lk:prop:cut} with $A=\lvert\Gamma\rvert$ an integer independent of $i$ and $w$. Write $\eta^{\rm cut}_j(w)=\sum_i\eta(\bar\Lambda_i,G_i)$. Assume
\begin{enumerate}[leftmargin=*,label=(\alph*)]
\item $\sum_i\sum_{x\in G_i}\operatorname{mult}(x)\le d^0_j$ for every twisted $w$;
\item $\log\lvert\Gamma\rvert\le C_1x^{C_1}L_\varepsilon$, $x^{C_1}\le b$, $L_\varepsilon\ge1$ and $f,m\ge1$;
\item $\sum_{w\ \mathrm{twisted}}\ell_we_w\,\eta^{\rm cut}_j(w)\le d^0_j\big(C(fmbL_\varepsilon^2+(\log x)^2N_c)+\sum_l\log l\sum_{1\le e<t_r}\len\Sha_l(E,e)\big)$.
\end{enumerate}
Then the inequality of Hypothesis~\ref{w4:thm:linkn} holds with the constant $C+C_1$.
\end{proposition}

\begin{proof}
At a component with $G_i=\emptyset$ both $\eta(\Lambda_i,G_i)$ and $\eta(\bar\Lambda_i,G_i)$ vanish. Propositions~\ref{lk:prop:ident} and~\ref{lk:prop:cut} and (a) give $\eta_j(w)\le\eta^{\rm cut}_j(w)+d^0_j\,v_{p_w}(\lvert\Gamma\rvert)$. With $\ell_w=\log N\mathfrak p_w/[\bK:\mathbb Q]=f_w\log p_w/[\bK:\mathbb Q]$ and $\sum_{w\mid p}e_wf_w=[\bK:\mathbb Q]$,
\[
\begin{aligned}
\sum_{w\ \mathrm{twisted}}\ell_we_w\,v_{p_w}(\lvert\Gamma\rvert)&\le\sum_pv_p(\lvert\Gamma\rvert)\log p=\log\lvert\Gamma\rvert\\ &\le C_1x^{C_1}L_\varepsilon\le C_1bL_\varepsilon^2\le C_1\big(fmbL_\varepsilon^2+(\log x)^2N_c\big),
\end{aligned}
\]
because $N_c\ge0$. Multiply by $d^0_j$ and add (c).
\end{proof}

\subsubsection*{The rank $d^0_j$}

Let $\theta_j\colon\widetilde M_{\mu_0,j}\to M_j$ be the theta map from the forms of level $j$ and standard type to the space $M_j$ of level $j$, let $\theta_j^*$ be its adjoint, and let $d^0_j=\dim Z_j$ with $Z_j=\theta_j^*(M_j)$, the space of lifted forms of level $j$ and standard type. Let $D_0$ be the number of base rows, so that $D_0=\sum_\sigma n^0_\sigma$ over the isotypic types $\sigma$ with multiplicities $n^0_\sigma$.

\begin{proposition}\label{lk:prop:d0}
Suppose that $\dim M_j=D_0$ for $1\le j\le m$. Then $d^0_j\le D_0$. If moreover $\theta_j$ is surjective, then $d^0_j=D_0$. Surjectivity holds when the level projection at the standard type is surjective, $\theta_j$ restricted to its image is the composite of that projection, a bijection $\theta_{j,t}$ of coordinate spaces and a coordinate isomorphism onto $M_j$.
\end{proposition}

\begin{proof}
$d^0_j=\dim\theta_j^*(M_j)\le\dim M_j=D_0$. If $\theta_j$ is surjective, $\theta_j^*$ is injective, because $\langle\theta_j^*z,u\rangle=\langle z,\theta_ju\rangle$ vanishes for all $u$ only if $z=0$, and then $d^0_j=\dim M_j$. The last statement is the composition of surjective maps.
\end{proof}

Theorem~\ref{r5:thm:DIVinN} only needs $d^0_j\le D_0$: the bracket in Hypothesis~\ref{w4:thm:linkn} is nonnegative once $C\ge0$, which may be assumed, so the bound is monotone in $d^0_j$.

\subsubsection*{The local estimates}

Propositions~\ref{lk:prop:ident} to~\ref{lk:prop:gd} reduce Hypothesis~\ref{w4:thm:linkn} to the following statement, with $\kappa=\kappa_{Z_i}$ for the set $Z_i\supset G_i$ of eigensystems which admit a nonzero map of Ribet type, so that after the trimming of Proposition~\ref{lk:prop:cut} one has $Z_i\cap S_i=G_i$ and Proposition~\ref{lk:prop:hlb} applies with $Z'=Z=G_i$ and $\alpha=1$.

\begin{hypothesis}[per-component congruence bound]\label{lk:hyp:block}
There is $C$ such that at every component $i$ of residue characteristic $l$ at a twisted place $w$ and every $\pi\in G_i$,
\[
\frac{\kappa(\pi)}{e_O}\le s_\pi+b_l+\sum_{0<e<t_r}\len\Sha_l(E,e)+v_l(h_E)+\varepsilon_l ,
\]
where $s_\pi$ is the local term of $\rho_\pi$ and $\varepsilon_l=l^{-2}$. Away from the two auxiliary primes one also needs
\[
 s_\pi\le\nu_l+\varepsilon_l,\qquad
 \frac{t_r}{e_O}\sum_{v\in Q_i}\len_O\bigl(O/(\zeta_{i,v}-1)\bigr)\le\varepsilon_l ,
\]
where $\nu_l$ is the cyclotomic exponent of the host at $l$. At $3$ and the small primes of the host take $b_l=C(fmb\log x+x^CL_\varepsilon^2)$, and at $q,d$ take $b_l=2Cfmb$. At the other nonauxiliary primes above $l_*$ take $b_l=0$ outside an exceptional set of at most $(N_T+2)(S_c+3\eta+t_r+8b/l_*+9)$ primes $l\le3b$. For each remaining small or exceptional prime the local argument requires a length estimate $b_l\le Cfn t_r^2n_VL_\varepsilon^2+\mathcal L(l)$, where $\mathcal L(l)$ is the contribution of the split or inert local blocks. The weighted sums of these lengths over the two sets of primes also require bounds. At the two auxiliary primes $l=g_k$ one requires the bound $\eta(\bar\Lambda_i,G_i)\le\sum_{\pi\in G_i}\operatorname{mult}(\pi)r_\pi$ with
\[
\begin{gathered}r_\pi=4fn_Ve_T+2fr^2+s_\pi+\sum_{0<e<t_r}\len\Sha_l(E,e)+v_l(h_E)+4n_V\nu_i\\ {}+t_r\sum_{q\in Q_i}\frac{v(\zeta_{i,q}-1)}{e_O},\qquad e_T=C n_V^2(n_V+\eta+t_r),\end{gathered}
\]
where $\nu_i=\sum_{0<k<t_r}\big(s_{0,i}+v_l(k)\big)$ with $s_{0,i}$ the local level exponent at $l$, $n_V=r+t_r$, and $\zeta_{i,q}$ is the value of the character $\chi^1$ on a tame generator at $q\in Q_i$. Together with these, the summation of the resulting bounds over $\pi$, $i$, $w$ and $l$ gives condition (c) of Proposition~\ref{lk:prop:gd}.
\end{hypothesis}

Lemma~\ref{lk:lem:descent} and Proposition~\ref{w4:prop:fitting} bound $\kappa(\pi)$ by the Fitting lengths of the cocycle module, once (F2) and (F3) hold for the blocks $\rho_\pi,\chi^1,\dots,\chi^{t_r}$. The bounds of the Fitting lengths are the block bounds, the lifting theorem with Wach modules at the places above $l$, the twist and the statements on twisted theta series listed after Proposition~\ref{w4:prop:fitting}. \notgiven{Hypothesis~\ref{lk:hyp:block}, the verification of (F2) and (F3), and the summation are not proved here.}

For the local lifting at $l=2$, first pass to
$K_N=\mathbb Q_l(\mu_{\operatorname{lcm}(l,N)})$ with $16\mid N$.
The cyclotomic group over $K_N$ is $1+2^{v_2(N)}\mathbb Z_2$, hence is
torsion free and procyclic. Herr's cyclotomic complex applies over this
field. For $K\subset K_N$ restriction followed by corestriction is
multiplication by $[K_N:K]$, with loss $v_l([K_N:K])$ in the lifting
bound. The quantitative estimate for the crystalline lattice over $K_N$
is still required.

\subsection{The bound for the congruence defects}\label{tw:sec:divinn}

\begin{theorem}\label{r5:thm:DIVinN}
Let $0<\varepsilon<\frac1{52}$, and suppose that the admissible data at the scale $x\ge x_0$ satisfy condition (D6) of Definition~\ref{r6:R6:def:adm}, that the conclusion of Hypothesis~\ref{w4:thm:linkn} holds at $x$, that $(\log x)^2\le x^\varepsilon$, and that $d^0_j\le D_0$ for $1\le j\le m$. Then for every $1\le j\le m$
\[
\sum_{w\ \mathrm{twisted}}\ell_w\,e_w\,\eta_j(w)\le(C_N+1)\,D_0\,x^{\varepsilon}\,(B_0+N_c).
\]
\end{theorem}

\begin{proof}
Put $y=x^\varepsilon=L_\varepsilon\ge1$. Since $B_0=fmby$, the first term in the bracket of Hypothesis~\ref{w4:thm:linkn} is $fmby^2=yB_0$, and $(\log x)^2N_c\le yN_c$ by hypothesis. By (D6), $V\le\sqrt b\le b\le B_0\le y(B_0+N_c)$, where $b\le B_0$ because $fmy\ge1$. Hence the bracket is at most $C_Ny(B_0+N_c)+y(B_0+N_c)$.
\end{proof}

Theorem~\ref{r5:thm:DIVinN} enters the bound for the heights $h_u$ of the period maps $B_u$ in \eqref{A4} through the local terms $E_b$ of Hypothesis~\ref{r4:R4:HT}, and the bound for the spread of the slopes of the family $(N_u)$ in \eqref{A3}. The local terms $E_u$ of Definition~\ref{r5:def:Eterms} do not depend on $\eta_j(w)$. Proposition~\ref{lk:prop:d0} gives the required inequality once $\dim M_j=D_0$.


\section{The proof of Theorem A}\label{sec:proofA}

This section assembles the proof of Theorem~A. It fixes the form of the admissible data that the arithmetic part of the argument uses, it constructs the weights $a_t$ and the nodes $w_{u_0,j}$, it states the inequalities that the chain of theta lifts satisfies, and it deduces from them the seven hypotheses \eqref{A1} to \eqref{A7} of Proposition~\ref{acct} together with the bound for the errors. The inequalities between slopes, heights and local terms are proved in the sections on the gain and the loss, on the slope inequality, on the archimedean side and on the twisted places, and Hypothesis~\ref{r6:R6:thm:records} below collects them in one statement.

\subsection{Admissible data on Selmer windows}\label{r6:R6:sec:adm}

The bound for the local terms at the twisted places uses a bound for Tate--Shafarevich groups of the host field. That bound holds on windows of scales whose position depends on the auxiliary primes $q$ and $d$. We therefore attach to admissible data a base scale $Y_0$, at which $q$ and $d$ are chosen, and a window $[Y,Y^{3/2}]$ above it, on which the Selmer bound holds.

For a number field $F$, a prime $l$ and an integer $i$, let $S_l$ be the set of places of $F$ above $l$ and the archimedean places, let $G_{F,S_l}$ be the Galois group of the maximal extension of $F$ unramified outside $S_l$, and let $\mathbb Q_l/\mathbb Z_l(-i)$ be the group $\mathbb Q_l/\mathbb Z_l$ on which $G_{F,S_l}$ acts through the $(-i)$-th power of the $l$-adic cyclotomic character. Put
\[
\Sha_l(F,i)=\ker\Big(H^1\big(G_{F,S_l},\mathbb Q_l/\mathbb Z_l(-i)\big)\to\prod_{v\mid l}H^1\big(F_v,\mathbb Q_l/\mathbb Z_l(-i)\big)\Big).
\]
For $F$ and $t_0\ge1$ put
\[
V(F,t_0)=\sum_{l}\log l\sum_{i=1}^{t_0-1}\operatorname{length}_{\mathbb Z_l}\Sha_l(F,i),
\]
the sum over all primes $l$, whenever every group $\Sha_l(F,i)$ with $1\le i<t_0$ has finite length.

Fix a primitive Dirichlet character $\psi_0$ of conductor $M_\psi$ and order $o$, put $N_0=8M_\psi$, and fix $0<\varepsilon<\frac1{52}$.

\begin{definition}\label{r6:R6:def:adm}
Let $J\ge2$, and let $x$, $Y$ and $Y_0$ be real numbers with $3\le Y_0\le Y\le x\le Y^{3/2}$ and $Y^2\le Y_0^J$. \emph{Admissible data at the scale $x$ for the window $[Y,Y^{3/2}]$ with base $Y_0$} consist of the following objects.
\begin{enumerate}
\item[(D1$'$)] Primes $q$ and $d$ that satisfy the conclusions of Theorems~\ref{stage1} and~\ref{stage2} with $N_0$, $K=J$, $\varepsilon$, $\tau=\varepsilon$, with $Y_0$ in place of $Y$, and with $z'=z\log^2Y_0$, where $z=Y_0^{\varepsilon/5}$. An integer $e\ge0$ with $x^{2+8\varepsilon}\le c\,2^e\le2x^{2+8\varepsilon}$, where $c$ is the integer of Proposition~\ref{field}(iii). The field $E$, the groups $H$ and $A_0$ and the integer $n=\lvert H\rvert$ of Proposition~\ref{field} for $\psi_0$, $e$, $q$ and $d$, with the identification of $\mathrm{Gal}(E/\mathbb Q)$ with $\mathcal G$ of Section~\ref{sec:field}.
\item[(D2)--(D5)] The objects (D2) to (D5) of Definition~\ref{def:admissible}.
\item[(D6)] With $t_0=\lfloor x^{400}\rfloor$, every group $\Sha_l(E,i)$ with $l$ prime and $1\le i<t_0$ has finite length, and $V(E,t_0)\le\exp(\frac12x^{\varepsilon})$.
\end{enumerate}
\end{definition}

Condition (D6) asks that $V(E,t_0)\le\sqrt b$ with $b=e^{L_\varepsilon}$ as in \eqref{scale}. When $Y_0=Y$, condition (D1$'$) is condition (D1) of Definition~\ref{def:admissible} with $K=J$ in place of $K=\frac32$. The conclusions (iii) of Theorems~\ref{stage1} and~\ref{stage2} at the scale $Y_0$ hold on the window $[Y,Y^{3/2}]$, because this interval lies in $[Y_0,Y_0^J]$.

The integer $t_r$ of Proposition~\ref{ranks} is the number of copies of the auxiliary space added at each step of the chain, and the twists $i<t_r$ are the twists at which the twisted places are examined. The next lemma shows that (D6) covers all of them.

\begin{lemma}\label{r6:R6:lem:tr}
Let $0<\varepsilon<\frac1{52}$. There is $x_0$ such that for every $x\ge x_0$ and all integers $n,m,k,t_r$ with $x^{2+8\varepsilon}\le n\le2x^{2+8\varepsilon}$ that satisfy the conclusions of Proposition~\ref{ranks}(i), we have $t_r\le x^{400}-1$.
\end{lemma}

\begin{proof}
Put $r=k+m-1$ and $f=2n$. Proposition~\ref{ranks}(i) gives $m\ge\sqrt{x\delta_1}=x^{1/2-\varepsilon/16}$, $r\le2k\le4x^{9/2+10\varepsilon}$, $k-1\le2x^{9/2+10\varepsilon}$ and $f\le4x^{2+8\varepsilon}$. Then
\[
\frac{f^2(k-1)^3L_\varepsilon^2}m\le16\cdot8\,x^{4+16\varepsilon+27/2+30\varepsilon+2\varepsilon-1/2+\varepsilon/16}=128\,x^{17+\frac{769}{16}\varepsilon},
\]
and $f\le4x^{17}$, so the definition of $n_W$ gives
\[
n_W\le8+10\,r\Big(f+\frac{f^2(k-1)^3L_\varepsilon^2}m\Big)\le8+5280\,x^{\frac{43}2+\frac{929}{16}\varepsilon}\le8+5280\,x^{22.625},
\]
because $\frac{929}{16}\varepsilon<\frac{929}{832}<\frac98$. With $T_O\le2x^{377}$ and $8fr\le128x^7$, the relation $r+t_r=8fr+4+T_On_W$ gives $t_r\le128x^7+4+2x^{377}(8+5280x^{22.625})\le x^{400}-1$ for $x\ge x_0$, because $377+22.625<400$.
\end{proof}

\begin{proposition}[the Selmer window]\label{r6:R6:thm:selmer}
Assume Hypothesis~\ref{w4:thm:vwin}. Let $\psi_0$ be a primitive Dirichlet character, $0<\varepsilon<\frac1{52}$ and $\alpha=2+8\varepsilon$. Let $x_e\ge1$, let $c_\ast\ge1$ be an integer, and let $e(x)$ be a function on $[x_e,\infty)$ with values in the nonnegative integers such that, for every $x\ge x_e$, the integer $c$ of Proposition~\ref{field}(iii) for $e=e(x)$ equals $c_\ast$ and $x^\alpha\le c_\ast2^{e(x)}\le2x^\alpha$. There are $J_1\ge1$ and $X_0\ge x_e^{J_1}$ with the following property. Let $X\ge X_0$, put $Y_0=X^{1/J_1}$ and $z=Y_0^{\varepsilon/5}$, and let $q$ and $d$ be primes with $q\equiv3\pmod8$, $e^{z/2}<q\le e^z$, $d\equiv7\pmod{16}$, $e^{z\log^2Y_0/2}<d\le e^{z\log^2Y_0}$, such that $q$ and every odd prime $p\le z$ divide $d+1$. Then there is $Y$ with $Y_0\le Y$ and $Y^2\le Y_0^{J_1+1}$ such that for every $x$ with $Y\le x\le Y^{3/2}$ the field $E$ of Proposition~\ref{field} for $\psi_0$, $e(x)$, $q$ and $d$ satisfies (D6) at the scale $x$.
\end{proposition}

\begin{proof}
Put $\theta=c_\ast^{1/\alpha}\ge1$, so that $\theta^\alpha=c_\ast$. Define $e'$ on $[1,\infty)$ by $e'(y)=e(\theta y)$ when $\theta y\ge x_e$ and $e'(y)=\lceil\alpha\log_2y\rceil$ otherwise. In the second case $y^\alpha\le2^{e'(y)}<2y^\alpha$. In the first case
\[
y^\alpha=\frac{(\theta y)^\alpha}{c_\ast}\le2^{e(\theta y)}\le\frac{2(\theta y)^\alpha}{c_\ast}=2y^\alpha .
\]
So $e'$ satisfies the normalization of Hypothesis~\ref{w4:thm:vwin}, which gives numbers $J\ge1$ and $X'_0$ for $e'$. Put $J_1=J$ and let $X_0\ge\max(X'_0,x_e^{J})$ be so large that $X^{1/J}\ge\theta^{400}$ for $X\ge X_0$.

Let $X\ge X_0$ and let $q,d$ be as in the statement. These are exactly the conditions of Definition~\ref{w4:def:winpair} for a window pair at scale $Y_0=X^{1/J}$. Hypothesis~\ref{w4:thm:vwin} gives $W_0$ with $Y_0\le W_0\le X^{1/2}$ such that for every $y\in[W_0,W_0^2]$, with $t=\lfloor y^{401}\rfloor$ and $E'=E(\psi_0,e'(y),q,d)$, every $\Sha_l(E',i)$ with $0<i<t$ is finite and $\sum_l\sum_{1\le i<t}\lambda(\Sha_l(E',i))\le\exp(\frac12y^\varepsilon)$.

Put $Y=\theta W_0$. Then $Y\ge W_0\ge Y_0$, and $Y^2=\theta^2W_0^2\le\theta^2X\le X^{1+1/J}=Y_0^{J+1}$ because $X^{1/J}\ge\theta^2$. Let $Y\le x\le Y^{3/2}$ and put $y=x/\theta$. Then $y\ge W_0$, and $y\le\theta^{1/2}W_0^{3/2}\le W_0^2$ because $W_0\ge Y_0\ge\theta$. Since $\theta y=x\ge Y_0\ge x_e$, we have $e'(y)=e(x)$.

The field $E(\psi_0,e(x),q,d)$ of Definition~\ref{w4:def:chifield} is the field $E$ of Proposition~\ref{field} for $\psi_0$, $e(x)$, $q$ and $d$. Indeed $q$ and $d$ are odd primes larger than the conductor $M_\psi$ of $\psi_0$, so $\operatorname{lcm}(2^{e+2},M_\psi,q,d)=\operatorname{lcm}(2^{e+2},M_\psi)\,qd$, and a class $a$ satisfies $(\frac aq)=(\frac ad)=1$ exactly when it is a square modulo $q$ and modulo $d$, so the two fixing groups coincide. By Lemma~\ref{w4:lem:indepS} with $S$ the set of places above $l$, the group $\Sha_l(E,i)$ of Definition~\ref{w4:def:sha} is the group $\Sha_l(E,i)$ of Subsection~\ref{r6:R6:sec:adm}, for every $i\ge1$, and $\lambda(\Sha_l(E,i))=\log l\cdot\operatorname{length}_{\mathbb Z_l}\Sha_l(E,i)$. Since $x^{400}=\theta^{400}y^{400}\le y^{401}$, the integer $t_0=\lfloor x^{400}\rfloor$ is at most $\lfloor y^{401}\rfloor$. Hence every $\Sha_l(E,i)$ with $1\le i<t_0$ is finite, and
\[
V(E,t_0)\le\sum_l\sum_{1\le i<\lfloor y^{401}\rfloor}\lambda\big(\Sha_l(E,i)\big)\le\exp\big(\tfrac12y^\varepsilon\big)\le\exp\big(\tfrac12x^\varepsilon\big),
\]
because the terms are nonnegative and $y\le x$. This is (D6) at the scale $x$.
\end{proof}

\begin{proposition}\label{r6:R6:prop:dataexist}
Assume Hypothesis~\ref{w4:thm:vwin}. For every primitive Dirichlet character $\psi_0$ and every $0<\varepsilon<\frac1{52}$ there is $J\ge2$ with the following property. For every $y$ there are real numbers $Y_0\le Y$ with $Y\ge y$ such that, for every $x$ with $Y\le x\le Y^{3/2}$, admissible data at the scale $x$ for the window $[Y,Y^{3/2}]$ with base $Y_0$ exist.
\end{proposition}

\begin{proof}
The integer $c$ of Proposition~\ref{field}(iii) divides $o$ and does not depend on $e$ once $2^{e+2}$ is divisible by the power of $2$ in $M_\psi$, as in the proof of Proposition~\ref{dataexist}. The interval from $x^{2+8\varepsilon}$ to $2x^{2+8\varepsilon}$ has ratio $2$, so for $x$ larger than a number $x_e$ it contains a number $c\,2^e$ with this constant $c$, and we fix a function $e(x)$ on $[x_e,\infty)$ as in Proposition~\ref{r6:R6:thm:selmer}. Let $J_1$ and $X_0$ be the numbers of Proposition~\ref{r6:R6:thm:selmer} for this function, and put $J=J_1+1$. Let $y$ be given, and let $X\ge X_0$ be so large that $Y_0=X^{1/J_1}$ satisfies $Y_0\ge\max(y,3)$, that $Y_0$ is larger than the numbers $Y_0$ of Theorems~\ref{stage1} and~\ref{stage2} for $N_0$, $K=J$, $\varepsilon$ and $\tau=\varepsilon$, and that every $x\ge Y_0$ is larger than the number $Y_0$ of Proposition~\ref{dataexist}. Put $z=Y_0^{\varepsilon/5}$. Theorem~\ref{stage1} at the scale $Y_0$ gives a prime $q\equiv3\pmod8$ with $e^{z/2}<q\le e^z$ that satisfies its conclusions, because its exceptional set has at most $Y_0^{-\varepsilon}\lvert P\rvert<\lvert P\rvert$ elements. Theorem~\ref{stage2} at the scale $Y_0$, with this $q$ and with $z'=z\log^2Y_0$, gives in the same way a prime $d$ with $e^{z'/2}<d\le e^{z'}$, $d\equiv7\pmod{16}$ and $d\equiv-1$ modulo $q$ and modulo every odd prime $p\le z$, that satisfies its conclusions. These primes satisfy the hypotheses of Proposition~\ref{r6:R6:thm:selmer}, which gives $Y$ with $Y_0\le Y$ and $Y^2\le Y_0^{J_1+1}=Y_0^J$. Thus $Y\ge y$.

Let $Y\le x\le Y^{3/2}$. The primes $q,d$ and $e=e(x)$ give (D1$'$) through Proposition~\ref{field}. Since $Y\le x\le Y^{3/2}\le Y^2$, Proposition~\ref{r6:R6:thm:selmer} gives (D6). The objects (D2) to (D5) depend on $x$, $n$, $E$ and the conductor of $E$, and the proof of Proposition~\ref{dataexist} constructs them for $x\ge Y_0$ from Propositions~\ref{ranks}, \ref{tops} and~\ref{behrend}, statement (C9), Lemma~\ref{disjoint}, Linnik's theorem and Lemma~\ref{three}.
\end{proof}

The proof of Theorem~\ref{mainthm} uses the admissible data through the properties of $q$ and $d$ and through (D2) to (D5). With (D1$'$) in place of (D1), the same proof applies after three changes. First, the bound $\lvert\ell_\psi\rvert\le401\,c_0K\log Y_0$ of Theorems~\ref{stage1}(ii) and~\ref{stage2}(ii) holds with $K=J$, and $Y_0\le x$ gives $\lvert\ell_\psi\rvert\le401\,c_0J\log x$. Second, $\log q\le Y_0^{\varepsilon/5}$ and $\log d\le Y_0^{\varepsilon/5}\log^2Y_0$, so $\log N_\psi\le\lambda=A\,x^{\varepsilon/5}\log^2x$ with a constant $A$ that depends on $\psi_0$ and $J$, and the argument with Proposition~\ref{thinprop} and Corollary~\ref{realbox} is unchanged. Third, Theorems~\ref{stage1}(iii) and~\ref{stage2}(iii) give $\lvert G^\ast_\chi(x)\rvert\le x^{1/2+\varepsilon}$ for $Y_0\le x\le Y_0^J$, an interval that contains the window $[Y,Y^{3/2}]$. The conclusion \eqref{finalbound} then holds on windows $[Y,Y^{3/2}]$ for arbitrarily large $Y$, which is the hypothesis of Corollary~\ref{corA}.

\subsection{The weights and the nodes}\label{r6:R6:sec:nodes}

Fix admissible data at the scale $x$ as in Definition~\ref{r6:R6:def:adm}, and let $a$, $b$, $b'$ and $L_\varepsilon$ be as in \eqref{scale}. Let $W$ be the number of roots of unity in $E$, and put $h_0=2\lceil x^{402}\rceil$.

\begin{definition}\label{r6:R6:def:weights}
\emph{Weights} for the data are integers $a_t$ and $b_t$, for $t\in T$, with
\[
a\le a_t\le a+\frac a{kx^{13}},\qquad b'\le b_t\le\Big(1+\frac1{x^2}\Big)b' .
\]
\end{definition}

Weights exist for large $x$, because $a/(kx^{13})\ge1$ and $b'/x^2\ge1$. The integers $a_t$ are the weights of Theorem~A, and the integers $b_t$ fix the nodes at the elements of $\Phi$ attached to $T$.

Under the identification of Section~\ref{sec:field}, an element $(e_1,e_2,h)\in\mathcal G$ with $e_1=0$ lies in $\Phi$. For $u_0\in A_0$ we divide $\Phi$ into three sets. The \emph{top} elements are the elements $(0,0,t)$ with $t\in T$. The \emph{thick} elements are the elements $(0,1,h)$ with $h\in T_2+u_0$, which are the elements thick with respect to $u_0$ for the set $T_2$ in the sense of Section~\ref{sec:step5}. The remaining elements of $\Phi$ are \emph{thin}. The top elements and the thick elements differ in the coordinate $e_2$, so the three sets are disjoint.

Fix an embedding $\iota\colon E\to\mathbb C$. The nodes are the archimedean parameters of algebraic Hecke characters of $E$, and the infinity type of such a character is trivial on the roots of unity of $E$. The next lemma makes this condition compatible with a shift of all parameters by a common integer.

\begin{lemma}\label{r6:R6:lem:cmroot}
Let the data be as above, and let $\zeta\in E$ be a root of unity. Then $\prod_{\rho\in\Phi}\iota(\rho(\zeta))=1$.
\end{lemma}

\begin{proof}
The field $E$ lies in $\mathbb Q(\zeta_N)$ with $N=N_1qd$, and $d$ divides $N$ exactly. The inertia group of $\mathbb Q(\zeta_N)/\mathbb Q$ at $d$ is the group of the classes $a\in(\mathbb Z/N)^\times$ with $a\equiv1\pmod{N_1q}$. Its image in $\mathrm{Gal}(E/\mathbb Q)=(\mathbb Z/N)^\times/U$ is the quotient of $(\mathbb Z/d)^\times$ by its squares, whose nontrivial element acts by $-1$ on $\sqrt{-d}$ and trivially on $\sqrt{-q}$ and on $H$. In the coordinates of Section~\ref{sec:field} the inertia group of $E$ at a prime $\mathfrak P$ above $d$ is therefore $\{0,\iota_d\}$ with $\iota_d=(1,0,0)$, and the ramification index of $d$ in $E$ is $2$.

Let $r$ be the order of $\zeta$. If $d$ divided $r$, then $E$ would contain $\mathbb Q(\zeta_d)$, in which $d$ has ramification index $d-1$, and $d-1$ would divide $2$. This is impossible because $d\ge7$. So $d\nmid r$. The element $\iota_d$ lies in the inertia group at $\mathfrak P$, so $\iota_d(\zeta)\equiv\zeta\pmod{\mathfrak P}$. Both $\zeta$ and $\iota_d(\zeta)$ are $r$-th roots of unity. If they were different, $1-\eta\in\mathfrak P$ for an $r$-th root of unity $\eta\ne1$, and $1-\eta$ divides $r=\prod_{\eta'\ne1}(1-\eta')$, where $\eta'$ runs over the $r$-th roots of unity different from $1$. Then $r\in\mathfrak P$, which contradicts $d\nmid r$. Hence $\iota_d(\zeta)=\zeta$.

Complex conjugation $\gamma_c$ corresponds to the class $a=-1$, so $\gamma_c=(1,1,h_c)$, where $h_c$ is the class of $-1$ in $H$. Since $E$ is a CM field, $\iota(\gamma_c(\zeta))=\overline{\iota(\zeta)}=\iota(\zeta)^{-1}$. Put $s_0=\gamma_c+\iota_d=(0,1,h_c)$. Then $s_0\in\Phi$, $s_0\ne0$ and $2s_0=0$, and $\iota(s_0(\zeta))=\iota(\gamma_c(\iota_d(\zeta)))=\iota(\zeta)^{-1}$. The set $\Phi$ is a subgroup of $\mathcal G$, and it is the disjoint union of the cosets $\{\rho,\rho+s_0\}$ for $\rho$ in a set $\Phi_0$ of representatives. For each $\rho$,
\[
\iota(\rho(\zeta))\,\iota\big((\rho+s_0)(\zeta)\big)=\iota\big(\rho(\zeta\,s_0(\zeta))\big)=\iota(\rho(1))=1,
\]
because $s_0(\zeta)=\zeta^{-1}$. The product over $\Phi$ is the product of these pairs.
\end{proof}

\begin{proposition}\label{r6:R6:prop:nodes}
Let the data be as above, let $u_0\in A_0$, and let $b_t$, for $t\in T$, be integers. Then there are integers $b_1(\rho)$, for $\rho\in\Phi$, with the following properties.
\begin{enumerate}
\item[(i)] $b_1(0,0,t)=b_t+\lfloor b'/x\rfloor$ for every top element $(0,0,t)$.
\item[(ii)] $\lceil b\rceil\le b_1(\rho)<\lceil b\rceil+W$ for every thick element $\rho$.
\item[(iii)] $b_1(\rho)=\lceil b'/(xn)\rceil$ for every thin element $\rho$.
\item[(iv)] For $1\le j\le m$ put $w_{u_0,j}(\rho)=b_1(\rho)+(j-\frac12)h_0$. Then each $w_{u_0,j}(\rho)$ is an integer, and $\prod_{\rho\in\Phi}\iota(\rho(\zeta))^{w_{u_0,j}(\rho)}=1$ for every root of unity $\zeta\in E$ and every $j$.
\end{enumerate}
\end{proposition}

\begin{proof}
The group of roots of unity of $E$ is cyclic of order $W$. Let $\zeta_W$ be a generator. For $\rho\in\mathcal G$ there is $\kappa(\rho)\in(\mathbb Z/W)^\times$ with $\rho(\zeta_W)=\zeta_W^{\kappa(\rho)}$, because $\rho(\zeta_W)$ is again a generator. For integers $c(\rho)$, $\rho\in\Phi$, the product $\prod_\rho\iota(\rho(\zeta))^{c(\rho)}$ equals $1$ for every root of unity $\zeta$ of $E$ exactly when it equals $1$ for $\zeta=\zeta_W$, that is, when $\sum_\rho\kappa(\rho)c(\rho)\equiv0\pmod W$. Lemma~\ref{r6:R6:lem:cmroot} with $\zeta=\zeta_W$ gives $\sum_{\rho\in\Phi}\kappa(\rho)\equiv0\pmod W$.

The set $T_2$ is nonempty by (D4), so there is a thick element $\rho_\ast$. Define $b_1^0(\rho)$ by (i) and (iii) on the top and thin elements, and put $b_1^0(\rho)=\lceil b\rceil$ on the thick elements. The integer $h_0$ is even, so $(j-\frac12)h_0$ is an integer for every $j$. Put $R_0=\sum_\rho\kappa(\rho)\big(b_1^0(\rho)+\frac12h_0\big)$, a class modulo $W$. Since $\kappa(\rho_\ast)$ is a unit modulo $W$, there is an integer $0\le s<W$ with $\kappa(\rho_\ast)s\equiv-R_0\pmod W$. Put $b_1(\rho_\ast)=\lceil b\rceil+s$ and $b_1(\rho)=b_1^0(\rho)$ for $\rho\ne\rho_\ast$. Then (i) to (iii) hold, and for $1\le j\le m$
\[
\sum_\rho\kappa(\rho)\,w_{u_0,j}(\rho)=R_0+\kappa(\rho_\ast)s+(j-1)h_0\sum_\rho\kappa(\rho)\equiv0\pmod W,
\]
which is (iv).
\end{proof}

\begin{corollary}\label{r6:R6:cor:nodes}
Let $c_h=402$. There is $x_0$ such that for every $x\ge x_0$, every choice of admissible data at the scale $x$, every $u_0\in A_0$ and all weights of Definition~\ref{r6:R6:def:weights}, the nodes $w_{u_0,j}$ of Proposition~\ref{r6:R6:prop:nodes} satisfy the bounds of Proposition~\ref{nodes}(i) at the thick elements and the bounds of Proposition~\ref{nodes}(ii) at the other elements of $\Phi$. In particular $1\le w_{u_0,j}(\rho)\le2b'$ when $\rho$ is not thick, and $b\le w_{u_0,j}(\rho)\le b+W+mh_0$ when $\rho$ is thick.
\end{corollary}

\begin{proof}
The nodes have the form $b_1(\rho)+(j'+\frac12)h_0$ with $j'=j-1$ and $0\le j'<m$, and $h_0=2\lceil x^{c_h}\rceil$. By (D1$'$), $n=c\,2^e\le2x^{2+8\varepsilon}\le2x^{9/4}$, Proposition~\ref{ranks} gives $m\le2\sqrt{x\delta_1}\le2\sqrt x$, and $W\le32n^2$ by the remark after Proposition~\ref{nodes}. At a thick element, Proposition~\ref{r6:R6:prop:nodes}(ii) is the hypothesis of Proposition~\ref{nodes}(i). At a thin element, (iii) is the first hypothesis of Proposition~\ref{nodes}(ii). At a top element, (i) gives $b_1-\lfloor b'/x\rfloor=b_t$, and $b'\le b_t\le(1+x^{-2})b'\le\frac54b'$ is the second hypothesis of Proposition~\ref{nodes}(ii). The node $w_{u_0,j}(\rho)\ge\lceil b'/(xn)\rceil\ge1$ in all three cases.
\end{proof}

The archimedean parameters of the $m$ lifts at an element $\rho\in\Phi$ are the integers $b_1(\rho)+j'h_0$ and $b_1(\rho)+(j'+\frac12)h_0$ for $0\le j'<m$, and the nodes are the second family. Consecutive parameters differ by $\frac12h_0$, and the choice $c_h=402$ together with Lemma~\ref{r6:R6:lem:tr} gives $\frac12h_0\ge x^{402}\ge\frac12(t_r+3)$. This is the gap between the archimedean parameters that Hypothesis~\ref{r6:R6:thm:records} uses.

\subsection{Chain records}\label{r6:R6:sec:records}

Fix admissible data at the scale $x$, weights, an element $u_0\in A_0$ and the nodes $w_{u_0,j}$ of Proposition~\ref{r6:R6:prop:nodes}. We use the following quantities. Let $N_T=\lvert T\rvert$, $N_2=\lvert T_2\rvert$, $A=\sum_{t\in T}a_t$ and $f=2n$. Put
\begin{gather*}
B_0=2\,n\,m\,b\,x^{\varepsilon},\qquad N_c=\frac{N_2\,k\,b\,x^{2\varepsilon}}{n^2},\qquad X''=x\,n\,X',\\
\operatorname{rd}_E=\lvert\operatorname{disc}E\rvert^{1/[E:\mathbb Q]},\qquad\Lambda_a=\log\big(a\operatorname{rd}_E\big),\qquad v=\tfrac12\operatorname{rd}_E,\qquad g=\log(g_0g_1),
\end{gather*}
with $X'$ as in \eqref{scale} and $g_0,g_1$ the primes of (D5). Then the number $\mathcal B$ of \eqref{calB} is
\begin{equation}\label{r6:R6:calB}
\mathcal B=N_T\,a\Big(1+\frac{m^2}k\Big)+m\,(B_0+N_c).
\end{equation}
Let $\lambda_1=\ell_{\chi_1}=L'(1,\chi_{-d})/L(1,\chi_{-d})$, which is real because $\chi_{-d}=(\frac{-d}{\cdot})$ is real, and put $\log\Omega=\gamma+\log(\pi/v)-\lambda_1$. Let $S$ be the sum over $t\in T$ and $1\le j\le m$ of the left side of \eqref{excessid} with $a=a_t$ and $w=w_{u_0,j}$.

For a hermitian lattice $N$ over $\mathcal O_E$ let $\rk N$, $\mu(N)$ and $\mu_{\max}(N)$ denote its rank, its slope and its maximal slope, and let $\mu_{\min}(N)$ be its minimal slope. For real numbers $u\ge0$ and $H>0$ we read $u\log(1+H/u)$ as $0$ when $u=0$.

\begin{definition}\label{r6:R6:def:record}
Let $C_\ast\ge1$, $C_H\ge0$, $e_D\ge0$ and $e_L\ge0$ be real numbers. A \emph{chain record} for the data, the weights and $u_0$, with constants $C_\ast$, $C_H$ and exponents $e_D$, $e_L$, consists of
\begin{enumerate}
\item[(a)] a finite nonempty set $\mathcal W$ with $\lvert\mathcal W\rvert\le2^{k-3}$,
\item[(b)] for every $u\in\mathcal W$ a prime $T_u$, hermitian $\mathcal O_E$-lattices $N_u$ with $\rk N_u\ge1$ and $M_u$, and a real number $h_u$,
\item[(c)] a hermitian $\mathcal O_E$-lattice $R$ with dual lattice $R^\vee$,
\item[(d)] real numbers $E_{u,j}$ for $u\in\mathcal W$ and $1\le j\le m$, and real numbers $\mathcal L$, $\Gamma$, $\mathcal I$, $\delta$, $\eta$ and $\epsilon$,
\end{enumerate}
such that the conditions (K1) to (K8) below hold. In them $u$ and $u'$ run over $\mathcal W$ and $j$ over $1\le j\le m$, and
\[
E_u=\sum_{j=1}^mE_{u,j},\qquad s=\frac{\sum_ud_u\big(\mu(N_u)-h_u\big)}{\sum_ud_u},\qquad u_F=\max_u\big(\mu_{\max}(N_u)-h_u\big)-s,\qquad d_u=\rk N_u .
\]
\begin{enumerate}
\item[(K1)] $S-\Gamma\le C_\ast N_T\,m\,b\,x^{e_L\varepsilon}$ and $0\le\mathcal I\le C_\ast N_T\,m\,b\,x^{e_L\varepsilon}$.
\item[(K2)] $\mu_{\max}(N_u)\le\mu_{\max}(M_u)+\Gamma+\mathcal L-E_u+C_\ast m B_0$ and $\mu(M_u)+\Gamma+\mathcal L-E_u-\mathrm{Lo}_u\le\mu(N_u)$, where
\[
\mathrm{Lo}_u=C_\ast\mathcal I+C_\ast N_T\,m\Big(\frac ak+b'\Big)\log T_u+C_\ast X''+1 .
\]
\item[(K3)] The numbers $\delta$, $\eta$ and $\epsilon$ satisfy
\[
0\le\delta\le\tfrac12,\qquad\eta\ge0,\qquad0\le\epsilon\le\tfrac14,\qquad(\eta+4\sqrt\epsilon)\,x^{C_H}\le6,
\]
and for every $H>0$ with $H\ge s-\mu_{\min}(R^\vee)$
\[
s-\log\lvert\mathcal W\rvert-\frac1{1-\delta}\Big(4u_F\log\Big(1+\frac H{u_F}\Big)+(\eta+4\sqrt\epsilon)H\Big)\le\mu(R^\vee).
\]
\item[(K4)] $\mu_{\min}(R^\vee)=-\mu_{\max}(R)$, $s+\mu_{\max}(R)\le x^{C_H}b$, and
\begin{multline*}
\mu(R^\vee)+h_u-\mathcal L\le\sum_{t\in T}\big(m\ph(a_t)-ma_t\log\Omega-a_t\log T_u\big)\\
+C_\ast N_T\,a\,\Lambda_a+C_\ast m\,(B_0+N_c)\,x^{e_D\varepsilon}.
\end{multline*}
\item[(K5)] $\lvert E_{u,j}-A\,g\rvert\le C_\ast B_0$.
\item[(K6)] $\big\lvert(h_u+A\log T_u)-(h_{u'}+A\log T_{u'})\big\rvert\le C_\ast\big(N_T\,a\,\Lambda_a+mB_0\big)$.
\item[(K7)] The auxiliary primes satisfy
\begin{gather*}
\log T_u\ge\sum_{j=1}^m\frac1f\sum_{\rho\in\Phi}\log\frac{w_{u_0,j}(\rho)\,v}{\pi}+mg-C_\ast\Big(\Lambda_a+\frac{m^2\Lambda_a}k\Big),\\
\log T_u\le C_\ast m L_\varepsilon,\qquad\lvert\log T_u-\log T_{u'}\rvert\le C_\ast\log k .
\end{gather*}
\item[(K8)] $-C_\ast N_T\,a\,\Lambda_a\big(1+\frac mk\big)\le\mu(M_u)$, $\mu_{\max}(M_u)\le C_\ast N_T\,a\,\Lambda_a$, and $\mu(N_u)\le\mu_{\max}(N_u)$.
\end{enumerate}
\end{definition}

In the application, $\mathcal W$ is a set of words that index the auxiliary primes $T_u$, $N_u$ is the rescaled lattice of images of the Hecke operators at $T_u$, $M_u$ is the corresponding base lattice, $R$ is the lattice of automorphic forms at the top of the chain, and $h_u$ bounds the height of the period map $B_u\colon N_u\to R^\vee$. The number $\Gamma$ is the total gain of the $m$ lifts at the inert primes, $\mathcal L$ is the sum of the local terms of the lifts, $E_{u,j}$ is the local term of the $j$-th lift at the twisted places, and $\mathcal I$ is the index loss at the inert primes.

\begin{hypothesis}[chain records]\label{r6:R6:thm:records}
There are absolute constants $e_D,e_L\ge0$ with the following property. For every primitive Dirichlet character $\psi_0$ there is $0<\varepsilon_1\le\frac1{52}$, and for every $0<\varepsilon<\varepsilon_1$ and every $J\ge2$ there are $C_\ast\ge1$, $C_H\ge0$ and $x_1$ such that the following holds. Let $x\ge x_1$, and let $Y_0\le Y$ be such that admissible data at the scale $x$ for the window $[Y,Y^{3/2}]$ with base $Y_0$ exist, for this $J$. Then there are admissible data at the scale $x$ for this window with this base, and weights for them, such that for every $u_0\in A_0$ and every choice of the nodes $w_{u_0,j}$ of Proposition~\ref{r6:R6:prop:nodes} there is a chain record with constants $C_\ast$, $C_H$ and exponents $e_D$, $e_L$. One can take $e_D=2$ and $e_L=0$.
\end{hypothesis}

Fix admissible data at the scale $x$, weights, $u_0\in A_0$ and the nodes $w_{u_0,j}$. Let $\mathcal W$, $T_u$, $N_u$, $M_u$, $h_u$, $R$, $E_{u,j}$, $\mathcal L$, $\Gamma$ and $\mathcal I$ be the objects of the construction, as described after Definition~\ref{r6:R6:def:record}, with $\mathcal I=\mathrm{IO}$ the index loss of Section~\ref{r2:sec:loss}.

\begin{proposition}\label{a8:prop:record}
Let $C_H\ge0$. Suppose that $x$ is so large that Lemma~\ref{r6:R6:lem:sizes} applies, and that the following hold for every $u\in\mathcal W$.
\begin{enumerate}
\item[(i)] The conclusions of Theorems~\ref{r2:A1} and~\ref{r2:A2}, with the constant $C_T$ of (L7), and $\mathrm{IO}\ge0$.
\item[(ii)] The hypotheses (a), (b), (d), (e) and (f) of Theorem~\ref{r3:A3}, with $\sigma=x^{-C_H}$, with numbers $\eta_0,\eta_1,e_0,\delta'$ that satisfy the conclusions of Lemma~\ref{r3:sizecond}, and with $\delta\le\frac12$.
\item[(iii)] $\mu_{\min}(R^\vee)=-\mu_{\max}(R)$ and $s+\mu_{\max}(R)\le x^{C_H}b$.
\item[(iv)] The conclusion of Theorem~\ref{r4:R4:A4} with a constant $C_0\ge1$.
\item[(v)] The conclusion of Theorem~\ref{r5:thm:Ebound} with the constant $C_E$, for every $1\le j\le m$.
\item[(vi)] $h_u$ is the number of Definition~\ref{r4:R4:hu}, and the numbers $e^{\max}_P$, $s_j$, $h_{j\rho}$, $e^{\max}_{\lambda,j}$ and $\langle\log\kappa\rangle$ in that definition do not depend on $u$.
\item[(vii)] The conclusion of Theorem~\ref{r4:R4:A6} with the constant $C_6$, and $\log T_u\le C_1mL_\varepsilon$ and $\lvert\log T_u-\log T_{u'}\rvert\le C_2\log k$ for $u'\in\mathcal W$.
\item[(viii)] The conclusion of Theorem~\ref{r4:R4:A7} with the constant $C_7$, and $\mu_{\max}(M_u)\le C_3N_Ta\Lambda_a$.
\item[(ix)] $1\le\lvert\mathcal W\rvert\le2^{k-3}$.
\end{enumerate}
Then the objects form a chain record with the constants $C_H$ and
\[
C_\ast=\max\big(6,\ C_T+2,\ 50C_0+10,\ C_E,\ C_6+2,\ C_1,\ C_2,\ C_3,\ C_7\big),
\]
and the exponents $e_D=1$ and $e_L=0$. Since $x^{\varepsilon}\le x^{2\varepsilon}$, they also form a chain record with $e_D=2$ and $e_L=0$.
\end{proposition}

\begin{proof}
We use $N_T\ge1$, $m\ge1$, $k\le a$, $\Lambda_a\ge L_\varepsilon\ge1$ and $B_0\ge0$ from Lemma~\ref{r6:R6:lem:sizes}, and $T_u\ge2$.

(a) is (ix).

(K1) Theorem~\ref{r2:A1} gives $S-\Gamma\le e_1\le(2\log4+1)N_Tmb$ and $0\le\mathrm{IO}\le N_Tmb$, and $2\log4+1<6\le C_\ast$.

(K2) Theorem~\ref{r2:A2}(ii) gives $\mu_{\max}(N_u)\le\mu_{\max}(M_u)+\Gamma+\mathcal L-E_u$, and $C_\ast mB_0\ge0$. Theorem~\ref{r2:A2}(i) gives $\mu(M_u)+\Gamma+\mathcal L-E_u-e_2\le\mu(N_u)$ with $e_2=4\,\mathrm{IO}+\mathrm{LT}+\mathrm{RX}+1$, where $\mathrm{LT}=4\lambda_T\log T_u=C_TN_Tm(\frac ak+b')\log T_u$ and $\mathrm{RX}\le6X''+2\log T_u$. Since $N_Tm(\frac ak+b')\ge1$ and $\log T_u\ge0$,
\[
e_2\le4\,\mathrm{IO}+(C_T+2)N_Tm\Big(\frac ak+b'\Big)\log T_u+6X''+1\le\mathrm{Lo}_u .
\]

(K3) Put $\eta=5\sigma+3\sqrt{\delta'}$ and $\epsilon=2e_0$. Then $\eta\ge0$ and $0\le\epsilon\le\frac14$ because $e_0\le\frac18$, and by Lemma~\ref{r3:sizecond}
\[
(\eta+4\sqrt\epsilon)\,x^{C_H}=5\sigma x^{C_H}+\big(3\sqrt{\delta'}+4\sqrt{2e_0}\big)x^{C_H}\le5+1=6 .
\]
Let $H>0$ with $H\ge s-\mu_{\min}(R^\vee)$. By (iii), $s+\mu_{\max}(R)\le H$, which is hypothesis (c) of Theorem~\ref{r3:A3}, and the other hypotheses do not involve $H$. Theorem~\ref{r3:A3} gives $s-e_3\le\mu(R^\vee)$ with $e_3=(1-\delta)^{-1}\big(4u_F\log(1+H/u_F)+(\eta+4\sqrt\epsilon)H\big)$, where $\Lambda(u,H)=u\log(1+H/u)$ as in Section~\ref{r3:sec}. Since $\log\lvert\mathcal W\rvert\ge0$, this gives (K3).

(K4) The first two statements are (iii). Theorem~\ref{r4:R4:A4} gives the inequality with $e_4=(50C_0+10)\big(N_Ta\Lambda_a+m\,x^\varepsilon(B_0+N_c)\big)$, which is at most $C_\ast N_Ta\Lambda_a+C_\ast m(B_0+N_c)x^{e_D\varepsilon}$ for $e_D=1$.

(K5) is (v), since $C_E\le C_\ast$.

(K6) By (vi),
\[
h_u+A\log T_u=e^{\max}_P+\sum_{j=1}^m\Big(s_j-\sum_\rho h_{j\rho}+e^{\max}_{\lambda,j}\Big)-\langle\log\kappa\rangle,
\]
with $A=\sum_{t\in T}a_t$, and the right side does not depend on $u$. So the left side of (K6) is $0$.

(K7) Theorem~\ref{r4:R4:A6} gives the first inequality with $e_6=C_6(L_\varepsilon+m^2L_\varepsilon/k)+2\le(C_6+2)(\Lambda_a+m^2\Lambda_a/k)$, because $L_\varepsilon\le\Lambda_a$ and $\Lambda_a\ge1$. The other two inequalities are (vii).

(K8) The first inequality is Theorem~\ref{r4:R4:A7}, with $C_7\le C_\ast$. The second is (viii). The third holds because $N_u$ is one of the sublattices over which $\mu_{\max}(N_u)$ is the maximum.
\end{proof}

The exponents $e_D=1$ and $e_L=0$ do not depend on $\psi_0$, $\varepsilon$ or $J$, so they are absolute, and Proposition~\ref{r6:R6:prop:errors} gives $\kappa=\max(2,e_D,e_L+1)=2$ and the constant $C=3$ of Theorem~A, whether one takes $e_D=1$ or $e_D=2$. The constant $C_\ast$ depends on $\varepsilon$ through $C_E$, which Hypothesis~\ref{r6:R6:thm:records} allows.

Condition (ix) is compatible with the lower bound $\lvert\mathcal W\rvert\ge C_0\Gamma_0^2x^{C_H}$ of Lemma~\ref{r3:sizecond}, with $\Gamma_0=\exp(C_\Gamma m^5fL_\varepsilon)$. Indeed $m\le2x^{1/2-\varepsilon/16}$ and $f\le4x^{2+8\varepsilon}$ give $\log\big(C_0\Gamma_0^2x^{C_H}\big)\le\log C_0+256\,C_\Gamma\,x^{9/2+139\varepsilon/16}+C_H\log x$, while $(k-3)\log2\ge\frac12x^{9/2+10\varepsilon}$ for large $x$, and $\frac{139}{16}<10$.



Hypothesis~\ref{r6:R6:thm:records} collects in one statement the inequalities proved in Sections~\ref{sec:lifts} to~\ref{sec:twisted}. Condition (a) is the choice of the words in Lemma~\ref{r3:sizecond}. Conditions (K1) and (K2) are Theorems~\ref{r2:A1} and~\ref{r2:A2} with the divisibility of Theorem~\ref{r1:thm:R1main}. Condition (K3) is Theorem~\ref{r3:A3} with Theorems~\ref{w2:S}, \ref{w2:I} and~\ref{w2:B}. Condition (K4) is Theorem~\ref{r4:R4:A4} with Theorem~\ref{r5:thm:DIVinN}. Condition (K5) is Theorem~\ref{r5:thm:Ebound}. Conditions (K6) and (K7) are proved in Section~\ref{r4:sec:R4}, and (K8) is Theorem~\ref{r4:R4:A7} with $\mu(N_u)\le\mu_{\max}(N_u)$. Each of these results is proved here from hypotheses stated in its section. The conditions (K1) to (K8) are stated with the exponents and constants of the formal proof, and some of the cited results give an inequality of a slightly different shape. Theorem~\ref{r3:A3} gives (K3) without the term $\log\lvert\mathcal W\rvert$, which only strengthens it, and Theorem~\ref{r2:A2} gives the two halves of (K2), the upper half without the term $C_\ast mB_0$. \notgiven{We have not checked line by line that the constants and exponents of Sections~\ref{sec:lifts} to~\ref{sec:twisted} give (K1) to (K8) with absolute $e_D$ and $e_L$, so Hypothesis~\ref{r6:R6:thm:records} is not proved here even modulo the hypotheses of those sections.}

The proof of Theorem~A below uses Hypothesis~\ref{r6:R6:thm:records} for one admissible datum at each scale. The conditions (K1) to (K8) depend on the window and the base through the field $E$, the ranks and the sets $T$ and $T_2$.

\subsection{The seven hypotheses of the accounting}\label{r6:R6:sec:A1A7}

We fix a chain record with constants $C_\ast$, $C_H$ and exponents $e_D$, $e_L$, and we write $y=x^{\varepsilon}=L_\varepsilon$. The first lemma bounds the logarithmic term of (K3).

\begin{lemma}\label{r6:R6:lem:phi}
Let $u\ge0$, $U\ge u$, $b>0$, $X>0$ and $0<H\le Xb$. Then $u\log(1+H/u)\le\log(1+X)\,(U+b)$.
\end{lemma}

\begin{proof}
For $u=0$ the left side is $0$ and the right side is nonnegative. Let $u>0$. The inequality $\log z\le z-1$ for $z>0$, with $z=(1+H/u)/(1+X)$, gives
\[
u\log\Big(1+\frac Hu\Big)\le u\log(1+X)+\frac{u+H}{1+X}-u=u\log(1+X)+\frac H{1+X}-\frac{uX}{1+X}.
\]
The last term is nonpositive, $u\log(1+X)\le U\log(1+X)$, and
\[
\frac H{1+X}\le\frac{Xb}{1+X}=b\Big(1-\frac1{1+X}\Big)\le b\log(1+X),
\]
by the inequality $1-1/z\le\log z$ for $z=1+X$.
\end{proof}

The second lemma bounds the spread $u_F$ of the slopes of the lattices $N_u$. Put
\begin{align*}
K_F&=C_\ast N_T\,a\,\Lambda_a\Big(3+\frac mk\Big)+4C_\ast mB_0+C_\ast A\log k,\\
e_2&=C_\ast^2N_T\,m\,b\,x^{e_L\varepsilon}+C_\ast^2N_T\,m^2\Big(\frac ak+b'\Big)L_\varepsilon+C_\ast X''+1,\\
U&=K_F+e_2 .
\end{align*}

\begin{lemma}\label{r6:R6:lem:uF}
For every $u\in\mathcal W$ we have $\mathrm{Lo}_u\le e_2$, and $0\le u_F\le U$.
\end{lemma}

\begin{proof}
By (K1), $\mathcal I\le C_\ast N_T\,m\,b\,x^{e_L\varepsilon}$, and by (K7), $\log T_u\le C_\ast mL_\varepsilon$. The coefficient $C_\ast N_T\,m(a/k+b')$ is nonnegative, so $\mathrm{Lo}_u\le e_2$.

By (K8), $\mu(N_u)-h_u\le\mu_{\max}(N_u)-h_u$ for every $u$, so the weighted mean $s$ is at most the maximum in the definition of $u_F$, and $u_F\ge0$.

Let $u,u'\in\mathcal W$. The first inequality of (K2) for $u$ and the second for $u'$ give
\begin{multline*}
\big(\mu_{\max}(N_u)-h_u\big)-\big(\mu(N_{u'})-h_{u'}\big)\\
\le\mu_{\max}(M_u)-\mu(M_{u'})+C_\ast mB_0+\mathrm{Lo}_{u'}+(E_{u'}-E_u)+(h_{u'}-h_u).
\end{multline*}
By (K8), $\mu_{\max}(M_u)-\mu(M_{u'})\le C_\ast N_T\,a\,\Lambda_a(2+m/k)$. By (K5), each $E_{u,j}$ lies within $C_\ast B_0$ of $Ag$, so $E_{u'}-E_u\le2C_\ast mB_0$. By (K6) and the third inequality of (K7),
\[
h_{u'}-h_u\le C_\ast\big(N_T\,a\,\Lambda_a+mB_0\big)+A\big(\log T_u-\log T_{u'}\big)\le C_\ast\big(N_T\,a\,\Lambda_a+mB_0\big)+C_\ast A\log k ,
\]
where we used $A\ge0$. Adding these bounds, and using $\mathrm{Lo}_{u'}\le e_2$,
\[
\big(\mu_{\max}(N_u)-h_u\big)-\big(\mu(N_{u'})-h_{u'}\big)\le K_F+e_2=U .
\]
The weighted mean over $u'$ with the weights $d_{u'}$ gives $\mu_{\max}(N_u)-h_u-s\le U$, and the maximum over $u$ gives $u_F\le U$.
\end{proof}

Define
\begin{gather*}
e_1=C_\ast N_T\,m\,b\,x^{e_L\varepsilon},\qquad e_3=\log\lvert\mathcal W\rvert+8\log\big(1+x^{C_H}\big)(U+b)+12\,b,\\
e_4=C_\ast N_T\,a\,\Lambda_a+C_\ast m\,(B_0+N_c)\,x^{e_D\varepsilon},\qquad e_5=C_\ast mB_0,\\
e_6=C_\ast\Lambda_a\Big(1+\frac{m^2}k\Big),\qquad e_7=C_\ast N_T\,a\,\Lambda_a\Big(1+\frac mk\Big),
\end{gather*}
and let $e_2$ be as above. All seven numbers are nonnegative, because $C_\ast\ge1$, $\log\lvert\mathcal W\rvert\ge0$ and $\Lambda_a\ge\log a\ge0$.

\begin{proposition}\label{r6:R6:prop:A1A7}
Let a chain record be given for admissible data at the scale $x$, weights and $u_0\in A_0$. Put $d_u=\rk N_u$, $\mu^N_u=\mu(N_u)$, $\mu^M_u=\mu(M_u)$, $\mu^C=\mu(R^\vee)$, $w_{j\rho}=w_{u_0,j}(\rho)$ for $\rho\in\Phi$, $v=\frac12\operatorname{rd}_E$, $\gamma$ Euler's constant, $\lambda_1=\ell_{\chi_1}$ and $g=\log(g_0g_1)$. Then the hypotheses \eqref{A1} to \eqref{A7} of Proposition~\ref{acct} hold for every $u\in\mathcal W$, with the index set $\mathcal W$ in place of $W$, with $S$, $\Gamma$, $\mathcal L$, $E_u$, $h_u$, $T_u$ and the numbers $e_1,\dots,e_7$ defined above.
\end{proposition}

\begin{proof}
The weights $a_t\ge a>0$, the nodes are positive by Corollary~\ref{r6:R6:cor:nodes}, $v>0$, $d_u\ge1$ and $T_u\ge2$, so the data are of the kind that Proposition~\ref{acct} admits.

\eqref{A1} is the first inequality of (K1). \eqref{A2} is the second inequality of (K2) with $\mathrm{Lo}_u\le e_2$ from Lemma~\ref{r6:R6:lem:uF}.

For \eqref{A3}, put $H=x^{C_H}b$. By (K4), $s-\mu_{\min}(R^\vee)=s+\mu_{\max}(R)\le H$, and $H>0$, so (K3) applies with this $H$. Lemma~\ref{r6:R6:lem:phi} with $u=u_F$, $U$ from Lemma~\ref{r6:R6:lem:uF} and $X=x^{C_H}$ gives $4u_F\log(1+H/u_F)\le4\log(1+x^{C_H})(U+b)$. By (K3), $(\eta+4\sqrt\epsilon)H\le6b$ and $(1-\delta)^{-1}\le2$. Both terms in the bracket of (K3) are nonnegative, so (K3) gives
\[
s-\log\lvert\mathcal W\rvert-8\log\big(1+x^{C_H}\big)(U+b)-12b\le\mu(R^\vee),
\]
which is \eqref{A3} with $e_3$, because $s$ is the weighted mean of $\mu^N_u-h_u$.

\eqref{A4} is the inequality of (K4) with $e_4$. For \eqref{A5}, summing (K5) over $j$ gives $E_u\le mAg+C_\ast mB_0=mAg+e_5$. \eqref{A6} is the first inequality of (K7) with $e_6$, and \eqref{A7} is the first inequality of (K8) with $e_7$.
\end{proof}

\subsection{The size of the errors}\label{r6:R6:sec:errors}

\begin{lemma}\label{r6:R6:lem:sizes}
Let $0<\varepsilon<\frac1{52}$, $J\ge2$ and $C_H\ge0$. There is $x_0$, depending on $\psi_0$, $\varepsilon$, $J$ and $C_H$, such that for every $x\ge x_0$, all admissible data at the scale $x$ with this $J$, all weights and every chain record with $\lvert\mathcal W\rvert\le2^{k-3}$ the following hold, with $y=x^{\varepsilon}$.
\begin{enumerate}
\item[(i)] $1\le b\le a$, $1\le m\le x$, $mb\le a$, $xm\le k\le x^5$ and $k\le a$.
\item[(ii)] $kb'\le a$ and $X''\le a$.
\item[(iii)] $0\le\Lambda_a\le2y$.
\item[(iv)] $\log x\le y$, $\log k\le5y$ and $\log(1+x^{C_H})\le(1+C_H)\,y$.
\item[(v)] $N_T\ge1$ and $0\le A\le2N_Ta$.
\item[(vi)] $\log\lvert\mathcal W\rvert\le k$, and $\max(a,\,k,\,b,\,N_Ta,\,mB_0,\,mN_c)\le\mathcal B$.
\end{enumerate}
\end{lemma}

\begin{proof}
(i) We have $b=e^{y}\ge1$ and $a=xb\ge b$. Proposition~\ref{ranks} gives $2\le m\le2\sqrt{x\delta_1}\le x$, so $mb\le xb=a$, and $xm\le k\le2x^{9/2+10\varepsilon}\le x^5$ because $10\varepsilon<\frac15$. For large $x$, $x^5\le xe^{y}=a$.

(ii) Since $b'=a/X'$, the inequality $kb'\le a$ is $k\le X'=\exp(x^{\varepsilon/3})$, which holds for large $x$ by (i). The inequality $X''\le a$ is $nX'\le b$, that is $n\le\exp(x^{\varepsilon}-x^{\varepsilon/3})$, which holds for large $x$ because $n\le2x^{2+8\varepsilon}$.

(iii) The integer $\lvert\operatorname{disc}E\rvert$ is at least $1$, so $\Lambda_a\ge\log a\ge0$. The field $E$ lies in $\mathbb Q(\zeta_N)$ with $N=N_1qd$ and $N_1=\operatorname{lcm}(2^{e+2},M_\psi)$. By the conductor--discriminant formula, $\lvert\operatorname{disc}E\rvert$ is the product of the conductors of the $[E:\mathbb Q]$ Dirichlet characters of $E$, and each of them divides $N$. Hence $\operatorname{rd}_E\le N$ and
\[
\log\operatorname{rd}_E\le\log\big(2^{e+2}M_\psi\big)+\log q+\log d\le\log(8M_\psi n)+Y_0^{\varepsilon/5}\big(1+\log^2Y_0\big),
\]
by (D1$'$) and $2^e\le n$. Since $Y_0\le x$ and $n\le2x^{2+8\varepsilon}$, the right side is at most $x^{\varepsilon/5}(1+\log^2x)+3\log x+\log(16M_\psi)$. Then $\Lambda_a=\log x+y+\log\operatorname{rd}_E\le2y$ for large $x$.

(iv) The first inequality holds for large $x$. The second is $\log k\le5\log x$, from $k\le x^5$. For the third, $\log(1+x^{C_H})\le\log2+C_H\log x\le(1+C_H)y$ because $y\ge1$.

(v) Proposition~\ref{tops} gives $N_T\ge\frac12\delta_1n\ge1$, and $0<a_t\le a+a/(kx^{13})\le2a$.

(vi) $\log\lvert\mathcal W\rvert\le(k-3)\log2\le k$. The first term of \eqref{r6:R6:calB} is at least $N_Ta\ge a$, and (i) gives $b\le a$ and $k\le a$. The second term of \eqref{r6:R6:calB} is $mB_0+mN_c$.
\end{proof}

\begin{proposition}\label{r6:R6:prop:errors}
Let $\kappa=\max(2,e_D,e_L+1)$. Let $0<\varepsilon<\frac1{52}$, $J\ge2$, $C_\ast\ge1$ and $C_H\ge0$. There is $x_0$ such that for every $x\ge x_0$, all admissible data at the scale $x$ with this $J$, all weights and every chain record with constants $C_\ast$, $C_H$ and exponents $e_D$, $e_L$, the numbers $e_1,\dots,e_7$ of Section~\ref{r6:R6:sec:A1A7} satisfy
\[
e_1+e_2+e_3+e_4+e_5+Ae_6+e_7\le x^{(\kappa+1)\varepsilon}\,\mathcal B .
\]
\end{proposition}

\begin{proof}
Let $x_0$ be as in Lemma~\ref{r6:R6:lem:sizes}, enlarged so that $y=x^\varepsilon\ge260(1+C_H)C_\ast^2$ for $x\ge x_0$. We bound each term by a multiple of $\mathcal B$, using Lemma~\ref{r6:R6:lem:sizes}, $y\ge1$, $C_\ast\ge1$ and $y^\alpha\le y^\kappa$ for $0\le\alpha\le\kappa$.

Since $mb\le a$, $e_1\le C_\ast y^{e_L}N_Ta\le C_\ast y^{e_L}\mathcal B$. In $e_2$, the first term is at most $C_\ast^2y^{e_L}\mathcal B$. Since $b'\le a/k$, the second term is at most $2C_\ast^2y\,N_Ta\,m^2/k\le2C_\ast^2y\,\mathcal B$. The third term is at most $C_\ast a\le C_\ast\mathcal B$, and $1\le\mathcal B$. Hence
\[
e_2\le5C_\ast^2y^{\max(1,e_L)}\mathcal B .
\]
Since $\Lambda_a\le2y$, $A\le2N_Ta$ and $m\le k$,
\[
e_4\le2C_\ast y\,N_Ta+C_\ast y^{e_D}m(B_0+N_c)\le3C_\ast y^{\max(1,e_D)}\mathcal B,\qquad e_5\le C_\ast\mathcal B,
\]
\[
Ae_6\le2N_Ta\cdot C_\ast\cdot2y\Big(1+\frac{m^2}k\Big)\le4C_\ast y\,\mathcal B,\qquad e_7\le C_\ast N_Ta\cdot2y\cdot2\le4C_\ast y\,\mathcal B .
\]
For $e_3$, $\log k\le5y$ and $A\le2N_Ta$ give
\[
K_F\le8C_\ast y\,N_Ta+4C_\ast\mathcal B+10C_\ast y\,N_Ta\le22C_\ast y\,\mathcal B,
\]
so $U+b\le22C_\ast y\,\mathcal B+5C_\ast^2y^{\max(1,e_L)}\mathcal B+\mathcal B\le28C_\ast^2y^{\max(1,e_L)}\mathcal B$. With $\log(1+x^{C_H})\le(1+C_H)y$, $\log\lvert\mathcal W\rvert\le\mathcal B$ and $b\le\mathcal B$,
\[
e_3\le\mathcal B+8(1+C_H)y\cdot28C_\ast^2y^{\max(1,e_L)}\mathcal B+12\mathcal B\le237(1+C_H)C_\ast^2y^{\kappa}\mathcal B,
\]
because $1+\max(1,e_L)\le\kappa$. Adding the seven bounds,
\[
e_1+e_2+e_3+e_4+e_5+Ae_6+e_7\le\big(18C_\ast^2+237(1+C_H)C_\ast^2\big)y^\kappa\mathcal B\le260(1+C_H)C_\ast^2y^\kappa\mathcal B\le y^{\kappa+1}\mathcal B .\qedhere
\]
\end{proof}

With the exponents $e_D=2$ and $e_L=0$ of Hypothesis~\ref{r6:R6:thm:records}, $\kappa=2$ and the bound is $x^{3\varepsilon}\mathcal B$.

\subsection{Proof of Theorem A}\label{r6:R6:sec:thmA}

Assume Hypotheses~\ref{w4:thm:vwin} and~\ref{r6:R6:thm:records}. We prove Theorem~A in the form stated in Section~\ref{sec:reduction}, which asks for one choice of admissible data at each scale of a window, for windows $[Y,Y^{3/2}]$ with $Y$ arbitrarily large.

\begin{proof}
Let $e_D$ and $e_L$ be the absolute constants of Hypothesis~\ref{r6:R6:thm:records}, and put $C=\max(2,e_D,e_L+1)+1$, which is $3$ for $e_D=2$ and $e_L=0$. Fix $\psi_0$, let $\varepsilon_1$ be as in Hypothesis~\ref{r6:R6:thm:records}, and fix $0<\varepsilon<\varepsilon_1$. Let $J\ge2$ be the number of Proposition~\ref{r6:R6:prop:dataexist}, and let $C_\ast$, $C_H$ and $x_1$ be the numbers of Hypothesis~\ref{r6:R6:thm:records} for $\varepsilon$ and this $J$. Let $x_0$ be the maximum of $x_1$ and of the numbers $x_0$ of Lemma~\ref{r6:R6:lem:tr}, Corollary~\ref{r6:R6:cor:nodes}, Lemma~\ref{r6:R6:lem:sizes} and Proposition~\ref{r6:R6:prop:errors} for these constants.

Let $y$ be given. Proposition~\ref{r6:R6:prop:dataexist} gives $Y_0\le Y$ with $Y\ge\max(y,x_0)$ such that admissible data with base $Y_0$ exist for every $x\in[Y,Y^{3/2}]$. Fix such an $x$. Then $x\ge x_1$, and Hypothesis~\ref{r6:R6:thm:records} gives admissible data at the scale $x$, for the window $[Y,Y^{3/2}]$ with base $Y_0$, and weights $a_t,b_t$, such that for every $u_0\in A_0$ the nodes $w_{u_0,j}$ of Proposition~\ref{r6:R6:prop:nodes} admit a chain record with constants $C_\ast$, $C_H$ and exponents $e_D$, $e_L$.

The weights satisfy $a\le a_t\le a+a/(kx^{13})$ by Definition~\ref{r6:R6:def:weights}, and Corollary~\ref{r6:R6:cor:nodes} gives (i).

Fix $u_0\in A_0$ and a chain record for it. Proposition~\ref{r6:R6:prop:A1A7} shows that the hypotheses \eqref{A1} to \eqref{A7} of Proposition~\ref{acct} hold with $\lambda_1=\ell_{\chi_1}$, with $w_{j\rho}=w_{u_0,j}(\rho)$ for $\rho\in\Phi$, and with $S$ the sum over $t\in T$ and $1\le j\le m$ of the left side of \eqref{excessid} with $a=a_t$ and $w=w_{u_0,j}$. Proposition~\ref{acct} gives the conclusion \eqref{acctconc}. The number $\ell_{\chi_1}=L'(1,\chi_{-d})/L(1,\chi_{-d})$ is real, so Corollary~\ref{acctexcess} applies and gives
\[
\sum_{t\in T}\sum_{j=1}^m\operatorname{Re}\mathrm{De}(a_t,t,w_{u_0,j})\le e_1+e_2+e_3+e_4+e_5+Ae_6+e_7 .
\]
Proposition~\ref{r6:R6:prop:errors} bounds the right side by $x^{C\varepsilon}\mathcal B$, which is (ii).
\end{proof}

The constant $C$ does not depend on $\varepsilon$ or on $\psi_0$, which is the uniformity that the last step of the proof of Theorem~\ref{mainthm} uses. With $C=3$ that proof gives $\zeta(s)\ne0$ and $L(s,\chi^{(i)})\ne0$ for $\operatorname{Re}s>\frac12+7\varepsilon$, for every $0<\varepsilon<\varepsilon_1$, and so for $\operatorname{Re}s>\frac12$.

The three primes of (D5) are congruent to $1$ modulo $15\,Q_Sqd$. The construction of the level structure of the lifts uses two primes congruent to $1$ modulo $\operatorname{lcm}(15,Q_Sqd)$ and a third prime, different from them, congruent to $1$ modulo $Q_Sqd$, all at most $(Q_Sqd)^{C_8}$ for a constant $C_8$. Since $\operatorname{lcm}(15,Q_Sqd)$ divides $15\,Q_Sqd$, the primes of (D5) have these properties with $C_8=2A_L(1+A_L)^2$, because $15\le Q_Sqd$ gives $(15\,Q_Sqd)^{A_L(1+A_L)^2}\le(Q_Sqd)^{2A_L(1+A_L)^2}$.

\subsection{Hypotheses in the final deduction}\label{sec:frontier}

\paragraph{The chain and the Hecke bounds.}\label{fr:sec:chain}\label{fr:sec:loss}
Proposition~\ref{a8:prop:N1} gives the frame valuation under the step-character
and determinant normalizations of Definition~\ref{a8:def:frame}.
Hypotheses~\ref{r1:hyp:H2} and~\ref{r1:hyp:H3} concern the odd-position
lattices and the normalized kernel. In the Hecke bound, Theorem~\ref{w1:thm:B7}
controls the isotypic projector, Proposition~\ref{w1:prop:C7} bounds the
window loss, and Lemma~\ref{w1:lem:A9} combines the final share estimates.
Hypothesis~\ref{r2:engines} still requires the individual share estimates,
the comparisons between levels and the construction of the test systems.
Corollary~\ref{r2:cor:L2} uses the Weyl lattice model. Proposition~\ref{r2:prop:Z1}
requires a polynomial bound for the constant of the base weights.
The colength at the auxiliary prime in (L7) is a further input.

\paragraph{Periods and multiplicities.}\label{fr:sec:frame}
Theorem~\ref{w2:Bproof} uses the seesaw relation, its nonzero pairing
and the period identifications. Theorem~\ref{w3s:thm:T2} proves the
double-step identity from the general and local Weil representation
laws. Lemma~\ref{w2:see:rat} specifies the three period identities
needed for rationality. The multiplicity calculation requires the
label identifications and the pre-trace hypotheses of
Subsections~\ref{w2x:sec:rows} and~\ref{w2x:sec:count}.
The estimates on individual rows in Theorem~\ref{w2:Sproof} also
require the kernel expansion, the main row, the loop rows and the
degenerate terms. Lemma~\ref{w39:sizes} proves the two exponential
size inequalities with $r_S=m/2-1$.

\paragraph{Heights and local terms.}\label{fr:sec:heights}\label{fr:sec:local}
The kernel norm uses the common integral Weber value and the period
identities in Section~\ref{r4:sec:R4}. The final pairing requires its
factorization through the branching map. For the covolume bound,
Theorem~\ref{cov:main} specifies the descent, flat twist and lattice
comparisons under its hypotheses (I1)--(I17).
Hypothesis~\ref{w3:W3:unif} gives a lower bound for the average of
component slopes with the term $q\,\operatorname{spread}(u)$ in the
error. The flat twist gives the main term in the covolume inequality.
The exact change of integral structure at good primes still enters the
uniformity in the level. At the auxiliary and inert places the degree
and orbit bounds enter Corollary~\ref{r5:cor:A5}.

\paragraph{Twisted places and assembly.}\label{fr:sec:twisted}\label{fr:sec:assembly}
Theorem~\ref{w4:thm:decomp} separates the plus and minus parts, including
the prime $2$. Proposition~\ref{w4:prop:pluschar} separates the character
summands when the group order is a unit. The generic estimate uses
(OneChar) and (Count). The two rank-four cases also require propagation
of automorphy and rigidity as in Subsection~\ref{tw:rankfour}.
The rank-two local comparison in Subsection~\ref{tw:levelpoints} applies
to Hilbert modular forms in the stated normalization.
For the twisted lattice, Propositions~\ref{lk:prop:ident} and~\ref{lk:prop:gd}
identify the congruence module and reduce its estimate to per-component
bounds. The Wach lifting estimate, the bounds for the other local rows,
and the integral theta-series comparisons still enter that argument.
Proposition~\ref{lk:prop:d0} shows why $\dim M_j=D_0$ suffices at the
last step. Theorem~\ref{r5:thm:DIVinN} uses only $d^0_j\le D_0$.

\paragraph{Published inputs.}\label{fr:sec:audit}
Section~\ref{sec:auxiliary} uses Theorem~1.4 of Thorner and Zaman
in a cyclotomic field for (C7) and in a ray class field, together with
Stark's bound, for (C9). For the local arguments Berger's de Rham theorem is
Th\'eor\`eme~0.6, equivalently Th\'eor\`eme~5.19, of~\cite{berger02}.
For algebraic Hecke characters the Hodge jumps at $\tau$ have sign
$-n_\tau$ in the local Artin convention. In the CM functional equation
the archimedean phase at infinity type $2m$ is
$\mathrm{i}^{\,\sum_i|2m_i|}$: substituting $2m$ in the local phase gives
this equality since the exponent is even.
The Rallis inner product identity still requires the comparison of the
Petersson norm with the operators of Proposition~\ref{w3:W3:HLprop}.
The mass and formal-degree formulas must use the measures of the
dimension estimate. The final implication to Theorem~\ref{thm:intro}
remains conditional on the hypotheses above.


\section*{Acknowledgments}
I thank Timothy Wesley Stone for raising the question of rationality for the period maps and for pointing me to the Lean declarations relevant to that comparison.
I thank \texttt{erythrofr} for proposing a sharper bound for the integrated deficit.
I thank \texttt{a-ratz} for identifying the original commit to cite, questioning my attribution of V17, and pointing out discrepancies between the declaration names I cited and those in the later community repository.

\bibliographystyle{plainnat}
\bibliography{references}

\end{document}
